% !TeX program = lualatex
% Fichier LaTEI généré depuis XML-TEI Métopes / Commons-Publishing.
% Zone technique (préambule, macros, mappings) : régénérable depuis le XML source.
% NE PAS MODIFIER la zone technique — elle sera écrasée à la prochaine génération.
% Zone éditoriale réversible : l'environnement lateiDocument dans ce fichier.
% Corriger uniquement dans la zone réversible.

\documentclass[11pt,twoside,openany]{book}

\usepackage[
  paperwidth=155mm,
  paperheight=230mm,
  top=30mm,
  bottom=19mm,
  inner=20mm,
  outer=30mm,
  headheight=14pt,
  headsep=8mm,
  footskip=10mm
]{geometry}
% Référentiel PURH v0.6 §6.2 : « le bandeau généré est situé environ 3,1 mm
% plus bas que dans le PDF imprimeur » — écart de rendu, pas une remise en
% cause de la marge du haut mesurée sur le maître InDesign (profile.
% margin_top_mm, qui reste la source de vérité pour la position du corps de
% texte). \topmargin remonte le bandeau de titre courant de 3,1 mm ;
% \headsep compense d'autant pour que le corps de texte, lui, démarre
% exactement à la même position qu'avant ce correctif.
\addtolength{\topmargin}{-3.1mm}
\addtolength{\headsep}{3.1mm}
\raggedbottom

\usepackage{fontspec}
\usepackage[french]{babel}
\usepackage{csquotes}
\usepackage{amsmath}
\usepackage{microtype}
\usepackage{indentfirst}
\usepackage{emptypage}
\usepackage[normalem]{ulem}
% [table] : nécessaire pour \rowcolor sur les lignes d'entête de tableau
% (référentiel §11.3, "fond foncé noir 30 %") — charge colortbl en plus des
% commandes \color habituelles (compatible avec l'usage \color[gray]{}
% déjà en place pour le titre courant).
\usepackage[table]{xcolor}

% Référentiel PURH v0.6 §12.1 : « le texte courant du PDF imprimeur
% correspond à un noir process à 90 % [K]. La sortie actuelle utilise du
% noir plein. » CMYK explicite (0,0,0,0.9), pas une valeur de gris RVB :
% "K" désigne le noir process de la quadrichromie, un concept distinct du
% gris \color[gray]{} utilisé ailleurs (titre courant) qui n'a pas de sens
% en CMJN. Appliqué globalement dès le début du document.
\definecolor{PURHBodyBlack}{cmyk}{0,0,0,0.9}
\AtBeginDocument{\color{PURHBodyBlack}}

\IfFontExistsTF{Chaparral Pro}
  {\setmainfont{Chaparral Pro}}
  {\setmainfont{TeX Gyre Pagella}}

\IfFontExistsTF{Josefin Sans}
  {\newfontfamily\PURHTitleFont{Josefin Sans}}
  {\newfontfamily\PURHTitleFont{TeX Gyre Heros}}

% Titraille PURH (référentiel §2.3-§2.5) : partie, article et section
% observés en Josefin Sans Thin, capitales, distinct du \PURHTitleFont
% "normal weight" ci-dessus utilisé ailleurs. Chargée comme famille à part
% (et non via \bfseries sur \PURHTitleFont) car "Thin" n'est pas une série
% NFSS standard que fontspec puisse sélectionner automatiquement — "Josefin
% Sans Thin" existe en revanche comme nom de famille indépendant portant
% elle-même ses propres graisses Thin (romain) et Thin Italic.
\IfFontExistsTF{Josefin Sans Thin}
  {\newfontfamily\PURHTitreFont{Josefin Sans Thin}}
  {\newfontfamily\PURHTitreFont{TeX Gyre Heros}}

\IfFontExistsTF{Latin Modern Mono}
  {\setmonofont{Latin Modern Mono}}
  {\setmonofont{TeX Gyre Cursor}}

% Romain, pas italique (référentiel PURH §2.3 : titres courants "Josefin
% Sans Thin/Light, 10 pt, romain" — l'italique systématique précédente était
% un défaut confirmé, pas un choix).
%
% Cinq vérifications humaines successives (2026-08-04/06) ont porté sur ce
% même réglage : la première jugeait le gris trop clair (corrigé en passant
% de \PURHTitreFont, famille Thin, à \PURHTitleFont, famille standard, sans
% \bfseries) ; la seconde a trouvé ce résultat trop noir et visuellement
% gras — le PDF imprimeur, lui, n'a « pas de graisse » — d'où un retour à la
% famille Thin avec un gris explicite (\color[gray]{0.25}, approximatif) ;
% la troisième a jugé ce gris encore trop clair, d'où un premier passage au
% même système que le fond d'entête de tableau et le texte courant
% (§11.3/§12.1) — une teinte CMJN noir X % plutôt qu'un gris RVB — à 50 %
% noir ; la quatrième a de nouveau jugé ce résultat trop clair au regard du
% PDF imprimeur, où le titre courant est « presque noir » : remonté à 85 %
% noir. La cinquième vérification, sur une page réelle complète d'un livre
% du corpus (pas seulement les liminaires), a de nouveau jugé le résultat
% trop clair à 85 % — poussé au maximum possible, 100 % noir (K plein),
% sans changer de famille (\PURHTitreFont, Josefin Sans Thin) : le résultat
% restait visiblement plus clair que le PDF imprimeur, confirmant qu'à ce
% stade le levier couleur seul était épuisé — un trait Thin, même à pleine
% saturation d'encre, couvre trop peu de surface pour lire aussi sombre
% qu'un trait plus épais à la même teinte.
%
% Sixième vérification (2026-08-06), comparaison directe côte à côte avec
% le PDF imprimeur (recadrage du bandeau à haute résolution) : la graisse
% du PDF imprimeur n'est PAS aussi fine qu'un trait Thin — plutôt un trait
% "Regular" (ni fin ni gras). Changement de levier plutôt que nouvel
% incrément de couleur : \PURHTitleFont (famille Josefin Sans complète,
% graisse Regular par défaut sans \bfseries — PAS \PURHTitreFont/Thin) à
% 75 % noir, un compromis délibéré entre les deux réglages précédents (ni
% le gris à 50 % jugé trop clair, ni l'aspect "gras" que la graisse
% Regular pourrait donner à 100 %) — comparé visuellement à un jeu
% d'échantillons Thin/Light/Regular/Medium × 75 %/100 % avant d'être
% retenu comme la combinaison la plus proche du PDF imprimeur.
%
% Septième vérification (2026-08-06) : le corps (Regular) validé comme
% correct par l'utilisateur — seule la teinte, encore un peu trop noire à
% 75 %, redescendue légèrement à 65 % (un ajustement fin, pas un nouveau
% changement de levier comme aux passes précédentes).
\newcommand{\PURHHeaderFont}{\PURHTitleFont\small\color[cmyk]{0,0,0,0.65}}

% Le corps et son pas de ligne sont fixés explicitement au lieu de dépendre
% de la table de tailles du \documentclass{book} choisi : elle donne un
% pas voisin mais pas identique au pas de grille cible (référentiel PURH
% §2.4, §5.3 : corps 11 pt sur une grille de 13,5 pt).
\renewcommand{\normalsize}{\fontsize{11pt}{13.5pt}\selectfont}
\normalsize

\newcommand{\PURHBookTitle}{Beautés vitales}
\newcommand{\PURHBookSubtitle}{Pour une approche contemporaine de la beauté}
\newcommand{\PURHBookAuthor}{Anne-Lise Worms ; Clélia Zernik}
\newcommand{\PURHPublisher}{Presses universitaires de Rouen et du Havre}
\newcommand{\PURHYear}{}
\newcommand{\PURHDOI}{}
\newcommand{\PURHISBN}{979-10-240-1926-0}

\title{\PURHBookTitle}
\author{\PURHBookAuthor}
\date{\PURHYear}

\setlength{\parindent}{5mm}
\setlength{\parskip}{0pt}
\linespread{1.0}
\pretolerance=100
\tolerance=500
\hyphenpenalty=500
\exhyphenpenalty=500
\emergencystretch=3em
\clubpenalty=10000
\widowpenalty=10000
\displaywidowpenalty=10000

\usepackage{enumitem}

\setlist[itemize,1]{
  label=\textendash,
  leftmargin=1.5em,
  itemsep=0.2\baselineskip,
  topsep=0.4\baselineskip
}

\setlist[enumerate,1]{
  leftmargin=1.8em,
  itemsep=0.2\baselineskip,
  topsep=0.4\baselineskip
}

\usepackage[nobottomtitles*]{titlesec}
\usepackage{titletoc}

\setcounter{secnumdepth}{0}
% Référentiel PURH v0.6 §9 ("Table des matières") : la cible exclut les
% sections internes (intertitres) de la TDM — seuls le niveau des parties
% (\part, niveau -1) et celui des ouvertures de contribution/front matter
% (\addcontentsline{toc}{chapter}, niveau 0) doivent y figurer. tocdepth=2
% incluait à tort les <div type="section1">/<div type="section2"> (niveaux
% \section=1, \subsection=2), gonflant la TDM à trois pages au lieu de deux.
\setcounter{tocdepth}{0}

% Josefin Sans Bold, 16 pt, capitales — le référentiel indiquait Josefin
% Sans Thin (§2.5, §5.3, §4.3), contredit par vérification humaine directe
% du PDF généré face au PDF imprimeur : les titres de partie y sont noirs et
% gras, alors que le Thin rendait un texte maigre et grisâtre, peu lisible.
% Observation directe suivie ici, comme pour le séparateur de note de bas de
% page défini plus bas dans ce même préambule. \PURHTitleFont (pas
% \PURHTitreFont, qui ne charge que la graisse Thin et ne peut donc pas
% produire de \bfseries réel) charge la famille Josefin Sans complète, dont
% fontspec sélectionne automatiquement la graisse Bold via NFSS.
\titleformat{\chapter}[display]
  {\PURHTitleFont\bfseries\fontsize{16pt}{19pt}\selectfont\raggedright}
  {\chaptertitlename~\thechapter}
  {10pt}
  {\MakeUppercase}

% Titre de partie : même correctif Bold que \chapter ci-dessus. Toujours
% \part* (pas de numéro affiché) : le label reste vide plutôt que d'exposer
% un numéro de partie non requis.
\titleformat{\part}[display]
  {\PURHTitleFont\bfseries\fontsize{16pt}{19pt}\selectfont\centering}
  {}
  {0pt}
  {\MakeUppercase}

% Titre de section (intertitre), 12 pt, capitales. Le référentiel §2.5
% indiquait Josefin Sans Thin ; corrigé en Bold (même correctif que la
% titraille partie/chapitre/contribution ci-dessus) après vérification
% humaine directe du PDF généré face au PDF imprimeur : les intertitres y
% apparaissent noirs et gras, pas maigres (chantier de parité v0.6,
% 2026-08-04). Alignement et espacement non chiffrés par le référentiel pour
% ce niveau : conservés tels quels (raggedright, séparations existantes).
\titleformat{\section}[block]
  {\PURHTitleFont\bfseries\fontsize{12pt}{14pt}\selectfont\raggedright}
  {}
  {0pt}
  {\MakeUppercase}

% Sous-section (référentiel v0.6 §4.3 : seul le niveau "section" — c'est-à-
% dire section1 — est chiffré ; aucune valeur propre à section2/section3
% n'est fournie). Un correctif antérieur avait délibérément retiré le gras
% à ce niveau (confirmé alors sur le livre réel — <div type="section2"> y
% porte de vrais sous-titres phrastiques, jamais des libellés courts) ; la
% vérification humaine directe du 2026-08-04 demande au contraire le même
% traitement noir et gras que les autres intertitres, suivie ici. Capitales
% non demandées à ce niveau, contrairement à \section : inchangé. Tailles
% conservées telles quelles (\large/\normalsize) faute de mesure référentiel.
\titleformat{\subsection}[block]
  {\PURHTitleFont\bfseries\large\raggedright}
  {}
  {0pt}
  {}

\titleformat{\subsubsection}[block]
  {\PURHTitleFont\bfseries\normalsize\raggedright}
  {}
  {0pt}
  {}

\titlespacing*{\part}{0pt}{0pt}{30pt}
\titlespacing*{\chapter}{0pt}{20pt}{18pt}
\titlespacing*{\section}{0pt}{18pt}{10pt}
\titlespacing*{\subsection}{0pt}{14pt}{8pt}
\titlespacing*{\subsubsection}{0pt}{12pt}{6pt}

\addto\captionsfrench{
  \renewcommand{\contentsname}{Table des matières}
}

% Entrées de contribution/front matter (référentiel PURH v0.6 §9.1,
% vérification humaine directe du 2026-08-04) : « titre de la communication
% sans graisse, bas de casse, Chaparral, calé à gauche, série de points
% puis numéro de ligne ». Chaparral Pro reste la fonte ambiante (aucune
% famille à sélectionner) : \PURHTitleFont\bfseries retiré, \hfill remplacé
% par le même filet pointillé que \section ci-dessous. \addvspace{8pt}
% inconditionnel retiré : « pas de saut de ligne entre les références sauf
% changement de section » — l'espacement avant une nouvelle partie vient
% désormais du bloc \part ci-dessous, pas d'ici.
%
% Filet pointillé + numéro de page : laissés dans ce 4e argument dédié (pas
% déplacés dans le texte de l'entrée transmis à \addcontentsline). Une
% première tentative avait déplacé \titlerule*/\contentspage dans ce texte
% pour les faire tomber sur la ligne du titre plutôt que sur celle de
% l'auteur (vérification humaine directe du 2026-08-04 : « les points de
% suite doivent être au niveau du titre, pas de l'auteur ») — abandonnée :
% le paquet bookmark, qui construit automatiquement les signets PDF depuis
% CE MÊME texte d'entrée, ne tolère pas \contentspage/\\ dans cet argument
% ("Token not allowed in a PDF string", puis désynchronisation de
% titlesec — bug réel constaté par compilation). Solution retenue : le nom
% d'auteur n'est plus concaténé DANS le texte de l'entrée de chapitre du
% tout — voir \latei_finish_contribution_toc_entry: (latei_macros.tex), qui
% l'écrit désormais comme une ligne de TDM séparée, via \addtocontents
% (jamais capturée par le mécanisme de signets de bookmark, qui n'observe
% que \addcontentsline/\contentsline).
\titlecontents{chapter}
  [0pt]
  {}
  {}
  {}
  {\titlerule*[0.5pc]{.}\contentspage}

% Entrées de partie (le référentiel dit « titres de section », mais
% désigne bien ici le niveau \part de ce document — les véritables
% intertitres/sections sont exclus de la TDM depuis §9/tocdepth=0) : Josefin
% Sans Bold, capitales (vérification humaine directe du 2026-08-04 : à la
% différence du nom d'auteur sous chaque entrée de contribution, qui doit
% rester bas de casse — voir \lateiTocAuthorLine dans latei_macros.tex —
% les titres de ce niveau doivent au contraire être en petites capitales).
% \scshape ici n'a AUCUN effet sur \PURHTitleFont (Josefin Sans, qui n'a pas
% de véritables petites capitales OpenType — confirmé par compilation
% isolée, "Font shape .../b/sc undefined", substitution silencieuse) : la
% mise en majuscules se fait donc en amont, à la source du texte, via
% \MakeUppercase{#1} dans \lateiRenderHead (latei_macros.tex) — seul le
% corps réduit ici (12 pt contre 16 pt sur la page de la partie elle-même)
% distingue ce niveau de véritables grandes capitales.
\titlecontents{part}
  [0pt]
  {\addvspace{1\baselineskip}\PURHTitleFont\bfseries\fontsize{12pt}{14pt}\selectfont\centering}
  {}
  {}
  {}
  [\addvspace{1\baselineskip}]

\titlecontents{section}
  [1.5em]
  {}
  {}
  {}
  {\titlerule*[0.5pc]{.}\contentspage}

\titlecontents{subsection}
  [3em]
  {\small}
  {}
  {}
  {\titlerule*[0.5pc]{.}\contentspage}

\usepackage{fancyhdr}

\pagestyle{fancy}
\fancyhf{}
% Folios à l'extérieur (référentiel PURH §2.3) : LE (verso) et RO (recto).
% Les titres courants intérieurs (RE, LO) et \chaptermark sont pris en
% charge par latei_macros.tex, qui distingue verso (livre/partie) et recto
% (contribution en cours) — les définir ici aussi serait dupliqué et, pire,
% écrasé silencieusement puisque ce fichier est \input après ce préambule.
\fancyhead[LE]{\PURHHeaderFont\thepage}
\fancyhead[RO]{\PURHHeaderFont\thepage}
\renewcommand{\headrulewidth}{0pt}
\renewcommand{\footrulewidth}{0pt}

\fancypagestyle{plain}{%
  \fancyhf{}%
  \renewcommand{\headrulewidth}{0pt}%
  \renewcommand{\footrulewidth}{0pt}%
}

\usepackage[flushmargin]{footmisc}

\setlength{\footnotesep}{0.6\baselineskip}
% Espace avant les notes (référentiel §5.1 : 3 mm) : \skip\footins régit
% l'espace entre la fin du corps de texte et le début de la zone de notes
% (filet inclus) — 1,2\baselineskip (~5,7 mm à 11/13,5 pt) en donnait trop.
\setlength{\skip\footins}{3mm}
% Filet de notes (référentiel §5.1 : 0,25 pt de large sur 72 pt = 25,4 mm de
% long). \footnoterule par défaut de book.cls fait 0,4 pt sur
% 0,4\columnwidth (~42 mm sur l'empagement de ce profil) — les deux
% constats "environ 0,40 pt" / "environ 42 mm" du §5.2 correspondent très
% exactement à cette valeur par défaut, jamais personnalisée jusqu'ici.
\renewcommand{\footnoterule}{%
  \kern-3pt
  \hrule width 72pt height 0.25pt
  \kern 2.6pt
}
% \footnotelayout (footmisc) n'a plus d'effet : la redéfinition plus bas du
% texte de note pour le retrait négatif de première ligne applique
% directement \fontsize{8.5pt}{10.2pt} et court-circuite
% le mécanisme normal de footmisc qui l'invoque.

\usepackage{etoolbox}

% Citations observées : 9/11 pt, retrait gauche 10 mm (pas de retrait
% droit), ~4 mm avant/après (référentiel PURH §5.3 ; état antérieur :
% 11/14 pt, retraits gauche ET droit ≈1,5em, 8pt avant/après).
\renewenvironment{quote}
  {%
    \par\begingroup
    \fontsize{9pt}{11pt}\selectfont
    \list{}{\leftmargin=10mm\rightmargin=0pt}%
    \item\relax
  }
  {%
    \endlist
    \endgroup
  }

\AtBeginEnvironment{quote}{\vspace*{4mm plus 1pt minus 1pt}}
\AtEndEnvironment{quote}{\vspace*{4mm plus 1pt minus 1pt}}

\usepackage{graphicx}
\usepackage{caption}

\graphicspath{{media/}{images/}{assets/images/}}

\captionsetup{
  font=small,
  labelfont=bf,
  labelsep=period,
  skip=11pt
}

% Référentiel PURH v0.6 §11.3 : titre de tableau 9/11 pt, centré, 10 mm
% avant, 3,5 mm après — réglage propre aux tableaux (\captionsetup[table]),
% distinct du réglage générique ci-dessus qui reste seul à s'appliquer aux
% figures (aucune cible équivalente donnée pour elles dans le référentiel).
\DeclareCaptionFont{PURHTableCaptionFont}{\fontsize{9pt}{11pt}\selectfont}
\captionsetup[table]{
  font=PURHTableCaptionFont,
  labelfont=bf,
  labelsep=period,
  justification=centering,
  aboveskip=10mm,
  belowskip=3.5mm
}

\addto\captionsfrench{
  \renewcommand{\figurename}{Figure}
  \renewcommand{\tablename}{Tableau}
}

\usepackage{array}
\usepackage{longtable}
\usepackage{tabularx}
\usepackage{booktabs}

\usepackage{verse}
\usepackage{ragged2e}
\usepackage{xurl}
\urlstyle{same}

\usepackage[
  unicode=true,
  hidelinks,
  pdfusetitle
]{hyperref}

\usepackage{bookmark}

\hypersetup{
  pdftitle={\PURHBookTitle},
  pdfauthor={\PURHBookAuthor},
  pdfsubject={\PURHPublisher},
  pdfcreator={Impressions},
  pdfproducer={LuaLaTeX}
}

% Numéro calé à gauche avec retrait négatif de première ligne (observé
% directement sur le PDF imprimeur — le référentiel indiquait un point
% après le numéro, contredit par cette observation directe, suivie ici) :
% \leftskip aligne toutes les lignes de la note sur la même marge gauche ;
% \parindent négatif ramène uniquement la première ligne (celle du numéro)
% en deçà de cette marge, jusqu'au bord. \@thefnmark seul, sans point.
% Surtout PAS de \noindent ici : \noindent annule précisément l'effet du
% \parindent négatif qu'il est censé appliquer à la première ligne — bug
% réel constaté après vérification humaine du PDF généré (la première ligne
% restait alignée sur \leftskip comme les suivantes, aucun retrait négatif
% visible) ; laisser LaTeX indenter naturellement la première ligne de
% \parindent (donc la ramener à 0, à la marge) est ce qui produit le retrait
% négatif recherché.
% \AtBeginDocument, pas une simple redéfinition de préambule : hyperref/
% bookmark redéfinissent eux-mêmes \@makefntext pour y ajouter leurs
% ancres, mais le font via leur propre \AtBeginDocument — un
% \renewcommand direct ici, même placé après leur \usepackage, se faisait
% donc encore écraser au début du document (bug réel rencontré et vérifié :
% le correctif n'avait aucun effet tant qu'il n'était pas, lui aussi,
% différé). Étant chargés plus haut, leur crochet s'exécute avant celui-ci.
% \fontsize{8.5pt}{10.2pt} explicite ici, pas laissé au
% \footnotelayout de footmisc (référentiel §5.1 : Chaparral Pro 8,5/10,2 pt) :
% cette redéfinition complète de \@makefntext remplace celle de footmisc et
% n'appelle donc plus \footnotelayout — la taille du corps du texte de note
% retombait silencieusement sur celle ambiante (~9/11 pt observés, cf. §5.2)
% tant que ce défaut n'était pas identifié.
\makeatletter
\AtBeginDocument{%
  \renewcommand{\@makefntext}[1]{%
    \fontsize{8.5pt}{10.2pt}\selectfont
    \setlength{\leftskip}{1.2em}%
    \setlength{\parindent}{-1.2em}%
    \@thefnmark\enskip#1%
  }%
}
\makeatother

% -----------------------------------------------------------------
% Macros utilitaires PURH
% -----------------------------------------------------------------
% Visibilité auteur/affiliation sur l'ouverture de contribution (référentiel
% PURH v0.6 §7.2/§17 P1 item 3) : piloté par le profil de mise en page, pas
% par une valeur fixe — « le profil doit distinguer conservation des
% métadonnées et visibilité sur la page ». Doit être déclaré ici, avant que
% les macros LaTEI ne soient chargées : elles lisent ce drapeau dans
% \lateiContributionAuthor / \lateiContributionAffiliation mais ne le
% déclarent pas elles-mêmes.
\newif\iflateiShowContributionAuthor
\lateiShowContributionAuthorfalse

\newcommand{\PURHSeparator}{%
  \par\addvspace{1.5\baselineskip}%
  \noindent\rule{5cm}{0.4pt}%
  \par\addvspace{1.5\baselineskip}%
}

% Titre principal de la page de titre (référentiel PURH v0.6 §8.1,
% vérification humaine directe du 2026-08-04) : même hauteur de départ que
% le faux-titre (\PURHTitlePage partage désormais \vspace*{0.25\textheight}
% avec \PURHFalseTitle), mais corps plus grand — Josefin Sans Bold
% capitales, comme le faux-titre, juste agrandi. Aucune mesure
% millimétrique donnée par l'utilisateur pour ce corps précis : 22/26 pt
% choisi pour rester nettement plus grand que les 12/14 pt du faux-titre.
\newcommand{\PurhTitleMain}[1]{%
  {\PURHTitleFont\bfseries\fontsize{22pt}{26pt}\selectfont\centering\MakeUppercase{#1}\par}
}

% Sous-titre : Josefin Sans Bold bas de casse (pas de \MakeUppercase, à la
% différence du titre), un peu plus petit, centré sur deux lignes — largeur
% de boîte pour forcer le retour à la ligne, même principe que
% \PURHContributionTitleWidth pour les titres d'ouverture de contribution
% (référentiel §7.3 : reproduire la coupure plutôt que l'exiger en dur).
\newcommand{\PurhSubtitle}[1]{%
  \par\vspace{0.6\baselineskip}%
  \begin{center}
  \parbox{88mm}{\PURHTitleFont\bfseries\fontsize{15pt}{18pt}\selectfont\centering #1}
  \end{center}
  \vspace{0.4\baselineskip}%
}

% Responsabilité éditoriale : Chaparral Pro (fonte principale, aucun
% changement de famille) gras bas de casse, plus petit que le sous-titre,
% sur deux lignes explicites — « sous la direction de » toujours seul sur
% sa propre ligne, puis les noms sur la suivante (référentiel §8.1,
% vérification humaine directe du 2026-08-04). #1 est déjà le texte complet
% des deux lignes, séparées par \\ côté appelant (latei_driver.py), pas
% reconstruites ici : cette macro ne fait que les mettre en forme.
\newcommand{\PurhContributors}[1]{%
  \par\vspace{0.5\baselineskip}%
  {\bfseries\fontsize{11pt}{13pt}\selectfont\centering #1\par}%
  \vspace{0.6\baselineskip}%
}

% Mention finale de la page de titre (référentiel v0.7, vérification
% humaine directe du 2026-08-04, complétée le 2026-08-05) : nom complet de
% l'éditeur, majuscules grasses, sur UNE SEULE ligne, Chaparral (fonte
% principale, aucun changement de famille) — pas \PURHTitleFont/Josefin,
% contrairement à toute la titraille (§2.4/§3 du référentiel v0.7) : cette
% mention n'est pas un niveau de titraille, c'est une mention
% institutionnelle calée en bas de page.
%
% \resizebox{0.95\linewidth}{!}{...} plutôt qu'une simple
% \fontsize{14pt}{16pt} fixe (première version, abandonnée le
% 2026-08-05) : à 14 pt, « PRESSES UNIVERSITAIRES DE ROUEN ET DU HAVRE »
% (44 caractères) ne tenait pas sur une seule ligne à la largeur
% d'empagement de ce profil (~105 mm) et retombait sur deux lignes — bug
% réel constaté par vérification humaine directe du PDF généré.
% \resizebox emballe le texte dans une boîte horizontale non coupable
% (empêchant tout retour à la ligne, contrairement à un simple changement
% de \fontsize) puis la met à l'échelle pour occuper exactement 95 % de la
% largeur de la page — garantit à la fois la ligne unique et le remplissage
% « une bonne partie de la largeur de la page » quel que soit le nom
% affiché ou le profil, sans avoir à calculer un corps de police à la main.
\newcommand{\PurhPublisherMention}[1]{%
  \noindent\resizebox{0.95\linewidth}{!}{\bfseries\MakeUppercase{#1}}\par
}

\newenvironment{PurhBlockQuote}
  {%
    \begin{quote}
  }
  {%
    \end{quote}
  }

\newenvironment{PurhBibliography}
  {%
    \par
    \setlength{\parindent}{0pt}%
    \setlength{\leftskip}{0pt}%
  }
  {%
    \par
  }

% -----------------------------------------------------------------
% Liminaires PURH (référentiel v0.6 §8.1, P0 — "construire la séquence
% complète des liminaires") : pages blanches, faux-titre, crédits, page de
% titre. Construites depuis les métadonnées du livre par latei_driver.py,
% jamais depuis le corps LaTEI réversible. Toutes en pagestyle empty (ni
% folio ni titre courant), mais comptées dans la pagination — jamais
% \begin{titlepage} (dont le mode de compatibilité LaTeX 2.09 remet parfois
% \c@page à 1, un risque inutile ici où le compte doit au contraire
% continuer sans interruption depuis la toute première page physique).
% -----------------------------------------------------------------
\newcommand{\PURHBlankPage}{%
  \clearpage
  \thispagestyle{empty}%
  \mbox{}%
  \clearpage
}

% Vérification humaine directe du 2026-08-04 (faux-titre) :
% remonté sur la page (0.35\textheight -> 0.25\textheight, approximatif —
% « un peu plus haut », aucune mesure millimétrique donnée) et repassé en
% Josefin Sans Bold capitales, même corps que les titres de section
% (\titleformat{\section}, 12/14 pt) — le référentiel ne donnait aucune
% cible pour ce niveau spécifique, seule cette vérification fait foi ici,
% comme pour les autres corrections de graisse de ce chantier.
\newcommand{\PURHFalseTitle}[1]{%
  \clearpage
  \thispagestyle{empty}%
  \begin{center}
  \vspace*{0.25\textheight}
  {\PURHTitleFont\bfseries\fontsize{12pt}{14pt}\selectfont\MakeUppercase{#1}\par}
  \end{center}
  \clearpage
}

% Chaparral Pro (fonte principale du document, aucun changement de famille
% nécessaire) 10 pt — vérification humaine directe du 2026-08-04 : \small
% ne correspond pas nécessairement à 10 pt exactement (dépend de l'échelle
% de tailles du \documentclass), remplacé par une taille explicite.
% \vspace*{\fill} (pas \vspace*{0.3\textheight}, ni un \vfill non étoilé)
% : deuxième vérification humaine directe — le colophon doit être calé en
% bas de page, pas centré verticalement. Un \vfill nu en tout début de page
% est silencieusement absorbé par l'algorithme de coupure de page de TeX
% (jamais visible, confirmé par reproduction : le contenu restait centré
% près du haut) — surtout avec \raggedbottom actif dans ce document, qui
% autorise justement les pages à ne pas s'étirer jusqu'à \textheight.
% \vspace* (étoilé) protège la glue même en tête de page.
\newcommand{\PURHCreditsPage}[1]{%
  \clearpage
  \thispagestyle{empty}%
  \begin{center}
  \vspace*{\fill}
  \fontsize{10pt}{11.5pt}\selectfont
  #1
  \end{center}
  \clearpage
}

% Vérification humaine directe du 2026-08-04 : même hauteur de départ que
% le faux-titre (0.25\textheight, pas 0.15) — seul le corps du titre change
% entre les deux pages, pas sa position verticale.
\newcommand{\PURHTitlePage}[1]{%
  \clearpage
  \thispagestyle{empty}%
  \begin{center}
  \vspace*{0.25\textheight}
  #1
  \end{center}
  \clearpage
}

% ============================================================
% Macros LaTEI — régénérables, ne pas modifier
% ============================================================

% Experimental LaTEI macros for the reversible TEI/PURH convergence path.
% These definitions are typographic only; the reversible source remains the
% controlled LaTEI body produced by purh_site.reversible.latex_writer.

\usepackage{xparse}
\usepackage{xstring}
\usepackage{newunicodechar}

% Unicode spaces remain documentary characters in the reversible body; these
% typographic mappings only prevent missing glyphs in direct LaTEI compilation.
\newunicodechar{ }{~}
\newunicodechar{ }{\nobreak\,}
\newunicodechar{ }{\,}
\newunicodechar{‑}{\nobreak\hbox{-}}
\newunicodechar{″}{\ensuremath{\prime\prime}}
% Titres courants recto/verso distincts (référentiel PURH v0.5, "Titres
% courants") : verso (pages paires, RE) porte le titre du livre ou, à
% l'intérieur d'une partie, celui de la partie ; recto (pages impaires, LO)
% porte le titre court de la contribution en cours. Les deux commandes sont
% ré-exécutées à chaque page par fancyhdr — elles doivent donc rester des
% noms de macro, jamais leur valeur expansée au moment de cet \input.
\newcommand{\lateiCurrentRunningTitle}{}
\newcommand{\lateiVersoRunningTitle}{}
\fancyhead[RE]{\PURHHeaderFont\nouppercase{\lateiVersoRunningTitle}}
\fancyhead[LO]{\PURHHeaderFont\nouppercase{\lateiCurrentRunningTitle}}
\AtBeginDocument{%
  \fancyhead[RE]{\PURHHeaderFont\nouppercase{\lateiVersoRunningTitle}}%
  \fancyhead[LO]{\PURHHeaderFont\nouppercase{\lateiCurrentRunningTitle}}%
}

% Ouverture de contribution (référentiel PURH v0.5, "Ouvertures de
% contribution") : titre, sous-titre, auteur et traducteur viennent des
% <p rend="..."> du <div type="titlePage">, seule source de vérité visible —
% ne pas les dupliquer ailleurs. Titre en Josefin Sans Bold 16 pt capitales
% centré (vérification humaine directe du PDF généré face au PDF imprimeur :
% noir et gras là où le Thin — indiqué par le référentiel §2.5/§5.3, mais
% contredit par cette observation — rendait un texte maigre et grisâtre,
% peu lisible ; même correctif que \titleformat{\part}/{\chapter} dans le
% préambule). Sous-titre d'abord laissé en Thin Italic 12 pt (référentiel
% §5.3, non signalé comme défaut lors de la première vérification), puis
% re-corrigé en Bold Italic bas de casse le 2026-08-04 : nouvelle
% vérification humaine directe montrant explicitement gras + italique sur
% le PDF imprimeur pour ce niveau — pas de majuscules forcées, déjà correct
% (aucun \MakeUppercase ici, contrairement au titre). Positions verticales
% exactes (mm) non calibrées ici — hors périmètre de cette passe.
%
% Coupures de ligne du titre (référentiel PURH v0.5/v0.6 §7.3, "coupures
% éditoriales" — P1 item 4, exécuté le 2026-08-04 sur demande explicite,
% objectif : reproduire la règle implicite du PDF imprimeur plutôt qu'exiger
% une annotation manuelle par titre) : le référentiel avertit qu'il ne faut
% pas reconstruire ces coupures « par une largeur de boîte arbitraire ».
% Vérification empirique sur 7 titres réels du PDF imprimeur de *Dissimuler
% pour mieux régner* (7, 9, 15, 27, 26, 51, 191 caractères environ) montre
% cependant que la largeur de boîte N'EST PAS arbitraire : elle correspond
% de très près (104 mm calculé contre 103,4 mm mesuré indépendamment par
% analyse de pixels sur le rendu réel) à la largeur d'empagement du profil
% 155×230 (105 mm), et un simple \parbox à cette largeur reproduit déjà
% correctement, via le même algorithme Knuth-Plass que tout paragraphe
% LaTeX centré, le nombre de lignes et la coupure exacte de plusieurs titres
% tests (ex. "Les lieux de la conjuration : société secrète et hétérotopie
% dans la littérature romantique", 4 lignes identiques au PDF imprimeur).
% Sur d'autres titres, la coupure obtenue diffère de celle du PDF imprimeur
% d'un mot (ex. "Les espaces du secret à Clarens" → "LES ESPACES DU SECRET
% À" / "CLARENS" ici, contre "LES ESPACES DU SECRET" / "À CLARENS" dans le
% PDF imprimeur) — écart probablement dû à une règle InDesign de
% conservation des mots courts (« à », « de »…) avec la ligne suivante, que
% l'algorithme de LaTeX ne reproduit pas. Valeur retenue : 104mm, meilleur
% compromis observé plutôt qu'une largeur arbitraire non vérifiée.
\newcommand{\PURHContributionTitleWidth}{104mm}
\newcommand{\lateiContributionTitle}[1]{%
  \par\vspace{0.5\baselineskip}%
  \begin{center}
  \parbox{\PURHContributionTitleWidth}{\PURHTitleFont\bfseries\fontsize{16pt}{19pt}\selectfont\centering\MakeUppercase{#1}}%
  \end{center}
  \vspace{0.3\baselineskip}%
}
\newcommand{\lateiContributionSubtitle}[1]{%
  {\PURHTitleFont\bfseries\fontsize{12pt}{14pt}\selectfont\itshape\centering #1\par}%
  \vspace{0.4\baselineskip}%
}
% \iflateiShowContributionAuthor (défini dans le préambule, valeur pilotée
% par le profil de mise en page — référentiel PURH v0.6 §7.2/§17 P1 item 3 :
% « auteur et affiliation non imprimés » sur l'ouverture de contribution,
% mais métadonnées conservées) : le contenu réversible <p rend="author-aut">
% / <p rend="authority_affiliation"> reste inchangé dans le corps LaTEI, seul
% son affichage PDF est basculé par ce drapeau, jamais son existence.
% Signature de fin d'article (référentiel PURH v0.6 §8/§17, vérification
% humaine directe du 2026-08-04) : « les articles sont signés à la fin »,
% calés à droite — un emplacement distinct de l'ouverture de contribution
% (où auteur/affiliation restent masqués, cf. \iflateiShowContributionAuthor
% ci-dessus). \lateiContributionAuthor/\lateiContributionAffiliation
% capturent donc leur contenu dans ces macros globales, INDÉPENDAMMENT du
% drapeau d'affichage en tête de page, pour que \lateiRenderContributionSignature
% (appelée après le corps de la contribution, voir \lateiRenderFrontGroup/
% \lateiRenderChapterGroup) puisse la réémettre à la fin. Réinitialisées à
% vide au début de chaque contribution (mêmes macros) pour ne jamais laisser
% fuiter la signature d'une contribution vers la suivante qui n'en a pas.
% \@empty (le sentinelle vide du noyau LaTeX) exige la catégorie de code 11
% pour "@", seulement active entre \makeatletter et \makeatother — ce
% fichier n'en a pas besoin ailleurs et n'ouvre pas ce bloc ici. \@empty y
% serait tokenisé comme "\@" (caractère isolé) suivi du texte ordinaire
% "empty", qui s'imprimerait littéralement (bug réel constaté par
% compilation : "empty empty" visible sur la page). \lateiSignatureEmpty
% propre, sans "@", sert de sentinelle à la place.
\newcommand{\lateiSignatureEmpty}{}
\newcommand{\lateiSignatureAuthor}{}
\newcommand{\lateiSignatureAffiliation}{}
\newcommand{\lateiResetContributionSignature}{%
  \global\let\lateiSignatureAuthor\lateiSignatureEmpty
  \global\let\lateiSignatureAffiliation\lateiSignatureEmpty
}
% Nom d'auteur sous une entrée de TDM (référentiel §9.1) : écrit comme une
% ligne de TDM SÉPARÉE de l'entrée de chapitre elle-même, via
% \addtocontents (voir \latei_finish_contribution_toc_entry: plus bas) —
% jamais concaténé dans le texte transmis à \addcontentsline{{toc}}{{chapter}}
% : le paquet bookmark, qui construit les signets PDF automatiquement
% depuis CE texte, ne tolère pas les macros qu'on y ajouterait pour
% positionner le filet pointillé/numéro de page au niveau du titre plutôt
% que de l'auteur (bug réel constaté par compilation, voir le commentaire
% sur \titlecontents{{chapter}} dans latei_preamble.py). \par au lieu de la
% précédente rupture "\\" à l'intérieur d'un même paragraphe : ici, il
% s'agit bien d'un paragraphe séparé, plus de retrait suspendu de
% \titlecontents à préserver.
%
% \renewcommand{\textsc}[1]{#1} ICI, dans le corps de \lateiTocAuthorLine
% (vérification humaine directe du 2026-08-05, référentiel v0.7) : le nom
% de l'auteur doit apparaître bas de casse dans la TDM, alors que
% \lateiSignatureAuthor — passé tel quel, sans capture intermédiaire — porte
% le nom de famille en petites capitales (<hi rend="small-caps">Nom</hi>,
% routé vers \teiHi[rend={small-caps}]{Nom} puis \textsc{Nom}).
%
% Une première tentative capturait une copie "texte brut" à l'avance, au
% moment de \lateiContributionAuthor, via \protected@xdef + un \textsc local
% — abandonnée : \teiHi (comme toute commande définie par \NewDocumentCommand
% de xparse) est \protected au sens eTeX, donc jamais développée par un
% \edef/\xdef quel que soit l'état de \textsc au même moment ; le fichier
% .toc obtenu contenait encore \teiHi[rend={small-caps}]{Nom} tel quel (bug
% réel constaté : petites capitales toujours visibles dans le PDF malgré la
% capture). \renewcommand fonctionne en revanche PARFAITEMENT ici : ce bloc
% s'exécute au moment où \tableofcontents \input le fichier .toc et exécute
% RÉELLEMENT \lateiTocAuthorLine{...} — une exécution normale de macro, pas
% un \edef — où \teiHi peut s'exécuter (pas seulement se développer) et
% appeler ce \textsc local, correctement redéfini à ce moment précis.
\newcommand{\lateiTocAuthorLine}[1]{%
  \par\noindent\hspace*{1em}%
  {\bfseries\renewcommand{\textsc}[1]{##1}#1\par}%
}
\newcommand{\lateiContributionAuthor}[1]{%
  \global\def\lateiSignatureAuthor{#1}%
  \iflateiShowContributionAuthor
    {\normalsize\bfseries\centering #1\par}%
    \vspace{0.3\baselineskip}%
  \fi
}
\newcommand{\lateiContributionAffiliation}[1]{%
  \global\def\lateiSignatureAffiliation{#1}%
  \iflateiShowContributionAuthor
    {\small\centering #1\par}%
    \vspace{0.3\baselineskip}%
  \fi
}
% Chaparral Pro (fonte principale, aucun changement de famille), sans
% graisse, corps un peu plus petit que le corps du texte (11 pt -> 10 pt),
% calé à droite. Le prénom en bas de casse et le nom en petites capitales
% viennent déjà tels quels du balisage source (<p rend="author-aut">Prénom
% <hi rend="small-caps">Nom</hi></p>, déjà routé vers \teiHi[rend={small-caps}]
% ailleurs) : #1 n'a besoin d'aucune reformulation ici, seulement d'un
% nouvel habillage.
\newcommand{\lateiRenderContributionSignature}{%
  \ifx\lateiSignatureAuthor\lateiSignatureEmpty\else
    \par\vspace{0.8\baselineskip}%
    {\fontsize{10pt}{12pt}\selectfont\raggedleft\lateiSignatureAuthor\par}%
    \ifx\lateiSignatureAffiliation\lateiSignatureEmpty\else
      {\fontsize{10pt}{12pt}\selectfont\raggedleft\lateiSignatureAffiliation\par}%
    \fi
  \fi
}
% Espace vertical entre le bloc de titre d'ouverture (titre, sous-titre,
% auteur/affiliation) et le début du corps : absent du généré alors que le
% PDF imprimeur montre un grand blanc, estimé à 5-6 lignes de titre (16 pt /
% 19 pt d'interligne) par vérification humaine directe (chantier de parité
% v0.6, 2026-08-04) — 5,5 × 19 pt ≈ 105 pt. Valeur volontairement fixe, pas
% pilotée par le profil (comme les autres corps de la titraille), faute de
% mesure millimétrique du référentiel à ce niveau précis (le §7.2 documente
% une position absolue du corps depuis le haut de page, pas cet écart local).
\newcommand{\PURHContributionBodyGap}{\vspace{105pt}}
\newcommand{\lateiContributionTranslator}[1]{%
  {\small\itshape #1\par}%
  \vspace{0.6\baselineskip}%
}
\NewDocumentCommand{\lateiRenderParagraph}{O{} +m}{%
  \iflateiinfootnote
    #2\unskip\space
  \else
    \IfStrEq{\lateiHeadContext}{titlePage}{%
      \IfSubStr{#1}{rend={title-main}}{\lateiContributionTitle{#2}}{%
        \IfSubStr{#1}{rend={title-sub}}{\lateiContributionSubtitle{#2}}{%
          \IfSubStr{#1}{rend={author-aut}}{\lateiContributionAuthor{#2}}{%
            \IfSubStr{#1}{rend={authority\_affiliation}}{\lateiContributionAffiliation{#2}}{%
              \IfSubStr{#1}{rend={editor-trl}}{\lateiContributionTranslator{#2}}{\par #2\par}%
            }%
          }%
        }%
      }%
    }{%
      \IfStrEq{\lateiHeadContext}{figure}{%
        \IfSubStr{#1}{rend={caption}}{\par\small #2\par}{%
          \IfSubStr{#1}{rend={credits}}{\par\footnotesize #2\par}{\par #2\par}%
        }%
      }{\par #2\par}%
    }%
  \fi
}
\NewDocumentCommand{\teiP}{O{} +m}{\lateiRenderParagraph[#1]{#2}}

% Nested footnotes are not valid LaTeX. If a teiNote appears inside another
% note, render a visible symbolic marker and inline parenthesized content.
\newif\iflateiinfootnote
\lateiinfootnotefalse
\NewDocumentCommand{\teiNote}{O{} +m}{%
  \iflateiinfootnote
    \textsuperscript{*}%
  \else
    \begingroup
      \lateiinfootnotetrue
      \footnote{#2}%
    \endgroup
  \fi
}
\NewDocumentCommand{\teiRef}{O{} +m}{\lateiRenderRef[#1]{#2}}
\NewDocumentCommand{\teiTitle}{O{} +m}{%
  \IfSubStr{#1}{level={m}}{\textit{#2}}{%
    \IfSubStr{#1}{level={j}}{\textit{#2}}{%
      \IfSubStr{#1}{level={a}}{\enquote{#2}}{#2}%
    }%
  }%
}

% <formula notation="latex"> : le contenu documentaire porte déjà ses propres
% délimiteurs de mode mathématique (\(...\) ou \[...\]) — passthrough tel
% quel, sans échappement, pour que LaTeX le compose réellement au lieu
% d'afficher le code source échappé comme texte visible.
\NewDocumentCommand{\teiFormula}{O{} +m}{#2}

% Matter switches keep book.cls's structural side effects (cleardoublepage,
% numbered vs unnumbered chapters via \@mainmatterfalse/true) but not its
% page-numbering side effect. PURH pagination is arabe continue from the
% first page to the last (référentiel PURH v0.5 §5.5/§5.6, P0 — état de
% pagination) : liminaries are never roman-numbered, and the switch to the
% body never restarts the counter at 1. \pagenumbering resets the page
% counter every time it is called, so it must be called at most once for
% the whole document, on whichever of front/main/back matter starts first.
% Every "started" flag is set with \global: each group/chapter is rendered
% from inside its own TeX group (teiElement's xparse "+b" body argument,
% IfSubStr's branch braces...), and a *local* \xxxtrue would revert the
% instant that group closes — silently re-arming the guard and making
% \lateiEnsureMainMatter re-run \cleardoublepage + reset the page counter
% on every single chapter instead of only the first one.
\newif\iflateifrontmatterstarted
\newif\iflateimainmatterstarted
\newif\iflateibackmatterstarted
\newif\iflateiinbibliography
\newif\iflateipagenumberinginitialized
\lateifrontmatterstartedfalse
\lateimainmatterstartedfalse
\lateibackmatterstartedfalse
\lateiinbibliographyfalse
\lateipagenumberinginitializedfalse
\NewDocumentCommand{\lateiEnsureContinuousArabicPagination}{}{%
  \iflateipagenumberinginitialized\else
    \pagenumbering{arabic}%
    \global\lateipagenumberinginitializedtrue
  \fi
}
\makeatletter
\NewDocumentCommand{\lateiEnsureFrontMatter}{}{%
  \iflateifrontmatterstarted\else
    \cleardoublepage
    \global\@mainmatterfalse
    \lateiEnsureContinuousArabicPagination
    \global\lateifrontmatterstartedtrue
  \fi
}
\NewDocumentCommand{\lateiEnsureMainMatter}{}{%
  \iflateimainmatterstarted\else
    \cleardoublepage
    \global\@mainmattertrue
    \lateiEnsureContinuousArabicPagination
    \global\lateimainmatterstartedtrue
  \fi
}
\NewDocumentCommand{\lateiEnsureBackMatter}{}{%
  \iflateibackmatterstarted\else
    \if@openright\cleardoublepage\else\clearpage\fi
    \global\@mainmatterfalse
    \lateiEnsureContinuousArabicPagination
    \global\lateibackmatterstartedtrue
  \fi
}
\makeatother

\ExplSyntaxOn
\tl_new:N \l_latei_option_type_tl
\tl_new:N \l_latei_option_page_title_tl
\tl_new:N \l_latei_option_target_tl
\tl_new:N \l_latei_option_url_tl
\tl_new:N \l_latei_option_n_tl
\tl_new:N \l_latei_option_rend_tl
\tl_new:N \l_latei_option_xmlid_tl
\tl_new:N \l_latei_option_internaltarget_tl
\tl_new:N \l_latei_option_numcols_tl
\tl_new:N \l_latei_graphic_path_tl
\tl_new:N \l_latei_graphic_local_path_tl
\tl_new:N \l_latei_running_title_tl
\tl_new:N \g_latei_current_running_title_tl
\tl_new:N \g_latei_verso_running_title_tl
\prop_new:N \g_latei_graphics_map_prop
\prop_new:N \g_latei_running_titles_map_prop
\cs_generate_variant:Nn \prop_get:NnNTF { NVNTF }
% Le verso par défaut est le titre du livre ; \latei_markboth_verso:n le
% remplace par le titre de la partie en cours tant qu'aucune autre partie
% ne l'écrase (référentiel : "verso : titre du livre ou de la partie").
% \PURHBookTitle est déjà défini à ce stade (préambule inséré avant ce
% fichier), mais on stocke le nom de macro, pas sa valeur expansée : la
% même robustesse que pour \lateiCurrentRunningTitle plus haut.
\tl_gset:Nn \g_latei_verso_running_title_tl { \PURHBookTitle }
\RenewDocumentCommand{\lateiCurrentRunningTitle}{}
  {
    \tl_use:N \g_latei_current_running_title_tl
  }
\RenewDocumentCommand{\lateiVersoRunningTitle}{}
  {
    \tl_use:N \g_latei_verso_running_title_tl
  }
\NewDocumentCommand{\lateiDeclareGraphic}{m m}
  {
    \prop_gput:Nnn \g_latei_graphics_map_prop { #1 } { #2 }
  }
\NewDocumentCommand{\lateiDeclareRunningTitle}{m m}
  {
    \prop_gput:Nnn \g_latei_running_titles_map_prop { #1 } { #2 }
  }
% Recto (titre court de la contribution) et verso (titre du livre ou de la
% partie) sont mis à jour indépendamment — jamais la même valeur des deux
% côtés (référentiel : "passe deux fois la même valeur à \markboth").
\cs_new_protected:Npn \latei_markboth_recto:n #1
  {
    \prop_get:NnNTF \g_latei_running_titles_map_prop { #1 } \l_latei_running_title_tl
      { \tl_gset_eq:NN \g_latei_current_running_title_tl \l_latei_running_title_tl }
      { \tl_gset:Nn \g_latei_current_running_title_tl { #1 } }
    \markboth{\tl_use:N \g_latei_verso_running_title_tl}{\tl_use:N \g_latei_current_running_title_tl}
  }
\cs_new_protected:Npn \latei_markboth_verso:n #1
  {
    \prop_get:NnNTF \g_latei_running_titles_map_prop { #1 } \l_latei_running_title_tl
      { \tl_gset_eq:NN \g_latei_verso_running_title_tl \l_latei_running_title_tl }
      { \tl_gset:Nn \g_latei_verso_running_title_tl { #1 } }
    \markboth{\tl_use:N \g_latei_verso_running_title_tl}{\tl_use:N \g_latei_current_running_title_tl}
  }
\NewDocumentCommand{\lateiMarkBoth}{m}
  {
    \latei_markboth_recto:n { #1 }
  }
\NewDocumentCommand{\lateiMarkBothVerso}{m}
  {
    \latei_markboth_verso:n { #1 }
  }
\RenewDocumentCommand{\chaptermark}{m}{\lateiMarkBoth{#1}}
\cs_new_protected:Npn \latei_chapter:n #1
  {
    \prop_get:NnNTF \g_latei_running_titles_map_prop { #1 } \l_latei_running_title_tl
      {
        \tl_gset_eq:NN \g_latei_current_running_title_tl \l_latei_running_title_tl
        \chapter{#1}
        \lateiMarkBoth{#1}
      }
      {
        \tl_gset:Nn \g_latei_current_running_title_tl { #1 }
        \chapter{#1}
        \lateiMarkBoth{#1}
      }
  }
\NewDocumentCommand{\lateiChapter}{m}
  {
    \latei_chapter:n { #1 }
  }
\keys_define:nn { latei / options }
  {
    type .tl_set:N = \l_latei_option_type_tl,
    data-page-title .tl_set:N = \l_latei_option_page_title_tl,
    target .tl_set:N = \l_latei_option_target_tl,
    url .tl_set:N = \l_latei_option_url_tl,
    n .tl_set:N = \l_latei_option_n_tl,
    rend .tl_set:N = \l_latei_option_rend_tl,
    xmlid .tl_set:N = \l_latei_option_xmlid_tl,
    internaltarget .tl_set:N = \l_latei_option_internaltarget_tl,
    numcols .tl_set:N = \l_latei_option_numcols_tl,
    unknown .code:n = {},
  }
\cs_new_protected:Npn \latei_extract_options:n #1
  {
    \tl_clear:N \l_latei_option_type_tl
    \tl_clear:N \l_latei_option_page_title_tl
    \tl_clear:N \l_latei_option_target_tl
    \tl_clear:N \l_latei_option_url_tl
    \tl_clear:N \l_latei_option_n_tl
    \tl_clear:N \l_latei_option_rend_tl
    \tl_clear:N \l_latei_option_xmlid_tl
    \tl_clear:N \l_latei_option_internaltarget_tl
    \tl_clear:N \l_latei_option_numcols_tl
    \keys_set:nn { latei / options } { #1 }
  }
% Pose une cible PDF nommée d'après le xmlid TEI, s'il est présent dans les
% options déjà extraites par \latei_extract_options:n. Utilisé pour que les
% ancres/citations restent atteignables depuis \lateiRenderRef (cf. cibles
% "#id" internes) même quand l'élément lui-même n'a pas de macro de
% sectionnement propre (qui créerait sinon sa propre destination).
\cs_new_protected:Npn \latei_hypertarget_if_xmlid:
  {
    \tl_if_empty:NF \l_latei_option_xmlid_tl
      { \hypertarget{\tl_use:N \l_latei_option_xmlid_tl}{} }
  }
% Ponts sans "_"/":" : \ExplSyntaxOff plus bas rend ces caractères inaptes
% à composer un nom de macro hors de cette zone ; \teiHead et \teiCit, définis
% après \ExplSyntaxOff, doivent donc passer par ces noms classiques.
\NewDocumentCommand{\lateiExtractOptions}{m}
  {
    \latei_extract_options:n { #1 }
  }
\NewDocumentCommand{\lateiHypertargetIfXmlid}{}
  {
    \latei_hypertarget_if_xmlid:
  }
% Nombre de colonnes du tableau en cours (calculé côté Python, cf.
% _table_column_count dans latex_writer.py — LaTeX ne peut pas le déduire
% à la volée sans avoir déjà lu toutes les lignes).
\NewDocumentCommand{\lateiTableNumCols}{}
  {
    \tl_use:N \l_latei_option_numcols_tl
  }
% Rupture structurante d'ouverture de contribution : saut de page, entrée de
% TDM et titre courant à partir de data-page-title — mais aucun titre visible
% ni numérotation de chapitre ici (référentiel PURH v0.5, "Ouvertures de
% contribution" : « sans libellé ni numérotation de chapitre »). Le titre
% réellement affiché vient du <div type="titlePage"> rendu par ailleurs
% (\lateiContributionTitle) ; le dupliquer ici romprait cette règle.
% \cleardoublepage (jamais \clearpage) : « parties et articles sur recto,
% blancs techniques si nécessaire » (référentiel §5.2) — indépendant de
% l'option de classe openany/openright, qui ne régit que le \chapter brut
% du chemin <div type="chapter">, distinct de celui-ci.
% \setcounter{footnote}{0} : référentiel PURH v0.6 §5/§17, « notes remise à 1
% par contribution » — chaque ouverture (front matter, chapitre/article, back
% matter) redémarre sa propre numérotation de notes plutôt que de poursuivre
% celle de la contribution précédente sur tout le livre.
% \thispagestyle{empty} : référentiel PURH v0.6 §7.2/§17 P1 item 2, « aucun
% en-tête ni folio » sur l'ouverture de contribution — seule cette première
% page bascule en pagestyle vide ; les pages suivantes de la contribution
% reprennent le titre courant normalement (référentiel §6).
% \g_latei_pending_toc_title_tl : l'entrée de TDM elle-même est différée
% jusqu'à \latei_finish_contribution_toc_entry: (appelée après le corps de
% la contribution, comme \lateiRenderContributionSignature) — au moment où
% cette rupture s'exécute, \lateiSignatureAuthor n'est pas encore connu (la
% page de titre, qui le capture, n'a pas encore été rendue : elle fait
% partie de #2, qui s'exécute après). Référentiel PURH v0.6 §9.1,
% vérification humaine directe du 2026-08-04 : « prénom et nom de l'auteur,
% en gras » sous chaque titre de contribution dans la TDM.
\tl_new:N \g_latei_pending_toc_title_tl
\cs_new_protected:Npn \latei_add_contribution_opening_break:
  {
    \tl_if_empty:NF \l_latei_option_page_title_tl
      {
        \cleardoublepage
        \thispagestyle{empty}
        \phantomsection
        \tl_gset_eq:NN \g_latei_pending_toc_title_tl \l_latei_option_page_title_tl
        \exp_args:NV \latei_markboth_recto:n \l_latei_option_page_title_tl
        \setcounter{footnote}{0}
      }
  }
\cs_new_protected:Npn \latei_finish_contribution_toc_entry:
  {
    \tl_if_empty:NF \g_latei_pending_toc_title_tl
      {
        % Titre seul ici : filet pointillé + numéro de page restent gérés
        % par le 4e argument de \titlecontents{chapter} (préambule), comme
        % avant — l'auteur, lui, est écrit séparément ci-dessous, jamais
        % concaténé dans CE texte (voir le commentaire sur
        % \lateiTocAuthorLine, plus haut, pour la raison : compatibilité
        % avec les signets PDF automatiques du paquet bookmark).
        \addcontentsline{toc}{chapter}{\tl_use:N \g_latei_pending_toc_title_tl}
        \ifx\lateiSignatureAuthor\lateiSignatureEmpty\else
          \addtocontents{toc}{\protect\lateiTocAuthorLine{\lateiSignatureAuthor}}
        \fi
        \tl_gclear:N \g_latei_pending_toc_title_tl
      }
  }
\NewDocumentCommand{\lateiRenderFrontGroup}{O{} +m}
  {
    \lateiEnsureFrontMatter
    \latei_extract_options:n { #1 }
    \lateiResetContributionSignature
    \latei_add_contribution_opening_break:
    #2
    \latei_finish_contribution_toc_entry:
    \lateiRenderContributionSignature
  }
\NewDocumentCommand{\lateiRenderChapterGroup}{O{} +m}
  {
    \lateiEnsureMainMatter
    \latei_extract_options:n { #1 }
    \lateiResetContributionSignature
    \latei_add_contribution_opening_break:
    #2
    \latei_finish_contribution_toc_entry:
    \lateiRenderContributionSignature
  }
\NewDocumentCommand{\lateiRenderPartGroup}{O{} +m}
  {
    \lateiEnsureMainMatter
    % Toujours forcé ici, pas seulement via \lateiEnsureMainMatter (dont le
    % \cleardoublepage ne joue qu'à la toute première bascule vers le corps) :
    % chaque partie, même la 2e ou la 3e du livre, doit ouvrir en recto.
    \cleardoublepage
    % référentiel PURH v0.6 §7.1 : « page de style vide » sur l'ouverture de
    % partie — aucun en-tête ni folio, comme pour l'ouverture de contribution.
    \thispagestyle{empty}
    \begingroup
      \lateiSetHeadContext{part}
      #2
    \endgroup
  }
\NewDocumentCommand{\lateiRenderBackGroup}{O{} +m}
  {
    \lateiEnsureBackMatter
    \latei_extract_options:n { #1 }
    \latei_add_contribution_opening_break:
    #2
  }
\NewDocumentCommand{\lateiRenderRef}{O{} +m}
  {
    \latei_extract_options:n { #1 }
    \tl_if_empty:NTF \l_latei_option_internaltarget_tl
      {
        \tl_if_empty:NTF \l_latei_option_target_tl
          { #2 }
          { \href{\tl_use:N \l_latei_option_target_tl}{#2} }
      }
      { \hyperlink{\tl_use:N \l_latei_option_internaltarget_tl}{#2} }
  }
\cs_new_protected:Npn \latei_figure_fallback:
  {
    \fbox{\parbox{0.8\linewidth}{\centering\footnotesize Image absente ou non fournie}}
  }
\cs_new_protected:Npn \latei_include_graphic:n #1
  {
    \tl_if_blank:nTF { #1 }
      { \latei_figure_fallback: }
      {
        \IfFileExists{\detokenize{#1}}
          { \includegraphics[width=0.95\linewidth,keepaspectratio]{\detokenize{#1}} }
          { \latei_figure_fallback: }
      }
  }
\NewDocumentCommand{\lateiRenderGraphic}{O{}}
  {
    \latei_extract_options:n { #1 }
    \tl_clear:N \l_latei_graphic_path_tl
    \tl_if_empty:NF \l_latei_option_url_tl
      { \tl_set_eq:NN \l_latei_graphic_path_tl \l_latei_option_url_tl }
    \tl_if_empty:NT \l_latei_graphic_path_tl
      {
        \tl_if_empty:NF \l_latei_option_target_tl
          { \tl_set_eq:NN \l_latei_graphic_path_tl \l_latei_option_target_tl }
      }
    \tl_if_empty:NT \l_latei_graphic_path_tl
      {
        \tl_if_empty:NF \l_latei_option_n_tl
          { \tl_set_eq:NN \l_latei_graphic_path_tl \l_latei_option_n_tl }
      }
    \prop_get:NVNTF \g_latei_graphics_map_prop \l_latei_graphic_path_tl \l_latei_graphic_local_path_tl
      { \exp_args:NV \latei_include_graphic:n \l_latei_graphic_local_path_tl }
      { \exp_args:NV \latei_include_graphic:n \l_latei_graphic_path_tl }
  }
\ExplSyntaxOff

% Référentiel PURH v0.6 §11.2 : Chaparral Pro 10 pt, retrait suspendu 5 mm.
% \hangindent en em dépendait de la taille de fonte ambiante (jamais fixée
% par cette macro) ; passé en mm, valeur absolue conforme à la cible.
\NewDocumentCommand{\lateiBibliographyEntry}{+m}{%
  \par\noindent\fontsize{10pt}{12pt}\selectfont\hangindent=5mm\hangafter=1 #1\par
}
\NewDocumentCommand{\lateiBibliographyBlock}{+m}{%
  \par
  \begingroup
    \lateiinbibliographytrue
    \begin{PurhBibliography}
    #1%
    \end{PurhBibliography}
  \endgroup
  \par
}

% Heads are contextual in TEI: the same <head> can name a chapter, section,
% figure, table, or list. The surrounding LaTEI environment sets this context.
\newcommand{\lateiHeadContext}{fallback}
\newcommand{\lateiSetHeadContext}[1]{\def\lateiHeadContext{#1}}
\NewDocumentCommand{\lateiRenderHead}{m}{%
  % \part* ajoute déjà lui-même son entrée de TDM sous le shape [display] de
  % \titleformat{\part} (vérifié empiriquement) : un \addcontentsline{toc}
  % {part}{...} manuel ici produirait un doublon, pas un renfort.
  %
  % \MakeUppercase{#1} ici (et non \scshape/\textsc dans
  % \titlecontents{part}, préambule) : vérification humaine directe du
  % 2026-08-05 a montré que \scshape n'a AUCUN effet visible sur
  % \PURHTitleFont (Josefin Sans) — confirmé par compilation isolée, le
  % journal LaTeX rapporte explicitement « Font shape .../b/sc undefined »,
  % silencieusement remplacée par la forme normale (bug réel, jamais une
  % erreur bloquante, donc invisible sans vérifier le journal). Cette fonte
  % n'a pas de véritables petites capitales OpenType. Palliatif : capitales
  % pleines (\MakeUppercase), à la taille réduite (12 pt) déjà fixée par
  % \titlecontents{part} — plus petites que le titre de partie affiché sur
  % sa propre page (16 pt), donc lisibles comme un niveau distinct dans la
  % TDM, sans être de véritables petites capitales typographiques.
  % Appliqué ici (à la source du texte de la partie, avant \part*) plutôt
  % que dans \titlecontents{part} : \part* réutilise ce même #1, déjà
  % majuscule, pour son propre \addcontentsline{{toc}}{{part}}{{...}} interne —
  % la mise en majuscules du titre affiché sur la page (via le \MakeUppercase
  % déjà présent dans \titleformat{{\part}}) n'est pas affectée puisqu'elle
  % s'applique une seconde fois, sans effet, sur un texte déjà majuscule.
  \IfStrEq{\lateiHeadContext}{part}{\part*{\MakeUppercase{#1}}\lateiMarkBothVerso{#1}}{%
  \IfStrEq{\lateiHeadContext}{chapter}{\lateiChapter{#1}}{%
    \IfStrEq{\lateiHeadContext}{unnumberedchapter}{\chapter*{#1}\lateiMarkBoth{#1}\addcontentsline{toc}{chapter}{#1}}{%
    \IfStrEq{\lateiHeadContext}{section}{\section{#1}}{%
      \IfStrEq{\lateiHeadContext}{subsection}{\subsection{#1}}{%
        \IfStrEq{\lateiHeadContext}{subsubsection}{\subsubsection{#1}}{%
          \IfStrEq{\lateiHeadContext}{figure}{\par\smallskip\noindent\textbf{#1}\par\smallskip}{%
            \IfStrEq{\lateiHeadContext}{table}{\par\smallskip\noindent\textbf{#1}\par\smallskip}{%
              \IfStrEq{\lateiHeadContext}{list}{\item[]\textbf{#1}}{%
                \par\smallskip\noindent\textbf{#1}\par\smallskip% fallback head
              }%
            }%
          }%
        }%
      }%
    }%
  }%
  }%
  }%
}
\NewDocumentCommand{\teiHead}{O{} +m}{%
  \lateiExtractOptions{#1}%
  \lateiHypertargetIfXmlid
  \lateiRenderHead{#2}%
}

% Apply rend values predictably, including combined values in any order. This
% mirrors the stable renderer order: small caps / position / weight / italic.
\NewDocumentCommand{\lateiIfRendSmallCaps}{m m m}{%
  \IfSubStr{\detokenize{#1}}{\detokenize{small-caps}}{#2}{%
    \IfSubStr{\detokenize{#1}}{\detokenize{small_caps}}{#2}{%
      \IfSubStr{\detokenize{#1}}{\detokenize{small caps}}{#2}{#3}%
    }%
  }%
}
\NewDocumentCommand{\lateiIfRendBold}{m m m}{%
  \IfSubStr{\detokenize{#1}}{\detokenize{bold}}{#2}{\IfSubStr{\detokenize{#1}}{\detokenize{gras}}{#2}{#3}}%
}
\NewDocumentCommand{\lateiIfRendSup}{m m m}{%
  \IfSubStr{\detokenize{#1}}{\detokenize{sup}}{#2}{\IfSubStr{\detokenize{#1}}{\detokenize{exposant}}{#2}{#3}}%
}
\NewDocumentCommand{\lateiIfRendSub}{m m m}{%
  \IfSubStr{\detokenize{#1}}{\detokenize{sub}}{#2}{\IfSubStr{\detokenize{#1}}{\detokenize{indice}}{#2}{#3}}%
}
\NewDocumentCommand{\lateiIfRendUnderline}{m m m}{%
  \IfSubStr{\detokenize{#1}}{\detokenize{underline}}{#2}{\IfSubStr{\detokenize{#1}}{\detokenize{souligne}}{#2}{#3}}%
}
\NewDocumentCommand{\lateiIfRendStrikethrough}{m m m}{%
  \IfSubStr{\detokenize{#1}}{\detokenize{strikethrough}}{#2}{\IfSubStr{\detokenize{#1}}{\detokenize{barre}}{#2}{#3}}%
}
\NewDocumentCommand{\lateiIfRendUppercase}{m m m}{%
  \IfSubStr{\detokenize{#1}}{\detokenize{uppercase}}{#2}{\IfSubStr{\detokenize{#1}}{\detokenize{majuscule}}{#2}{#3}}%
}
\NewDocumentCommand{\lateiHiSmallCaps}{m +m}{%
  \lateiIfRendSmallCaps{#1}{\textsc{#2}}{#2}%
}
\NewDocumentCommand{\lateiHiPosition}{m +m}{%
  \lateiIfRendSup{#1}{\textsuperscript{#2}}{%
    \lateiIfRendSub{#1}{$_{#2}$}{#2}%
  }%
}
\NewDocumentCommand{\lateiHiBold}{m +m}{%
  \lateiIfRendBold{#1}{\textbf{#2}}{#2}%
}
\NewDocumentCommand{\lateiHiItalic}{m +m}{%
  \IfSubStr{\detokenize{#1}}{\detokenize{italic}}{\textit{#2}}{#2}%
}
\NewDocumentCommand{\lateiHiUnderline}{m +m}{%
  \lateiIfRendUnderline{#1}{\underline{#2}}{#2}%
}
\NewDocumentCommand{\lateiHiStrikethrough}{m +m}{%
  \lateiIfRendStrikethrough{#1}{\sout{#2}}{#2}%
}
\NewDocumentCommand{\lateiHiUppercase}{m +m}{%
  \lateiIfRendUppercase{#1}{\MakeUppercase{#2}}{#2}%
}
\NewDocumentCommand{\teiHi}{O{} +m}{%
  \lateiHiItalic{#1}{\lateiHiBold{#1}{\lateiHiPosition{#1}{\lateiHiSmallCaps{#1}{\lateiHiUnderline{#1}{\lateiHiStrikethrough{#1}{\lateiHiUppercase{#1}{#2}}}}}}}%
}

\NewDocumentCommand{\teiForeign}{O{} +m}{\textit{#2}}
\NewDocumentCommand{\teiTerm}{O{} +m}{#2}
\NewDocumentCommand{\teiName}{O{} +m}{#2}
\NewDocumentCommand{\teiPersName}{O{} +m}{#2}
\NewDocumentCommand{\teiPlaceName}{O{} +m}{#2}
\NewDocumentCommand{\teiOrgName}{O{} +m}{#2}
\NewDocumentCommand{\teiDate}{O{} +m}{#2}
\NewDocumentCommand{\teiNum}{O{} +m}{#2}
\NewDocumentCommand{\teiLabel}{O{} +m}{#2}
\NewDocumentCommand{\teiQ}{O{} +m}{\enquote{#2}}
\NewDocumentCommand{\teiSaid}{O{} +m}{\enquote{#2}}
\NewDocumentCommand{\teiAuthor}{O{} +m}{#2}
\NewDocumentCommand{\teiEditor}{O{} +m}{#2}
\NewDocumentCommand{\teiPublisher}{O{} +m}{#2}
\NewDocumentCommand{\teiBiblScope}{O{} +m}{#2}
\NewDocumentCommand{\teiIdno}{O{} +m}{#2}

\NewDocumentCommand{\teiGraphic}{O{}}{\lateiRenderGraphic[#1]}
\NewDocumentCommand{\teiPtr}{O{}}{}
\NewDocumentCommand{\teiLb}{O{}}{\linebreak{}}
\NewDocumentCommand{\teiPb}{O{}}{\ignorespaces}
% <anchor xml:id="..."/> : marqueur de position TEI, sans contenu visible —
% seule sa cible PDF (pour les <ref target="#..."> qui pointent dessus) compte.
\NewDocumentCommand{\teiAnchor}{O{}}{%
  \lateiExtractOptions{#1}%
  \lateiHypertargetIfXmlid
}

\NewDocumentEnvironment{teiDiv}{O{} +b}{%
  \IfSubStr{#1}{type={titlePage}}{%
    \begingroup
    \lateiSetHeadContext{titlePage}%
    #2%
    \PURHContributionBodyGap
    \endgroup
  }{%
    \begingroup
    \lateiSetHeadContext{fallback}%
    \IfSubStr{#1}{type={chapter}}{\lateiSetHeadContext{chapter}}{}%
    \IfSubStr{#1}{type={section}}{\lateiSetHeadContext{section}}{}%
    \IfSubStr{#1}{type={section1}}{\lateiSetHeadContext{section}}{}%
    \IfSubStr{#1}{type={section2}}{\lateiSetHeadContext{subsection}}{}%
    \IfSubStr{#1}{type={subsection}}{\lateiSetHeadContext{subsection}}{}%
    \IfSubStr{#1}{type={section3}}{\lateiSetHeadContext{subsubsection}}{}%
    #2%
    \endgroup
  }%
}{}
\NewDocumentEnvironment{teiList}{O{} +b}{%
  \begingroup
  \lateiSetHeadContext{list}%
  \IfSubStr{#1}{type={ordered}}{\begin{enumerate}}{\begin{itemize}}%
  #2%
  \IfSubStr{#1}{type={ordered}}{\end{enumerate}}{\end{itemize}}%
  \endgroup
}{}
\NewDocumentCommand{\teiItem}{O{} +m}{\item #2}
\NewDocumentEnvironment{teiQuote}{O{} +b}{\begin{PurhBlockQuote}#2\end{PurhBlockQuote}}{}
% <quote> utilisée en cours de phrase (dans <p>/<item>, ou n'importe où
% dans <note> — cf. _write_quote côté Python, qui décide du choix entre
% teiQuote et teiQuoteInline) : \teiQuote/PurhBlockQuote force un
% \begin{quote}, qui coupe le paragraphe même en plein milieu d'une phrase.
% \teiQuoteInline reste une macro simple (aucun \begingroup ni saut de
% paragraphe) et ajoute les guillemets via \enquote (csquotes), comme
% \teiQ déjà.
\NewDocumentCommand{\teiQuoteInline}{O{} +m}{\enquote{#2}}
% Poésie (référentiel PURH v0.5, "Poésie") : de vrais <lg>/<l> TEI passent
% par l'environnement `verse` (aligné à gauche, une ligne par \\, jamais
% justifié) — pas le bloc de citation, qui produisait des blancs excessifs
% en forçant des \linebreak dans un paragraphe justifié. <lg> imbriqués
% (strophes) s'emboîtent naturellement, `verse` gère leur espacement.
\NewDocumentEnvironment{teiLg}{O{} +b}{%
  \begin{verse}
  #2%
  \end{verse}
}{}
\NewDocumentCommand{\teiL}{O{} +m}{#2\\}
\NewDocumentEnvironment{teiFigure}{O{} +b}{%
  \begingroup
  \lateiSetHeadContext{figure}%
  \begin{center}#2\end{center}%
  \endgroup
}{}
\NewDocumentEnvironment{teiCit}{O{} +b}{%
  \lateiExtractOptions{#1}%
  \lateiHypertargetIfXmlid
  #2%
}{}
\NewDocumentEnvironment{teiBibl}{O{} +b}{\lateiBibliographyEntry{#2}}{}

% Tableau réel (longtable/booktabs), pas une simple suite de paragraphes
% centrés : le nombre de colonnes (numcols, calculé côté Python — cf.
% _table_column_count) fixe le gabarit de colonnes. Les teiCell d'une ligne
% sont déjà séparées par un "&" littéral dans la source (écrit côté Python,
% cf. _write_row), et teiRow/teiCell sont des macros (pas des environnements) :
% \multicolumn doit être le tout premier jeton lu par le scanner d'alignement
% de TeX pour une cellule fusionnée (@cols>1) — même non protégée, l'envelopper
% dans une macro (et a fortiori un environnement, qui insère un \begingroup
% avant le corps) le rend invisible à la détection de \omit ("Misplaced
% \omit", confirmé par reproduction isolée). \multicolumn est donc écrit tel
% quel côté Python (cf. _write_cell) pour les cellules fusionnées ; teiCell ne
% gère plus que la mise en forme (gras d'entête), jamais \multicolumn. Une
% ligne d'entête (role=label/header sur la ligne, cf. TEI) reçoit
% \toprule/\midrule ; \bottomrule ferme le tableau.
% Référentiel PURH v0.6 §11.3 : texte 8/9,5 pt, marges internes 2 mm,
% filets 0,25 pt, fond d'entête noir 30 % (un gris clair : 30 % de couverture
% d'encre, pas un fond sombre). Chaparral Pro déjà la fonte principale du
% document (\setmainfont) : aucun changement de famille nécessaire, seule la
% taille change ici. \heavyrulewidth/\lightrulewidth/\cmidrulewidth
% (booktabs) valent par défaut 0,08em/0,05em/0,05em — proportionnels à la
% fonte ambiante, donc jamais 0,25 pt de façon fiable sans réglage explicite.
\newif\iflateirowisheader
\lateirowisheaderfalse
\NewDocumentEnvironment{teiTable}{O{} +b}{%
  \lateiExtractOptions{#1}%
  \begingroup
  \lateiSetHeadContext{table}%
  \fontsize{8pt}{9.5pt}\selectfont
  \setlength{\tabcolsep}{2mm}
  \setlength{\heavyrulewidth}{0.25pt}
  \setlength{\lightrulewidth}{0.25pt}
  \setlength{\cmidrulewidth}{0.25pt}
  \par\begin{center}%
  \begin{longtable}{*{\lateiTableNumCols}{l}}
  \toprule
  #2%
  \bottomrule
  \end{longtable}%
  \end{center}\par%
  \endgroup
}{}
% \rowcolor pour les lignes d'entête est écrit directement par
% purh_site/reversible/latex_writer.py (_write_row), avant même ce \teiRow,
% jamais ici : même contrainte que \multicolumn sur une cellule fusionnée
% (cf. teiCell) — un \rowcolor conditionnel à l'intérieur de cette macro
% n'est plus le premier jeton vu par le scanner d'alignement de TeX et
% produit "Misplaced \noalign" (confirmé par reproduction isolée).
\NewDocumentCommand{\teiRow}{O{} +m}{%
  \IfSubStr{#1}{role={label}}{\lateirowisheadertrue}{%
    \IfSubStr{#1}{role={header}}{\lateirowisheadertrue}{\lateirowisheaderfalse}%
  }%
  #2%
  \\%
  \iflateirowisheader\midrule\lateirowisheaderfalse\fi
}
\NewDocumentCommand{\teiCell}{O{} +m}{%
  \lateiExtractOptions{#1}%
  \iflateirowisheader\bfseries\fi
  \IfSubStr{#1}{role={label}}{\bfseries}{\IfSubStr{#1}{role={header}}{\bfseries}{}}%
  #2%
}
% teiHeader is metadata, not running text. Document wrappers pass through;
% unknown non-header elements keep their content visible rather than vanish.
\NewDocumentEnvironment{teiElement}{O{} +b}{%
  \IfSubStr{#1}{name={teiHeader}}{}{%
    \IfSubStr{#1}{name={group}}{%
      \IfSubStr{#1}{type={section1}}{\lateiRenderPartGroup[#1]{#2}}{%
      \IfSubStr{#1}{type={article}}{\lateiRenderChapterGroup[#1]{#2}}{%
      \IfSubStr{#1}{type={chapter}}{\lateiRenderChapterGroup[#1]{#2}}{%
      \IfSubStr{#1}{type={acknowledgments}}{\lateiRenderFrontGroup[#1]{#2}}{%
      \IfSubStr{#1}{type={abbreviations}}{\lateiRenderFrontGroup[#1]{#2}}{%
      \IfSubStr{#1}{type={introduction}}{\lateiRenderFrontGroup[#1]{#2}}{%
      \IfSubStr{#1}{type={foreword}}{\lateiRenderFrontGroup[#1]{#2}}{%
      \IfSubStr{#1}{type={preface}}{\lateiRenderFrontGroup[#1]{#2}}{%
      \IfSubStr{#1}{type={dedication}}{\lateiRenderFrontGroup[#1]{#2}}{%
      \IfSubStr{#1}{type={conclusion}}{\lateiRenderBackGroup[#1]{#2}}{%
      \IfSubStr{#1}{type={bibliography}}{\lateiRenderBackGroup[#1]{#2}}{#2}%
      }}}}}}}}}}%
    }{%
      \IfSubStr{#1}{name={back}}{\lateiEnsureBackMatter #2}{%
      \IfSubStr{#1}{name={listBibl}}{\lateiBibliographyBlock{#2}}{%
      \IfSubStr{#1}{name={biblStruct}}{\lateiBibliographyEntry{#2}}{#2}%
      }}%
    }%
  }%
}{}

% ---------------------------------------------------------------------------
% Commandes \latei... — corrections de mise en page PDF uniquement.
% Ces commandes ne sont pas restaurées vers XML lors du retour Métopes.
% Le parseur LaTEI parcourt leur contenu sans créer d'élément TEI.
% ---------------------------------------------------------------------------

% Macro interne — traduit les tailles normalisées en dimension PDF.
% Valeurs acceptées : small (~0,5 ligne), medium (~1 ligne), large (~2 lignes).
\NewDocumentCommand{\lateiApplyVSpace}{m}{%
  \IfStrEqCase{#1}{%
    {small}{\addvspace{0.5\baselineskip}}%
    {medium}{\addvspace{\baselineskip}}%
    {large}{\addvspace{2\baselineskip}}%
  }[\PackageError{latei}{%
    Valeur inconnue : #1. Valeurs acceptees : small, medium, large.%
  }{}]%
}

% Indentation — alinéas
% --------------------------------------------------------------------------

% \lateiNoIndent{...} — supprime le retrait de première ligne du contenu.
\NewDocumentCommand{\lateiNoIndent}{+m}{{\parindent=0pt #1}}

% \lateiIndent{...} — force le retrait de première ligne même si LaTeX le supprime.
\newlength{\lateiParindent}
\AtBeginDocument{\setlength{\lateiParindent}{\parindent}}
\NewDocumentCommand{\lateiIndent}{+m}{{\setlength{\parindent}{\lateiParindent}#1}}

% Espacement vertical autonome
% --------------------------------------------------------------------------

% \lateiVSpace{small|medium|large} — insère un blanc vertical entre deux éléments.
\NewDocumentCommand{\lateiVSpace}{m}{\par\lateiApplyVSpace{#1}}

% Espacement vertical enveloppant
% --------------------------------------------------------------------------

% \lateiSpaceBefore{small|medium|large}{...} — ajoute un blanc avant le contenu.
\NewDocumentCommand{\lateiSpaceBefore}{m +m}{\par\lateiApplyVSpace{#1}#2}

% \lateiSpaceAfter{small|medium|large}{...} — ajoute un blanc après le contenu.
\NewDocumentCommand{\lateiSpaceAfter}{m +m}{#2\par\lateiApplyVSpace{#1}}

% Sauts de page autonomes
% --------------------------------------------------------------------------

% \lateiPageBreak — force un saut de page.
\NewDocumentCommand{\lateiPageBreak}{}{\newpage}

% \lateiClearPage — force une nouvelle page en vidant les flottants.
\NewDocumentCommand{\lateiClearPage}{}{\clearpage}

% Sauts de page enveloppants
% --------------------------------------------------------------------------

% \lateiPageBreakBefore{...} — force un saut de page avant le contenu.
\NewDocumentCommand{\lateiPageBreakBefore}{+m}{\newpage#1}

% \lateiPageBreakAfter{...} — force un saut de page après le contenu.
\NewDocumentCommand{\lateiPageBreakAfter}{+m}{#1\newpage}

% \lateiClearPageBefore{...} — vide les flottants et force une nouvelle page avant.
\NewDocumentCommand{\lateiClearPageBefore}{+m}{\clearpage#1}

% \lateiClearPageAfter{...} — force une nouvelle page avec vidage des flottants après.
\NewDocumentCommand{\lateiClearPageAfter}{+m}{#1\clearpage}

% Contrôle des coupures de page
% --------------------------------------------------------------------------

% \lateiKeepWithNext{...} — interdit une coupure de page immédiatement après.
\NewDocumentCommand{\lateiKeepWithNext}{+m}{#1\nopagebreak[4]}

% \lateiKeepTogether{...} — maintient le contenu entier sur une même page.
\NewDocumentCommand{\lateiKeepTogether}{+m}{\begin{samepage}#1\end{samepage}}

% \lateiNoPageBreakBefore{...} — interdit une coupure de page immédiatement avant.
\NewDocumentCommand{\lateiNoPageBreakBefore}{+m}{\nopagebreak[4]#1}

% \lateiNoPageBreakAfter{...} — interdit une coupure de page immédiatement après.
\NewDocumentCommand{\lateiNoPageBreakAfter}{+m}{#1\nopagebreak[4]}

% lateiDocument marks the reversible editorial zone in a monofile.
% Completely transparent at compilation — only serves as a delimiter for the parser.
\newenvironment{lateiDocument}{}{}

% ============================================================
% Mappings graphiques — régénérables, ne pas modifier
% ============================================================

% Experimental LaTEI graphic mapping.
% This file is a compilation artifact, not a reversible source.

% ============================================================
% Mappings titres courants — régénérables, ne pas modifier
% ============================================================

% Experimental LaTEI running-title mapping.
% This file is a compilation artifact, not a reversible source.
\lateiDeclareRunningTitle{Plotin contre Platon ? Chercher la beauté dans le vivant et dans la vie}{Plotin contre Platon ? Chercher la beauté dans le vivant…}
\lateiDeclareRunningTitle{Vieux et laid, mais sage : éloges romains de la laideur}{Vieux et laid, mais sage : éloges romains de la laideur}
\lateiDeclareRunningTitle{Retour à « l’aisthétique »}{Retour à « l’aisthétique »}
\lateiDeclareRunningTitle{Richard Owen : la beauté de l’adaptation}{Richard Owen : la beauté de l’adaptation}
\lateiDeclareRunningTitle{Ernst Haeckel (1834-1919) et les « formes artistiques de la nature »}{Ernst Haeckel (1834-1919) et les « formes artistiques…}
\lateiDeclareRunningTitle{Les métamorphoses de la morphologie dans la paléontologie évolutive du xx e  siècle}{Les métamorphoses de la morphologie dans la paléontologie…}
\lateiDeclareRunningTitle{« Étranges merveilles »}{« Étranges merveilles »}
\lateiDeclareRunningTitle{« Art fossile »}{« Art fossile »}
\lateiDeclareRunningTitle{Beautés animales : ordre ou expression ?}{Beautés animales : ordre ou expression ?}
\lateiDeclareRunningTitle{« Paraître est une fonction vitale »}{« Paraître est une fonction vitale »}
\lateiDeclareRunningTitle{Les « bestioles »}{Les « bestioles »}
\lateiDeclareRunningTitle{Herbes et « bestioles » ou le paradigme de la beauté vitale}{Herbes et « bestioles » ou le paradigme de la beauté…}
\lateiDeclareRunningTitle{« Créer c’est résister »}{« Créer c’est résister »}
\lateiDeclareRunningTitle{Poétique du jardin japonais : la villa détachée et les jardins de Katsura}{Poétique du jardin japonais : la villa détachée…}
\lateiDeclareRunningTitle{« Le » jardin japonais, idéal, acmè , parangon occidental de beauté}{« Le » jardin japonais, idéal, acmè , parangon occidental…}
\lateiDeclareRunningTitle{Un impossible mur de feuilles : la haie Katsura gaki}{Un impossible mur de feuilles : la haie Katsura gaki}
\lateiDeclareRunningTitle{Le « chemin de rosée »}{Le « chemin de rosée »}
\lateiDeclareRunningTitle{Quels paradigmes de beauté dans les arts ?}{Quels paradigmes de beauté dans les arts ?}
\lateiDeclareRunningTitle{Notes pour une philosophie de la tension de surface ou La beauté comme relation}{Notes pour une philosophie de la tension de surface…}
\lateiDeclareRunningTitle{Sauver les phénomènes : la beauté vitale en photographie}{Sauver les phénomènes : la beauté vitale en photographie}
\lateiDeclareRunningTitle{Le grain : image/information/technique}{Le grain : image/information/technique}
\lateiDeclareRunningTitle{L’invention photographique perpétuée : le point clé}{L’invention photographique perpétuée : le point clé}
\lateiDeclareRunningTitle{La beauté, une expérience socialement définie ?}{La beauté, une expérience socialement définie ?}
\lateiDeclareRunningTitle{Enculturation musicale : apprendre une nouvelle culture}{Enculturation musicale : apprendre une nouvelle culture}
\lateiDeclareRunningTitle{Pourquoi aimons-nous la musique (que nous aimons) ?}{Pourquoi aimons-nous la musique (que nous aimons) ?}
\lateiDeclareRunningTitle{Beauté musicale ?}{Beauté musicale ?}
\lateiDeclareRunningTitle{Approche sociologique de la beauté en situation de vulnérabilité}{Approche sociologique de la beauté en situation…}
\lateiDeclareRunningTitle{Les évolutions des critères de beauté, entre représentations individuelles et valeurs sociétales}{Les évolutions des critères de beauté, entre…}
\lateiDeclareRunningTitle{Définition de la vulnérabilité et place de la beauté dans une situation de vulnérabilité}{Définition de la vulnérabilité et place de la beauté…}
\lateiDeclareRunningTitle{L’origine vitale de la beauté : la vie comme relation}{L’origine vitale de la beauté : la vie comme relation}
\lateiDeclareRunningTitle{Le paradoxe de la beauté : la contemplation vitale}{Le paradoxe de la beauté : la contemplation vitale}
\lateiDeclareRunningTitle{L’enjeu vital de la beauté : joie et horreur, dans toutes les dimensions de nos vies}{L’enjeu vital de la beauté : joie et horreur, dans toutes…}
\lateiDeclareRunningTitle{« Là, tout n’est qu’ordre et beauté / Luxe, calme et volupté »}{« Là, tout n’est qu’ordre et beauté / Luxe, calme…}
\lateiDeclareRunningTitle{Quelques facéties sur le colloque « Beautés vitales » de Cerisy-la-Salle}{Quelques facéties sur le colloque « Beautés vitales »…}

\begin{document}

\lateiEnsureContinuousArabicPagination
\PURHBlankPage
\PURHBlankPage
\PURHFalseTitle{Beautés vitales}
\PURHCreditsPage{© Presses universitaires de Rouen et du Havre.\par\vspace{0.1\baselineskip}2 place Émile Blondel – 76821 Mont-Saint-Aignan Cedex\par\vspace{0.1\baselineskip}\url{http://purh.univ-rouen.fr}\par\vspace{0.1\baselineskip}ISBN : 979-10-240-1926-0\par}
\PURHTitlePage{\PurhTitleMain{Beautés vitales}
\PurhSubtitle{Pour une approche contemporaine de la beauté}
\vspace{2\baselineskip}
\PurhContributors{sous la direction de\\Anne-Lise Worms et Clélia Zernik}
\vspace*{\fill}
\PurhPublisherMention{Presses universitaires de Rouen et du Havre}}
\PURHBlankPage

% ============================================================
% Zone éditoriale réversible — corrections autorisées ici
% ============================================================

\begin{lateiDocument}
\begin{teiElement}[name={TEI},change={commons\_edition}]

  \begin{teiElement}[name={teiHeader}]

    \begin{teiElement}[name={fileDesc}]

      \begin{teiElement}[name={titleStmt}]

        \teiTitle[type={main}]{Beautés vitales}
        \teiTitle[type={sub}]{Pour une approche contemporaine de la beauté}
        \teiAuthor[role={pbd}]{
          \teiPersName{
            \begin{teiElement}[name={forename}]
Anne-Lise
\end{teiElement}
            \begin{teiElement}[name={surname}]
Worms
\end{teiElement}
          }
          \begin{teiElement}[name={affiliation}]

            \teiOrgName{Université de Rouen Normandie, ÉRIAC – UR 4705, France}
          
\end{teiElement}
        }
        \teiAuthor[role={pbd}]{
          \teiPersName{
            \begin{teiElement}[name={forename}]
Clélia
\end{teiElement}
            \begin{teiElement}[name={surname}]
Zernik
\end{teiElement}
          }
          \begin{teiElement}[name={affiliation}]

            \teiOrgName{École nationale supérieure des Beaux-arts de Paris, France Université PSL, chaire Beauté·s, France }
          
\end{teiElement}
        }
      
\end{teiElement}
      \begin{teiElement}[name={editionStmt}]

        \begin{teiElement}[name={edition}]

          \teiDate{}
        
\end{teiElement}
        \begin{teiElement}[name={respStmt}]

          \begin{teiElement}[name={resp}]

\end{teiElement}
          \teiName{}
        
\end{teiElement}
        \begin{teiElement}[name={sponsor}]

\end{teiElement}
      
\end{teiElement}
      \begin{teiElement}[name={publicationStmt}]

        \begin{teiElement}[name={ab},type={expression}]

          \begin{teiBibl}[xmlid={bibl-001}]

            \teiPublisher{}
            \begin{teiElement}[name={address}]

              \begin{teiElement}[name={addrLine}]

\end{teiElement}
            
\end{teiElement}
            \begin{teiElement}[name={pubPlace}]

\end{teiElement}
            \begin{teiElement}[name={availability}]

              \begin{teiElement}[name={licence},target={-}]

\end{teiElement}
              \teiP[xmlid={p-001}]{}
            
\end{teiElement}
          
\end{teiBibl}
        
\end{teiElement}
        \begin{teiElement}[name={ab},type={book}]

          \begin{teiBibl}[xmlid={bibl-002}]

            \teiIdno[type={ISBN-13}]{979-10-240-1926-0}
            \teiIdno[type={ISSN}]{}
            \teiDate[type={publishing}]{}
            \teiDate[type={legal-deposit}]{}
            \teiDate[type={imprint}]{}
            \teiBiblScope[unit={page}]{}
            \begin{teiElement}[name={availability}]

              \begin{teiElement}[name={licence},target={-}]

\end{teiElement}
              \teiP[xmlid={p-002}]{}
            
\end{teiElement}
            \begin{teiElement}[name={distributor}]

\end{teiElement}
            \teiName[type={sales-agent}]{}
            \teiRef[type={site},target={-}]{}
          
\end{teiBibl}
          \teiName[type={printer}]{
            \teiName{}
            \teiNum{}
          }
          \begin{teiElement}[name={address}]

            \begin{teiElement}[name={addrLine}]

\end{teiElement}
          
\end{teiElement}
          \begin{teiElement}[name={dim},type={printingvolume}]

\end{teiElement}
          \begin{teiElement}[name={dimensions}]

            \begin{teiElement}[name={width},unit={mm}]

\end{teiElement}
            \begin{teiElement}[name={height},unit={mm}]

\end{teiElement}
            \begin{teiElement}[name={depth},unit={mm}]

\end{teiElement}
          
\end{teiElement}
          \begin{teiElement}[name={dim},type={weight},unit={gr}]

\end{teiElement}
          \begin{teiElement}[name={extent}]

\end{teiElement}
          \begin{teiElement}[name={measure},type={price},unit={EUR}]

\end{teiElement}
        
\end{teiElement}
        \begin{teiElement}[name={ab},subtype={PDF},type={digital\_download}]

          \begin{teiBibl}[xmlid={bibl-003}]

            \teiIdno[type={ISBN}]{}
            \teiIdno[subtype={\#\#},type={DOI}]{}
            \teiRef[type={DOI}]{}
            \teiRef[type={site},target={-}]{}
            \begin{teiElement}[name={distributor}]

\end{teiElement}
            \teiDate[type={publishing}]{}
            \begin{teiElement}[name={availability}]

              \begin{teiElement}[name={licence}]

\end{teiElement}
              \teiP[xmlid={p-003}]{}
            
\end{teiElement}
            \teiName[type={sales-agent}]{}
          
\end{teiBibl}
          \begin{teiElement}[name={dim},type={weight},unit={Mo}]

\end{teiElement}
          \begin{teiElement}[name={measure},type={price},unit={EUR}]

\end{teiElement}
        
\end{teiElement}
        \begin{teiElement}[name={ab},subtype={EPUB},type={digital\_download}]

          \begin{teiBibl}[xmlid={bibl-004}]

            \teiIdno[type={ISBN-13}]{}
            \teiIdno[type={DOI}]{}
            \teiRef[type={DOI}]{}
            \teiRef[type={site},target={-}]{}
            \begin{teiElement}[name={distributor}]

\end{teiElement}
            \teiDate[type={publishing}]{}
            \begin{teiElement}[name={availability}]

              \begin{teiElement}[name={licence}]

\end{teiElement}
              \teiP[xmlid={p-004}]{}
            
\end{teiElement}
            \teiName[type={sales-agent}]{}
          
\end{teiBibl}
          \begin{teiElement}[name={dim},extent={},type={weight},unit={Mo}]

\end{teiElement}
          \begin{teiElement}[name={measure},type={price},unit={EUR}]

\end{teiElement}
          \begin{teiElement}[name={desc},type={accessMode}]
textual
\end{teiElement}
          \begin{teiElement}[name={desc},type={accessibilityFeature}]
ARIA
\end{teiElement}
          \begin{teiElement}[name={desc},type={accessibilityFeature}]
highContrastDisplay
\end{teiElement}
          \begin{teiElement}[name={desc},type={accessibilityFeature}]
readingOrder
\end{teiElement}
          \begin{teiElement}[name={desc},type={accessibilityFeature}]
structuralNavigation
\end{teiElement}
          \begin{teiElement}[name={desc},type={accessibilityFeature}]
tableOfContents
\end{teiElement}
          \begin{teiElement}[name={desc},type={accessibilityFeature}]
unlocked
\end{teiElement}
          \begin{teiElement}[name={desc},subtype={soundHazard},type={accessibilityHazard}]
noSoundHazard
\end{teiElement}
          \begin{teiElement}[name={desc},subtype={flashingHazard},type={accessibilityHazard}]
noFlashingHazard
\end{teiElement}
          \begin{teiElement}[name={desc},subtype={motionSimulationHazard},type={accessibilityHazard}]
noMotionSimulationHazard
\end{teiElement}
          \begin{teiElement}[name={desc},type={accessModeSufficient}]
textual
\end{teiElement}
          \begin{teiElement}[name={desc},type={accessibilitySummary}]

\end{teiElement}
        
\end{teiElement}
        \begin{teiElement}[name={ab},subtype={HTML},type={digital\_online}]

          \begin{teiBibl}[xmlid={bibl-005}]

            \begin{teiElement}[name={distributor}]

\end{teiElement}
            \teiName[type={CMS}]{}
            \teiIdno[type={ISBN}]{}
            \teiIdno[type={DOI}]{}
            \teiRef[type={DOI}]{}
            \teiRef[type={site},internaltarget={}]{}
            \teiDate[type={publishing}]{}
            \teiDate[type={embargoend}]{}
            \begin{teiElement}[name={availability}]

              \begin{teiElement}[name={licence}]

\end{teiElement}
              \teiP[xmlid={p-005}]{}
            
\end{teiElement}
          
\end{teiBibl}
        
\end{teiElement}
      
\end{teiElement}
      \begin{teiElement}[name={sourceDesc}]

        \begin{teiBibl}[xmlid={bibl-006}]

\end{teiBibl}
        \begin{teiBibl}[type={edition1},xmlid={bibl-007}]

          \teiDate{}
          \begin{teiElement}[name={extent},rend={format}]

\end{teiElement}
          \teiPublisher{}
          \teiIdno[type={ISBN}]{}
          \begin{teiElement}[name={edition},rend={type}]

\end{teiElement}
        
\end{teiBibl}
      
\end{teiElement}
    
\end{teiElement}
    \begin{teiElement}[name={profileDesc}]

      \begin{teiElement}[name={langUsage}]

        \begin{teiElement}[name={language},ident={fr-FR}]

\end{teiElement}
      
\end{teiElement}
      \begin{teiElement}[name={textClass}]

        \begin{teiElement}[name={keywords},scheme={keywords},xmllang={fr}]

          \begin{teiList}

            \teiItem{}
            \teiItem{}
            \teiItem{}
          
\end{teiList}
        
\end{teiElement}
      
\end{teiElement}
      \begin{teiElement}[name={abstract},rend={4e-couv}]

        \teiP[xmlid={p-006}]{}
      
\end{teiElement}
      \begin{teiElement}[name={abstract},rend={resume},xmllang={fr}]

        \teiP[xmlid={p-007}]{}
      
\end{teiElement}
      \begin{teiElement}[name={abstract},rend={abstract},xmllang={fr}]

        \teiP[xmlid={p-008}]{}
      
\end{teiElement}
    
\end{teiElement}
    \begin{teiElement}[name={revisionDesc}]

      \begin{teiElement}[name={change},when={2025},who={\#\#}]

\end{teiElement}
    
\end{teiElement}
  
\end{teiElement}
  \begin{teiElement}[name={text},type={book},xmlid={book}]

    \begin{teiElement}[name={front}]

      \begin{teiDiv}[type={ack},xmlid={limin_01}]

        \teiP[xmlid={p-009}]{}
      
\end{teiDiv}
    
\end{teiElement}
    \begin{teiElement}[name={group},type={book},xmlid={book-001}]

      \begin{teiElement}[name={group},type={acknowledgments},data-page-title={Remerciements},data-include-href={Ch00\_Remerciements.xml},xmlid={acknowledgments-001}]

        \begin{teiElement}[name={front}]

          \begin{teiDiv}[type={titlePage},xmlid={titlepage-001}]

            \teiP[rend={title-main},xmlid={p-010}]{Remerciements}
          
\end{teiDiv}
        
\end{teiElement}
        \begin{teiElement}[name={body},xmlid={body-001}]

          \teiP[xmlid={p1}]{Nous tenons à remercier pour leur accueil et leur soutien, patients et bienveillants, toute l’équipe du Centre culturel international de Cerisy : Édith Heurgon et Dominique Peyrou, co-directeurs du CCIC ; Arnaud Chauvel, Patricia Doyère, Arthur Dumas, Nada Essid, Chantal Gosselin, Pascal Hédouin, Michaël Morel, Ghislaine Paysant et Jean-Christophe Tournière. }
          \teiP[xmlid={p2}]{Sans les compétences, l’inventivité et la vigilance constante de Justin Jaricot, coordinateur scientifique de la chaire Beauté·s, la conception et l’organisation du colloque n’auraient tout simplement pas été possibles : qu’il en soit ici amicalement remercié. }
          \teiP[xmlid={p3}]{Nous sommes également très reconnaissantes, pour leur soutien institutionnel et financier, à l’université PSL, à la chaire Beauté·s, à la fondation L’Oréal, à l’université de Rouen Normandie ainsi qu’aux Presses universitaires de Rouen et du Havre. }
          \teiP[xmlid={p4}]{Nous remercions la direction du festival Les Heures musicales de l’Abbaye de Lessay, pour le beau concert de l’ensemble La Tempête, et les champagnes Dom Perignon qui nous ont offert une dégustation mémorable !}
          \teiP[xmlid={p5}]{Merci aux musiciens Umadevi Nageswara Rao, chanteuse, et Nicolas Worms, compositeur et pianiste, pour la soirée musicale \teiHi[rend={italic}]{La Ligne et le cercle}.}
          \teiP[xmlid={p6}]{Merci à Keiko Courdy, cinéaste, pour la projection de son film \teiHi[rend={italic}]{L’Île invisible} et pour le film\teiHi[rend={italic}]{ Le Goût de la beauté}, qui capte si bien « l’esprit de Cerisy », tourné pendant le colloque avec la collaboration de Donatienne Berthereau que nous remercions aussi. }
          \teiP[xmlid={p7}]{Tous nos remerciements vont également à Thania Birem-Perez, chargée de communication et de pilotage pour la chaire Beauté·s, qui nous a beaucoup aidées pour l’édition de ce volume. }
          \teiP[xmlid={p8}]{Enfin, nous tenons à témoigner notre immense gratitude à chaque contributeur et contributrice, personnellement, pour leurs communications stimulantes et enthousiastes, sans oublier tous les auditeurs et auditrices de ces journées.}
        
\end{teiElement}
      
\end{teiElement}
      \begin{teiElement}[name={group},type={foreword},data-page-title={Avant-propos},data-page-authors={Anne-Lise Worms@@Université de Rouen Normandie, ÉRIAC – UR 4705, France||Clélia Zernik@@École nationale supérieure des Beaux-arts de Paris, France@@Université PSL, chaire Beauté·s, France},data-include-href={Ch01\_avantpropos.xml},xmlid={foreword-001}]

        \begin{teiElement}[name={front}]

          \begin{teiDiv}[type={titlePage},xmlid={titlepage-002}]

            \teiP[rend={title-main},xmlid={p-011}]{Avant-propos}
            \teiP[rend={author-aut},xmlid={p-012}]{Anne-Lise \teiHi[rend={small-caps}]{Worms} }
            \teiP[rend={authority\_affiliation},xmlid={p-013}]{Université de Rouen Normandie, ÉRIAC – UR 4705, France}
            \teiP[rend={author-aut},xmlid={p-014}]{Clélia \teiHi[rend={small-caps}]{Zernik}}
            \teiP[rend={authority\_affiliation},xmlid={p-015}]{École nationale supérieure des Beaux-arts de Paris, France }
            \teiP[rend={authority\_affiliation},xmlid={p-016}]{Université PSL, chaire Beauté·s, France}
          
\end{teiDiv}
        
\end{teiElement}
        \begin{teiElement}[name={body},xmlid={body-002}]

          \teiP[xmlid={p1-2}]{La notion de beauté s’inscrit dans le temps long de l’histoire : histoire des civilisations et des cultures, histoire des sociétés, histoire politique, histoire de la philosophie et histoire des idées, histoire des arts, qui n’est pas seulement celle des « beaux-arts », et histoire de la littérature. Et pourtant, aucune des « Histoires de la beauté » qui ont été publiées – que l’on songe, pour ne prendre que deux exemples, à celles de Umberto Eco\teiNote[xmlid={ftn1},n={1},place={foot}]{\teiP[xmlid={p-017}]{Umberto Eco (dir.), \teiHi[rend={italic}]{Histoire de la beauté}, Myriem Bouzaher (trad.), Paris, Flammarion, 2004.}} ou de Georges Vigarello\teiNote[xmlid={ftn6},n={2},place={foot}]{\teiP[xmlid={p-018}]{Georges Vigarello, \teiHi[rend={italic}]{Histoire de la beauté. Le corps et l’art d’embellir de la Renaissance à nos jours}, Paris, Éditions du Seuil, 2004.}} –, n’a jamais épuisé, ni même prétendu épuiser, l’infinie variété de ce qui est ou fut jugé comme « beau », ici et ailleurs, dans le monde entier. Ainsi, la notion de beauté serait intimidante, aliénante, parfaitement relative, historiquement et culturellement située, et donc impraticable, voire poussiéreuse. À quoi bon, pourquoi donc parler – encore – de « beauté » ? }
          \teiP[xmlid={p2-2}]{L’hypothèse que nous soutenons réaffirme à l’inverse les enjeux d’universalité, d’humanisme, de plaisirs et d’émotions qui sont spécifiques à la beauté. En nous appuyant sur une réinterprétation de la notion de beauté, nous pensons possible, et même nécessaire, de retrouver dans celle-ci une vitalité, une valeur, un pouvoir de transformation et d’orientation capables de nous interroger sur la manière de nous confronter aux enjeux d’aujourd’hui, car, loin d’être mortifère, la beauté peut être, selon nous, l’opérateur d’un futur commun, d’un monde habitable. }
          \teiP[xmlid={p3-2}]{La légende raconte que le mouvement écologique est né lorsque fut diffusée l’image iconique du « lever de Terre » prise par l’astronaute de la Nasa William Anders, le 24 décembre 1968, pendant la mission Apollo 8. C’est la beauté de notre planète, toute bleue et flottant dans l’espace infini, qui nous en a fait percevoir sa vulnérabilité et notre attachement – sa préciosité. Cette photographie, d’une beauté sidérante, nous fait soudain comprendre que notre planète n’est pas seulement une ressource dans laquelle piocher, mais une existence pour elle-même à laquelle nous tenons. Elle a été utilisée par les organismes et les associations militant pour la protection de l’environnement, au point d’être classée en 2003, par le magazine \teiHi[rend={italic}]{Life}\teiNote[xmlid={ftn7},n={3},place={foot}]{\teiP[xmlid={p-019}]{« 100 Photographs That Changed the World »,\teiHi[rend={italic}]{ Life}, n\teiHi[rend={sup}]{o} 312, 1\teiHi[rend={sup}]{er} août 2003.}}\teiHi[rend={italic}]{,} parmi les photos ayant le plus changé le monde. }
          \teiP[xmlid={p4-2}]{On a souvent caricaturé le désintéressement du jugement de goût kantien, en le réduisant à un rapport aux objets désincarné et abstrait. Mais Kant lui-même le dit\teiNote[xmlid={ftn8},n={4},place={foot}]{\teiP[xmlid={p-020}]{Emmanuel Kant, \teiHi[rend={italic}]{Critique de la faculté de juger} \lbrack{}1790\rbrack{}, Alain Renaut (trad.), Paris, Flammarion, 1995, § 41, « De l’intérêt empirique s’attachant au beau », p. 281 et suiv. }} : si le jugement sur la beauté doit être désintéressé, c’est-à-dire ne pas concerner notre intérêt propre, notre consommation, notre confort, nos conditions de vie, il demeure cependant \teiHi[rend={italic}]{intéressant}. Quand je dis « c’est beau », ce jugement me rend loquace et joyeux, j’échange avec mes amis et souhaite partager mon enthousiasme, enthousiasme qui porte sur mon intérêt pour la chose elle-même et non la chose telle qu’elle me concerne ou m’intéresse. En voyant ainsi pour la première fois la terre vue de loin, l’homme conçoit la terre non plus comme utile pour l’homme, mais comme belle pour elle-même, et il porte sur elle un jugement non plus intéressé, mais intéressant – qui a de l’intérêt pour la planète elle-même. C’est le début de l’écologie. L’époque contemporaine fourmille de ce type de retournements : de l’intérêt pour l’homme à un intérêt pour autre chose que l’homme, pour la vie elle-même, considérée non plus pour servir l’homme, mais comme valeur au-delà de lui. On peut donc voir au fondement de l’approche contemporaine et de ses prises de conscience – climatiques, décoloniales, politiques –, ce même retournement que la beauté nous enseigne : d’un jugement où le sujet est central à un jugement où le sujet découvre l’intérêt de la chose pour elle-même. Aussi l’ensemble des textes réunis ici trouve-t-il son unité dans ce mouvement de l’attention qui fait retour vers la vie, et qui découvre la beauté dans cet attachement pour la vie elle-même, comme le symbolise la planète Terre toute bleue, flottant dans l’immensité.}
          \teiP[xmlid={p5-2}]{Telle fut l’hypothèse à l’origine du colloque \teiHi[rend={italic}]{Beautés vitales. Pour une approche contemporaine de la beauté}, initié par la chaire Beauté·s de l’université PSL avec le soutien de L’Oréal et organisé en partenariat avec l’Université de Rouen Normandie, qui s’est tenu du 15 au 21 juillet 2022 au château de Cerisy. Cette rencontre avait pour but d’explorer la relation entre beauté et vie à travers la notion de « beauté vitale », entendue dans toutes ses dimensions : « vitale » au sens d’une beauté qui « s’ancre dans la vie », qui exprime et affirme\teiHi[rend={italic}]{ la vie}, mais « vitale » aussi au sens de « nécessaire pour la vie ». Notre approche est également « contemporaine » à plusieurs titres. Nous revendiquons en effet l’actualité de la notion de beauté, par opposition à ce qu’elle peut sembler avoir d’inactuel, notamment dans le domaine de l’art contemporain, tout en nous inscrivant dans le temps long de l’histoire de la beauté – ou des beautés –, afin de revisiter les conceptions héritées du passé et de les confronter aux enjeux de l’époque et du monde dans lequel nous vivons. }
          \teiP[xmlid={p6-2}]{C’est dans cette perspective que nous avons réuni, sans restriction en termes de discipline ou de méthodologie, chercheurs et artistes autour de tables rondes, communications, débats et propositions artistiques. Les approches relevaient de domaines de recherche volontairement variés (langues et cultures de l’Antiquité, philosophie, sociologie du soin, musique, sciences cognitives, architecture, sciences de l’innovation), accompagnées d’expériences esthétiques participatives émouvantes : concert de Nicolas Worms et Uma Dewi, au château, autour de la musique indienne, concert à l’abbaye de Lessay, tournage du film de Keiko Courdy. Le rapport d’étonnement des doctorants à la fin du séjour, composé d’approches artistiques sensibles, parfois drôles\teiHi[rend={italic}]{,} fut la preuve de l’intensité et de la fécondité des échanges de la semaine. C’est dans ce contexte que la notion de beauté a pu retrouver une centralité et une dimension prospective qui ont semblé essentielles à tous les participants. }
          \teiP[xmlid={p7-2}]{Les contributions que nous avons réunies pour ce volume sont issues des communications présentées au cours du colloque, dont on trouvera le programme exhaustif à la fin de l’ouvrage. Si le volume retrace le déroulement de cette semaine à Cerisy, il ne le suit pas non plus strictement. Il partage cependant cette intuition forte que la beauté est d’abord une expérience vivante, et nous retrouverons ici, comme entre parenthèses, des moments d’écoute musicale ou de suspension contemplative, des promenades dans des jardins exotiques, etc. Nous avons voulu laisser les auteurs relativement libres de la forme qu’ils souhaitaient donner à leur texte, tout en respectant, bien sûr, les consignes de l’éditeur. L’ensemble de l’ouvrage reflète ainsi la variété de tons et de voix des interventions et, d’une certaine manière, redit le caractère bigarré, chatoyant de la notion de beauté qui, sans cadre ni norme, nous relie au monde avec toujours la même émotion.}
          \teiP[xmlid={p8-2}]{La première partie du volume s’attache aux anciens et nouveaux discours. Contrairement à ce que l’on pourrait croire, les discours sur la beauté dans l’Antiquité ne tendent pas nécessairement à dissocier beauté et vie. De fait, si, pour Platon, la beauté est une Idée au-delà du monde sensible, elle est cependant intimement liée à la vie, par l’intermédiaire du désir amoureux. C’est ce que montre Plotin, qui voit dans la vie l’expression même de la beauté (Anne-Lise Worms). En contexte romain, la beauté peut également aller au-delà des caractérisations physiques pour devenir le propre de l’homme sage (Catherine Baroin). Dans les discours contemporains, c’est dans une même attention extrême à la vie que l’approche esthétique trouve sa plus grande pertinence (Rémi Mermet).}
          \teiP[xmlid={p9}]{Dans la deuxième partie de l’ouvrage, la beauté, rapatriée dans le champ du vivant, s’exprime à toutes les échelles : celles de l’évolution des espèces (Claudine Cohen), de l’animal (Bertrand Prévost), de la bestiole ou de la mauvaise graine (Claudine Sagaert) en passant par le monde minéral (Didier Nectoux). C’est cette même inscription de la beauté dans le vivant que l’on retrouve dans la nature travaillée par le jardin japonais (Philippe Bonnin).}
          \teiP[xmlid={p10}]{De même, dans les arts – et c’est ce que tend à montrer notre troisième partie – c’est un paradigme mobile, plutôt que figé et surplombant, qui domine : un paradigme de la relation (Chiara Vecchiarelli), un dépassement de l’instant figé de la photographie (Carole Maigné) vers un modèle musical qui retient toute la capacité d’improvisation de la vie (Pierre Sauvanet) ou de son caractère socialement situé (Guilhem Marion).}
          \teiP[xmlid={p11}]{Enfin, si la beauté est aussi essentielle à la vie, c’est qu’elle nous permet de traverser les épreuves intimes et de faire société (Gisèle Dambuyant) ou de surmonter les catastrophes (Clélia Zernik). La beauté est vitale en tant qu’elle est nécessaire à la vie et, réciproquement, elle nous fait voir ce qu’est véritablement la vie (Frédéric Worms). Ne pourrait-on pas, dès lors, soutenir qu’aujourd’hui encore « la beauté sauvera le monde\teiNote[xmlid={ftn9},n={5},place={foot}]{\teiP[xmlid={p-021}]{Fiodor Dostoïevski, \teiHi[rend={italic}]{L’Idio}t, André Markowicz (trad.), Arles, Actes Sud, « Babel », 1993, t. II, p. 102.}} » ?}
          \teiP[xmlid={p12}]{Nous avons regroupé dans le dernier chapitre de courts extraits des « rapports d’étonnement » de doctorants qui ont assisté au colloque ainsi que des images du film que la réalisatrice Keiko Courdy a tourné à Cerisy pendant cette même semaine.}
        
\end{teiElement}
      
\end{teiElement}
      \begin{teiElement}[name={group},type={section1},xmlid={section1-001}]

        \teiHead[xmlid={beaute-et-vie-anciens-et-nouveaux-discours-001}]{Beauté et vie : anciens et nouveaux discours}
        \begin{teiElement}[name={group},type={chapter},data-page-title={Plotin contre Platon ? Chercher la beauté dans le vivant et dans la vie},data-page-authors={Anne-Lise Worms@@Université de Rouen Normandie, ÉRIAC – UR 4705, France},data-include-href={Ch02\_Plotin\_contre\_Platon\_Worms.xml},xmlid={chapter-001}]

          \begin{teiElement}[name={front}]

            \begin{teiDiv}[type={titlePage},xmlid={titlepage-003}]

              \teiP[rend={title-main},xmlid={p-022}]{Plotin contre Platon ? Chercher la beauté dans le vivant et dans la vie}
              \teiP[rend={author-aut},xmlid={p-023}]{Anne-Lise \teiHi[rend={small-caps}]{Worms}}
              \teiP[rend={authority\_affiliation},xmlid={p-024}]{Université de Rouen Normandie, ÉRIAC – UR 4705, France}
            
\end{teiDiv}
            \begin{teiElement}[name={epigraph}]

              \begin{teiCit}

                \begin{teiQuote}
J’ai toujours eu dans l’esprit, sans bien m’en rendre compte, une sorte de balance. Sur un plateau il y avait la douleur, la mort, sur l’autre la beauté de la vie. Le premier portait un poids toujours plus lourd, le second, presque rien que d’impondérable. Mais il m’arrivait de croire que l’impondérable pût l’emporter, par moments.
\end{teiQuote}
                \begin{teiBibl}[xmlid={bibl001}]
Philippe Jaccottet,\teiHi[rend={italic}]{ À travers un verger}\teiNote[xmlid={ftn1-2},n={1},place={foot}]{\teiP[xmlid={p-025}]{Philippe Jaccottet, \teiHi[rend={italic}]{À travers un verger} \lbrack{}1975\rbrack{}, Paris, Gallimard, « Bibliothèque de la Pléiade », 2014, t. II, p. 558. }}.
\end{teiBibl}
              
\end{teiCit}
            
\end{teiElement}
          
\end{teiElement}
          \begin{teiElement}[name={body},xmlid={body-003}]

            \teiP[xmlid={p1-3}]{L’expression « beauté grecque » renvoie généralement à des normes et à des critères de beauté d’ordre mathématique : symétrie (en grec \teiHi[rend={italic}]{summetria}, « bonnes proportions »), mesure, harmonie, règle (\teiHi[rend={italic}]{kanôn}) sont les termes qui viennent spontanément à l’esprit. De fait, nombreux sont les textes, philosophiques ou non, mais aussi les œuvres d’art de l’antiquité grecque et romaine, et leurs multiples réceptions, qui défendent ou expriment cette idée. Et c’est bien souvent en nous référant à ces critères que nous disons de telle personne ou de telle sculpture qu’elles sont belles. }
            \teiP[xmlid={p2-3}]{Nous voulons cependant interroger cette représentation, elle-même normée, de « la » beauté grecque ou de la beauté \teiHi[rend={italic}]{pour} les Grecs. Il m’est impossible d’évoquer ici la variété des conceptions et des représentations de la beauté en Grèce ancienne. Nous avons donc choisi de prendre un exemple, celui de Plotin, philosophe du \teiHi[rend={small-caps}]{iii}\teiHi[rend={sup}]{e} siècle après J.-C., qui se revendiquait comme simple « exégète » de Platon, mais dont l’interprétation singulière et originale des écrits de son maître a conduit les historiens de la philosophie à le qualifier de « néoplatonicien ». Les textes sur lesquels nous nous appuierons nous paraissent en effet significatifs en ce qu’ils révèlent sinon une rupture, du moins une différence réelle et singulière entre deux conceptions de la beauté. La première est largement admise dans la philosophie grecque ancienne. Elle se caractérise par les notions évoquées à l’instant – harmonie, \teiHi[rend={italic}]{summétria}, mesure – et renvoie à une réalité statique, qui s’opposerait au mouvement et aurait la stabilité, voire l’immobilité du principe et de la raison. Elle s’exprime tout particulièrement dans le terme grec de \teiHi[rend={italic}]{kosmos}, qui a le sens à la fois d’ordre et de beauté. À l’inverse, la seconde conception de la beauté active des notions différentes, non pas celles de désordre, d’absence de mesure ou de symétrie, mais, précisément, celles de mouvement ou encore de lumière : la beauté serait avant tout ce qui, \teiHi[rend={italic}]{au-delà de} la symétrie, chatoie, brille, et qui seul peut susciter en retour, chez celui qui la contemple, un élan, une émotion, c’est-à-dire le désir, l’amour (\teiHi[rend={italic}]{érôs}). La beauté y est donc étroitement associée à ce qui est de l’ordre de la vie, du vivant, car elle est, on le verra, une forme vivante. Cette expression constitue presque un oxymore : à la stabilité et l’immobilité de la forme s’opposent l’instabilité et la mobilité de la vie. C’est la tension entre ces deux conceptions que nous voudrions ici explorer à travers la lecture de Plotin.}
            \teiP[xmlid={p3-3}]{Le premier texte est un passage du \teiHi[rend={italic}]{Traité 1 }(I, 6), qui a reçu pour titre « Du beau\teiNote[xmlid={ftn41},n={2},place={foot}]{\teiP[xmlid={p-026}]{Lorsque nous citons un traité pour la première fois, nous indiquons d’abord son rang dans l’ordre chronologique de rédaction des traités par Plotin, ordre que nous indique Porphyre aux chapitres 4 à 6 de sa \teiHi[rend={italic}]{Vie de Plotin}. Puis nous mentionnons entre parenthèses sa référence dans l’édition établie par Porphyre lui-même de l’œuvre de son maître : les \teiHi[rend={italic}]{Ennéades}. Par la suite, nous nous référons à chaque traité en indiquant son numéro dans l’ordre chronologique, puis la référence du chapitre et enfin celle des lignes. Pour la traduction, nous nous fondons sur l’édition de référence des écrits de Plotin, publiée par Paul Henry et Hans-Rudolf Schwyzer : \teiHi[rend={italic}]{Plotini Opera}. \teiHi[rend={italic}]{Editio minor}, Oxford, Oxford Clarendon Press, t. 1, 1964, t. 2, 1977, t. 3, 1982. Précisons également que les titres ne sont pas de Plotin. Ils nous ont été transmis par Porphyre, qui indique ceux « qui ont prévalu » (Porphyre, \teiHi[rend={italic}]{Vie de Plotin}, chap. 4, l. 19).}} », et dans lequel Plotin cherche à apporter une réponse à cette question fondamentale : « D’où vient le beau ? » }
            \teiP[xmlid={p4-3}]{Au premier chapitre du traité, il énumère différentes sortes de beauté, en se référant de façon implicite à Platon : il y a le beau « qui existe surtout dans la vue », mais aussi « dans les objets perçus par l’ouïe et selon l’arrangement des mots », ou encore « dans la musique, et la musique tout entière ». Ce sont là des beautés sensibles. Mais le beau existe également dans ce qui se situe au-delà de la sensation : les « belles occupations, dispositions et sciences, ainsi que la beauté des vertus ». Puis il ajoute, de façon quelque peu sibylline : « Et s’il existe quelque chose encore avant tout ceci, cela se montrera de soi\teiNote[xmlid={ftn42},n={3},place={foot}]{\teiP[xmlid={p-027}]{\teiHi[rend={italic}]{Traité 1}, chap. 1, l. 1-6. Sauf mention contraire, les traductions sont les nôtres (ici : Plotin, \teiHi[rend={italic}]{Traité 1}, introduction, traduction, commentaire et notes par Anne-Lise Darras-Worms, Paris, Éditions du Cerf, 2007). }}. » L’influence, ici, d’une lecture conjointe du \teiHi[rend={italic}]{Banquet} et d’\teiHi[rend={italic}]{Hippias Majeur}, se traduit notamment par le choix des différentes sortes de beauté et par leur présentation selon un schéma ascendant. On se rappelle en effet que Socrate, dans \teiHi[rend={italic}]{Hippias Majeur}, dit des plaisirs « qui nous viennent de l’ouïe et de la vue » que « c’est cela qui est beau\teiNote[xmlid={ftn43},n={4},place={foot}]{\teiP[xmlid={p-028}]{\teiHi[rend={italic}]{Hippias Majeur}, 297e-298a ; voir aussi \teiHi[rend={italic}]{Philèbe}, 51b-e.}} ». Dans le \teiHi[rend={italic}]{Banquet}, Diotime conseille à qui veut s’initier aux mystères de l’amour de passer de la beauté des corps à celle de l’âme et de « contempler ce qui est beau dans les occupations, les règles », puis dans « les sciences » pour enfin parvenir à la contemplation de la Beauté en soi, terme ultime de l’ascension\teiNote[xmlid={ftn44},n={5},place={foot}]{\teiP[xmlid={p-029}]{Voir \teiHi[rend={italic}]{Banquet}, 210b-d.}}. La reprise, partielle, des définitions platoniciennes conduit Plotin à s’interroger sur \teiHi[rend={italic}]{l’origine commune} de ces beautés et, plus précisément, sur ce qui fait que nous \teiHi[rend={italic}]{percevons} tel ou tel objet comme beau, en éprouvons du plaisir et \teiHi[rend={italic}]{disons} que cet objet est beau. L’on s’attendrait à ce qu’il explique alors, conformément à la thèse du \teiHi[rend={italic}]{Banquet}, que la parenté entre les beautés sensibles et non sensibles résulte de leur commune participation à la Forme en soi du Beau. Plotin souscrit bien sûr à cette thèse platonicienne fondamentale, qu’il reprend à son compte dans d’autres chapitres de ce traité-ci comme dans d’autres traités, mais il va plus loin et s’en démarque dans une certaine mesure. D’une part, il montre que c’est leur participation non pas à la seule Forme du Beau qui rend les choses belles, mais à la Forme en soi (ou Idée, au sens platonicien du terme), qui est la Beauté même. D’autre part, il affirme que les émotions et le plaisir que suscite en nous toute beauté, quelle qu’elle soit, provient d’autre chose encore, qui est lié à la nature même de la Forme. Et il insiste : « Qu’est-ce donc cela qui est présent aux corps ? \lbrack{}…\rbrack{} Qu’est-ce donc, ce qui meut les regards de ceux qui contemplent, les tourne vers soi, les attire et fait qu’ils se réjouissent de leur contemplation\teiNote[xmlid={ftn45},n={6},place={foot}]{\teiP[xmlid={p-030}]{\teiHi[rend={italic}]{Traité 1}, chap. 1, lignes 16-19. }} ? » Notons, dès à présent, deux notions importantes, celles de « présence » et de « mouvement ». Nous y reviendrons.}
            \teiP[xmlid={p5-3}]{En réponse aux questions qu’il vient de formuler, Plotin évoque une explication générale et consensuelle, car admise, selon lui, par « tous, à peu près » : « De fait, il est dit chez tous, à peu près, que c’est la bonne proportion des parties les unes par rapport aux autres et par rapport au tout ainsi que l’ajout de la belle couleur qui produisent la beauté pour la vue et, pour ces objets comme pour tous les autres en général, être beau, c’est être bien proportionné et mesuré\teiNote[xmlid={ftn46},n={7},place={foot}]{\teiP[xmlid={p-031}]{\teiHi[rend={italic}]{Traité 1}, chap. 1, lignes 20-26. }}. »}
            \teiP[xmlid={p6-3}]{Cette conception est, de fait, largement partagée non seulement par les artistes eux-mêmes, mais aussi par bien des philosophes antiques, ce qui explique pourquoi la « beauté grecque » ou la beauté pour les Grecs est bien souvent identifiée à cette seule définition. On peut citer ici le fameux \teiHi[rend={italic}]{Canon} de Polyclète, célèbre sculpteur grec du \teiHi[rend={small-caps}]{v}\teiHi[rend={sup}]{e} siècle avant J.-C., qui désigne à la fois un traité, dans lequel sont définis les critères déterminants de la beauté d’une œuvre d’art, et une statue réalisée par Polyclète lui-même, dont nous possédons un témoignage grâce au Doryphore, copie d’époque romaine. Chez les philosophes, la théorie de la \teiHi[rend={italic}]{summetria} est soutenue majoritairement par les stoïciens, qui l’appliquent à l’art comme au corps, à l’âme, aux sciences ou au monde lui-même\teiNote[xmlid={ftn47},n={8},place={foot}]{\teiP[xmlid={p-032}]{Voir, par exemple, \teiHi[rend={italic}]{Stoicorum Veterum Fragmenta}, Hans Friedrich von Arnim (éd.), Leipzig, Teubner, 1903, vol. III, § 472. Cette thèse est reprise par Cicéron dans les \teiHi[rend={italic}]{Tusculanes}, IV, 31. }}. Mais on la trouve aussi chez Platon, dans le \teiHi[rend={italic}]{Sophiste}, le \teiHi[rend={italic}]{Philèbe}, le\teiHi[rend={italic}]{ Timée }et les\teiHi[rend={italic}]{ Lois}\teiNote[xmlid={ftn48},n={9},place={foot}]{\teiP[xmlid={p-033}]{Platon, \teiHi[rend={italic}]{Sophiste}, 235d-e ; \teiHi[rend={italic}]{Philèbe}, 63e-64e ; \teiHi[rend={italic}]{Timée}, 35b-37 a et 87c-90e ; \teiHi[rend={italic}]{Lois}, Livre II, 668d-e. }} ou encore chez Aristote, pour qui « les formes les plus hautes du beau » sont « l’ordre, la bonne proportion, le défini\teiNote[xmlid={ftn49},n={10},place={foot}]{\teiP[xmlid={p-034}]{Aristote, \teiHi[rend={italic}]{Métaphysique}, Livre VII, chap. 3, 1078a-b. Voir également \teiHi[rend={italic}]{Physique}, Livre VII, chap. 3, 246b et \teiHi[rend={italic}]{Topiques}, VI, 2, 139b.}} ».}
            \teiP[xmlid={p7-3}]{Or, si cette explication n’est certes pas erronée pour Plotin, qui l’assume lui-même dans d’autres traités\teiNote[xmlid={ftn50},n={11},place={foot}]{\teiP[xmlid={p-035}]{Ainsi dans le \teiHi[rend={italic}]{Traité 47} (III, 2), chap. 3, lignes 6-13 et chap. 17, ou dans le \teiHi[rend={italic}]{Traité 5} (V, 9), chap. 11, lignes 11-15. }}, elle est, selon lui, très insuffisante. En effet, elle ne peut s’appliquer qu’à un tout composé de plusieurs parties, mais ne rend nullement compte de la beauté de ce qui est « simple » (\teiHi[rend={italic}]{haplous}) : « pour eux, rien de simple, seulement ce qui est composé, sera beau\teiNote[xmlid={ftn51},n={12},place={foot}]{\teiP[xmlid={p-036}]{\teiHi[rend={italic}]{Traité 1}, chap. 1, lignes 25 et 26. }} » Comment, dans ces conditions, justifier la beauté d’une chose simple ? D’où provient la beauté de chaque partie dans un tout, des « belles couleurs » elles-mêmes, de la lumière du soleil, de l’or, de l’éclair, des astres dans la nuit, d’un son seul ou de chaque son considéré en lui-même séparément dans une composition musicale\teiNote[xmlid={ftn52},n={13},place={foot}]{\teiP[xmlid={p-037}]{\teiHi[rend={italic}]{Ibid}., lignes 30-36. }} ? Il en va de même pour la beauté des choses non sensibles : les belles occupations, les belles lois, les belles connaissances, les belles sciences et, au-delà, la beauté de l’âme en elle-même ou celle de l’Intellect lui-même, « lui qui est isolé et un\teiNote[xmlid={ftn53},n={14},place={foot}]{\teiP[xmlid={p-038}]{\teiHi[rend={italic}]{Ibid}., lignes 40-44 ; 49 et 50 ; 53 et 54.}} ». La cause unique de toutes ces beautés ne peut, par conséquent, être réduite à la seule \teiHi[rend={italic}]{summetria}. En d’autres termes, elle ne peut être uniquement rationnelle. }
            \teiP[xmlid={p8-3}]{La preuve que, lorsque nous sommes face à la beauté, ce n’est pas à elle seule la participation à la Forme qui nous émeut, réside dans la nature même du sentiment qui nous saisit alors : le sentiment amoureux. Dans les chapitres suivants du traité, Plotin s’attache à décrire la spécificité et la puissance de ce sentiment en utilisant des termes empruntés à Platon, notamment dans le \teiHi[rend={italic}]{Banquet} et le \teiHi[rend={italic}]{Phèdre} : « stupeur, étonnement délicieux, désir, amour et effroi accompagné de plaisir\teiNote[xmlid={ftn54},n={15},place={foot}]{\teiP[xmlid={p-039}]{\teiHi[rend={italic}]{Ibid.}, chap. 4, lignes 16 et 17.}} ». Cet amour que nous éprouvons à l’égard des objets sensibles, nous l’éprouvons à un degré supérieur à l’égard de l’âme, l’âme d’un autre ou notre propre âme, lorsqu’elle est séparée du corps, purifiée de ses passions par la pratique des vertus et que, sur elle, « répand sa lumière l’Intellect de forme divine\teiNote[xmlid={ftn55},n={16},place={foot}]{\teiP[xmlid={p-040}]{\teiHi[rend={italic}]{Ibid.}, chap. 5, lignes 16 et 17.}} ». Selon Plotin, c’est avant tout cette lumière qui attire l’âme, mais l’Intellect lui-même la tient d’ailleurs : cette lumière provient de ce qui est \teiHi[rend={italic}]{au-delà} de lui, Beauté au-delà de la Beauté, Principe premier, c’est-à-dire le Bien. Plotin définit le Bien de plusieurs façons dans les trois derniers chapitres du traité. L’une d’entre elles est particulièrement intéressante ici : le Bien est « sans mélange, simple, pur, lui dont tout dépend et vers quoi tout regarde et existe et vit et pense, car il est cause de la vie, de l’Intellect et de l’être\teiNote[xmlid={ftn56},n={17},place={foot}]{\teiP[xmlid={p-041}]{\teiHi[rend={italic}]{Ibid.}, chap. 7, lignes 10 et 11.}} ». Nous reviendrons sur le sens de cette définition, qui renvoie à un point de doctrine important dans l’histoire du platonisme – la triade être, vie, pensée – et permet de mieux comprendre le lien entre vie et beauté. Mais notons d’emblée, à propos du premier chapitre du traité, que ce qui « meut les regards de ceux qui contemplent » et induit l’élan, le désir amoureux est donc la présence, certes plus ou moins intense selon la nature de l’objet contemplé, mais réelle, du Bien. Dans ces conditions, l’enjeu majeur de l’objection adressée par Plotin à ceux qui réduisent la beauté à la seule symétrie et les raisons pour lesquelles il insiste à ce point sur la beauté de ce qui est simple peuvent être formulés ainsi : une belle couleur, un éclair, un son unique, une vertu tiennent leur beauté précisément de leur nature même de choses simples et non composées – simples comme l’est, à un niveau supérieur, \teiHi[rend={italic}]{le }simple, le Bien ou l’Un\teiNote[xmlid={ftn57},n={18},place={foot}]{\teiP[xmlid={p-042}]{Trois « niveaux de réalité », ou hypostases, forment la structure de la métaphysique plotinienne : l’Âme, l’Intellect et l’Un ou le Bien. L’Un-Bien, que Plotin identifie au Bien-soleil de la \teiHi[rend={italic}]{République} (508-509), est « au-delà de » l’Intellect. }}. }
            \teiP[xmlid={p9-2}]{Lisons maintenant, en regard du \teiHi[rend={italic}]{Traité 1}, ce passage du \teiHi[rend={italic}]{Traité 38} (VI, 7) dans lequel Plotin reprend sa thèse selon laquelle la beauté ne réside pas, ou pas seulement, dans la \teiHi[rend={italic}]{summetria} : }
            \begin{teiCit}[xmlid={cit01}]

              \begin{teiQuote}
L’âme s’élève au-dessus de l’Intellect, mais elle ne peut courir au-delà du Bien, parce qu’il n’est rien qui soit au-delà de lui. Et si elle reste dans l’Intellect, elle contemple des objets beaux, certes, et vénérables, mais elle ne possède pas encore entièrement ce qu’elle cherche. C’est, en effet, comme si elle s’approchait d’un visage beau, certes, mais qui ne peut encore mouvoir les regards, car en lui ne se distingue pas la grâce (\teiHi[rend={italic}]{charis}) qui court sur la beauté. Il faut donc dire qu’ici-bas aussi, la beauté se trouve dans la lumière qui brille au-dessus de la symétrie, plutôt que dans la symétrie, et que c’est cela qui suscite l’amour. Pourquoi donc, en effet, est-ce sur un visage vivant que resplendit davantage la lumière de la beauté, alors que sur un visage mort, il n’en reste qu’un vestige, même si ce visage n’est pas encore flétri dans sa chair et dans sa symétrie ? Et parmi les statues, les plus vivantes ne sont-elles pas les plus belles, même si les autres ont plus de symétrie ? Et un homme laid, s’il est vivant, n’est-il pas plus beau que l’homme beau représenté dans une statue ? N’est-ce pas parce que cet homme-là est plus désirable ? Mais cela, c’est parce qu’il a une âme ; mais cela, c’est parce qu’il a davantage la forme du Bien ; mais cela, c’est parce que son âme a pris la couleur de la lumière du Bien et qu’ainsi colorée, elle s’est éveillée, elle s’est élevée, toute légère, et qu’elle élève et rend léger ce qui lui appartient, et, autant que possible, le rend bon et l’éveille\teiNote[xmlid={ftn58},n={19},place={foot}]{\teiP[xmlid={p-043}]{Plotin, \teiHi[rend={italic}]{Traité 38} (VI, 7), chap. 22, lignes 24-36, traduction de Pierre Hadot légèrement modifiée (voir note suivante).}}.
\end{teiQuote}
            
\end{teiCit}
            \teiP[xmlid={p10-2}]{Le \teiHi[rend={italic}]{Traité 38} a reçu pour titre « Comment la multiplicité des Idées a été produite et sur le Bien ». Cependant, comme le montre Pierre Hadot dans son commentaire, « \lbrack{}ce\rbrack{} traité est un traité sur le Bien. Mais c’est aussi une description de l’itinéraire de l’âme vers le Bien ». Ceci, pour Pierre Hadot, « n’est pas étonnant \lbrack{}…\rbrack{}, car selon la \teiHi[rend={italic}]{République} (509d), le Bien, c’est ce que toute âme poursuit et \lbrack{}…\rbrack{}, en un certain sens, s’il n’y a pas de définition abstraite du Bien, il y a en quelque sorte une définition concrète : le Bien est ce “toujours au-delà” vers quoi s’élève l’âme, ce “toujours au-delà” qu’elle désire\teiNote[xmlid={ftn59},n={20},place={foot}]{\teiP[xmlid={p-044}]{Plotin, \teiHi[rend={italic}]{Traité 38}, introduction, traduction, commentaire et notes par Pierre Hadot, Paris, Éditions du Cerf, p. 50.}} ». }
            \teiP[xmlid={p11-2}]{L’itinéraire de l’âme vers le Bien constitue une expérience de contemplation, qui peut également prendre la forme d’un « exercice spirituel\teiNote[xmlid={ftn60},n={21},place={foot}]{\teiP[xmlid={p-045}]{Nous empruntons cette expression à Pierre Hadot. Voir notamment son ouvrage, \teiHi[rend={italic}]{Exercices spirituels et philosophie antique}, Paris, Albin Michel, 2002. }} » et comporte trois étapes. Dans un premier temps, lorsqu’elle perçoit la beauté des choses sensibles, l’âme comprend que cette beauté provient des Formes, ou Idées, qui composent le monde intelligible et ne font qu’un avec l’Intellect. Ensuite, dans un mouvement de conversion, elle se tourne alors et remonte vers l’Intellect qui est, pour Plotin, la Beauté en soi, non pas la seule Forme du Beau mais la Beauté même, mais qui est \teiHi[rend={italic}]{aussi} la Vie\teiNote[xmlid={ftn61},n={22},place={foot}]{\teiP[xmlid={p-046}]{On utilise généralement une majuscule pour désigner la Vie quand cette notion s’applique à l’Intellect ou au Bien-Un. }} en tant qu’émanation du Bien ou, plus précisément, l’\teiHi[rend={italic}]{énergéia} (l’énergie\teiNote[xmlid={ftn62},n={23},place={foot}]{\teiP[xmlid={p-047}]{Le terme grec \teiHi[rend={italic}]{énergéia} a plusieurs sens chez Plotin, selon le niveau de réalité auquel elle s’applique : « activité » ou « actuation » d’une virtualité dans et par l’Intellect, « émanation » pour ce qui concerne le Bien. Voir Pierre Hadot, \teiHi[rend={italic}]{Traité 38}, \teiHi[rend={italic}]{op. cit.}, p. 395. }}) du Bien. Plotin dit ainsi : « Cet Intellect, cette Vie, sont “semblables au Bien” (\teiHi[rend={italic}]{agathoeidès}) et on les désire aussi dans la mesure où ils sont “semblables au Bien”. Et si je les appelle “semblables au Bien”, c’est parce que la Vie est “l’énergie” du Bien, ou plutôt “l’énergie” qui émane du Bien et que l’Intellect est cette même “énergie”, après qu’elle a été délimitée\teiNote[xmlid={ftn63},n={24},place={foot}]{\teiP[xmlid={p-048}]{\teiHi[rend={italic}]{Traité 38}, chap. 21, lignes 4-6.}}. » Cependant, l’amour que l’âme éprouve pour la Vie, \teiHi[rend={italic}]{c’est-à-dire} l’Intellect et la Beauté, ne lui suffit pas, car elle ne peut se contenter de ce qui est « semblable au Bien ». En effet, son désir la pousse pour ainsi dire naturellement à rechercher \teiHi[rend={italic}]{l’origine}, la « source », dirait Plotin, de cette Beauté, car ce qui attise ce désir est, en réalité, toujours « quelque chose d’autre », qui « s’ajoute à eux, venant d’en-haut\teiNote[xmlid={ftn64},n={25},place={foot}]{\teiP[xmlid={p-049}]{\teiHi[rend={italic}]{Traité 38}, chap. 21, lignes 13 et 14.}} ». Cette « autre chose » est précisément le Bien, qui est Beauté \teiHi[rend={italic}]{au-delà de} l’Intellect-Beauté-Vie et dont la lumière « court sur eux\teiNote[xmlid={ftn65},n={26},place={foot}]{\teiP[xmlid={p-050}]{\teiHi[rend={italic}]{Traité 38}, chap. 22, ligne 3.}} ». La remontée vers le Bien constitue, pour l’âme, la troisième et dernière étape de l’ascension. }
            \teiP[xmlid={p12-2}]{Avant de lire de plus près ce passage du \teiHi[rend={italic}]{Traité 38}, revenons tout d’abord sur la définition de l’Intellect comme Vie. Comme l’indique Pierre Hadot dans son article sur « Être, vie, pensée chez Plotin et avant Plotin\teiNote[xmlid={ftn66},n={27},place={foot}]{\teiP[xmlid={p-051}]{Pierre Hadot, « Être, vie, pensée chez Plotin et avant Plotin », dans\teiHi[rend={italic}]{ Entretiens sur l’Antiquité classique}, t. V, \teiHi[rend={italic}]{Les sources de Plotin}, Vandœuvres-Genève, Fondation Hardt, 1960, p. 105-142. }} », auquel nous renvoyons pour davantage de précisions, « la triade \lbrack{}être-vie-pensée\rbrack{} se situe au plan de la seconde hypostase (l’Intellect), de ce monde intelligible qu’est l’Intelligence <ou Intellect>. Elle définit en quelque sorte l’Intelligence, révèle sa structure et sa genèse\teiNote[xmlid={ftn67},n={28},place={foot}]{\teiP[xmlid={p-052}]{\teiHi[rend={italic}]{Ibid}., p. 130.}} ». Or, au sein de cette triade, « la vie précède les deux autres termes\teiNote[xmlid={ftn68},n={29},place={foot}]{\teiP[xmlid={p-053}]{\teiHi[rend={italic}]{Ibid.}, p. 139.}} », car elle « sort immédiatement de l’Un et c’est à partir d’elle que l’Intelligence se constitue\teiNote[xmlid={ftn69},n={30},place={foot}]{\teiP[xmlid={p-054}]{\teiHi[rend={italic}]{Ibid.}, p. 133. }} ». }
            \teiP[xmlid={p13}]{Le passage du \teiHi[rend={italic}]{Traité 38} que nous allons lire maintenant met en évidence la prééminence de la vie et explicite également la relation qui peut être établie entre vie et beauté. }
            \teiP[xmlid={p14}]{Plotin procède ici par comparaisons. La première est celle de l’Intellect avec un visage beau. On se rappelle que, dans le \teiHi[rend={italic}]{Traité 1}, il demandait : « Qu’est-ce donc cela qui est présent aux corps ? \lbrack{}…\rbrack{} Qu’est-ce donc, ce qui meut les regards de ceux qui contemplent ? » La réponse est donnée ici, dans le \teiHi[rend={italic}]{Traité 38}. Si un visage beau peut « mouvoir les regards » et « susciter l’amour », c’est parce que l’âme sent en lui la \teiHi[rend={italic}]{présence} du Bien, qui se manifeste par la « grâce » ou la « lumière » du Bien qui resplendissent sur lui, ces deux notions étant ici posées comme équivalentes par Plotin. C’est à ces conditions que l’âme peut dire que ce visage, ou tout autre objet sensible, est beau. L’emploi du verbe « courir » est significatif : il implique le mouvement et, en ce sens, il est le signe de la « vie », cette vie dont a vu que Plotin la définit comme \teiHi[rend={italic}]{l’énergéia} du Bien : vie de l’Intellect en tant qu’hypostase, tout d’abord, dont procède l’Âme hypostase, vie de l’intellect individuel dans chaque âme individuelle ensuite. On comprend mieux, ainsi, pourquoi la \teiHi[rend={italic}]{summetria} n’est qu’un critère secondaire de beauté : « Il faut \teiHi[rend={italic}]{donc} dire qu’ici-bas aussi, la beauté se trouve dans la lumière qui brille au-dessus de la symétrie, plutôt que dans la symétrie, et que c’est cela qui suscite l’amour. » Dans son commentaire, Pierre Hadot explique : « La lumière du Bien brille donc sur l’Intellect dans la mesure où il est Vie, \teiHi[rend={italic}]{énergéia} émanant du Bien. C’est d’ailleurs pourquoi Plotin dit que la lumière du Bien est sur l’Intellect comme la grâce de la \teiHi[rend={italic}]{vie} sur un visage\teiNote[xmlid={ftn70},n={31},place={foot}]{\teiP[xmlid={p-055}]{Pierre Hadot, \teiHi[rend={italic}]{Traité 38}, \teiHi[rend={italic}]{op. cit.}, p. 290.}}. » Et il poursuit : « Ce qui correspond ici-bas à la lumière du Bien, c’est la Vie, précisément parce que, dans le monde spirituel, la lumière du Bien chatoyant sur l’Intellect se trouve dans la Vie. Partout, c’est la Vie qui attire l’amour de l’âme, parce qu’elle est la trace du Bien\teiNote[xmlid={ftn71},n={32},place={foot}]{\teiP[xmlid={p-056}]{Pierre Hadot, \teiHi[rend={italic}]{Traité 38}, \teiHi[rend={italic}]{op. cit.,} p. 290 et 291.}}. »}
            \teiP[xmlid={p15}]{La comparaison saisissante qui suit, entre un visage mort et un visage vivant, confirme cette thèse. De quelle vie s’agit-il ? L’on serait tenté de dire que c’est avant tout la vie propre à l’intellect de l’être humain dont on contemple le visage, mais ce peut être aussi, à un niveau inférieur, la vie organique, qui est un produit de la vie intellectuelle. Or, ce qui est vrai du visage humain « réel », pourrait-on dire, l’est aussi du visage représenté dans l’art, comme le montre l’exemple des statues : « Et parmi les statues, les plus vivantes ne sont-elles pas les plus belles, même si les autres ont plus de symétrie ? » Qu’entend-on par « statue vivante » ? Une statue dans laquelle se reflète, là aussi, la vie de l’intellect, qui ne s’exprime pas nécessairement, ou pas seulement, par de bonnes proportions (\teiHi[rend={italic}]{summetria}), mais qui est comme « animée » (nous employons ce terme à dessein) par la vie de l’intellect et \teiHi[rend={italic}]{par conséquent}, par la présence, l’éclat, la lumière ou la grâce du Bien. }
            \teiP[xmlid={p16}]{On notera cependant que Plotin distingue la vie, et \teiHi[rend={italic}]{donc} la beauté, qui se manifestent dans l’art et celles qui se manifestent dans l’être humain vivant : « Et un homme laid, s’il est vivant, n’est-il pas plus beau que l’homme beau représenté dans une statue ? » La différence tient à l’émotion, au sentiment que provoque la vision de ces deux types de beauté. L’une est plus désirable que l’autre. Lorsque Plotin évoque un homme « laid », mais vivant, on ne peut pas ne pas penser à la figure de Socrate, tel qu’il est décrit par Alcibiade à la fin du \teiHi[rend={italic}]{Banquet} (215b) : laid à l’extérieur, mais beau à l’intérieur, et objet de désir pour Alcibiade comme pour d’autres, précisément à cause de l’intérieur, c’est-à-dire à cause de la beauté de son âme d’homme sage, philosophe. Les dernières lignes du passage le confirment. Si un homme laid vivant peut être beau, c’est, dit Plotin, « parce qu’il a une âme ; or cela, c’est parce qu’il a davantage la forme du bien ». Et c’est effectivement la vie, la vie de l’âme que décrit alors Plotin. Une vie lumineuse, colorée de la lumière du Bien, dont elle manifeste l’\teiHi[rend={italic}]{énergéia} dans toutes ses actions : l’éveil, le mouvement d’élévation et d’allègement qu’elle imprime à elle-même comme à « ce qui lui appartient », la capacité à « rendre bon ». }
            \teiP[xmlid={p17}]{Il est donc indéniable, à la lecture de ce passage du \teiHi[rend={italic}]{Traité 38}, en regard des interrogations et des définitions formulées dans le \teiHi[rend={italic}]{Traité 1}, que la beauté entretient, pour Plotin, un lien étroit avec le vivant et avec la vie, qu’elle est « vitale » au sens où elle s’ancre dans la vie. Certes, l’on pourrait objecter que la vie dont il parle ici est une vie débarrassée de la pesanteur du corps et de la matière, une vie « pure » car purifiée, la « vraie » vie de l’âme, qui est la vie dans l’Intellect. Mais son souci constant d’expliciter les raisons profondes des expériences esthétiques et amoureuses que nous faisons dans le monde sensible, lorsque nous sommes saisis de stupeur et de désir à la vue d’un corps humain, d’une œuvre d’art ou d’une action morale, montre que, même dans ce monde-ci, ce qui rend beaux ces différents objets est aussi une forme de vie. C’est ainsi que l’on assiste, chez ce philosophe « platonicien », à une réhabilitation, certes mesurée, mais réelle, de la beauté sensible.}
            \teiP[xmlid={p18}]{On en donnera ici deux exemples dans lesquels le travail de l’artiste et, en l’occurrence, celui du sculpteur, est ainsi valorisé. }
            \teiP[xmlid={p19}]{Dans le \teiHi[rend={italic}]{Traité 31} – l’autre traité explicitement consacré par Plotin à la question du Beau, et qui a reçu pour titre « De la beauté intelligible » –, Plotin critique de façon allusive une conception de l’art présente dans plusieurs dialogues de Platon, notamment au Livre X de la \teiHi[rend={italic}]{République}. Il s’agit de la thèse bien connue selon laquelle l’art constitue une simple imitation de la nature, puisque l’artiste se contente de reproduire un modèle naturel, qui imite lui-même un modèle intelligible. Or, dit Plotin :}
            \begin{teiCit}[xmlid={cit02}]

              \begin{teiQuote}
Si quelqu’un méprise les arts parce qu’ils produisent en imitant la nature, il faut dire tout d’abord que les choses naturelles aussi imitent d’autres choses. Ensuite, il faut savoir que les arts n’imitent pas tout simplement ce qui est vu, mais qu’ils remontent aux principes rationnels, dont provient la nature. Puis il faut savoir aussi qu’ils produisent beaucoup de choses par eux-mêmes et même suppléent, parce qu’ils possèdent la beauté, à toute déficience en quelque objet que ce soit : ainsi Phidias, qui produisit son Zeus sans se référer à aucun objet sensible, mais en le concevant tel qu’il serait, s’il consentait à apparaître à nos regards\teiNote[xmlid={ftn72},n={33},place={foot}]{\teiP[xmlid={p-057}]{\teiHi[rend={italic}]{Traité 31}, chap. 1, lignes 34-40. }}. 
\end{teiQuote}
            
\end{teiCit}
            \teiP[xmlid={p20}]{L’exemple du Zeus sculpté par Phidias constitue, de fait, une objection solide puisque, par définition, il n’existe pas de modèle sensible de la divinité. Si l’art n’imite pas la nature, mais agit \teiHi[rend={italic}]{comme} elle\teiNote[xmlid={ftn73},n={34},place={foot}]{\teiP[xmlid={p-058}]{Plotin suit ici Aristote. Voir par exemple \teiHi[rend={italic}]{Physique}, Livre II, chap. 8, 199a. Nous nous permettons de renvoyer ici, pour davantage de précisions, à notre commentaire : Plotin, \teiHi[rend={italic}]{Traité 31}, introduction, traduction, commentaire et notes par Anne-Lise Darras-Worms, Paris, Vrin, 2018. }}, c’est parce qu’il est le produit d’une âme et d’un intellect agissant, en d’autres termes, de la vie de l’intellect dans l’artiste. Or, ce qui est vrai d’une statue l’est aussi du corps humain non figuré, réel. Plotin s’interroge ainsi, un peu plus loin dans le même traité : « De quelle source a donc rayonné la beauté d’Hélène qui fut l’objet de tant de combats, ou de toutes ces femmes dont la beauté est semblable à celle d’Aphrodite\teiNote[xmlid={ftn74},n={35},place={foot}]{\teiP[xmlid={p-059}]{\teiHi[rend={italic}]{Traité 31}, chap. 2, lignes 9-11.}} ? » On notera ici l’emploi du verbe « rayonner », qui renvoie à la notion de lumière, indissociable de celle de beauté et qui fait signe vers l’\teiHi[rend={italic}]{energeia} du Bien, comme nous l’avons vu précédemment. }
            \teiP[xmlid={p21}]{Dans le dernier chapitre du \teiHi[rend={italic}]{Traité 1}, l’on retrouve, à titre de comparaison, l’exemple de la statue. À la question : « Comment pourrais-tu voir quelle est la sorte de beauté que possède une âme bonne ? », Plotin répond : }
            \begin{teiCit}[xmlid={cit03}]

              \begin{teiQuote}
Reviens en toi-même et vois : et si tu ne t’es pas encore vu toi-même comme beau, fais comme le sculpteur d’une statue qui doit devenir belle – il enlève ceci, il a gratté cela, il a rendu lisse tel endroit, pur tel autre, jusqu’à ce qu’il ait fait apparaître un beau visage sur la statue – de la même manière, enlève, toi aussi, tout ce qui est superflu et redresse tout ce qui est tordu, purifie tout ce qui est obscur puis travaille à le rendre brillant et ne cesse pas de \teiHi[rend={italic}]{sculpter} ta propre \teiHi[rend={italic}]{statue}, jusqu’à ce que resplendisse pour toi la splendeur divine de la vertu, jusqu’à ce que tu voies la \teiHi[rend={italic}]{tempérance installée sur son socle sacré}\teiNote[xmlid={ftn75},n={36},place={foot}]{\teiP[xmlid={p-060}]{\teiHi[rend={italic}]{Traité 1}, chap. 9, lignes 6-15.}}.
\end{teiQuote}
            
\end{teiCit}
            \teiP[xmlid={p22}]{L’expression « sculpter une statue » est empruntée au \teiHi[rend={italic}]{Phèdre} de Platon, où Socrate dit, de celui qui est amoureux d’un beau garçon, qu’il « le sculpte pour lui-même et l’orne comme une statue, dans l’intention de l’honorer et de le célébrer en des mystères\teiNote[xmlid={ftn76},n={37},place={foot}]{\teiP[xmlid={p-061}]{\teiHi[rend={italic}]{Phèdre}, 252d, 1-9.}} ». Le parallèle peut être fait avec l’exemple de Phidias dans le \teiHi[rend={italic}]{Traité 31}. L’action du sculpteur est comparable à une action morale, telle la purification des passions de l’âme par la pratique des vertus. Il s’agit donc de faire de soi, en un certain sens, dans un acte créateur, une statue vivante. Or, c’est là ce qui transparaît sur le visage du sage : « Ainsi, par exemple, pour un homme de bien, douce est la trace de la vertu qui se manifeste clairement chez un jeune homme, car elle concorde avec la vérité qui est à l’intérieur\teiNote[xmlid={ftn77},n={38},place={foot}]{\teiP[xmlid={p-062}]{\teiHi[rend={italic}]{Traité 1}, chap. 3, lignes 15 et 16. }} » ou encore ce passage du chapitre 5 dans lequel Plotin évoque la contemplation, en soi-même ou chez autrui, d’une « grandeur d’âme, un caractère juste, une tempérance pure, un courage au “visage viril”, une gravité et une pudeur qui se diffusent dans une disposition sans trouble, sans vague et sans émotion et, répandant sa lumière sur tout cela, l’Intellect de forme divine\teiNote[xmlid={ftn78},n={39},place={foot}]{\teiP[xmlid={p-063}]{\teiHi[rend={italic}]{Traité 1}, chap. 5, lignes 13-17. L’expression « visage viril » est employée par Homère dans l’\teiHi[rend={italic}]{Iliade} (chant VII, v. 212) à propos d’Ajax.}} ».}
            \teiP[xmlid={p23}]{Notons ici que ce texte du \teiHi[rend={italic}]{Traité 31}, de même que celui du \teiHi[rend={italic}]{Traité 38}, illustrent bien les connotations à la fois esthétiques et morales de l’adjectif \teiHi[rend={italic}]{kalos} qui, on le sait, signifie à la fois « beau » et « bon » en grec. Peu importe, donc, que le visage d’un être humain, vivant ou figuré, soit plus ou moins symétrique, car l’essentiel est que l’on y contemple la vie, c’est-à-dire la trace du Bien. Cela seul peut nous faire dire que ce visage est beau. L’acte de contemplation est en effet l’acte fondamental de qui cherche à savoir d’où provient la beauté, quelle qu’elle soit. Nous qui contemplons, contemplons des êtres qui eux-mêmes contemplent ou sont, dans le cas des œuvres d’art, mais pas seulement, des produits d’une contemplation. C’est que, comme le montre Plotin, notamment dans le \teiHi[rend={italic}]{Traité 30} (III, 8, « De la contemplation »), toute vie, depuis la vie organique jusqu’à la Vie dans l’Intellect, est contemplation. C’est ainsi que la beauté vit et \teiHi[rend={italic}]{se} vit. }
            \teiP[xmlid={p24}]{Nous conclurons avec la description émouvante que donne Porphyre, dans la \teiHi[rend={italic}]{Vie de Plotin}, de son maître lorsqu’il enseignait, car elle nous paraît tout à fait révélatrice, à elle seule, du lien indissoluble qui unit vie et beauté : }
            \begin{teiCit}[xmlid={cit04}]

              \begin{teiQuote}
Quand il parlait se manifestait l’intellect, qui faisait briller sa lumière jusque sur son visage ; lui qui était agréable à voir, il apparaissait particulièrement beau surtout dans ces moments-là ; une légère sueur couvrait son visage, sa douceur rayonnait et face aux questions se montrait sa bienveillance, ainsi que sa vigueur\teiNote[xmlid={ftn79},n={40},place={foot}]{\teiP[xmlid={p-064}]{Porphyre, \teiHi[rend={italic}]{Vie de Plotin}, chap. 13, lignes 5-10. Traduction : Porphyre, \teiHi[rend={italic}]{Vie de Plotin}, Paris, Vrin, 1992, vol. II, \teiHi[rend={italic}]{Études d’introduction, texte grec et traduction française, commentaire, notes complémentaires, bibliographie}, par Luc Brisson, Jean-Louis Cherlonneix, Marie-Odile Goulet-Cazé \teiHi[rend={italic}]{et al.}, p. 155. }}.
\end{teiQuote}
            
\end{teiCit}
          
\end{teiElement}
        
\end{teiElement}
        \begin{teiElement}[name={group},type={chapter},data-page-title={Vieux et laid, mais sage : éloges romains de la laideur},data-page-authors={Catherine Baroin@@Université de Rouen Normandie, ÉRIAC – EA 4705, France},data-include-href={Ch03\_vieux\_laid\_sage\_Baroin.xml},xmlid={chapter-002}]

          \begin{teiElement}[name={front}]

            \begin{teiDiv}[type={titlePage},xmlid={titlepage-004}]

              \teiP[rend={title-main},xmlid={p-065}]{Vieux et laid, mais sage : éloges romains de la laideur}
              \teiP[rend={author-aut},xmlid={p-066}]{Catherine \teiHi[rend={small-caps}]{Baroin}}
              \teiP[rend={authority\_affiliation},xmlid={p-067}]{Université de Rouen Normandie, ÉRIAC – EA 4705, France}
            
\end{teiDiv}
          
\end{teiElement}
          \begin{teiElement}[name={body},xmlid={body-004}]

            \teiP[xmlid={p1-4}]{Dans les textes latins, sur le temps long (du \teiHi[rend={small-caps}]{ii}\teiHi[rend={sup}]{e} siècle avant notre ère au \teiHi[rend={small-caps}]{ii}\teiHi[rend={sup}]{e} siècle de notre ère au moins), il y a une association, voire une équivalence, entre la beauté physique et la beauté morale et, symétriquement, entre la laideur physique et la laideur morale. Cette association est exprimée par la langue : dans la conclusion de son ouvrage \teiHi[rend={italic}]{Beau et laid en latin. Étude de vocabulaire}, Pierre Monteil indique que, dans la langue latine, « très rares sont les termes qui expriment, à l’exclusion de toute autre, une notion esthétique de “beauté” ou de “laideur” », car le lexique de la beauté et de la laideur est à la fois esthétique et moral, et il relève un peu plus loin la « très fréquente interférence de la beauté physique et de la beauté morale ; de la laideur physique et de la laideur morale\teiNote[xmlid={ftn1-3},n={1},place={foot}]{\teiP[xmlid={p-068}]{Pierre Monteil, \teiHi[rend={italic}]{Beau et laid en latin. Étude de vocabulaire}, Paris, Klincksieck, 1964, p. 340 et 341.}} ».}
            \teiP[xmlid={p2-4}]{La valorisation de la beauté physique est également à l’œuvre dans les images plastiques, en particulier dans la statuaire, lorsque les hommes et les femmes qui se font représenter dans des sculptures, bustes ou statues en pied, choisissent d’apparaître avec un corps et un visage dotés des caractéristiques traditionnelles de la beauté : des membres bien proportionnés, une peau lisse, une allure juvénile… C’est particulièrement vrai, sous l’Empire, des Julio-Claudiens, figurés avec un visage jeune, selon un modèle qui rappelle celui d’Alexandre le Grand et des portraits hellénistiques\teiNote[xmlid={ftn77-2},n={2},place={foot}]{\teiP[xmlid={p-069}]{Voir notamment Peter Stewart, \teiHi[rend={italic}]{Roman art}, Oxford University Press, 2004 ; Christopher H. Hallett, \teiHi[rend={italic}]{The Roman Nude. Heroic Portrait Statuary 200 BC-300 AD}, Oxford University Press, 2005, chap. 5.}}. Cependant, à côté de ces portraits idéalisés, d’autres sont qualifiés par les historiens d’art de « véristes » ou de « naturalistes », lorsqu’ils figurent des rides, des défauts tels que les verrues, peut-être dans un souci de rendre une ressemblance avec l’individu représenté, ou bien pour marquer une qualité morale\teiNote[xmlid={ftn78-2},n={3},place={foot}]{\teiP[xmlid={p-070}]{Voir Peter Stewart, \teiHi[rend={italic}]{op. cit.} ; Véronique Dasen, « Autour du portrait romain : marques identitaires et anomalies physiques », dans Agostino Paravicini Bagliani, Jean-Michel Spieser et Jean Wirth (dir.), \teiHi[rend={italic}]{Le portrait. La représentation de l’individu}, Florence, Editions del Galluzzo, 2007, p. 17-33 ; Susan B. Matheson et J. J. Pollitt, \teiHi[rend={italic}]{Old Age in Greek and Roman Art}, New Haven, Yale University Art Gallery, 2022.}}. En outre, des études récentes s’intéressent aux terres cuites ou aux bronzes de taille plus réduite qui représentent des personnages laids, contrefaits, renvoyant à des types « physiques » (les nains), ethniques (les pygmées), comiques, ou bien à des pathologies\teiNote[xmlid={ftn79-2},n={4},place={foot}]{\teiP[xmlid={p-071}]{Voir Christophe Vendries, « La caricature et l’Antiquité dans les arts visuels. Regard historiographique et épistémologique », dans Anne Gangloff, Valérie Huet et Christophe Vendries (dir.), \teiHi[rend={italic}]{La notion de caricature dans l’Antiquité. Textes et images}, Rennes, PUR, 2021, p. 25-47 (en particulier p. 42-44).}}.}
            \teiP[xmlid={p3-4}]{Cette question de la représentation, littéraire et iconographique, de la beauté et de la laideur ne se conçoit pas sans une distinction essentielle : celle du statut des individus et de la ligne de partage fondamentale dans le monde ancien qui passe entre individus libres et non libres. La plupart du temps, en effet, les représentations de corps déformés, contrefaits, tordus, de traits grossiers, sont celles d’individus, masculins ou féminins, de statut social modeste ou bien de rang servile.}
            \teiP[xmlid={p4-4}]{C’est sur ce « fond » de concepts et de représentations, mentales et plastiques, que se dessine dans certains textes la figure de l’homme laid et/ou vieux, mais vertueux moralement, et pouvant même être qualifié de sage. On peut considérer que le modèle de ce personnage est celui du philosophe Socrate, étudié notamment par Paul Zanker\teiNote[xmlid={ftn80},n={5},place={foot}]{\teiP[xmlid={p-072}]{Paul Zanker, \teiHi[rend={italic}]{The Mask of Socrates. The Image of the Intellectual in Antiquity }\lbrack{}\teiHi[rend={italic}]{Die Maske des Sokrates : Das Bild des Intellektuellen in der antiken Kunst}, Münich, C. H. Bech, 1995\rbrack{}, Alan Shapiro (trad.), Berkeley, Los Angeles et Oxford, University of California Press, 1995. Sur les portraits littéraires et iconographiques de Socrate, voir Mélanie Lucciano, « Les représentations iconographiques de Socrate », \teiHi[rend={italic}]{Camenae}, n\teiHi[rend={sup}]{o} 10, 2012.}}. Mais d’autres figures sont concernées, en particulier des figures romaines, que l’on peut mettre en relation avec des conceptions romaines de la laideur et de la beauté et avec des pratiques romaines du soin de soi, afin de voir comment se déploient dans les textes ces portraits de philosophes à la fois laids physiquement, ou du moins éloignés de l’apparence traditionnelle du citoyen homme de bien, mais beaux moralement. À cette fin, on s’intéressera rapidement, dans un premier temps, à l’articulation entre laideur physique et laideur morale, puis au divorce noté et critiqué par certains textes entre apparence physique et mauvaises mœurs, et l’on s’attardera enfin sur deux figures de philosophes vieux et pouvant sembler laids, mais véritablement sages.}
            \begin{teiDiv}[type={section1},xmlid={div01}]

              \teiHead[xmlid={laideur-physique-et-morale-001}]{Laideur physique et morale}
              \teiP[xmlid={p5-4}]{De même qu’il y a, dans les textes latins, une association entre la beauté convenable et honorable du citoyen – sa \teiHi[rend={italic}]{dignitas} en termes cicéroniens (\teiHi[rend={italic}]{De officiis}, I, 130) – et sa moralité, sa qualité d’homme de bien\teiNote[xmlid={ftn81},n={6},place={foot}]{\teiP[xmlid={p-073}]{Voir Catherine Baroin, « La beauté du corps masculin dans le monde romain : état de la recherche récente et pistes de réflexion », \teiHi[rend={italic}]{DHA}, suppl. 14, 2015, p. 31-51 (en particulier p. 32-39), \teiRef[target={https://www.cairn.info/revue-dialogues-d-histoire-ancienne-2015-Supplement14-page-31.htm}]{https://www.cairn.info/revue-dialogues-d-histoire-ancienne-2015-Supplement14-page-31.htm}. Dans les références aux textes anciens, sont indiqués successivement les livres en chiffres romains, les chapitres le cas échéant ou les numéros des lettres ou poèmes, puis les paragraphes ou les vers pour la poésie. Sauf indication contraire, les textes anciens sont cités dans la collection « CUF » des Belles Lettres et les traductions sont personnelles.}}, il existe, en parallèle, une association entre la laideur physique et la laideur morale, défauts qui se disent avec les mêmes termes latins : \teiHi[rend={italic}]{deformitas},\teiHi[rend={italic}]{ foeditas}, \teiHi[rend={italic}]{turpitudo}\teiNote[xmlid={ftn82},n={7},place={foot}]{\teiP[xmlid={p-074}]{Le substantif \teiHi[rend={italic}]{deformitas}, qui signifie laideur ou difformité, est le contraire de la beauté, qui peut se dire \teiHi[rend={italic}]{forma} ; l’adjectif « \teiHi[rend={italic}]{formosus} », qui signifie beau, s’oppose à « \teiHi[rend={italic}]{informis} » ou à « \teiHi[rend={italic}]{deformis} » : les termes qui expriment la laideur désignent l’absence, la privation de beauté (avec les préfixes \teiHi[rend={italic}]{in}- et \teiHi[rend={italic}]{de}-). Sur ces termes, voir Pierre Monteil, \teiHi[rend={italic}]{op. cit}., p. 60-69, 343 et 344 ; Jean-François Thomas, \teiHi[rend={italic}]{Déshonneur et honte en latin : étude sémantique}, Louvain, Peeters, 2007, p. 124 et suiv. (sur \teiHi[rend={italic}]{turpis} et \teiHi[rend={italic}]{turpitudo}) ; Danièle Conso, \teiHi[rend={italic}]{Forma. Étude sémantique et étymologique}, vol. I, Besançon, Presses universitaires de Franche-Comté, 2015, p. 32 et 33 et p. 371-381 (sur \teiHi[rend={italic}]{deformis }et \teiHi[rend={italic}]{deformitas}).}}. C’est ainsi que des hommes jugés laids sont considérés comme mauvais moralement, ou l’inverse, sans que l’on puisse discerner ce qui est premier dans cette estimation. L’un des exemples les plus représentatifs à ce titre est celui de Vatinius. Vatinius est un homme politique romain (consul en 47), qui a été l’allié de Jules César et l’ennemi politique de Cicéron, avant que les deux hommes ne se réconcilient\teiNote[xmlid={ftn83},n={8},place={foot}]{\teiP[xmlid={p-075}]{Sur sa carrière, voir Antonio Pistellato, « Historiographie des guerres civiles et guerre civile des historiographies. Publius Vatinius », dans Robinson Baudry et Sylvain Destephen (dir.), \teiHi[rend={italic}]{La société romaine et ses élites. Hommages à Élisabeth Deniaux}, Paris, Picard, 2012, p. 43-52.}}. Plusieurs témoignages textuels, de Cicéron d’abord, puis d’autres auteurs qui l’ont visiblement repris, indiquent que Vatinius souffrait de \teiHi[rend={italic}]{strumae}, ce qu’on traduit par scrofules, écrouelles ou goitre, en tout cas d’une affection visible du visage et du cou\teiNote[xmlid={ftn84},n={9},place={foot}]{\teiP[xmlid={p-076}]{Voir Cicéron, \teiHi[rend={italic}]{Contre Vatinius}, 4 et 39.}}. Certains textes mentionnent aussi des problèmes aux membres inférieurs (pieds ou jambes), peut-être la goutte\teiNote[xmlid={ftn85},n={10},place={foot}]{\teiP[xmlid={p-077}]{Sénèque, \teiHi[rend={italic}]{De constantia sapientis}, 17, 3 ; Quintilien, \teiHi[rend={italic}]{Institution oratoire}, VI, 3, 77 ; Plutarque, \teiHi[rend={italic}]{Cicéron}, 9, 3 et 26, 3 ; Macrobe, \teiHi[rend={italic}]{Saturnales}, II, 4, 16. Sur les défauts physiques attribués à Vatinius, voir Catherine Baroin, « Le corps du prêtre romain dans le culte public : début d’une enquête », dans Lydie Bodiou, Véronique Mehl et Myriam Soria-Audebert (dir.), \teiHi[rend={italic}]{Corps outragés, corps ravagés de l’Antiquité au Moyen Âge}, Turnhout, Brepols, 2011, p. 291-316 (en particulier p. 306-308).}}. Ces défauts physiques alimentent l’invective politique des discours de Cicéron ainsi que ses lettres\teiNote[xmlid={ftn86},n={11},place={foot}]{\teiP[xmlid={p-078}]{Voir Cicéron, \teiHi[rend={italic}]{Ad Atticum}, II, 9, 2 (Paris, Les Belles Lettres, « CUF », t. I, lettre 36, 1934).}} et ont influencé la tradition historiographique postérieure sur la laideur, morale et physique, de ce personnage. Ainsi l’historien Velléius Paterculus écrit-il :}
              \begin{teiCit}[xmlid={cit01-2}]

                \begin{teiQuote}
Il n’y avait personne à qui Vatinius ne parût inférieur : la difformité de son corps \lbrack{}\teiHi[rend={italic}]{deformitas corporis}\rbrack{} rivalisait avec la laideur honteuse de son esprit \lbrack{}\teiHi[rend={italic}]{turpitudo ingenii}\rbrack{}, au point que son âme \lbrack{}\teiHi[rend={italic}]{animus}\rbrack{} semblait enfermée dans une demeure bien digne d’elle \lbrack{}\teiHi[rend={italic}]{eius dignissimo domicilio inclusus}\teiNote[xmlid={ftn87},n={12},place={foot}]{\teiP[xmlid={p-079}]{Velléius Paterculus, \teiHi[rend={italic}]{Histoire romaine}, II, 69, 3-4.}}\rbrack{}.
\end{teiQuote}
              
\end{teiCit}
              \teiP[xmlid={p6-4}]{De plus, comme l’indique une remarque du philosophe Sénèque, cette laideur provoque le rire et la moquerie : Sénèque écrit en effet que Vatinius était « un homme né pour qu’on s’en moque et pour qu’on le haïsse\teiNote[xmlid={ftn88},n={13},place={foot}]{\teiP[xmlid={p-080}]{Sénèque, \teiHi[rend={italic}]{De constantia sapientis}, 17, 3.}} ». De façon générale, la laideur et les défauts physiques sont une source de plaisanteries (\teiHi[rend={italic}]{facetiae}) dans les discours politiques et judiciaires\teiNote[xmlid={ftn89},n={14},place={foot}]{\teiP[xmlid={p-081}]{Voir Cicéron, \teiHi[rend={italic}]{De oratore}, II, 236 (repris par Quintilien, \teiHi[rend={italic}]{Institution oratoire}, VI, 3, 8).}}, plaisanteries qui visent à faire rire le public et/ou les juges, afin de les faire pencher en faveur de celui qui parle et de décrédibiliser l’adversaire. La laideur peut faire rire, elle peut faire peur aussi\teiNote[xmlid={ftn90},n={15},place={foot}]{\teiP[xmlid={p-082}]{Voir Sénèque, \teiHi[rend={italic}]{De constantia sapientis}, 5, 2 (à propos des enfants).}}.}
              \teiP[xmlid={p7-4}]{L’adéquation entre vices moraux et laideur physique est aussi visible dans un autre domaine : celui des allégories, qu’on rencontre en particulier dans les textes poétiques. Ainsi, dans les \teiHi[rend={italic}]{Métamorphoses} d’Ovide, les défauts moraux (\teiHi[rend={italic}]{uitia}) sont des personnages laids et répugnants. L’Envie (\teiHi[rend={italic}]{Inuidia}), par exemple, est décrite ainsi (par opposition à la beauté de la déesse Athéna, qui est venue lui demander son intervention) : « La pâleur siège sur son visage, la maigreur, sur son corps entier ; son regard n’est jamais droit, ses dents sont jaunies par le tartre, son cœur (\teiHi[rend={italic}]{pectora}) est vert de fiel, sa langue est imprégnée de poison\teiNote[xmlid={ftn91},n={16},place={foot}]{\teiP[xmlid={p-083}]{Ovide, \teiHi[rend={italic}]{Métamorphoses}, II, 775-777.}}. » La pâleur, la maigreur, le regard torve, les dents sales sont des signes de mauvaise hygiène physique et mentale à la fois. On retrouve les mêmes signes, qui sont ici accumulés, dans d’autres contextes littéraires pour exprimer tout ensemble la laideur, la mauvaise santé et l’absence de moralité.}
            
\end{teiDiv}
            \begin{teiDiv}[type={section1},xmlid={div02}]

              \teiHead[xmlid={divorce-entre-apparence-physique-et-mentalite-001}]{Divorce entre apparence physique et mentalité}
              \teiP[xmlid={p8-4}]{Si un aspect physique honorable est censé correspondre à une bonne moralité – un caractère vertueux et un genre de vie respectable –, par un effet de symétrie inversée, dans d’autres textes, est dénoncé le divorce entre l’apparence de la beauté, le caractère convenable, voire séduisant de l’aspect extérieur, et une laideur morale dissimulée. C’est l’un des objets de la satire latine. D’après le poète Horace, Lucilius, le créateur romain de ce genre littéraire, a ainsi osé, le premier, écrire des satires et « arracher la peau brillante dont chacun, paradant sous les regards, recouvrait sa laideur intérieure\teiNote[xmlid={ftn92},n={17},place={foot}]{\teiP[xmlid={p-084}]{Horace, \teiHi[rend={italic}]{Satires}, II, 1, 64-65 : « \teiHi[rend={italic}]{detrahere et pellem, nitidus qua quisque per ora/cederet, introrsum turpis }\lbrack{}…\rbrack{} » (François Villeneuve trad., Paris, Les Belles Lettres, « CUF », 1932, sauf pour \teiHi[rend={italic}]{pellem}, traduit par « enveloppe », mais que nous préférons traduire littéralement). D’après Marie-Claire Rolland, « La peau-objet chez les Romains. Indifférence, curiosité ou répulsion ? », \teiHi[rend={italic}]{La Peaulogie}, n\teiHi[rend={sup}]{o} 2, 2018, \teiRef[target={https://lapeaulogie.fr/la-peau-objet-chez-les-romains-indifference-curiosite-ou-repulsion/}]{https://lapeaulogie.fr/la-peau-objet-chez-les-romains-indifference-curiosite-ou-repulsion/}, le terme \teiHi[rend={italic}]{pellis} désigne plutôt la peau de bête, comme un vêtement, par distinction avec \teiHi[rend={italic}]{cutis, }employé pour la peau humaine et vulgarisé par Celse. Voir aussi, notamment sur la peau comme « enveloppe qui ne correspond pas à l’intériorité de l’individu », Floriane Sanfilippo, « Les métaphores du vêtement dans le lexique de la peau en latin », dans Juliette Delalande, Barthélémy Enfrein, Misel Jabin \teiHi[rend={italic}]{et al. }(dir.), \teiHi[rend={italic}]{Himation}.\teiHi[rend={italic}]{ Métaphores du vêtement dans l’Antiquité classique et tardive}, Pessac, Ausonius Éditions, 2024, \teiRef[target={https://una-editions.fr/les-metaphores-du-vetement-dans-le-lexique-de-la-peau-en-latin/}]{https://una-editions.fr/les-metaphores-du-vetement-dans-le-lexique-de-la-peau-en-latin/}.}} ». Ailleurs, Horace dénonce lui aussi l’hypocrisie morale et sociale de certains :}
              \begin{teiCit}[xmlid={cit02-2}]

                \begin{teiQuote}
L’homme de bien \lbrack{}\teiHi[rend={italic}]{uir bonus}\rbrack{}, quel est-il ? « Celui qui observe les décrets du sénat, les lois et le droit, le juge qui tranche des projets nombreux et importants, l’homme dont la parole fait foi quand il est répondant dans une affaire et témoin dans une cause. » Cependant, toute sa maison et le voisinage entier voient sa laideur intérieure \lbrack{}\teiHi[rend={italic}]{introrsum turpem}\rbrack{}, mais qui a un bel aspect sous une jolie peau \lbrack{}\teiHi[rend={italic}]{speciosum pelle decora}\teiNote[xmlid={ftn93},n={18},place={foot}]{\teiP[xmlid={p-085}]{ Horace, \teiHi[rend={italic}]{Épîtres}, I, 16, 40-45 (François Villeneuve trad., Paris, Les Belles Lettres, « CUF », 1934, modifiée ici aussi pour \teiHi[rend={italic}]{pelle}).}}\rbrack{}.
\end{teiQuote}
              
\end{teiCit}
              \teiP[xmlid={p9-3}]{Il y a ici une opposition entre le rôle public, que tous peuvent voir, et la vie privée, qui n’est visible que de la maisonnée et du voisinage. Un homme peut être laid et honteux (\teiHi[rend={italic}]{turpis}) « à l’intérieur \lbrack{}\teiHi[rend={italic}]{introrsum}\rbrack{} » tout en conservant une belle apparence (\teiHi[rend={italic}]{speciosum}), grâce à une « peau » à la fois élégante et convenable, bienséante (\teiHi[rend={italic}]{decora}). Plus nettement que dans les vers d’Horace à propos de Lucilius, cet « intérieur » est ici un espace domestique, privé, et pas nécessairement une intériorité psychologique ou morale\teiNote[xmlid={ftn94},n={19},place={foot}]{\teiP[xmlid={p-086}]{Sur cette question, dans le champ de la philosophie, voir Frédérique Ildefonse, \teiHi[rend={italic}]{Le multiple dans l’âme. Sur l’intériorité comme problème}, Paris, Vrin, 2022 (en particulier l’introduction p. 9-28).}}. Cette opposition entre l’aspect extérieur, désigné ici par le mot « peau », et l’intérieur distingue moins une intériorité mentale et une extériorité physique que ce qui n’est pas visible de tous (la vie privée, la moralité, le genre de vie dans la maison) et ce qui est montré dans la vie publique.}
              \teiP[xmlid={p10-3}]{\teiAnchor[xmlid={_gjdgxs}]Cette lecture peut s’appuyer sur les exemples qu’on trouve chez un autre auteur de satires, Juvénal, ou dans les épigrammes satiriques de Martial\teiNote[xmlid={ftn95},n={20},place={foot}]{\teiP[xmlid={p-087}]{Voir Martial, \teiHi[rend={italic}]{Épigrammes}, IX, 47 sur Pannychus, prétendu philosophe.}}, qui dénoncent la contradiction entre, d’une part, une apparence virile, parfois associée à une allure de philosophe, ainsi qu’un discours de moraliste, et, d’autre part, une vie privée de débauché. Ainsi, une \teiHi[rend={italic}]{Satire} (2, 9-21) de Juvénal raille ceux qui affectent d’imiter les anciens Romains champions de la moralité traditionnelle, ceux qui collectionnent les bustes des philosophes grecs\teiNote[xmlid={ftn96},n={21},place={foot}]{\teiP[xmlid={p-088}]{Sont mentionnés (v. 5-8) : les Stoïciens Chrysippe et Cléanthe, Pittacos (l’un des Sept Sages) et Aristote : le rattachement à telle ou telle école n’a pas d’importance ici.}}, qui ont les « cheveux plus courts que le sourcil » (v. 15), qui prennent un air austère, alors qu’ils sont obscènes (comme le souligne l’oxymore du v. 9 : « \teiHi[rend={italic}]{tristibus obscaenis }») et débauchés (v. 10 : ce sont des « cinèdes socratiques », \teiHi[rend={italic}]{Socratici cinaedi}\teiNote[xmlid={ftn97},n={22},place={foot}]{\teiP[xmlid={p-089}]{Les \teiHi[rend={italic}]{cinaedi}, du grec \teiHi[rend={italic}]{kinaidoi}, sont des débauchés, des efféminés, accusés de cultiver les plaisirs érotiques, en particulier avec des hommes libres (voir Florence Dupont et Thierry Éloi, \teiHi[rend={italic}]{L’érotisme masculin dans la Rome antique}, Paris, Belin, 2001, p. 14, 15, 165, 166 et p. 220-223). Ici, ils sont qualifiés de « socratiques » parce qu’ils affectent de se présenter comme des philosophes.}}). D’autres signes (leur visage, (\teiHi[rend={italic}]{uultus}), leur démarche (\teiHi[rend={italic}]{incessus}), v. 17) montrent à qui sait les déchiffrer leur nature véritable. De telles lectures du corps sont proches de celles que pratiquent les physiognomonistes, qui prétendent déceler la vraie nature des hommes sous les apparences\teiNote[xmlid={ftn98},n={23},place={foot}]{\teiP[xmlid={p-090}]{Pour une présentation générale de ces traités, voir Jérôme Wilgaux, « La physiognomonie antique : bref état des lieux », dans Véronique Dasen et Jérôme Wilgaux (dir.), \teiHi[rend={italic}]{Langages et métaphores du corps dans le monde antique}, Rennes, PUR, 2008, p. 185-195.}}.}
              \teiP[xmlid={p11-3}]{On trouve aussi des dénonciations de faux philosophes en dehors du discours satirique, en particulier à l’époque de la seconde sophistique. Ceux-là prennent l’allure et les attributs traditionnels des philosophes à la grecque : le bâton, la besace, le manteau (\teiHi[rend={italic}]{tribôn}, en grec, et \teiHi[rend={italic}]{pallium}, en latin), les sandales et la barbe\teiNote[xmlid={ftn99},n={24},place={foot}]{\teiP[xmlid={p-091}]{Sur le manteau des philosophes, voir Catherine Baroin et Emmanuelle Valette-Cagnac, « Quand les Romains s’habillaient à la grecque, ou les divers usages du \teiHi[rend={italic}]{pallium} », \teiHi[rend={italic}]{Revue historique}, vol. CCCIX ; n\teiHi[rend={sup}]{o} 3, 2007, p. 517-551 (en particulier p. 537-540). Sur leurs cheveux et leur barbe, voir Bruno Sudan, « Favorinus d’Arles : corps ingrat, prodigieux destin », dans Véronique Dasen et Jérôme Wilgaux (dir.), \teiHi[rend={italic}]{op. cit.}, p. 169-182 (en particulier p. 177 et suiv\teiHi[rend={italic}]{.}) ; Pierre Brulé,\teiHi[rend={italic}]{ Les sens du poil (grec)}, Paris, Les Belles Lettres, 2015, p. 134 et 135. Paul Zanker, \teiHi[rend={italic}]{op. cit}., p. 266, note que le bâton était caractéristique des Cyniques, parce qu’il faisait allusion à Héraklès, le patron de cette « secte », et parce qu’il constituait un appui, voire une protection dans leur vie errante.}}, mais ils n’en ont pas la sagesse. Aulu-Gelle décrit ainsi un personnage « chevelu » (\teiHi[rend={italic}]{crinitus}), portant un \teiHi[rend={italic}]{pallium} et une barbe lui descendant jusqu’à l’aine, qui se présente comme philosophe auprès d’Hérode Atticus, mais qui n’est qu’un méchant mendiant\teiNote[xmlid={ftn100},n={25},place={foot}]{\teiP[xmlid={p-092}]{Aulu-Gelle, \teiHi[rend={italic}]{Nuits attiques}, IX, 2, 1 et suiv. Voir aussi Épictète, \teiHi[rend={italic}]{Entretiens}, IV, 8, 12-21 et 32 sur le fait que la connaissance de la raison et la pratique d’une vie vertueuse sont plus importantes pour définir le philosophe que le manteau, la longue barbe et la chevelure abondante. Sur les discours concernant l’apparence des philosophes à l’époque de la seconde sophistique en relation avec les différentes Écoles, voir Michael Trapp, « Visibly different? Looking at \teiHi[rend={italic}]{philosophi} in the Roman Imperial period », dans Pierre Vesperini (dir.), \teiHi[rend={italic}]{Philosophari. Usages romains des savoirs grecs sous la République et sous l’Empire}, Paris, Classiques Garnier, 2017, p. 353-369 (notamment sur Apulée, \teiHi[rend={italic}]{Apologie}, 4 et 22).}}.}
              \teiP[xmlid={p12-3}]{Contrairement à ces personnages souvent décrits comme porteurs de cheveux longs, on note que, chez Juvénal, ce sont les cheveux courts, voire rasés, qui font partie de l’apparence du philosophe et qui contribuent à donner une impression de \teiHi[rend={italic}]{tristitia}, morosité, pouvant servir à cacher des mœurs dépravées. C’est ce que critique aussi Quintilien :}
              \begin{teiCit}[xmlid={cit03-2}]

                \begin{teiQuote}
\lbrack{}…\rbrack{} à notre époque, chez la plupart d’entre eux \lbrack{}ceux qui ont enseigné la sagesse\rbrack{}, c’est sous le nom de philosophe qu’ont été dissimulés les plus grands vices. On ne travaillait pas, en effet, à être considéré comme philosophe grâce à la vertu et à l’étude \lbrack{}\teiHi[rend={italic}]{uirtute ac studiisi}\rbrack{}, mais l’expression du visage \lbrack{}\teiHi[rend={italic}]{uultum}\rbrack{}, la morosité \lbrack{}\teiHi[rend={italic}]{tristitiam}\rbrack{} et une allure en désaccord avec celle de tous les autres \lbrack{}\teiHi[rend={italic}]{dissentientem a ceteris habitum}\rbrack{} servaient à couvrir \lbrack{}\teiHi[rend={italic}]{praetendebant}\rbrack{} les mœurs les plus dépravées \lbrack{}\teiHi[rend={italic}]{pessimis moribus}\teiNote[xmlid={ftn101},n={26},place={foot}]{\teiP[xmlid={p-093}]{Quintilien, \teiHi[rend={italic}]{Institution oratoire}, I, pr. 15. Voir \teiHi[rend={italic}]{idem}, XI, 1, 34 (\teiHi[rend={italic}]{barba} et \teiHi[rend={italic}]{tristitia}) ; XII, 3, 12 sur ceux qui, se prétendant philosophes, ont laissé pousser leur barbe et prennent un air austère en public (\teiHi[rend={italic}]{in publico tristes}). Voir aussi Pline le Jeune, \teiHi[rend={italic}]{Lettres}, I, 22 : la lettre fait l’éloge de Titius Aristo, éminent jurisconsulte, qui est « une image de la frugalité antique » (§ 4), qui est prêt à se suicider si sa maladie est sans espoir de guérison (§ 8), et qui n’est pas (§ 6) comme « ceux qui étalent leur goût pour la sagesse dans leur allure physique \lbrack{}\teiHi[rend={italic}]{istis qui sapientiae studium habitu corporis praeferunt}\rbrack{} ». }}\rbrack{}.
\end{teiQuote}
              
\end{teiCit}
              \teiP[xmlid={p13-2}]{C’est non seulement le désaccord entre l’apparence austère, l’affichage de \teiHi[rend={italic}]{uirtus}, et les \teiHi[rend={italic}]{mores} qui est critiqué ici, mais encore une expression, la \teiHi[rend={italic}]{tristitia}, plus particulièrement attribuée aux philosophes stoïciens, et une mise qui se veulent différentes de celles du citoyen homme de bien.}
              \teiP[xmlid={p14-2}]{Autre exemple d’allure de philosophe qui n’est pas conforme aux mœurs romaines, celle des cyniques. Voici le portrait que fait Martial d’un philosophe qui se réclame de cette école :}
              \begin{teiCit}[xmlid={cit04-2}]

                \begin{teiQuote}
\lbrack{}…\rbrack{} ce vieux portant bâton et besace \lbrack{}\teiHi[rend={italic}]{cum baculo peraque}\rbrack{}, dont la chevelure blanche et gâtée \lbrack{}\teiHi[rend={italic}]{cana/putrisque… coma}\rbrack{} se hérisse et dont la barbe sale descend sur sa poitrine ; ce vieux qu’abrite un long manteau crasseux \lbrack{}\teiHi[rend={italic}]{cerea… abolla}\rbrack{}, qui lui tient lieu d’épouse sur son misérable grabat, et auquel les passants donnent une nourriture qu’il mendie par des aboiements, tu le prends pour un cynique, trompé que tu es par son image fabriquée \lbrack{}\teiHi[rend={italic}]{imagine ficta}\rbrack{}. Ce n’est pas un philosophe cynique, Cosmus. Qu’est-ce donc ? Un chien\teiNote[xmlid={ftn102},n={27},place={foot}]{\teiP[xmlid={p-094}]{Martial, \teiHi[rend={italic}]{Épigrammes}, IV, 53, 3-8 (H.-J. Izaac trad., Paris, Les Belles Lettres, « CUF », 1930, modifiée). Le jeu repose sur le fait que le mot grec \teiHi[rend={italic}]{kunikos} vienne du mot qui signifie « chien », \teiHi[rend={italic}]{kuôn}, \teiHi[rend={italic}]{kunos}.}}.
\end{teiQuote}
              
\end{teiCit}
              \teiP[xmlid={p15-2}]{Outre les attributs traditionnels du philosophe, énumérés plus haut, la saleté et l’activité de mendiant caractérisent ce personnage. Réduit à son apparence visuelle et sonore, il est comme la caricature de lui-même : il n’est même plus un philosophe cynique, il est réduit à un chien. Les philosophes, du moins ceux qui prétendent à cette réputation ou qui se disent sages (\teiHi[rend={italic}]{sapientes}), sont donc critiqués à plusieurs titres : pour le divorce entre leur apparence visible et leur vie privée, mais aussi pour leur défaut de soins du corps, de propreté corporelle et de bienséance vestimentaire, toutes caractéristiques qui contreviennent à l’éthique du \teiHi[rend={italic}]{uir bonus}.}
              \teiP[xmlid={p16-2}]{Cependant, on peut repérer une deuxième figure de distorsion entre l’apparence, d’un côté, et les \teiHi[rend={italic}]{mores} (le caractère) et la \teiHi[rend={italic}]{uita} (le genre de vie), de l’autre. Décrivant la décadence morale de son époque – selon un \teiHi[rend={italic}]{topos} moraliste présent à chaque génération –, Sénèque évoque, parmi les manifestations des vices, les « soins excessifs du corps et le souci de la beauté physique, qui montrent ostensiblement la laideur morale\teiNote[xmlid={ftn103},n={28},place={foot}]{\teiP[xmlid={p-095}]{Sénèque, \teiHi[rend={italic}]{De beneficiis}, I, 10, 2 : « \teiHi[rend={italic}]{cultus corporum nimius et formae cura prae se ferens animi deformitatem} »\teiHi[rend={italic}]{.}}} ». Cela correspond à une préoccupation qui apparaît plus généralement dans les textes latins de l’époque classique : pour être moralement et socialement convenable, il faut un minimum de soins apportés au corps et à la tenue vestimentaire, c’est-à-dire de \teiHi[rend={italic}]{cultus}, mais ce \teiHi[rend={italic}]{cultus} ne doit pas être excessif, sinon il existe un risque de tomber dans l’excès opposé à la saleté\teiNote[xmlid={ftn104},n={29},place={foot}]{\teiP[xmlid={p-096}]{Voir Catherine Baroin, « La beauté du corps masculin… », art. cité, p. 38 et 39.}}. Ici, en outre, la réflexion de Sénèque prend la forme d’un paradoxe typiquement stoïcien : un aspect extérieur trop soigné montre que l’on est laid sur le plan moral, parce que l’on est trop peu préoccupé de son \teiHi[rend={italic}]{animus}, de sa vertu, de sa sagesse\teiNote[xmlid={ftn105},n={30},place={foot}]{\teiP[xmlid={p-097}]{Voir Épictète, \teiHi[rend={italic}]{Entretiens}, III, 1, 9.}}.}
            
\end{teiDiv}
            \begin{teiDiv}[type={section1},xmlid={div03}]

              \teiHead[xmlid={vieux-et-laid-mais-sage-001}]{Vieux et laid, mais sage}
              \teiP[xmlid={p17-2}]{C’est donc ce cadre général qu’il faut avoir à l’esprit lorsqu’on rencontre dans les textes les portraits d’hommes âgés, dont le corps est chétif, dont le physique et l’apparence présentent des caractéristiques qui sont associées à la laideur dans d’autres contextes, mais qui sont paradoxalement valorisées. Le modèle de cette figure, comme on l’a dit plus haut, est celui de Socrate, dont la laideur s’oppose à la \teiHi[rend={italic}]{kalokagathia} aristocratique : à l’alliance de beauté et de qualité morale qui est un idéal athénien\teiNote[xmlid={ftn106},n={31},place={foot}]{\teiP[xmlid={p-098}]{Voir Paul Zanker, \teiHi[rend={italic}]{op. cit}., p. 38 et p. 112. De même, sur la figure d’Ésope qui associe laideur et sagesse, voir François Lissarrague, « Aesop, Between Man and Beast: Ancient Portraits and Illustrations », dans Beth Cohen (éd.), \teiHi[rend={italic}]{Not the Classical Ideal. Athens and the Construction of the Other in Greek Art}, Leyde-Boston-Cologne, Brill, p. 132-149 (ici p. 136) ; Jeremy B. Lefkowitz, « Ugliness and value in the \teiHi[rend={italic}]{Life of Aesop} », dans Ineke Sluiter et Ralph M. Rosen, \teiHi[rend={italic}]{Kakos. Badness and Anti-Value in Classical Antiquity}, Leyde-Boston, Brill, 2008, p. 59-81.}}. La laideur de Socrate s’accompagne de l’exceptionnelle vertu morale que lui attribuent, en particulier, les dialogues platoniciens. Dans le fameux éloge qu’Alcibiade fait de lui dans le \teiHi[rend={italic}]{Banquet}, il dit vouloir utiliser des images pour dire la vérité sur lui :}
              \begin{teiCit}[xmlid={cit05}]

                \begin{teiQuote}
Je déclare donc qu’il est tout pareil à ces silènes qu’on voit exposés dans les ateliers des sculpteurs, et que les artistes représentent un pipeau ou une flûte à la main ; si on les ouvre en deux \lbrack{}\teiHi[rend={italic}]{dioichthentes}\rbrack{}, on voit qu’ils contiennent, à l’intérieur, des statues de dieux \lbrack{}\teiHi[rend={italic}]{endothen agalmata… theôn}\rbrack{}. Je déclare ensuite qu’il a l’air du satyre Marsyas. Une chose est certaine, de figure, tu es leur pareil, Socrate : toi-même tu ne le contesteras pas. \lbrack{}…\rbrack{} il ignore tout, il ne sait rien, c’est du moins l’air qu’il se donne. N’est-ce point la façon d’un silène ? Si, tout à fait, car tels sont les dehors \lbrack{}\teiHi[rend={italic}]{exôthen}\rbrack{} du personnage, dont il est revêtu à la manière du silène sculpté. Une fois le silène ouvert \lbrack{}\teiHi[rend={italic}]{anoichtheis}\rbrack{}, avez-vous idée de toute la sagesse dont il regorge \lbrack{}\teiHi[rend={italic}]{endothen… sôphrosunês}\rbrack{}, ô buveurs mes amis ? \lbrack{}…\rbrack{} quand il est sérieux et qu’il s’ouvre \lbrack{}\teiHi[rend={italic}]{anoichthentos}\rbrack{}, je ne sais si quelqu’un a vu les images fascinantes qu’il contient \lbrack{}\teiHi[rend={italic}]{ta entos agalmata}\rbrack{}. Moi, je les ai vues déjà, et elles m’ont paru si divines \lbrack{}\teiHi[rend={italic}]{theia}\rbrack{}, et précieuses comme l’or, et parfaitement belles \lbrack{}\teiHi[rend={italic}]{pankala}\rbrack{}, et extraordinaires \lbrack{}\teiHi[rend={italic}]{thaumasta}\rbrack{}, qu’il me fallait en un mot exécuter toutes les volontés de Socrate\teiNote[xmlid={ftn107},n={32},place={foot}]{\teiP[xmlid={p-099}]{Platon, \teiHi[rend={italic}]{Banquet}, 215a-215b et 216d-216e (Paul Vicaire et Jean Laborderie trad., Paris, Les Belles Lettres, « CUF », 1992, modifiée). C’est la puissance de la parole socratique qui est comparée au charme de la musique jouée par Marsyas. Ces Silènes tiennent leur nom du vieux Satyre qui a élevé Dionysos et qui était connu pour sa laideur (nez camus, lèvres épaisses). Les traducteurs de la collection « CUF » justifient la traduction par « images fascinantes » du nom « \teiHi[rend={italic}]{agalmata} » (traduit une première fois par « statues ») en renvoyant à d’autres emplois du terme chez Platon (\teiHi[rend={italic}]{Charmide}, 154c ; \teiHi[rend={italic}]{Phèdre}, 251a, 252d). Sur ce passage, voir notamment Paul Zanker, \teiHi[rend={italic}]{op. cit}., p. 38 et 39.}}.
\end{teiQuote}
              
\end{teiCit}
              \teiP[xmlid={p18-2}]{L’opposition entre la laideur de Socrate et sa sagesse, qui relève du divin et de la beauté parfaite, est construite comme une distinction entre extérieur et intérieur, soulignée par le lexique de l’ouverture, à partir de l’image du Silène. Cependant, là aussi, la beauté de Socrate est moins « intérieure » que cachée à ceux qui ne la savent pas la voir, à ceux qui ne le fréquentent ni ne l’écoutent\teiNote[xmlid={ftn108},n={33},place={foot}]{\teiP[xmlid={p-100}]{Épictète (\teiHi[rend={italic}]{Entretiens}, IV, 11, 19), par une sorte de raccourci, considère que Socrate était « plein de grâce et de charme » et recherché par « les plus beaux et les plus nobles » (Joseph Souilhé et Amand Jagu trad., Paris, Les Belles Lettres, « CUF », 1965).}}.}
              \teiP[xmlid={p19-2}]{Y a-t-il des figures romaines qui montrent la même opposition entre une apparence physique peu séduisante et une sagesse de philosophe, belle et attractive\teiNote[xmlid={ftn109},n={34},place={foot}]{\teiP[xmlid={p-101}]{Voir Francesca Romana Berno, « Claranus, Héraclès, Mucius Scaevola : paradigmes de la persuasion dans la lettre 66 de Sénèque », dans Élisabeth Gavoille et François Guillaumont (dir.), \teiHi[rend={italic}]{Conseiller, diriger par lettre}, Tours, Presses universitaires François-Rabelais, 2017, p. 265-281, ici p. 268 et 269 : la figure socratique sert de référence pour souligner la rupture des philosophes, manifestée par leur aspect physique, avec l’opinion commune ; c’est particulièrement vrai des stoïciens et, parmi eux, de Chrysippe de Soles. Sur l’aspect attribué par la tradition à Chrysippe, voir Paul Zanker, \teiHi[rend={italic}]{op. cit}., p. 96-113 ; Susan B. Matheson et J. J. Pollitt, \teiHi[rend={italic}]{op. cit.}, p. 127 et fig. 75.}} ? On peut tout d’abord citer la figure du philosophe Euphratès, dans le portrait qu’en donne une \teiHi[rend={italic}]{Lettre} (I, 10) de Pline le Jeune. Rappelons que, d’après l’historiographie, des philosophes ont été expulsés de Rome par Vespasien en 71 et de nouveau par Domitien en 89-90 et 93-94\teiNote[xmlid={ftn110},n={35},place={foot}]{\teiP[xmlid={p-102}]{Sur ces épisodes complexes, voir Yann Rivière, « L’exil des mages et des sages. Un empire sans philosophes ? (I\teiHi[rend={sup}]{er} siècle ap. J.-C.) », dans Pierre Vesperini (dir.), \teiHi[rend={italic}]{op. cit.}, p. 265-352 ; Clément Bady, « L’expulsion des philosophes de 93-94 p. C. Philosophie et sociabilité aristocratique dans la Rome des Flaviens », \teiHi[rend={italic}]{REA}, vol. 122, n\teiHi[rend={sup}]{o} 1, 2020, p. 107-125.}}, puis sont revenus sous Nerva (qui devient empereur en 96\teiNote[xmlid={ftn111},n={36},place={foot}]{\teiP[xmlid={p-103}]{Pour Hubert Zehnacker (Pline le Jeune, \teiHi[rend={italic}]{Lettres}, I, 10, Paris, Les Belles Lettres, « CUF », 2009, n. 1 \teiHi[rend={italic}]{ad loc}.), l’éloge d’Euphratès a pour but de justifier la mesure de Nerva. }}). La lettre de Pline le Jeune (écrite peut-être en 98, sous Nerva ou au tout début du principat de Trajan\teiNote[xmlid={ftn112},n={37},place={foot}]{\teiP[xmlid={p-104}]{Adrian N. Sherwin-White (\teiHi[rend={italic}]{The Letters of Pliny. A Historical and Social Commentary}, Oxford, Clarendon Press, 1966, p. 108) pense que la lettre est écrite après janvier 1998.}}) prend donc place dans un contexte où faire l’éloge d’un philosophe est un signe de la liberté politique retrouvée, liberté qui permet l’épanouissement des « études libérales » (\teiHi[rend={italic}]{liberalia studia}, I, 10, 1). Euphratès est né en Orient et a été un élève de Musonius Rufus, chevalier romain et philosophe stoïcien\teiNote[xmlid={ftn113},n={38},place={foot}]{\teiP[xmlid={p-105}]{Voir Paul Zanker, \teiHi[rend={italic}]{op. cit}., p. 257 et suiv. ; Patrick Robiano, dans Richard Goulet (dir.),\teiHi[rend={italic}]{ Dictionnaire des philosophes antiques, }Paris, Éditions du CNRS, 2000, vol. III, p. 337-342, entrée « Euphratès » (Vespasien ne le renvoie pas de Rome, d’après Philostrate, \teiHi[rend={italic}]{Vie d’Apollonios de Tyane}, V, 37). Sur Musonius Rufus, voir Paul Zanker, \teiHi[rend={italic}]{op. cit}., p. 258 et 259 ; Marie-Odile Goulet-Cazé, dans Richard Goulet (dir.),\teiHi[rend={italic}]{ op. cit}., 2005, vol. IV, p. 555-572, entrée « Musonius Rufus » ; Valéry Laurand\teiHi[rend={italic}]{, Stoïcisme et lien social : enquête autour de Musonius Rufus}, Paris, Classiques Garnier, 2014 (en particulier l’introduction).}}. D’après le témoignage de la \teiHi[rend={italic}]{Vie d’Apollonios de Tyane} de Philostrate et de lettres attribuées à Apollonios, Euphratès a été considéré comme un faux philosophe, un sophiste cupide. Ces critiques sont dues à la rivalité entre lui et Apollonios, qui apparaît comme une rivalité entre deux écoles, le stoïcisme et le pythagorisme\teiNote[xmlid={ftn114},n={39},place={foot}]{\teiP[xmlid={p-106}]{Voir Patrick Robiano, \teiHi[rend={italic}]{ibid}.}}. }
              \teiP[xmlid={p20-2}]{À l’opposé de ces attaques attribuées à Apollonios, la lettre de Pline, qu’il a adressée à Attius Clemens (inconnu en dehors de cette correspondance\teiNote[xmlid={ftn115},n={40},place={foot}]{\teiP[xmlid={p-107}]{Voir Michèle Ducos, dans Richard Goulet (dir.),\teiHi[rend={italic}]{ op. cit}., 1994, vol. II, p. 426, entrée « Clemens (Attius-) ».}}) pour le convaincre d’écouter Euphratès, est un éloge de celui-ci, qualifié de \teiHi[rend={italic}]{philosophus }(I, 10, 2), et de sage (\teiHi[rend={italic}]{sapiens}, § 4). Pline (§ 2) explique avoir connu Euphratès quand, tout jeune homme, il était tribun militaire en Syrie. Il vante d’abord son affabilité (Euphratès est \teiHi[rend={italic}]{obuius}), son caractère accessible, ouvert à tous (\teiHi[rend={italic}]{expositus}) et plein de l’\teiHi[rend={italic}]{humanitas }(bonté, humanité) qu’il enseigne (§ 2). Pline fait un peu plus loin (§ 5) l’éloge de son enseignement et de son style, qui convainc même ses adversaires. Son physique est aussi séduisant que sa parole :}
              \begin{teiCit}[xmlid={cit06}]

                \begin{teiQuote}
§ 6. Avec cela une haute taille \lbrack{}\teiHi[rend={italic}]{proceritas corporis}\rbrack{}, un beau visage \lbrack{}\teiHi[rend={italic}]{decora facies}\rbrack{}, des cheveux longs \lbrack{}\teiHi[rend={italic}]{demissus capillus}\rbrack{}, une barbe immense et blanche \lbrack{}\teiHi[rend={italic}]{ingens et barba cana}\rbrack{}, des caractéristiques qui peuvent sembler fortuites et sans importance, mais qui lui valent pourtant beaucoup de respect \lbrack{}\teiHi[rend={italic}]{plurimum uenerationis}\rbrack{}. § 7. Dans le soin de lui-même, rien d’hirsute \lbrack{}\teiHi[rend={italic}]{Nullus horror in cultu}\rbrack{}, rien d’austère \lbrack{}\teiHi[rend={italic}]{nulla tristitia}\rbrack{}, mais beaucoup de sérieux \lbrack{}\teiHi[rend={italic}]{multum seueritatis}\rbrack{}\teiHi[rend={italic}]{ }; son abord inspire le respect, pas la peur \lbrack{}\teiHi[rend={italic}]{reuearis occursum, non reformides}\rbrack{}. La pureté \lbrack{}\teiHi[rend={italic}]{sanctitas}\rbrack{} de sa vie est très grande, à l’égal de son amabilité \lbrack{}\teiHi[rend={italic}]{comitas}\rbrack{} ; il fait la chasse aux défauts, non aux hommes, et il ne châtie pas ceux qui font erreur, mais il les corrige.
\end{teiQuote}
              
\end{teiCit}
              \teiP[xmlid={p21-2}]{Pline ajoute (§ 8) qu’Euphratès a pour beau-père un Romain de bonne réputation, Pompeius Julianus, qui est « le premier de sa province », ce qui est une façon de souligner la bonne intégration du philosophe dans la société romaine\teiNote[xmlid={ftn116},n={41},place={foot}]{\teiP[xmlid={p-108}]{Sur cet aspect, voir Pierre Grimal, « Deux figures de la \teiHi[rend={italic}]{Correspondance} de Pline : le philosophe Euphratès et le rhéteur Isée », \teiHi[rend={italic}]{Latomus}, n\teiHi[rend={sup}]{o} 14, 1955, p. 370-383. Adrian N. Sherwin-White\teiHi[rend={italic}]{, op. cit}., p. 110, note que ce personnage n’est pas autrement connu et qu’il n’était pas nécessairement d’un rang social élevé.}}. }
              \teiP[xmlid={p22-2}]{D’après cette lettre, Euphratès a certains des attributs physiques traditionnels du philosophe : les cheveux longs et la longue barbe. Pour Paul Zanker, ce sont là les caractéristiques des philosophes qu’il appelle « Charismatiques », à cause de leur aspect et de leur pouvoir d’attraction, comme Dion de Pruse, Ælius Aristide et Pérégrinus\teiNote[xmlid={ftn117},n={42},place={foot}]{\teiP[xmlid={p-109}]{Paul Zanker, \teiHi[rend={italic}]{op. cit}., p. 261-263.}}. Cette caractéristique peut aussi renvoyer aux Pythagoriciens (dont Apollonios de Tyane se réclame). Il y a en tout cas chez les philosophes de cette période une préoccupation particulière pour la barbe et les cheveux longs, qui sont des signes identitaires\teiNote[xmlid={ftn118},n={43},place={foot}]{\teiP[xmlid={p-110}]{Voir Paul Zanker, \teiHi[rend={italic}]{op. cit}., p. 108-113 sur la barbe de Chrysippe et celle des philosophes du III\teiHi[rend={sup}]{e} s. av. J.-C. et p. 259 et 260 sur la période de l’empire romain ; Bruno Sudan, art. cité, p. 175 et 176 : Musonius et Dion Chrysostome sont les auteurs d’ouvrages intitulés respectivement \teiHi[rend={italic}]{Sur la taille des cheveux} et \teiHi[rend={italic}]{Éloge de la chevelure}. Sur les cheveux longs des Pythagoriciens, voir Philostrate, \teiHi[rend={italic}]{Vie d’Apollonios de Tyane}, I, 8 et 32 ; VII, 15.}}. Domitien, qui chasse les philosophes de Rome, fait couper (ou raser) la barbe et la chevelure d’Apollonios de Tyane\teiNote[xmlid={ftn119},n={44},place={foot}]{\teiP[xmlid={p-111}]{Philostrate, \teiHi[rend={italic}]{Vie d’Apollonios de Tyane}, VII, 34-36. }}. Notons cependant qu’il est souvent difficile d’évaluer la longueur de la barbe, même si, ici, l’adjectif \teiHi[rend={italic}]{ingens} marque une barbe vraiment visible, probablement longue ; il en va de même pour les cheveux\teiNote[xmlid={ftn120},n={45},place={foot}]{\teiP[xmlid={p-112}]{Paul Zanker, \teiHi[rend={italic}]{op. cit}., p. 261, parle de cheveux allant jusqu’aux épaules.}}. D’autre part, il est possible que cette barbe et ces cheveux longs soient aussi un indice de grécité, à une époque où les empereurs romains et les membres de l’élite sont représentés imberbes et avec des cheveux courts. Quoi qu’il en soit, d’après Pline, dans le cas d’Euphratès, les apparences ne sont pas trompeuses.}
              \teiP[xmlid={p23-2}]{Pline précise que ces caractéristiques peuvent sembler inutiles, comme pour signifier qu’elles sont moins importantes que sa sagesse, mais il souligne par là même que cet aspect, qui va à l’inverse de l’\teiHi[rend={italic}]{habitus} (aspect et mise) du citoyen romain homme de bien, n’est pas tant l’attribut habituel des philosophes qu’une apparence qui produit, ici, un profond respect : les cheveux longs et la barbe blanche sont connotés positivement, parce que le visage est qualifié de « beau » et « convenable » à la fois (\teiHi[rend={italic}]{decora facies}), et parce que ces caractéristiques inspirent un respect mêlé d’admiration, \teiHi[rend={italic}]{ueneratio}, que les hommes ont pour les dieux, les magistrats ou les empereurs\teiNote[xmlid={ftn121},n={46},place={foot}]{\teiP[xmlid={p-113}]{Voir \teiHi[rend={italic}]{Oxford Latin Dictionary}, Oxford, Clarendon Press, 1968, entrée « \teiHi[rend={italic}]{Ueneratio }».}}. L’âge d’Euphratès, qui peut avoir plus de 60 ans quand la lettre est écrite, autorise la blancheur de la barbe (et des cheveux probablement). En outre, le physique d’Euphratès est dans son ensemble en accord avec ses qualités morales et sociales, « pédagogiques » même (\teiHi[rend={italic}]{humanitas}, \teiHi[rend={italic}]{uitae sanctitas}, \teiHi[rend={italic}]{comitas}). Pline précise encore qu’il ne négligeait pas le soin de son corps et de sa mise (\teiHi[rend={italic}]{nullus horror in cultu}), ce qui l’oppose aux cyniques, à leurs cheveux et barbes sales blâmés par Martial. D’autre part, la \teiHi[rend={italic}]{tristitia}, qu’on reproche généralement aux philosophes stoïciens\teiNote[xmlid={ftn122},n={47},place={foot}]{\teiP[xmlid={p-114}]{Nous ne croyons pas que cette remarque sur l’expression d’Euphratès le distingue des cyniques, comme l’indique Paul Zanker, \teiHi[rend={italic}]{op. cit}., p. 260 et 261. Sur la \teiHi[rend={italic}]{tristitia} attribuée à ceux qui se réclament du stoïcisme, voir Cicéron, \teiHi[rend={italic}]{De Finibus}, IV, 79 : « C’est de cette humeur difficile et de cette rudesse \lbrack{}\teiHi[rend={italic}]{illorum tristitiam atque asperitatem}\rbrack{} que Panétius \lbrack{}de Rhodes, représentant du Moyen Stoïcisme\rbrack{} ne put s’accommoder ; voilà pourquoi il critiqua l’âpreté \lbrack{}\teiHi[rend={italic}]{acerbitatem}\rbrack{} de leurs idées et les épines \lbrack{}\teiHi[rend={italic}]{spinas}\rbrack{} de leur argumentation et se montra, dans la doctrine plus moelleux \lbrack{}\teiHi[rend={italic}]{mitior}\rbrack{}, dans la discussion plus clair \lbrack{}…\rbrack{} » (Jules Martha trad., Paris, Les Belles Lettres, « CUF », 1930). Voir aussi Tacite, \teiHi[rend={italic}]{Annales}, XVI, 22, 4, sur ceux qui soutiennent le sénateur Thraséa Paetus et qui sont comme lui \teiHi[rend={italic}]{rigidi} et \teiHi[rend={italic}]{tristes}.}}, est remplacée ici par une qualité bien romaine, la \teiHi[rend={italic}]{seueritas}. Enfin, les atouts de l’apparence physique et de la mise d’Euphratès, qui à la fois lui donnent un aspect de philosophe et participent de la dignité de son apparence, se retrouvent pour certains d’entre eux dans le portrait élogieux de Trajan que fait ce même Pline le Jeune\teiNote[xmlid={ftn123},n={48},place={foot}]{\teiP[xmlid={p-115}]{Voir Lilja Saara, « Descriptions of Human Appearance in Pliny’s Letters », \teiHi[rend={italic}]{Arctos}, n\teiHi[rend={sup}]{o} 12, 1978, p. 55-62 (p. 58 et 59 sur cette lettre).}} :}
              \begin{teiCit}[xmlid={cit07}]

                \begin{teiQuote}
Sa robustesse \lbrack{}\teiHi[rend={italic}]{firmitas}\rbrack{}, sa grande taille \lbrack{}\teiHi[rend={italic}]{proceritas corporis}\rbrack{}, la noblesse de sa tête \lbrack{}\teiHi[rend={italic}]{honor capitis}\rbrack{} et la beauté digne de son visage \lbrack{}\teiHi[rend={italic}]{dignitas oris}\rbrack{}, outre cela, son âge mûr mais qui ne le courbe pas \lbrack{}\teiHi[rend={italic}]{aetatis indeflexa maturitas}\rbrack{}, sa chevelure ornée des marques précoces de la vieillesse \lbrack{}\teiHi[rend={italic}]{festinatis senectutis insignibus }(…)\teiHi[rend={italic}]{ ornata caesaries}\rbrack{}, ce qui est une sorte de cadeau des dieux, destiné à accroître sa majesté \lbrack{}\teiHi[rend={italic}]{maiestatem}\rbrack{}, tout cela ne révèle-t-il pas \lbrack{}\teiHi[rend={italic}]{ostentant}\rbrack{} déjà à tous points de vue un prince\teiNote[xmlid={ftn124},n={49},place={foot}]{\teiP[xmlid={p-116}]{Pline le Jeune, \teiHi[rend={italic}]{Panégyrique de Trajan}, 4, 7. Voir aussi, à propos de Claude, Suétone, \teiHi[rend={italic}]{Claude}, 30, 1 : « L’autorité et la beauté digne de son aspect physique \lbrack{}\teiHi[rend={italic}]{Auctoritas dignitasque formae}\rbrack{} ne lui manquaient pas, en tout cas quand il était assis ou debout et surtout quand il était au repos, car il était élancé sans être maigre et il avait de beaux cheveux blancs \lbrack{}\teiHi[rend={italic}]{specie canitique pulchra}\rbrack{}, avec un cou bien plein \lbrack{}…\rbrack{}. »}} ?
\end{teiQuote}
              
\end{teiCit}
              \teiP[xmlid={p24-2}]{Pline ne parle pas de la façon dont Euphratès est vêtu, pour mieux souligner d’autres aspects, et peut-être parce que celui-ci ne portait pas le traditionnel manteau de philosophe. En effet, d’après une lettre attribuée à Apollonios de Tyane, Euphratès serait un « nouveau-riche » (\teiHi[rend={italic}]{Épîtres}, 6), qui fait montre de sa richesse et exhibe ses vêtements, parce qu’il a renoncé au \teiHi[rend={italic}]{tribôn }qu’il portait à ses débuts, reniant ainsi l’idéal de Zénon de Citium, le fondateur du stoïcisme (\teiHi[rend={italic}]{Épîtres}, 3). À l’inverse, on peut considérer que l’abandon du manteau de philosophe est une romanisation bienvenue de l’apparence. C’est ce que suggère un passage d’Épictète à son propos :}
              \begin{teiCit}[xmlid={cit08}]

                \begin{teiQuote}
§ 15. Mais les prétendus philosophes \lbrack{}\teiHi[rend={italic}]{hoi kaloumenoi philosophoi}\rbrack{} briguent, eux aussi, le titre par des moyens qui sont indifférents : à peine ont-ils pris le manteau \lbrack{}\teiHi[rend={italic}]{tribôna}\rbrack{} et laissé croître leur barbe \lbrack{}\teiHi[rend={italic}]{pôgôna}\rbrack{} qu’ils proclament : « Je suis philosophe ». \lbrack{}…\rbrack{} § 17. Aussi, Euphratès avait bien raison de dire : « Pendant longtemps j’ai essayé de cacher que j’étais philosophe, et cela m’a servi. Premièrement, je savais que tout ce que je faisais de bien n’était pas dû à ceux qui me voyaient, mais à moi-même : c’est pour moi-même que je mangeais décemment \lbrack{}\teiHi[rend={italic}]{kalôs}\rbrack{}, que je gardais la modestie dans mon regard, dans ma démarche ; tout cela pour moi-même et pour Dieu\teiNote[xmlid={ftn125},n={50},place={foot}]{\teiP[xmlid={p-117}]{\teiHi[rend={italic}]{Entretiens}, IV, 8, 15-17 (Joseph Souilhé et Amand Jagu trad., Paris, Les Belles Lettres, « CUF », 1963).}}.
\end{teiQuote}
              
\end{teiCit}
              \teiP[xmlid={p25}]{Euphratès a renoncé au manteau (mais pas à la barbe, d’après la lettre de Pline) et il manifeste un souci de convenance dans ses activités quotidiennes, afin de ne pas se comporter de façon choquante dans le milieu où il vit : il dit le faire pour lui-même et pour Dieu, mais ses repas, ses regards, sa démarche sont visibles par les autres hommes – le philosophe n’est pas seulement sous son propre regard et sous celui de la divinité des stoïciens. Or les qualités de décence dans la façon de manger, de regarder, de marcher peuvent être considérées comme des qualités proprement romaines. }
              \teiP[xmlid={p26}]{Autre figure, dans un portrait plus développé qui articule plusieurs de ces traits, celui de Claranus, dans l’une des \teiHi[rend={italic}]{Lettres à Lucilius} (66) de Sénèque\teiNote[xmlid={ftn126},n={51},place={foot}]{\teiP[xmlid={p-118}]{Sur cette lettre, voir Brad Inwood, \teiHi[rend={italic}]{Seneca. Selected Philosophical Letters}, Oxford University Press, 2007, p. 155-181 ; Francesca Romana Berno, art. cité\teiHi[rend={italic}]{.}}}. Cette épître porte sur la question du classement des biens selon les stoïciens\teiNote[xmlid={ftn127},n={52},place={foot}]{\teiP[xmlid={p-119}]{Les biens du troisième ordre concernent précisément l’apparence : « une démarche modeste, une expression de visage qu’on a voulu calme et qui est honnête, des gestes qui conviennent à un homme sage \lbrack{}\teiHi[rend={italic}]{modestus incessus et compositus ac probus uultus et conueniens prudenti uiro gestus}\rbrack{} » (Sénèque, \teiHi[rend={italic}]{Lettres à Lucilius}, 66, 5). }}. Sénèque y réaffirme que la vertu est le souverain bien, indépendamment du niveau de richesse et de l’aspect physique\teiNote[xmlid={ftn128},n={53},place={foot}]{\teiP[xmlid={p-120}]{Voir Francesca Romana Berno, art. cité, p. 272 et 273 : « Le corps et sa condition (donc, tous les biens et maux corporels) relèvent ensemble de la catégorie de l’apparence, donc de l’indifférence sur le plan des valeurs ; au contraire, l’âme et son bien (la vertu) constituent l’unique réalité authentique. De la sorte, Sénèque réintroduit dans le discours la taxinomie traditionnelle des biens \lbrack{}…\rbrack{}. »}}. La description de Claranus, qui n’est pas autrement connu\teiNote[xmlid={ftn129},n={54},place={foot}]{\teiP[xmlid={p-121}]{Voir Richard Goulet, dans Richard Goulet (dir.), \teiHi[rend={italic}]{op. cit.}, 1994, vol. II, p. 400 et 401, \teiHi[rend={italic}]{s. v. }Claranus. }}, ouvre la lettre et sert de point de départ à la démonstration. La voici : }
              \begin{teiCit}[xmlid={cit09}]

                \begin{teiQuote}
J’ai revu Claranus mon camarade d’études, après de nombreuses années. Tu n’attends pas, je pense, que j’ajoute : un vieillard. Il a en tout cas une âme jeune et vigoureuse \lbrack{}\teiHi[rend={italic}]{uiridem animo ac uigentem}\rbrack{}, qui lutte contre son corps chétif \lbrack{}\teiHi[rend={italic}]{cum corpusculo suo}\rbrack{}. En effet, la nature s’est conduite injustement avec lui et elle a mal logé une telle âme \lbrack{}\teiHi[rend={italic}]{talem animum male collocauit}\rbrack{}. À moins peut-être qu’elle n’ait voulu nous montrer \lbrack{}\teiHi[rend={italic}]{ostendere}\rbrack{} précisément que l’esprit le plus courageux, le plus heureux peut se cacher sous n’importe quelle peau. \lbrack{}…\rbrack{} § 2. Il s’est trompé, je crois, celui qui a dit : « le mérite est plus beau quand il vient d’un beau corps » \lbrack{}\teiHi[rend={italic}]{gratior et pulchro ueniens e corpore uirtus}\teiNote[xmlid={ftn130},n={55},place={foot}]{\teiP[xmlid={p-122}]{Virgile, \teiHi[rend={italic}]{Énéide}, V, 344 (à propos d’Euryale, pendant les courses de chevaux qui font partie des Jeux funéraires donnés en l’honneur d’Anchise).}}\rbrack{}. Le mérite n’a besoin d’aucun ornement honorable \lbrack{}\teiHi[rend={italic}]{ullo honestamento}\rbrack{} : il est à lui seul un honneur considérable et il consacre le corps où il se trouve \lbrack{}\teiHi[rend={italic}]{ipsa magnum sui decus est et corpus suum consecrat}\rbrack{}. Quoi qu’il en soit, je commence à regarder \lbrack{}\teiHi[rend={italic}]{intueri}\rbrack{} autrement mon cher Claranus : il me semble aussi beau et droit de corps qu’il l’est d’âme \lbrack{}\teiHi[rend={italic}]{formosus mihi uidetur et tam rectus corpore quam est animo}\rbrack{}. § 3. Un grand homme peut sortir d’une cabane ; une âme belle et grande peut sortir même d’un corps laid et humble \lbrack{}\teiHi[rend={italic}]{potest et ex deformi humilique corpusculo formosus animus ac magnus}\rbrack{}. \lbrack{}…\rbrack{} Claranus a été, je crois, mis au monde pour nous faire comprendre par son exemple \lbrack{}\teiHi[rend={italic}]{exemplar}\rbrack{} que la laideur du corps ne défigure pas l’âme, mais que c’est de la beauté de l’âme que le corps reçoit un ornement \lbrack{}\teiHi[rend={italic}]{non deformitate corporis foedari animum, sed pulchritudine animi corpus ornari}\teiNote[xmlid={ftn131},n={56},place={foot}]{\teiP[xmlid={p-123}]{Sénèque, \teiHi[rend={italic}]{Lettres à Lucilius}, 66, 1-4.}}\rbrack{}.
\end{teiQuote}
              
\end{teiCit}
              \teiP[xmlid={p27}]{Ce passage sur Claranus repose sur la distinction entre \teiHi[rend={italic}]{corpus} et \teiHi[rend={italic}]{animus} qui ne sont pas en accord : le personnage est âgé (il a plus de 60 ans\teiNote[xmlid={ftn132},n={57},place={foot}]{\teiP[xmlid={p-124}]{Sénèque est né vers 1 av. J.-C., il écrit les \teiHi[rend={italic}]{Lettres à Lucilius} à partir de 62 et, si Claranus a été son condisciple, ils doivent être à peu près du même âge.}}), il est fragile (il a un petit corps : \teiHi[rend={italic}]{corpusculum}), alors qu’il est « vert et vigoureux » en esprit, dans l’âme (\teiHi[rend={italic}]{animo}\teiNote[xmlid={ftn133},n={58},place={foot}]{\teiP[xmlid={p-125}]{Peut-être y a-t-il ici un renvoi à Chrysippe, dont Diogène Laërce (VII, 182 ;\teiHi[rend={italic}]{ Diogenes Laertius : Vitae philosophorum libri}, Miroslav Marcovich éd., Stuttgart-Leibniz, Teubner, 1999) dit qu’il avait un « petit corps ordinaire, comme le montre sa statue du Céramique ». Sur ce texte, voir Paul Zanker, \teiHi[rend={italic}]{op. cit}., p. 135.}}). C’est sans doute grâce à sa sagesse et à la philosophie que Claranus a cette force d’âme\teiNote[xmlid={ftn134},n={59},place={foot}]{\teiP[xmlid={p-126}]{Francesca Romana Berno (art. cité, p. 268) renvoie à la \teiHi[rend={italic}]{Lettre} 30 sur Aufidius Bassus, où l’on trouve aussi l’idée de la lutte entre l’\teiHi[rend={italic}]{animus }et un corps qui a toujours été « faible et desséché », qui est usé par la vieillesse, mais plein d’énergie (\teiHi[rend={italic}]{alacer animo est}) grâce à la philosophie (§ 3) ; Sénèque voulait voir si la diminution de sa force physique (\teiHi[rend={italic}]{corporis uires}) diminuerait la vigueur de son âme (\teiHi[rend={italic}]{uigor animi}) ; au contraire, celle-ci grandit (§ 13).}}. On peut dire que la beauté de Claranus est vitale, dans la mesure où sa sagesse, qui lui donne de la beauté, est en même temps force de vie.}
              \teiP[xmlid={p28}]{On retrouve dans l’évocation de ce personnage l’idée du divorce entre l’apparence physique et l’\teiHi[rend={italic}]{animus} évoquée précédemment, avec une opposition qui est autant entre extérieur et intérieur qu’entre ce qui est visible et ce qui ne l’est pas au premier abord ou aux yeux de qui n’est pas philosophe : ici, c’est l’\teiHi[rend={italic}]{ingenium} (l’esprit, le tempérament, le talent) courageux et parfaitement heureux (\teiHi[rend={italic}]{beatissimum}) qui se dissimule (\teiHi[rend={italic}]{latere}) sous une peau ordinaire, la peau du corps (\teiHi[rend={italic}]{cutis}). Cette image est prolongée par l’analogie entre la formule disant qu’un grand homme peut sortir d’une cabane et celle de l’âme logée dans le corps, exprimée avec le verbe \teiHi[rend={italic}]{collocare}, ce qui constitue là aussi une métaphore spatiale\teiNote[xmlid={ftn135},n={60},place={foot}]{\teiP[xmlid={p-127}]{Voir Brad Inwood, \teiHi[rend={italic}]{op. cit.}, p. 157 pour des rapprochements textuels à propos de ces deux images. Sur le fait que Sénèque, après Cicéron, construise l’idée de l’intériorité comme espace de l’\teiHi[rend={italic}]{animus}, voir Marion Bourbon, « De l’“intérieur” à l’intime : la construction sénéquienne de l’intériorité », dans Géraldine Puccini (dir.), \teiHi[rend={italic}]{L’intime de l’Antiquité à nos jours. }I. \teiHi[rend={italic}]{L’espace de l’intime}, Pessac, PUB, 2019, p. 93-106. Il est bien question, en effet, dans d’autres lettres de Sénèque, de « regarder à l’intérieur de sa poitrine/son cœur » (83, 1 : \teiHi[rend={italic}]{in pectus intimum introspicere}), de regarder à l’intérieur de soi (80, 10 : \teiHi[rend={italic}]{intus te ipse considera}). Quelques images rencontrées précédemment, pas seulement dans des textes philosophiques, par exemple celle du corps comme « domicile de \teiHi[rend={italic}]{l’animus} » chez Velléius Paterculus à propos de Vatinius, semblent précéder le concept d’intériorité chez Sénèque.}}. }
              \teiP[xmlid={p29}]{La citation de Virgile montre bien que cette figure prend le contre-pied des valeurs épiques, selon lesquelles le héros est à la fois jeune, beau et vaillant, mais aussi de la tradition romaine qui associe prestance physique et vertu morale. Plus encore, la prise en compte de la qualité de l’\teiHi[rend={italic}]{animus} et de l’\teiHi[rend={italic}]{ingenium} modifie l’appréciation esthétique du corps : le corps est vu comme doté de la beauté et de la rectitude de l’âme. L’adjectif \teiHi[rend={italic}]{rectus}, droit, est associé, de façon générale, à la beauté (par opposition à un corps « tordu », de travers : \teiHi[rend={italic}]{prauus}), et aussi à la moralité, à un comportement social et moral convenable. On note le jeu étymologique \teiHi[rend={italic}]{formosus}/\teiHi[rend={italic}]{deformis} : le corps est \teiHi[rend={italic}]{deformis}, mais l’\teiHi[rend={italic}]{animus} est \teiHi[rend={italic}]{formosus}\teiNote[xmlid={ftn136},n={61},place={foot}]{\teiP[xmlid={p-128}]{Voir Francesca Romana Berno, art. cité, p. 272 : Sénèque rapporte à l’âme des adjectifs qui concernent le corps.}}. La conclusion de Sénèque est que la beauté de l’âme est un ornement pour le corps\teiNote[xmlid={ftn137},n={62},place={foot}]{\teiP[xmlid={p-129}]{Voir aussi Pline le Jeune, \teiHi[rend={italic}]{Lettres}, I, 10, 5 sur Titius Aristo : sa « grandeur d’âme (\teiHi[rend={italic}]{magnitudo animi}) ajoute un ornement (\teiHi[rend={italic}]{ornat}) » à sa frugalité et à la modestie de son train de vie.}}, comme les vêtements, les bijoux ou les insignes du pouvoir peuvent en être les ornements directement visibles. }
              \teiP[xmlid={p30}]{Une autre lettre de Sénèque (\teiHi[rend={italic}]{Lettres à Lucilius}, 115) vante de même la beauté de l’âme, qui n’est visible qu’au regard d’une autre âme :}
              \begin{teiCit}[xmlid={cit10}]

                \begin{teiQuote}
§ 3. S’il pouvait nous être donné de pénétrer du regard \lbrack{}\teiHi[rend={italic}]{inspicere}\rbrack{} l’âme de l’homme de bien \lbrack{}\teiHi[rend={italic}]{animum boni uiri}\rbrack{}, quel air de beauté en elle \lbrack{}\teiHi[rend={italic}]{o quam pulchram faciem}\rbrack{}, de sainteté, quel effet lumineux de majesté paisible nous verrions \lbrack{}\teiHi[rend={italic}]{uideremus}\rbrack{}. \lbrack{}…\rbrack{} § 6. \lbrack{}…\rbrack{} si nous formons le propos de libérer de tout obstacle le regard de notre âme \lbrack{}\teiHi[rend={italic}]{acies animi}\rbrack{}, alors nous pourrons voir clairement \lbrack{}\teiHi[rend={italic}]{perspicere}\rbrack{} la vertu, même si elle est dans un corps en ruine \lbrack{}\teiHi[rend={italic}]{etiam obrutam corpore}\rbrack{}, même sous le manteau de la pauvreté \lbrack{}\teiHi[rend={italic}]{etiam paupertate opposita}\rbrack{}, à travers l’abjection et l’infamie \lbrack{}\teiHi[rend={italic}]{etiam humilitate et infamia obiacentibus}\rbrack{}. § 7. Oui, nous apercevrons \lbrack{}\teiHi[rend={italic}]{Cernemus}\rbrack{} sa beauté \lbrack{}\teiHi[rend={italic}]{pulchritudinem illam}\rbrack{}, fût-elle présentée sous les dehors les plus repoussants \lbrack{}\teiHi[rend={italic}]{quamuis sordido obiectam}\teiNote[xmlid={ftn138},n={63},place={foot}]{\teiP[xmlid={p-130}]{Sénèque, \teiHi[rend={italic}]{Lettres à Lucilius}, 115, 3 et 6-7 (Henri Noblot trad., Paris, Les Belles Lettres, « CUF », 1964, modifiée).}}\rbrack{} \lbrack{}…\rbrack{}.
\end{teiQuote}
              
\end{teiCit}
              \teiP[xmlid={p31}]{Le délabrement du corps et la pauvreté de l’apparence ne doivent pas empêcher de voir la beauté de l’âme. Mais cela n’est possible que si l’on cherche à regarder au-delà des apparences, comme le souligne ici le lexique de la vision. Ainsi, le portrait de Claranus, qui a valeur de modèle, \teiHi[rend={italic}]{exemplar}\teiNote[xmlid={ftn139},n={64},place={foot}]{\teiP[xmlid={p-131}]{Sur \teiHi[rend={italic}]{exemplar} au § 4, voir Élisabeth Gavoille, « La force de l’exemple dans les \teiHi[rend={italic}]{Lettres à Lucilius} », dans É. Gavoille et F. Guillaumont (éd.), \teiHi[rend={italic}]{op. cit.}, p. 283-300 : « \lbrack{}…\rbrack{} Claranus vaut comme vivante leçon, modèle supérieur qui participe manifestement d’un plan divin (\teiHi[rend={italic}]{in exemplar editus}). » (p. 288).}}, illustre un paradoxe typiquement stoïcien, dont on peut trouver la formulation, teintée d’ironie, dans un discours de Cicéron, le \teiHi[rend={italic}]{Pro Murena}, où Cicéron critique Caton, accusateur de Muréna que lui-même défend. Cicéron cite les \teiHi[rend={italic}]{sententiae} (dogmes) et les \teiHi[rend={italic}]{praecepta} (règles morales) de Zénon : « Seul le sage, fût-il le plus contrefait des hommes, est beau \lbrack{}\teiHi[rend={italic}]{solos sapientes esse, si distortissimi sint, formosos}\rbrack{} ; fût-il le dernier des gueux, il est riche ; fût-il de condition servile, il est roi\teiNote[xmlid={ftn140},n={65},place={foot}]{\teiP[xmlid={p-132}]{Cicéron, \teiHi[rend={italic}]{Pro Murena}, 61 (André Boulanger trad., Paris, Les Belles Lettres, « CUF », 1943).}}. » Ce paradoxe trouve une explication dans le portrait que ce même Caton fait du sage stoïcien dans le traité\teiHi[rend={italic}]{ De Finibus }de Cicéron : « \lbrack{}…\rbrack{} oui, il méritera d’être appelé beau \lbrack{}\teiHi[rend={italic}]{pulcher}\rbrack{}, les traits de l’âme étant plus beaux que ceux du corps \lbrack{}\teiHi[rend={italic}]{animi etiam liniamenta sunt pulchriora quam corporis}\teiNote[xmlid={ftn141},n={66},place={foot}]{\teiP[xmlid={p-133}]{Cicéron, \teiHi[rend={italic}]{De Finibus}, III, 75 (Jules Martha trad., Paris, Les Belles Lettres, « CUF », 1930). Cicéron reprend plus loin cette idée, en la critiquant (\teiHi[rend={italic}]{idem}, IV, 74) : « Et puis, seul le sage est beau, seul il est libre, seul il est vraiment citoyen \lbrack{}\teiHi[rend={italic}]{Solum praeterea formosum, solum liberum, solum ciuem}\rbrack{} ; au lieu que tous les autres, c’est le contraire ; vous dites même d’eux qu’ils sont des insensés \lbrack{}\teiHi[rend={italic}]{insanos}\rbrack{}. C’est ce que les Grecs appellent des “paradoxes” \lbrack{}\teiHi[rend={italic}]{paradoxa}\rbrack{}, terme que nous pourrions rendre par “choses qui émerveillent” \lbrack{}\teiHi[rend={italic}]{admirabilia}\rbrack{}. Qu’ont-elles pourtant de si merveilleux \lbrack{}\teiHi[rend={italic}]{quid… admirationis}\rbrack{}, quand on les regarde de près ? » (\teiHi[rend={italic}]{ibid.}).}}\rbrack{} \lbrack{}…\rbrack{} ». Pour les stoïciens, il y a en effet un primat de l’âme vertueuse sur le corps\teiNote[xmlid={ftn142},n={67},place={foot}]{\teiP[xmlid={p-134}]{Voir Cicéron, \teiHi[rend={italic}]{De Finibus}, IV, 26. Cicéron demande « comment, en quel endroit, il se fait <néanmoins> que vous ayez tout à coup laissé le corps en route \lbrack{}\teiHi[rend={italic}]{quonam modo aut quo loco corpus subito deserueritis}\rbrack{} » (\teiHi[rend={italic}]{ibid.}). Voir cependant Valéry Laurand, \teiHi[rend={italic}]{op. cit}., p. 95 et suiv. sur la place donnée au corps et aux exercices du corps chez Musonius et chez Sénèque.}}. La suite de la lettre de Sénèque sur Claranus le dit clairement : « La vertu \lbrack{}\teiHi[rend={italic}]{uirtus}\rbrack{} n’est pas plus estimable dans un corps vigoureux et libre d’entraves \lbrack{}\teiHi[rend={italic}]{in corpore ualido ac libero}\rbrack{} que dans un corps infirme et garrotté \lbrack{}\teiHi[rend={italic}]{in morbido et uincto}\teiNote[xmlid={ftn143},n={68},place={foot}]{\teiP[xmlid={p-135}]{Sénèque, \teiHi[rend={italic}]{Lettres à Lucilius}, 66, 22. Voir aussi \teiHi[rend={italic}]{idem}, 24-25 ; 34 : les gens de bien sont tous égaux, quels que soient leur âge, leur aspect physique (beau : \teiHi[rend={italic}]{formosus}, ou laid : \teiHi[rend={italic}]{deformis}), leur degré de richesse et leur situation sociale.}}\rbrack{}. »}
              \teiP[xmlid={p32}]{Cependant, il faut aussi insister sur le fait que, pour Sénèque, le sage ne doit pas avoir une apparence contraire à ce que Cicéron appelle la \teiHi[rend={italic}]{dignitas}. Ce souci apparaît dans une autre \teiHi[rend={italic}]{Lettre à Lucilius} : }
              \begin{teiCit}[xmlid={cit11}]

                \begin{teiQuote}
Tenue grossière \lbrack{}\teiHi[rend={italic}]{Asperum cultum}\rbrack{}, cheveux non coupés, barbe trop négligée, haine déclarée de l’argenterie, matelas posé à même le sol, et toute autre chose qui vise à se faire remarquer par un moyen pervers, évite tout cela. \lbrack{}…\rbrack{} Au-dedans \lbrack{}\teiHi[rend={italic}]{Intus}\rbrack{}, que tout soit différent, mais que notre apparence soit commune à celle du peuple \lbrack{}\teiHi[rend={italic}]{frons populo nostra conueniat}\rbrack{}. § 3. Que notre toge ne soit pas éclatante, mais qu’elle ne soit pas sale non plus \lbrack{}\teiHi[rend={italic}]{Non splendeat toga, ne sordeat quidem}\teiNote[xmlid={ftn144},n={69},place={foot}]{\teiP[xmlid={p-136}]{Sénèque, \teiHi[rend={italic}]{Lettres à Lucilius}, 5, 2. Sur ce texte et d’autres du même type, voir Michel Blonski, \teiHi[rend={italic}]{Se nettoyer à Rome. }\teiHi[rend={ small-caps italic}]{ii}\teiHi[rend={italic sup}]{e}\teiHi[rend={italic}]{ s. avant J.-C.- }\teiHi[rend={ small-caps italic}]{ii}\teiHi[rend={italic sup}]{e}\teiHi[rend={italic}]{ s. après J.-C. Pratiques et enjeux}, Paris, Les Belles Lettres, 2014, p. 123-125 (qui renvoie notamment à Épictète, \teiHi[rend={italic}]{Entretiens}, IV, 11, 9-18 : à rebours d’une apparence négligée, la propreté distingue l’homme des animaux et elle est indispensable à sa vie en société). Sur la distinction entre l’apparence et la disposition d’esprit du philosophe stoïcien, Michael Trapp, art. cité, p. 366-368.}}\rbrack{} \lbrack{}…\rbrack{}.
\end{teiQuote}
              
\end{teiCit}
              \teiP[xmlid={p33}]{Sénèque critique ceux qui affectent de porter la tenue de philosophe en arborant une apparence physique, en particulier pileuse, proche de celle des cyniques. Cela ne sert qu’à attirer l’attention (\teiHi[rend={italic}]{notabilia}), dans la mise (\teiHi[rend={italic}]{habitus}) et le genre de vie (§ 1). Au contraire, le philosophe stoïcien, écrit-il, ne doit pas se distinguer des autres hommes par son apparence physique et vestimentaire ; c’est à l’intérieur (\teiHi[rend={italic}]{intus}) qu’il se distingue d’autrui\teiNote[xmlid={ftn145},n={70},place={foot}]{\teiP[xmlid={p-137}]{Sur ce point, voir Nicolas Davieau, « Montrer le corps : prouver le philosophe ? Le corps des philosophes dans la statuaire antique », \teiHi[rend={italic}]{Cahiers Mondes Anciens}, n\teiHi[rend={sup}]{o} 8, \teiHi[rend={italic}]{Montrer, démontrer. La validation des informations, des savoirs et des statuts dans l’Antiquité}, 2016, DOI : \teiRef[target={https://doi.org/10.4000/mondesanciens.1683}]{\teiHi[rend={underline}]{10.4000/mondesanciens.1683}}, § 24, qui renvoie à Lucien : « On peut en effet trouver chez Lucien (\teiHi[rend={italic}]{Hermotime}, 18) une description des stoïciens dans laquelle les soins du corps ne sont pas dédaignés : “Je voyais les stoïciens, modestes dans leur démarche, simples dans leurs vêtements, l’air réfléchi, la figure mâle, rasés presque tous jusqu’à la peau, rien qui annonçât la mollesse, rien qui trahît l’excès contraire et qui sentît le cynique, mais ce juste milieu que tout le monde regarde comme ce qu’il y a de meilleur.” En réalité, les Stoïciens incarnent ici une forme de juste milieu, entre l’excès de mollesse, attribuable aux épicuriens ou autres cyrénaïques, et l’excès cynique, tout d’âpreté \lbrack{}…\rbrack{} ».}}. Il ne faut pas s’employer à faire le contraire de la foule (\teiHi[rend={italic}]{uulgus}), sinon, les philosophes feront fuir les autres, dont ils veulent pourtant l’amélioration morale.}
              \teiP[xmlid={p34}]{Sénèque poursuit en disant qu’il est contre nature de préférer la saleté corporelle (\teiHi[rend={italic}]{squalor}) à une toilette simple (\teiHi[rend={italic}]{faciles munditiae}) et d’avoir du goût pour des aliments dégoûtants, alors que les stoïciens souhaitent vivre conformément à la nature, et il explique que la sobriété (\teiHi[rend={italic}]{frugalitas}), obligatoire chez le philosophe, peut ne pas être « mal peignée » : \teiHi[rend={italic}]{incompta}\teiNote[xmlid={ftn146},n={71},place={foot}]{\teiP[xmlid={p-138}]{Sénèque, \teiHi[rend={italic}]{Épîtres à Lucilius}, 5, 4.}}. Cet usage du participe \teiHi[rend={italic}]{incomptus} est figuré, mais le sens propre est présent aussi dans ce terme comme l’atteste le début de la lettre. Le soin des cheveux fait partie d’un ensemble de soins du corps et d’un \teiHi[rend={italic}]{cultus} minimal. L’ensemble de ces prescriptions vise, comme on l’a vu plus haut, à bien différencier le genre de vie stoïcien de celui des cyniques.}
              \teiP[xmlid={p35}]{Au terme de ce parcours, on ne peut pas dire exactement que les éloges d’Euphratès et de Claranus sont des éloges de la laideur, des éloges paradoxaux. En effet, Euphratès, pour âgé, barbu et chevelu qu’il soit, a un beau visage et un aspect qui inspirent le respect. Quant à la laideur de Claranus, due à son âge et à son corps chétif, elle n’en est pas une : ce n’est pas tant qu’elle est compensée par la beauté de son âme, c’est qu’il est beau lui-même, transfiguré par sa vertu et sa qualité de sage. C’est cela qui fait sa parure, son ornement et sa décence tout à la fois. Là où, le plus souvent, la laideur de Socrate subsiste, dans l’éloge qu’Alcibiade fait de lui et dans l’ensemble de la tradition littéraire et iconographique, le regard porté sur les hommes par le philosophe ou l’apprenti philosophe doit savoir dépasser les apparences, physiques et vestimentaires, pour voir la beauté de l’âme, qui seule compte et qui orne le corps, pour reprendre les mots de Sénèque. Nul besoin, pour le sage romain, d’être beau, il suffit d’être convenable et apprécié comme tel. Dans les lettres que nous avons citées, la question du regard et de la vision est fondamentale, autant ou peut-être plus que celle d’une opposition entre extériorité et intériorité. Parvenir à voir le sage sous une apparence de laideur et/ou de pauvreté demande une clairvoyance analogue à celle qu’il faut pour percer à jour les hypocrites qui jouent aux hommes de bien, mais qui ne sont que des hommes mauvais ou débauchés.}
              \teiP[xmlid={p36}]{La question de l’évaluation de la beauté, physique et/ou morale, par le regard du spectateur\teiNote[xmlid={ftn147},n={72},place={foot}]{\teiP[xmlid={p-139}]{Sur cette question, voir notamment Robert Garland, \teiHi[rend={italic}]{The Eye of the Beholder. Deformity and Disability in the Graeco-Roman World} \lbrack{}1995\rbrack{}, Bristol Classical Press, 2010.}} se pose aussi pour les Modernes. Ainsi, on peut penser à une statue antique bien connue, conservée au Musée du Louvre, où la tradition a voulu voir une image de Sénèque mourant, et qui est aujourd’hui considérée comme un exemplaire du type du « vieux pêcheur\teiNote[xmlid={ftn148},n={73},place={foot}]{\teiP[xmlid={p-140}]{Sur cette statue (MR 314), voir \teiRef[target={https://collections.louvre.fr/en/ark:/53355/cl010275490}]{https://collections.louvre.fr/en/ark:/53355/cl010275490}.}} ». Cette statue, découverte à Rome en 1594, a été restaurée quelques années plus tard de façon à recréer une figure fameuse de l’Antiquité, celle de Sénèque, et à donner l’image frappante d’un comportement stoïcien, le suicide\teiNote[xmlid={ftn149},n={74},place={foot}]{\teiP[xmlid={p-141}]{Sur l’histoire de cette statue, voir Mary Beard et John Henderson, \teiHi[rend={italic}]{Classical Art from Greece to Rome}, Oxford, Oxford University Press, 2001, fig. 1 et p. 1-3.}}. Elle a inspiré un tableau de Rubens\teiNote[xmlid={ftn150},n={75},place={foot}]{\teiP[xmlid={p-142}]{Peter Paul Rubens, \teiHi[rend={italic}]{La mort de Sénèque} (entre 1608 et 1613), Munich, Alte Pinacothek (Mary Beard et John Henderson \teiHi[rend={italic}]{op. cit.}, fig. 2). Une autre version du même sujet se trouve au Musée du Prado. Plusieurs éléments de la scène sont tirés du récit du suicide de Sénèque chez Tacite (\teiHi[rend={italic}]{Annales}, XV, 61-65).}}. En 1700, un commentateur (François Raguenet, dans \teiHi[rend={italic}]{Les monuments de Rome} \teiNote[xmlid={ftn151},n={76},place={foot}]{\teiP[xmlid={p-143}]{Paris, v\teiHi[rend={sup}]{ve} Claude Barbin et v\teiHi[rend={sup}]{ve} Daniel Horthemels, 1700, p. 44.}}) jugea que si les artistes savaient faire un Christ aussi expressif, celui-ci ferait pleurer tous les chrétiens. Dans les années 1760, Winckelmann pensait qu’il ne s’agissait pas de Sénèque, mais d’un esclave comique. Enfin, quand Napoléon, en 1807, fit installer cette statue au Louvre, le corps décharné et le visage épuisé du vieillard suscitèrent non l’admiration, mais le dégoût des spectateurs.}
              \teiP[xmlid={p37}]{Indépendamment de l’identification du personnage représenté, Sénèque ou vieux pêcheur, et de la compréhension de la statue, toutes les époques ne jugent évidemment pas la représentation de la nudité et de la vieillesse de la même façon et n’y lisent pas les mêmes valeurs esthétiques, sociales et morales.}
            
\end{teiDiv}
          
\end{teiElement}
        
\end{teiElement}
        \begin{teiElement}[name={group},type={chapter},data-page-title={La beauté contemporaine au miroir de l’esthétique},data-page-authors={Rémi Mermet@@UCLouvain Saint-Louis Bruxelles, FRS-FNRS, Belgique},data-include-href={Ch04\_beaute\_contemporaine\_Mermet.xml},xmlid={chapter-003}]

          \begin{teiElement}[name={front}]

            \begin{teiDiv}[type={titlePage},xmlid={titlepage-005}]

              \teiP[rend={title-main},xmlid={p-144}]{La beauté contemporaine au miroir de l’esthétique}
              \teiP[rend={author-aut},xmlid={p-145}]{Rémi \teiHi[rend={small-caps}]{Mermet}}
              \teiP[rend={authority\_affiliation},xmlid={p-146}]{UCLouvain Saint-Louis Bruxelles, FRS-FNRS, Belgique}
            
\end{teiDiv}
          
\end{teiElement}
          \begin{teiElement}[name={body},xmlid={body-005}]

            \teiP[xmlid={p1-5}]{Caractériser la beauté propre à notre époque est une tâche ardue\teiNote[xmlid={ftn1-4},n={1},place={foot}]{\teiP[xmlid={p-147}]{Par fidélité à « l’esprit » de Cerisy, autant qu’à la définition de l’esthétique comme philosophie circonstancielle et appliquée, nous avons reproduit ici, avec le minimum de corrections et après une courte introduction, le texte prononcé le 16 juillet 2022 sous le titre « Le style, la forme, la vie : sur l’actualité de l’esthétique », dans le cadre de la session « Beautés \teiHi[rend={italic}]{dans} la vie ». Quelques notes complémentaires ont toutefois été ajoutées ici ou là, lorsqu’il nous a paru nécessaire d’affiner ou d’actualiser le propos. Nous remercions l’ensemble des participants au colloque pour leurs remarques généreuses et avisées, qui continueront encore longtemps à nourrir notre pensée.}}. Dans notre vie de tous les jours, l’expérience de la beauté semble relever d’une évidence qui se passe de toute forme de conceptualisation. Combien de fois ne s’écrie-t-on pas : « c’est beau ! » face à une œuvre, un paysage ou un événement, sans qu’il ne soit jamais besoin de fournir davantage d’explications. Inversement, dans les discours savants les plus récents, la beauté tend à apparaître – pour autant qu’elle fasse l’objet d’une véritable réflexion critique – comme une notion surannée, dépassée, voire franchement réactionnaire, hypothéquant par avance toute tentative de renouvellement théorique\teiNote[xmlid={ftn65-2},n={2},place={foot}]{\teiP[xmlid={p-148}]{Voir à ce propos Crispin Sartwell, « Beauty », dans Edward N. Zalta (dir.), \teiHi[rend={italic}]{The Stanford Encyclopedia of Philosophy}, édition de l’été 2022, \teiRef[target={https://plato.stanford.edu/archives/sum2022/entries/beauty/}]{\teiHi[rend={underline}]{https://plato.stanford.edu/archives/sum2022/entries/beauty/}}, en particulier le chapitre intitulé « The Politics of Beauty ». La question de la beauté n’a bien évidemment pas totalement disparu du champ scientifique. Cependant, l’emploi de la notion, principalement dans les sciences naturelles, mais aussi certaines sciences sociales, se fait trop souvent sans réelle problématisation critique. Pour une présentation succincte des dernières recherches en la matière, voir Jacques Morizot, « Beauté (A) », dans Maxime Kristanek (dir.), \teiHi[rend={italic}]{L’Encyclopédie philosophique}, mai 2020, \teiRef[target={https://encyclo-philo.fr/beaute-a}]{\teiHi[rend={underline}]{https://encyclo-philo.fr/beaute-a}}, en particulier la section « Nouvelles explorations méthodologiques ».}}. Pour parer cette double difficulté, nous nous proposons ici d’opérer un détour épistémologique, c’est-à-dire de revenir sur les derniers développements de notre propre discipline, l’esthétique philosophique, qui est la discipline traditionnellement dévolue à l’étude de la beauté. Loin de constituer un simple état des lieux (qui n’est en soi pas inutile au spécialiste comme à l’amateur), ce retour de l’esthétique sur sa propre actualité nous tend en effet un miroir sur le sort présent de la beauté, et son éventuelle vitalité. Car s’il est possible de montrer comment l’esthétique est en mesure de surmonter les critiques dont elle ne cesse de faire l’objet depuis plusieurs années, il est probable que le problème de la beauté, dans la conjonction séculaire qui continue malgré tout à l’unir à l’art et au sensible, trouvera par ricochet une pertinence renouvelée pour la pensée de notre temps. La prétendue « crise » de l’esthétique contemporaine – telle est au fond la question qui nous occupe – n’est-elle pas avant tout celle d’une beauté en attente de sa redéfinition ?}
            \begin{teiDiv}[type={section1},xmlid={div01-2}]

              \teiHead[xmlid={malaise-dans-lesthetique-001}]{Malaise dans l’esthétique}
              \teiP[xmlid={p2-5}]{Répondre, ou plutôt esquisser une réponse à cette question, implique en premier lieu de revenir aux sources mêmes de cette « crise », dont Jacques Rancière dressait en 2004 le constat sans concession. On peut lire, en ouverture de son \teiHi[rend={italic}]{Malaise dans l’esthétique} :}
              \begin{teiCit}[xmlid={cit01-3}]

                \begin{teiQuote}
L’esthétique a mauvaise réputation. Il ne se passe guère une année sans qu’un ouvrage nouveau proclame soit la fin de son temps, soit la perpétuation de ses méfaits. Dans l’un ou l’autre cas, l’accusation est la même. L’esthétique serait le discours captieux par lequel la philosophie ou une certaine philosophie détourne à son profit le sens des œuvres de l’art et des jugements de goût.
\end{teiQuote}
              
\end{teiCit}
              \begin{teiCit}[xmlid={cit02-3}]

                \begin{teiQuote}
Si l’accusation est constante, ses attendus varient. Il y a vingt ou trente ans, le sens du procès pouvait se résumer dans les termes de Bourdieu. Le jugement esthétique « désintéressé », tel que Kant en avait fixé la formule, était le lieu par excellence de la « dénégation du social ». La distance esthétique servait à dissimuler une réalité sociale marquée par la radicale séparation entre les « goûts de nécessité » propres à l’\teiHi[rend={italic}]{habitus} populaire et les jeux de la distinction culturelle réservés à ceux-là seuls qui en avaient les moyens. Une même inspiration animait, dans le monde anglo-saxon, les travaux de l’histoire sociale ou culturelle de l’art. \lbrack{}…\rbrack{}
\end{teiQuote}
              
\end{teiCit}
              \begin{teiCit}[xmlid={cit03-3}]

                \begin{teiQuote}
Cette forme de critique n’est plus guère à la mode. Depuis vingt ans, l’opinion intellectuelle dominante n’en finit pas de dénoncer dans toute forme d’explication « sociale » une complicité ruineuse avec les utopies de l’émancipation déclarées responsables de l’horreur totalitaire. \lbrack{}…\rbrack{} On aurait pu penser que l’esthétique sortirait blanchie de ce cours nouveau de la pensée. Apparemment, il n’en est rien. L’accusation s’est simplement renversée. L’esthétique est devenue le discours pervers qui interdit \lbrack{}le pur\rbrack{} face à face \lbrack{}avec l’événement inconditionné de l’œuvre\rbrack{} en soumettant les œuvres, ou nos appréciations, à une machine de pensée conçue pour d’autres fins\teiNote[xmlid={ftn66-2},n={3},place={foot}]{\teiP[xmlid={p-149}]{Jacques Rancière, \teiHi[rend={italic}]{Malaise dans l’esthétique}, Paris, Galilée, 2004, p. 9 et 10.}}.
\end{teiQuote}
              
\end{teiCit}
              \teiP[xmlid={p3-5}]{En d’autres termes, l’esthétique participerait, pour parler comme Arthur Danto, d’un « assujettissement philosophique de l’art\teiNote[xmlid={ftn67-2},n={4},place={foot}]{\teiP[xmlid={p-150}]{Arthur Danto, \teiHi[rend={italic}]{L’assujettissement philosophique de l’art}, Claude Hary-Schaeffer (trad.), Paris, Éditions du Seuil, 1993.}} », qui serait lui-même le symptôme d’une illusion spéculative nous empêchant d’accéder à l’expérience concrète du monde en général, et des œuvres en particulier.}
              \teiP[xmlid={p4-5}]{Vingt ans après la parution de l’essai de Rancière, la situation de l’esthétique semble plus paradoxale encore. Si personne ne lui conteste sa place au sein de l’édifice académique, elle n’est pas totalement parvenue à surmonter les objections successives qui lui ont été adressées. Pire, elle fait désormais face à leur juxtaposition, au point de se retrouver sous le feu croisé de critiques opposées. Du côté des sciences humaines, tout d’abord, l’esthétique philosophique se voit reprocher sa cécité idéologique. Bien loin d’avoir disparu, la critique de type bourdieusien\teiNote[xmlid={ftn68-2},n={5},place={foot}]{\teiP[xmlid={p-151}]{Pierre Bourdieu, \teiHi[rend={italic}]{La distinction : critique sociale du jugement}, Paris, Éditions de Minuit, 1979.}} s’est au contraire renforcée avec l’essor des études culturelles. Sans vouloir unifier artificiellement un champ dont la fertilité tient justement à l’extrême diversité de ses questionnements, force est en effet de constater la persistance, peut-être moins frontale, peut-être plus diffuse, mais par là même plus prégnante, de la dénonciation du « désintéressement » esthétique comme paravent théorique des structures de domination, dont le concept d’art, cette création récente de la pensée occidentale, serait l’un des bras armés\teiNote[xmlid={ftn69-2},n={6},place={foot}]{\teiP[xmlid={p-152}]{Deborah Knight, « Aesthetics and Cultural Studies », dans Jerrold Levinson (dir.), \teiHi[rend={italic}]{The Oxford Handbook of Aesthetics}, Oxford, Oxford University Press, 2003, p. 783-795.}}. Du côté des sciences naturelles, ensuite, l’esthétique philosophique se voit reprocher son impuissance à trouver un ancrage objectif à ses démonstrations : personne n’a encore jamais pu observer un sujet à ce point « désintéressé » qu’il pourrait se poser comme l’illustration parfaite du sujet transcendantal kantien\teiNote[xmlid={ftn70-2},n={7},place={foot}]{\teiP[xmlid={p-153}]{Voir Bruno Trentini, « Que fait la cognition incarnée à l’esthétique ? », dans Giuseppe Di Liberti et Pierre Léger (dir.), \teiHi[rend={italic}]{La cognition incarnée : un programme de recherche entre psychologie et philosophie}, Sesto San Giovanni, Mimesis, 2022, p. 193-207.}}. D’où le développement de différents programmes de « naturalisation » de l’esthétique, censés lui redonner une assise scientifique – programmes que Jacques Morizot a dernièrement décrits et regroupés selon quatre grandes lignes de recherche\teiNote[xmlid={ftn71-2},n={8},place={foot}]{\teiP[xmlid={p-154}]{Jacques Morizot (dir.), \teiHi[rend={italic}]{Naturaliser l’esthétique ? Questions et enjeux d’un programme philosophique}, Rennes, PUR, 2014, p. 18-27.}} :}
              \begin{teiList}[xmlid={list1},rendition={\#list-decimal},type={ordered}]

    \teiItem{une approche empirique de l’expérience esthétique, fondée sur la psychologie expérimentale ;}
    \teiItem{la recherche, en deçà des variations culturelles du goût et de la création, d’universaux renvoyant à « des fonctions d’optimisation de nos ressources cognitives et neuronales\teiNote[xmlid={ftn72-2},n={9},place={foot}]{\teiP[xmlid={p-155}]{\teiHi[rend={italic}]{Ibid.}, p. 22.}} » ;}
    \teiItem{l’émergence de la neuro-esthétique, fournissant de nombreuses données sur l’ancrage physiologique de notre expérience sensorielle et affective ;}
    \teiItem{l’application au domaine esthétique du modèle évolutionniste darwinien, sur un double plan biologique et anthropologique. De fait, critique naturaliste et critique culturelle du désintéressement ne s’excluent pas toujours mutuellement : elles peuvent parfois s’étayer l’une l’autre, comme dans le projet d’esthétique bioévolutive d’Ellen Dissanayake, dont l’article séminal vient tout juste de regagner en visibilité à l’occasion de sa traduction en français\teiNote[xmlid={ftn73-2},n={10},place={foot}]{\teiP[xmlid={p-156}]{Ellen Dissanayake, « Expérience esthétique et évolution humaine », Pierre Léger (trad.), \teiHi[rend={italic}]{Proteus}, n\teiHi[rend={sup}]{o} 16, 2020, p. 52-65, \teiRef[target={http://www.revue-proteus.com/articles/Proteus16-7.pdf}]{\teiHi[rend={underline}]{http://www.revue-proteus.com/articles/Proteus16-7.pdf}}. Voir également Lorenzo Bartalesi, \teiHi[rend={italic}]{Histoire naturelle de l’esthétique}, Sophie Burdet (trad.) et Jean-Marie Schaeffer (présentation), Paris, CNRS Éditions, 2021.}}.}
   
\end{teiList}
              \teiP[xmlid={p5-5}]{Contestée sur plusieurs fronts, et pas forcément à tort, l’esthétique philosophique fait encore trop souvent le choix de l’indifférence et du repli, comme si le recentrage sur un corpus sanctionné par la tradition suffisait à la protéger d’attaques considérées comme purement extérieures, et donc nulles et non avenues dans son champ de compétences propre. Une position qui, à mon sens, légitime plus qu’elle ne pare les coups qu’elle reçoit. Une position peu apte, surtout, à lui assurer un réel avenir théorique, d’autant que les incompréhensions entre esthétique dite analytique et esthétique dite continentale sont encore légion. Bien sûr, de nombreuses voix se sont élevées, au sein même de la discipline, pour tenter une autocritique de ses biais sous-jacents. Le risque, avec ce genre d’exercice, est néanmoins d’aboutir à une dissolution de cela même qu’on cherchait à sauver. C’est notamment le cas du dernier essai de la regrettée Jacqueline Lichtenstein, qui, du souci tout à fait fondé de sortir la réflexion esthétique de son autarcie philosophique, de la mettre à l’épreuve effective des œuvres et des artistes, en est venu à lui dénier presque toute possibilité d’appréhender son objet\teiNote[xmlid={ftn74-2},n={11},place={foot}]{\teiP[xmlid={p-157}]{Jacqueline Lichtenstein, \teiHi[rend={italic}]{Les raisons de l’art : essai sur les théories de la peinture}, Paris, Gallimard, 2014.}}. En cela, cette approche se révèle exemplaire de la situation contemporaine de l’esthétique philosophique. Ce n’est en effet pas la moindre des contradictions de notre temps que de voir se multiplier les essais d’esthétique consacrés aux limites, réelles ou supposées, de l’esthétique elle-même – contradiction à laquelle, il me faut bien l’avouer, n’échappe pas totalement le présent exposé. Le dynamisme de l’esthétique, sur le plan éditorial tout du moins, semble donc en partie se nourrir de son propre dénigrement. Cet étrange état de fait vient peut-être de ce que nous vivons, ou pensons dorénavant vivre à l’ère du « triomphe de l’esthétique », ou plutôt de « l’hyper-esthétique », comme l’a tout récemment proposé Yves Michaud\teiNote[xmlid={ftn75-2},n={12},place={foot}]{\teiP[xmlid={p-158}]{Yves Michaud, \teiHi[rend={italic}]{L’art à l’état gazeux : essai sur le triomphe de l’esthétique}, Paris, Stock, 2003 et \teiHi[rend={italic}]{« L’art, c’est bien fini » : essai sur l’hyper-esthétique et les atmosphères}, Paris, Gallimard, 2021.}}. Notre quotidien – c’est-à-dire nos corps, nos objets, nos environnements, notre économie même – semble tellement travaillé par une injonction de beauté, d’agrément et de divertissement que le discours esthétique, au même titre que la production artistique, en viendrait irrémédiablement à se vider de toute substance critique, et donc à tourner en boucle sur lui-même.}
              \teiP[xmlid={p6-5}]{On pourrait même aller jusqu’à dire, avec Danièle Cohn et Giuseppe Di Liberti\teiNote[xmlid={ftn76-2},n={13},place={foot}]{\teiP[xmlid={p-159}]{Danièle Cohn et Giuseppe Di Liberti (dir.), \teiHi[rend={italic}]{Esthétique : connaissance, art, expérience}, Paris, Vrin, 2012.}}, que ce sentiment contemporain de crise est en réalité consubstantiel à l’esthétique elle-même, et plus précisément à l’histoire laborieuse de sa naissance comme discipline philosophique durant la seconde moitié du \teiHi[rend={small-caps}]{xviii}\teiHi[rend={sup}]{e} siècle. Quoi de commun, en effet, entre :}
              \begin{teiList}[xmlid={list2},rendition={\#list-decimal},type={ordered}]

    \teiItem{la science de la connaissance sensible et l’art de la beauté du penser que Baumgarten, inventeur du terme « esthétique », appelait de ses vœux\teiNote[xmlid={ftn77-3},n={14},place={foot}]{\teiP[xmlid={p-160}]{Alexandre Gottlieb Baumgarten, \teiHi[rend={italic}]{Esthétique, précédée des Méditations philosophiques sur quelques sujets se rapportant à l’essence du poème et de la Métaphysique}, Jean-Yves Pranchère (trad. et présentation), Paris, L’Herne, 1988.}} ; }
    \teiItem{la critique du jugement esthétique, c’est-dire du beau et du sublime, dont Kant précisait qu’il ne fournit aucune connaissance à proprement parler de ses objets\teiNote[xmlid={ftn78-3},n={15},place={foot}]{\teiP[xmlid={p-161}]{Emmanuel Kant, \teiHi[rend={italic}]{Critique de la faculté de juger}, Alexis Philonenko (trad.), Paris, Vrin, 1993.}} ; }
    \teiItem{l’\teiHi[rend={italic}]{Esthétique} de Hegel, qui, en dépit de son nom, ne relève ni d’une science autonome du sensible, ni d’une analytique du jugement, mais d’une histoire spirituelle de la beauté artistique aboutissant à son nécessaire dépassement philosophique\teiNote[xmlid={ftn79-3},n={16},place={foot}]{\teiP[xmlid={p-162}]{Georg Wilhelm Friedrich Hegel, \teiHi[rend={italic}]{Cours d’esthétique}, Heinrich Gustav Hotho (éd.), Jean-Pierre Lefebvre et Veronika von Schenck (trad.), Paris, Aubier, 1995-1997.}} ? }
   
\end{teiList}
              \teiP[xmlid={p7-5}]{Face à cette valse-hésitation, un penseur comme Konrad Fiedler a pu être tenté de rejeter en bloc le lien établi bon an, mal an entre sensibilité, beauté et création artistique. Il soutient en effet, dès la fin du \teiHi[rend={small-caps}]{xix}\teiHi[rend={sup}]{e} siècle, que « l’erreur première de l’esthétique et de la réflexion sur l’art est d’associer art et beauté comme si le besoin d’art de l’homme visait la constitution d’un monde du beau\teiNote[xmlid={ftn80-2},n={17},place={foot}]{\teiP[xmlid={p-163}]{Konrad Fiedler, \teiHi[rend={italic}]{Aphorismes}, Danièle Cohn (éd.) et Sacha Zilberfarb (trad.), Paris, Éditions Rue d’Ulm, 2013, p. 5, § 2. }} », anticipant ainsi de plus d’un siècle les problématiques les plus brûlantes qui nous animent aujourd’hui.}
            
\end{teiDiv}
            \begin{teiDiv}[type={section1},xmlid={div02-2}]

              \teiHead[xmlid={retour-a-laisthetique-001}]{Retour à « l’aisthétique »}
              \teiP[xmlid={p8-5}]{Mais, s’il est vrai qu’écrire l’histoire d’une discipline est déjà une manière de définition, alors la mise au jour actuelle d’une \teiHi[rend={italic}]{autre} histoire de l’esthétique, qui ne se limite pas à quelques auteurs canoniques, est un indéniable signe de renouveau. Cette histoire alternative – ouverte à des champs tels que l’histoire, la théorie et la critique d’art, la psychophysiologie de la perception et des émotions, la sociologie et l’anthropologie de la création et de la réception, l’épistémologie et la théorie de la connaissance, ou encore l’histoire des sensibilités\teiNote[xmlid={ftn81-2},n={18},place={foot}]{\teiP[xmlid={p-164}]{Sur le champ en pleine croissance de l’histoire des sensibilités, voir notamment \teiHi[rend={italic}]{Sensibilités : entre histoire et anthropologie}, dossier thématique coordonné par Camille Chamois, Quentin Deluermoz et Hervé Mazurel, \teiHi[rend={italic}]{L’Homme}, n\teiHi[rend={sup}]{o} 247-248, 2023, p. 5-40.}} – s’est entre autres cristallisée, depuis une trentaine d’années, autour d’une constellation de pensée ayant fleuri au tournant des \teiHi[rend={small-caps}]{xix}\teiHi[rend={sup}]{e} et \teiHi[rend={small-caps}]{xx}\teiHi[rend={sup}]{e} siècles : nous pensons ici à la \teiHi[rend={italic}]{Kunstwissenschaft} ou « science de l’art » germanophone\teiNote[xmlid={ftn82-2},n={19},place={foot}]{\teiP[xmlid={p-165}]{Dossier \teiHi[rend={italic}]{Problème de la }Kunstwissenschaft, Holger Schmid (dir.) avec la collaboration de Eliane Escoubas, \teiHi[rend={italic}]{La part de l’œil}, n\teiHi[rend={sup}]{o}\teiHi[rend={italic}]{ }15/16, 1999-2000.}} et, dans une moindre mesure, à sa cousine francophone\teiNote[xmlid={ftn83-2},n={20},place={foot}]{\teiP[xmlid={p-166}]{Jacqueline Lichtenstein, Carole Maigné et Arnauld Pierre (dir.), \teiHi[rend={italic}]{Vers la science de l’art : l’esthétique scientifique en France, 1857-1937}, Paris, PUPS, 2013.}}. Là n’est pas le lieu de s’attarder sur ces théories, qui ont été l’occasion d’une collaboration entre histoire, philosophie, anthropologie et psychologie dont la fécondité est jusqu’à présent restée inégalée. Ce qu’il nous faut ici retenir, c’est l’ambition de cette science de l’art d’articuler étroitement fondation critique de son propre discours, abandon des carcans convenus du bon goût, prise en compte de l’historicité des œuvres et de nos interprétations, et réflexion générique sur le sensible, d’un double point de vue perceptif et affectif. La \teiHi[rend={italic}]{Kunstwissenschaft} part en effet du principe que notre sensibilité n’est jamais passive, mais bel et bien active, qu’elle est un processus de mise en forme de notre expérience non absolument séparé de notre intellect, que cette activité est toujours incarnée dans une culture et une histoire particulière, et que cette incarnation nous est justement donnée de façon éclatante dans les œuvres d’art. Ce faisant, le lien entre l’esthétique et l’artistique se voit réaffirmé, non plus sur le mode de l’assujettissement, comme le pensait Danto, ni sur celui de la vaporisation, comme le soutient Michaud, mais sur celui de l’exemplarité. L’artistique n’y est pas conçu comme l’objet exclusif ou définitivement perdu de l’esthétique, mais comme son objet de réflexion privilégié, à partir duquel il est possible de penser le caractère formateur, et donc cognitif, de \teiHi[rend={italic}]{toute} sensibilité. Telle est, selon Danièle Cohn, « la voie qui reconduit l’esthétique à elle-même\teiNote[xmlid={ftn84-2},n={21},place={foot}]{\teiP[xmlid={p-167}]{Danièle Cohn et Giuseppe Di Liberti (dir.), \teiHi[rend={italic}]{op. cit.}, p. 26.}} ».}
              \teiP[xmlid={p9-4}]{Et pour cause : la piste probablement la plus prometteuse de l’esthétique philosophique contemporaine, l’une de ses tendances les plus vigoureuses, consiste en un retour assumé à la question de l’\teiHi[rend={italic}]{aisthesis}, c’est-à-dire à l’esthétique comme science des percepts et des affects. Cette esthétique, qui, comme le remarque Mildred Galland-Szymkowiak, « vaut comme savoir sur le beau et l’art pour autant qu’elle est, d’abord, science de la sensation\teiNote[xmlid={ftn85-2},n={22},place={foot}]{\teiP[xmlid={p-168}]{Mildred Galland-Szymkowiak, « L’esthétique comme science de la sensation. Jalons pour une histoire », dans Gernot Böhme, \teiHi[rend={italic}]{Aisthétique : pour une esthétique de l’expérience sensible}, Martin Kaltenecker et Franck Lemonde (trad.), Emmanuel Alloa et Céline Flécheux (préface), Mildred Galland-Szymkowiak (postface), Dijon, Les Presses du réel, 2020, p. 237. Voir également Giovanni Matteucci, \teiHi[rend={italic}]{Estetica e natura umana. La mente estesa tra percezione, emozione ed espressione}, Rome, Carocci, 2019.}} », se veut une « esthétique au-delà de l’esthétique », dans le sens où, soutient cette fois-ci Wolfgang Welsch, les problématiques traditionnelles de l’esthétique, restreintes à des questions telles que l’art, le beau ou le goût, gagnent nécessairement en complexité et en profondeur une fois replacées au sein d’une définition élargie – on pourrait ajouter : accueillante, impure et décomplexée – de la discipline\teiNote[xmlid={ftn86-2},n={23},place={foot}]{\teiP[xmlid={p-169}]{Wolfgang Welsch, « Aesthetics Beyond Aesthetics », dans Francis Halsall, Julia Jansen et Tony O’Connor (dir.), \teiHi[rend={italic}]{Rediscovering Aesthetics: Transdisciplinary Voices from Art History, Philosophy, and Art Practice}, Stanford, Stanford University Press, 2009, p. 178-192. Voir également Thomas Mercier-Bellevue et Thomas Morisset (dir.), \teiHi[rend={italic}]{D’une beauté trouble et mêlée : sur l’impureté de l’expérience esthétique}, Paris, Vrin, 2024.}}. C’est ainsi que le philosophe allemand Gernot Böhme, décédé en 2022, a développé sa théorie des atmosphères, c’est-à-dire de la spatialité affective au sein de laquelle baigne notre expérience des choses, en partant d’une étude phénoménologique générale de la perception, pour se déployer ensuite en direction d’une esthétique écologique de la nature, d’une esthétique critique du design en régime capitaliste, et d’une esthétique des arts contemporains et de leur rejet des formes traditionnelles de la représentation\teiNote[xmlid={ftn87-2},n={24},place={foot}]{\teiP[xmlid={p-170}]{Gernot Böhme, \teiHi[rend={italic}]{op. cit}. Nous remercions Pierre Sauvanet de nous avoir ici suggéré la lecture de l’essai de Bruce Bégout, \teiHi[rend={italic}]{Le concept d’ambiance} (Paris, Éditions du Seuil, 2020), qui défend, à l’encontre de Böhme, une appréhension foncièrement pathique (passive ?) de l’expérience du monde. Cette appréhension manque justement selon nous de la richesse de la pensée esthétique, un peu trop rapidement évacuée au nom d’un primat du quotidien sur l’art, comme s’il s’agissait là de domaines mutuellement (et irrémédiablement) exclusifs.}}. Dans un tout autre genre, Bence Nanay, qui s’inscrit dans la lignée de la philosophie analytique, a entrepris la fondation d’une esthétique en tant que « philosophie de la perception ». Celle-ci traite de questions aussi diverses que les modalités de l’attention perceptive ou de la représentation iconique, la possibilité d’une historicité de la vision, ou encore le problème de l’unicité de l’expérience sensible\teiNote[xmlid={ftn88-2},n={25},place={foot}]{\teiP[xmlid={p-171}]{Bence Nanay, \teiHi[rend={italic}]{L’esthétique, une philosophie de la perception}, Jacques Morizot (trad.), Rennes, PUR, 2021.}}. Ce projet de recherche puise d’ailleurs assez largement dans la boîte à outils des sciences cognitives, dont la variante dite de la « cognition incarnée », centrée sur une réévaluation du corps dans les processus cognitifs, a récemment trouvé un écho favorable – mais est-ce bien surprenant ? – chez plusieurs tenants de l’esthétique de l’expérience sensible (en particulier sous sa forme pragmatiste, comme chez Richard Shusterman\teiNote[xmlid={ftn89-2},n={26},place={foot}]{\teiP[xmlid={p-172}]{Voir notamment Giuseppe Di Liberti et Pierre Léger (dir.), \teiHi[rend={italic}]{op. cit}.}}).}
              \teiP[xmlid={p10-4}]{Ce retour à l’\teiHi[rend={italic}]{aisthesis}, nous insistons, ne signe en aucun cas le rejet ou l’oubli du problème de la signification des œuvres, des expériences et des jugements. N’en déplaise à Jean-Marie Schaeffer, la conception « aisthétique » de l’esthétique, qui fait effectivement de l’art le lieu par excellence du sensible, ne « sous-estime » en rien « la complexité herméneutique des œuvres\teiNote[xmlid={ftn90-2},n={27},place={foot}]{\teiP[xmlid={p-173}]{Jean-Marie Schaeffer, « Postface à l’édition de 2016 », dans \teiHi[rend={italic}]{Adieu à l’esthétique}, Sesto San Giovanni, Mimesis, 2016, p. 92.}} ». Au contraire, et pour en rester aux exemples précédents, Böhme consacre de nombreuses analyses à la question des signes et des symboles, c’est-à-dire au problème du rapport entre esthétique, sémiotique et herméneutique. Quant à Nanay, ses réflexions sur l’historicité de la vision le conduisent à s’interroger sur l’influence possible des valeurs culturelles sur les processus visuels. Au fond, on pourrait dire avec Jacques Morizot que « \lbrack{}l\rbrack{}’esthétique en est arrivée aujourd’hui à un point où cela refait sens de croiser le perceptuel et le symbolique, c’est-à-dire de coupler la psychologie en acte et le sens en contexte\teiNote[xmlid={ftn91-2},n={28},place={foot}]{\teiP[xmlid={p-174}]{Jacques Morizot (dir.), \teiHi[rend={italic}]{op. cit.}, p. 30. Il eût été probablement plus judicieux de parler ici de « sensible » plutôt que de « perceptuel », car ce dernier terme tend implicitement à éliminer toute dimension proprement affective du champ de l’\teiHi[rend={italic}]{aisthesis}.}} ». Sans nous étendre à nouveau sur tout ce que les penseurs de la \teiHi[rend={italic}]{Kunstwissenschaft} pourraient encore nous apprendre en la matière (ce qui rend d’autant plus urgente la réécriture en cours de l’histoire de l’esthétique\teiNote[xmlid={ftn92-2},n={29},place={foot}]{\teiP[xmlid={p-175}]{Voir à ce propos Andreas Beyer, Danièle Cohn et Tania Vladova, « Introduction » au numéro thématique \teiHi[rend={italic}]{Esthétique et science de l’art}, \teiHi[rend={italic}]{Trivium}, n\teiHi[rend={sup}]{o} 6, 2010, DOI : \teiRef[target={https://doi.org/10.4000/trivium.3664}]{10.4000/trivium.3664}.}}), il est clair qu’un tel programme, par son ambition, son ouverture d’esprit également, nécessitera de l’ensemble des acteurs du champ esthétique le dépassement d’un certain nombre de postures théoriques – postures qui, malgré l’indéniable intensification des échanges interdisciplinaires, continuent de favoriser la fragmentation artificielle des savoirs.}
              \teiP[xmlid={p11-4}]{De la part des philosophes, il s’agira d’éviter, dans un réflexe de défense humaniste qui, nous devons bien l’admettre, a longtemps été mien, de faire preuve d’une défiance chronique envers toute velléité d’explication causale ou expérimentale de l’esprit humain. Non qu’il faille abandonner par-là l’exigence critique propre à la philosophie. Au contraire, c’est en faisant preuve d’une saine curiosité envers les tentatives de naturalisation de l’esthétique que la philosophie pourra pleinement jouer son rôle de garde-fou épistémologique\teiNote[xmlid={ftn93-2},n={30},place={foot}]{\teiP[xmlid={p-176}]{Voir par exemple Denis Forest, \teiHi[rend={italic}]{Neuroscepticisme : les sciences du cerveau sous le scalpel de l’épistémologue}, Montreuil-sous-Bois, Ithaque, 2014 et \teiHi[rend={italic}]{Neuropromesses : une enquête philosophique sur les frontières des neurosciences}, Paris, Ithaque, 2022.}}, quitte à prononcer, en dernier ressort, l’impossibilité même de cette entreprise de naturalisation.}
              \teiP[xmlid={p12-4}]{De façon symétrique, les scientifiques devront se préserver de l’ornière d’une « esthétique scientiste », d’un réductionnisme naturaliste aussi tentant que contre-productif, dont Fernando Vidal a bien montré les impasses\teiNote[xmlid={ftn94-2},n={31},place={foot}]{\teiP[xmlid={p-177}]{Fernando Vidal, « La neuroesthétique, un esthétisme scientiste », \teiHi[rend={italic}]{Revue d’histoire des Sciences humaines}, vol. 25, n\teiHi[rend={sup}]{o} 2, 2011, p. 239-264, DOI : 10.3917/rhsh.025.0239.}}. Mettre l’esthétique à l’épreuve des faits scientifiques présuppose une élaboration conceptuelle préalable (sur le beau, le plaisir, l’art, etc.), qui aura tout à gagner à passer au tamis d’une clarification philosophique rigoureuse, mais aussi à être remise dans une perspective historique plus large – c’est-à-dire une perspective qui ne se limite pas à la seule invocation rituelle de Gustav Fechner comme père de l’esthétique expérimentale au \teiHi[rend={small-caps}]{xix}\teiHi[rend={sup}]{e} siècle\teiNote[xmlid={ftn95-2},n={32},place={foot}]{\teiP[xmlid={p-178}]{Gustav Theodor Fechner, \teiHi[rend={italic}]{Vorschule der Aesthetik}, Leipzig, Breitkopf \& Härtel, 1876.}}, mais inclut de véritables connaissances en matière d’histoire des arts, des sciences, des pratiques et des concepts\teiNote[xmlid={ftn96-2},n={33},place={foot}]{\teiP[xmlid={p-179}]{Voir à ce propos Whitney Davis, « Neurovisuality », \teiHi[rend={italic}]{nonsite.org}, n\teiHi[rend={sup}]{o} 2, 2011, \teiRef[target={https://nonsite.org/neurovisuality/}]{https://nonsite.org/neurovisuality/}.}}.}
              \teiP[xmlid={p13-3}]{C’est là, bien évidemment, le rôle pour ainsi dire constitutif des sciences humaines, sans lesquelles la philosophie tournerait à vide, et les sciences naturelles procéderaient hors sol, ou plutôt hors contexte. Non que les savoirs établis par ces dernières, notamment en matière de création artistique et de jugement esthétique, soient à relativiser absolument. Comme le soutient le philosophe Paul Crowther, une position entièrement relativiste n’est, en esthétique pas plus qu’ailleurs, logiquement tenable\teiNote[xmlid={ftn97-2},n={34},place={foot}]{\teiP[xmlid={p-180}]{Paul Crowther, \teiHi[rend={italic}]{Defining Art, Creating the Canon: Artistic Value in an Era of Doubt}, Oxford, Oxford University Press, 2007.}}. Mais entre l’idée, défendue par Jean-Marie Schaeffer, selon laquelle l’expérience esthétique est une donnée cognitive naturelle, qui transcende toutes les époques et toutes les cultures\teiNote[xmlid={ftn98-2},n={35},place={foot}]{\teiP[xmlid={p-181}]{Jean-Marie Schaeffer, \teiHi[rend={italic}]{L’expérience esthétique}, Paris, Gallimard, 2015.}}, et celle, défendue par Philippe Descola, selon laquelle cette même expérience perd toute efficacité cognitive en dehors d’un moment très restreint de l’histoire occidentale\teiNote[xmlid={ftn99-2},n={36},place={foot}]{\teiP[xmlid={p-182}]{Philippe Descola, \teiHi[rend={italic}]{Les formes du visible : une anthropologie de la figuration}, Paris, Éditions du Seuil, 2021.}}, il est peut-être possible d’ouvrir une troisième voie, en faisant fructifier, en toute conscience, l’apport spécifique de la tradition esthétique dans l’interprétation d’objets et d’expériences qui n’appartiennent pas forcément à notre canon culturel. La véritable « neutralité axiologique », ainsi que l’a montré Isabelle Kalinowski dans ses commentaires sur Max Weber, ne relève jamais d’une position de surplomb impossible à atteindre dans les faits. Elle relève au contraire d’un usage éclairé du caractère nécessairement situé, localisé, engagé de tout point de vue sur le monde, y compris le point de vue philosophique ou scientifique\teiNote[xmlid={ftn100-2},n={37},place={foot}]{\teiP[xmlid={p-183}]{Isabelle Kalinowski, « Leçon wébériennes sur la science et la propagande », dans Max Weber, \teiHi[rend={italic}]{La science, profession et vocation}, Isabelle Kalinowski (trad.), Marseille, Agone, 2005, p. 191 et suiv. L’autrice y démontre combien la \teiHi[rend={italic}]{Wertfreiheit} wébérienne n’est en réalité jamais affaire de « neutralité » – d’où l’inadéquation de la traduction consacrée par l’usage.}}. C’est là, probablement, la plus grande leçon de « l’aisthétique », dont les études actuellement menées par Guilhem Marion sur les rapports entre culture et perception montrent, me semble-t-il, tout le potentiel pour les sciences expérimentales\teiNote[xmlid={ftn101-2},n={38},place={foot}]{\teiP[xmlid={p-184}]{Voir le texte de Guilhem Marion dans le présent volume.}}.}
            
\end{teiDiv}
            \begin{teiDiv}[type={section1},xmlid={div03-2}]

              \teiHead[xmlid={beautes-inframinces-001}]{Beautés inframinces}
              \teiP[xmlid={p14-3}]{Ainsi parvenus à cette série de propositions méthodologiques, que nous apprend, plus concrètement, ce bref portrait de « l’aisthétique » contemporaine sur la beauté dans notre vie, ici et maintenant ? Si l’on accepte un tant soit peu le constat d’une « hyper-esthétique », ne devrait-on pas tout bonnement abandonner l’espoir de pouvoir encore faire une véritable expérience de la beauté ? Car, quand tout tend à devenir esthétique, quand tout devient l’occasion de vivre une expérience sensationnelle, le risque n’est-il pas justement celui de l’anesthésie globale ? Bien sûr, la critique n’est pas neuve, et a émané de l’ensemble du spectre idéologique, ainsi que le rappelle à juste titre l’un des derniers numéros de la \teiHi[rend={italic}]{Nouvelle Revue d’esthétique}, précisément consacré au problème de « l’esthétisation\teiNote[xmlid={ftn102-2},n={39},place={foot}]{\teiP[xmlid={p-185}]{\teiHi[rend={italic}]{Esthétisation}, dossier thématique coordonné par Emmanuel Alloa et Christoph Haffter, \teiHi[rend={italic}]{Nouvelle revue d’esthétique}, vol. 28, n\teiHi[rend={sup}]{o} 2, 2021.}} ». Cette critique, peut-on lire, remonte au moins à l’entre-deux-guerres, et la prise de conscience liée aux emplois possiblement délétères du programme moderne d’embellissement généralisé de la vie. On pourrait même remonter plus loin encore dans le temps, puisque l’historien de l’art suisse Heinrich Wölfflin écrivait dès 1909, à propos de l’appétence croissante du public pour les beautés artistiques du passé :}
              \begin{teiCit}[xmlid={cit04-3}]

                \begin{teiQuote}
On se sent obligé de tout voir, et c’est fort dommage, car ce faisant on voit trop, ce qui revient à ne rien voir du tout.
\end{teiQuote}
              
\end{teiCit}
              \begin{teiCit}[xmlid={cit05-2}]

                \begin{teiQuote}
Le zèle du touriste « cultivé » moderne est une monstruosité psychologique. En Italie, les églises et les musées sont passés au peigne fin de l’histoire de l’art, avec une exhaustivité dont chacun, cependant, a tôt fait de s’apercevoir qu’elle ne peut laisser aucune impression profonde : un clou chasse l’autre. Voyager devient un supplice, mais on aurait mauvaise conscience de se rendre la vie plus commode. En avant, marche ! Les guides de voyage prescrivent la quantité de ce qui doit être « fait » chaque jour, et pas question de rester à la traîne !
\end{teiQuote}
              
\end{teiCit}
              \begin{teiCit}[xmlid={cit06-2}]

                \begin{teiQuote}
\lbrack{}…\rbrack{} Il est vrai que nos grands musées incitent eux-mêmes à cette boulimie du voir, et lorsqu’on observe la masse des visiteurs qui se bousculent à travers les salles au risque de l’épuisement psychologique, il nous vient des pensées bien pessimistes quant à l’utilité de telles institutions\teiNote[xmlid={ftn103-2},n={40},place={foot}]{\teiP[xmlid={p-186}]{Heinrich Wölfflin, « Du mauvais usage de l’histoire de l’art », dans \teiHi[rend={italic}]{L’explication des œuvres d’art}, Sacha Zilberfarb (trad.) et Danièle Cohn (postface), Paris, L’Écarquillé, 2018, p. 10 et 11.}}.
\end{teiQuote}
              
\end{teiCit}
              \teiP[xmlid={p15-3}]{Si l’accusation est donc tout sauf récente, et, sous la plume de Wölfflin, pas nécessairement dénuée d’élitisme, il faut avouer que la donne a quelque peu changé en un siècle. Pour nous en tenir à la seule question du musée, quiconque a déjà tenté d’approcher \teiHi[rend={italic}]{La Joconde }un jour de grande fréquentation saura à quoi nous faisons ici référence. La « boulimie » du voir qui est aujourd’hui la nôtre concerne désormais moins une volonté « cultivée » d’exhaustivité historique, qu’un souci de voir ce que tant d’autres ont vu avant nous, tout en faisant voir à ces autres (dans une mise en abyme dont le selfie est l’expression la plus emblématique\teiNote[xmlid={ftn104-2},n={41},place={foot}]{\teiP[xmlid={p-187}]{Voir à ce propos Danièle Cohn, « On voit trop », dans Heinrich Wölfflin, \teiHi[rend={italic}]{op. cit.}, p. 50-56.}}) que l’on a enfin vu ce que l’on « devait » voir. Dans de telles conditions – et c’est là un enjeu majeur pour les services de médiation culturelle\teiNote[xmlid={ftn105-2},n={42},place={foot}]{\teiP[xmlid={p-188}]{Sans parler des problèmes politiques et écologiques soulevés, plus généralement, par le tourisme de masse et sa mise en scène sur les réseaux sociaux.}} – il est pour le moins difficile de faire une réelle expérience des œuvres, et encore moins de leur éventuelle beauté, qui constituait pourtant la source, ou du moins le prétexte, de cette prescription du voir\teiNote[xmlid={ftn106-2},n={43},place={foot}]{\teiP[xmlid={p-189}]{Sur la déception éprouvée par beaucoup de visiteurs du Louvre face à \teiHi[rend={italic}]{La Joconde}, voir « Les catégories au musée. Entretien avec Sébastien Allard », dans Danièle Cohn et Rémi Mermet (dir.), \teiHi[rend={italic}]{L’Histoire de l’art et ses concepts. Autour de Heinrich Wölfflin}, Paris, Éditions Rue d’Ulm, 2020, p. 145 et 146.}}. Car pour qu’il y ait expérience, et partage de cette expérience, pour qu’une expérience vécue devienne réellement formatrice pour les sujets, il faut, comme le soutenait Goethe en son temps, que ces mêmes sujets s’affrontent à l’objectalité des choses – à des choses qui, d’une manière ou d’une autre, leur offre une résistance, leur offre matière à sentir et à penser, plutôt que de servir d’excuse à la satisfaction immédiate d’un désir autocentré\teiNote[xmlid={ftn107-2},n={44},place={foot}]{\teiP[xmlid={p-190}]{Johann Wolfgang von Goethe, « L’expérimentation comme médiation entre le sujet et l’objet », dans Danièle Cohn et Giuseppe Di Liberti (dir.), \teiHi[rend={italic}]{op. cit.}, p. 101-113. Cette formation du sujet par l’expérience qu’il fait des choses est l’un des enjeux majeurs de la tradition pédagogique de la \teiHi[rend={italic}]{Bildung}, dont Goethe est un éminent représentant. L’idée, ici, n’est bien évidemment pas d’opposer l’usage « populaire » des réseaux sociaux à un usage « légitime » du musée, qui serait l’apanage des seules élites culturelles. Il s’agit au contraire de promouvoir le musée comme espace démocratique par excellence, c’est-à-dire comme lieu d’une expérience émancipatrice, tout à la fois individuelle et collective, au contact de ces formes de réflexivité sensible et culturelle que sont les œuvres d’arts. }}. « Un fait de notre vie ne vaut pas dans la mesure où il est vrai », disait encore Goethe à Eckermann, c’est-à-dire dans la mesure où nous pouvons simplement attester l’avoir vécu, mais bien « dans la mesure où il a quelque chose à signifier\teiNote[xmlid={ftn108-2},n={45},place={foot}]{\teiP[xmlid={p-191}]{\teiHi[rend={italic}]{Conversation de Goethe avec Eckermann}, mercredi 30 mars 1831 (la traduction française de Jean Chuzeville, Paris, Gallimard, 1988, p. 413, est légèrement différente).}} ».}
              \teiP[xmlid={p16-3}]{Sur le plan théorique, ce dépassement de la particularité du vécu en direction d’une généralisation de sa signification n’est pas sans évoquer un autre concept goethéen, étroitement lié à la notion de beauté : le concept de style. Sans entrer dans le détail, le style occupe pour Goethe l’échelon le plus abouti de la création artistique, en ce qu’en lui l’œuvre atteint la dignité de l’universel dans l’éclat même de sa manifestation singulière\teiNote[xmlid={ftn109-2},n={46},place={foot}]{\teiP[xmlid={p-192}]{Johann Wolfgang von Goethe, « Simple imitation de la nature, manière, style », dans \teiHi[rend={italic}]{Écrits sur l’art}, Jean-Marie Schaeffer (trad.), Tzvetan Todorov (intro.), Paris, Flammarion, 1996, p. 95-101.}} – une sorte de beauté kantienne objectivée, pourrait-on dire. Le style goethéen est autant affaire de beauté que de vérité, ou plutôt : il est affaire de beauté pour autant qu’il est affaire de vérité, d’une vérité sensible, non pas abstraite, mais en acte, qui ne s’offre à nous que « sous forme de figures visibles et tangibles\teiNote[xmlid={ftn110-2},n={47},place={foot}]{\teiP[xmlid={p-193}]{\teiHi[rend={italic}]{Ibid.}, p. 98.}} ». Le style, chez Goethe, n’est donc pas la forme opposée au contenu, le sensible opposé au sens : il est le contenu sensible en tant qu’il fait sens dans sa mise en forme même – mise en forme dont la beauté nous révèle la justesse intrinsèque\teiNote[xmlid={ftn111-2},n={48},place={foot}]{\teiP[xmlid={p-194}]{Sur la justesse, voir Danièle Cohn, \teiHi[rend={italic}]{L’artiste, le vrai et le juste : sur l’esthétique des Lumières}, Paris, Éditions Rue d’Ulm, 2014.}}. Contre toute la rhétorique superficielle et commerciale du style « comme style de vie », pour citer le slogan d’un célèbre grand magasin parisien, c’est dans ce rapport intime au sens et au vrai que réside selon nous la puissance critique du concept de style, et donc la vitalité plus que jamais actuelle de la beauté\teiNote[xmlid={ftn112-2},n={49},place={foot}]{\teiP[xmlid={p-195}]{Nous renvoyons ici à Rémi Mermet, « S comme style », dans Clélia Zernik et Justin Jaricot (dir.), \teiHi[rend={italic}]{Abécédaire de la beauté}, Paris, Éditions B42, 2022, p. 190-193.}}. C’est, me semble-t-il, ce que prouve le retour en force de la question stylistique, après quelques décennies de disgrâce, dans les sciences contemporaines (histoire littéraire, de l’art, des sciences, mais aussi sociologie ou anthropologie\teiNote[xmlid={ftn113-2},n={50},place={foot}]{\teiP[xmlid={p-196}]{Voir par exemple le projet EPISTYLE, dirigé par Matteo Vagelli, qui a pour ambition de cartographier les usages du concept de style en histoire des sciences, \teiRef[target={https://pric.unive.it/projects/epistyle/home}]{https://pric.unive.it/projects/epistyle/home} (consulté le 19 juin 2025).}}), retour qui s’accompagne, dans le champ de la philosophie, d’un déploiement tout aussi remarquable de la réflexion autour de la problématique des « formes de vie » et de la stylistique de l’existence\teiNote[xmlid={ftn114-2},n={51},place={foot}]{\teiP[xmlid={p-197}]{Voir Estelle Ferrarese et Sandra Laugier (dir.), \teiHi[rend={italic}]{Formes de vie}, Paris, CNRS Éditions, 2018 et Martin Mees, « La vie comme œuvre d’art ? », \teiHi[rend={italic}]{Nouvelle Revue d’esthétique}, n\teiHi[rend={sup}]{o} 28, 2021, p. 61-68.}}.}
              \teiP[xmlid={p17-3}]{Mais plus encore que sa fécondité épistémologique, ce qui distingue le paradigme contemporain du style, c’est qu’il n’est plus le « grand » style à portée normative des classiques, ni le style unitaire (national, d’époque) que les historiens comme Wölfflin ont tenté de retrouver dans les arts du passé, ni même encore la tentation dandyste d’un repli esthétisant sur soi-même. Le style – telle est l’hypothèse que nous sommes de plus en plus nombreux à formuler – marque désormais plus sobrement ce qu’on pourrait nommer l’opérativité générique de la vie, en tant que celle-ci n’existe que dans la variabilité de ses formes, c’est-à-dire de formes qui revendiquent pour elles-mêmes une valeur, ou plutôt une justesse propre, dans un refus conjoint de l’hégémonie comme de l’indifférence. Un pluralisme sans relativisme, en somme, qui rend par là même possible la compréhension comme la controverse, l’échange comme le rapport de force, expressions vivantes de la portée non seulement collective, mais transformatrice – inhéremment politique – du champ esthétique\teiNote[xmlid={ftn115-2},n={52},place={foot}]{\teiP[xmlid={p-198}]{On retrouve ce faisant l’un des autres grands motifs de l’esthétique ranciérienne, à savoir la question du « partage du sensible ». Voir à ce propos Jacques Rancière, \teiHi[rend={italic}]{Le partage du sensible : esthétique et politique}, Paris, La Fabrique, 2000. Nous renvoyons également à Rémi Mermet, « Le style, encore et toujours », \teiHi[rend={italic}]{Revue de synthèse}, vol. 145, n\teiHi[rend={sup}]{o} 3-4, 2024, p. 311-323, DOI : 10.1163/19552343-14234060.}}. Aussi, plutôt que de dénoncer, avec Carole Talon-Hugon, l’« art documentaire » comme le présent symptôme d’une confusion fatale, pour l’un comme pour l’autre, de l’art et de la science\teiNote[xmlid={ftn116-2},n={53},place={foot}]{\teiP[xmlid={p-199}]{Carole Talon-Hugon, \teiHi[rend={italic}]{L’artiste en habit de chercheur}, Paris, PUF, 2021. Sur l’art documentaire, voir notamment Aline Caillet et Frédéric Pouillaude (dir.), \teiHi[rend={italic}]{Un art documentaire : enjeux esthétiques, politiques et éthiques}, Rennes, PUR, 2017.}}, nous préférons retenir la proposition de Marielle Macé, selon laquelle le documentaire serait justement l’art qui nous est le plus contemporain, car il « ne prend pas pour objets des vies, mais des formes de vie \lbrack{}…\rbrack{}, les \teiHi[rend={italic}]{formes} qui traversent ces vies réelles ou possibles, faisant émerger en elles autant d’idées du vivre. \lbrack{}…\rbrack{} C’est un parti pris sur les formes, en même temps qu’un parti pris sur la vie : car le documentaire est ce qui engage sa propre forme dans une attention au formel de la vie et à ce qui, de la vie et pour elle, s’en dégage\teiNote[xmlid={ftn117-2},n={54},place={foot}]{\teiP[xmlid={p-200}]{Marielle Macé, \teiHi[rend={italic}]{Styles : critique de nos formes de vie}, Paris, Gallimard, 2016, p. 310 et 312.}} ». Il y a là d’après nous reprise, réactualisation, sur un mode beaucoup moins monumental, beaucoup moins prométhéen, de l’articulation chère à Goethe entre poésie et vérité\teiNote[xmlid={ftn118-2},n={55},place={foot}]{\teiP[xmlid={p-201}]{Johann Wolfgang von Goethe, \teiHi[rend={italic}]{Poésie et Vérité : souvenirs de ma vie}, Pierre du Colombier (trad.), Paris, Aubier, 1941. Voir à ce propos Danièle Cohn, \teiHi[rend={italic}]{La Lyre d’Orphée : Goethe et l’esthétique}, Paris, Flammarion, 1999.}}, dont témoigne à sa manière Annie Ernaux lorsqu’elle place en exergue de son dernier ouvrage les mots suivants : « Si je ne les écris pas, les choses ne sont pas allées jusqu’à leur terme, elles ont été seulement vécues\teiNote[xmlid={ftn119-2},n={56},place={foot}]{\teiP[xmlid={p-202}]{Annie Ernaux, \teiHi[rend={italic}]{Le Jeune Homme}, Paris, Gallimard, 2022, p. 9.}}. »}
              \teiP[xmlid={p18-3}]{Or, puisque la beauté serait, toujours d’après Talon-Hugon, le vecteur ayant progressivement conduit l’art à son délitement « documentaire », en le soumettant aux caprices d’une subjectivité ne respectant plus aucune règle\teiNote[xmlid={ftn120-2},n={57},place={foot}]{\teiP[xmlid={p-203}]{Carole Talon-Hugon, \teiHi[rend={italic}]{L’art victime de l’esthétique}, Paris, Hermann, 2014. }}, une lecture inverse pourra peut-être déceler dans cet art « documentaire » une nouvelle forme de beauté – une beauté subtile, discrète, surgissant au ras même de nos vies, ou plutôt dans les interstices et les craquelures de nos anesthésies. Comment qualifier cette beauté ? Nous opterons volontiers pour l’adjectif « inframince », ce néologisme créé par Marcel Duchamp dans les années 1930 pour désigner « \lbrack{}l\rbrack{}a chaleur d’un siège (qui vient d’être quitté) », le « sifflotement » produit « dans la marche » par le frottement des deux jambes d’un pantalon de velours, ou encore les « \lbrack{}r\rbrack{}eflets de la lumière sur diff\lbrack{}érentes\rbrack{} surfaces plus ou moins polies\teiNote[xmlid={ftn121-2},n={58},place={foot}]{\teiP[xmlid={p-204}]{Marcel Duchamp, « Inframince », dans \teiHi[rend={italic}]{Duchamp du signe suivi de Notes}, Paris, Flammarion, 2008, p. 277-292.}} ». À rebours du récit mythique de l’art contemporain, qui a érigé Duchamp – le contempteur autoproclamé de l’art « rétinien » – en père de l’art conceptuel, Thierry Davila a bien montré l’importance de la question de « l’inframince » pour comprendre l’ensemble de son œuvre. Loin d’exclure toute dimension esthétique, l’art de Duchamp attend au contraire du spectateur qu’il « bris\lbrack{}e\rbrack{} l’anesthésie pour découvrir \lbrack{}…\rbrack{} ce qui se présente à lui et qu’il a sous les yeux sans pour autant en relever d’emblée la configuration ténue. La tâche de l’art et des artistes \lbrack{}…\rbrack{} consiste alors \lbrack{}…\rbrack{} à se donner les moyens d’étendre le champ de la perception, la question du beau devenant un sujet dérivé (et mineur) au regard des enjeux phénoménologiques – et politiques – de cette quête des variations infimes\teiNote[xmlid={ftn122-2},n={59},place={foot}]{\teiP[xmlid={p-205}]{Thierry Davila, \teiHi[rend={italic}]{De l’inframince : brève histoire de l’imperceptible, de Marcel Duchamp à nos jours}, Paris, Éditions du Regard, 2010, p. 274. Ce qui est ici suggéré, nous semble-t-il, c’est que l’« anesthésie » invoquée par Duchamp dans ses discours (Marcel Duchamp, \teiHi[rend={italic}]{op. cit}., p. 182) doit être avant tout comprise comme une neutralisation du « bon goût » et des préjugés qui le sous-tendent. En nous libérant des injonctions d’une esthétique conventionnelle, Duchamp parvient à nous rendre sensible l’expérience de l’inframince dans toute sa contingence, et à renouer ainsi avec l’esthétique au sens étymologique du terme. L’interprétation de Davila nous éloigne donc drastiquement de celle, nominaliste, d’un Thierry de Duve (\teiHi[rend={italic}]{Nominalisme pictural. Marcel Duchamp, la peinture et la modernité}, Paris, Éditions de Minuit, 1984) ou de celle, ontologique, d’un Danto (\teiHi[rend={italic}]{La transfiguration du banal. Une philosophie de l’art}, Claude Hary-Schaeffer trad., Jean-Marie Schaeffer préf., Paris, Éditions du Seuil, 1989). À noter que c’est précisément dans le grand magasin évoqué plus haut que Duchamp a acheté son fameux porte-bouteilles, destiné à devenir le premier « véritable » readymade. Voir à ce propos la notice de Didier Ottinger, « Marcel Duchamp, \teiHi[rend={italic}]{Porte-bouteilles}, 1914/1964 », sur le site du Centre Pompidou, https://www.centrepompidou.fr/fr/ressources/oeuvre/cLz5nr (consulté le 19 juin 2025).}} ». Si le régime contemporain de l’art consiste bien, comme l’affirme Davila, en une extension du « champ de la perception », la place de la beauté n’en est selon nous en aucun cas minorée\teiNote[xmlid={ftn123-2},n={60},place={foot}]{\teiP[xmlid={p-206}]{Elle en ressort au contraire renforcée, face à ce qui était jusqu’alors le paradigme dominant du sublime dans la pensée sur l’art contemporain. Ce « retour » inframince de la beauté est à la fois plus secret et plus radical que celui qu’a fini par reconnaître Arthur Danto, après moult circonvolutions, dans \teiHi[rend={italic}]{The Abuse of Beauty: Aesthetics and the Concept of Art}, Chicago, Open Court, 2003.}}, puisqu’elle se niche précisément dans cette « quête des variations infimes » qui dévoile, dans son ordinaire même, l’extraordinaire puissance poétique de notre vie sensible – dont « l’aisthétique » cherche en retour à percer le sens\teiNote[xmlid={ftn124-2},n={61},place={foot}]{\teiP[xmlid={p-207}]{C’est ce que confirme le défrichement actuel, en relation directe avec « l’aisthétisation » de l’esthétique, du champ de « l’esthétique du quotidien ». Voir à ce propos Yuriko Saito, « Aesthetics of the Everyday », dans Edward N. Zalta (dir.), \teiHi[rend={italic}]{The Stanford Encyclopedia of Philosophy}, édition du printemps 2021, \teiRef[target={https://plato.stanford.edu/archives/spr2021/entries/aesthetics-of-everyday/}]{\teiHi[rend={underline}]{https://plato.stanford.edu/archives/spr2021/entries/aesthetics-of-everyday/}} ; Barbara Formis, \teiHi[rend={italic}]{Esthétique de la vie ordinaire}, Paris, PUF, 2010 (en particulier p. 38-45) ; ou encore Laurent Jenny, \teiHi[rend={italic}]{La vie esthétique : stases et flux}, Lagrasse, Verdier, 2013. Cette attention à la poétique ordinaire de notre vie sensible – loin, donc, de tout élitisme de distinction – constitue selon nous un déplacement analogue, dans le domaine esthétique, au déplacement opéré en éthique par la question du \teiHi[rend={italic}]{care} (voir par exemple Marie Garrau, Care\teiHi[rend={italic}]{ et attention}, Paris, PUF, 2014, mais également le texte de Gisèle Dambuyant dans le présent volume).}}.}
              \teiP[xmlid={p19-3}]{Cette exploration de l’inframince, pour celles et ceux qui ont suivi les activités de la chaire Beauté·s, nous semble trouver dans les travaux de l’artiste plasticienne Ann Veronica Janssens ou dans la musique \teiHi[rend={italic}]{noise} de Romain Perrot, des exemples de choix\teiNote[xmlid={ftn125-2},n={62},place={foot}]{\teiP[xmlid={p-208}]{Nous pensons notamment à la Working Lecture « Structural Colors » du 14 juin 2022 (https://psl.eu/agenda/structural-colors, consulté le 19 juin 2025), ainsi qu’à la séance « Beautés de l’abject » du séminaire de la chaire Beauté·s, le 24 avril 2022 (https://psl.eu/agenda/seminaire-chaire-beautes, consulté le 19 juin 2025). Si le phénomène d’iridescence, au principe des couleurs structurelles de Janssens, est explicitement qualifié d’inframince par Duchamp, le caractère saturé de la \teiHi[rend={italic}]{harsh noise} semble \teiHi[rend={italic}]{a priori} s’opposer de toutes ses forces à cette poursuite de l’« imperceptible ». À y regarder de plus près, cependant, cette musique extrême conduit paradoxalement à un affinement de l’oreille, à une écoute minutieuse du matériau sonore, qui étend de fait le « champ de la perception » vers des régions inattendues et négligées, redéfinissant par là-même les limites de la justesse comme de la beauté musicales. Sur le rapport global de l’art de Janssens à l’inframince, voir Marianne Massin, \teiHi[rend={italic}]{Expérience esthétique et art contemporain}, Rennes, PUR, 2013, p. 104 et 105. Sur l’éducation de l’ouïe induite par la \teiHi[rend={italic}]{harsh noise}, voir Catherine Guesde et Pauline Nadrigny, \teiHi[rend={italic}]{The Most Beautiful Ugly Sound in the World. À l'écoute de la }noise, Paris, Éditions MF, 2018.}}. }
              \begin{teiFigure}[xmlid={figure01}]

                \teiHead[xmlid={figure-1-ann-veronica-janssens-en-collaboration-avec-maria-boto-ordonez-structural-colors-illustration-de-la-conference-du-14-juin-2022-paris-ecole-nationale-superieure-des-beaux-arts-laboratorium-kask-conservatorium-hogent-001}]{Figure 1. Ann Veronica Janssens, en collaboration avec María Boto Ordoñez, \teiHi[rend={italic}]{Structural Colors}, illustration de la conférence du 14 juin 2022, Paris, École nationale supérieure des beaux-arts. LABORATORIUM, KASK \& Conservatorium, HOGENT.}
              
\end{teiFigure}
              \teiP[xmlid={p20-3}]{Des exemples qui montrent combien l’art et la beauté ne sont en rien choses du passé, ou plutôt combien l’art, en dépit et peut-être même grâce à ses récentes métamorphoses, assume plus que jamais sa fonction de propédeutique à l’appréhension de la beauté du monde, naturel ou culturel, hors de toute norme préétablie. Si, en effet, il n’existe pas de beauté « en soi », celle-ci ne dépend-elle pas nécessairement, pour son émergence, sa reconnaissance et son partage, d’une « attitude caractéristique de l’esprit humain\teiNote[xmlid={ftn126-2},n={63},place={foot}]{\teiP[xmlid={p-209}]{Ernst Cassirer, « Qu’est-ce que la beauté ? », Fabien Capeillères (trad.), dans \teiHi[rend={italic}]{Écrits sur l’art}, Fabien Capeillères (éd. et postface) et John M. Krois (présentation), Paris, Cerf, 1995, p. 123.}} », dont l’art (en Occident tout du moins) a historiquement constitué le principal laboratoire ? Aussi conclurons-nous, avec Danièle Cohn, sur la nécessité pour l’esthétique d’approcher « en lumière rasante le travail des artistes ». Car c’est, dit-elle, en se mettant à leur école que l’esthétique, « poétique modeste », pourra poursuivre « la tâche qui est sa donne de départ : assurer la force formatrice de la culture, renouant ainsi avec l’idée d’une éducation esthétique dont le but est de constituer une sagesse sensible. La largesse de l’expérience esthétique, sa générosité, tient à l’édification d’une telle sagesse, et les artistes par leurs œuvres nous aident grandement à faire front contre une pauvreté en expérience\teiNote[xmlid={ftn127-2},n={64},place={foot}]{\teiP[xmlid={p-210}]{Danièle Cohn, « L’esthétique, tout simplement », \teiHi[rend={italic}]{Archives de Philosophie,} vol. 80, n\teiHi[rend={sup}]{o} 2, 2017, p. 230. Il y a là une proximité patente avec la volonté de Baptiste Morizot et Estelle Zhong Mengual de lutter contre une certaine « crise de la sensibilité » contemporaine, qui serait à l’origine de la crise environnementale et de la destruction systémique du vivant. Il nous incombera toutefois de montrer à l’avenir en quoi notre approche se distingue de leur Esthétique de la rencontre (Paris, Éditions du Seuil, 2018), du fait notamment de prémisses philosophiques différentes.}} ». Partant de ce constat, ce n’est autre que la figure du \teiHi[rend={italic}]{felix æstheticus}, telle que célébrée par Baumgarten, que cette attention renouvelée à la beauté de, ou plutôt \teiHi[rend={italic}]{dans} nos vies, contribue à ressusciter.}
            
\end{teiDiv}
          
\end{teiElement}
        
\end{teiElement}
      
\end{teiElement}
      \begin{teiElement}[name={group},type={section1},xmlid={section1-002}]

        \teiHead[xmlid={la-beaute-au-coeur-de-la-vie-et-du-vivant-001}]{La beauté au cœur de la vie et du vivant}
        \begin{teiElement}[name={group},type={chapter},data-page-title={Formes, diversité et beauté du vivant},data-page-subtitle={Morphologie, science et art de Richard Owen à Stephen Jay Gould},data-page-authors={Claudine Cohen@@EHESS, Centre de recherches sur les arts et le langage, France@@EPHE-PSL, 3e section Sciences de la vie et de la terre, France@@UBE-CNRS-EPHE, laboratoire Biogéosciences – UMR 6282, France},data-include-href={Ch05\_formes\_diversit\%C3\%A9\_beaut\%C3\%A9\_Cohen.xml},xmlid={chapter-004}]

          \begin{teiElement}[name={front}]

            \begin{teiDiv}[type={titlePage},xmlid={titlepage-006}]

              \teiP[rend={title-main},xmlid={p-211}]{Formes, diversité et beauté du vivant}
              \teiP[rend={title-sub},xmlid={p-212}]{Morphologie, science et art de Richard Owen à Stephen Jay Gould}
              \teiP[rend={author-aut},xmlid={p-213}]{Claudine \teiHi[rend={small-caps}]{Cohen}}
              \teiP[rend={authority\_affiliation},xmlid={p-214}]{EHESS, Centre de recherches sur les arts et le langage, France}
              \teiP[rend={authority\_affiliation},xmlid={p-215}]{EPHE-PSL, 3\teiHi[rend={sup}]{e} section Sciences de la vie et de la terre, France}
              \teiP[rend={authority\_affiliation},xmlid={p-216}]{UBE-CNRS-EPHE, laboratoire Biogéosciences – UMR 6282, France}
            
\end{teiDiv}
          
\end{teiElement}
          \begin{teiElement}[name={body},xmlid={body-006}]

            \teiP[xmlid={p1-6}]{La pensée biologique issue de la \teiHi[rend={italic}]{Naturphilosophie} allemande a eu, à partir de la fin du \teiHi[rend={small-caps}]{xviii}\teiHi[rend={sup}]{e} siècle, une incidence majeure sur la représentation du vivant et de son devenir\teiNote[xmlid={ftn1-5},n={1},place={foot}]{\teiP[xmlid={p-217}]{Voir Claudine Cohen, « De la biologie au roman : le modèle morphologique et ses variations », \teiHi[rend={italic}]{Romantisme}, n\teiHi[rend={sup}]{o }4, 2007, p. 47-59. }}. Approche à la fois philosophique, scientifique et artistique des êtres, de leur formation et des processus qui les animent, la \teiHi[rend={italic}]{Naturphilosophie} propose la notion d’un modèle morphologique qui rendrait compte, à partir d’une matrice idéale, de toute leur variabilité. Au moment de l’émergence de la biologie comme science, cette conception engage la notion d’un ordre caché de la nature qui régit l’organisation et les transformations du vivant et relie entre eux tous les êtres du monde. }
            \teiP[xmlid={p2-6}]{Au cœur de cette biologie romantique\teiNote[xmlid={ftn34},n={2},place={foot}]{\teiP[xmlid={p-218}]{Voir notamment Timothy Lenoir, « The Göttingen School and the development of Transcendantal Naturphilosophie in the Romantic Era », \teiHi[rend={italic}]{Studies in History of Biology}, n\teiHi[rend={sup}]{o} 5, 1981, p. 111-205.}}, la morphologie transcendantale repose en effet sur la notion d’un « plan commun de composition », et celle d’un « archétype » qui permet de désigner, tout en la rendant intelligible, cette organisation morphologique commune aux êtres vivants. On doit à Goethe d’avoir formulé les principes de cette \teiHi[rend={italic}]{morphologie} et théorisé la notion de forme, à la rencontre des sciences de la nature, de la philosophie et de l’art : }
            \begin{teiCit}[xmlid={cit01-4}]

              \begin{teiQuote}
On rencontre \lbrack{}…\rbrack{} dans le cheminement de l’art, du savoir et de la science, plusieurs tentatives pour fonder et développer une connaissance que nous aimerions appeler la morphologie. \lbrack{}…\rbrack{}
\end{teiQuote}
            
\end{teiCit}
            \begin{teiCit}[xmlid={cit02-4}]

              \begin{teiQuote}
Pour désigner dans son ensemble l’existence d’un être réel, l’allemand dispose du mot « forme » \lbrack{}\teiHi[rend={italic}]{Gestalt}\rbrack{}. En employant ce terme, il fait abstraction de ce qui est mobile, il admet que les éléments formant un tout sont établis, achevés et fixés dans leurs caractères.
\end{teiQuote}
            
\end{teiCit}
            \begin{teiCit}[xmlid={cit03-4}]

              \begin{teiQuote}
Mais si nous observons toutes les formes, et en particulier les formes organiques, nous constatons qu’il ne se trouve nulle part de constance, d’immobilité, d’achèvement, et qu’au contraire tout oscille dans un mouvement incessant. C’est pourquoi notre langue se sert à fort juste titre du mot formation \lbrack{}\teiHi[rend={italic}]{Bildung}\rbrack{}, tant pour désigner ce qui est produit que ce qui est en voie de l’être.
\end{teiQuote}
            
\end{teiCit}
            \begin{teiCit}[xmlid={cit04-4}]

              \begin{teiQuote}
Si donc nous voulons introduire une morphologie, nous n’avons pas à parler de forme, mais si nous employons ce terme, nous pouvons penser tout au plus l’idée, le concept, ou un élément fixé pour un instant seulement dans l’expérience\teiNote[xmlid={ftn35},n={3},place={foot}]{\teiP[xmlid={p-219}]{Johann Wolfgang von Goethe, \teiHi[rend={italic}]{La métamorphose des plantes, Objet et méthode de la morphologie}, Henriette Bideau (trad.), Laboissière-en-Thelle, Éditions Triades, 2013.}}.
\end{teiQuote}
            
\end{teiCit}
            \teiP[xmlid={p3-6}]{En définissant la morphologie, Goethe met ainsi en rapport les formes vivantes avec la notion de leur transformation. La tradition de la morphologie parcourt l’histoire de la biologie et de la paléontologie évolutive, depuis Goethe jusqu’à ses théorisations contemporaines. À l’admiration de la beauté statique des formes vivantes s’oppose une approche à la fois biologique et esthétique de leur plasticité. Il s’agit alors de comprendre les formes vivantes comme des formes en mouvement, en devenir, en insistant sur la beauté qui naît de la dynamique de leurs mouvements, de leur croissance et de leurs transformations, ces phénomènes dynamiques étant inscrits dans les structures mêmes du vivant. }
            \teiP[xmlid={p4-6}]{Tout au long de l’histoire de la biologie évolutive, cet intérêt pour les formes et leurs transformations connaît des éclipses et des résurgences. Il est remarquable qu’à chacun de ces épisodes, l’approche morphologique s’articule à un regard esthétique, voire artistique. Nous explorerons ici trois grands moments de cette histoire, particulièrement significatifs à cet égard : }
            \begin{teiList}[xmlid={list1-2},rendition={\#},type={unordered}]

    \teiItem{la première moitié du \teiHi[rend={small-caps}]{xix}\teiHi[rend={sup}]{e} siècle en Angleterre, autour de la figure du grand naturaliste Richard Owen qui développe une conception du vivant visant à concilier structuralisme biologique et morphologie transcendantale, et à concilier théologie et esthétique ;}
    \teiItem{les dernières décennies du \teiHi[rend={small-caps}]{xix}\teiHi[rend={sup}]{e} siècle, qui voient s’articuler, chez l’embryologiste allemand Ernst Haeckel, une philosophie évolutive reposant sur un vitalisme mystique et un regard esthétique porté sur les formes vivantes ;}
    \teiItem{enfin, le dernier tiers du \teiHi[rend={small-caps}]{xx}\teiHi[rend={sup}]{e} siècle, au cours duquel, après avoir été longtemps oubliées au profit des approches microévolutives, la morphologie des organismes et la « macro-évolution » sont réintégrées à l’évolutionnisme néodarwinien, chez des paléontologues tels que l’Américain Stephen Jay Gould ou l’Allemand Adolf Seilacher. Ce dernier développe cette approche par la proposition de la notion d’un « art fossile ». }
   
\end{teiList}
            \teiP[xmlid={p5-6}]{De Goethe à Owen, de Haeckel à Gould et à Seilacher, se perpétue une remarquable insistance sur cette double approche, vitale et esthétique, des formes vivantes. Nous montrerons comment la notion de beauté se trouve chaque fois au cœur de ces réflexions et s’articule, dans ces différents contextes scientifiques et philosophiques, à une approche renouvelée du vivant et de son devenir. }
            \begin{teiDiv}[type={section1},xmlid={div01-3}]

              \teiHead[xmlid={richard-owen-la-beaute-de-ladaptation-001}]{Richard Owen : la beauté de l’adaptation}
              \teiP[xmlid={p6-6}]{Figure passablement oubliée aujourd’hui, le naturaliste anglais Richard Owen (1804-1892) fut pourtant une autorité scientifique de premier plan dans la science anglaise du \teiHi[rend={small-caps}]{xix}\teiHi[rend={sup}]{e} siècle. Brillant anatomiste et paléontologue, c’est à lui que Darwin confie les fossiles rapportés de son voyage du Beagle pour les étudier et les publier\teiNote[xmlid={ftn36},n={4},place={foot}]{\teiP[xmlid={p-220}]{Richard Owen, \teiHi[rend={italic}]{Fossil Mammalia. Part I. The Zoology of the Voyage of HMS Beagle: Under the Command of Captain Fitzroy-During the Years 1832 to 1836} \lbrack{}1839\rbrack{}, Redditch, Read Books Ltd, 2019.}}. Il est l’inventeur du nom des dinosaures (1842), identifie l’espèce du gorille (1848) et écrit la première monographie sur l’\teiHi[rend={italic}]{archæopteryx }(1863\teiNote[xmlid={ftn37},n={5},place={foot}]{\teiP[xmlid={p-221}]{Richard Owen, « On the Archeopteryx of von Meyer, with a description of the fossil remains of a long-tailed species, from the lithographic stone of Solenhofen », \teiHi[rend={italic}]{Philosophical Transactions of the Royal Society of London}, 1863, p. 33-47.}})\teiHi[rend={italic}]{.} C’est aussi un homme d’institutions : membre du Royal College of Surgeons, professeur de physiologie à la Royal Institution, surintendant des collections d’Histoire naturelle de Bloomsbury, il fonde le Natural History Museum de Londres en 1863.}
              \teiP[xmlid={p7-6}]{Owen est en son temps un des plus importants théoriciens des formes du vivant et de leur devenir\teiNote[xmlid={ftn38},n={6},place={foot}]{\teiP[xmlid={p-222}]{Voir Claudine Cohen, « Richard Owen : paléontologie, embryologie et morphologie transcendantale vers 1840 », dans Olivier Bloch (dir.), \teiHi[rend={italic}]{Philosophies de la nature}, Paris, Publications de la Sorbonne, 2000, p. 301-309 et \teiHi[rend={italic}]{La Méthode de Zadig. La trace, le fossile, la preuve.} Paris, Éditions du Seuil, 2011.}}. À partir de 1840, il propose d’intégrer la morphologie à la paléontologie, et tente de concilier la conception de l’unité fonctionnelle et adaptative de l’organisme avec la thèse, inspirée par les morphologistes allemands, d’un devenir des espèces rapporté à un plan d’organisation unique qui se diversifie dans la multitude des espèces connues. Owen vise ainsi à accorder les « deux grandes lumières » de l’histoire naturelle française, Georges Cuvier et Étienne Geoffroy Saint-Hilaire, qui s’étaient affrontées lors d’un débat célèbre à l’Académie des Sciences de Paris en 1830\teiNote[xmlid={ftn39},n={7},place={foot}]{\teiP[xmlid={p-223}]{Voir Toby Appel,\teiHi[rend={italic}]{ The Cuvier-Geoffroy Debate, French Biology in the decades before Darwin}, Oxford, Oxford University Press, 1987.}}. }
              \teiP[xmlid={p8-6}]{À côté de ses grands traités de paléontologie, deux ouvrages de « philosophie anatomique »\teiHi[rend={italic}]{ – Report on the Archetype and Homologies of the Vertebrate Skeleton}, 1847 (discours prononcé en 1846 devant l’Association britannique pour le progrès des sciences) et \teiHi[rend={italic}]{On the Nature of Limbs}, 1849 – posent les bases de cette « morphologie transcendantale » (le mot est introduit en France par Étienne Serre dès 1820). }
              \teiP[xmlid={p9-5}]{Au centre de sa construction théorique et de sa pratique d’anatomiste\teiNote[xmlid={ftn40},n={8},place={foot}]{\teiP[xmlid={p-224}]{Voir Claudine Cohen, \teiHi[rend={italic}]{op. cit.}, chap. IV.}}, Owen théorise deux notions essentielles pour penser les formes vivantes et leur devenir, celle d’\teiHi[rend={italic}]{homologie} et celle d’\teiHi[rend={italic}]{archétype}.}
              \begin{teiDiv}[type={section2},xmlid={div02-3}]

                \teiHead[xmlid={le-concept-dhomologie-001}]{Le concept d’Homologie }
                \teiP[xmlid={p10-5}]{Owen définit le concept morphologique d’\teiHi[rend={italic}]{homologie}, inspiré du « principe de connexion » mis en place par Étienne Geoffroy Saint-Hilaire, et défini comme la variation, rapportée à un modèle unique, d’un même type de structure dans deux ou plusieurs organismes différents. Owen distingue soigneusement l’analogie et l’homologie : dans des espèces différentes, l’\teiHi[rend={italic}]{analogie} qualifie des éléments ayant une fonction identique, tandis que l’\teiHi[rend={italic}]{homologie} désigne le statut d’un même organe sous toutes les variétés possibles de ses formes et de ses fonctions : ainsi, l’aile du reptile \teiHi[rend={italic}]{Draco volans} est \teiHi[rend={italic}]{analogue} à celle de l’oiseau (elle remplit la même fonction sans avoir aucun lien anatomique avec elle), alors que les pattes antérieures du lézard sont les \teiHi[rend={italic}]{homologues} des ailes de l’oiseau, car elles dérivent d’une même structure anatomique (le membre antérieur des vertébrés). }
              
\end{teiDiv}
              \begin{teiDiv}[type={section2},xmlid={div03-3}]

                \teiHead[xmlid={larchetype-001}]{L’archétype}
                \teiP[xmlid={p11-5}]{Owen définit d’autre part la notion d’« Archétype » – type idéal rendant compte à la fois de l’unité et de la diversité du vivant\teiNote[xmlid={ftn41-2},n={9},place={foot}]{\teiP[xmlid={p-225}]{Pour une analyse de cette notion, voir Nicolaas A. Rupke, « Richard Owen’s Vertebrate Archetype », \teiHi[rend={italic}]{ISIS}, n\teiHi[rend={sup}]{o} 84, 1993, p. 231-251.}}. Il se réfère à Platon : « Quel est l’animal qui servit de modèle à Celui qui a fait ce grand animal, le Monde ? \lbrack{}…\rbrack{} L’animal qui servit d’Archétype est celui qui contient en lui-même tous les autres animaux comme ses parties constituantes\teiNote[xmlid={ftn42-2},n={10},place={foot}]{\teiP[xmlid={p-226}]{Richard Owen, \teiHi[rend={italic}]{Principes d’ostéologie comparée ou Recherches sur l’archétype et les homologies du squelette vertébré}, Paris, Baillère, 1855.}}. » Un « plan », un « modèle », une « idée » uniques président à l’organisation anatomique de tous les vertébrés. Cet exemplaire idéal, qui est à la base de l’organisation des animaux, peut être reconstitué à travers \teiHi[rend={italic}]{les homologies} que l’on reconnaît dans leur organisation : ainsi, « les belles et nombreuses évidences d’unité de plan que l’étude de la structure des membres locomotifs a mises en lumière ». L’archétype est la forme fondamentale (la plus simple et la plus « primitive ») qui rend raison de la structure de tous les organismes qui en dérivent. Ainsi, l’archétype des vertébrés selon Richard Owen est l’animal à la structure simple à partir de laquelle se sont diversifiées les autres espèces de la même classe : il s’agit d’un poisson osseux, le \teiHi[rend={italic}]{lepidosiren}, dont la structure du squelette dessine le plan auquel on peut rapporter par homologie ceux de tous les vertébrés, les poissons, les reptiles, les oiseaux, les quadrupèdes et l’homme lui-même. Il est possible dès lors de penser l’\teiHi[rend={italic}]{adaptation} comme le processus par lequel des parties ou organes d’un être vivant se transforment en rapport avec un milieu donné. La beauté du vivant réside précisément dans la perfection de l’adaptation : }
                \begin{teiCit}[xmlid={cit05-3}]

                  \begin{teiQuote}
Perfection de l’adaptation : La satisfaction ressentie par un esprit bien constitué est grande lorsqu’il reconnaît l’adaptation des parties à leurs fonctions appropriées. Mais lorsque cette adaptation est atteinte, comme dans le gros orteil du pied de l’Homme et de l’Autruche, par une structure qui manifeste en même temps une concordance harmonieuse avec le type commun, le pouvoir de la Grande Cause de toute l’organisation est apprécié aussi pleinement qu’il est possible par notre intelligence limitée\teiNote[xmlid={ftn43-2},n={11},place={foot}]{\teiP[xmlid={p-227}]{Richard Owen, \teiHi[rend={italic}]{On the Nature of the Limbs}, Londres, John Van Voorst, 1849, p. 9. Toutes les traductions sont de notre fait.}}.
\end{teiQuote}
                
\end{teiCit}
                \teiP[xmlid={p12-5}]{« Chez l’Homme », ajoute Owen, « le principe de l’adaptation spéciale va plus loin. Tandis qu’une paire de membres est expressément organisée pour la locomotion et la station redressée, l’autre paire est laissée libre pour exécuter les multiples ordres de sa volonté rationnelle et inventive, et elle est remarquablement organisée pour le toucher délicat et la préhension, que l’on désigne du terme de “manipulation\teiNote[xmlid={ftn44-2},n={12},place={foot}]{\teiP[xmlid={p-228}]{\teiHi[rend={italic}]{Ibid}., p. 38.}}” ».}
                \teiP[xmlid={p13-4}]{Ainsi, chez Owen, la pensée de la forme rend compte à la fois de l’unité et de la diversité du vivant. L’adaptation de l’être vivant à son milieu par ses fonctions est admirable, mais elle n’est pas la réalisation directe d’une finalité. Elle résulte de la \teiHi[rend={italic}]{réutilisation} de structures préexistantes dans une adaptation anatomique donnée. }
                \teiP[xmlid={p14-4}]{Owen croit au devenir des espèces, sous la forme d’une « tendance innée » des êtres à se métamorphoser au cours de leur histoire. Pour lui, les faunes se renouvellent par l’action d’une puissance créatrice continue dans la nature : il y a un « devenir ordonné des choses vivantes » qui se traduit dans la genèse successive des espèces. « La création de chaque classe animale, Reptile, Oiseaux, Quadrupèdes, a été successive et continue depuis les temps les plus anciens\teiNote[xmlid={ftn45-2},n={13},place={foot}]{\teiP[xmlid={p-229}]{Richard Owen, \teiHi[rend={italic}]{Presidential Address at the 28th Meeting of the British Association}, Londres, 22 septembre 1858, p. 3.}}. » Cette création progressive d’espèces toujours bien adaptées à des environnements changeants ne peut se comprendre que par l’intervention d’une « intelligence anticipatrice », d’une puissance divine. }
                \teiP[xmlid={p15-4}]{Ainsi la notion d’une « dérivation des formes vivantes » se trouve-t-elle encadrée par une théologie. Cette approche de la nature, qui mêle morphologie transcendantale, notion d’une « harmonie de la nature » et continuité des espèces, est donc à la fois esthétique et théologique. « La dérivation \lbrack{}…\rbrack{} est une manifestation du pouvoir créatif vers la variété et la beauté du résultat », écrit Owen\teiNote[xmlid={ftn46-2},n={14},place={foot}]{\teiP[xmlid={p-230}]{Richard Owen, \teiHi[rend={italic}]{On the Anatomy of Vertebrates}, Londres, Longmans, Green \& Co, 1868, vol. 3, p. 808. }}.}
                \teiP[xmlid={p16-4}]{Les fortes implications esthétiques, voire artistiques, de cette approche conjuguée de la forme et du mouvement se traduisent notamment dans l’œuvre du sculpteur Benjamin Waterhouse Hawkins (1807-1894), qui collabore avec Owen pour mettre au point ses reconstitutions animalières selon les principes anatomiques et morphologiques du maître. Hawkins est le sculpteur et le metteur en scène des fameuses reconstitutions grandeur nature des animaux fossiles étudiés par Owen, et notamment des dinosaures (mégalosaure, iguanodon) présentées en 1854 au public londonien du Crystal Palace. }
                \teiP[xmlid={p17-4}]{À ses étudiants en art, Waterhouse Hawkins enseigne les principes de la morphologie animale directement inspirés d’Owen. Dans l’introduction de son ouvrage intitulé \teiHi[rend={italic}]{Une vision comparée du corps humain et animal}\teiNote[xmlid={ftn47-2},n={15},place={foot}]{\teiP[xmlid={p-231}]{Benjamin Waterhouse Hawkins, \teiHi[rend={italic}]{A comparative view of Human and Animal Frame,} Londres, Chapman and Hall, 1860.}}, Hawkins déclare son intention de transmettre à ses étudiants « la conviction forte de l’unité de dessin et de l’unicité de plan sur laquelle tous les animaux sont construits : une unité toujours si apparente et importante pour le naturaliste lorsqu’il examine et compare n’importe laquelle des grandes classes du règne animal\teiNote[xmlid={ftn48-2},n={16},place={foot}]{\teiP[xmlid={p-232}]{\teiHi[rend={italic}]{Ibid}., « Introduction », n. p. }} ». Dans les planches qui accompagnent le texte, Hawkins met en évidence :}
                \begin{teiCit}[xmlid={cit06-3}]

                  \begin{teiQuote}
\lbrack{}P\rbrack{}ar la répétition des formes, le fait qu’un modèle élémentaire a été créé et fixé par l’architecte divin au commencement, et qu’il s’est perpétué dans toute la durée des temps, jusqu’à aujourd’hui. Or, ce modèle était si parfait (car élaboré et préconçu par la sagesse omnisciente) que de légères modifications de parties secondaires permettent d’adapter la totalité d’un organisme à toutes les circonstances changeantes qui ont été, ou peuvent devenir, les conditions de vie spécifiques aux différents groupes d’êtres qui constituent les embranchements du règne animal\teiNote[xmlid={ftn49-2},n={17},place={foot}]{\teiP[xmlid={p-233}]{\teiHi[rend={italic}]{Ibid}.}}.
\end{teiQuote}
                
\end{teiCit}
                \begin{teiFigure}[xmlid={figure01-2}]

                  \teiHead[xmlid={figure-2-benjamin-waterhouse-hawkins-plate-two-man-and-the-lion-dans-a-comparative-view-of-the-human-and-animal-frame-londres-chapman-and-hall-1860-p-12-et-13-university-of-wisconsin-madison-cc-by-001}]{Figure 2. Benjamin Waterhouse Hawkins, « Plate two – Man, and the lion », dans \teiHi[rend={italic}]{A comparative view of the human and animal frame}, Londres, Chapman and Hall, 1860, p. 12 et 13. University of Wisconsin-Madison, CC-BY.}
                
\end{teiFigure}
                \teiP[xmlid={p18-4}]{Pour mettre en évidence les homologies des structures squelettiques qui constituent les architectures animales, il sélectionne « certains animaux parmi les plus importants ou les plus familiers, qui doivent faire partie de manière presque quotidienne de chaque composition artistique\teiNote[xmlid={ftn50-2},n={18},place={foot}]{\teiP[xmlid={p-234}]{\teiHi[rend={italic}]{Ibid}.}} ». Le cheval, le chien, le lion, l’éléphant, le chameau sont tour à tour convoqués, comparés : l’anatomie du détail de leurs parties (le pied, la main, le cou) met en évidence leurs homologies. L’Homme apparaît comme le point de comparaison ultime et la structure la plus parfaite. On ne saurait mieux montrer l’usage des principes mis en place par Owen pour l’étude des formes vivantes et de leur devenir au service de l’art et de la beauté. }
              
\end{teiDiv}
            
\end{teiDiv}
            \begin{teiDiv}[type={section1},xmlid={div04}]

              \teiHead[xmlid={ernst-haeckel-1834-1919-et-les-formes-artistiques-de-la-nature-001}]{Ernst Haeckel (1834-1919) et les « formes artistiques de la nature »}
              \teiP[xmlid={p19-4}]{L’œuvre du naturaliste et embryologiste allemand Ernst Haeckel élargit le champ de la morphologie, pour lui donner une dimension naturaliste, culturelle et artistique. Chez lui la morphologie comparée fonde l’étude des formes vivantes, elle s’applique aussi à celles des sociétés, des langues, des objets techniques ; elle contribue ainsi à une visée unificatrice, par laquelle les significations naturelles et culturelles de la « forme » se répondent. }
              \teiP[xmlid={p20-4}]{Professeur d’embryologie et d’anatomie comparée à l’université d’Iéna, passionné d’art et de dessin, Ernst Haeckel fait côtoyer dans son œuvre ses intérêts esthétiques avec une vision anthropocentriste de l’évolution du vivant, et une approche fortement raciste de l’évolution humaine\teiNote[xmlid={ftn51-2},n={19},place={foot}]{\teiP[xmlid={p-235}]{Voir Robert Richards, \teiHi[rend={italic}]{The Tragic Sense of Life: Ernst Haeckel and the Struggle over Evolutionary Thought,} Chicago, The University of Chicago Press, 2008.}}. Spécialiste des invertébrés marins, voyageur et travailleur infatigable, il sillonne les océans tropicaux, les mers du Nord des côtes de la Norvège à la Finlande et la Russie, du Mexique, voyage en Palestine et en Asie mineure. Ses recherches en zoologie et en biologie marine le conduisent à décrire et à classer des milliers d’espèces d’invertébrés aquatiques (il est l’auteur d’importantes monographies consacrées aux radiolaires, aux éponges calcaires, aux méduses et aux siphonophores), dont il magnifie la beauté par le dessin. Dans la mouvance de la morphologie allemande, Haeckel souligne et illustre ainsi de manière éclatante la richesse esthétique du monde vivant. Entre 1899 et 1904, il publie \teiHi[rend={italic}]{Kunstformen der Natur}, une série de cent lithographies en couleur, destinée au grand public, et visant à mettre en évidence « le domaine infini des formes de vie primaires vivant dans les profondeurs de l’océan », révélant leur symétrie, leur organisation et tout le merveilleux chatoiement de leurs couleurs. }
              \teiP[xmlid={p21-3}]{« La nature génère une abondance inépuisable de formes merveilleuses, dont la beauté et la diversité dépassent de loin toutes les formes d’art produites par l’homme », écrit Haeckel dans l’introduction de ses \teiHi[rend={italic}]{Kunstformen der Natur}\teiNote[xmlid={ftn52-2},n={20},place={foot}]{\teiP[xmlid={p-236}]{Ernst Haeckel,\teiHi[rend={italic}]{ Kunstformen der Natur}, Leipzig-Vienne, Bibliographisches Institut, 1904.}}, y affirmant aussi sa conviction qu’un art réformé, naturaliste, peut aider le peuple à s’émanciper des autorités religieuses et politiques qui le dominent en favorisant l’ignorance et la superstition.}
              \begin{teiFigure}[xmlid={figure02}]

                \teiHead[xmlid={figure-3-ernst-haeckel-kunstformen-der-natur-leipzig-vienne-bibliographisches-institut-1904-planche-8-discomedusa-bnf-gallica-departement-estampes-et-photographie-pet-fol-hd-77-001}]{Figure 3. Ernst Haeckel, \teiHi[rend={italic}]{Kunstformen der Natur}, Leipzig-Vienne, Bibliographisches Institut, 1904, planche 8, \teiHi[rend={italic}]{Discomedusa}. BNF-Gallica, département Estampes et photographie, PET FOL-HD-77.}
              
\end{teiFigure}
              \teiP[xmlid={p22-3}]{Son livre illustré de planches somptueuses donne à voir « ces organismes aux belles formes \lbrack{}jusque-là\rbrack{} dissimulées dans des ouvrages rares et coûteux et peu accessibles aux amateurs » et connaît un très large succès public. Cette vision des « formes artistiques de la nature » rejoint, en ce tournant du \teiHi[rend={small-caps}]{xx}\teiHi[rend={sup}]{e} siècle, le programme artistique du \teiHi[rend={italic}]{Jugendstil }ou Art nouveau – un art ornemental « où triomphaient la libellule, le marronnier, l’iris, les lys et l’orchidée\teiNote[xmlid={ftn53-2},n={21},place={foot}]{\teiP[xmlid={p-237}]{Louis Aragon, préface à Roger-Henri Guerrand, \teiHi[rend={italic}]{L’Art nouveau en Europe}, Paris, Perrin, 2009.}} » et qui, prônant une libération par rapport aux styles rigides du passé, s’inspire des formes de la nature, reproduisant à l’infini dans les œuvres d’art, les affiches, l’architecture et les décors de la vie quotidienne, leurs élégantes courbes et leurs libres volutes. }
              \teiP[xmlid={p23-3}]{Tout en se proclamant disciple de Darwin, Haeckel professe une vénération quasi mystique des formes naturelles et de la force vitale qui les anime. Il exalte la beauté des formes vivantes comme l’effet d’une \teiHi[rend={italic}]{vis vitalis} (force vitale) qui organise leur symétrie, leurs proportions, leur exubérance et leur diversité. Son approche du vivant s’articule à une philosophie moniste, dont le « Beau », le « Bien » et le « Vrai » constitue les piliers et pour laquelle il n’y a pas de différence fondamentale entre la nature organique et la nature inorganique : la vie ne diffère de la nature inorganique que par le degré de son organisation. }
            
\end{teiDiv}
            \begin{teiDiv}[type={section1},xmlid={div05}]

              \teiHead[xmlid={les-metamorphoses-de-la-morphologie-dans-la-paleontologie-evolutive-du-xx-e-siecle-001}]{Les métamorphoses de la morphologie dans la paléontologie évolutive du \teiHi[rend={small-caps}]{xx}\teiHi[rend={sup}]{e} siècle}
              \teiP[xmlid={p24-3}]{Le naturaliste écossais D’Arcy Wentworth Thompson (1860-1948) est animé par des principes philosophiques tout différents lorsqu’il élabore sa pensée profondément empreinte d’une méditation sur la morphologie des êtres vivants. Auteur d’un grand livre intitulé \teiHi[rend={italic}]{Forme et croissance} (qui connut deux éditions, 1917 et 1942, la seconde étant considérablement augmentée\teiNote[xmlid={ftn54-2},n={22},place={foot}]{\teiP[xmlid={p-238}]{D’Arcy Wentworth Thompson, \teiHi[rend={italic}]{On Growth and Form}, Cambridge, Cambridge University Press, 1917 ;\teiHi[rend={italic}]{ Forme et Croissance} \lbrack{}\teiHi[rend={italic}]{On Growth and Form}, Cambridge, Cambridge University Press, 1942\rbrack{}, Paris, Éditions du Seuil, 1994.}}), D’Arcy Thompson est à la fois physicien et zoologiste, et c’est en pur mécaniste qu’il considère les structures vivantes, modelées selon lui par les forces qui agissent dans le monde physique. Ces structures sont façonnées de façon optimale par des règles physiques et mathématiques simples, pour l’usage et le mode de vie de l’animal\teiNote[xmlid={ftn55-2},n={23},place={foot}]{\teiP[xmlid={p-239}]{Voir Claudine Cohen, « Gould et D’Arcy Thompson », \teiHi[rend={italic}]{Comptes Rendus Palevol}, 2004, vol. 3, n\teiHi[rend={sup}]{o} 5, p. 421-431.}}.}
              \teiP[xmlid={p25-2}]{L’organisme est un solide, sa structure et son adaptation peuvent être pensées sur des modèles mécaniques : ainsi la structure du squelette d’un diplodocus trouve son élucidation dans la construction et la statique des ponts, tandis que la forme d’une méduse évoque la chute d’une goutte d’eau. Les considérations de taille, de surface, de volume, de poids, de vitesse et d’économie d’espace jouent un rôle essentiel dans cette approche de l’organisme considéré « avec l’œil de l’ingénieur ». }
              \teiP[xmlid={p26-2}]{Si D’Arcy admet l’idée d’une évolution du vivant, il refuse la sélection naturelle comme facteur d’innovation biologique : ce sont les forces physiques qui selon lui s’imposent directement aux organismes et les transforment. On comprend que dans ce cadre il soit difficile de penser une évolution continue : il existe à ses yeux une discontinuité nécessaire entre les « types » vivants, de même qu’il est impossible de passer de manière continue d’une forme géométrique à une autre. Ses « grilles de coordonnées transformées » rendent compte de l’effet de l’imposition des forces physiques pour modeler la forme de l’animal.}
              \teiP[xmlid={p27-2}]{La pensée de D’Arcy Thompson reste marginale dans les développements de la biologie de la première moitié du \teiHi[rend={small-caps}]{xx}\teiHi[rend={sup}]{e} siècle, en partie à cause de son refus du darwinisme. Cependant, à la fin des années 1960, le jeune Stephen Jay Gould (1942-2002) trouve chez D’Arcy Thompson des éléments qui lui permettent de réintroduire la notion de forme et d’organisme dans la pensée évolutionniste\teiNote[xmlid={ftn56-2},n={24},place={foot}]{\teiP[xmlid={p-240}]{Stephen Jay Gould, « D’Arcy Thompson and the Science of Form », \teiHi[rend={italic}]{New Literary History}, n\teiHi[rend={sup}]{o} 2, 1971, p. 229-258.}} et ainsi d’engager un renouveau de la réflexion sur la macroévolution en un moment où les tenants de la synthèse néodarwinienne se concentrent surtout sur l’élucidation des mécanismes évolutifs au niveau microscopique du gène. }
              \teiP[xmlid={p28-2}]{C’est d’abord en s’inspirant de D’Arcy que Gould prône la réhabilitation de la morphologie. Il se propose d’étudier les transformations des organismes sur le plan morphologique de la croissance – en tenant compte des hétérochronies (retards et accélérations) du développement\teiNote[xmlid={ftn57-2},n={25},place={foot}]{\teiP[xmlid={p-241}]{Voir Stephen Jay Gould, \teiHi[rend={italic}]{Ontogeny and phylogeny}, Cambridge, Harvard University Press, 1977.}}. Il montre aussi comment, par « convergence », des animaux d’origine phylogénétique différente acquièrent des formes semblables lorsqu’ils sont adaptés à des conditions identiques. Ainsi, « le pouce du panda » est en fait une réutilisation secondaire d’une structure (le sésamoïde du radius) qui mime la fonction d’un « pouce » préhensile.}
              \begin{teiDiv}[type={section2},xmlid={div06}]

                \teiHead[xmlid={architectures-vivantes-001}]{Architectures vivantes }
                \teiP[xmlid={p29-2}]{Comme D’Arcy, Gould n’hésite pas à utiliser des modèles architecturaux pour mettre en lumière la notion d’exaptation\teiNote[xmlid={ftn58-2},n={26},place={foot}]{\teiP[xmlid={p-242}]{Stephen Jay Gould et Elizabeth Vrba, « Exaptation: a missing term in the science of form », \teiHi[rend={italic}]{Paleobiology}, n\teiHi[rend={sup}]{o} 8, 1982, p. 4-15.}} visant à élucider la manière dont évoluent les structures vivantes. Il observe que les « médaillons de la basilique Sant Marc à Venise » sont des éléments architecturaux qui n’ont d’existence que par rapport à la structure d’ensemble. Ils prennent cependant leur fonction propre dans l’édifice final par l’utilisation secondaire qui en est faite, et par la décoration qui leur est apposée\teiNote[xmlid={ftn59-2},n={27},place={foot}]{\teiP[xmlid={p-243}]{Stephen Jay Gould et Richard Lewontin, « The Spandrels of San Marco and the Panglossian Paradigm: a critique of the Panglossian Programme », \teiHi[rend={italic}]{Proceedings of the Royal Society of London}, vol. 205, n\teiHi[rend={sup}]{o} 1161, 1989, p. 581-598.}}. De même dans l’ordre du vivant, l’adaptation fonctionnelle résulte souvent d’un « bricolage » opéré à partir d’éléments présents par construction, mais fonctionnellement inutilisés. Ainsi Gould propose-t-il de repenser la notion de l’adaptation comme \teiHi[rend={italic}]{exaptation}, dégagée de tout substrat finaliste. Il s’éloigne ici de D’Arcy Thompson, considérant que l’adaptation ne consiste pas dans l’imposition directe, par les forces physiques, d’une forme biomécanique optimale sur la matière organique. Il emprunte alors au paléoichnologiste allemand Adolf Seilacher la notion de « contraintes architecturales\teiNote[xmlid={ftn60-2},n={28},place={foot}]{\teiP[xmlid={p-244}]{Adolf Seilacher, « Arbeitkonzept zur Konstruktionmorphologie », \teiHi[rend={italic}]{Letheia}, n\teiHi[rend={sup}]{o} 3, 1970, p. 393-396.}} » : chez tout être vivant, l’adaptation doit tenir compte des spécificités propres à la structure de l’organisme, qui résultent de son histoire. }
              
\end{teiDiv}
              \begin{teiDiv}[type={section2},xmlid={div07}]

                \teiHead[xmlid={etranges-merveilles-001}]{« Étranges merveilles »}
                \teiP[xmlid={p30-2}]{À l’imprédictibilité de l’histoire de la vie répond l’imprévisibilité des formes vivantes elles-mêmes. Ainsi, la faune des Invertébrés cambriens des schistes de Burgess traduit une radiation adaptative d’animaux aux formes étranges, inconnues dans notre monde, et pour cause : ils sont presque tous éteints. Ces « étranges merveilles » sont intelligibles sur le fond d’une histoire contingente du vivant. }
                \begin{teiFigure}[xmlid={figure03}]

                  \teiHead[xmlid={figure-4-reconstitution-d-anomalocaris-canadensis-etrange-merveille-anomalocaris-appartient-a-un-groupe-dinvertebres-cambriens-eteints-dont-la-structure-na-rien-de-commun-avec-les-animaux-existant-aujourdhui-jun-ni075-creative-commons-attribution-share-alike-4-0-international-001}]{Figure 4. Reconstitution d’\teiHi[rend={italic}]{anomalocaris canadensis}, « étrange merveille ». \teiHi[rend={italic}]{Anomalocaris} appartient à un groupe d’invertébrés cambriens éteints dont la structure n’a rien de commun avec les animaux existant aujourd’hui. © Jun (@ni075), Creative Commons Attribution-Share Alike 4.0 International.}
                
\end{teiFigure}
                \teiP[xmlid={p31-2}]{Adolf Seilacher, quant à lui, s’est particulièrement attaché à étudier les organismes précambriens de la faune d’Ediacara\teiNote[xmlid={ftn61-2},n={29},place={foot}]{\teiP[xmlid={p-245}]{Adolf Seilacher, « Vendozoa: organismic construction in the Pterozoic biosphere », \teiHi[rend={italic}]{Letheia}, n\teiHi[rend={sup}]{o} 22, 1989, p. 229-239.}}, les « vendobiontes », des animaux à corps mous qui ne sont guère connus que par leurs empreintes conservées dans les sédiments grâce à l’action des tapis microbiens rigidifiant leur surface. À partir de l’étude des formes et les reliefs de ces traces, Seilacher propose des hypothèses sur la structure et le mode de vie de ces organismes énigmatiques.}
                \teiP[xmlid={p32-2}]{Certains de ces animaux archaïques à corps mou et de forme aplatie présenteraient un plan d’organisation en ruban, en galettes ou en feuilles pouvant évoquer des « matelas pneumatiques », souvent segmentés selon des motifs fractals. Ils n’auraient pas de tube digestif visible et devraient absorber les nutriments par la surface de leur corps. }
              
\end{teiDiv}
              \begin{teiDiv}[type={section2},xmlid={div08}]

                \teiHead[xmlid={art-fossile-001}]{« Art fossile »}
                \teiP[xmlid={p33-2}]{En 1997, Seilacher présente à l’université de Yale (États-Unis) et de Tübingen (Allemagne) une exposition intitulée \teiHi[rend={italic}]{Fossil Art. }Dans l’introduction au catalogue de l’exposition, Seilacher se démarque de la mystique prônée par Ernst Haeckel, et de ses « formes artistiques de la nature » : « Nous savons aujourd’hui, énonce-t-il, qu’il n’y a pas de mystérieuse \teiHi[rend={italic}]{vis vitalis} responsable de la symétrie des radiolaires et des coquilles de mollusques, mais plutôt qu’une synergie de règles élémentaires préside à la construction des formes\teiNote[xmlid={ftn62-2},n={30},place={foot}]{\teiP[xmlid={p-246}]{Adolf Seilacher, \teiHi[rend={italic}]{Fossil Art, An Exhibition of the Geologisches Institut}, catalogue d’exposition, Tübingen, Universität Tübingen, 1997, p. 7.}}. »}
                \begin{teiFigure}[xmlid={figure04}]

                  \teiHead[xmlid={figure-5-couverture-du-catalogue-de-lexposition-fossil-art-preparee-par-adolf-seilacher-et-presentee-a-tubingen-allemagne-en-1997-collection-paleontologique-de-luniversite-de-tubingen-001}]{Figure 5. Couverture du catalogue de l’exposition \teiHi[rend={italic}]{Fossil Art}, préparée par Adolf Seilacher et présentée à Tübingen (Allemagne) en 1997. © Collection paléontologique de l’université de Tübingen.}
                
\end{teiFigure}
                \teiP[xmlid={p34-2}]{Le titre de cette exposition, explique-t-il, est en rapport avec la conception moderne de l’Art : « ce n’est pas le simple plaisir esthétique, mais “la fascination”, qui attire le public. » Or, « les fossiles sont des objets fascinants à tous égards par leur histoire, mais aussi par l’évidente rencontre de processus biologiques et physiques qui président à leurs formes\teiNote[xmlid={ftn63-2},n={31},place={foot}]{\teiP[xmlid={p-247}]{\teiHi[rend={italic}]{Ibid}.}} ». Ces objets fascinants, ce sont des empreintes : des traces de limaces cambriennes géantes à la recherche de nourriture, ou bien, sur des grès cambriens d’Espagne, des formes en courbes de lasso traduisant une quête primitive de nourriture (qui couvre une large surface sans l’explorer systématiquement) ; ou encore des traces d’un organisme se déplaçant sous terre, mais pourvu d’une sorte de tuba oscillant fonctionnant comme un aspirateur pour inhaler les particules nutritives. « Longtemps, les tentatives pour faire de l’art à partir d’objets fossiles ont utilisé les spécimens eux-mêmes comme objets de beauté\teiNote[xmlid={ftn64-2},n={32},place={foot}]{\teiP[xmlid={p-248}]{Stephen J. Gould, « Foreword », dans Adolf Seilacher, \teiHi[rend={italic}]{op. cit}., p. 8.}} », écrit Gould dans son commentaire introductif au catalogue de cette exposition. Mais les œuvres d’art ici désignées « représentent \teiHi[rend={italic}]{les activités et les comportements des organismes} \lbrack{}nous soulignons\rbrack{} – terriers, traces, pistes et sentiers\teiNote[xmlid={ftn65-3},n={33},place={foot}]{\teiP[xmlid={p-249}]{\teiHi[rend={italic}]{Ibid}. }} ».}
                \teiP[xmlid={p35-2}]{Il s’agit bien ici de comprendre les formes vivantes comme des formes en mouvement, en devenir, en insistant sur la beauté qui naît de leur croissance, de leurs transformations, de leurs adaptations, de leurs mouvements, ces phénomènes dynamiques étant inscrits dans les structures mêmes des êtres.}
                \begin{teiFigure}[xmlid={figure05}]

                  \teiHead[xmlid={figures-6-planche-extraite-du-catalogue-de-lexposition-fossil-art-preparee-par-adolf-seilacher-et-presentee-a-tubingen-allemagne-en-1997-natural-history-museum-in-aarhus-danemark-001}]{Figures 6. Planche extraite du catalogue de l’exposition \teiHi[rend={italic}]{Fossil Art}, préparée par Adolf Seilacher et présentée à Tübingen (Allemagne) en 1997. © Natural History Museum in Aarhus, Danemark. }
                
\end{teiFigure}
                \begin{teiFigure}[xmlid={figure06}]

                  \teiHead[xmlid={figure-7-planche-extraite-du-catalogue-de-lexposition-fossil-art-preparee-par-adolf-seilacher-et-presentee-a-tubingen-allemagne-en-1997-collection-paleontologique-de-luniversite-de-tubingen-001}]{Figure 7. Planche extraite du catalogue de l’exposition \teiHi[rend={italic}]{Fossil Art}, préparée par Adolf Seilacher et présentée à Tübingen (Allemagne) en 1997. © Collection paléontologique de l’université de Tübingen.}
                
\end{teiFigure}
                \teiP[xmlid={p36-2}]{Depuis plus de deux siècles, l’attention aux formes vivantes portée par la tradition de la morphologie a été plusieurs fois revisitée, offrant des clés chaque fois renouvelées pour décrire et comprendre l’unité et la diversité des êtres actuels et fossiles. La beauté est chaque fois invoquée, et redéfinie de manière différente : chez Owen, la beauté dans la nature consiste dans la perfection de l’adaptation, elle est l’effet d’une impulsion divine. Chez Haeckel, la beauté naturelle est exubérance des formes, effet de la force vitale, tandis que chez Gould et Seilacher la beauté est définie comme émergence de formes imprévisibles, résultant de l’ad/ex-aptation, des contraintes de construction et de l’historicité du vivant, et résultant aussi des mouvements de la vie même. }
                \teiP[xmlid={p37-2}]{Ainsi, au nom de la morphologie sont proposées trois visions divergentes de l’organisme, de la structure du vivant, de son évolution et de sa place dans la nature – mais dans chaque cas, la recherche scientifique s’articule à une quête esthétique. La réflexion sur les formes vivantes en mouvement et en transformation (\teiHi[rend={italic}]{bildung}), ouvrant sur une convergence de l’art et des sciences, apparaît comme une des réponses possibles à la question de la beauté vitale.}
              
\end{teiDiv}
            
\end{teiDiv}
          
\end{teiElement}
        
\end{teiElement}
        \begin{teiElement}[name={group},type={chapter},data-page-title={Beautés animales : ordre ou expression ?},data-page-authors={Bertrand Prévost@@Université Bordeaux Montaigne, laboratoire MICA – EA 4426, France},data-include-href={Ch06\_beautes\_animales\_Prevost.xml},xmlid={chapter-005}]

          \begin{teiElement}[name={front}]

            \begin{teiDiv}[type={titlePage},xmlid={titlepage-007}]

              \teiP[rend={title-main},xmlid={p-250}]{Beautés animales : ordre ou expression ?}
              \teiP[rend={author-aut},xmlid={p-251}]{Bertrand \teiHi[rend={small-caps}]{Prévost}}
              \teiP[rend={authority\_affiliation},xmlid={p-252}]{Université Bordeaux Montaigne, laboratoire MICA – EA 4426, France}
            
\end{teiDiv}
          
\end{teiElement}
          \begin{teiElement}[name={body},xmlid={body-007}]

            \teiP[xmlid={p1-7}]{Le registre de la vie, pas moins que celui de la beauté, a ses ambiguïtés ; et associer les deux revient peut-être à une double complication. Le vivant, en effet, du moins si on le prend dans sa réalité physiologique et organique, serait plutôt aux antipodes de la beauté : que produit donc chez la plupart d’entre nous la vision des entrailles, des fluides corporels, des organes, du sang sinon une sorte de dégoût voire d’horreur ; pour ne rien dire des sons et des odeurs ? }
            \teiP[xmlid={p2-7}]{Que viennent faire les bêtes dans cette affaire ? Elles semblent accentuer davantage encore, d’une part, ce contraste entre le vivant organique et sa destinée informe (odeur, bruits, sons), et de l’autre la vie génératrice de formes toutes tendues vers une souveraine beauté. C’est bien sûr de ce côté-là qu’on range ces sublimes beautés animales, qui participeraient (conjointement avec les Vénus sculptées) des débuts de l’art figuratif. }
            \teiP[xmlid={p3-7}]{Mais précisément, dès que l’on tente d’arrimer beauté, esthétique ou art aux animaux, tout se passe comme si cette même beauté avait pour fonction de conjurer le \teiHi[rend={italic}]{vivant} animal : on va le voir, tant dans ses productions qu’en lui-même. Tout se passe comme si la beauté animale ne pouvait exister qu’en contrevenant à l’informe organique, à sa mollesse et à sa morphologie tout en sinuosités, en courbes, sans orthogonalité… Bref, la beauté animale se confondrait avec l’imposition d’un \teiHi[rend={italic}]{ordre. }Comment en effet ne pas succomber devant la beauté des structures régulières qui organisent la toile d’araignée, la rosace dessinée dans le sable par un poisson-globe du Japon, la forme de nombreux nids, des termitières, des alvéoles des rayons de miel, etc. ? Toutes ces régularités symptomatisent en réalité le développement d’un \teiHi[rend={italic}]{plan}. L’intelligence, l’ingéniosité ou l’habileté prêtée à l’animal industrieux est toujours idéalement celle d’un architecte qui planifie un \teiHi[rend={italic}]{projet.} Et la poïétique animale s’entendra ici comme un pouvoir de projection par lequel l’animal, à l’instar de l’homme, est capable de se poser en sujet devant une objectivité formelle et matérielle. }
            \teiP[xmlid={p4-7}]{Mais ce qui vaut pour des productions vaut autant pour la forme même de l’animal, et révèlera cette fois non plus un plan subjectif, mais un \teiHi[rend={italic}]{plan de la nature} : de D’Arcy Thompson à Roger Caillois, nombreux sont ceux à s’être émerveillés devant la perfection spiralée des coquilles de mollusques, devant la précision des enroulements des cornes de certaines espèces de caprinés, et à avoir cherché, il faudrait presque dire, leur géométrie secrète. Pour ne rien dire de toutes les organisations symétriques qui donneraient au vivant quelque chose d’un ordre minimal ! À quoi avons-nous ici affaire, sinon à la convocation d’un concept de beauté conçue comme ordre, symétrie, ou plutôt « symmétrie » (avec deux M pour traduire le grec \teiHi[rend={italic}]{symmetria}), eurythmie, harmonie ? Bref, un concept de beauté somme toute fort classique, qui remonterait au moins aux Grecs, mais sans doute plus loin encore, et qui se confondrait avec l’art lui-même, si tant est qu’une même racine indo-européenne lie art et ordre\teiNote[xmlid={ftn1-6},n={1},place={foot}]{\teiP[xmlid={p-253}]{Voir Émile Benvéniste, \teiHi[rend={italic}]{Le vocabulaire des Institutions européennes}, Paris, Éditions de Minuit, 1969, t. II, p. 100 et 101.}}, toutes choses que de très nombreuses théories « modernes » de l’ornement et de la beauté cosmétique auront repris à leur compte (Gottfried Semper, Charles Blanc ou même Ernst Gombrich).}
            \begin{teiDiv}[type={section1},xmlid={div01-4}]

              \teiHead[xmlid={paraitre-est-une-fonction-vitale-001}]{« Paraître est une fonction vitale »}
              \teiP[xmlid={p5-7}]{On doit à un zoologue suisse du \teiHi[rend={small-caps}]{xx}\teiHi[rend={sup}]{e} siècle d’avoir pris la question des beautés animales d’un point de vue radicalement différent : il s’agit d’Adolf Portmann. Sous la notion d’« apparences authentiques », il entendait faire droit à un spectre de phénomènes beaucoup plus étendu que les seules régularités formelles, qu’une seule beauté pensée comme ordre : tout le monde des couleurs et leurs qualités singulières, l’infinie complexité des motifs, la richesse ornementale de tous les appendices, crêtes, queue, etc. Bref, ce ne sont pas tant les formes animales qui intéressaient Portmann que quelque chose de plus profond qui les transfigure en \teiHi[rend={italic}]{formes expressives}, en apparences. On peut même dire que tout le travail de biologie morphologique du zoologue aura consisté à rejouer l’opposition d’un plan d’organisation, par lequel les espèces sont des composés organisés en formes et fonctions, avec un plan expressif par quoi les espèces apparaissent et se donnent à voir en deçà de toute fonctionnalité. }
              \teiP[xmlid={p6-7}]{Cela supposait bien entendu une critique radicale du schème utilitariste ou fonctionnaliste darwinien : toute forme sert à quelque chose, qu’il s’agisse de la sélection naturelle ou de la sélection sexuelle, même si le sens de cette fonction peut provisoirement nous échapper ; mais cela supposait tout autant une critique du schème morphogénétique et des mécanismes physico-chimiques (voire géométrico-mathématiques) qui président à son développement : la formation des couleurs, par exemple, peut bien s’expliquer par le double jeu de couleurs pigmentaires (les rouges, jaunes ou noirs par exemple) et de couleurs structurales ou optiques (le bleu par exemple), il s’agit toujours de déterminer une extériorité causale, une explication extérieure des formes.}
              \teiP[xmlid={p7-7}]{Pour accéder à une dimension expressive, qui aille au-delà de ces causalités extérieures et qui embrasse une dynamique non plus explicative, mais implicative, non plus de développement, mais d’enveloppement, il aura fallu que Portmann en passe par la proposition d’une notion : l’autoprésentation (\teiHi[rend={italic}]{Selbstdarstellung}) : « L’autoprésentation est la manifestation d’un soi (\teiHi[rend={italic}]{die Manifestation eines Selbst}), dont l’essence nous reste toujours cachée et dont le fondement matériel dans le protoplasme nous est pareillement inaccessible pour le moment\teiNote[xmlid={ftn16},n={2},place={foot}]{\teiP[xmlid={p-254}]{Adolf Portmann, \teiHi[rend={italic}]{Aufbruch der Lebensforschung}, Franckfort-sur-le-Main, Suhrkamp Verlag, 1965, p. 213. Voir encore : « L’organisme a aussi à apparaître, \lbrack{}…\rbrack{} il doit se présenter dans sa spécificité », dans « L’autoprésentation, motif de l’élaboration des formes vivantes » \lbrack{}« Selbstdarstellung als Motiv der lebendigen Formbildung », \teiHi[rend={italic}]{Geist und Werk. Aus der Werkstatt unserer Autoren. Zum 75, Geburtstag von Dr. Daniel Brody}, Zurich, Rhein Verlag, 1958\rbrack{}, Jacques Dewitte (trad.), \teiHi[rend={italic}]{Études phénoménologiques}, vol. 12, n\teiHi[rend={sup}]{o} 23-24, 1996, p. 131-164, ici p. 150.}}. » S’ouvre maintenant une autre dimension, qui est portée par le préfixe \teiHi[rend={italic}]{selbst}- :}
              \begin{teiCit}[xmlid={cit01-5}]

                \begin{teiQuote}
Tous les êtres doués de relation au monde possèdent aussi le caractère de l’autoprésentation, qui a été le plus souvent méconnu jusqu’ici. Les parties nécessaires à la relation au monde sont à chaque fois façonnées selon une particularité typique du groupe, une singularité qui se manifeste dans de nombreuses structures et de nombreux modes de comportement dont la spécificité ne peut être expliquée uniquement à partir des fonctions de la simple conservation de l’individu et de l’espèce\teiNote[xmlid={ftn17},n={3},place={foot}]{\teiP[xmlid={p-255}]{\teiHi[rend={italic}]{Ibid}., p. 158.}}.
\end{teiQuote}
              
\end{teiCit}
              \teiP[xmlid={p8-7}]{Ce serait aller bien vite en besogne que d’imaginer que cette ipséité, ce soi, \teiHi[rend={italic}]{selbst}-, concerne une subjectivité animale qui, après coup, aurait à apparaître, comme si l’animal se tenait derrière le masque de son image. Ce serait précisément en revenir à l’idée d’une secondarité de l’apparence expressive au regard de l’existence subjective. Il faut au contraire affirmer que ce « soi » est celui de l’apparence elle-même. Ce soi apparaissant vise d’abord le soi d’une image, et il constitue sans aucun doute la première pierre vers l’idée d’une image \teiHi[rend={italic}]{en-soi}, que viendra couronner plus tard le concept d’apparence inadressée. Exactement comme avec la volonté de puissance nietzschéenne, où c’est la puissance qui veut, et non la volonté, il faut dire : c’est l’apparence qui se présente, ce qui n’implique pas du tout que cette apparence soit indéterminée. Bien au contraire : Portmann a toujours insisté sur le haut degré de différenciation de la présentation, et il est à cet égard tout à fait significatif qu’il ait conçu l’autoprésentation dans les termes d’une \teiHi[rend={italic}]{héraldique naturelle }:}
              \begin{teiCit}[xmlid={cit02-5}]

                \begin{teiQuote}
Les motifs signalétiques présentent une caractéristique importante : ils ne peuvent se confondre avec d’autres, un peu comme les bannières et les blasons d’autrefois. D’ailleurs la science héraldique à son apogée offre bien des traits communs avec les règles auxquelles ces formations signalétiques sont assujetties. \lbrack{}…\rbrack{} Ces ensembles signalétiques nous font songer à des drapeaux, surtout les « coiffures » de certains vertébrés, un vol de canards, l’aspect de certains pluviers limicoles ou la livrée des poissons familiers des récifs de corail\teiNote[xmlid={ftn18},n={4},place={foot}]{\teiP[xmlid={p-256}]{Adolf Portmann, \teiHi[rend={italic}]{La forme animale} \lbrack{}\teiHi[rend={italic}]{Die Tiergestalt}, Bâle, F. Reinhardt, 1960, 2\teiHi[rend={sup}]{e} éd.\rbrack{}, Georges Rémy (trad.), Paris, Payot, 1962, p. 114.}}.
\end{teiQuote}
              
\end{teiCit}
              \teiP[xmlid={p9-6}]{Un drapeau, un blason, une bannière ne sont pas autre chose que la mise en forme expressive, qu’une formule plastique de l’expression d’une individualité, ce qui ne signifie pas que cette individualité renvoie à une identité personnelle. Là encore, il faut se montrer prudent ; que gagne-t-on en effet à introduire l’idée d’une différence ou d’une distinction individuelle ? Pas grand-chose, s’il s’agit de penser l’animal comme individu apparaissant. Il nous resterait des individus subjectifs, des personnes animales, auxquelles parfois, dans un zoo voire dans la nature, on va jusqu’à donner un nom. Mais de quelle individualité parlons-nous ? Tout l’enjeu de l’autoprésentation est justement de mettre l’individualité dans l’apparence, quitte à ce que cette individualité vienne contredire l’existence individuelle de l’animal, selon des modalités qu’il faudra préciser. }
              \teiP[xmlid={p10-6}]{Car en vérité, en posant que « paraître est une fonction vitale\teiNote[xmlid={ftn19},n={5},place={foot}]{\teiP[xmlid={p-257}]{Adolf Portmann, préface à Theo Jahn, \teiHi[rend={italic}]{La vie et ses formes}, Paris, Bordas, 1968, p. 15.}} », Portmann non seulement ancrait les apparences dans la vie même, mais en les appréhendant comme expressives, il les sortait ou les extrayait d’une vie étroitement organique au profit d’une vie supérieure, d’une vitalité. De fait, c’est clairement le principe de conservation dans ses rapports à la physiologie – comme système organique – qui se trouvait critiqué : }
              \begin{teiCit}[xmlid={cit03-5}]

                \begin{teiQuote}
Le métabolisme sert la conservation de l’individu, mais, quelle que soit l’importance qu’on lui accorde, nous devons nous faire à l’idée que l’organisme n’est pas là pour faire fonctionner le métabolisme, mais que ce dernier est pratiqué pour que cette forme de vie particulière (\teiHi[rend={italic}]{diese besondere Lebensform}) se réalise dans les individus. L’organisme fait fonctionner le métabolisme pour que son mode d’être spécifique individuel \teiHi[rend={italic}]{(seine spezifische Seinsweise in Einzelwesen}) puisse s’affirmer un temps dans l’apparence\teiNote[xmlid={ftn20},n={6},place={foot}]{\teiP[xmlid={p-258}]{Adolf Portmann,\teiHi[rend={italic}]{ Neue Wege der Biologie}, Munich, Piper, 1961, p. 217\teiHi[rend={italic}]{.}}}.
\end{teiQuote}
              
\end{teiCit}
              \teiP[xmlid={p11-6}]{D’une certaine façon, Portmann retrouvait là une forme d’inspiration nietzschéenne extrêmement méfiante face à un darwinisme étroit, pensant la vie comme adaptation et « lutte pour la survie » (\teiHi[rend={italic}]{struggle for life}) : « Les fonctions de conservation, ajoutait le zoologue, sont au service de quelque chose qui est à conserver et qui est davantage que la simple somme des opérations de conservation.\teiNote[xmlid={ftn21},n={7},place={foot}]{\teiP[xmlid={p-259}]{Adolf Portmann, « L’autoprésentation, motif de l’élaboration des formes vivantes », art. cité, p. 158.}} »}
            
\end{teiDiv}
            \begin{teiDiv}[type={section1},xmlid={div02-4}]

              \teiHead[xmlid={le-plan-expressif-de-lapparaitre-001}]{Le plan expressif de l’apparaître}
              \teiP[xmlid={p12-6}]{Dans toute son œuvre, Portmann n’aura de cesse de souligner l’autonomie de ce plan expressif des apparences, son indépendance tant au regard du plan de l’individualité physiologique (de ses fonctions et de ses formes) que de l’organicité elle-même. Or, c’est justement ici, dans cette indépendance, ou mieux, dans ce décollement d’un plan expressif que se niche toute la différence entre la beauté d’une forme et l’expressivité d’une apparence.}
              \teiP[xmlid={p13-5}]{D’abord, au niveau ontogénétique, celui du développement de l’individu : penser le détachement d’un plan expressif, ce sera éprouver un plan ni organique, ni inorganique, mais bien plutôt \teiHi[rend={italic}]{anorganique}. Les ocelles du jaguar, par exemple, ou encore les couleurs chatoyantes des perroquets ne sont pas en effet tout à fait inorganiques, puisque chimiquement, ce sont toujours des cellules vivantes qui les composent, puisque ce sont toujours des processus morphogénétiques qui président à leur formation ; mais elles ne sont pas pour autant tout à fait organiques puisqu’elles ne se décalquent pas sur l’anatomie ou les divisions de l’organisme, et surtout, puisqu’elles demeurent le plus souvent a-fonctionnelles au regard de la conservation de l’espèce. }
              \teiP[xmlid={p14-5}]{C’est donc bien quelque chose comme un plan qui se dessine, un plan cosmétique. On a dit à quel point Portmann a toujours insisté sur « \teiHi[rend={italic}]{l’autonomie} relative » de ce plan. Or, c’est bien l’indépendance totale de l’apparence au regard de toute organisation métabolique, de tout organisme, qui autorise à parler d’élégance. Marbrures, zébrures, taches, ocelles… : ces formes se détachent de toute structuration organique, et les biologistes ont depuis longtemps abandonné l’idée de leur trouver une base anatomique. Ce qui importe ici tient fondamentalement à cette façon de faire se décoller un plan esthétique ou expressif dans la mesure où l’apparence anorganique n’est plus attachée à un sujet vivant – exactement comme une parure ou un vêtement est amovible – mais sans pour autant constituer un corps objectif au-delà ou par-delà la surface subjective – toujours comme un vêtement ou une parure n’est amovible que dans la mesure où il consiste en un objet. À la lettre, il faudrait dire que les animaux \teiHi[rend={italic}]{gèrent leur apparence, }si tant est que l’apparence fasse toujours l’objet d’une « gestion », au sens étymologique latin du port : façon de signifier le détachement ou le décollement de cette apparence de son substrat corporel ou organique. Pour autant, les animaux ne revêtent pas l’apparence comme on revêt un vêtement ; ou plutôt, c’est notre conception du vêtement comme corps, chose ou objet qu’il faudrait revoir. C’est d’ailleurs bien le cadeau extraordinaire que nous fait l’idée d’une cosmétique animale, par différence avec une cosmétique humaine. Car revêtir une apparence n’est pas porter un objet. Et l’on sait tout le poids symbolique que l’on a pu donner, en histoire, en sociologie comme en anthropologie, au vêtement et plus généralement aux atours corporels, au titre d’un second corps « social », « culturel », etc. Mais si l’apparence est cosmétique, c’est en tant qu’elle fait se lever \teiHi[rend={italic}]{une} parure, tout à fait singulière, bien que cette singularité ne se confonde ni avec une subjectivité organique ni avec l’objectivité d’un corps autre.}
              \teiP[xmlid={p15-5}]{C’est une façon de dire que l’extériorité de l’apparaître contraint de penser une expressivité irréductible au jeu de la subjectivité et de l’objectivité. La condition est en effet double : il faut que la forme se détache ou se décolle du substrat organique (subjectif), mais il faut dans le même temps qu’elle ne se substantialise pas dans un vêtement (objectif). Portmann allait même jusqu’à entrevoir la possibilité théorique d’un renversement total entre le corps et l’apparence :}
              \begin{teiCit}[xmlid={cit04-5}]

                \begin{teiQuote}
Dans cette perspective, d’innombrables caractères négligés du vivant acquièrent une dignité nouvelle, alors qu’ils étaient considérés jusque-là comme accessoires. Et s’ils étaient l’essentiel ? Si les êtres vivants n’étaient pas là afin que soit pratiqué le métabolisme, mais pratiquaient le métabolisme afin que la particularité qui se réalise dans le rapport au monde et l’autoprésentation ait pendant un certain temps une durée dans le monde\teiNote[xmlid={ftn22},n={8},place={foot}]{\teiP[xmlid={p-260}]{Adolf Portmann, « L’autoprésentation, motif de l’élaboration des formes vivantes », art. cité, p. 157 (nous soulignons). }} ? 
\end{teiQuote}
              
\end{teiCit}
              \teiP[xmlid={p16-5}]{Que l’apparence ait besoin de durer « un certain temps dans le monde » signifie bien qu’elle ne peut se confondre avec des formes corporelles actuelles, avec une individualité organique qui va lui donner son ici et son maintenant extensif et qualitatif. « Je veux un corps ! » : tout se passe en effet comme si l’apparence expressive réclamait un corps, un substrat organique pour exister actuellement. Ce qui ne signifie pas qu’elle ne soit pas réelle ; elle a au contraire la réalité d’une virtualité intensive qui se distingue en droit de toute forme d’organisation : image singulière « idéale sans être abstraite, réelle sans être actuelle ». Autant dire que cette singularité (Portmann parle bien d’une \teiHi[rend={italic}]{particularité} qui se réalise) possède une étrange individualité, qu’il faudrait davantage nommer avec Gilbert Simondon, « pré-individualité », pour dire le potentiel vital qui l’anime et qui précède toute forme de vie organique individuelle. Il ne faudrait donc pas s’imaginer des corps nus revêtant tel ou tel habit : ce serait encore confronter deux individus, alors que ce qui fait la vitalité de l’apparence tient davantage à une puissance d’individuation, préorganique, qui vient soumettre l’organique à son insistance, qui vient en faire l’instrument de son actualisation comme passage à l’individu.}
              \teiP[xmlid={p17-5}]{Cependant, ce détachement expressif de l’apparence de l’individualité ne s’entend pas seulement au seul niveau des individus vivants ; il faut encore l’entendre à un niveau phylogénétique, au niveau du développement de l’espèce tout entière ; c’est ainsi que va se lever une vie des apparences qui se distinguera en droit de la vie des individus. En effet, cette antériorité d’un plan vital expressif sur un plan organique métabolique, \teiHi[rend={italic}]{quitte précisément à perdre l’individu organique}, est encore confirmée au niveau de la macroévolution par ce que l’on appelait jadis les « phénomènes d’hypertélie », à savoir le développement extrême d’ornements mettant en péril la propre survie de l’espèce par le handicap qu’il génère\teiNote[xmlid={ftn23},n={9},place={foot}]{\teiP[xmlid={p-261}]{Cet argument est au cœur de la théorie du handicap développée par Amotz et Avishag Zahavi : voir \teiHi[rend={italic}]{The Handicap Principle. A Missing Piece of Darwin’s Puzzle}, New York-Oxford, Oxford University Press, 1997. }}. Le paon fait souvent les frais d’une étrange contradiction entre la sélection sexuelle et la sélection naturelle. En effet, les femelles seraient portées à préférer les mâles arborant les plumes les plus longues et les plus vivement colorées, trait qui en croissant par l’hérédité va mettre en péril le devenir même de l’espèce, puisque ces mêmes mâles seront de plus en plus handicapés par leur longue traîne et particulièrement exposés aux prédateurs par leurs vives couleurs\teiNote[xmlid={ftn24},n={10},place={foot}]{\teiP[xmlid={p-262}]{Voir cependant les réserves de Roslyn Dakin et Robert Montgomerie, « Peahens prefer peacocks displaying more eyespots, but rarely », \teiHi[rend={italic}]{Animal Behaviour}, vol. 82, n\teiHi[rend={sup}]{o} 1, juillet 2011, p. 21-28. }}. Il en irait de même avec l’élan d’Irlande (le \teiHi[rend={italic}]{megaleceros giganteus}) dont l’accroissement démesuré des bois aurait entraîné sa disparition il y a onze mille ans. Stephen Jay Gould a certes montré qu’il n’en était rien, et que la taille de ses bois était tout à fait adéquate à sa masse et son volume corporels\teiNote[xmlid={ftn25},n={11},place={foot}]{\teiP[xmlid={p-263}]{Voir Stephen Jay Gould, « L’élan d’Irlande mal nommé, mal traité, mal compris », dans \teiHi[rend={italic}]{Darwin et les grandes énigmes de la vie}, Paris, Éditions du Seuil, 1997, p. 81-93.}}. Mais d’autres cas semblent résister, comme les gigantesques plumes caudales du faisan argus. Il est fascinant de voir ici l’éthologue constater « un aspect vraiment étrange et, toute réflexion faite, inquiétant de la phylogenèse\teiNote[xmlid={ftn26},n={12},place={foot}]{\teiP[xmlid={p-264}]{Konrad Lorenz, \teiHi[rend={italic}]{L’agression. Une histoire naturelle du mal}, Vilma Fritsch (trad.), Paris, Flammarion, 1969, p. 47.}} ». Pire, quand on fait de cette superbe livrée l’un des « produit\lbrack{}s\rbrack{} les plus stupide\lbrack{}s\rbrack{} de la sélection uniquement intra-espèce\teiNote[xmlid={ftn27},n={13},place={foot}]{\teiP[xmlid={p-265}]{\teiHi[rend={italic}]{Ibid}.}} », on comprend vite que l’histoire naturelle de la morale qu’on s’était donnée comme horizon s’est transmuée en une moralisation de l’histoire naturelle beaucoup plus dommageable\teiNote[xmlid={ftn28},n={14},place={foot}]{\teiP[xmlid={p-266}]{Allusion au sous-titre de l’ouvrage de Konrad Lorenz : \teiHi[rend={italic}]{L’agression. Une histoire naturelle du mal}, \teiHi[rend={italic}]{op. cit}. Pour une juste critique de ces positions, voir notamment Vinciane Despret, L\teiHi[rend={italic}]{a danse du cratérope écaillé. Naissance d’une théorie éthologique}, Paris, La Découverte, 2021, p. 70-73.}}. Car ce qui se rend manifeste dans les phénomènes dits d’hypertélie, c’est bien que le plan expressif de l’apparence non seulement ne se confond pas avec un développement organique individuel, mais encore qu’il peut le contredire jusque dans la mort. Tout se passe en effet comme si les apparences voulaient se maintenir en vie coûte que coûte dans leurs déterminations les plus expressives ; tout se passe comme si la vitalité des apparences cherchait à l’emporter sur la vie organique, les individus qui les portent dussent-ils disparaître sous le poids d’une contrainte trop forte pour l’organisme, pour son métabolisme, pour des conditions de vie actuelle. }
              \teiP[xmlid={p18-5}]{Finalement, il se peut que ce qui se dessine ici, dans cette réhabilitation de l’expressivité des formes, revienne à faire revivre une expérience de l’image fort ancienne en vérité et somme toute presque banale, une expérience qui pense l’image comme extériorité, réalité détachée des choses, arrachée aux corps, mais qui ne saurait se loger dans une représentation subjective. Et il paraît évidemment que le tournant de la modernité n’aura pas peu fait pour intérioriser l’expressivité. }
              \teiP[xmlid={p19-5}]{Les apparences animales sont très précisément des images au sens où Lucrèce développait, dans un passage fameux du livre IV du \teiHi[rend={italic}]{De rerum natura}, sa théorie des simulacres : une théorie qui pensait l’image comme toujours soutirée aux corps, comme un arrachement, une soustraction des corps. Or, c’est justement le monde des animaux qui donnait au poète-philosophe latin la matière expressive de sa pensée, plus précisément sous l’espèce de la \teiHi[rend={italic}]{dépouille animale}. Très schématiquement, rappelons que les simulacres naissent par le détachement de l’enveloppe superficielle des corps qui vient heurter nos sens :}
              \begin{teiCit}[xmlid={cit05-4}]

                \begin{teiQuote}
Il existe des images \lbrack{}\teiHi[rend={italic}]{simulacra}\rbrack{}, comme nous les nommons ; sorte de membranes détachées de la surface des corps, elles voltigent de tous côtés à travers les airs. \lbrack{}…\rbrack{} Je dis que les choses envoient de leur surface des effigies \lbrack{}\teiHi[rend={italic}]{effigias}\rbrack{}, formes ténues \lbrack{}\teiHi[rend={italic}]{tenuisque figuras}\rbrack{} d’elles-mêmes, des membranes en quelque sorte ou des écorces, puisque l’image revêt l’aspect, la forme exacte de n’importe quel corps dont, vagabonde, elle émane. \lbrack{}…\rbrack{} Mainte chose visible émet des corpuscules. Certains se dissipent et s’évaporent, la fumée du bois vert ou la chaleur du feu, d’autres sont plus serrées, plus denses, telles en été les tuniques rondes que déposent les cigales, les membranes que les veaux dès leur naissance abandonnent, la robe dont le serpent furtif se dévêt dans les ronces ; nous voyons souvent sa dépouille \lbrack{}\teiHi[rend={italic}]{spoliis}\rbrack{} flottante garnir les broussailles\teiNote[xmlid={ftn29},n={15},place={foot}]{\teiP[xmlid={p-267}]{Lucrèce, \teiHi[rend={italic}]{De la nature – De rerum natura}, Livre IV, l. 34-62, José Kany-Turpin (trad.), Paris, Garnier Flammarion, 1998, p. 245-247.}}.
\end{teiQuote}
              
\end{teiCit}
              \teiP[xmlid={p20-5}]{Une telle appréhension de l’image ne se contente pas de célébrer son extériorité ; elle en démontre la genèse expressive : un arrachement superficiel, un décollement subtil, qui trouve sa forme poétique dans les figures d’une peau délaissée, d’une surface abandonnée.}
            
\end{teiDiv}
          
\end{teiElement}
        
\end{teiElement}
        \begin{teiElement}[name={group},type={chapter},data-page-title={L’herbe et la bête},data-page-subtitle={Philosophie de la « mauvaise » herbe à la bestiole},data-page-authors={Claudine Sagaert@@Université de Toulon, laboratoire Babel, France},data-include-href={Ch07\_lherbe\_labete\_Sagaert.xml},xmlid={chapter-006}]

          \begin{teiElement}[name={front}]

            \begin{teiDiv}[type={titlePage},xmlid={titlepage-008}]

              \teiP[rend={title-main},xmlid={p-268}]{L’herbe et la bête }
              \teiP[rend={title-sub},xmlid={p-269}]{Philosophie de la « mauvaise » herbe à la bestiole}
              \teiP[rend={author-aut},xmlid={p-270}]{Claudine \teiHi[rend={small-caps}]{Sagaert}}
              \teiP[rend={authority\_affiliation},xmlid={p-271}]{Université de Toulon, laboratoire Babel, France}
            
\end{teiDiv}
            \begin{teiElement}[name={epigraph}]

              \begin{teiCit}

                \begin{teiQuote}
J’aime l’araignée et j’aime l’ortie,\teiLb{}Parce qu’on les hait ;\teiLb{}Et que rien n’exauce et que tout châtie\teiLb{}Leur morne souhait ;\teiLb{}\lbrack{}…\rbrack{}\teiLb{}Passants, faites grâce à la plante obscure,\teiLb{}Au pauvre animal.\teiLb{}Plaignez la laideur, plaignez la piqûre,\teiLb{}\teiLb{}Oh ! Plaignez le mal !\teiLb{}Il n’est rien qui n’ait sa mélancolie ;\teiLb{}Tout veut un baiser.\teiLb{}Dans leur fauve horreur, pour peu qu’on oublie\teiLb{}De les écraser,\teiLb{}\teiLb{}Pour peu qu’on leur jette un œil moins superbe,\teiLb{}Tout bas, loin du jour,\teiLb{}La vilaine bête et la mauvaise herbe\teiLb{}Murmurent : Amour !
\end{teiQuote}
                \begin{teiBibl}[xmlid={bibl001-2}]
Victor Hugo, « J’aime l’araignée », \teiHi[rend={italic}]{Les Contemplations}, Livre III, XXVII, 1856.
\end{teiBibl}
              
\end{teiCit}
            
\end{teiElement}
          
\end{teiElement}
          \begin{teiElement}[name={body},xmlid={body-008}]

            \teiP[xmlid={p1-8}]{L’artiste plasticien américain Tony Matelli dialogue avec des pissenlits et autres graminées. Dans sa série d’œuvres intitulée \teiHi[rend={italic}]{Mauvaises herbes}\teiNote[xmlid={ftn1-7},n={1},place={foot}]{\teiP[xmlid={p-272}]{Tony Matelli, \teiHi[rend={italic}]{Weed }504, 517, 518, 519 et 524, galerie Andréhn-Schiptjenko, Paris, bronzes peints, juin 2020.}}\teiHi[rend={italic}]{, }il transfigure dans un style hyperréaliste celles que l’on a coutume de juger inesthétiques, envahissantes et nuisibles. Dans cette galerie d’art parisienne, les mauvaises herbes s’imposent. Infiltrées entre sol et mur, ces « sculptures de brindilles moulées, puis réalisées en bronze et minutieusement peintes à la main\teiNote[xmlid={ftn73-3},n={2},place={foot}]{\teiP[xmlid={p-273}]{Joy des Horts, « Tony Matelli, éloge de la fragilité »,\teiHi[rend={italic}]{ Les hardis}, 5 juin 2020, \teiRef[target={https://www.leshardis.com/2020/06/tony-matelli/}]{\teiHi[rend={underline}]{https://www.leshardis.com/2020/06/tony-matelli/}} (consulté le 25 juin 2025).}} » suscitent l’étonnement du spectateur. Elles captent les regards. }
            \teiP[xmlid={p2-8}]{Dans un autre contexte, l’écrivain cubain Antón Arrufat, dans la nouvelle intitulée \teiHi[rend={italic}]{Éloge du cocuyo}\teiNote[xmlid={ftn74-3},n={3},place={foot}]{\teiP[xmlid={p-274}]{Antȯn Arrufat, \teiHi[rend={italic}]{Fracture et autres histoires/Fractura y otras historias}, Marc Sagaert (éd. et trad.), Marseille, L’atinoir, 2015.}}, nous parle de la danse nocturne de ce type de cafards nommé pyrophore. Scintillante, la beauté de sa parade nuptiale « embellit la noirceur de la nuit\teiNote[xmlid={ftn75-3},n={4},place={foot}]{\teiP[xmlid={p-275}]{\teiHi[rend={italic}]{Ibid.,} p. 37.}} » et, tel un « flambeau des bois », sa lumière « sème d’émeraude l’horizon\teiNote[xmlid={ftn76-3},n={5},place={foot}]{\teiP[xmlid={p-276}]{\teiHi[rend={italic}]{Ibid.}}} ». Jadis, l’exceptionnel éclat dégagé par les pyrophores en faisait d’étonnantes parures féminines. Lors des fêtes nocturnes, les Cubaines les portaient en chapelet autour du cou. « La lumière huilée du cocuyo offrait à ces femmes un charme fantastique\teiNote[xmlid={ftn77-4},n={6},place={foot}]{\teiP[xmlid={p-277}]{\teiHi[rend={italic}]{Ibid.}, p. 38.}}. »}
            \teiP[xmlid={p3-8}]{Le travail de Tony Matelli et le texte d’Antón Arufat interpellent. Ils invitent à un autre regard sur ces espèces généralement résumées à de « moindres-êtres\teiNote[xmlid={ftn78-4},n={7},place={foot}]{\teiP[xmlid={p-278}]{Voir François Dagognet, \teiHi[rend={italic}]{Des détritus, des déchets, de l’abject}, Paris, Les Empêcheurs de penser en rond, 1998.}} ». De la « mauvaise » herbe aux plantes adventices, de l’insecte répugnant au nuisible grouillant, de la bestiole à la vermine, dans l’esprit de nombre d’individus, ces vivants souillent, détruisent et contaminent. La fabrication de leur laideur est indissociable de ce qu’on leur impute. Considérés comme des nuisibles, on n’hésite pas à les arracher ou à les écraser quand on n’essaie pas de les anéantir. Élevés au rang d’ennemis, on utilise même des dérivés d’arme chimique pour les combattre\teiNote[xmlid={ftn79-4},n={8},place={foot}]{\teiP[xmlid={p-279}]{Fritz Haber, chimiste, juif allemand, synthétise le gaz de combat à base de chlore dit « gaz moutarde » et participera par ses recherches à la mise au point du zyklon B. Après la guerre de 1914-1918, le « gaz moutarde » sera transformé en DDT, et après la guerre du Vietnam, l’agent orange de Monsanto sera recyclé en agriculture comme herbicide. Insecticides et herbicides sont issus de l’industrie militaire. Voir Fabrice Nicolino et François Veillerette, \teiHi[rend={italic}]{Pesticides, révélations sur un scandale français}, Paris, Fayard, 2007. }}. Certes, depuis quelques décennies, des voix s’élèvent contre les atteintes portées aux vivants, mais trop rares sont encore celles qui défendent la vie des dites « bestioles » et mauvaises herbes. Considérant ainsi que le mépris et la dévalorisation de ce type d’espèces sont dus à une vision réductrice et utilitariste de l’homme, en déconstruire les présupposés ne permettrait-il pas d’augurer d’autres types d’alliances entre les espèces et d’en révéler la beauté vitale ? }
            \teiP[xmlid={p4-8}]{En traçant une cartographie des qualités négatives attribuées aux mauvaises herbes comme aux petites bêtes, il apparaît que cette conception n’est pas sans lien avec une impuissance de l’homme à les contrôler et à les maîtriser. Les connotations peu élogieuses à partir desquelles elles ont été pensées, nommées, décrites, ont contribué à en fabriquer la laideur et par là même à les considérer comme inutiles et nuisibles.}
            \begin{teiDiv}[type={section1},xmlid={div01-5}]

              \teiHead[xmlid={etranges-etrangeres-001}]{Étranges étrangères}
              \teiP[xmlid={p5-8}]{Les « mauvaises herbes », appelées adventices et rudérales, sont associées à des étrangères ou des importunes. Adventices\teiNote[xmlid={ftn80-3},n={9},place={foot}]{\teiP[xmlid={p-280}]{Adventices, selon le latin \teiHi[rend={italic}]{adventus}, indique ce qui vient du dehors.}}, elles s’introduisent sur n’importe quel territoire. Rudérales\teiNote[xmlid={ftn81-3},n={10},place={foot}]{\teiP[xmlid={p-281}]{Rudérale vient du latin \teiHi[rend={italic}]{rudus}, décombre.}}, elles poussent dans les décombres et s’étendent sur des espaces devenus inoccupés. Elles font preuve d’une vitalité extrême, d’un dynamisme et d’une puissance de vie propres. On les retrouve dans les villes, dans les jardins, dans les champs, sur les chemins, sur les terrains de jeux, entre deux pavés, dans les pots de fleurs comme dans les fissures des revêtements ou au beau milieu d’espaces insalubres. Nomades, elles s’invitent là où elles ne sont pas conviées. Dotées également d’une grande plasticité écologique, elles résistent à la sécheresse comme à l’excès d’humidité. Du fait de leur pouvoir germinatif élevé, elles infestent les cultures et sont étiquetées comme étant des nuisibles. De Pline l’Ancien\teiNote[xmlid={ftn82-3},n={11},place={foot}]{\teiP[xmlid={p-282}]{Voir Pline L’Ancien, \teiHi[rend={italic}]{Histoire naturelles}, livre XVIII, chapitre XLIV. Il parle de la rapidité avec laquelle l’ivraie se répand dans un champ de blé et le contamine.}} au discours scientifique de 1992 à Rio de Janeiro, des propos similaires les qualifient. « Pestes étrangères colonisatrices\teiNote[xmlid={ftn83-3},n={12},place={foot}]{\teiP[xmlid={p-283}]{Liliana Motta, « Éloge du dehors », \teiHi[rend={italic}]{Chimères}, vol. 82, n\teiHi[rend={sup}]{o} 1, 2014, p. 11-18. Certaines d’entre elles comme les polygonum ont reçu le titre de « pestes étrangères colonisatrices ».}} » ou éléments perturbateurs, elles sont celles qui dévastent, ravagent et souillent\teiNote[xmlid={ftn84-3},n={13},place={foot}]{\teiP[xmlid={p-284}]{Voir Mary Douglas, \teiHi[rend={italic}]{De la souillure. Essai sur les notions de pollution et de tabou}, Paris, La Découverte, 2005. Ce qui n’est pas à sa place est de l’ordre de la souillure.}}. }
            
\end{teiDiv}
            \begin{teiDiv}[type={section1},xmlid={div02-5}]

              \teiHead[xmlid={inesthetisme-et-souillure-001}]{Inesthétisme et souillure}
              \teiP[xmlid={p6-8}]{L’idée de souillure n’est pas sans lien avec leur prétendu inesthétisme. Non seulement on ne leur reconnaît aucune beauté, mais on considère qu’elles créent de la dysharmonie. Elles fabriquent de la laideur. Elles la propagent en vérolant les jardins. Leur supposé inesthétisme n’est jamais exempt de connotation morale comme l’indique le terme « mauvais » dans l’appellation « mauvaises herbes ». Ce terme opposé à « bon » renvoie étonnamment à une projection de la théorie des vertus et des vices appliquée à la nature. Tels sont appréhendés entre autres le chiendent, le pissenlit, l’ortie et l’ivraie. Virgile, dans les \teiHi[rend={italic}]{Bucoliques,} ne parle-t-il pas de la « funeste ivraie\teiNote[xmlid={ftn85-3},n={14},place={foot}]{\teiP[xmlid={p-285}]{Virgile, \teiHi[rend={italic}]{Bucoliques}, V, 36-37. « Et la funeste ivraie avec la folle avoine / envahissent nos champs dont l’orge était si belle », Paul Valéry (trad.), Paris, Gallimard, 1966. }} » ? Funeste, l’ivraie, comme l’indique son étymologie grecque \teiHi[rend={italic}]{zizanion}, sème la zizanie. Utilisée comme arme et semée dans le champ de l’ennemi, la loi romaine dut en réprimer la pratique\teiNote[xmlid={ftn86-3},n={15},place={foot}]{\teiP[xmlid={p-286}]{Voir Justinien, \teiHi[rend={italic}]{Les cinquante livres du Digeste,} livre IX, loi 27, alinéa 14, Henri Hulot (trad.), Paris, Chez Rondonneau, 1804, t. 6, p.114.}}. }
              \teiP[xmlid={p7-8}]{Connotée négativement, l’ivraie est, dans la \teiHi[rend={italic}]{Bible}, la marque du malin\teiNote[xmlid={ftn87-3},n={16},place={foot}]{\teiP[xmlid={p-287}]{Parabole de Mathieu, 13, 24-30.}}. L’Église l’utilise pour transmettre un message spirituel : Le Seigneur sème le bon grain, la nuit le diable sème l’ivraie\teiNote[xmlid={ftn88-3},n={17},place={foot}]{\teiP[xmlid={p-288}]{Proverbes, 6, 19.}}. Si l’une et l’autre croissent, seul le Seigneur au moment opportun est apte à séparer le bon grain de l’ivraie. L’ivraie est alors brûlée comme le sont les fils du mal au jour du jugement dernier. }
              \teiP[xmlid={p8-8}]{L’ivraie comme les autres mauvaises herbes sont symboles du mal, des vices et des péchés\teiNote[xmlid={ftn89-3},n={18},place={foot}]{\teiP[xmlid={p-289}]{Isaïe, 34, 13.}}. Les Proverbes l’indiquent : « Auprès du champ d’un paresseux et auprès de la vigne d’un homme privé de cœur \lbrack{}…\rbrack{} les orties poussent partout, les chardons en couvrent la surface\teiNote[xmlid={ftn90-3},n={19},place={foot}]{\teiP[xmlid={p-290}]{Proverbes, 24, 30-31.}}. »}
            
\end{teiDiv}
            \begin{teiDiv}[type={section1},xmlid={div03-4}]

              \teiHead[xmlid={inutiles-et-nuisibles-001}]{Inutiles et nuisibles}
              \teiP[xmlid={p9-7}]{Les « mauvaises herbes », différant des herbes médicinales comme des herbes potagères, sont non seulement décrétées inutiles, mais aussi dangereuses pour la santé. L’ivraie en est de nouveau la parfaite illustration. En latin \teiHi[rend={italic}]{ebriaca herba, }du bas latin \teiHi[rend={italic}]{ebriacus, }du latin classique \teiHi[rend={italic}]{ebrius}, elle signifie l’ivresse, l’ébriété. L’ivraie est réputée provoquer une sorte d’enivrement. Renommée herbe d’ivrogne, Galien parle des maux de tête que sa consommation engendre\teiNote[xmlid={ftn91-3},n={20},place={foot}]{\teiP[xmlid={p-291}]{Paul-Victor Fournier, \teiHi[rend={italic}]{Dictionnaire des plantes médicinales et vénéneuses de France}, Paris, Omnibus, 1947, p. 525. Il faudra attendre le \teiHi[rend={small-caps}]{xx}\teiHi[rend={sup}]{e} siècle pour qu’il soit établi que la toxicité ne dépendait pas de l’ivraie, mais d’un champignon qui l’infestait en produisant des alcaloïdes toxiques comme la témuline.}}. }
              \teiP[xmlid={p10-7}]{Au nom de dimensions hygiénistes, esthétiques et économiques, des pesticides sont toujours employés pour éradiquer lesdites indésirables. }
              \teiP[xmlid={p11-7}]{Immaîtrisables, indomptables, les mauvaises herbes sont considérées comme inutiles, nuisibles et laides. Générant de la dysharmonie, elles sont également symbole du mal. Ces critères sont proches de ceux que l’on attribue à celles que l’on nomme familièrement les « bestioles ». }
            
\end{teiDiv}
            \begin{teiDiv}[type={section1},xmlid={div04-2}]

              \teiHead[xmlid={les-bestioles-001}]{Les « bestioles »}
              \teiP[xmlid={p12-7}]{Du fait de leur petite taille qui peut varier de ¼ de millimètre à 30 centimètres, comme le note Jean-Marc Drouin dans son ouvrage \teiHi[rend={italic}]{La philosophie de l’insecte}\teiNote[xmlid={ftn92-3},n={21},place={foot}]{\teiP[xmlid={p-292}]{Jean-Marc Drouin, \teiHi[rend={italic}]{La philosophie de l’insecte}, Paris, Éditions du Seuil, 2008.}}, les insectes et les arachnides sont le plus souvent méprisés. Michelet nous interpelle à ce propos. « Quelle taille faut-il pour mériter votre estime ? » écrit-il en 1867 dans son ouvrage intitulé \teiHi[rend={italic}]{L’insecte, l’infini vivant}\teiNote[xmlid={ftn93-3},n={22},place={foot}]{\teiP[xmlid={p-293}]{Jules Michelet, \teiHi[rend={italic}]{L’insecte}, \teiHi[rend={italic}]{l’infini vivant}. \teiHi[rend={italic}]{Étude par Marcellin Berthelot}, Paris, Calmann-Lévy, 1857, p. 359. }}. À la petitesse desdites « bestioles », s’ajoute leur infinie variété qui les rend insaisissables sous des caractéristiques fixes. Plus de 1,3 million d’espèces sont décrites, mais d’autres sont inventoriées au fil des ans. Sans compter qu’on dénombre 50 000 espèces d’arachnides. Innombrables et difficilement catégorisables, les « petites bêtes » dérangent. }
            
\end{teiDiv}
            \begin{teiDiv}[type={section1},xmlid={div05-2}]

              \teiHead[xmlid={du-pullulement-au-degout-de-la-laideur-a-la-salete-001}]{Du pullulement au dégoût, de la laideur à la saleté}
              \teiP[xmlid={p13-6}]{Immaîtrisables et incontrôlables, elles le sont du fait de leur pullulement. Si leur multiplication hors de tout contrôle peut être signe de vitalité, cependant ce pullulement est généralement associé à l’idée d’envahissement. Il génère répulsion et répugnance. Aurel Kolnai l’indique. Outre la pourriture, les excréments, les sécrétions, il faut compter parmi ce qui dégoûte les bêtes rampantes, la vermine, ce qui grouille, adhère, s’agglutine et pullule\teiNote[xmlid={ftn94-3},n={23},place={foot}]{\teiP[xmlid={p-294}]{Aurel Kolnai, \teiHi[rend={italic}]{Le dégoût,} 1929, Olivier Cossé (trad.), Paris, Agalma, 1997, p. 55.}} – conception que l’on retrouve dans l’étude sociologique conduite par Nathalie Blanc\teiNote[xmlid={ftn95-3},n={24},place={foot}]{\teiP[xmlid={p-295}]{Nathalie Blanc, « La blatte, ou le monde en images », dans Stéphane Frioux et Émilie-Anne Pépy (dir.), \teiHi[rend={italic}]{L’animal sauvage entre nuisance et patrimoine. France, }\teiHi[rend={ small-caps italic}]{xvi}\teiHi[rend={italic sup}]{e}\teiHi[rend={italic}]{-}\teiHi[rend={ small-caps italic}]{xxi}\teiHi[rend={italic sup}]{e}\teiHi[rend={italic}]{ siècle}, Lyon, ENS Éditions, 2009, p. 103-114, \teiRef[target={http://books.openedition.org/enseditions/6498}]{http://books.openedition.org/enseditions/6498} (consulté le 25 juin 2025). }}.}
              \teiP[xmlid={p14-6}]{Si le dégoût n’implique pas nécessairement la laideur, cependant l’un et l’autre ne sont pas sans lien. À peine évoque-t-on une prolifération de cafards ou de punaises qu’une déplaisance nous gagne. Inesthétisme et répugnance non seulement se mélangent dans notre appréhension, mais ce type de répulsion, comme l’indique Descartes, aux paragraphes 85 et 89 du \teiHi[rend={italic}]{Traité des passions, }n’est pas sans engendrer une aversion profonde, c’est-à-dire de la haine. Haine qui, comme dans \teiHi[rend={italic}]{Ninive} d’Henrietta Rose-Innes\teiNote[xmlid={ftn96-3},n={25},place={foot}]{\teiP[xmlid={p-296}]{Henrietta Rose-Innes, \teiHi[rend={italic}]{Ninive}, Élisabeth Gilles (trad.), Genève, Éditions ZOÉ, « Écrits d’ailleurs », 2014.}}, ou \teiHi[rend={italic}]{La Cinquième Histoire} de Clarisse Lispector\teiNote[xmlid={ftn97-3},n={26},place={foot}]{\teiP[xmlid={p-297}]{Clarisse Lispector, « La Cinquième Histoire », Claire Varin (trad.), \teiHi[rend={italic}]{Dérives}, n\teiHi[rend={sup}]{os} 37-38-39, 1983.}}, conduit à l’extermination des cafards. Le rapprochement entre haine et laideur est du reste marqué dans certaines langues. Notons ainsi qu’en allemand,\teiHi[rend={italic}]{ häβlich }désigne ce qui est laid, \teiHi[rend={italic}]{der Haβ }caractérise la haine et \teiHi[rend={italic}]{die Häβlichkeit} la laideur\teiNote[xmlid={ftn98-3},n={27},place={foot}]{\teiP[xmlid={p-298}]{Dominique Joseph Mozin et Adolphe Peschier, \teiHi[rend={italic}]{Dictionnaire complet des langues française et allemande}, Stuttgart-Tübingen, J. G. Cotta, 1844, vol. 3, p. 846. }}. Ce mixte de répugnance et d’antipathie conduit à écraser ou à anéantir les « bestioles » pour éviter notamment d’être contaminé par la saleté dont on les croit porteuses. Saleté et laideur sont dans ce cadre en accointance comme elles le sont du reste dans de nombreuses langues. En italien\teiHi[rend={italic}]{ bruttamento,} (la souillure), \teiHi[rend={italic}]{bruttare} (salir), \teiHi[rend={italic}]{bruttato} (souillé), ont la même racine que le terme qui caractérise le laid \teiHi[rend={italic}]{brutto}, ou la laideur, \teiHi[rend={italic}]{bruttura}\teiNote[xmlid={ftn99-3},n={28},place={foot}]{\teiP[xmlid={p-299}]{Antonio Bruttura et Angelo Maria Renzi, \teiHi[rend={italic}]{Dictionnaire italien-français}, Paris, Baudry Librairie européenne, 1861, p. 121.}}\teiHi[rend={italic}]{. }Laideur et saleté sont également présentes dans le terme corse \teiHi[rend={italic}]{bruttézzaria} signifiant tout autant la laideur, la répugnance que l’immondice. Dans son ouvrage \teiHi[rend={italic}]{De la souillure,} Mary Douglas précise que selon le code de la pureté du Lévitique sont considérés comme animaux impurs ceux qui grouillent ou rampent. Si la question de l’impureté a fait débat, elle est à mettre au compte de « leur apparence répugnante ou de leurs habitudes malpropres\teiNote[xmlid={ftn100-3},n={29},place={foot}]{\teiP[xmlid={p-300}]{Mary Douglas, \teiHi[rend={italic}]{De la souillure}, \teiHi[rend={italic}]{op. cit}., p. 65.}} ». Autrement dit, les petites bêtes sont pour l’homme une abomination\teiNote[xmlid={ftn101-3},n={30},place={foot}]{\teiP[xmlid={p-301}]{Lévitique, XI, 20.}}. }
              \teiP[xmlid={p15-6}]{Inutiles, insaisissables, incontrôlables, hideuses, sales, répugnantes, nocives et symbole du mal, telles sont les qualités négatives attribuées tant aux petites bêtes qu’aux mauvaises herbes. Par là même, elles ont été rejetées hors de tous les paradigmes de beauté et de ce qu’ils impliquent. Et pourtant, si par un changement de regard il est possible de leur reconnaître de la beauté, ne sont-elles pas une condition de possibilité même de la beauté vitale ? }
            
\end{teiDiv}
            \begin{teiDiv}[type={section1},xmlid={div06-2}]

              \teiHead[xmlid={beaute-et-laideur-en-question-001}]{Beauté et laideur en question}
              \teiP[xmlid={p16-6}]{Trois conceptions de la beauté et de la laideur peuvent être distinguées dans cette analyse. La première est relative à la beauté métaphysique qui caractérise l’Être, fixe, immuable, éternel. À une verticalité du beau auquel le sensible renvoie par participation, correspond la laideur comme insaisissable, immaîtrisable, incontrôlable. Au beau corrélé au bien, au \teiHi[rend={italic}]{kaloskagathos}, qui associe beau et bon, renvoie le laid, signe du mal. Le terme grec \teiHi[rend={italic}]{aischros} spécifie le difforme et le vil, le terme laid \teiHi[rend={italic}]{kakos} le mauvais et le malsain. En conséquence, si la beauté est la forme visible du bien et la laideur la forme visible du mal, l’apparence n’est jamais seulement qu’un inesthétisme. Elle est toujours le signe d’autre chose. Autrement dit, le caractère de malignité projeté sur ce type d’espèces vivantes contribue à fabriquer la laideur de leur apparence. Dans le cas des petites bêtes et des mauvaises herbes, l’apparence n’est que l’allégorie du mal.}
              \teiP[xmlid={p17-6}]{Une seconde détermination opposant la beauté et la laideur n’implique pas une dimension métaphysique, mais une subjectivation du beau d’ordre anthropologique. Le beau est référé alors à des critères qui répondent à l’ordre, à l’harmonie et à la mesure. Ce qui est considéré comme « laid » y contrevient. Désordre, dysharmonie, prolifération, souillure sont autant de caractéristiques qui perturbent les normes organisationnelles mises en place par l’homme. L’art des jardins à la française en est une des illustrations. Les mauvaises herbes échappant à la structuration esthétique de l’espace élaborée par l’homme ne peuvent qu’être assimilées à de l’enlaidissement. La laideur est alors le signe de l’indomptable. Ajoutons que si le bon et l’utile n’ont que rarement été associés au beau, par contre l’inutile et le nuisible ont conduit à en fabriquer la laideur. Cafard, araignée, pissenlit ou ortie n’en sont que les exemples les plus communs. Cependant, en considérant que les catégories de l’inutile et du nuisible sont des constructions relatives non seulement à la connaissance que les hommes ont de ces espèces, mais à leurs intérêts à une époque et dans une culture données, on comprend d’autant mieux la fragilité de ce type de propriété et, avec elle, l’idée de laideur qui en émane. André Georges Haudricourt et Louis Hédin l’indiquent à propos du maïs, du riz et du millet. Si les céréales cultivées descendent des « mauvaises herbes\teiNote[xmlid={ftn102-3},n={31},place={foot}]{\teiP[xmlid={p-302}]{André G. Haudricourt et Louis Hédin, \teiHi[rend={italic}]{L’homme et les plantes cultivées}, Auguste Chevalier (préf.), Paris, Gallimard, 1943, p. 87.}} », alors, comme le dit Ralph Waldo Emerson, la mauvaise herbe n’est qu’une plante « dont on n’a pas encore trouvé les vertus\teiNote[xmlid={ftn103-3},n={32},place={foot}]{\teiP[xmlid={p-303}]{Ralph Waldo Emerson, \teiHi[rend={italic}]{Essais politiques et sociaux} \lbrack{}1878\rbrack{}, Marie Dugard (trad.), Librairie Armand Colin, 1912, p. 274-275.}} ». De même, nombre de qualités n’ont que récemment été reconnues aux petites bêtes, comme le ver de terre ou le bousier. Linné en synthétise l’idée : « ce qui nous paraît nuisible nous est le plus utile\teiNote[xmlid={ftn104-3},n={33},place={foot}]{\teiP[xmlid={p-304}]{Carl Von Linné, \teiHi[rend={italic}]{L’équilibre de la nature}, Bernard Jasmin (trad.) et Camille Limoges (introduction et notes), Paris, Vrin, 1972, p. 165.}} ». Un problème cependant se pose, l’utilité entretient-elle un quelconque rapport avec la beauté ? L’utile n’a-t-il pas été selon les mots de Théophile Gautier considéré comme laid\teiNote[xmlid={ftn105-3},n={34},place={foot}]{\teiP[xmlid={p-305}]{Théophile Gautier, préface à \teiHi[rend={italic}]{Mademoiselle de Maupin} (1835). }} ?}
              \teiP[xmlid={p18-6}]{Chaque espèce, énonce Lévi-Strauss\teiNote[xmlid={ftn106-3},n={35},place={foot}]{\teiP[xmlid={p-306}]{Claude Lévi-Strauss, dans l’émission de Jacques Chancel \teiHi[rend={italic}]{Radioscopie}, 15 avril 1981, \teiRef[target={https://www.radiofrance.fr/franceculture/extinctions-animales-125-000-ans-d-influence-du-genre-humain-5996803}]{https://www.radiofrance.fr/franceculture/extinctions-animales-125-000-ans-d-influence-du-genre-humain-5996803} (consulté le 25 juin 2025).}}, est irremplaçable dans la chaîne des êtres vivants où chacun remplit sa fonction et chaque espèce est aussi irremplaçable d’un point de vue esthétique parce que chacune constitue une sorte de chef-d’œuvre, avec une beauté particulière. L’existence d’une espèce est aussi importante que l’existence de l’œuvre d’un grand peintre. Pour une œuvre d’art cependant, nous employons tous nos efforts pour la protéger dans des musées alors que pour une espèce vivante nous la traitons avec une désinvolture, un mépris et une ignorance tout à fait incroyables\teiNote[xmlid={ftn107-3},n={36},place={foot}]{\teiP[xmlid={p-307}]{\teiHi[rend={italic}]{Ibid.}}}.}
              \teiP[xmlid={p19-6}]{À partir de cette approche, une troisième conception de la beauté et de la laideur pointe. Elle est de l’ordre de la vitalité de la vie. Empruntant à Plotin\teiNote[xmlid={ftn108-3},n={37},place={foot}]{\teiP[xmlid={p-308}]{« Pourquoi donc, en effet, la splendeur de la beauté resplendit-elle plus sur un visage vivant, alors qu’il n’y en a qu’une trace sur le visage d’un mort, même s’il n’est pas encore détruit dans sa chair et dans sa symétrie ? Pourquoi un homme vivant, même s’il est laid est-il plus beau qu’un homme beau représenté dans une statue ? C’est pour cette raison qu’il est plus désirable, mais s’il est plus désirable c’est parce qu’il a une âme, c’est-à-dire qu’il possède davantage la forme du Bien ; et il en va ainsi parce qu’il est en quelque sorte coloré par la lumière du Bien et que son âme, ainsi colorée s’éveille, qu’elle est soulevée et emporte avec elle ce qui lui appartient en lui donnant la bonté et l’éveil dont il est capable », Plotin, \teiHi[rend={italic}]{Traité 38 }(VI, 7), chap. 22, l. 23-36, Francesco Fronterotta (trad.), Paris, Flammarion, « GF », 2007. Voir également l’article d’Anne-Lise Worms dans le présent volume.}} l’idée selon laquelle un vivant est beau du simple fait qu’il est vivant, il apparaît alors que le dynamisme de la vie est une condition de possibilité de la beauté. En ce sens, si laideur il y a, elle est le signe non seulement d’une extinction de la vie, mais de tout ce qui participe à engendrer prématurément la mort. Là où il y a anéantissement, la beauté ne peut survivre. En somme, ce n’est pas seulement ce que l’on juge beau qui est beau, mais ce qui en toutes les espèces comme dans leurs interactions est l’émanation d’une force de germination. Si la beauté vitale est besoin premier, elle l’est en tant que possible révélation, et mise en relation qui ouvre sur l’éblouissement. Sans la capacité de s’émouvoir, de sentir, de ressentir, de parler la beauté, nous ne serions pas vraiment vivant\teiNote[xmlid={ftn109-3},n={38},place={foot}]{\teiP[xmlid={p-309}]{« L’homme qui a le plus vécu n’est pas celui qui a compté le plus d’années ; mais celui qui a le plus senti la vie », Jean-Jacques Rousseau, \teiHi[rend={italic}]{L’Émile ou De l’éducation} dans \teiHi[rend={italic}]{Œuvres Complètes}, Paris, Gallimard, « Bibliothèque de la Pléiade », 1969, t. IV, p. 25.}}. La laideur quant à elle n’est pas celle projetée sur lesdites mauvaises herbes et « bestioles », elle est celle de tous ceux pour qui aucune épiphanie au sens étymologique du terme n’advient. La laideur est précipitation dans la mort, elle est en lien avec ce qui provoque l’exténuation, l’asthénie, la destruction.}
            
\end{teiDiv}
            \begin{teiDiv}[type={section1},xmlid={div07-2}]

              \teiHead[xmlid={herbes-et-bestioles-ou-le-paradigme-de-la-beaute-vitale-001}]{Herbes et « bestioles » ou le paradigme de la beauté vitale}
              \teiP[xmlid={p20-6}]{« Mauvaises » herbes et petites bêtes révèlent autre chose. De par leur puissance de germination, elles sont les paradigmes mêmes de toute potentialité inscrite au cœur du vivant. Et c’est cette potentialité qui rend pleinement vivant et contient en germe la beauté.}
              \teiP[xmlid={p21-4}]{Adventice et rudérale, la vitalité de la « mauvaise » herbe est puissance de vie. Elle s’infiltre dans n’importe quel espace. « Elle pousse entre et parmi les autres choses\teiNote[xmlid={ftn110-3},n={39},place={foot}]{\teiP[xmlid={p-310}]{Henry Miller, \teiHi[rend={italic}]{Hamlet}, Paris, Corrêa, 1956, p. 48 et 49, cité dans Gilles Deleuze et Félix Guattari, \teiHi[rend={italic}]{Mille plateaux. Capitalisme et schizophrénie}, Paris, Éditions de Minuit, 1980, p. 29.}} », elle n’a pas d’enracinement,\teiNote[xmlid={ftn111-3},n={40},place={foot}]{\teiP[xmlid={p-311}]{Voir notamment Gilles Deleuze et Félix Guattari, \teiHi[rend={italic}]{Mille plateaux, op. cit.}, p. 11-17.}} mais une ligne de fuite. Nomade, elle est le modèle même d’une pensée qui surgit. Kafka l’écrit dans son \teiHi[rend={italic}]{Journal} : « Les choses qui me viennent à l’esprit se présentent à moi non par leur racine, mais par un point quelconque situé vers leur milieu. Essayez donc de les retenir, essayez donc de retenir un brin d’herbe qui ne commence à croître qu’au milieu de la tige, et de vous tenir à lui\teiNote[xmlid={ftn112-3},n={41},place={foot}]{\teiP[xmlid={p-312}]{Franz Kafka, \teiHi[rend={italic}]{Journal}, Paris, Grasset, 1954, p. 4, cité dans Gilles Deleuze et Félix Guattari, \teiHi[rend={italic}]{Mille plateaux}, \teiHi[rend={italic}]{op. cit}., p. 34.}}. » L’herbe pousse ici et là, elle fait rhizome avec le vent, avec un animal, avec l’homme. Elle est telle une pensée qui entre en rapport avec une autre pensée, telle une pensée qui se connecte à des éléments hétérogènes. Si « le cerveau lui-même est une herbe\teiNote[xmlid={ftn113-3},n={42},place={foot}]{\teiP[xmlid={p-313}]{\teiHi[rend={italic}]{Ibid.}}} », c’est précisément parce qu’il permet des connexions improbables. En somme, l’herbe qui infiltre sa vigueur non policée dans les plantations est bien le paradigme de la pensée créatrice. Pierre Boulez le traduit ainsi : une « forme musicale, jusque dans ses ruptures et proliférations, est comparable à de la mauvaise herbe\teiNote[xmlid={ftn114-3},n={43},place={foot}]{\teiP[xmlid={p-314}]{Pierre Boulez, \teiHi[rend={italic}]{Par volonté et par hasard}, Éditions du Seuil, 1975, p. 14, cité dans Gilles Deleuze et Félix Guattari, \teiHi[rend={italic}]{Mille plateaux}, \teiHi[rend={italic}]{op. cit}., p. 19.}} ». Il en va de même de l’écriture ou de l’art qui font réseaux ou rhizome. Vitale, la beauté est semblable à ces flux. Elle se révèle. Elle ne maintient pas en vie, elle rend vivant. On comprend alors Gilles Deleuze et Félix Guattari quand ils écrivent : « rien n’est beau, rien n’est amoureux, rien n’est politique, sauf les tiges souterraines et les racines aériennes, l’adventice et le rhizome\teiNote[xmlid={ftn115-3},n={44},place={foot}]{\teiP[xmlid={p-315}]{\teiHi[rend={italic}]{Ibid.}, p. 24.}} ».}
              \teiP[xmlid={p22-4}]{De même, les petites bêtes, empreintes de vitalité et de résistance, explorent, investissent, parasitent, marquent des territoires et se nourrissent de multiples matières. L’araignée tisse sa toile, elle range « les filets, les uns tout droits pour soutenir l’ouvrage, les autres en ronds\teiNote[xmlid={ftn116-3},n={45},place={foot}]{\teiP[xmlid={p-316}]{Sénèque, \teiHi[rend={italic}]{Lettres à Lucilius}, lettre 121, Joseph Baillard (trad.), Paris, Hachette, 1914.}} ». Modèle même du fil de la pensée qui se dénoue, l’homme qui réfléchit s’inspire selon Démocrite\teiNote[xmlid={ftn117-3},n={46},place={foot}]{\teiP[xmlid={p-317}]{Plutarque, \teiHi[rend={italic}]{L’intelligence des animaux}, Myrto Gondicas (trad.), Paris, Arléa, 2012, § 20.}} de l’ouvrage de l’araignée. Pour Plutarque, « poètes et prosateurs tissent et tendent comme les araignées font leur toile\teiNote[xmlid={ftn118-3},n={47},place={foot}]{\teiP[xmlid={p-318}]{Plutarque, \teiHi[rend={italic}]{Isis et Osiris}, § 20 (358d), dans \teiHi[rend={italic}]{Œuvres morales}, Christian Froidefond (éd. et trad.), Paris, Les Belles Lettres, 1988, t. 5, p. 194.}} ». Paul Valéry indique également qu’il écrit « ces notes comme une araignée sur sa toile ». Quant à Mallarmé, à partir de différentes idées, il réalise « aux points de rencontre de merveilleuses dentelles\teiNote[xmlid={ftn119-3},n={48},place={foot}]{\teiP[xmlid={p-319}]{Stéphane Mallarmé, « À Théodore Aubanel, 28 juillet 1866 », dans \teiHi[rend={italic}]{Correspondance 1862-1871}, Paris, Gallimard, 1995, p. 315.}} ». Que l’on soit mauvaises herbes ou bestioles, écrivains ou créateurs, dans la vitalité de cet élan s’affirme une puissance de vie et cette puissance de vie en tant qu’elle est invitation à l’échange entre les espèces comme entre les hommes est génératrice en puissance de beauté. Autrement dit, la beauté est vitale car elle est rencontre. }
            
\end{teiDiv}
            \begin{teiDiv}[type={section1},xmlid={div08-2}]

              \teiHead[xmlid={rencontre-et-devenir-001}]{Rencontre et devenir}
              \teiP[xmlid={p23-4}]{Pour qu’il y ait rencontre, il ne suffit pas d’être confronté à quelqu’un ou quelque chose, il faut que des signes nous obligent à sentir autrement, que quelque chose nous percute et que nous soyons comme renvoyés hors de nous-mêmes. Que notre être, nos convictions, nos repères soient ébranlés. La rencontre n’est pas nécessairement rencontre d’une personne. Deleuze l’indique. On peut rencontrer une idée, un livre, un animal, une œuvre d’art, un paysage, un brin d’herbe. La rencontre n’implique pas communion mais « communication sans mise en commun », « bloc de devenir\teiNote[xmlid={ftn120-3},n={49},place={foot}]{\teiP[xmlid={p-320}]{Gilles Deleuze et Claire Parnet, \teiHi[rend={italic}]{Dialogues}, Paris, Flammarion, « Champs Essais », 1996, p. 88.}} ». Dans le devenir animal notamment, il y a rencontre entre deux règnes, celle d’un homme et d’un animal\teiNote[xmlid={ftn121-3},n={50},place={foot}]{\teiP[xmlid={p-321}]{Gilles Deleuze et Félix Guattari, \teiHi[rend={italic}]{Mille plateaux, op. cit.,} p. 336.}}. Il n’y a pas de fusion ou de ressemblance. Il s’agit d’intégrer en soi la manière de sentir de l’autre sans pour autant être l’autre. À partir de là, il n’y a plus ni homme ni animal, mais déterritorialisation, « chacun déterritorialise l’autre, dans une conjonction de flux, dans un continuum d’intensités réversibles\teiNote[xmlid={ftn122-3},n={51},place={foot}]{\teiP[xmlid={p-322}]{Gilles Deleuze et Félix Guattari, \teiHi[rend={italic}]{Kafka. Pour une littérature mineure}, Paris, Éditions de Minuit, 1975, p. 40.}} ». Ainsi s’invente une relation autre, qu’elle soit à l’herbe ou à la bestiole, elle est rencontre\teiNote[xmlid={ftn123-3},n={52},place={foot}]{\teiP[xmlid={p-323}]{Gilles Deleuze, \teiHi[rend={italic}]{Critique et clinique}, Paris, Éditions de Minuit, 1993, p. 79. En se référant au texte de Whitman \teiHi[rend={italic}]{Les chênes et moi}, dans \teiHi[rend={italic}]{Critique et clinique}, Deleuze parle de la rencontre entre le poète et de jeunes chênes comme d’un corps à corps. Le poète ne se fond pas, ne se confond pas avec les chênes, mais quelque chose se passe entre eux, entre le corps humain et l’arbre. Le corps recevant comme l’écrit Whitman « un peu de sève claire et de fibre élastique » et l’arbre de son côté recevant un peu de conscience. Peut-être faisons-nous un échange, écrit Whitman.}}, résonnance avec un dehors. Dans la \teiHi[rend={italic}]{Métamorphose}, Kafka n’est pas cette bête immonde, « ce monstrueux » insecte jamais nommé, mais il le « devient » en ressentant ce qu’il pressent de ce que ressent l’animal\teiNote[xmlid={ftn124-3},n={53},place={foot}]{\teiP[xmlid={p-324}]{Voir Gilles Deleuze et Félix Guattari, \teiHi[rend={italic}]{Mille plateaux. Capitalisme et schizophrénie}, Paris, Éditions de Minuit, 1980, p. 293.}}. Écrire comme un cafard ce n’est pas imiter le cafard, c’est rendre vivant le cafard. Il ne s’agit pas de décrire avec exactitude, mais d’écrire comme si l’écriture ressentait ce que le cafard éprouve. Comme si le ressenti de l’animal passait dans les mots, se glissait entre les mots, comme si les mots enveloppaient l’animal. On comprend alors que les mots, la syntaxe en soient eux-mêmes bouleversés. Il n’est plus possible de distinguer ce qui revient au cafard ou à la langue. On assiste en quelque sorte à une animalité de l’écriture et à une écriture de l’animalité. }
              \teiP[xmlid={p24-4}]{Écrire et ressentir sont une seule et même chose dans cette zone d’indiscernabilité\teiNote[xmlid={ftn125-3},n={54},place={foot}]{\teiP[xmlid={p-325}]{Voir \teiHi[rend={italic}]{Ibid., }p. 638\teiHi[rend={italic}]{.}}}. Ainsi, « en écrivant on donne toujours l’écriture à ceux qui n’en ont pas, mais ceux-ci donnent à l’écriture un devenir sans lequel elle ne serait pas\teiNote[xmlid={ftn126-3},n={55},place={foot}]{\teiP[xmlid={p-326}]{Gilles Deleuze et Claire Parnet, \teiHi[rend={italic}]{Dialogues}, \teiHi[rend={italic}]{op. cit}., p. 55.}} ».}
              \teiP[xmlid={p25-3}]{Dans une approche distincte, Catherine Chalmers, artiste américaine, nous montre ce que nous n’avons pas envie de voir. Dans la série \teiHi[rend={italic}]{Execution}, pour figurer ce que subissent certaines bestioles, elle n’hésite pas à mettre en scène des cafards gazés, électrocutés, brûlés. Si elle choisit comme personnage principal de son œuvre le cafard, c’est précisément parce qu’il est pour nombre d’individus « le vrai nom de la laideur », selon les termes qu’emploie Clarisse Lispector dans \teiHi[rend={italic}]{La Cinquième Histoire}\teiNote[xmlid={ftn127-3},n={56},place={foot}]{\teiP[xmlid={p-327}]{Clarisse Lispector, art. cité.}}. Lors d’un entretien avec Lesley Marin, Catherine Chalmers insiste sur le fait que « cette détestation est totalement disproportionnée et injustifiable. Les blattes ne piquent pas, ne mordent pas, ne transportent pas d’agents pathogènes extrêmement dangereux pour l’homme ». « Le cafard », dit-elle, « est remarquablement subtil et beau\teiNote[xmlid={ftn128-2},n={57},place={foot}]{\teiP[xmlid={p-328}]{« Excerpts from a Conversation, with Lesley A. Martin », \teiHi[rend={italic}]{American Cockroach}, n\teiHi[rend={sup}]{o} 20, 2004, \teiRef[target={https://www.catherinechalmers.com/more\#/american-cockroach-interview/}]{\teiHi[rend={underline}]{htts://www.catherinechalmers.com/more\#/american-cockroach-inteview/}} (consulté le 25 juin 2025).}} ». Ses ailes sont d’un ambre brillant et translucide, ses pattes souples sont bien dessinées et ses antennes explorent le monde avec la grâce des bras d’une ballerine. Se défaire de nos présupposés permet de ressentir différemment non seulement le cafard, mais les sublimes couleurs des tarentules après leur mue ou la vivacité des couleurs de certaines araignées. Ainsi, voir avec d’autres yeux, comme le défendait Marcel Proust, peut conduire à « sortir de notre condition humaine \lbrack{}…\rbrack{} \lbrack{}à\rbrack{} s’imaginer la vie d’une blatte, \lbrack{}à\rbrack{} explorer le monde avec de longues antennes\teiNote[xmlid={ftn129-2},n={58},place={foot}]{\teiP[xmlid={p-329}]{Voir \teiRef[target={https://www.catherinechalmers.com/more}]{https://www.catherinechalmers.com/more} (consulté le 25 juin 2025).}} ». Tomas Saraceno, plasticien argentin, nous y invite par-delà sa rencontre avec des araignées. Ces œuvres ne sont nullement des toiles, ou plutôt si, mais ce sont celles des petites ouvrières patientes et talentueuses qu’il élève dans son atelier de Berlin et expose sous le feu de délicats projecteurs. Comme Catherine Chalmers avec les cafards, Tomas Saraceno avec les araignées révèle dans cette rencontre d’une singulière altérité une vitalité persistante. « Quand je suis arrivé sur les lieux de l’exposition », confie-t-il, « il y avait déjà des araignées qui étaient ici chez elles. Certaines sont venues coloniser les toiles de celles que j’avais amenées avec moi\teiNote[xmlid={ftn130-2},n={59},place={foot}]{\teiP[xmlid={p-330}]{Emmanuel Grandjean, « Saraceno, l’homme araignée », \teiHi[rend={italic}]{Le Temps}, 31 août 2019, \teiRef[target={https://www.letemps.ch/societe/tomas-saraceno-lartiste-araignee}]{https://www.letemps.ch/societe/tomas-saraceno-lartiste-araignee} (consulté le 26 juin 2025).}}\teiHi[rend={sup}]{ }». Ce tissage de concert, de fils et de vibrations est pour l’artiste « une forme de beauté ». Le travail de Tomas Saraceno révèle également que les toiles d’araignées sont une métaphore du lien. Elles traduisent ce que nous sommes au travers des relations que nous entretenons avec les autres espèces. }
            
\end{teiDiv}
            \begin{teiDiv}[type={section1},xmlid={div09}]

              \teiHead[xmlid={creer-cest-resister-001}]{« Créer c’est résister »}
              \teiP[xmlid={p26-3}]{Contre tous les actes abjects, « devant la bassesse et la vulgarité d’existence », il ne reste que la résistance\teiNote[xmlid={ftn131-2},n={60},place={foot}]{\teiP[xmlid={p-331}]{Gilles Deleuze et Félix Guattari, \teiHi[rend={italic}]{Qu’est-ce que la philosophie ?}, Paris, Éditions de Minuit, 1991, p. 103. }}. Si « créer, c’est résister\teiNote[xmlid={ftn132-2},n={61},place={foot}]{\teiP[xmlid={p-332}]{\teiHi[rend={italic}]{Ibid.,} p. 106.}} », tout créateur qui résiste « à la servitude, à l’intolérable, à la honte\teiNote[xmlid={ftn133-2},n={62},place={foot}]{\teiP[xmlid={p-333}]{\teiHi[rend={italic}]{Ibid., }p. 105.}} » libère de « nouvelles possibilités de vie\teiNote[xmlid={ftn134-2},n={63},place={foot}]{\teiP[xmlid={p-334}]{Voir Gilles Deleuze, \teiHi[rend={italic}]{Nietzsche et la philosophie}, Paris, PUF, 1962, p. 115. }} ». Il « libère une vie, une vie puissante, une vie plus que personnelle. Ce n’est pas “sa” vie qu’il libère, mais la vie\teiNote[xmlid={ftn135-2},n={64},place={foot}]{\teiP[xmlid={p-335}]{Gilles Deleuze et Claire Parnet, « R comme résistance », dans \teiHi[rend={italic}]{Abécédaire}, documentaire, Pierre-André Boutang (écriture et réalisation), France, Sodaperaga, 1996, 450 minutes.}} ». En transformant notre pouvoir d’affecter et d’être affecté, l’écriture porte « la vie à sa puissance absolue\teiNote[xmlid={ftn136-2},n={65},place={foot}]{\teiP[xmlid={p-336}]{Gilles Deleuze et Claire Parnet, \teiHi[rend={italic}]{Dialogues}, \teiHi[rend={italic}]{op. cit}., p. 12.}}\teiHi[rend={sup}]{ }», elle est santé\teiNote[xmlid={ftn137-2},n={66},place={foot}]{\teiP[xmlid={p-337}]{Gilles Deleuze, \teiHi[rend={italic}]{Critique et clinique}, Paris, Éditions de Minuit, 1993, p. 9.}} ou « grande santé\teiNote[xmlid={ftn138-2},n={67},place={foot}]{\teiP[xmlid={p-338}]{Gilles Deleuze et Claire Parnet, \teiHi[rend={italic}]{Dialogues}, \teiHi[rend={italic}]{op. cit}., p. 12.}} ». Subversive en puissance, elle est respiration, nouveau souffle. Ouverte à tous les possibles, elle est créatrice. Vitale elle rend pleinement vivant car elle est révélation et cette révélation est beauté. }
              \teiP[xmlid={p27-3}]{En inaugurant une nouvelle manière de percevoir et de sentir bestioles grouillantes ou « mauvaises » herbes poussant dans la pensée, l’œuvre nous invite à voir autrement, à dissoudre les clichés, les opinions, les schèmes de pensée fixes. Elle ouvre des zones inexplorées, inexpérimentées. En nettoyant les perceptions communes, les affections ordinaires et opinions toujours préétablies, l’écriture et l’art créent un langage fait d’affects, de percepts et de blocs de sensations\teiNote[xmlid={ftn139-2},n={68},place={foot}]{\teiP[xmlid={p-339}]{Gilles Deleuze et Félix Guattari, \teiHi[rend={italic}]{Qu’est-ce que la philosophie ?}, \teiHi[rend={italic}]{op. cit.}, p. 166.}}. Par affects il faut entendre des principes d’intensité, des puissances d’affirmations, et par percepts un ensemble de perceptions et de sensations qui vont survivre à ceux qui les ont éprouvées\teiNote[xmlid={ftn140-2},n={69},place={foot}]{\teiP[xmlid={p-340}]{\teiHi[rend={italic}]{Ibid.,} p. 159.}}. Par son langage, l’artiste comme l’écrivain extraient des forces et, comme l’écrit Paul Klee, « il ne reproduit pas le visible, il rend visible\teiNote[xmlid={ftn141-2},n={70},place={foot}]{\teiP[xmlid={p-341}]{Paul Klee, \teiHi[rend={italic}]{Théorie de l’art moderne}, Paris, Gonthier, 1964, p. 34.}} ». Créateur d’agrégats sensibles, l’artiste « nous les donne et nous fait devenir avec eux, il nous prend dans le composé\teiNote[xmlid={ftn142-2},n={71},place={foot}]{\teiP[xmlid={p-342}]{Gilles Deleuze et Félix Guattari, \teiHi[rend={italic}]{Qu’est-ce que la philosophie ?}, \teiHi[rend={italic}]{op. cit}., p. 166.}} ». Il y atteint en parvenant à faire émerger l’« incommensurable nouveauté qu’on ne savait plus voir\teiNote[xmlid={ftn143-2},n={72},place={foot}]{\teiP[xmlid={p-343}]{\teiHi[rend={italic}]{Ibid.}}} ». La beauté n’est plus alors forme ou être, elle est la force du ressenti et par là même devient vitale.}
            
\end{teiDiv}
          
\end{teiElement}
        
\end{teiElement}
        \begin{teiElement}[name={group},type={chapter},data-page-title={Beautés minérales, sources d’inspiration et de création},data-page-authors={Didier Nectoux@@École Mines Paris-PSL, musée de Minéralogie, France},data-include-href={Ch08\_beautes\_minerales\_Nectoux.xml},xmlid={chapter-007}]

          \begin{teiElement}[name={front}]

            \begin{teiDiv}[type={titlePage},xmlid={titlepage-009}]

              \teiP[rend={title-main},xmlid={p-344}]{Beautés minérales, sources d’inspiration et de création}
              \teiP[rend={author-aut},xmlid={p-345}]{Didier \teiHi[rend={small-caps}]{Nectoux}}
              \teiP[rend={authority\_affiliation},xmlid={p-346}]{École Mines Paris-PSL, musée de Minéralogie, France}
            
\end{teiDiv}
          
\end{teiElement}
          \begin{teiElement}[name={body},xmlid={body-009}]

            \teiP[xmlid={p1-9}]{Pour appréhender la beauté minérale, il suffit d’aller visiter le musée de Minéralogie de l’École des mines de Paris. C’est un lieu préservé du temps où les minéraux sont présentés dans un décor du \teiHi[rend={small-caps}]{xix}\teiHi[rend={sup}]{e} siècle. C’est en effet Pierre Ours Armand Dufrenoy (1792-1857), alors directeur de cette prestigieuse école, mais aussi professeur de minéralogie, qui a commandé la réalisation de cet écrin de vitrines, placards et tiroirs en chêne de Hongrie afin d’abriter ce qui est aujourd’hui considéré comme l’une des plus importantes collections mondiales. }
            \begin{teiFigure}[xmlid={figure01-3}]

              \teiHead[xmlid={figure-8-galerie-de-mineralogie-musee-de-mineralogie-mines-paris-psl-photographie-didier-nectoux-001}]{Figure 8. Galerie de minéralogie. Musée de Minéralogie, Mines Paris-PSL, photographie Didier Nectoux.}
            
\end{teiFigure}
            \teiP[xmlid={p2-9}]{La volonté esthétique est évidente, mais elle n’est pas la raison première de l’exposition des milliers de minéraux qui la composent. La collection a été créée en 1794 par le Comité de salut public, avec pour objectif de faire l’inventaire des ressources minérales de la République et du globe dans son intégralité. Elle constitue un objet stratégique sur lequel les scientifiques vont se pencher. Elle servira aussi à l’enseignement de la minéralogie et de la géologie à destination des futurs ingénieurs des mines, chargés de l’exploitation et la transformation de ces matières minérales, indispensables au développement de l’industrie dans l’intégralité de ses secteurs. À l’occasion de ces recherches et pendant l’exploitation minière, les ingénieurs et les géologues mettront à jour des minéraux exceptionnels par la taille, la couleur et la forme. Ces pièces uniques seront mises de côté pour leur beauté et acheminées vers la collection, au même titre que les dizaines de milliers d’autres échantillons nécessitant analyse et classement. }
            \teiP[xmlid={p3-9}]{Lorsqu’au \teiHi[rend={small-caps}]{xx}\teiHi[rend={sup}]{e} siècle la collection deviendra musée et s’ouvrira au grand public, la dimension esthétique s’imposera comme une évidence pour susciter émerveillement et curiosité chez le visiteur. La salle d’entrée du musée et le salon central sont exemplaires à cet égard.}
            \begin{teiFigure}[xmlid={figure02-2}]

              \teiHead[xmlid={figure-9-salon-central-du-musee-musee-de-mineralogie-mines-paris-psl-photographie-didier-nectoux-001}]{Figure 9. Salon central du musée. Musée de Minéralogie, Mines Paris-PSL, photographie Didier Nectoux.}
            
\end{teiFigure}
            \teiP[xmlid={p4-9}]{Les plus beaux spécimens ont été sélectionnés pour leur beauté et placés dans des alvéoles éclairées. L’effet est saisissant de beauté et fait vibrer le visiteur, étonné par la diversité des formes, des couleurs, des éclats. Pour certains échantillons, c’est même l’incrédulité qui surgit : il ne semble pas croyable que ces objets soient naturels.}
            \begin{teiFigure}[xmlid={figure03-2}]

              \teiHead[xmlid={figure-10-sepiolite-ensmp-6280-eskisehir-turquie-musee-de-mineralogie-mines-paris-psl-photographie-didier-nectoux-001}]{Figure 10. Sépiolite, ENSMP 6280. Eskisehir, Turquie. Musée de Minéralogie, Mines Paris – PSL, photographie Didier Nectoux.}
            
\end{teiFigure}
            \teiP[xmlid={p5-9}]{Pour illustrer ce point, citons un mot d’enfant inscrit dans le livre d’or du musée : « Les minéraux sont tellement beaux qu’on dirait des faux ! » On ne peut mieux traduire l’impression ressentie par celle ou celui qui pénètre ce lieu.}
            \teiP[xmlid={p6-9}]{La beauté intrinsèque du minéral peut s’explorer en s’appuyant sur des critères descriptifs (forme, texture et couleur) qui sont ceux-là mêmes qui ont permis aux scientifiques de les caractériser et de les classer. Ils sont la source d’inspiration et de création qui touche tous les domaines artistiques. Architectes, peintres, sculpteurs et designers s’en empareront pour proposer des œuvres originales, à leur tour sources d’étonnement et d’émotions.}
            \begin{teiDiv}[type={section1},xmlid={div01-6}]

              \teiHead[xmlid={la-forme-001}]{La forme}
              \begin{teiDiv}[type={section2},xmlid={div02-6}]

                \teiHead[xmlid={rappels-historiques-001}]{Rappels historiques }
                \teiP[xmlid={p7-9}]{Parmi les cinq solides platoniciens, le cube, le tétraèdre, l’octaèdre sont des formes que l’on retrouve facilement en observant les cristaux comme la pyrite, le diamant, la fluorine et de nombreux minéraux du système cubique. Ils témoignent, selon les anciens, d’une organisation du monde fait de symétrie. La recherche d’une harmonie fut constante au cours des siècles. La Renaissance vit sa formalisation dans des ouvrages comme le \teiHi[rend={italic}]{De Divina proportione} de Pacioli et Vinci (Venise, Paganino Paganini, 1509) ou, plus tard, dans les \teiHi[rend={italic}]{Harmonices Mundi} de Kepler (Linz, Johann Planck pour \teiHi[rend={italic}]{Gottfried Tampach}, 1619). La volonté d’expliquer la forme des cristaux en y associant des tentatives de classification y est permanente, mais toujours incomplète et insatisfaisante. }
                \teiP[xmlid={p8-9}]{Il faut attendre le \teiHi[rend={small-caps}]{xviii}\teiHi[rend={sup}]{e} siècle pour que surgissent des propositions plus élaborées et satisfaisantes. À partir d’observations et de mesures, des scientifiques comme le Suédois Torbern Olof Bergman (\teiHi[rend={italic}]{Manuel du minéralogiste}, Paris, Chez Cuchet, 1784) ou le français Jean-Baptiste Romé de L’Isle (\teiHi[rend={italic}]{Cristallographie} ou \teiHi[rend={italic}]{Description des formes propres à tous les corps du règne minéral dans l’état de combinaison saline, pierreuse ou métallique}, Paris, Imprimerie de Monsieur, 1783) énoncent des principes et des lois qui constituent les éléments précurseurs de la cristallographie. Mais c’est l’abbé René Just Haüy (\teiHi[rend={italic}]{Traité de minéralogie}, Paris, Imprimerie de Delance, 1802 et \teiHi[rend={italic}]{Traité de cristallographie}, Paris, Bachelier et Huzard, 1817) qui apportera une contribution déterminante et fondatrice de la minéralogie moderne et de la cristallographie. Ses définitions de la molécule intégrante et des formes primitives ne sont pas autre chose que celles de la maille cristalline et des systèmes cristallins.}
                \teiP[xmlid={p9-8}]{Grâce à ses travaux, on arrive très rapidement à la conclusion selon laquelle le monde minéral se répartit en sept systèmes d’organisation des atomes qui sont les systèmes cubique, quadratique, orthorhombique, rhomboédrique, hexagonal, monoclinique et triclinique. }
                \begin{teiFigure}[xmlid={figure04-2}]

                  \teiHead[xmlid={figure-11-les-systemes-cristallins-schema-musee-de-mineralogie-mines-paris-psl-001}]{Figure 11. Les systèmes cristallins, schéma. Musée de Minéralogie, Mines Paris – PSL.}
                
\end{teiFigure}
                \teiP[xmlid={p10-8}]{Ces formes géométriques de base se retrouvent exprimées par les minéraux lorsqu’ils sont bien formés. Les cristaux trouvés dans la nature, exposés dans les musées, deviennent alors sources d’inspiration pour les artistes.}
                \teiP[xmlid={p11-8}]{Ceux-ci empruntent directement formes et symétrie aux minéraux, comme on le verra ci-après. Mais il arrive que les œuvres s’affranchissent des lois de symétrie cristalline. L’artiste Mylène Guermont crée ainsi \teiHi[rend={italic}]{Cristal A}, une œuvre en béton pour laquelle elle a pris soin de ne construire aucun axe de symétrie. }
                \teiP[xmlid={p12-8}]{Cette création n’a pas d’équivalent dans la nature, et fait cependant penser immanquablement à un minéral. Mœbius, dans ses dessins, représente aussi des formes cristallines impossibles. On peut se demander s’il s’agit de la même volonté de détournement lorsqu’on découvre des représentations de cristaux (souvent proches d’un quartz) dans différentes œuvres picturales, y compris les tatouages. Dans de nombreux cas, il est plus probable que la représentation erronée du minéral soit le fruit d’une méconnaissance des règles géométriques cristallines.}
              
\end{teiDiv}
              \begin{teiDiv}[type={section2},xmlid={div03-5}]

                \teiHead[xmlid={systeme-cubique-001}]{Système cubique}
                \teiP[xmlid={p13-7}]{Le système cubique fournit les formes les plus évidentes. Des minéraux courants comme le sel gemme (halite), la pyrite ou la fluorine offrent des formes et des compositions parfaites géométriquement. Les architectes ne s’y sont pas trompés. Citons les \teiHi[rend={italic}]{Maisons cubes }de Harry Reijnders et László Vákár en 1977 à Rotterdam, ou encore les \teiHi[rend={italic}]{Maisons cubes} de MVRDV Architectes en 2016 à Logrono (Espagne), directement inspirées des assemblages de pyrite.}
                \begin{teiFigure}[xmlid={figure05-2}]

                  \teiHead[xmlid={figure-12-pyrite-ensmp-16759-huanuco-perou-musee-de-mineralogie-mines-paris-psl-photographie-didier-nectoux-001}]{Figure 12. Pyrite ENSMP 16759, Huanuco, Perou. Musée de Minéralogie, Mines Paris – PSL, photographie Didier Nectoux.}
                
\end{teiFigure}
                \teiP[xmlid={p14-7}]{Les concepteurs de meubles s’emparent de ces compositions, ainsi Laurence Manders et son\teiHi[rend={italic}]{ Cabinet magique}. La Maison Rapin propose même une « Collection pyrite » dans sa gamme de meubles. Le nom de ce minéral se retrouve également dans les œuvres de Marie Guerrier. David Jameson, architecte, préfère, quant à lui, NaCl (ou halite) House pour sa maison construite à Bethesda dans le Maryland (États-Unis).}
                \teiP[xmlid={p15-7}]{La forme la plus insolite de construction est au Japon. Elle est celle de \teiHi[rend={italic}]{Mineral House} en 2013, maison individuelle de John Perreault qui s’inspire directement du pyritoèdre aux faces à 5 côtés.}
              
\end{teiDiv}
              \begin{teiDiv}[type={section2},xmlid={div04-3}]

                \teiHead[xmlid={systeme-hexagonal-001}]{Système hexagonal}
                \teiP[xmlid={p16-7}]{Pour le système hexagonal, on ne peut s’empêcher de penser au Kinemax du Futuroscope de Poitiers, édifié en 1987 par Denis Laming. Ce bâtiment est l’exacte reproduction d’un ensemble de cristaux de quartz enfumés. L’empilage d’hexagones est une façon parfaite d’occuper l’espace en équilibrant les contraintes dans toutes les directions. Pour s’en assurer, il suffit d’observer la construction par les abeilles des alvéoles dans une ruche.}
                \teiP[xmlid={p17-7}]{Dans un autre domaine, la maison Chanel n’hésite pas, pour la présentation de sa collection automne-hiver au Grand Palais, à ériger un ensemble de cristaux de quartz tels qu’ils peuvent se présenter dans une géode. C’est toujours le principe de la géode qui a présidé à la conception de The Amethyst Hotel par NL Architect en Chine.}
              
\end{teiDiv}
              \begin{teiDiv}[type={section2},xmlid={div05-3}]

                \teiHead[xmlid={systeme-monoclinique-001}]{Système monoclinique}
                \teiP[xmlid={p18-7}]{Le système de cristallisation monoclinique est déjà beaucoup plus difficile à conceptualiser. C’est pourtant le parti pris qu’a adopté l’architecte Jean Nouvel en 2019 pour son ouvrage « rose des sables » qui abrite le Musée national du Qatar. La rose des sables est en effet l’une des formes de cristallisation du gypse qui présente des lames cristallines enchevêtrées. Tout le monde a à l’esprit ces cristaux relativement courants dans le désert saharien.}
              
\end{teiDiv}
              \begin{teiDiv}[type={section2},xmlid={div06-3}]

                \teiHead[xmlid={systeme-orthorhombique-001}]{Système orthorhombique}
                \teiP[xmlid={p19-7}]{Dernier exemple qui associe une forme cristalline d’un minéral à une construction architecturale, la \teiHi[rend={italic}]{Crystal Cathedral} (Philippe Johnson, 1977) à Garden Grove (États-Unis), dont le campanile semble être la reproduction géante d’un cristal de stibnite (sulfure d’antimoine) qui cristallise dans le système orthorhombique.}
              
\end{teiDiv}
            
\end{teiDiv}
            \begin{teiDiv}[type={section1},xmlid={div07-3}]

              \teiHead[xmlid={la-texture-et-la-couleur-001}]{La texture et la couleur}
              \teiP[xmlid={p20-7}]{Il serait présomptueux de faire un inventaire exhaustif des différentes textures des minéraux. Elles sont liées aux caractéristiques intrinsèques de ceux-ci, mais également aux modalités de leur formation et croissance.}
              \teiP[xmlid={p21-5}]{Les variations de couleurs sont infinies. Elles diffèrent souvent pour une même espèce minérale en fonction d’impuretés, d’inclusions, de substitutions ou de lacunes à l’échelle atomique, et de la déformation du réseau cristallin sous l’effet des contraintes. }
              \teiP[xmlid={p22-5}]{On peut s’arrêter à titre d’exemple sur l’agate. La richesse et la diversité des couleurs des cernes concentriques qui la caractérisent ont très fréquemment servi à la création de décors muraux ou de sols, mais aussi à celle d’objets, qu’ils soient ou non fabriqués à partir de ce minéral. Les puzzles proposés par Nervous System offrent un excellent exemple de ce que l’imagination créatrice peut produire à partir de ce minéral.}
              \begin{teiFigure}[xmlid={figure06-2}]

                \teiHead[xmlid={figure-13-orbicular-geode-puzzle-par-nervous-system-photographie-de-jessica-rosenkrantz-001}]{Figure 13. Orbicular Geode Puzzle, par Nervous System. Photographie de Jessica Rosenkrantz.}
              
\end{teiFigure}
              \teiP[xmlid={p23-5}]{Certaines espèces minéralogiques sont sélectionnées pour leur couleur afin de fabriquer des objets décoratifs. Parmi les plus classiques, on peut citer la malachite pour sa couleur verte, l’azurite pour le bleu, la rhodochrosite pour le rouge, la tourmaline ou le quartz pour le rose, le quartz améthyste pour le violet. Ces minéraux sont relativement courants et leurs formes cristallines suffisamment dimensionnées pour la réalisation d’objets de grandes dimensions comme des vases ou des statuettes par exemple. Les figures de croissances cristallines qui provoquent des nuances ou des variations de teintes sont fréquemment utilisées par l’artiste pour donner une singularité d’aspect à sa création.}
              \teiP[xmlid={p24-5}]{Dans l’histoire des civilisations, la transparence, la couleur et l’éclat des minéraux sont des critères qui ont très souvent attiré l’attention des femmes et des hommes. L’idée de parure et d’ornementation est présente dès la préhistoire, comme en témoignent les artéfacts découverts dans les fouilles archéologiques. }
              \teiP[xmlid={p25-4}]{Si les termes de pierre fine, de pierre précieuse, semi-précieuse et de gemme ont vu leurs définitions changer au cours des âges, ils regroupent cependant un certain nombre de minéraux dont les critères d’aspect satisfont lapidaires, joailliers et orfèvres.}
            
\end{teiDiv}
            \begin{teiDiv}[type={section1},xmlid={div08-3}]

              \teiHead[xmlid={symbolique-des-mineraux-001}]{Symbolique des minéraux}
              \teiP[xmlid={p26-4}]{Dans la plupart des cultures antiques, l’Homme a attribué aux minéraux et à leurs variétés de gemmes des vertus apotropaïques et des significations symboliques portées par leurs couleurs ou leurs formes. Les pouvoirs politiques et religieux se sont emparés de ces symboles pour asseoir leur domination. La beauté minérale a été mise au service de doctrines et de croyances. Leurs propriétaires se voyaient ainsi distingués, mais également détenteurs des pouvoirs supposés des pierres arborées. Aucune civilisation n’a échappé à ce mécanisme.}
              \begin{teiDiv}[type={section2},xmlid={div09-2}]

                \teiHead[xmlid={symbole-religieux-001}]{Symbole religieux}
                \teiP[xmlid={p27-4}]{Dans les civilisations orientales, le jade tient une place particulière. Il est associé aux rituels funèbres de la Chine ancienne. Les objets et costumes en jade étaient censés épargner au corps la putréfaction. Le jade était l’élément intermédiaire entre les humains et les dieux. Les noms mêmes des minéraux du jade, jadéite et néphrite, possèdent une origine liée aux pouvoirs miraculeux de celles-ci. Le jade néphrite, avec une composition minéralogique de la famille des amphiboles, tire son nom du grec \teiHi[rend={italic}]{néphros}, le rein, organe qu’il est censé soigner. La signification est identique (\teiHi[rend={italic}]{piedra de ijada} ou « pierre du flanc ») pour la jadéite, dont la composition minéralogique est celle d’un pyroxène.}
                \teiP[xmlid={p28-3}]{Les exemples de pierres aux pouvoirs divins peuvent se décliner dans toutes les époques et toutes les civilisations. Ces croyances sont encore très présentes et se retrouvent dans ce qu’on appelle aujourd’hui communément la « lithothérapie ».}
              
\end{teiDiv}
              \begin{teiDiv}[type={section2},xmlid={div10}]

                \teiHead[xmlid={pouvoir-politique-et-societe-001}]{Pouvoir politique et société}
                \teiP[xmlid={p29-3}]{Gemmes et pierres ornementales sont synonymes d’ostentation. On affiche sa puissance, sa richesse, son rayonnement en se parant de pierres précieuses et en construisant sa maison avec des matériaux rares et esthétiques. Diamants, rubis, émeraudes, marbres et porphyres ont là aussi traversé l’Histoire et les civilisations.}
                \teiP[xmlid={p30-3}]{Pour ne citer qu’un seul exemple, prenons celui de la vente des joyaux de la couronne de France en 1887. L’association de cette institution à la monarchie est telle que la III\teiHi[rend={sup}]{e} République, par un vote de l’assemblée, décide de sa destruction par le démantèlement de nombreux bijoux et la vente de cet ensemble patrimonial exceptionnel en le bradant. Seuls quelques objets et quelques gemmes sont épargnés de la vente et mis en dépôt au Louvre, au Muséum national d’histoire naturelle et à l’École des mines de Paris. }
                \begin{teiFigure}[xmlid={figure07}]

                  \teiHead[xmlid={figure-14-vitrine-des-joyaux-de-la-couronne-de-france-musee-de-mineralogie-mines-paris-psl-photographie-didier-nectoux-001}]{Figure 14. Vitrine des joyaux de la couronne de France. Musée de Minéralogie, Mines Paris – PSL, photographie Didier Nectoux.}
                
\end{teiFigure}
                \teiP[xmlid={p31-3}]{Le monde de la culture et du spectacle s’est lui aussi emparé de la beauté des pierres. On ne peut s’empêcher de penser aux rares et somptueuses parures portées par l’actrice Elizabeth Taylor, décrite par les médias comme une « icône de la beauté éternelle » (\teiHi[rend={italic}]{Paris Match}, 23 mars 2011). Les stars se sont appropriées les privilèges et le statut des reines et rois en se couvrant de pierres précieuses.}
              
\end{teiDiv}
              \begin{teiDiv}[type={section2},xmlid={div11}]

                \teiHead[xmlid={le-nom-des-mineraux-001}]{Le nom des minéraux}
                \teiP[xmlid={p32-3}]{La symbolique des minéraux s’étend à leur désignation. Les maisons de luxe donnent à leurs parfums et produits de beauté des noms de minéraux. Topaze, hématite, saphir, or, jade, azurite, diamant, émeraude, opale, améthyste, mica, et bien d’autres encore sont devenus des produits commerciaux permettant à celles et ceux qui les emploient de s’approprier les beautés, vertus et propriétés des gemmes et minéraux.}
                \teiP[xmlid={p33-3}]{Le quartz, synonyme d’éternité, de dureté, de transparence et de pureté (cristal de roche) est aujourd’hui le nom d’entreprises et de sociétés aussi variées qu’un laboratoire d’analyses, un bureau d’étude, une radio, un théâtre national, une agence immobilière, un horloger, etc.}
                \begin{teiFigure}[xmlid={figure08}]

                  \teiHead[xmlid={figure-15-logos-dorganismes-et-societes-portant-le-nom-du-mineral-quartz-musee-de-mineralogie-mines-paris-psl-composition-didier-nectoux-001}]{Figure 15. Logos d’organismes et sociétés portant le nom du minéral Quartz. Musée de Minéralogie, Mines Paris – PSL, composition Didier Nectoux.}
                
\end{teiFigure}
                \teiP[xmlid={p34-3}]{Dans la société d’aujourd’hui, traversée par des problématiques liées à l’environnement, mais aussi à l’éthique des conditions d’exploitation, les noms de certains minéraux revêtent une signification nouvelle. Le mica, par exemple, dont l’appellation a souvent été reprise par des sociétés de cosmétiques, symbolise la lutte contre l’exploitation des enfants employés par des entreprises productrices de cette matière première. Le coltan ou colombo-tantalite source de tantale, métal nécessaire à la fabrication des téléphones portables, s’est vu attribuer la désignation de « minerai de sang ».}
                \begin{teiFigure}[xmlid={figure09}]

                  \teiHead[xmlid={figure-16-panneau-coltan-musee-de-mineralogie-mines-paris-psl-photographie-cyrille-benhamou-omniscience-001}]{Figure 16. Panneau coltan. Musée de Minéralogie, Mines Paris– PSL, photographie Cyrille Benhamou OMNISCIENCE.}
                
\end{teiFigure}
                \teiP[xmlid={p35-3}]{La République démocratique du Congo, l’un des premiers pays producteurs et détenteurs de réserves pour ce minerai, mais également pour le cobalt, le tungstène et l’or, est ravagée depuis plus de dix années par des conflits guerriers meurtriers. Dans ce contexte, les noms de ces métaux et minerais peuvent parfois être associés à des destructions et massacres des êtres humains, des animaux et des systèmes écologiques. Les minéraux apparaissent alors comme des beautés fatales. }
                \teiP[xmlid={p36-3}]{Les minéraux, matières premières indispensables à toutes les activités humaines sont aussi quand ils cristallisent, des symboles de beauté. Formes et couleurs stimulent l’imagination de chacun. Les artistes et les créateurs s’en emparent pour produire des œuvres originales comme des produits de consommation. Il s’agit également de s’emparer de leurs propriétés et vertus. Cette appropriation va jusqu’à l’utilisation des noms et des imaginaires symboliques qui leur ont été attribués au cours des âges depuis l’aube de l’humanité. Comme les minéraux, la beauté minérale a été domestiquée, souvent instrumentalisée dans une recherche de rayonnement et de pouvoir. Raconter l’histoire des minéraux et de leurs utilisations est une façon singulière de raconter l’histoire des civilisations, faite de grandeur et de désastres\teiNote[xmlid={ftn1-8},n={1},place={foot}]{\teiP[xmlid={p-347}]{Voir également Cyrille Benhamou, \teiHi[rend={italic}]{Curiosités minérales}, Paris, Omniscience, 2017 ; Didier Nectoux et Éloïse Gaillou, \teiHi[rend={italic}]{Le musée de Minéralogie de l’École des mines de Paris}, Paris, Gallimard-L’École des arts joailliers, 2022. }}.}
              
\end{teiDiv}
            
\end{teiDiv}
          
\end{teiElement}
        
\end{teiElement}
        \begin{teiElement}[name={group},type={chapter},data-page-title={Poétique du jardin japonais : la villa détachée et les jardins de Katsura},data-page-authors={Philippe Bonnin@@CNRS-Université Paris-8, Laboratoire LAVUE – UMR 7218, France},data-include-href={Ch09\_poetique\_jardin\_japonais\_Bonnin.xml},xmlid={chapter-008}]

          \begin{teiElement}[name={front}]

            \begin{teiDiv}[type={titlePage},xmlid={titlepage-010}]

              \teiP[rend={title-main},xmlid={p-348}]{Poétique du jardin japonais : la villa détachée et les jardins de Katsura}
              \teiP[rend={author-aut},xmlid={p-349}]{Philippe \teiHi[rend={small-caps}]{Bonnin}}
              \teiP[rend={authority\_affiliation},xmlid={p-350}]{CNRS-Université Paris-8, Laboratoire LAVUE – UMR 7218, France}
            
\end{teiDiv}
          
\end{teiElement}
          \begin{teiElement}[name={body},xmlid={body-010}]

            \teiP[xmlid={p1-10}]{Quoi de commun entre la beauté que nous éprouvons et pensons en Occident depuis la Grèce antique, et celle d’un univers qui nous était totalement étranger, l’univers sino-japonais, surtout lorsqu’il fut fermé pendant des siècles à tout échange, tel le Japon d’Edo ? L’expérience de la beauté : nous n’en revenons pas ; notre pensée n’en revient pas indemne. Il faut sans doute réviser nos catégories de fond en comble, ce qui avait été commencé dès le japonisme, continué par l’aventure du paysage, mais qu’il faut poursuivre par l’épreuve de la vie quotidienne et l’expérience de la spatialité, au cœur de l’\teiHi[rend={italic}]{ethos} nippon. Il y faudra certainement réviser nos concepts, et l’idée même de conceptualisation ou de sa pratique. Et monter en généralité, dans une généralité sensible et expérientielle si l’on peut dire. Certes nous pensons déjà que le goût de la beauté s’éprouve dans une relation plus que dans un ou des états du sujet et de l’objet, en particulier dans une lecture des signes ou des dispositions que nous avons déposées, encodées dans les objets et dans le monde, ou que nous croyons y trouver. Y retrouver pour le plus grand soulagement d’une angoisse latente – vitale ou constitutive devrait-on dire – de perdition, d’abandon, de chaos ou d’anomie. C’est probablement cette intuition de mise en coïncidence de l’attente fébrile, inassouvie, inquiète, et du paysage visuel, auditif, tactile et situationnel qui génère la jubilation. De là, la conviction d’une intelligibilité, lorsque la correspondance entre registres du réel peut être entendue, ressentie, souvent au prix d’un long apprentissage des codages engravés. Mais on pourrait dire que le propre de la jubilation esthétique réside précisément dans l’instantané, le raccourci, le court-circuit, l’immédiateté de cette relation, par-delà le lent travail d’apprentissages et de déchiffrage qu’on oublie alors.}
            \teiP[xmlid={p2-10}]{On s’en persuadera peut-être à travers l’expérience de la beauté d’une œuvre qui use de procédés inconnus de notre univers culturel, dans le jardin japonais de Katsura. }
            \begin{teiDiv}[type={section1},xmlid={div01-7}]

              \teiHead[xmlid={le-jardin-japonais-ideal-acme-parangon-occidental-de-beaute-001}]{« Le » jardin japonais, idéal, \teiHi[rend={italic}]{acmè}, parangon occidental de beauté}
              \teiP[xmlid={p3-10}]{L’un des sommets des expériences de la beauté, telle qu’imaginée par les Occidentaux, est cette image qu’ils se font du jardin japonais, ou ce qu’ils en croient connaître. Et le zénith en est sans doute occupé par le jardin sec, singulièrement le jardin aux quinze pierres du temple Ryōan-ji (龍安寺), à Kyôto\teiNote[xmlid={ftn1-9},n={1},place={foot}]{\teiP[xmlid={p-351}]{On en trouvera une image dans Philippe Bonnin, \teiHi[rend={italic}]{Vocabulaire de la spatialité japonaise}, Paris, CNRS Éditions, 2014, planche 52.}}, suivi de ses semblables, ceux des temples Daisen-in, du Daitoku-ji et de multiples autres\teiNote[xmlid={ftn35-2},n={2},place={foot}]{\teiP[xmlid={p-352}]{Emmanuel Mares et Emmanuel Marès, \teiHi[rend={italic}]{Les concepteurs de jardins et de parcs japonais}, \teiHi[rend={italic}]{Projets de paysage}, dossier thématique, n\teiHi[rend={sup}]{o} 8, 2012, DOI : 10.4000/paysage.14776.}}, plus ou moins célèbres et réussis. Plus particulièrement ceux du temple Rinzai Myōshin-ji\teiNote[xmlid={ftn36-2},n={3},place={foot}]{\teiP[xmlid={p-353}]{Temple beaucoup moins fréquenté, quoique plus impressionnant encore et plus abstrait, mais dont Augustin Berque dit quelques mots dans \teiHi[rend={italic}]{Poétique de la terre}, Paris, Belin, 2014 : temple dont le « jardin » de méditation principal est formé d’un grand rectangle totalement nu et vide, n’étaient les longues stries rectilignes d’un sol de gravier (voir Philippe Bonnin, \teiHi[rend={italic}]{Vocabulaire…}, \teiHi[rend={italic}]{op. cit}., photographies p. 227 et 282 et planche 55).}} : le second de ses trois jardins, le plus petit, intérieur, dispose sept pierres agencées autour de l’axe médian, flottant sur une mer d’ondes concentriques, et ne le rend en rien au Ryōan-ji.}
              \teiP[xmlid={p4-10}]{Ces jardins paradoxaux, pour peu végétalisés qu’ils soient, contrairement à la notion européenne de jardin\teiNote[xmlid={ftn37-2},n={4},place={foot}]{\teiP[xmlid={p-354}]{Encore qu’il existe une version inversée du Ryôanji au Shôden-ji, à Kyoto : les quinze pierres y sont remplacées par autant de massifs végétaux. Voir Emmanuel Mares, \teiHi[rend={italic}]{Kare sansui}, dans Philippe Bonnin, \teiHi[rend={italic}]{Vocabulaire…}, \teiHi[rend={italic}]{op. cit}., p. 227 et photographie planche 53.}}, sont plus proches en cela de la sculpture de pierres et de graviers, si ces pierres avaient été taillées et non pas seulement choisies et disposées, ou de la catégorie plus récente du \teiHi[rend={italic}]{land-art}. Mais ils ne sont en rien l’alpha et l’oméga du jardin sino-japonais. Ils n’en sont qu’un épisode très particulier dans une longue histoire.}
            
\end{teiDiv}
            \begin{teiDiv}[type={section1},xmlid={div02-7}]

              \teiHead[xmlid={l-ideal-du-jardin-pour-les-japonais-001}]{L'idéal du jardin pour les Japonais}
              \teiP[xmlid={p5-10}]{En fait, pour les Japonais, l’idéal de beauté du jardin fut toujours, depuis le \teiHi[rend={small-caps}]{xvii}\teiHi[rend={sup}]{e} siècle, le jardin de la villa retirée Katsura Rikyû (villa impériale détachée, sur le bord de la rivière Katsura). La villa fut édifiée par un prince de la lignée impériale, reprise par son fils, puis transmise de génération en génération par cette branche parallèle de la famille impériale, qui prit le nom de Katsura.}
              \teiP[xmlid={p6-10}]{Mais quelle serait donc cette sorte de beauté différente, exotique, fascinante, formulée dans un code qui nous est étranger ? C’est ce qu’il s’agirait d’esquisser ici. Mission presque impossible il est vrai, tant la tâche est complexe, et tant se lèvent d’obstacles : comment pourrait-il se loger une beauté durable dans un jardin fait de végétaux labiles, à la forme transformée d’une année sur l’autre, chaque printemps, et même à chaque saison ? Quand la beauté s’approche ainsi d’une situation, d’une relation esthétique au lieu, comment la saisir sans l’éprouver par le corps, par tous nos sens ? Pour parler d’un jardin, il faut d’abord le voir, le sentir, le parcourir. Au moins, le décrire, en partager la découverte, le parcours. Lequel est bien trop riche, bien trop long et complexe pour ces quelques lignes\teiNote[xmlid={ftn38-2},n={5},place={foot}]{\teiP[xmlid={p-355}]{Qu’on nous permette ici de renvoyer le lecteur à la tentative pour satisfaire cette exigence, dans \teiHi[rend={italic}]{Éloge d'un jardin japonais}, réédition de \teiHi[rend={italic}]{Katsura et ses jardins} : Philippe Bonnin, \teiHi[rend={italic}]{Éloge d’un jardin japonais}, Paris, Arléa, 2022, sur laquelle s’appuie cette étude.}}. Décrire donc, d’abord et avant tout.}
              \teiP[xmlid={p7-10}]{En effet, contrairement aux apparences, on ne saurait parler des choses japonaises en termes trop abstraits, malgré la sensation de condensation de sens qu’elles nous procurent. Les Japonais y répugnent, se méfient des généralités, de la théorisation. Ils préfèrent le silence, la simplicité de quelques mots peignant la réalité, un énoncé plus modeste. Ils savent y lire le suave et le profond. Ils savent prendre le temps d’un recul intérieur, d’une suspension, d’une attente réceptive sur ce que les choses nous font, sur ce qu’elles ébranlent en nous, la poignante émotion des choses.}
              \teiP[xmlid={p8-10}]{Mais au moins faut-il tenter de dépasser le simple étonnement fasciné, l’émerveillement qui coupe le souffle, laisse muet après trop peu de temps accordé devant ces beautés flagrantes, à peine entraperçues. Quelques arrêts sur image, pris arbitrairement dans le chapelet de lieux et de spectacles qui s’égrène infiniment, quelques essais de redéploiement des paysages si serrés, nous permettront peut-être d’apercevoir les codes, les procédés, les dispositifs de prédilection de ce jardin et de ses concepteurs.}
              \teiP[xmlid={p9-9}]{Nous tenterons d’égrener un chapelet de lieux qui ne valent pas seulement pour eux-mêmes, quoiqu’ils soient agencés avec grand art, grand soin et grande subtilité, mais simultanément comme évocations, allusions à d’autres lieux célèbres, réels ou mythiques, comme réminiscences, comme citations historiques ou poétiques. Ces lieux sont des récits, et des affirmations de valeurs, de choix esthétiques dûment élaborés et discutés, tels que les références à la simplicité rustique, à la nature brute plutôt qu’élaborée.}
              \teiP[xmlid={p10-9}]{Ce sont en même temps des lieux propices aux jeux, aux joutes poétiques, aux récréations, aux divertissements musicaux, à la délectation du moment et de l’heure, des lumières, des floraisons, des saisons, de l’épanouissement des cerisiers et de la lune avec leurs rites propres transmis au cours des siècles.}
              \teiP[xmlid={p11-9}]{Ces lieux déploient un large éventail de qualités, par leurs dispositions propres : en montée escarpée à gravir lentement, en sombre et humide forêt, en large prairie ouverte, en rivage déchiré ou en sentier délicat, en surplomb d’un fossé, ou appelant le geste, le pas ou l’arrêt, qui fait déjà corps avec nous. Et chaque fois, ce sont des agencements savants et soignés de minéraux et de végétaux choisis, mises en forme pensées selon des dispositions espérées, des lumières et des ambiances conjuguées aux heures et aux saisons, aux intempéries, à l’éclatement des bourgeons et aux vols des libellules, des lucioles, des papillons noirs, aux bruissements des oiseaux, des carpes, du vent, formant des espaces expressifs, eux-mêmes matière du paysage, composant le moment, accueillant et organisant la présence, ici et maintenant.}
            
\end{teiDiv}
            \begin{teiDiv}[type={section1},xmlid={div03-6}]

              \teiHead[xmlid={lhistoire-de-lhistoire-dune-histoire-001}]{L’histoire de l’histoire d’une histoire}
              \teiP[xmlid={p12-9}]{C’est bien l’histoire d’une ancienne réécriture jardinière, et réécriture du roman médiéval \teiHi[rend={italic}]{Genji Monogatari}, qu’il s’agit d’écouter si l’on veut pleinement apprécier ce jardin, comme il en serait de tout autre. Peut-être un peu plus ici. }
              \teiP[xmlid={p13-8}]{En janvier 1579 naît Kosamaru, le sixième fils du prince impérial Sanehito, fils aîné de l’empereur Ōgimachi (106\teiHi[rend={sup}]{e} empereur, 1517-1593). Le généralissime Toyotomi Hideyoshi, n’ayant pas encore de fils, adopte Kosamaru en 1586 – jusqu’en 1589, quand lui vient enfin un héritier. Il lui octroie le nom de Toshihito et, après divers changements, cinq villages sur les deux rives de la rivière Katsura. Le prince Toshihito, dont la résidence est à Kyoto, fait élever à Katsura un premier pavillon, le Koshoin, en 1615, sur l’emplacement d’un ancien « palais détaché », villégiature aristocratique d’un lointain passé. Il fait recreuser l’ancien étang de style \teiHi[rend={italic}]{shinden}, reformer les îles, etc.}
              \teiP[xmlid={p14-8}]{Le prince Toshihito enseigne à l’empereur régnant successeur, Go-Mizunoo (108\teiHi[rend={sup}]{e} empereur, 1611-1629), l’\teiHi[rend={italic}]{Anthologie de la poésie ancienne et moderne}, le \teiHi[rend={italic}]{Kokin wakashū}, dont il est devenu fin connaisseur, et reconnu en cela par son maître Hosokawa Yūsai qui lui en délivre le titre, en 1600, à l’âge de vingt ans.}
              \teiP[xmlid={p15-8}]{La tradition retient que le 20 juillet 1616, le prince Toshihito fit une excursion le long de la rivière Katsura pour « y voir ses melons\teiNote[xmlid={ftn39-2},n={6},place={foot}]{\teiP[xmlid={p-356}]{Toshihito le désigne dans une lettre ultérieure (1618) comme le « pavillon de thé dans le plant de melons de la basse Katsura ».}} »\teiHi[rend={sup}]{,} accompagné d’une joyeuse troupe, et que deux jours plus tard, Konoe Sakiko, l’épouse de l’empereur retiré Go-Yōsei (107\teiHi[rend={sup}]{e} empereur de 1586 à 1611, petit-fils de Ōgimachi, et donc frère aîné de Toshihito), visitant le temple Senshō-ji, se rendit également au « pavillon de thé dans le plant de melons de la basse Katsura ». Le pavillon se devait d’être suffisamment aménagé pour recevoir cet hôte de marque.}
              \teiP[xmlid={p16-8}]{Toshihito fut probablement l’artisan du mariage de Tokugawa Kazuko avec l’empereur Go-mizunoo (1596-1680) – c’est-à-dire l’alliance de la famille des \teiHi[rend={italic}]{shōguns} Tokugawa et de la lignée impériale –, célébré le 17 juillet 1620. Prenant désormais une place importante auprès du \teiHi[rend={italic}]{shōgunat}, Toshihito devint une figure publique, et la villa le lieu de visite de nombre d’hôtes. Il en fut ainsi jusqu’à sa mort, en 1629, à l’âge de cinquante ans. }
              \teiP[xmlid={p17-8}]{En juin 1624, la description qu’a donnée de Katsura le moine Kinshuku Kentaku (1580-1658) du Shōkoku-ji, qui y fut invité, décrit collines, lac, îles et barques : « J’ai visité la maison de campagne du prince Hachijō, dans la basse Katsura. Les bateaux voguaient sur l’étang creusé dans le jardin, passant sous les ponts qui reliaient les îles. Il y a une colline artificielle et un pavillon d’été dans le jardin. De ce pavillon, on voit les montagnes, au-delà des jardins, de tous côtés, c’est le plus joli paysage du Japon. »}
              \teiP[xmlid={p18-8}]{L’automne de l’année suivante 1625, c’est le puissant supérieur Ishin Sūden (1569-1633) du Konchi-in, qui relate ses \teiHi[rend={italic}]{Souvenirs du pavillon de Katsura} : le prince a « construit une belle maison de campagne dans les jardins de laquelle coule l’eau des montagnes ». Sur une grande barque richement décorée, on donnait les banquets, à l’imitation des descriptions de la résidence du puissant Fujiwara no Michinaga (966-1027), ministre de la Gauche, puis ministre des Affaires suprêmes et modèle du personnage du prince dans le \teiHi[rend={italic}]{Genji monogatari}\teiNote[xmlid={ftn40-2},n={7},place={foot}]{\teiP[xmlid={p-357}]{Héros de l’antique roman du \teiHi[rend={italic}]{Genji Monogatari}, \teiHi[rend={small-caps}]{xi}\teiHi[rend={sup}]{e} siècle, attribué à Murasaki Shikibu.}}, dont Toshihito avait recopié de sa main tous les passages qui concernaient Katsura, et qui possédait une villa sur les bords de la rivière Katsura, en ce lieu même.}
              \teiP[xmlid={p19-8}]{À la mort de Toshihito, à l’âge de 51 ans, douze années d’abandon voient le vieux \teiHi[rend={italic}]{shoin} se détériorer considérablement. Les catastrophes naturelles sont fréquentes au Japon : séismes, inondations… et c’est l’un de ces cataclysmes qui endommage gravement la villa. Au point qu’en 1631, le moine Kinshuku Kentaku du Shōkoku-ji est désolé par l’état de délabrement, et note : « Rien n’a été remis en état depuis la mort de Keiko-in \lbrack{}le prince Toshihito\rbrack{}, et tout s’est dégradé. » }
              \teiP[xmlid={p20-8}]{Il faut attendre 1641, et le mariage de son fils Toshitada (1619-1662), élevé dans l’étude de la poésie et des classiques chinois, y dépassant même son père, pour que commencent de nouveaux travaux dans les jardins, qui font sans cesse référence à la culture poétique.}
              \teiP[xmlid={p21-6}]{Les revenus issus de son alliance avec le clan Kaga vont permettre à Toshitada de rénover la villa. Le prince était fasciné par le chapitre du \teiHi[rend={italic}]{Dit du Genji} qui décrivait un jardin aux fleurs et senteurs merveilleuses, à l’île centrale qu’on accostait dans des barques pilotées par de jeunes gens en costumes chinois, dans les reflets mauves des glycines et des corètes jaunes de montagne. C’est cette image qu’il s’attacha à rendre à Katsura. Et, quand les travaux de restauration furent accomplis, les pierres redressées, les plantations remises en ordre, redessiné le cours du ruisselet qui murmurait de nouveau, leur vinrent des poèmes nostalgiques. Ils firent tout comme dans les descriptions anciennes des festivités qui nous sont parvenues : « Là, de nouveaux rafraîchissements lui furent servis et l’on récita des poèmes. Lorsque la barque revint au rivage, la compagnie s’assit autour d’une table basse pour réciter des poèmes en répons. »}
              \teiP[xmlid={p22-6}]{Il y intègre progressivement les formes du style \teiHi[rend={italic}]{sukiya}, plus libre, utilisé pour la construction des pavillons de thé de type rustique \teiHi[rend={italic}]{sōan} (style des huttes de chaume), face à l’architecture \teiHi[rend={italic}]{shoin}, celle des cabinets des lettrés. Le Koshoin est le plus ancien témoignage du mariage de ces styles, que l’on qualifiera de \teiHi[rend={italic}]{kirei zashiki }ou « séjour de beauté », en un « beau \teiHi[rend={italic}]{sabi} » : le \teiHi[rend={italic}]{kirei sabi}, le raffinement et la distinction dans la retenue. Y domine l’inspiration du maître de thé Kobori Enshū, associant l’élégance de la maison impériale avec l’austérité rustique que prônait Sen no Rikyū. Les références aux chaumières et au dénuement rural des pavillons de thé y sont surtout une posture poétique, littéraire, philosophique, savante : une référence au retrait érémitique des sages chinois anciens, \teiHi[rend={italic}]{inkyo}.}
              \teiP[xmlid={p23-6}]{Ainsi les références et les allusions qui plongeaient déjà dans plusieurs siècles de culture traversent les dispositions des jardins, amplifiant la délectation des lieux, par-delà leur beauté intrinsèque. Nous ne pourrons, dans ces courtes pages, suivre avec fidélité et de manière exhaustive le parcours qui fonde le jardin, faisant le tour de son étang central comme il était déjà de coutume au \teiHi[rend={small-caps}]{xvii}\teiHi[rend={sup}]{e} siècle, et sur le même chemin qui demeure tracé par les pierres disposées à l’époque. Pourtant, pas un recoin qui ne ménage des surprises et des émerveillements, tant furent concentrés l’inventivité et le soin du ou des concepteurs\teiNote[xmlid={ftn41-3},n={8},place={foot}]{\teiP[xmlid={p-358}]{Nous n’entrerons pas ici dans la discussion sur la paternité de la conception, à laquelle nous avons déjà consacré quelques pages : « Le mythe de Kobori Enshû », dans : \teiHi[rend={italic}]{Éloge d’un jardin japonais}, \teiHi[rend={italic}]{op. cit}, p. 252.}}. Nous n’en retiendrons que quelques stations, quelques figures, quelques dispositifs ou procédés des jardiniers japonais pour accéder à la beauté.}
              \begin{teiFigure}[xmlid={figure01-4}]

                \teiHead[xmlid={figure-17-katsura-no-miya-go-besso-zenzu-anonyme-releve-aquarelle-des-jardins-de-la-villa-katsura-vers-1688-bibliotheque-nationale-de-la-diete-tokyo-001}]{Figure 17. \teiHi[rend={italic}]{Katsura no miya go-besso zenzu,} anonyme, relevé aquarellé des jardins de la villa Katsura, vers 1688. Bibliothèque nationale de la Diète, Tokyo.}
              
\end{teiFigure}
            
\end{teiDiv}
            \begin{teiDiv}[type={section1},xmlid={div04-4}]

              \teiHead[xmlid={un-impossible-mur-de-feuilles-la-haie-katsura-gaki-001}]{Un impossible mur de feuilles : la haie \teiHi[rend={italic}]{Katsura gaki}  }
              \teiP[xmlid={p24-6}]{C’est ainsi qu’est accueilli le visiteur qui vient de la capitale, ayant traversé à gué la rivière Katsura. Non pas un rempart comme ceux qui protègent le palais impérial, non pas des murs d’énormes roches ou crépis de plâtre blanc comme les châteaux forts, pas même des murs d’adobe aux lits de vieilles tuiles des puissants temples, mais un écran fait de simples feuilles de bambou, bien qu’on ne voie ni les tiges ni les cimes ou les épis de ceux-ci. Étrange haie en vérité, qui déjoue l’entendement et intrigue le regard attentif, dès le départ. Haie unique en son genre, dans tout le pays.}
              \teiP[xmlid={p25-5}]{Par un savant artifice, les jardiniers ont ployé les jeunes bambous qui poussent à l’intérieur du domaine, le long de la haie, faisant courber la tête et tissant étroitement les tiges les unes aux autres en un rideau de feuillage, velours vert sans trame ni chaîne.}
              \teiP[xmlid={p26-5}]{De beaux sujets, des ormes \teiHi[rend={italic}]{keyaki}, apportent leur ombre à la scène, tout au long du parcours. Assurément par la volonté du maître, mais sans qu’ils obéissent à un alignement rectiligne. Et tout au contraire leur implantation vaticine, tantôt hors de la limite, tantôt au-dedans, tantôt au milieu, comme d’une liberté espiègle, se jouant des règles qu’un pouvoir supérieur aurait voulu imposer. Et il en était bien un en ces temps.}
              \teiP[xmlid={p27-5}]{Un mur de feuilles pour défendre un prince de sang ? Singulière confiance, tout de même.}
              \teiP[xmlid={p28-4}]{Pour continuer la clôture, c’est une palissade de bambous secs, \teiHi[rend={italic}]{O-gaki}, aux fûts verticaux tissés non de planches, mais de simples ramilles, le tout serré de croix en corde rêche, noire, qui ferme le domaine. Plus qu’à une clôture défensive, celle-ci fait penser à toutes ces déclinaisons de palissades de bambou développées par chaque temple et portant son propre nom, selon que les bambous sont par paires ou isolés, refendus ou pleins, horizontaux ou verticaux, et toutes les infimes variations que l'on peut imaginer\teiNote[xmlid={ftn42-3},n={9},place={foot}]{\teiP[xmlid={p-359}]{Voir la troisième partie de \teiHi[rend={italic}]{La beauté du seuil}, consacrée aux clôtures, particulièrement p. 185-189 et la liste infinie des appellations dans Philippe Bonnin, \teiHi[rend={italic}]{La beauté du seuil : esthétique japonaise de la limite}, Paris, CNRS Éditions, 2021, p. 145-189. La première partie de l’ouvrage est consacrée à la notion de \teiHi[rend={italic}]{kekkai}.}}. Pas plus de murailles, là encore. C’est bien d’un \teiHi[rend={italic}]{kekkai} dont il s’agit : une limite marquée davantage pour que l’esprit s’en imprègne et s’en convainque que pour diviser absolument l’espace.}
            
\end{teiDiv}
            \begin{teiDiv}[type={section1},xmlid={div05-4}]

              \teiHead[xmlid={entrer-dans-le-jardin-001}]{Entrer dans le jardin}
              \teiP[xmlid={p29-4}]{Les étapes suivantes de l’entrée se développent dans le même esprit. La porte d’accès ouverte donne vue sur un écran végétal, comme un mur bahut en angle droit, fait de camélias taillés à la perfection. Dans toute entrée de temple, dans toute demeure traditionnelle, un écran peint, tissé ou une somptueuse tranche d’arbre géant se dressent face au regard du visiteur : il serait malséant, sinon malveillant, de pénétrer trop vite dans la demeure. Seuls les mauvais esprits, l’œil diabolique, parce qu’il ne sait qu’aller tout droit, sont arrêtés par ce dispositif du \teiHi[rend={italic}]{genkan}.}
              \teiP[xmlid={p30-4}]{Ce qui frappera l’esprit, quelques pas plus loin, ayant foulé du pied les galets violines d’aragonite qui tapissent la chaussée, c’est un coude qui forme un dispositif étonnant : un embranchement est offert, dirigé vers le centre du domaine, encadré de deux haies de camélias taillées à hauteur d’homme pour conduire la vue, mais formant un couloir bloqué par un pin, comme si l’on avait voulu que le regard ne porte pas, ou pas distinctement, qu’il ne puisse qu’apercevoir des bribes, deviner l’étendue du jardin, imaginer et espérer ce qu’il visitera bientôt, s’y préparer, sans en épuiser d’avance la teneur. C’est un procédé classique, le \teiHi[rend={italic}]{miegakure}, une sorte de « montré-caché » qui amplifie la beauté du paysage en filigrane.}
              \teiP[xmlid={p31-4}]{Une pierre d’arrêt, ronde et nouée en croix d’une corde noire, est disposée à la limite du pavage, juste avant l’herbe du promontoire en presque-île herbue sur lequel s’élève le pin. Parfois c’est plus explicitement une grande perche de bambou montée et nouée sur deux courts pieds de travers, qui est disposée en arrêt. Les deux sont aussi des signes de \teiHi[rend={italic}]{kekkai}, encore plus dénués de coercition : « veuillez ne pas aller au-delà ».}
              \teiP[xmlid={p32-4}]{L’esprit serait déjà comblé, pourrait se satisfaire de ces premières subtilités. Mais il nous faut dire qu’un nom est donné à ce pin : \teiHi[rend={italic}]{Sumiyoshi no matsu}. Ce pin n’est pas seulement un arbre planté ici pour faire écran. Son nom renvoie à tout l’univers de contes, de poèmes et du théâtre nô qui surgissent à l’esprit dès qu’on l’évoque. }
              \teiP[xmlid={p33-4}]{\teiHi[rend={italic}]{Sumiyoshi monogatari }est un vieux conte, anonyme, de la période de Kamakura (1185-1333), peut-être ancré dans un fonds plus ancien encore. Il ressemble à notre « Cendrillon » : le prince héros du conte recherche sa bien-aimée, écartée par une marâtre, et qui a dû fuir dans un monastère. Il est au bord des flots, au sanctuaire Sumiyoshi (près de Sakai, Osaka maintenant), vénérant le grand dieu marin, protecteur des navigateurs : Sumiyoshi Myōjin est un \teiHi[rend={italic}]{kami} en trois personnes (3 frères), fils de Susanoo-no-Mikoto : les \teiHi[rend={italic}]{kamis} de la mer, des poissons et des coquillages. Les deux pins inséparables – Sumiyoshi et Takasago – tourmentés par les vents et les vagues, sur cette grève sableuse et désolée \teiHi[rend={italic}]{suhama}, métaphorisent son état d’âme. Les canards figurés sur l’estampe, symboles de fidélité des sentiments, les apparitions qu’il a eues en rêve et qui montrent que la bien-aimée pense à lui, que son tourment est partagé, lui font poursuivre sa quête. Ce moment du conte est rapporté dans le poème Tanka : }
              \begin{teiLg}

                \teiL{De Suminoé (autre nom de Sumiyoshi)}
                \teiL{Vague qui bat le rivage}
                \teiL{Fût-ce en pleine nuit}
                \teiL{Et sur la route des songes}
                \teiL{Fuiriez-vous donc les regards\teiNote[xmlid={ftn43-3},n={10},place={foot}]{\teiP[xmlid={p-360}]{\teiHi[rend={italic}]{Fujiwara no Toshiyuki no Ason}, dans \teiHi[rend={italic}]{De cent poètes un poème}, René Sieffert (trad.), Aurillac, POF, 1993, poème n\teiHi[rend={sup}]{o} 18, p. 42.}}}
              
\end{teiLg}
              \begin{teiFigure}[xmlid={figure02-3}]

                \teiHead[xmlid={figure-18-fujiwara-no-toshiyuki-ason-n-o-18-hyakunin-isshu-no-uchi-utagawa-kuniyoshi-vers-1842-gravure-sur-bois-oban-tate-e-fujiwara-no-toshiyuki-avec-un-page-et-un-serviteur-contemplant-le-pont-a-tambours-du-temple-sumiyoshi-parmi-les-pins-dans-une-bande-de-brume-the-trustees-of-the-british-museum-001}]{Figure 18. \teiHi[rend={italic}]{Fujiwara no Toshiyuki Ason }(n\teiHi[rend={sup}]{o} 18), \teiHi[rend={italic}]{Hyakunin isshu no uchi}, Utagawa Kuniyoshi, vers 1842. Gravure sur bois, oban tate-e. Fujiwara no Toshiyuki avec un page et un serviteur, contemplant le pont à tambours du temple Sumiyoshi parmi les pins dans une bande de brume. © The Trustees of the British Museum.}
              
\end{teiFigure}
              \teiP[xmlid={p34-4}]{Le thème du conte, maintes et maintes fois représenté, dans un nô de Zeami, et repris dans la littérature, dans les \teiHi[rend={italic}]{Notes de chevet }de Sei Shônagon et le \teiHi[rend={italic}]{Dit du Genji }de Murasaki Shikibu – au livre quatorzième, \teiHi[rend={italic}]{Miwotsukushi}, par exemple –, fait évidemment partie du bagage littéraire des princes Hachijō no miya : « De Takasago, navigant sur la baie, la lune sort avec la marée, devant la silhouette de l’île d’Awaji, loin sur la mer jusqu’à Naruo, arrivée à Suminoe. » Le prince Genji, qui en pince pour la dame d’Akashi, reprend un célèbre poème :}
              \begin{teiLg}

                \teiL{Depuis que je vis}
                \teiL{Voilà bien longtemps déjà}
                \teiL{à Sumiyoshi}
                \teiL{sur le rivage}
                \teiL{le jeune pin\teiNote[xmlid={ftn44-3},n={11},place={foot}]{\teiP[xmlid={p-361}]{\teiHi[rend={italic}]{Nô et Kyôgen. Théâtre du Moyen Âge}, René Sieffert (présent. et trad.), Aurillac, POF, 1979, 2 t., vol. 1, p. 44.}}.}
              
\end{teiLg}
              \teiP[xmlid={p35-4}]{Le poème décrit Takasago et Suminoe, les pins mariés qui, selon la légende, resteront ensemble pour l’éternité. }
              \teiP[xmlid={p36-4}]{Ces deux pins étaient plantés dès l’origine du jardin, de part et d’autre de l’étang : « Depuis les barques, ils admirèrent les pins Sumiyoshi et Takasago, ces deux arbres qui inspirèrent tant de poètes. Enfin ils grimpèrent jusqu’au pavillon de thé pour voir la lune se refléter dans l’étang\teiNote[xmlid={ftn45-3},n={12},place={foot}]{\teiP[xmlid={p-362}]{Récit de visites, sous le septième chef de la famille Hachijo no Miya.}}. »}
              \teiP[xmlid={p37-3}]{On les trouve figurés sur le plan aquarellé du \teiHi[rend={small-caps}]{xvii}\teiHi[rend={sup}]{e} siècle, plus grands qu’aujourd’hui – combien de fois auront-ils été remplacés depuis ?}
              \teiP[xmlid={p38}]{Il faudrait évoquer d’autres pensées qui nous traversent au long du parcours, et qui renvoient à des concepts majeurs de l’esthétique japonaise : les marques du temps, \teiHi[rend={italic}]{sabishi}, que montre la couverture végétale d’un faîte de mur, et qu’on s’est volontairement abstenu de rénover trop rapidement – on l’apprécie telle qu’elle ; le coup d’œil, à la dérobée,\teiHi[rend={italic}]{ Nozoku}, que l’on jette pour apercevoir un détail, un lieu écarté ; la transparence de certaines haies vives, \teiHi[rend={italic}]{Kakine}, diaphanes, savamment entretenues pour délimiter l’endroit tout en lui attachant ses alentours ; encore une sorte de \teiHi[rend={italic}]{kekkai}, qu’Itoh Teiji, le grand historien de l’architecture, détaille à loisir et dont les traités des jardins décrivent le soin jaloux qu’on leur porte, l’ordre des éclaircissements et des tailles de structure qu’il faut respecter pour atteindre au résultat parfait.}
            
\end{teiDiv}
            \begin{teiDiv}[type={section1},xmlid={div06-4}]

              \teiHead[xmlid={momijiyama-la-montagne-aux-erables-001}]{\teiHi[rend={italic}]{Momijiyama, }la montagne aux érables}
              \teiP[xmlid={p39}]{Mais arrêtons-nous à ce monticule, que l’on contournerait si facilement sans y prendre garde. Il a pour nom \teiHi[rend={italic}]{Momijiyama Kōyō san,} « la montagne aux érables rouges ». Cette si haute éminence, qui culmine à 12 m 80, est couverte de jeunes érables aux feuilles délicatement étoilées, promesses d’un incendie de couleurs pour l’automne prochain. Là aussi, le jardin se réfère à une antique et plaisante légende. Le nom, \teiHi[rend={italic}]{Momijiyama}, provient en fait d’un autre poème chinois de Bai Juyi (772-846), rappelé dans \teiHi[rend={italic}]{Le Dit des Heike}. Dans les années de Jūan (vers 1171-1175), l’empereur Takakura « avait une vive prédilection pour le feuillage rutilant de l’automne, si bien qu’il avait \lbrack{}…\rbrack{} fait construire une colline sur laquelle il avait fait planter \lbrack{}…\rbrack{} des érables \lbrack{}…\rbrack{} ; il l’avait appelée la Colline du rouge feuillage\teiNote[xmlid={ftn46-3},n={13},place={foot}]{\teiP[xmlid={p-363}]{ Poème chinois de Bai Juyi (772-846), rappelé dans\teiHi[rend={italic}]{ Le Dit des Heike}, René Seffert (trad.), Aurillac, POF, 1988, p. 249.}} ». Or, une nuit, le vent éparpille toutes les feuilles. Au matin, les valets les ramassent toutes, et comme il fait frais, les brûlent pour réchauffer leur saké. Le secrétaire particulier est atterré : « Pour vous, ce peut être la prison ou l’exil sur l’heure, et moi aussi je vais encourir sa colère\teiNote[xmlid={ftn47-3},n={14},place={foot}]{\teiP[xmlid={p-364}]{ \teiHi[rend={italic}]{Ibid.}}}. » Sur ce, l’empereur arrive et les questionne. De bonne humeur, il rit et récite :}
              \begin{teiLg}

                \teiL{Au fond des bois pour chauffer mon vin}
                \teiL{Le rutilant feuillage j’ai brûlé\teiNote[xmlid={ftn48-3},n={15},place={foot}]{\teiP[xmlid={p-365}]{ \teiHi[rend={italic}]{Ibid.}}}.}
              
\end{teiLg}
              \teiP[xmlid={p40}]{« Qui donc a pu leur enseigner ces vers ? Ce qu’ils ont fait là est du meilleur goût\teiNote[xmlid={ftn49-3},n={16},place={foot}]{\teiP[xmlid={p-366}]{ \teiHi[rend={italic}]{Ibid.}}} ! », leur prodiguant des éloges quand ils attendaient de terribles reproches.}
            
\end{teiDiv}
            \begin{teiDiv}[type={section1},xmlid={div07-4}]

              \teiHead[xmlid={le-chemin-de-rosee-001}]{Le « chemin de rosée »}
              \teiP[xmlid={p41}]{Il faut le dire clairement, quoiqu’aucun ouvrage consacré à Katsura ne semble en prendre pleinement conscience : il n’y a pas UN jardin de Katsura, mais DES jardins, un assemblage, composé comme une histoire des jardins japonais. Jardins qui s’interpénètrent et se superposent même parfois : un antique jardin-paradis des anciennes résidences aristocratiques, \teiHi[rend={italic}]{shinden-zukuri}, un jardin de thé, \teiHi[rend={italic}]{roji}, un jardin de promenade, \teiHi[rend={italic}]{kaiyūshiki-teien}, et un jardin sec, \teiHi[rend={italic}]{kare-sansui}. Toutes ces formes historiques sont ici présentes, avec leurs vocabulaires propres, leurs procédés et leurs ruses, mêlées à des évocations de paysages maritimes, de sentiers de montagne, de huttes d’ermites, etc. Mais surtout, ce « jardin de thé » est un chemin savamment construit qui mène au pavillon où sera partagé le divin breuvage : le \teiHi[rend={italic}]{roji} ou « chemin de rosée ». L’appellation fait allusion au sutra qui décrit la voie où renaît celui qui a échappé aux désirs matériels.}
              \teiP[xmlid={p42}]{L’entrée dans ce « jardin de thé » était marquée autrefois (le plan aquarellé et les descriptions historiques en attestent) par une porte très légère, diaphane, quasi symbolique, aujourd’hui disparue\teiNote[xmlid={ftn50-3},n={17},place={foot}]{\teiP[xmlid={p-367}]{On peut s’en faire une idée par les photographies des planches 42, 53 et 54 dans Philippe Bonnin, \teiHi[rend={italic}]{La beauté du seuil}, \teiHi[rend={italic}]{op. cit.}}}. Encore un archétype de \teiHi[rend={italic}]{kekkai} : faire comprendre que l’on franchit un seuil, que l’on entre dans un autre monde, et en construire les règles dans son esprit.}
              \teiP[xmlid={p43}]{Sur ce seuil le sol même change : on quitte la large allée pavée pour un étroit sentier, creux et sinueux, où sont offerts les pas calculés des pierres plates, \teiHi[rend={italic}]{tobiishi} (littéralement « pierres volantes », « sauts »), arrangées comme le pas du pluvier balançant de droite et de gauche. À pas comptés, rythmés, toutes les trois à quatre pierres, une station est marquée par une aire plus large, pour un arrêt, une contemplation du paysage renouvelé : une lanterne de pierre, basse, au bonnet carré, qui éclaire le chemin et la matière même des dalles, veinées de blanc et de violine ou de vert jade. Chaque pierre a été savamment choisie dans la montagne, accueillie en cérémonie nuitamment au jardin pour la protéger des mauvais esprits, et agencée avec soin. Pour un peu, on ne regarderait que ses pieds, sur quoi ils sont posés, foulant des merveilles, oubliant le jardin et ses mille compositions. Katsura est particulièrement renommé pour ses pierres parmi les amateurs de jardins. Elles n’ont pas bougé depuis quatre siècles, mémoire du pas ralenti, chemin dense, préparation du corps et de l’esprit avant que d’atteindre au \teiHi[rend={italic}]{chashitsu}, la pièce pour le thé. }
              \teiP[xmlid={p44}]{Au gré d’un détour, c’est une vasque d’ablutions, \teiHi[rend={italic}]{chōzubachi}, savamment taillée, géométriquement, comme un moderne ou un artiste contemporain n’aurait su mieux faire. Éclairée d’une nouvelle lanterne, elle précède le banc et le préau \teiHi[rend={italic}]{soto-koshikake}, où les invités attendront que l’hôte vienne les chercher. Petit abri rustique fait de branches pas même écorcées, il est en fait d’une rusticité on ne peut plus sophistiquée, ce qui caractérise le style \teiHi[rend={italic}]{sukiya zukuri}. }
              \teiP[xmlid={p45}]{Devant eux se dressent sur un talus quinze sagoutiers \teiHi[rend={italic}]{sotetsu} aux troncs noirs et courbés, stipes plus que troncs, aux palmes hirsutes, natifs des contrées méridionales de Kyūshū. À leur pied s’étend une longue bande minérale\teiNote[xmlid={ftn51-3},n={18},place={foot}]{\teiP[xmlid={p-368}]{Voir la planche 46, dans Philippe Bonnin, \teiHi[rend={italic}]{Vocabulaire…}, \teiHi[rend={italic}]{op. cit.}}}, rectiligne, stricte, un agencement de pierres étonnant, semi-ordonné, qui forme un long rectangle très contradictoire avec le tracé du sentier sinueux que l’on vient de parcourir. Et c’est bien une volonté de s’opposer esthétiquement aux règles édictées par le grand maître de thé que l’on révère pourtant. Un défi à la Furuta Oribe pour se libérer de l’emprise de ses modèles, tout en conservant leur niveau d’exigence. Un exégète de l’esthétique japonaise y reconnaîtrait la fusion et l’articulation des trois styles \teiHi[rend={italic}]{shin-gyō-sō} : officiel, relâché, libre.}
            
\end{teiDiv}
            \begin{teiDiv}[type={section1},xmlid={div08-4}]

              \teiHead[xmlid={la-dalle-qui-devient-pont-001}]{La dalle qui devient pont}
              \teiP[xmlid={p46}]{Plus loin, le chemin devient dalle, puis pont au-dessus d’un menu torrent, une dramatisation du passage, sans garde-corps, vertigineux, au-dessus de la « cascade du tambour », et là encore, la mémoire des poètes est sollicitée, tant le pont fait figure du lien, tout comme chez nous\teiNote[xmlid={ftn52-3},n={19},place={foot}]{\teiP[xmlid={p-369}]{On se souviendra de Georg Simmel, \teiHi[rend={italic}]{La tragédie de la culture}, Paris, Rivages, 1988, p. 161-168 : « Pont et porte ». }}, et bien plus encore lorsqu’il est chanté par la femme amoureuse : }
              \begin{teiLg}

                \teiL{Large est le cours}
                \teiL{Au gué de la \lbrack{}rivière\rbrack{} Sahogawa}
                \teiL{Où crie le pluvier}
                \teiL{J’y jette donc un pont}
                \teiL{Dans l’espoir que vous viendrez\teiNote[xmlid={ftn53-3},n={20},place={foot}]{\teiP[xmlid={p-370}]{\teiHi[rend={italic}]{Man’yōshū}, 528, dans \teiHi[rend={italic}]{Nō et Kyōgen}…, \teiHi[rend={italic}]{op. cit}.}}.}
              
\end{teiLg}
              \teiP[xmlid={p47}]{Le son de la cascade évoque la présence d’un torrent, d’une source de montagne, d’une eau pure venant du nord-est qui va baigner le jardin, selon les règles immuables. Le procédé est ici celui de la transposition : jamais on ne voit la source, jamais on ne voit la cascade. On l’entend. Et près de la dalle du pont est dressée une haute pierre à la face droite comme une chute d’eau. C’est ainsi seulement qu’on la voit, transposée : un \teiHi[rend={italic}]{mitate}.}
              \teiP[xmlid={p48}]{Ce sont ensuite des évocations de rivages maritimes, paysages du pays natal de beaucoup de ces nobles japonais qui ont dû quitter leur domaine pour venir à la cour, et gardent la nostalgie des flots se déchirant sur les récifs sauvages, s’épuisant sur les plages de galets et de sable \teiHi[rend={italic}]{suhama}. Tel était le pays de l’épouse du prince, les abords de Miyazu, et particulièrement cette évocation du cordon lagunaire d’Amanohashidate – figuré par un fin monolithe de granit lancé entre deux îles, à la taille stricte, à peine arqué comme un \teiHi[rend={italic}]{koto} – lequel fut lui-même appelé ainsi parce qu’il était l’image du « pont lancé sur le ciel » dans le récit mythologique du \teiHi[rend={italic}]{Kojiki}, la cosmologie du Japon, le récit de la création d’une terre stable. Un \teiHi[rend={italic}]{mitate} au second degré, pourrait-on dire. }
              \teiP[xmlid={p49}]{Parfois, les rivages côtiers, les rochers, les bancs de sable, les pins, les ponts, tout s’estompe, tout est noyé dans la bruine, le brouillard, ou même les trombes des pluies équatoriales. Ils n’en deviennent que mystérieusement plus beaux, d’une beauté à deviner plutôt qu’elle ne sera offerte violemment en évidence, beauté à construire dans son esprit, qu’on dit \teiHi[rend={italic}]{kakurenbaï}\teiNote[xmlid={ftn54-3},n={21},place={foot}]{\teiP[xmlid={p-371}]{\teiHi[rend={italic}]{Bai} est le prunier, summum de beauté dans les classifications japonaises \teiHi[rend={italic}]{shô-chiku-baï }: \teiHi[rend={italic}]{shô} le pin, \teiHi[rend={italic}]{chiku} le bambou, \teiHi[rend={italic}]{baï} le prunier.}}, la beauté cachée, qui s’efface, la beauté profonde et mystérieuse, \teiHi[rend={italic}]{yūgen}, « à travers l’élégance et le rêve », dit l’historien Itoh Teiji.}
            
\end{teiDiv}
            \begin{teiDiv}[type={section1},xmlid={div09-3}]

              \teiHead[xmlid={le-luth-et-le-pin-001}]{Le luth et le pin}
              \teiP[xmlid={p50}]{C’est après cent détours, mille et une extases que l’on parvient en vue du « pavillon du luth et du pin », le Shōkin-Tei \lbrack{}pavillon \teiHi[rend={italic}]{tei}, du luth \teiHi[rend={italic}]{kin/koto}, et du pin s\teiHi[rend={italic}]{hō/matsu}\rbrack{}. Ce nom s’inspire du livre dix-huitième du \teiHi[rend={italic}]{Genji monogatari} : \teiHi[rend={italic}]{matsukaze}, « le vent dans les pins ». Le prince Toshihito, grand joueur de \teiHi[rend={italic}]{koto}, lors d’une soirée du 14\teiHi[rend={sup}]{e} jour de la neuvième lune de l’an 18 de l’ère Keichō (1613), où fut proposé le thème : « Le \teiHi[rend={italic}]{koto} de nuit avec le vent parmi les pins », composa un poème impromptu :}
              \begin{teiLg}

                \teiL{Comme il souffle sur les cordes du koto, }
                \teiL{Le vent parmi les pins chante une mélodie}
                \teiL{De printemps et d’automne.}
              
\end{teiLg}
              \teiP[xmlid={p51}]{On n’en finirait pas de décrire et d’analyser les dispositions de ce pavillon, de sa pièce de thé, de sa facture et de son décor peint, de sa simple couverture de chaume quand ses communs le sont de tuiles, comme les plus riches temples bouddhistes. Face au lever de la lune, on y poursuivait les fêtes nocturnes. Puis on s’approchait de la petite vallée voisine, la « vallée des lucioles » ; \teiHi[rend={italic}]{hotaru}, dans l’ombre des soirées aux premières chaleurs caniculaires, à la moitié du septième mois. Elles réveillaient la mémoire de ces multitudes de poèmes qu’elles ont suscités. Les lucioles sont devenues un de ces \teiHi[rend={italic}]{kigo}, mots marqueurs obligés de la composition des haïkus. }
              \begin{teiLg}

                \teiL{Chasse aux lucioles}
                \teiL{Le batelier est ivre}
                \teiL{Quelle catastrophe\teiNote[xmlid={ftn55-3},n={22},place={foot}]{\teiP[xmlid={p-372}]{Bashō, \teiHi[rend={italic}]{Cent cinq haïkaï}, Koumiko Muraoka et Fouad El-Etr (trad.), Paris, La Délirante, 1979, p. 20. }}.}
              
\end{teiLg}
              \teiP[xmlid={p52}]{Les lucioles évoquent les âmes des défunts qui hantent la nuit, les rêves, les inquiétudes et les souvenirs. O-bon, la grande fête du retour des âmes mortes n'est pas très loin : dans un mois seulement.}
              \begin{teiLg}

                \teiL{De l’obscurité}
                \teiL{À l’appel de l’homme obscur}
                \teiL{Une luciole\teiNote[xmlid={ftn56-3},n={23},place={foot}]{\teiP[xmlid={p-373}]{Fūteki (répondant au poème d’Izumi Sikibu, p. 137), dans Bashō, \teiHi[rend={italic}]{Friches}, René Sieffert (trad.), Paris, Verdier, « Verdier poche », 2006, p. 123 et 128. }}.}
              
\end{teiLg}
            
\end{teiDiv}
            \begin{teiDiv}[type={section1},xmlid={div10-2}]

              \teiHead[xmlid={lile-du-grand-mont-001}]{L’île du grand mont }
              \teiP[xmlid={p53}]{Dans cette fantasmagorie d’un monde totalisant contenu dans ces jardins, c’est ensuite une indispensable montagne qui complète le paysage japonais, une montagne de quinze mètres à vrai dire, \teiHi[rend={italic}]{ōyamashima}, qui réunit tous les attributs nécessaires : la forêt sombre et humide, les sentiers escarpés traversés d'entrelacs de racines serpentines, et jusqu'au pavillon du col, tout comme sur la \teiHi[rend={italic}]{Tokaidô}, pour prendre repos et se rafraîchir. Là est le « pavillon de thé du col », le pavillon « pour contempler les fleurs ».}
              \begin{teiLg}

                \teiL{Plus haut que les alouettes}
                \teiL{Je me repose en plein ciel}
                \teiL{Au col de la montagne\teiNote[xmlid={ftn57-3},n={24},place={foot}]{\teiP[xmlid={p-374}]{Bashō, \teiHi[rend={italic}]{Cent cinq haikai}, \teiHi[rend={italic}]{op. cit}., non paginé.}}.}
              
\end{teiLg}
              \teiP[xmlid={p54}]{On y contemplait la glycine, dit-on, si vénérée pour ses grappes au violet pâle, rosées, mauves ou blanches, parfois presque d’un bleu de pervenche un peu passé, que chaque jardin public conduit aujourd’hui sur une vaste pergola. Mais pourquoi n’y en a-t-il plus ici, aujourd’hui, plus une seule ?}
              \begin{teiFigure}[xmlid={figure03-3}]

                \teiHead[xmlid={figure-19-a-linterieur-du-sanctuaire-kameido-tenjin-de-la-serie-cent-vues-dedo-utagawa-hiroshige-c-1857-gravure-sur-bois-brooklyn-museum-gift-of-dr-and-mrs-frank-l-babbott-jr-62-79-5-photographie-brooklyn-museum-001}]{Figure 19.\teiHi[rend={italic}]{ À l’intérieur du sanctuaire Kameido Tenjin}, de la série \teiHi[rend={italic}]{Cent vues d’Edo}, Utagawa Hiroshige, c. 1857. Gravure sur bois. Brooklyn Museum, Gift of Dr. and Mrs. Frank L. Babbott, Jr., 62.79.5. Photographie : Brooklyn Museum. }
              
\end{teiFigure}
              \teiP[xmlid={p55}]{À défaut de glycines, le regard est comblé d’un \teiHi[rend={italic}]{shakkei} (un « emprunt de paysage »), ce procédé qui, effaçant le second plan derrière quelque haie, capte la silhouette bleutée des montagnes lointaines pour les composer dans le jardin. Ainsi le monde est-il au complet, parfait.}
              \teiP[xmlid={p56}]{Plus loin, le prince Toshitada s’était fait construire une thébaïde, pour se retirer plus encore de l’agitation et des mondanités, comme un pur ermite et poète. C’est le pavillon Shōi-Ken, « Le cœur au calme » (ou : le pavillon du sourire), dénomination qui fait référence à la poésie chinoise antique, celle du très fameux poète Li Bai (701-762). Souhaitant se retirer dans un ermitage montagnard, selon la tradition des lettrés, pour regarder le monde avec distance et ironie, il écrivait alors : }
              \begin{teiCit}[xmlid={cit01-6}]

                \begin{teiQuote}
Quand ils demandent ce que dans l’azur des montagnes je pense de la vie, je ris et ne réponds rien. Que mon cœur ici trouve la tranquillité, les pêchers en fleurs et les ondes courantes ont disparu au loin, bien loin. Autre est mon univers. Les hommes n’y sont point\teiNote[xmlid={ftn58-3},n={25},place={foot}]{\teiP[xmlid={p-375}]{Un poème, sans doute le plus connu de Li Bai, est appris par tous les enfants en Chine (\teiHi[rend={italic}]{Tangshi}, VII, 1, n\teiHi[rend={sup}]{o} 233, dans \teiHi[rend={italic}]{Anthologie de la poésie chinoise}, Paris, Gallimard, « Bibliothèque de la Pléiade », n\teiHi[rend={sup}]{o} 602, 2015).}}.
\end{teiQuote}
              
\end{teiCit}
              \begin{teiLg}

                \teiL{Pensées d’une nuit calme}
                \teiL{La lumière de la lune baigne mon chevet}
                \teiL{Comme si elle en couvrait le sol de givre}
                \teiL{La tête aux étoiles je contemple le disque lumineux}
                \teiL{Mais baissant les yeux je songe au lieu où je suis né\teiNote[xmlid={ftn59-3},n={26},place={foot}]{\teiP[xmlid={p-376}]{\teiHi[rend={italic}]{Ibid}.}}.}
              
\end{teiLg}
              \teiP[xmlid={p57}]{Que de sagesse ! Pourtant… d’après la légende, Li Bai serait mort d’avoir voulu, ivre, sur un bateau, attraper le reflet de la lune dans l'eau.}
              \teiP[xmlid={p58}]{On ne saurait quitter ces jardins, même à grandes enjambées, même en feignant d’ignorer mille merveilles à y détailler, sans évoquer l’estrade « pour contempler la lune », \teiHi[rend={italic}]{Tsukimi-dai}, qui prolonge le sol du palais au sud-est, en avancée vers l’étang, où se reflètera la lune d’automne montante.}
              \teiP[xmlid={p59}]{Les courtisans de l’époque Heian avaient coutume d’aller admirer la lune dans ce lieu idéal de la basse Katsura. Ce rituel de contemplation de la lune \teiHi[rend={italic}]{tsukimi }de la nuit du quinze de la huitième lune de l’ancien calendrier, donc d’automne, fut observé par la cour jusqu’au\teiHi[rend={small-caps}]{ xiii}\teiHi[rend={sup}]{e} siècle environ, puis se répandit parmi les esthètes et poètes des classes aisées.}
              \begin{teiLg}

                \teiL{En votre séjour}
                \teiL{Près la rivière où limpide}
                \teiL{Se mire la lune}
                \teiL{Bien agréable est sans doute }
                \teiL{l’ombrage du Katsura\teiNote[xmlid={ftn60-3},n={27},place={foot}]{\teiP[xmlid={p-377}]{\teiHi[rend={italic}]{Le Dit du Genji}, René Sieffert (trad.), Paris, Éditions Diane de Selliers, 2007, Ch. XIV, p. 33.}}.}
              
\end{teiLg}
              \begin{teiFigure}[xmlid={figure04-3}]

                \teiHead[xmlid={figure-20-katsura-lestrade-pour-contempler-la-lune-tsukimi-dai-photographie-philippe-bonnin-001}]{Figure 20. Katsura : L’estrade « pour contempler la lune », Tsukimi-dai (月見台). Photographie Philippe Bonnin.}
              
\end{teiFigure}
              \teiP[xmlid={p60}]{Une légende ancienne raconte qu’un jour, le roi de toutes les divinités descend sur terre sous les traits d’un honorable vieillard (ou du Bouddha). Ayant faim, il frappe à la porte et demande qu’on lui serve un repas. Le singe lui apporte les fruits qu’il a cueillis dans les arbres, le renard lui présente un poisson, mais le lièvre n’a rien à lui offrir. Le dieu le blâme sévèrement. Le lièvre demande alors à ses amis de préparer un bon feu. Lorsque les flammes s’élèvent et que les braises sont suffisantes, le lièvre se jette sur le bûcher en sacrifice et offre son corps comme nourriture au vieillard. Ce dernier, déclinant son titre de roi des dieux, réunit les restes du lièvre et les place sur la lune afin que son sacrifice soit visible de tous. Depuis cette époque, dit la légende, il y a un lièvre dans la lune. Et chaque année, le 15\teiHi[rend={sup}]{e} jour du 8\teiHi[rend={sup}]{e} mois lunaire (qui tombe en septembre ou en octobre en fonction du calendrier), se tient la fête \teiHi[rend={italic}]{Tsukimi}, la fête de la pleine lune. }
              \begin{teiLg}

                \teiL{Vers le sud-ouest incline}
                \teiL{La fleur de katsura}
                \teiL{Encore en bouton\teiNote[xmlid={ftn61-3},n={28},place={foot}]{\teiP[xmlid={p-378}]{Uritsu, dans \teiHi[rend={italic}]{Bashō, jours d’hiver}, René Sieffert (trad.), Aurillac, POF, 1987, p. 55.}}.}
              
\end{teiLg}
              \teiP[xmlid={p61}]{La « fleur de Katsura », c’est évidemment la lune, « encore en bouton » – à son premier quartier – au début du mois, qui se couche tôt, déclinant vers l’ouest.}
              \teiP[xmlid={p62}]{Une autre légende, d’origine chinoise également, raconte qu’un homme avait percé le secret taoïste de l’immortalité, mais qu’il se rendit coupable d’une offense grave. Il fut alors exilé sur la lune, et condamné pour l’éternité à débiter un arbre \teiHi[rend={italic}]{katsura} de cinq mille pieds de haut. Ainsi l’y aperçoit-on. Le galant poète en tire une image :}
              \begin{teiLg}

                \teiL{De mes yeux vous vois}
                \teiL{Mais la main ne puis vous prendre}
                \teiL{Vous qui êtes pareille}
                \teiL{Au katsura de la lune,}
                \teiL{M’amie que dois-je vous faire\teiNote[xmlid={ftn62-3},n={29},place={foot}]{\teiP[xmlid={p-379}]{\teiHi[rend={italic}]{Man’yōshū}, 632, René Sieffer (trad.), Aurillac, POF, 1998.}} ?}
              
\end{teiLg}
              \teiP[xmlid={p63}]{Une autre version, d’origine chinoise, rapporte encore qu’un arbre fabuleux est effectivement planté au Palais de la lune. L’être divin Wugang qui s’acharne à l’abattre n’y parvient jamais. Par la violence de ses coups, une nuit de pleine lune d’automne, il en fait tomber une pluie de graines que ramassa, surpris, le moine Deming. Les montrant à son maître Zhiyi, il en reçoit l’explication, et l’autorisation de les semer autour de leur temple Ligyin, au bord du lac de l’ouest, sur le « pic aux osmanthes de la lune ». Cela se passe sous le règne de l’Empereur Xuanzong de la dynastie des T’ang.}
              \teiP[xmlid={p64}]{Pour les Chinois, c’est un osmanthe (le cassier des Chinois, \teiHi[rend={italic}]{olea fragrans,} olivier odoriférant). Pour les Japonais, c’est un \teiHi[rend={italic}]{katsura}, dit « arbre caramel » (le gainier japonais\teiHi[rend={italic}]{ cercidiphyllum japonicum}) en raison du parfum délicat de ses feuilles, à la fin de l’été. Certains critiques européens l’assimilent même à l’arbre de Judée (\teiHi[rend={italic}]{cercis siliquastrum}).}
              \teiP[xmlid={p65}]{Au linteau du dernier pavillon de ce jardin, « la tour de la lune et de la vague », dit aussi « Le pavillon de thé du prunier sous la lune », est placée une planchette qui porte deux caractères peints de blanc, de la main de l’empereur Reigen (1654-1732, 112\teiHi[rend={sup}]{e} \teiHi[rend={italic}]{tennō}) : on y lit \teiHi[rend={italic}]{gekka} 月歌, c’est-à-dire « la poésie et la lune », dit Mori Osamu, ou bien \teiHi[rend={italic}]{getsu no uta}, la chanson de la lune.}
              \begin{teiLg}

                \teiL{Que je voie la lune}
                \teiL{Et la tristesse me saisit}
                \teiL{Bien que l’automne}
                \teiL{Ne vienne pas seulement pour moi\teiNote[xmlid={ftn63-3},n={30},place={foot}]{\teiP[xmlid={p-380}]{\teiHi[rend={italic}]{Oe-No-Chisato}, chant 193 : \teiHi[rend={italic}]{Kokinshū}, Alain Kervern (trad.) dans Alain Kervern, \teiHi[rend={italic}]{À l’ouest blanchit la lune. Grand Almanach poétique japonais}, Bédée, Folle Avoine, 1992, Livre IV, p. 64.}}.}
              
\end{teiLg}
              \teiP[xmlid={p66}]{Alain Kervern, qui transcrit ce haïku ancien, du début de l’époque de Heian\teiNote[xmlid={ftn64-3},n={31},place={foot}]{\teiP[xmlid={p-381}]{Alain Kervern, \teiHi[rend={italic}]{op. cit}.}}, rapporte également ces mots de Bashō : « L’homme raffiné vit en amitié avec les saisons et se conforme aux grands principes de la nature. Il ne regarde que les fleurs et ne pense qu’à la lune\teiNote[xmlid={ftn65-4},n={32},place={foot}]{\teiP[xmlid={p-382}]{\teiHi[rend={italic}]{Ibid}., p. 69.}}. »}
              \begin{teiFigure}[xmlid={figure05-3}]

                \teiHead[xmlid={figures-21a-et-b-lestrade-pour-contempler-la-lune-tsukimi-dai-photographie-philippe-bonnin-et-lune-dautomne-a-ishiyama-hiroshige-c-1834-gravure-sur-bois-domaine-public-via-wikimedia-commons-001}]{Figures 21a et b. L’estrade « pour contempler la lune », Tsukimi-dai (月見台). Photographie Philippe Bonnin et \teiHi[rend={italic}]{Lune d’automne à Ishiyama}, Hiroshige, c. 1834. Gravure sur bois. Domaine public, via Wikimedia Commons.}
              
\end{teiFigure}
              \teiP[xmlid={p67}]{Les peintures et gravures d’Andō Hiroshige (1797-1851) incluent dans chaque série de vues paysagères une lune d’automne à son lever : la « lune d’automne de la rivière Takara », la « lune d’automne sur le mont Hase », etc.}
            
\end{teiDiv}
            \begin{teiDiv}[type={section1},xmlid={div11-2}]

              \teiHead[xmlid={dans-le-gout-japonais-001}]{Dans le goût japonais}
              \teiP[xmlid={p68}]{Certes, il est des auteurs, japonais eux-mêmes, qui critiquent rudement ce fait que les Japonais ne sauraient reconnaître la beauté des choses, des huit paysages d’Omi par exemple\teiNote[xmlid={ftn66-3},n={33},place={foot}]{\teiP[xmlid={p-383}]{On retrouve la lune à l’honneur dans les « huit beaux paysages de Omi », \teiHi[rend={italic}]{Omi hakkei,} choisis par Konoe Masaie sur le pourtour du lac Biwa à l’imitation des « huit beaux paysages de Xiâo Xiâng » du modèle chinois, telle la « lune d’automne au temple d’Ishikawa ». Voir « Esthétique et urbanité », dans Philippe Bonnin et Jacques Pezeu-Massabuau, \teiHi[rend={italic}]{Façons d’habiter le Japon, maisons, villes et seuils}, Paris, CNRS Éditions, 2017, p. 63-89.}}, que s’ils ont été préalablement désignés à cette dignité par une autorité, par l’Autre, par les peintres chinois en ce qui concerne particulièrement les huit beaux paysages de la Xiao Xiang\teiNote[xmlid={ftn67-3},n={34},place={foot}]{\teiP[xmlid={p-384}]{Voir à ce sujet Augustin Berque, « La naissance du paysage en Chine », \teiHi[rend={italic}]{Xoana}, n\teiHi[rend={sup}]{o} 5, 1977, p. 23-28.}}. Souvent les premières apparences leur donneraient raison, même si l’on ne peut nier la création d’une esthétique spécifique au Japon dans les phases où celui-ci s’est refermé au monde. Mais alors, si les œuvres japonaises ne pouvaient se comprendre autrement que par référence, pourquoi sommes-nous, nous Occidentaux généralement ignorants de ces références, pourquoi sommes-nous captivés par leur beauté ? La dimension que nous reconnaissons comme créative, originale, se superpose-t-elle toujours à une dimension référentielle ? L’interréférentialité des objets esthétiques est-elle une loi du genre, obligée, plus générale qu’on ne le croit, qui s’impose à toute création, aussi libre se croit-elle ?}
            
\end{teiDiv}
          
\end{teiElement}
        
\end{teiElement}
      
\end{teiElement}
      \begin{teiElement}[name={group},type={section1},xmlid={section1-003}]

        \teiHead[xmlid={quels-paradigmes-de-beaute-dans-les-arts-001}]{Quels paradigmes de beauté dans les arts ?}
        \begin{teiElement}[name={group},type={chapter},data-page-title={Notes pour une philosophie de la tension de surface ou La beauté comme relation},data-page-authors={Chiara Vecchiarelli@@ENS-PSL, Centre international d’étude de la philosophie française contemporaine, France},data-include-href={Ch10\_note\_pour\_une\_philosophie\_Vecchiarelli.xml},xmlid={chapter-009}]

          \begin{teiElement}[name={front}]

            \begin{teiDiv}[type={titlePage},xmlid={titlepage-011}]

              \teiP[rend={title-main},xmlid={p-385}]{Notes pour une philosophie de la tension de surface ou La beauté comme relation}
              \teiP[rend={author-aut},xmlid={p-386}]{Chiara \teiHi[rend={small-caps}]{Vecchiarelli}}
              \teiP[rend={authority\_affiliation},xmlid={p-387}]{ENS-PSL, Centre international d’étude de la philosophie française contemporaine, France}
            
\end{teiDiv}
          
\end{teiElement}
          \begin{teiElement}[name={body},xmlid={body-011}]

            \teiP[xmlid={p1-11}]{Qu’auraient en commun une pierre et une bulle de savon ?}
            \teiP[xmlid={p2-11}]{L’une demeure impassible au toucher, l’autre éclate. L’une semble tendre vers l’intemporel par sa promesse de durabilité ; l’autre participe de la plus brève des durées, celle des choses qui apparaissent et disparaissent aussitôt. Le ferme et l’éphémère, l’aérien et le pesant semblent s’opposer, dans ces êtres qui paraissent si différents lorsqu’on les considère dans leur en-soi.}
            \teiP[xmlid={p3-11}]{Si de nombreux témoignages existent quant à la beauté de telles réalités prises individuellement – qu’il suffise de songer à l’art japonais du \teiHi[rend={italic}]{suiseki}, ou à l’iconographie occidentale des bulles de savon dans la tradition picturale – c’est pourtant vers une tout autre conception de la beauté que peuvent nous conduire la pierre et la bulle : une dimension qui se situe au-delà de l’individu, qui excède toute compréhension de ces entités comme des substances closes sur elles-mêmes.}
            \teiP[xmlid={p4-11}]{Nous proposerons de penser une telle dimension comme étant relationnelle et, précisément en tant que relationnelle, vitale.}
            \teiP[xmlid={p5-11}]{Assigner à la relation une teneur esthétique, lui reconnaître une fonction vitale, est un geste qui ne va pas de soi. Cela suppose de repenser nos catégories esthétiques et de faire émerger un champ de pensée qui ne se confond ni avec celui de l’objet perçu, ni avec celui du sujet qui imagine ou perçoit ; un lieu qui ne se situe ni dans l’œil du regardeur ni dans la chose en soi.}
            \teiP[xmlid={p6-11}]{Nous allons suggérer qu’une pierre et une bulle peuvent avoir suffisamment de force pour nous conduire vers cette dimension au-delà de l’objet et du sujet, un au-delà qui est, en dernière instance, un en deçà, car logiquement antérieur à l’un comme à l’autre.}
            \teiP[xmlid={p7-11}]{Cette force porte un nom dissimulé dans le fil de l’eau, dans la peau d’une bulle.}
            \teiP[xmlid={p8-11}]{Tôt ou tard dans la vie, nous avons tous pu être fascinés par la beauté éphémère des bulles de savon – par leur flottement à mi-chemin entre le ciel et la terre, toujours sur le point de disparaître. Les bulles de savon sont des entités singulières en ceci que leur existence coïncide avec leur surface : elles sont, tout entières, surface. Cette face extérieure – qui représente, pour chaque être, sa limite immédiate et définit en même temps le seuil à partir duquel un échange, une rencontre, peut advenir – les constitue sans reste. Une bulle est alors à la fois la possibilité d’un échange et sa visibilité ; la possibilité tendue de la rencontre, et sa fragilité, toujours menacée par un autre toucher.}
            \teiP[xmlid={p9-10}]{Ce qui rend une telle existence possible, c’est une force physique : la tension de surface\teiHi[rend={italic}]{.}}
            \teiP[xmlid={p10-10}]{Notion physique et chimique à la fois, la tension de surface se définit comme la grandeur positive caractérisant la relation entre deux surfaces. Lorsqu’un contact se produit entre des corps fluides, leur géométrie se modifie, l’aire de l’interface diminue et une augmentation d’énergie accompagne cette modification : la rencontre devient alors un événement énergétique.}
            \teiP[xmlid={p11-10}]{Ces êtres qui se touchent connaissent alors une intensification.}
            \teiP[xmlid={p12-10}]{La tension de surface désigne précisément cette force qui maintient l’équilibre d’une énergie accrue à l’interface entre deux fluides. Elle agit à l’endroit – et dans l’événement – de leur rencontre\teiNote[n={1},place={foot},xmlid={ftn1-10}]{\teiP[xmlid={p-388}]{Voir Étienne Guyon, Jean-Pierre Hulin et Luc Petit, Hydrodynamique physique, Les Ulis, EDP Sciences, 1996 ; Charles Vernon Boys, Soap-Bubbles and the Forces Which Mould Them \lbrack{}The Society for Promoting Christian Knowledge, 1902\rbrack{}, New York, Doubleday Anchor Books, 1959 ; John H. Lienhard IV, John H. Lienhard V, A Heat Transfer Textbook \lbrack{}5e éd.\rbrack{}, Cambridge, Phlogiston Press, 2019.}}. C’est une tension qui donne réalité à la relation, qui en soutient l’existence sans pour autant la substantialiser, sans la faire précipiter dans une forme figée.}
            \teiP[xmlid={p13-9}]{Toute bulle visqueuse existe par l’équilibre entre deux forces : la tension de surface, qui tend à minimiser l’air de la bulle – car l’énergie de surface diminue avec l’augmentation de cette dernière – et la pression exercée par l’air intérieur, qui s’oppose à la diminution du volume. On dira donc de la bulle visqueuse qu’elle existe entièrement sur la limite, transparente et tendue, de la relation. Et que cette limite possède tout le potentiel nécessaire pour devenir une notion philosophique : elle est, pour nous, à la fois l’expérience même d’une rencontre et un événement philosophique.}
            \teiP[xmlid={p14-9}]{La vision d’une bulle qui flotte se dédouble en intensité – et en potentiel – dans l’événement plus rare encore de la rencontre entre deux bulles. Le contact engendre, chez ces êtres flottants, un être nouveau. Lorsque deux bulles visqueuses se joignent, l’entre-deux se réinvente à l’interface, et un réel inédit voit le jour. Là où deux tensions de surface se rencontrent, un événement topologique a lieu : une nouvelle forme d’existence apparaît. On pourrait dire que l’interface tendue entre les deux bulles devient partie intégrante de l’être de l’une et de l’autre. La surface tendue de la rencontre réunit alors deux forces : la force physique qu’est la tension de surface, et la force relationnelle, qui se manifeste dans l’invention d’une topologie partagée, par laquelle deux êtres entrent en coïncidence.}
            \teiP[xmlid={p15-9}]{Cette force permet également, dans certaines conditions, à certains organismes de se mouvoir à la surface de l’eau. Un exemple du rôle de la tension de surface est observable sur les étendues d’eau calmes, telles que les étangs, sur lesquels des insectes appelés « patineurs d’eau » (\teiHi[rend={italic}]{gerris lacustris}) se déplacent comme sur une surface élastique. Ce phénomène est rendu possible par la combinaison de la tension superficielle de l’eau – qui agit comme une membrane élastique en repoussant les substances hydrophobes – et de la légère déformation de la surface, induite par le poids de l’insecte, laquelle génère une force de sustentation suffisante pour le maintenir à flot et lui permettre de se mouvoir à la surface de l’eau comme sur une toile tendue.}
            \teiP[xmlid={p16-9}]{Dans son ouvrage \teiHi[rend={italic}]{Imagination et invention}\teiNote[n={2},place={foot},xmlid={ftn14}]{\teiP[xmlid={p-389}]{Gilbert Simondon, Imagination et invention. 1965-1966 \lbrack{}2008\rbrack{}, Paris, PUF, 2014. }}, le philosophe Gilbert Simondon propose une analogie entre un autre phénomène aquatique lié à la tension de surface et les alvéoles des nids d’abeille. Ces dernières s’apparentent, dans son analyse, à la surface sphérique d’une goutte d’eau. Ce que ces formes ont en commun, c’est l’optimisation d’une relation : obtenir un maximum de surface avec un minimum de cire, ou – en tant que liquide – se contenir dans une sphère à la plus petite surface possible. Dans les deux cas, il s’agit de solutions que le réel a trouvées pour articuler un seuil – une relation – entre l’élément et le tout, entre le volume et le poids. Dans les mots de Simondon :}
            \begin{teiCit}[xmlid={cit01-7}]

              \begin{teiQuote}
\lbrack{}C\rbrack{}es minima deviennent des \teiHi[rend={italic}]{optima}, non en tant que systèmes dont l’Énergie est dégradée, mais quand ils sont pris comme solutions d’un problème, donc quand, au lieu d’être la formule du système d’ensemble, ils sont des moyens, des sous-ensembles apparaissant comme jouant un rôle de médiateurs et d’intermédiaires, entre une configuration plus vaste et une matière élémentaire\teiNote[n={3},place={foot},xmlid={ftn15}]{\teiP[xmlid={p-390}]{Gilbert Simondon, L’individuation à la lumière des notions de forme et d’information \lbrack{}1958\rbrack{}, Grenoble, Éditions Jérôme Millon, 2013, p. 86. Nous soulignons.}}.
\end{teiQuote}
            
\end{teiCit}
            \teiP[xmlid={p17-9}]{Par sa structure hexagonale, l’alvéole devient compatible avec toutes les autres ; par sa tension de surface, la goutte optimise la relation à son milieu\teiNote[n={4},place={foot},xmlid={ftn16-2}]{\teiP[xmlid={p-391}]{La structure hexagonale des alvéoles ne se limite pas à rendre compatibles les alvéoles entre eux, elle permet aussi de créer une surface continue au travers du plus grand nombre possible de relations entre chaque alvéole et les autres alvéoles. Chacun des hexagones est en contact avec six autres hexagones alors qu’une surface composée de carrés ne permet de travailler que quatre proximités. }}. Ces médiateurs sembleraient avoir pour fonction d’affiner la limite – de la rendre subtile, pourrions-nous dire, à la manière d’un \teiHi[rend={italic}]{inframince duchampien}. Ils sont des formes de la relation, la limite même de l’être tendu dans la relation.}
            \teiP[xmlid={p18-9}]{Un bloc de marbre blanc de Carrare posé au sol, un parallélépipède de quelques centimètres de hauteur et d’un demi-mètre de largeur dont la surface se tend sous nos yeux : la face supérieure de l’œuvre \teiHi[rend={italic}]{Milkstone }(1978) de Wolfgang Laib\teiNote[n={5},place={foot},xmlid={ftn17-2}]{\teiP[xmlid={p-392}]{Wolfgang Laib est un artiste allemand né en 1950, connu pour employer, dans ses œuvres, de grandes quantités de pollen qu’il tamise au sol pour créer des monochromes dont la tonalité est donnée par la couleur du pollen, faisant signe vers ce que l’artiste estime constituer « le début potentiel de la vie d’une plante » (« Pollen is the potential beginning of the life of the plant. It is as simple, as beautiful, and as complex as this. And of course, it has so many meanings. I think everybody who lives knows that pollen is important », site internet du MoMA, Musée d’art moderne de New York, \teiRef[target={https://www.moma.org/calendar/exhibitions/1315}]{https://www.moma.org/calendar/exhibitions/1315}, consulté le 1\teiHi[rend={sup}]{er} juillet 2025).}} possède la même blancheur immaculée que les autres faces de la sculpture, mais contrairement à celles-ci, elle présente une tension. La surface de l’œuvre est liquide. Légèrement creusée, elle est remplie d’une couche de lait dont le travail consiste à tendre, presque imperceptiblement, cette face supérieure. Toute entière tendue, la surface anime le marbre par ailleurs intemporel. Versé sur l’intemporalité de cette roche âgée de millions d’années, le lait amène une autre notion de temps, celle d’un temps éphémère qui s’évapore en quelques heures. Ces deux temporalités disparates – apparemment incompatibles, et pourtant composant le corps de l’œuvre – sont maintenues ensemble dans son unité, précisément par le phénomène de tension de surface.}
            \teiP[xmlid={p19-9}]{Dans un court commentaire sous forme d’enregistrement audio réalisé par le Museum of Modern Art de New York, l’artiste s’exprime à leur sujet en ces termes :}
            \begin{teiCit}[xmlid={cit02-6}]

              \begin{teiQuote}
Cette pièce dont nous parlons est une pierre de lait (\teiHi[rend={italic}]{Milkstone}). C’est une de mes toutes premières œuvres. C’est une petite dalle de marbre blanc pur avec une indentation très, très fine sur le dessus, entourée d’un rebord. Je commence à travailler avec une petite machine, mais la plupart du rebord et du reste, je les ponce à la main avec du papier de verre et de l’eau, ce qui demande beaucoup de travail. Mais c’est un travail méditatif d’une grande beauté, que j’aime toujours faire. Je participe à quelque chose que je ressens comme très indépendant de moi-même et qui possède aussi une signification universelle. On verse du lait à la surface, mais il ne reste que quelques heures, puis il faut le remplacer. Cette pierre a des millions d’années, et le lait n’est là que durant quelques heures. C’est une chose très, très simple, mais la surface du lait peut contenir tout ce que vous pouvez penser\teiNote[n={6},place={foot},xmlid={ftn18-2}]{\teiP[xmlid={p-393}]{« This piece we are talking about is a Milkstone. It’s one of my very early pieces. It’s a small slab of pure white marble which has a very, very fine indention on the top with a rim. I work first with a small machine, but most of the rim and everything I grind by hand with sandpaper and water, which is a lot of work. But it’s a very beautiful, meditative work, which I always like to do. I participate in something which I feel is very independent of myself and also has a universal meaning. And you pour milk into this on the surface, which is only for some hours, and then it has to be replaced. I mean the stone is like millions of years old and the milk is just there for some hours. It’s a very very simple thing, but the milk surface can contain everything you can think of. » Pour la transcription et la captation audio de la déclaration de Wolfgang Laib, voir le site internet du MoMA : \teiRef[target={https://www.moma.org/audio/playlist/50/754}]{https://www.moma.org/audio/playlist/50/754} (consulté le 1er juillet 2025). Notre traduction.}}.
\end{teiQuote}
            
\end{teiCit}
            \teiP[xmlid={p20-9}]{Du film de lait qui tend la surface du marbre du \teiHi[rend={italic}]{Milkstone}, Wolfgang Laib estime qu’il contient \teiHi[rend={italic}]{tout ce que l’on peut penser }: tout le potentiel, tout ce qui n’est pas encore dit, mais relève de l’ordre du pensable – de la puissance de penser. Cette surface devient l’image même du potentiel – non pas celui de l’œuvre en soi, mais de l’œuvre en train de se faire, à travers ce travail de tension qui s’exerce à la surface, devenue agent de la limite.}
            \teiP[xmlid={p21-7}]{Pour que le potentiel puisse agir, il faut lui faire une place : creuser le marbre, y inscrire une indentation – quelque chose comme un alinéa, cet espace qui signale le début d’un nouveau paragraphe, entendu ici comme un commencement inscrit dans la continuité de l’existence.}
            \teiP[xmlid={p22-7}]{Le creux empli de lait permet au marbre de revenir à un état de potentialité, de se « repotentialiser » dans une relation au sein de laquelle deux temporalités disparates deviennent co-présentes. Une relation qui permet à la pierre de se liquéfier – au sens de la \teiHi[rend={italic}]{liquefactio} des alchimistes chère à Jung. Image de la vie qui se tend dans le devenir, cette relation s’avère être un état de liquéfaction où tout devient co-présent.}
            \teiP[xmlid={p23-7}]{On pourrait dire, en mobilisant vocabulaire simondonien, que la surface tendue du lait de \teiHi[rend={italic}]{Milkstone} – cette peau liquide sur un corps de pierre – est l’image même de la \teiHi[rend={italic}]{limite tendue de l’être.}}
            \teiP[xmlid={p24-7}]{L’histoire de l’art ne manque pas d’exemples d’artistes ayant tenté d’animer le marbre, de lui insuffler une forme de vitalité. Aux galeries de la Villa Borghese à Rome, une sculpture commandée par le cardinal Scipione Borghese est installée au centre d’une salle, disposée de manière à permettre la déambulation autour d’elle : \teiHi[rend={italic}]{Le Rapt de Proserpine,} sculptée entre 1621 et 1622 par un tout jeune Gian Lorenzo Bernini – âgé de seulement vingt-trois ans lors de son achèvement. Celui qui, plus tard, animera les façades des palais romains en modulant les volumes comme nul architecte avant lui, représente ici l’enlèvement de Proserpine par Pluton, le maître des Enfers. Divinité chthonienne, premier fils de Cronos et Rhéa, Pluton porte plusieurs noms : \teiHi[rend={italic}]{Ploutôn}, \teiHi[rend={italic}]{Hadès} ou encore \teiHi[rend={italic}]{Haïdès} – « celui que l’on ne voit pas » ou bien, selon d’autres interprétations, « celui qui rend invisible ». C’est lui qui rend Proserpine invisible en l’enlevant – fille de Zeus et de Déméter – pour l’emmener dans son royaume souterrain. Chaque année, avec le printemps, Proserpine refait surface, accueillie sur Terre par un tapis de fleurs. Elle symbolise la saison du renouveau, le retour de la vie végétale, à laquelle les œuvres de Laib rendent hommage à travers le pollen.}
            \teiP[xmlid={p25-6}]{Elle redevient visible sous la forme d’un potentiel : celui d’une vie qui recommence après le long sommeil de l’hiver – comme après une mue imaginale. Un détail dans la sculpture de Bernini ne cesse d’attirer le regard : l’emprise de Pluton sur le corps de la jeune femme, ses doigts qui s’enfoncent dans la chair de Proserpine, modulant sa surface. Dans ce détail, qui s’apparente à la tension du lait de \teiHi[rend={italic}]{Milkstone}, se tiennent ensemble trois tensions : celle de Pluton pour le corps de Proserpine, celle de Proserpine pour la surface de la Terre et celle de Bernini pour la surface du marbre.}
            \teiP[xmlid={p26-6}]{La dynamique qui s’annonce dans le marbre de Bernini, et qui devient conscience matérielle dans le lait de \teiHi[rend={italic}]{Milkstone}, est de même nature que celle qui rend possible la formation d’une bulle de savon. Ensemble, ces tensions font signe vers cette force superficielle – en laquelle réside ce qui importe le plus dans la vie\teiNote[n={7},place={foot},xmlid={ftn19-2}]{\teiP[xmlid={p-394}]{Un développement ultérieur de ce texte pourra préciser en quoi le sujet mythologique choisi dans l’œuvre de Bernini demeure partiellement inachevé pour les besoins de notre analyse, en raison de sa mise en scène d’une relation non consensuelle. C’est précisément ce désaccord que l’œuvre d’art contemporaine engage, assumant ce que le sujet mythologique de Bernini laisse en suspens.}}.}
            \teiP[xmlid={p27-6}]{Elle constitue aussi le dispositif qui est au cœur d’une proposition esthétique d’une grande originalité, une proposition qui refuse de cantonner la question esthétique au seul niveau de la perception pour la déplacer sur le terrain ouvert de la relation. C’est ce déplacement qu’opère Gilbert Simondon dans son ouvrage \teiHi[rend={italic}]{Du mode d’existence des objets techniques}\teiNote[n={8},place={foot},xmlid={ftn20-2}]{\teiP[xmlid={p-395}]{Gilbert Simondon, Du mode d’existence des objets techniques \lbrack{}1958\rbrack{}, Paris, Aubier, 1989.}}, où il élabore une théorie esthétique novatrice : la beauté y est pensée non comme une qualité perçue, mais comme un mode d’être relationnel.}
            \teiP[xmlid={p28-5}]{Mais c’est d’abord la réalité esthétique elle-même qui se trouve redéfinie, envisagée comme une médiation capable de créer un monde – un monde logiquement antérieur à l’objet autant qu’au sujet, où la relation prime. Cette réflexion s’inscrit, chez Simondon, dans le cadre plus large d’une ontologie de la relation. Dans le prolongement de cette pensée, la beauté se trouve requalifiée : non plus comme une propriété intrinsèque à un objet, mais comme le fruit d’un agencement relationnel. Cette reconfiguration prend forme notamment dans l’étude des individus techniques, que Simondon envisage à travers le prisme de l’esthétique. Selon lui, « les objets techniques ne sont pas directement beaux en eux-mêmes\teiNote[n={9},place={foot},xmlid={ftn21-2}]{\teiP[xmlid={p-396}]{\teiHi[rend={italic}]{Ibid}, p. 254.}} », mais le deviennent lorsqu’ils « sont insérés dans un monde, soit géographique, soit humain\teiNote[n={10},place={foot},xmlid={ftn22-2}]{\teiP[xmlid={p-397}]{\teiHi[rend={italic}]{Ibid}.}} ». Autrement dit, la beauté n’émerge que lorsque l’objet dépasse sa fonction pure pour devenir un geste – un vecteur de relation. Si certains objets, y compris techniques, peuvent être considérés beaux, ce n’est donc qu’en raison de leur insertion dans un univers, de leur capacité à se transformer en points de convergence dans une trame qui organise le monde en un lieu vivant, un espace ouvert à la relation :}
            \begin{teiCit}[xmlid={cit03-6}]

              \begin{teiQuote}
\lbrack{}U\rbrack{}n outil, une machine ou un ensemble technique sont beaux quand ils s’insèrent dans un monde humain et le recouvrent en l’exprimant ; si l’alignement des tableaux d’un central téléphonique est beau, ce n’est pas en lui-même ni par sa relation au monde géographique, car il peut être n’importe où ; c’est parce que ces voyants lumineux qui tracent d’instant en instant des constellations multicolores et mouvantes représentent des gestes réels d’une multitude d’êtres humains, rattachés les uns aux autres par l’entrecroisement des circuits. Le central téléphonique est beau en action, parce qu’il est à tout instant l’expression et la réalisation d’un aspect de la vie d’une cité et d’une région ; une lumière, c’est une attente, une intention, un désir, une nouvelle imminente, une sonnerie qu’on n’entendra pas, mais qui va retentir au loin dans une autre maison\teiNote[n={11},place={foot},xmlid={ftn23-2}]{\teiP[xmlid={p-398}]{\teiHi[rend={italic}]{Ibid}., p. 256 et 257.}}.
\end{teiQuote}
            
\end{teiCit}
            \teiP[xmlid={p29-5}]{Cette beauté réside non pas dans l’objectualité, dont elle s’écarte, ni, a fortiori, dans le paradigme de l’imitation, que Simondon remplace par celui de l’insertion. Elle réside dans le potentiel relationnel du réel, une réalité qui est esthétique dans la mesure où elle incarne « le sens du devenir » car c’est en elle que tout ce qui est stable peut se mettre en mouvement, bourgeonner, rayonner, devenir relation créatrice – et non plus simplement terme d’une relation inessentielle, ou fonctionnelle à l’établissement d’un sujet fixe, voire à la représentation inadéquate d’un idéal inatteignable.}
            \teiP[xmlid={p30-5}]{C’est ainsi qu’il convient d’entendre la proposition selon laquelle tout geste médiateur est un geste esthétique\teiNote[n={12},place={foot},xmlid={ftn24-2}]{\teiP[xmlid={p-399}]{\teiHi[rend={italic}]{Ibid.}, p. 268.}}, et d’admettre également que « tout ce qui intervient comme intermédiaire entre sujet et objet peut prendre la valeur d’image\teiNote[n={13},place={foot},xmlid={ftn25-2}]{\teiP[xmlid={p-400}]{Gilbert Simondon, \teiHi[rend={italic}]{Imagination et invention},\teiHi[rend={italic}]{ op. cit.}, p. 12. Si tant est qu’il soit possible de désigner Proserpine comme « la jeune fille indicible », c’est uniquement en ce sens : dans la mesure où elle incarne une tension qui ne relève pas de l’ordre du discours, mais de celui de l’image – image au sein de laquelle, et d’elle seule, la relation devient première. Voir, pour une analyse de la dimension du mystère, entendu ici comme indicible, dans le mythe de Proserpine, Giorgio Agamben et Monica Ferrando, \teiHi[rend={italic}]{La ragazza indicibile. Mito e mistero di Kore}, Milan, Mondadori Electa, 2010.}} ». Cette perspective s’applique pleinement à l’œuvre d’art, à condition toutefois de cesser de la concevoir comme une substance figée, réduite à l’un des termes d’une relation sujet/objet.}
            \teiP[xmlid={p31-5}]{Certaines œuvres, en effet, fonctionnent davantage comme des champs relationnels : elles fondent leur principe sur la primauté de la relation. Ce sont ces œuvres qui nous invitent à repenser à la fois nos catégories esthétiques et nos priorités philosophiques.}
            \teiP[xmlid={p32-5}]{Ainsi, une œuvre d’art peut être saisie comme une relation vivante, à l’image d’une bulle de savon, toute entière tension de surface. Elle incarne la promesse de reconfigurer, dans la rencontre, tout réel – et ne sera belle que dans la mesure où elle portera la tension même de cette possible rencontre.}
          
\end{teiElement}
        
\end{teiElement}
        \begin{teiElement}[name={group},type={chapter},data-page-title={Sauver les phénomènes : la beauté vitale en photographie},data-page-subtitle={À partir de Simondon et de Marker},data-page-authors={Carole Maigné@@Université de Lausanne, Faculté des lettres, France},data-include-href={Ch11\_sauver\_les\_phenomenes\_Maigne.xml},xmlid={chapter-010}]

          \begin{teiElement}[name={front}]

            \begin{teiDiv}[type={titlePage},xmlid={titlepage-012}]

              \teiP[rend={title-main},xmlid={p-401}]{Sauver les phénomènes : la beauté vitale en photographie}
              \teiP[rend={title-sub},xmlid={p-402}]{À partir de Simondon et de Marker}
              \teiP[rend={author-aut},xmlid={p-403}]{Carole \teiHi[rend={small-caps}]{Maigné}}
              \teiP[rend={authority\_affiliation},xmlid={p-404}]{Université de Lausanne, Faculté des lettres, France}
            
\end{teiDiv}
          
\end{teiElement}
          \begin{teiElement}[name={body},xmlid={body-012}]

            \teiP[xmlid={p1-12}]{Quel sens peut avoir l’idée de beauté vitale en photographie ? La photographie pose d’emblée problème : enregistrer, fixer, sont des gestes sans lesquels la photographie n’est pas, mais vont contre la vitalité, interrompant le flux de la vie, figeant son mouvement. Il y a arrêt sur image, et cet arrêt semble sans fin, définitif. Placée ainsi hors temps à force de fixité, l’image photographique est hors flux. La dimension morbide et mortifère de la beauté en photographie et son association intime avec l’embaumement sont ainsi des \teiHi[rend={italic}]{topoi} répandus. Faut-il saisir une vitalité dans l’inépuisable production et l’extrême diffusion de l’image photographique, comme si la démultiplication de l’image et de son usage, voire le marché des images en tant que tel, autorisait à confondre circulation et flux vital, omniprésence et beauté ? Ce serait rabattre la beauté sur le stéréotype, ce flux étant an-esthétique jusqu’à l’anesthésie. Ce n’est pas en s’additionnant à d’autres images que la photographie trouve sa vitalité, elle s’y dissout plutôt, le flux étant à sa façon mortifère. Faut-il se résigner à identifier la beauté vitale de l’image photographique à celle de l’objet photographié ? En externalisant les raisons du beau hors de l’image, on réduirait l’image photographique à n’être qu’une simple fenêtre sur le monde : un médium si transparent qu’on en effacerait \teiHi[rend={italic}]{in fine} le cadre et la démarche, abolissant du même coup l’acte qui la permet. « L’idée de photographie », pour reprendre le titre du puissant livre de Brunet\teiNote[xmlid={ftn1-11},n={1},place={foot}]{\teiP[xmlid={p-405}]{François Brunet, \teiHi[rend={italic}]{La naissance de l’idée de photographie}, Paris, PUF, 2012. }}, tiendrait lieu de photographie, sans plus s’interroger sur l’acte de photographier. }
            \teiP[xmlid={p2-12}]{C’est plutôt dans la texture même de l’image photographique qu’il faut déceler une beauté vitale, autrement dit repartir de sa singularité propre. Le pari sera ici de tenir le syntagme « beauté vitale » en gardant la spécificité de cette image-là, c’est-à-dire la spécificité de sa procédure. Pour ouvrir cette voie, il faudra cesser de circonscrire la photographie à la saisie d’un objet, à une « prise » de vue, et y voir le travail d’une forme inventive, parce qu’imaginative, qui dans l’enregistrement même, ouvre au-delà de l’objet photographié. Questionner la beauté comme forme vivante impose de revisiter son réalisme supposé : l’image photographique n’est belle ni parce que son objet est beau, ni parce qu’elle est un beau linceul, ni parce que prise dans le flux du visuel, mais parce qu’elle est une forme vivante qui ne s’est jamais réellement soumise au régime naturaliste de l’empreinte figée qu’on lui a prêté. Réviser ce réalisme suppose de réviser ce que l’on peut entendre par image technique : l’image enregistrée ne nie pas nécessairement la vitalité du vivant. En termes simondoniens, la critique de l’hylémorphisme, à savoir la critique de l’image comme moule ou empreinte, sera au cœur de notre propos, contre un des \teiHi[rend={italic}]{topoi} les plus tenaces sur la photographie\teiNote[xmlid={ftn83-4},n={2},place={foot}]{\teiP[xmlid={p-406}]{Carole Maigné, « Philosophizing from and with the Camera: The wild invention of photography », dans Jamil Alioui, Matthieu Amat et Carole Maigné (dir.), \teiHi[rend={italic}]{The Idea and Practice of Philosophy in Gilbert Simondon}, Bâle, Schwabe, « Philosophia », 2023. }}. }
            \teiP[xmlid={p3-12}]{En pariant sur le fait que la philosophie de Gilbert Simondon est trop attachée à la technique pour dissoudre l’image photographique dans le régime ultra général du « tout image », et en pariant sur la manière dont il fait de l’image non un support, mais une opération, il nous a semblé que Chris Marker serait un excellent « technologue » simondonien, lui le « super-bricoleur\teiNote[xmlid={ftn84-4},n={3},place={foot}]{\teiP[xmlid={p-407}]{Julien Gester, « La seconde vie de Chris Marker », \teiHi[rend={italic}]{Les Inrockuptibles}, 29 avril 2008, \teiRef[target={https://www.lesinrocks.com/cinema/la-seconde-vie-de-chris-marker-13720-29-04-2008/}]{https://www.lesinrocks.com/cinema/la-seconde-vie-de-chris-marker-13720-29-04-2008/} (consulté le 1\teiHi[rend={sup}]{er} juillet 2025). }} ». Marker engrène des images les unes aux autres en jouant justement de leur écart technique (photographie, film, vidéo, ordinateur), créant une beauté vitale par résonance et rémanence des photographies les unes aux autres, refusant d’en faire une \teiHi[rend={italic}]{prise} du réel pour mieux en interroger la fugacité et la beauté transitive, assumant une \teiHi[rend={italic}]{déprise} par l’acte photographique qui n’est pourtant pas du tout un renoncement au monde tel qu’il est. Ouvrant très tôt la technique photographique à l’image numérique, il n’a de cesse d’interroger l’obsolescence programmée des images contemporaines. Un autre parti pris de notre lecture sera de considérer que, comme Kracauer dans \teiHi[rend={italic}]{Théorie du film}, il fait de l’image photographique le cœur incandescent du film, y revenant d’ailleurs constamment\teiNote[xmlid={ftn85-4},n={4},place={foot}]{\teiP[xmlid={p-408}]{Christophe Chazalon, « L’art du bricolage numérique », dans Christine Van Assche, Raymond Bellour et Jean-Michel Frodon (dir.), \teiHi[rend={italic}]{Chris Marker. Catalogue de l’exposition à la cinémathèque française}, Paris, La Cinémathèque française, 2018, p. 372-375, insistant sur un retour de la photographie dans l’œuvre après \teiHi[rend={italic}]{Chats perchés}. }}, ce pour quoi \teiHi[rend={italic}]{Sans soleil}\teiNote[xmlid={ftn86-4},n={5},place={foot}]{\teiP[xmlid={p-409}]{Chris Marker (réal.),\teiHi[rend={italic}]{ Sans Soleil} \lbrack{}1983\rbrack{}, DVD vidéo, Potemkine Film, 2020.}} peut être traité photographiquement. Le régime de visibilité qu’il y instaure sera une boussole poétique de réflexion technologique, car c’est dans ce film qu’il mêle la photographie analogique et numérique pour la première fois, jouant, au sein même de la photographie, des écarts techniques. Cela n’interdira évidemment pas d’ouvrir le spectre à d’autres œuvres. L’enjeu n’est une exégèse définitive ni de Marker ni de Simondon, mais une confrontation théorético-esthétique de l’un par l’autre. }
            \teiP[xmlid={p4-12}]{Il sera possible de cerner l’image photographique comme beauté vitale parce qu’étant artifice et artefact, « point-clé » au sens de Simondon. Pour cela, il faudra d’abord cerner un travail de l’image photographique saisie depuis son grain ou son pixel, comme modulation et concrétisation. C’est alors la vitalité des écarts entre images, techniquement différentes, qui porte et perpétue la beauté, bien plus que le flux indifférencié du visuel, du « tout image », « la » technique comme telle n’ayant aucun sens et ne se comprenant qu’au pluriel en ses points clés. Cette invention photographique techniquement perpétuée débouche sur une beauté vivante de la réparation, sans jamais se complaire dans la consolation. Sur chacun de ces points, Marker et Simondon nous semblent dialoguer\teiNote[xmlid={ftn87-4},n={6},place={foot}]{\teiP[xmlid={p-410}]{Nous remercions vivement Jean-Michel Frodon et Jamil Alioui pour leur relecture suggestive et acérée.}}. }
            \begin{teiDiv}[type={section1},xmlid={div01-8}]

              \teiHead[xmlid={le-grain-image-information-technique-001}]{Le grain : image/information/technique}
              \teiP[xmlid={p5-12}]{L’articulation entre technique et vitalité est au cœur de la pensée de Simondon : « on peut parler de \teiHi[rend={italic}]{vie technique}, comme étant ce qui réalise en l’homme cette mise en relation des deux fonctions. L’homme est capable d’assumer la relation entre le vivant qu’il est et la machine qu’il fabrique ; l’opération technique exige \teiHi[rend={italic}]{une vie technique et naturelle}\teiNote[xmlid={ftn88-4},n={7},place={foot}]{\teiP[xmlid={p-411}]{Gilbert Simondon, \teiHi[rend={italic}]{Du mode d’existence des objets techniques} \lbrack{}désormais \teiHi[rend={italic}]{MEOT}\rbrack{}, Paris, Aubier, 1989, p. 125, nous soulignons.}} ». Cette vie technique et naturelle engage une philosophie de la photographie comme philosophie de la genèse technique de l’image photographique, qui s’attache donc non à l’image une fois tirée et exposée devant nous, mais à celle qui se fait dans l’appareil\teiNote[xmlid={ftn89-4},n={8},place={foot}]{\teiP[xmlid={p-412}]{Voir Carole Maigné, « Philosophizing from and with the Camera », art. cité.}}. S’il est connu que Simondon exige que nous rentrions dans la machine pour mieux en décrypter le fonctionnement, ce serait réduire son propos que de se contenter de ce regard venu du dedans : le boîtier est lui-même pris dans un monde transductif et métastable. Une logique de l’information plus que de la forme guide son analyse : la photographie n’est ni une découpe du monde, ni une trace fixée de ce qui fut. Le boîtier n’enregistre pas ce qui est devant lui comme si un face à face entre appareil et monde pouvait faire sens : l’image photographique s’inscrit dans une logique de la relation et de l’individuation bien plus que dans une interrogation sur ce qu’elle « représente\teiNote[xmlid={ftn90-4},n={9},place={foot}]{\teiP[xmlid={p-413}]{Si l’on part de Simondon, la photographie comme ressemblance par contact, comme adhérence de l’image au réel ne tient plus. Si Barthes est congédié, ce n’est pas parce que l’image est muette, mais parce qu’elle se situe dans un système d’informations et de relations. }} ». }
              \teiP[xmlid={p6-12}]{C’est bien pourquoi Simondon pense l’image photographique depuis son grain, depuis sa texture, informative, car elle est opération. Contre l’hylémorphisme, on connaît bien ce passage où l’argile ne prend pas forme grâce à un moule qui lui donnerait forme : }
              \begin{teiCit}[xmlid={cit01-8}]

                \begin{teiQuote}
On ne peut pas dire que le moule donne forme ; c’est la terre qui prend forme selon le moule, parce qu’elle communique avec l’ouvrier. La positivité de cette prise de forme appartient à la terre et à l’ouvrier ; elle est cette résonance interne, le travail de cette résonance interne. Le moule intervient comme condition de fermeture, limite, arrêt d’expansion, direction de médiation. L’opération technique institue la résonance interne dans la matière prenant forme, au moyen de conditions énergétiques et de conditions topologiques\teiNote[xmlid={ftn91-4},n={10},place={foot}]{\teiP[xmlid={p-414}]{Gilbert Simondon, \teiHi[rend={italic}]{L’individuation à la lumière des notions de forme et d’information }\lbrack{}désormais \teiHi[rend={italic}]{ILFI}\rbrack{}, Paris, Millon, 2017, p. 45. }}.
\end{teiQuote}
              
\end{teiCit}
              \teiP[xmlid={p7-12}]{Ce qui vaut de l’argile vaut de l’image : la genèse de l’objet technique fait partie de son être, le processus technique implique une opération de transduction, où peu à peu s’individue un objet, liant de manière indissociable technique et vie par échange d’énergie\teiNote[xmlid={ftn92-4},n={11},place={foot}]{\teiP[xmlid={p-415}]{\teiHi[rend={italic}]{MEOT}, \teiHi[rend={italic}]{op. cit}., p. 20. }}. Il n’y a pas matière inerte et forme venue du dehors, il y a opération, ce qu’ignore l’hylémorphisme en effaçant la médiation opérative, « comme un homme qui reste à l’extérieur de l’atelier et ne considère que ce qui entre et ce qui en sort\teiNote[xmlid={ftn93-4},n={12},place={foot}]{\teiP[xmlid={p-416}]{Gilbert Simondon, \teiHi[rend={italic}]{ILFI}, \teiHi[rend={italic}]{op. cit}., p. 46.}} ». }
              \teiP[xmlid={p8-12}]{Contre un lieu commun qui assène qu’une image photographique est la révélation d’une image latente, Simondon pense une résonance « pendant la prise de forme\teiNote[xmlid={ftn94-4},n={13},place={foot}]{\teiP[xmlid={p-417}]{\teiHi[rend={italic}]{Ibid}., p. 47.}} ». Plaidant la modulation contre le moule\teiNote[xmlid={ftn95-4},n={14},place={foot}]{\teiP[xmlid={p-418}]{\teiHi[rend={italic}]{Ibid}.}}, l’image photographique devient modulation plutôt qu’empreinte. Contre toute une tradition qui fait de l’image photographique un embaumement, une trace de ce qui fut et n’est plus, voire de ce qui fut et est toujours déjà congédié\teiNote[xmlid={ftn96-4},n={15},place={foot}]{\teiP[xmlid={p-419}]{Voir par exemple Georges Didi-Huberman, \teiHi[rend={italic}]{La ressemblance par contact}, Paris, Éditions de Minuit, 2008. }}, Simondon réfléchit à une continuité entre image, support et réel, traversée par une vitalité qui conjugue sans jamais séparer technique et individuation du monde. Il fait résonner ensemble le chimique et l’optique, sans privilégier l’un sur l’autre, l’image photographique se cristallisant par transformation de proche et proche des micrograins chimiques sous l’action d’une lumière optique orientée par l’appareil, le déclic de la prise de vue étant un système de résonance porté par l’opération de la « prise\teiNote[xmlid={ftn97-4},n={16},place={foot}]{\teiP[xmlid={p-420}]{Gilbert Simondon, \teiHi[rend={italic}]{ILFI}, \teiHi[rend={italic}]{op. cit}.,p. 48. }} ». L’image photographique se conçoit donc comme des grains modulés : « une surface photographique, dans sa partie active, support des signaux, est constituée par une émulsion contenant une multitude de grains d’argent, primitivement sous forme de combinaison chimique\teiNote[xmlid={ftn98-4},n={17},place={foot}]{\teiP[xmlid={p-421}]{\teiHi[rend={italic}]{Ibid}., p. 236. }}. » Elle ne reproduit pas une forme, elle informe\teiNote[xmlid={ftn99-4},n={18},place={foot}]{\teiP[xmlid={p-422}]{Simondon, dans \teiHi[rend={italic}]{MEOT}, \teiHi[rend={italic}]{op. cit}., p. 123, oppose la mémoire vivante de l’homme et la mémoire de la machine pour constater que la machine ne conserve pas des formes, mais une traduction des formes au moyen d’un codage, ce que sont les grains d’argent de la pellicule photographique. }}. En enregistrant, elle « module alors point par point la lumière comme la modulaient les objets photographiés\teiNote[xmlid={ftn100-4},n={19},place={foot}]{\teiP[xmlid={p-423}]{Gilbert Simondon, \teiHi[rend={italic}]{ILFI}, \teiHi[rend={italic}]{op. cit}., p. 223.}} ». Et c’est bien parce qu’elle enregistre et module qu’elle n’est pas un automatisme, là aussi contre une confusion récurrente concernant l’appareil photographique et contre un \teiHi[rend={italic}]{topos} de la réflexion technique : Simondon considère l’automatisme comme le plus bas degré d’opérativité de la machine. }
              \teiP[xmlid={p9-11}]{Marker ne cesse de montrer l’image technique qui techniquement se fabrique. Ce n’est pas là une tautologie, c’est lutter contre la fascination du visuel, de l’image toute faite qui semble tout dire. Il met en scène constamment des mains qui agissent, qui font et défont l’image, avec des techniques variées. Dans le documentaire \teiHi[rend={italic}]{On vous parle de Prague. Le deuxième procès d’Artur London} (1971), sur le tournage\teiNote[xmlid={ftn101-4},n={20},place={foot}]{\teiP[xmlid={p-424}]{Ce qui dénonce l’image fabriquée de l’idéologie stalinienne lors du procès Slansky filmé, c’est une contre-image, celle du tournage de \teiHi[rend={italic}]{L’Aveu }de Costa Gavras, elle-même filmée par Marker qui rend compte de ce tournage. C’est donc un film de Marker sur le film de Costa Gavras en train de se faire et contre un film ancien mensonger, caméra filmant une caméra qui agit contre une caméra passée, et sûrement pas une absence de caméra pour mieux dire « authentiquement » le réel. }}, nous dit la voix off, il y a beaucoup de monde, beaucoup de mains, des « mains intelligentes qui rendent en fin de compte cette soi-disant industrie plus proche de la poterie que de l’automation ». Marker montre des câbles, des interrupteurs, la caméra et les mécaniciens de l’image pour signifier le cinéma qui se fait. Contre les \teiHi[rend={italic}]{topoi} de l’automatisme et de l’empreinte, Marker joue la mise en abyme de la fabrication de l’image par la machine\teiNote[xmlid={ftn102-4},n={21},place={foot}]{\teiP[xmlid={p-425}]{Les mains des mécanos de l’image se mêlent aux mains menottées de Montand. Il est clair qu’il s’agit ici de « démonter » un discours idéologique. Dans \teiHi[rend={italic}]{Le fond de l’air est rouge}, ce n’est pas un hasard si les deux chapitres se nomment « les mains fragiles » et « les mains coupées ». Je reste ici délibérément sur le plan de la philosophie de la technique. }}. C’est pourquoi il faut interroger le lieu d’émergence de l’image, sa genèse : }
              \begin{teiCit}[xmlid={cit02-7}]

                \begin{teiQuote}
Information. La séquence du pilote américain ne montre pas que les Américains napalmaient le Vietnam, mais un américain « se montrant » en train de napalmer le Vietnam. Le mode d’information fait partie de l’information et l’enrichit. C’est un des principes du choix des documents : chaque fois que c’était possible (écrans de télévision, lignes du kinescope, citations d’actualité, lettre enregistrée sur « mini-cassette », images tremblées, voix de radios, commentaires des images à la première personne par ceux qui les ont captées, rappel des conditions de tournage, caméra-clandestine, ciné-tract…), \teiHi[rend={italic}]{rapprocher le document des circonstances concrètes de son élaboration, faire en sorte que l’information n’apparaisse pas comme }cosa mentale\teiHi[rend={italic}]{, mais comme une matière – avec son grain, ses aspérités, quelquefois ses échardes}\teiNote[xmlid={ftn103-4},n={22},place={foot}]{\teiP[xmlid={p-426}]{Chris Marker, \teiHi[rend={italic}]{Le Fond de l’air est rouge}, Paris, François Maspero, 1978, p. 10, nous soulignons. }}.
\end{teiQuote}
              
\end{teiCit}
              \teiP[xmlid={p10-11}]{Ce travail du grain est décisif et va croissant, à mesure que Marker s’interroge certes sur l’histoire qui se fait, se raconte, se transmet, mais de plus en plus sur les modalités de production de l’image. Alors que la poterie nous renvoyait tout droit à l’hylémorphisme que dénonce Simondon, Marker s’avère de plus en « technologue » et pas seulement « technophile\teiNote[xmlid={ftn104-4},n={23},place={foot}]{\teiP[xmlid={p-427}]{Alice Leroy et Gabriel Bortzmeyer, « Raymond Bellour (2018) », entretien, \teiHi[rend={italic}]{Débordements}, 25 juillet 2018, \teiHi[rend={underline}]{https://www.debordements.fr/Raymond-Bellour-2018} (consulté le 1\teiHi[rend={sup}]{er} juillet 2025).}} » : il questionne toujours davantage la poétique de l’image technique, il n’est pas seulement dans la machine, il invente avec elle, il vit avec elle, au point de se construire une vie en elle (voir, sur Second Life, le site Gorgomancy.net). Le terme d’information nous semble changer peu à peu de sens, signifiant toujours plus le branchement des images techniques au monde et non le seul contenu politique et historique, au sens des « actualités », des « faits », des « récits », des « combats » (ce qui ne veut pas dire que le sens problématique de l’histoire est pour autant évacué, il ne l’est pas !). En son sens élargi, l’information inclut un rapport au monde qui ne cesse de se modifier, une labilité des contenus qui touche à la vie dans ce qu’elle a de plus modeste, de plus ténu\teiNote[xmlid={ftn105-4},n={24},place={foot}]{\teiP[xmlid={p-428}]{Marker semble reprendre de plus en plus à son compte cette analyse qu’il a faite de Giraudoux : « Depuis le début du siècle, à travers les poèmes-conversations d’apollinaire, les journaux collés de Picasso et de Braque, les calques de Max Ernst, jusqu’aux mobiles de Calder et aux tic-toc-chocs de Mac Laren en passant par d’autres moins prévus, se déroule un complot qui est probablement la seule entreprise réussie de notre temps, mais qui le rachète amplement, et qui consiste à relever les choses les plus humbles du dédain où les abandonnait l’art des époques égoïstes, humanistes et mégalomanes. Tout ce qui se cachait à l’ombre des statues et des temples vient à la lumière, la rédemption s’étend à la création tout entière, on convie au même repas l’or et le plomb (avec une préférence pour le plomb, on sauve de l’indifférence et du dégoût le caillou, le graffito, la flaque, la tâche, le clou » (Chris Marker, \teiHi[rend={italic}]{Giraudoux par lui-même}, Paris, Éditions du Seuil, 1952, p. 26 et 27). \teiHi[rend={italic}]{Sans soleil} insiste sur la « banalité » qu’il a traquée comme un chasseur de prime. }}. Cette capacité informative de l’appareil n’est pas d’engranger des contenus figés, mais de créer un ordre dans le chaos du monde, un ordre appareillé, parce que notre rapport au monde est intimement lié aux machines : }
              \begin{teiCit}[xmlid={cit03-7}]

                \begin{teiQuote}
La machine, œuvre d’organisation, d’information, est, comme la vie et avec la vie, ce qui s’oppose au désordre, au nivellement de toutes choses tendant à priver l’univers de pouvoirs de changement. La machine est ce par quoi l’homme s’oppose à la mort de l’univers ; elle ralentit, comme la vie, la dégradation de l’énergie et devient stabilisatrice du monde\teiNote[xmlid={ftn106-4},n={25},place={foot}]{\teiP[xmlid={p-429}]{\teiHi[rend={italic}]{MEOT}, \teiHi[rend={italic}]{op. cit}., p. 15. }}. 
\end{teiQuote}
              
\end{teiCit}
              \teiP[xmlid={p11-11}]{Le geste à la fin de \teiHi[rend={italic}]{Sans soleil} de planter des fiches électroniques pour faire circuler l’énergie informative est ici décisif tout autant qu’il est novateur, Marker étant un des premiers, si ce n’est le premier, à introduire l’image électronique dans ses œuvres\teiNote[xmlid={ftn107-4},n={26},place={foot}]{\teiP[xmlid={p-430}]{Je remercie Jean-Michel Frodon pour cette information et plus largement pour m’avoir recommandé \teiHi[rend={italic}]{Sans Soleil} dans ce parcours Simondon/Marker. }}. En fichant, la main opère plus qu’elle ne façonne, ce qui est là à notre avis le cœur de la technicité simondonienne. Ce geste induit une réflexion théorique sur la nature technique de l’image photographique, car la main qui opère est la même que celle qui opérait le déclic. La technique change certes, mais la photographie demeure même si elle se modifie. La texture depuis le grain modulé ne crée pas de rupture \teiHi[rend={italic}]{de jure} entre argentique et numérique chez Simondon, et Marker tranche \teiHi[rend={italic}]{de facto} le débat entre les deux modalités de production de l’image photographique. \teiHi[rend={italic}]{Sans soleil} propose une continuité dans la dissemblance entre l’aile de l’avion électronique de Pearl Harbour et le vol d’avion filmé. Toute la phénoménalité du monde est traversée par la technique, il n’y a pas de monde humain sans machine, la technique ne désigne pas un après inauthentique à quoi s’opposerait une virginité a-technique : l’aile de l’avion qui passe du vol vu, photographié, filmé au vol simulé, et inversement, constitue ce qui pour nous fait réalité. }
              \teiP[xmlid={p12-11}]{\teiHi[rend={italic}]{Sans soleil} convoque ainsi des mains opérantes (celles des Japonais qui dansent, qui prient, qui manient le bâton en ronde… mais aussi les pattes levées des chats \teiHi[rend={italic}]{maneki neko}) qui font et défont l’image photographique. Le grain s’affiche aussi dans le travail sur le flou, puisque Marker fait du fixe avec du flou et inversement\teiNote[xmlid={ftn108-4},n={27},place={foot}]{\teiP[xmlid={p-431}]{Vincent Jacques, « Images floues, frontières troubles. La pratique photographique de Chris Marker », dans Vincent Jacques (dir.), \teiHi[rend={italic}]{Chris Marker. Photographie}, Grâne, Creaphis éditions, 2018, p. 86. }}, exposant la trame de l’image puisque l’image est matière\teiNote[xmlid={ftn109-4},n={28},place={foot}]{\teiP[xmlid={p-432}]{\teiHi[rend={italic}]{Ibid}., p. 99 : « réfléchir sur le médium en tant que tel, une habitude chez Marker : mettre en avant le support, l’image est aussi matière chez lui ».}}. Et cette matière du grain est vitale : une superbe séquence, d’autant plus marquante qu’elle est simple, montre une eau de pluie ruisselant sur une vitre comme ruisselait l’image électronique quelques instants auparavant. Le grain fait le lien entre l’organique et le minéral. À l’inverse, c’est aussi une séquence terrible de \teiHi[rend={italic}]{Level 5} où le ralenti excessif de la femme sautant de la falaise à Okinawa, lors du suicide collectif d’une partie des habitants, transforme l’image argentique en image quasi électronique et pixelisée, dans un granuleux qui affole le regard du spectateur en écho à celui terrifié de la femme. Le granuleux s’exhibe aussi quand il s’agit de photographie de photographie ou de film de film, qui renversent le rapport entre regard et appareil : on est regardé par la télévision japonaise plus qu’on ne la regarde, nous dit Marker, exhibant alors des visages en gros plan, qui font vaciller la limite entre le réel granuleux et l’imaginaire flou, ou plutôt, en termes simondoniens, entre le réel et le potentiel. Or, « le potentiel est une des formes du réel, aussi complète que l’actuel\teiNote[xmlid={ftn110-4},n={29},place={foot}]{\teiP[xmlid={p-433}]{\teiHi[rend={italic}]{MEOT}, \teiHi[rend={italic}]{op. cit}., p. 155. }} ». \teiHi[rend={italic}]{Zapping zone (proposary for an imaginary television)}, de 1990, métisse les genres de l’image technique\teiNote[xmlid={ftn111-4},n={30},place={foot}]{\teiP[xmlid={p-434}]{Voir l’article de Jean-Michel Frodon, \teiRef[target={https://aoc.media/critique/2021/11/03/reapparitions-et-avatars-de-zapping-zone-installation-fantome-de-chris-marker/}]{https://aoc.media/critique/2021/11/03/reapparitions-et-avatars-de-zapping-zone-installation-fantome-de-chris-marker/} (consulté le 3 septembre 2025).}}. La simple liste des techniques qui sont comme tissées entre elles est en soi une démarche de technologue : un chat \teiHi[rend={italic}]{maneki neko} à l’entrée, des images fixes (photographies, photogrammes, diapositives), des images en mouvement qui reprennent souvent des séquences d’anciens films, des collages, sept ordinateurs et treize moniteurs, des sons de tous types (naturels, enregistrés, électroniques), le tout organisé par arborescence en « zones » (« séquences », « Bosnie », « photos », « théorie des ensembles », etc.\teiNote[xmlid={ftn112-4},n={31},place={foot}]{\teiP[xmlid={p-435}]{Je reprends des éléments de la description de Christine Van Assche : « \teiHi[rend={italic}]{Zapping zone (Proposals for an imaginary television)}, Le métissage des genres », dans Christine Van Assche, Raymond Bellour et Jean-Michel Frodon (dir.), \teiHi[rend={italic}]{Chris Marker},\teiHi[rend={italic}]{ op. cit.}, p. 330-337. Il serait fécond de prolonger ici la réflexion sur la vidéo et la télévision à partir des études de Raymond Bellour, \teiHi[rend={italic}]{L’entre-images. Photo. Cinéma. Vidéo}, Paris, La Différence, 2002. La vidéo est selon Bellour une « passeuse » entre images, elle provoque recouvrements et configurations inédites. « L’entre-images » qu’elle provoque est analysée depuis l’esthétique du cinéma et de la perception des images alors que nous essayons ici de proposer le prisme d’une philosophie de la technique et une critique du réalisme supposé mortifère de l’image photographique (voir les analyses sur \teiHi[rend={italic}]{La Jetée}, p. 146-148). Il reste que la thèse d’un arrêt sur image, et donc d’un retour du photographique, qui se développe au tournant des années 1960, quand se développent la vidéo et l’image électronique, est passionnante. }}). }
              \teiP[xmlid={p13-10}]{Ces mains opérantes qui se retrouvent sur écran ou sur papier (les livres composés et imprimés de Marker) sont aussi celles de Marker à l’ouvrage. Car le grain est le grain que l’on manipule, que l’on choisit, en déléguant le moins possible l’acte technique qui l’informe : faire primer le grain contre l’automatisme impose de maîtriser la machine à enregistrer. Simondon raconte une contre-histoire de la photographie comme régression : l’histoire de la photographie est celle d’une dislocation de l’acte photographique en des instances séparées (boîtier, film, développement), absorbées par « l’univers industriel concentrationnaire » du marché\teiNote[xmlid={ftn113-4},n={32},place={foot}]{\teiP[xmlid={p-436}]{Gilbert Simondon, \teiHi[rend={italic}]{Imagination et invention, 1965-1966}, Paris, PUF, 2014, p. 287. }} qui s’emploie à compartimenter les tâches et destituer l’importance du photographe. Marker semble avoir constamment joué \teiHi[rend={italic}]{contre} cette contre-histoire : que ce soit dans le train-cinéma de Medvedkine, où le wagon concentrait toutes les opérations techniques dans sa boîte mobile, retrouvant en quelque sorte la caméra polyvalente des frères Lumières, celle des origines du cinéma\teiNote[xmlid={ftn114-4},n={33},place={foot}]{\teiP[xmlid={p-437}]{Je remercie Agnès Devictor pour cette référence : il y a un parallèle à faire entre la chambre de Daguerre, la caméra des premiers temps des Lumières, le ciné-train de Medvedkine vu par Marker et le Polaroïd vu par Simondon : les deux derniers rejouent leur origine, décalée dans le temps, une rétro-action historique et inventive et donc émancipatrice. }}, ou que ce soit dans \teiHi[rend={italic}]{Chats perchés} en 2004 : « pouvoir faire tout un film, \teiHi[rend={italic}]{Chats perchés} (2004) avec mes dix doigts sans aucun appui ni intervention extérieurs… et ensuite aller vendre moi-même le DVD que j’ai enregistré à la braderie de Saint-Blaise… là j’avoue que j’ai eu un sentiment de triomphe : du producteur au consommateur, direct. Pas de plus-value. J’avais accompli le rêve de Marx\teiNote[xmlid={ftn115-4},n={34},place={foot}]{\teiP[xmlid={p-438}]{Julien Gester, art. cité. }}. »}
            
\end{teiDiv}
            \begin{teiDiv}[type={section1},xmlid={div02-8}]

              \teiHead[xmlid={linvention-photographique-perpetuee-le-point-cle-001}]{L’invention photographique perpétuée : le point clé }
              \teiP[xmlid={p14-10}]{« La présence de l’homme aux machines est une invention perpétuée. Ce qui réside dans les machines, c’est de la réalité humaine, du geste humain fixé et cristallisé en structures qui fonctionnent. Ces structures ont besoin d’être soutenues au cours de leur fonctionnement, et la plus grande perfection coïncide avec la plus grande ouverture, avec la plus grande liberté du fonctionnement\teiNote[xmlid={ftn116-4},n={35},place={foot}]{\teiP[xmlid={p-439}]{\teiHi[rend={italic}]{MEOT}, \teiHi[rend={italic}]{op. cit}., p. 12.}}. » Il faut évidemment prendre le syntagme « invention perpétuée » de Simondon au sens fort. Il induit une poétique de l’image technique qui se concentre dans la notion de « point clé », quelque peu énigmatique de prime abord, mais féconde. }
              \teiP[xmlid={p15-10}]{Le grain est actif, car il est technicisé. Simondon ouvre le réel au potentiel, ouvre l’appareil au devenir, conteste l’achèvement de l’acte technique : « la machine est enfermée dans cette vision réductrice qui la considère comme achevée en elle-même et parfaite, qui la fait coïncider avec son état actuel, avec ses déterminations matérielles. Envers l’objet d’art, une pareille attitude consisterait à réduire un tableau à une certaine étendue de peinture séchée et fendillée sur une toile tendue\teiNote[xmlid={ftn117-4},n={36},place={foot}]{\teiP[xmlid={p-440}]{\teiHi[rend={italic}]{MEOT}, \teiHi[rend={italic}]{op. cit.}, p. 146. }}. » Moduler le grain, c’est refuser de stabiliser l’image comme image à l’arrêt : Marker fait de la photographie un film, en jouant sur la potentialité et la densité de l’image (\teiHi[rend={italic}]{La Jetée}, constituée de photographies étant ici exemplaire de ce procédé photographique « filmique »). Parce qu’elle est opératoire, l’image photographique ne cesse jamais son activité, non au sens d’un mystère insondable ou d’un pouvoir transcendant, mais parce que le geste technique lui-même ne s’achève pas. La photographie brûle quand elle est à l’arrêt, quand le potentiel s’efface : « Un instant arrêté grillerait comme l’image d’un film bloqué devant la fournaise du projecteur », nous dit \teiHi[rend={italic}]{Sans soleil}. Marker n’a ainsi cessé de jouer et de prolonger ses machines : Raymond Bellour raconte que lors de l’exposition \teiHi[rend={italic}]{Passages de l’image} au centre Pompidou, Marker ne jurait que par son Apple 2C, « tellement augmenté qu’il était devenu un monstre proliférant\teiNote[xmlid={ftn118-4},n={37},place={foot}]{\teiP[xmlid={p-441}]{Albane Penaranda, « Chris Marker, l’écrivain cinéaste », entretiens avec Raymond Bellour et Christine Van Assche, \teiHi[rend={italic}]{France Culture}, épisode 8, \teiHi[rend={underline}]{https://www.radiofrance.fr/franceculture/podcasts/les-nuits-de-france-culture/nuitchris-marker-entretien-avec-raymond-bellour-et-christine-van-assche-5775796} (consulté le 1\teiHi[rend={sup}]{er} juillet 2025).}} », il n’a cessé de collectionner les ordinateurs sa vie durant… }
              \teiP[xmlid={p16-10}]{Le devenir le plus profond, le plus vital, c’est ce que Simondon appelle un déphasage permanent\teiNote[xmlid={ftn119-4},n={38},place={foot}]{\teiP[xmlid={p-442}]{\teiHi[rend={italic}]{ILFI}, \teiHi[rend={italic}]{op. cit}., Note complémentaire, p. 341 : « L’activité constructive donne à l’homme l’image réelle de son acte, par ce que ce qui est actuellement objet de la construction devient moyen d’une construction ultérieure, grâce à une \teiHi[rend={italic}]{permanente} médiatisation » (nous soulignons). }} : c’est notre mode d’être vital au monde qui ne cesse de se déstabiliser, de se saturer et de décompenser, de créer de la forme sur un fond pour la moduler. Or, Marker ne cesse de déphaser, non seulement les images et le réel, mais aussi les images techniques les unes aux autres. S’il est un « as du montage » (\teiHi[rend={italic}]{Level 5}), il semble aussi jouer sur autre chose : il y a certes la découpe, le choc, la confrontation, l’assemblage, mais aussi des effets de différenciation, de déprise, une présence du fond qui ne se réduit pas à l’image agencée. L’image saisie depuis son grain est photographie mise en réseau avec d’autres types d’images elles-mêmes granuleuses et informatives, toutes autant de potentiels d’images : dessinée, imprimée, pixelisée. Le grain imprimé d’une gravure d’enfance se module en grain photographique puis en grain filmique\teiNote[xmlid={ftn120-4},n={39},place={foot}]{\teiP[xmlid={p-443}]{Guy Gauthier, « Images d’enfance », dans Philippe Dubois (dir.), \teiHi[rend={italic}]{Théorème. Recherches sur Chris Marker}, Paris, Presses Sorbonne Nouvelle, 2002, p. 46-59. }}. Simondon nous disait ci-dessus : « l’opération technique institue la résonance interne dans la matière prenant forme, au moyen de conditions énergétiques et de conditions topologiques\teiNote[xmlid={ftn121-4},n={40},place={foot}]{\teiP[xmlid={p-444}]{\teiHi[rend={italic}]{ILFI}, \teiHi[rend={italic}]{op. cit}., p. 45.}}. » Marker branche constamment des images techniques les unes sur les autres, crée des résonances inédites entre elles, travaille des topologies (le « Japon », le « dépays », le « chat », la « chouette » qui ne sont pas que des lieux, mais des nœuds informatifs) et des énergies (saturer, fluidifier, contraster, etc.). Car ce n’est pas seulement une thématique qui les relie, aussi puissante que soit celle de la mémoire et du temps dans l’œuvre de Marker, c’est aussi leur caractère opératoire. La vitalité de l’image technique n’en finit pas, car technique et vie ne se dissocient pas. }
              \teiP[xmlid={p17-10}]{« La compatibilité réside dans la mise en suspens de l’activité chimique du matériau sensible entre la fabrication et le développement, ce qui permet d’insérer la prise de vue dans cet intervalle temporel\teiNote[xmlid={ftn122-4},n={41},place={foot}]{\teiP[xmlid={p-445}]{\teiHi[rend={italic}]{Imagination et invention}, \teiHi[rend={italic}]{op. cit}., p. 286. }}. » La portée de l’invention réside dans le suspens du temps (chimique) plus que dans le \teiHi[rend={italic}]{graphein} et son lien à la lumière (\teiHi[rend={italic}]{photo}\teiNote[xmlid={ftn123-4},n={42},place={foot}]{\teiP[xmlid={p-446}]{Carole Maigné, art. cité. }}) : suspendue et différée, l’action chimique perd justement son caractère impérieux et immédiat. Simondon qualifie la photographie non pas comme « révélation » de l’image latente, mais comme insertion d’un geste « entre-temps » : il ne s’agit pas tant de révéler que d’insérer « juste à temps\teiNote[xmlid={ftn124-4},n={43},place={foot}]{\teiP[xmlid={p-447}]{C’est là l’importance de ce que Simondon thématise dans le supplément « Allagmatique » de \teiHi[rend={italic}]{L’individuation}, \teiHi[rend={italic}]{op. cit}. : « l’opération est un \teiHi[rend={italic}]{metaxu} entre deux structures et est pourtant d’une autre nature que toute structure », p. 531. Ce qui est « intervalle », « au milieu de », est caractéristique de tout acte opératoire. }} ». Il se dégage alors une temporalité particulière : non pas celle de l’œuvre comme \teiHi[rend={italic}]{poiesis}, création \teiHi[rend={italic}]{ex nihilo}, ce qui reviendrait à promouvoir l’hylémorphisme ou encore le démiurge. Simondon propose en quelque sorte une poétique sans \teiHi[rend={italic}]{poiesis}, une poétique de l’image qui fait de la technicité un foyer innovant et non une instrumentation. Le temps de l’image photographique est une suspension du temps dans l’instant : « pour la photographie, il s’agit du processus physique de formation d’image réelle et du processus photochimique, compatibilisés par le phénomène de l’image latente ; cette compatibilité appartient à la catégorie des états d’équilibre, autorisant la succession temporelle des phases par la mise en suspens d’une activité\teiNote[xmlid={ftn125-4},n={44},place={foot}]{\teiP[xmlid={p-448}]{\teiHi[rend={italic}]{Imagination et invention}, \teiHi[rend={italic}]{op. cit}., p. 287. }} ». L’instantané de la prise n’est pas conjuguée au présent, il est pris dans un devenir en équilibre qui tient ensemble toutes les dimensions du temps (passé, présent, futur), il constitue un moment opératoire (\teiHi[rend={italic}]{metaxu}) qui ouvre les possibles : l’intervalle suspendu permet la bifurcation, l’insertion \teiHi[rend={italic}]{dans} et \teiHi[rend={italic}]{par} le temps de l’image d’un aléa fécond. Surtout, cette suspension du temps permet d’insérer quelque chose qui ne pourrait être sans l’appareil : il faut l’artifice technique de la suspension pour que l’image photographique se fasse. }
              \teiP[xmlid={p18-10}]{L’art est inventif, car il crée de la convergence, contre l’usage, l’ignorance, l’automatisme : }
              \begin{teiCit}[xmlid={cit04-6}]

                \begin{teiQuote}
L’art est une réaction profonde contre la perte de signification et de rattachement à l’ensemble de l’être dans sa destinée ; il n’est pas ou ne doit pas être compensation, réalité advenant après coup, mais au contraire unité primitive, préface à un développement selon l’unité ; l’art annonce, préfigure, introduit, ou achève, mais ne réalise pas : il est l’inspiration profonde et unitaire qui amorce et consacre\teiNote[xmlid={ftn126-4},n={45},place={foot}]{\teiP[xmlid={p-449}]{\teiHi[rend={italic}]{MEOT}, \teiHi[rend={italic}]{op. cit}., p. 200.}}.
\end{teiQuote}
              
\end{teiCit}
              \teiP[xmlid={p19-10}]{Cette citation est ici décisive : créer du sens n’est pas tant faire œuvre, au sens où l’œuvre serait une clôture sur soi, que placer un geste, ouvert et sans terme indispensable. « Il n’y a pas de centre de l’acte, il n’y a pas de limites de l’acte. Chaque acte est centré, mais infini ; la valeur d’un acte est sa largeur, sa capacité d’étalement transductif\teiNote[xmlid={ftn127-4},n={46},place={foot}]{\teiP[xmlid={p-450}]{\teiHi[rend={italic}]{ILFI},\teiHi[rend={italic}]{ op. cit}., p. 324. }}. » Le geste n’est ni anodin ni sûr de sa finalité : il ne s’agit ni d’agir pour agir, ni de faire pour faire, mais de mettre en branle et d’orienter. Le schéma n’est pas ici dialectique. Il ne s’agit pas de poser une contradiction et de la résoudre en la dépassant, il n’y a pas de succession téléologisée. Comme le dit clairement Simondon, le déphasage de l’être et du monde n’est pas dialectique, car il est en quelque sorte scission du dedans, et en cela vitale. Il est « dédoublement » interminable, différenciation constante, car chaque phase n’est qu’un point de bascule pour une autre. En ce sens, l’art ne compense rien du monde, ne résout rien non plus, n’a pas de parole définitive à délivrer, parce qu’il n’arrête rien de la marche des choses et qu’il ne fait pas système. Par contre, il achève et consacre, en créant des points clés : des nœuds de sens, des lieux de circulation impossibles à fermer, ce que Simondon appelle aussi des « réticulations ». }
              \teiP[xmlid={p20-10}]{Les points de réticulation sont des points d’insertion de l’action humaine dans le monde\teiNote[xmlid={ftn128-3},n={47},place={foot}]{\teiP[xmlid={p-451}]{\teiHi[rend={italic}]{MEOT}, \teiHi[rend={italic}]{op. cit}., p. 165. }}. Ce qui importe est leur réversibilité et leur pluralité : ils créent quelque chose d’indissociable entre homme et monde, des chemins qui se parcourent en deux sens, créant ce que Simondon appelle un « fond continu\teiNote[xmlid={ftn129-3},n={48},place={foot}]{\teiP[xmlid={p-452}]{\teiHi[rend={italic}]{Ibid}., p. 167. }} », on pourrait dire une basse continue si on prenait une métaphore musicale. Car ce que Simondon a en vue, c’est l’articulation forme/fond et le geste même de figurer : ce déphasage où le geste s’insère dans le monde et crée une articulation, un point clé ; il menace toujours de se détacher et devenir instrumental, mais invite aussi à en créer un autre\teiNote[xmlid={ftn130-3},n={49},place={foot}]{\teiP[xmlid={p-453}]{\teiHi[rend={italic}]{Ibid}., p. 168. }}. Un point clé objective et subjective à la fois\teiNote[xmlid={ftn131-3},n={50},place={foot}]{\teiP[xmlid={p-454}]{\teiHi[rend={italic}]{Ibid}.}}, menace toujours de devenir abstraction, de fragmenter la réticulation, de morceler ce qui est joint. L’outil est justement une forme de détachement, ce moment où l’articulation devient une action répétable, ajustée, libérée du fond du monde pour servir partout et en tous lieux\teiNote[xmlid={ftn132-3},n={51},place={foot}]{\teiP[xmlid={p-455}]{\teiHi[rend={italic}]{Ibid}., p. 170. }}. La pensée esthétique recrée un sentiment de totalité, d’unification : « L’œuvre d’art refait un univers réticulaire au moins pour la perception », et c’est par son intensité qu’elle devient « un nœud de la réalité éprouvée\teiNote[xmlid={ftn133-3},n={52},place={foot}]{\teiP[xmlid={p-456}]{\teiHi[rend={italic}]{Ibid}., p. 180. }} ». Et parce qu’elle est « point clé », elle reste dans cet intervalle entre subjectivation et objectivation. L’intervalle importe : Simondon considère que c’est là que se « concrétise » un monde intermédiaire, un monde qui justement articule l’homme et le monde\teiNote[xmlid={ftn134-3},n={53},place={foot}]{\teiP[xmlid={p-457}]{\teiHi[rend={italic}]{Ibid}., p. 182. }}. On pourrait dire que l’art crée donc un « entremonde », une réalité insérée, dont la beauté vient justement de ce caractère ni fusionné ni détaché. Le déphasage y est donc en quelque sorte contenu, maintenu et voulu dans une tension, une polarité entre sujet et objet sans jamais céder à l’un ni à l’autre, et sans jamais rêver de leur fusion. }
              \teiP[xmlid={p21-8}]{Ainsi, « c’est bien l’insertion qui définit l’objet esthétique, et non l’imitation : un morceau de musique qui imite des bruits ne peut s’insérer dans le monde, parce qu’il remplace certains éléments de l’univers (par exemple le bruit de la mer) au lieu de les compléter\teiNote[xmlid={ftn135-3},n={54},place={foot}]{\teiP[xmlid={p-458}]{\teiHi[rend={italic}]{Ibid}., p. 184. }} ». L’art ne saurait être une anthropologisation du monde, comme l’enfant qui se mire dans un étang chez Hegel, il n’est pas non plus une esthétisation de la technique, il est un travail de l’écart, de l’intervalle, une manière de « faire bourgeonner le monde » en multipliant les points clés. La beauté vient justement de la différenciation, de l’impossibilité d’unifier. L’insertion récuse une esthétique de l’un, du tout, elle promeut le couplage, le câblage, le branchement. Décrivant un central téléphonique, Simondon précise : « Une lumière, c’est une attente, une intention, un désir, une nouvelle imminente, une sonnerie qu’on n’entendra pas, mais qui va retentir au loin dans une autre maison\teiNote[xmlid={ftn136-3},n={55},place={foot}]{\teiP[xmlid={p-459}]{\teiHi[rend={italic}]{Ibid}., p. 187. }}. » C’est là que réside ce qu’il nomme l’achèvement. Le central est beau, le sémaphore est beau, on pense à Fernand Léger, mais aussi au pont transbordeur de Marseille photographié par Germaine Krull, et plus encore à son œuvre extraordinaire, \teiHi[rend={italic}]{Metal}, de 1928. }
              \teiP[xmlid={p22-8}]{Marker réfléchit comme Simondon à la photographie comme enregistrement et non comme automatisme, ce pour quoi résonance et concrétisation, termes si cruciaux dans la démonstration de Simondon, font étonnamment sens ici. La concrétisation concerne chez Simondon la manière dont un objet technique devient concret, ce qui veut dire qu’il ne l’est pas au départ, il est d’ailleurs une « unité en devenir\teiNote[xmlid={ftn137-3},n={56},place={foot}]{\teiP[xmlid={p-460}]{\teiHi[rend={italic}]{Ibid}., p. 20. }} », devenir concret qui accentue la « résonance interne », à savoir l’interdépendance fonctionnelle de ses parties. Ceci concerne \teiHi[rend={italic}]{stricto sensu} l’objet technique, mais non pas l’objet artistique. Néanmoins, le caractère fonctionnel peut se faire beau, et plus encore le devenir concret technique est inachevable : « l’objet technique n’est jamais complètement connu ; pour cette raison même, il n’est jamais non plus complètement concret\teiNote[xmlid={ftn138-3},n={57},place={foot}]{\teiP[xmlid={p-461}]{\teiHi[rend={italic}]{Ibid}., p. 35. }} ». Il y a là peut-être une indécision entre technique et art. La création de la « zone Tarkowski » dans \teiHi[rend={italic}]{Sans soleil} étonne, comme si se trouvait là aussi une traduction artistique et esthétique de la philosophie technique de Simondon. Car au fond, ce que crée Marker est un « milieu associé » au sens simondonien, à la fois technique et naturel, artificiel et naturel (les deux aspects s’imbriquent toujours), produit d’une imagination créatrice ici esthétique, milieu qui tient sa consistance de milieu parce qu’il s’insère, s’impose, reconfigure les temps, opère un saut fécond, « un conditionnement du présent par l’avenir, par ce qui n’est pas encore\teiNote[xmlid={ftn139-3},n={58},place={foot}]{\teiP[xmlid={p-462}]{\teiHi[rend={italic}]{Ibid}., p. 57. }} ». Si « inventer, c’est faire fonctionner sa pensée comme pourra fonctionner une machine\teiNote[xmlid={ftn140-3},n={59},place={foot}]{\teiP[xmlid={p-463}]{\teiHi[rend={italic}]{Ibid}., p. 138. }} », sans se soumettre à une causalité ni une finalité préétablies, en s’appuyant sur des « formes fonctionnantes », Marker explore lui aussi une invention machinique qui n’est ni machinale, ni automatisée. Travailler la machine, c’est inventer avec elle et par elle. Marker précise : « mon ami Hayao Yamaneko a trouvé une solution : si les images du présent ne changent pas, changer les images du passé… il m’a montré les bagarres des sixties traitées par son synthétiseur. Des images moins menteuses, dit-il avec la conviction des fanatiques, que celles que tu vois à la télévision. \teiHi[rend={italic}]{Au moins elles se donnent pour ce qu’elles sont, des images, pas la forme transportable et compacte d’une réalité déjà inaccessible}. Hayao appelle le monde de sa machine : la zone – en hommage à Tarkovski » (\teiHi[rend={italic}]{Sans soleil}, nous soulignons). L’image numérique que Marker insère est innovante, car elle décompose et recompose, elle crée des points clés sur un fond de réalité devenu opaque à force d’être vu. Les luttes mille fois photographiées de Sanziruka en 1968 autour de l’aéroport de Naruta doivent trouver une forme nouvelle, comme les images de Pearl Harbour au cinéma. Mais outre le contenu politique, le grain numérique relance artistiquement le grain argentique, la zone joue sur leur écart, défait et refait l’image, casse la compacité phénoménale en réinstallant le spectateur dans l’image fabriquée, en train de se faire. }
            
\end{teiDiv}
            \begin{teiDiv}[type={section1},xmlid={div03-7}]

              \teiHead[xmlid={beaute-vitale-et-reparation-001}]{Beauté vitale et réparation }
              \teiP[xmlid={p23-8}]{Puisque la transduction est une opération de proche en proche, transposée à la photographie, cela signifie que toute image photographique résonne avec une autre : l’irruption de la photographie dans le champ sensible relève par sa technicité d’une « structure réticulaire amplifiante\teiNote[xmlid={ftn141-3},n={60},place={foot}]{\teiP[xmlid={p-464}]{\teiHi[rend={italic}]{ILFI}, \teiHi[rend={italic}]{op. cit.}, p. 33. }} ». Elle l’est d’autant plus qu’elle est transindividuelle, l’usage démultiplié de cette image atteste qu’une nouvelle forme perceptive s’ancre : elle s’insère, et donc s’invite, s’impose et circule. Si l’on prend acte que « ce qui réside dans les machines, c’est de la réalité humaine, du geste humain fixé et cristallisé en structures qui fonctionnent\teiNote[xmlid={ftn142-3},n={61},place={foot}]{\teiP[xmlid={p-465}]{\teiHi[rend={italic}]{MEOT}, \teiHi[rend={italic}]{op. cit}., p. 12.}} », la dimension transductive et métastable de l’image-technique en fait peut-être alors, ce sera l’objet de ce troisième moment de notre réflexion, une réparation. L’invention technique étant acquisition pour toujours et pour tous, l’artifice et le naturel compâtissent pour sauver le monde phénoménal : « la prise de conscience ne suffit pas, car les organismes n’ont pas seulement une structure connaissable, ils tendent et se développent. C’est une tâche philosophique, psychologique, sociale, de \teiHi[rend={italic}]{sauver les phénomènes} en les réinstallant dans le devenir, en les remettant en invention, par l’approfondissement de l’image qu’ils recèlent\teiNote[xmlid={ftn143-3},n={62},place={foot}]{\teiP[xmlid={p-466}]{\teiHi[rend={italic}]{Imagination et invention}, \teiHi[rend={italic}]{op. cit}., p. 14. Nous soulignons. }}. » En quoi s’agit-il de sauver les phénomènes ? Comment réparer le monde sans l’illusion d’une consolation ? }
              \teiP[xmlid={p24-8}]{Simondon topologise l’art à l’horizontale, prolongeant esthétiquement et artistiquement le modèle de la transduction : « l’art desserre les liens de l’eccéité ; il multiplie l’eccéité, donnant à l’identité le pouvoir de se répéter sans cesser d’être identité\teiNote[xmlid={ftn144-2},n={63},place={foot}]{\teiP[xmlid={p-467}]{\teiHi[rend={italic}]{MEOT}, \teiHi[rend={italic}]{op. cit}., p. 200. }}. » En ce sens, l’art sauve les phénomènes, et nous avec, car il fait que « toute réalité, singulière dans l’espace et dans le temps, est pourtant une réalité en réseau : ce point est homologue d’une infinité d’autres qui lui répondent et qui sont lui-même sans pourtant anéantir l’eccéité de chaque nœud du réseau : là, en cette structure réticulaire du réel, réside ce qu’on peut nommer mystère esthétique\teiNote[xmlid={ftn145-2},n={64},place={foot}]{\teiP[xmlid={p-468}]{\teiHi[rend={italic}]{Ibid}., p. 201. }} ». L’art sauve les phénomènes, car il relance la réticulation, il impose de lutter contre toute réification. « L’art n’éternise pas, mais rend transductif, donnant à une réalité localisée et accomplie le pouvoir de passer à d’autres lieux et à d’autres moments. Il ne rend pas éternel, mais donne le pouvoir de renaître et de se réaccomplir ; il laisse des semences de quiddité\teiNote[xmlid={ftn146-2},n={65},place={foot}]{\teiP[xmlid={p-469}]{\teiHi[rend={italic}]{Ibid.}, p. 200. }}. » Ce qui sauve le phénomène n’est pas son absolutisation, mais sa métamorphose possible. Ceci suppose que l’invention, rupture par insertion, soit aussi un geste d’intrusion, d’irruption, bref, soit pour partie « sauvage\teiNote[xmlid={ftn147-2},n={66},place={foot}]{\teiP[xmlid={p-470}]{Gilbert Simondon, \teiHi[rend={italic}]{L’invention dans les techniques}, Paris, Éditions du Seuil, 2005, p. 284 et 292 ; \teiHi[rend={italic}]{Imagination et invention}, \teiHi[rend={italic}]{op. cit}., p. 181 et 182. }} ». Dans sa contre-histoire de la photographie, Simondon insiste sur la domestication du noyau inventif, sur sa progressive aliénation au nom de l’usage, du commerce, bref, sur la manière dont l’externalisation des tâches est démonétisation et extinction de l’invention\teiNote[xmlid={ftn148-2},n={67},place={foot}]{\teiP[xmlid={p-471}]{Voir Carole Maigné, art. cité.}}. À ce titre, parler de « créativité » est tuer l’invention. Simondon relègue la créativité au souci d’habiller un objet pour mieux le vendre\teiNote[xmlid={ftn149-2},n={68},place={foot}]{\teiP[xmlid={p-472}]{\teiHi[rend={italic}]{L’invention dans les techniques}, \teiHi[rend={italic}]{op. cit}., p. 340. }}, la créativité est du même ordre que ce qu’est un concept pour une agence de publicité ou de communication. La relance phénoménale en art est sauvage, car elle installe un « enjambement\teiNote[xmlid={ftn150-2},n={69},place={foot}]{\teiP[xmlid={p-473}]{\teiHi[rend={italic}]{Imagination et invention}, \teiHi[rend={italic}]{op. cit}., p. 182 et 185. }} », elle sauve, car elle fait des sauts, dilatant le réel vers l’avenir. }
              \teiP[xmlid={p25-7}]{L’invention technique est sauvage, car indomptable, renversant les temps : }
              \begin{teiCit}[xmlid={cit05-5}]

                \begin{teiQuote}
Cette évocation présente des idéaux, véhicule des valeurs, et se projette vers l’avenir comme exemple à suivre pour les autres générations : l’image-souvenir veut se réincarner et se perpétuer, elle apporte avec elle la sous-jacence d’une anticipation, et fait en une certaine mesure violence au présent pour l’amener à s’ouvrir vers un avenir de reviviscence. L’anticipation, à son tour, reprend de vieux rêves, contient l’écho d’aspirations anciennes, déjà matérialisées dans d’anciennes images-objets\teiNote[xmlid={ftn151-2},n={70},place={foot}]{\teiP[xmlid={p-474}]{\teiHi[rend={italic}]{Ibid.}, p. 16. }}.
\end{teiQuote}
              
\end{teiCit}
              \teiP[xmlid={p26-7}]{S’il y a chez Simondon un devenir image de tout objet, au point de parvenir à ce syntagme explicite, l’objet-image, ce n’est pas pour tout dissoudre dans le visuel, dans une jouissance du voir, mais au contraire parce que tout ce qui se crée, se transforme, devient sens :}
              \begin{teiCit}[xmlid={cit06-4}]

                \begin{teiQuote}
Presque tous les objets produits par l’homme sont en quelque mesure des objets-images ; ils sont porteurs de significations latentes, non pas seulement cognitives, mais aussi conatives et affectivo-émotives ; les objets-images sont presque des organismes, ou tout au moins des germes capables de revivre et de se développer dans le sujet. Même en dehors du sujet, à travers les échanges et l’activité des groupes, ils se multiplient, se propagent et se reproduisent à l’état néoténique, jusqu’à ce qu’ils trouvent l’occasion d’être réassumés et déployés jusqu’au stade imaginal en se trouvant réincorporés à une invention nouvelle\teiNote[xmlid={ftn152},n={71},place={foot}]{\teiP[xmlid={p-475}]{\teiHi[rend={italic}]{Ibid.}, p. 13. }}.
\end{teiQuote}
              
\end{teiCit}
              \teiP[xmlid={p27-7}]{Si nous maintenons que la latence ici n’est pas celle de la « révélation » sur la surface photosensible, mais un potentiel inventif qui ne peut s’effectuer que par l’artifice de l’appareil, Simondon suggère une articulation originale des images techniques et mentales, assumant le monisme qui crée le lien entre nature et technique et y associant l’imagination. À ce titre, on pourrait défendre l’idée que si « les objets-images – œuvres d’art, vêtements, machines – entrent en obsolescence et deviennent des souvenirs larvaires, fantômes du passé qui s’amenuisent avec les vestiges des civilisations disparues\teiNote[xmlid={ftn153},n={72},place={foot}]{\teiP[xmlid={p-476}]{\teiHi[rend={italic}]{Ibid.}, p. 14. }} », il y a aussi possibilité de réparer, réinventer, casser l’obsolescence programmée grâce à la technicité des images elles-mêmes, grâce à leur beauté. }
              \teiP[xmlid={p28-6}]{Marker ne donne pas le sentiment de croire à l’obsolescence des images, mais bien plus à leur pouvoir de réticulation\teiNote[xmlid={ftn154},n={73},place={foot}]{\teiP[xmlid={p-477}]{On sera sensible au fait que \teiHi[rend={italic}]{Sans soleil} s’ouvre par cette citation de Racine (seconde préface de \teiHi[rend={italic}]{Bajazet} : « l’éloignement des pays \teiHi[rend={italic}]{répare} la trop grande proximité des temps » (nous soulignons). }}, même son œuvre est \teiHi[rend={italic}]{de facto} confrontée à l’obsolescence des techniques\teiNote[xmlid={ftn155},n={74},place={foot}]{\teiP[xmlid={p-478}]{Voir l’article cité de Jean-Michel Frodon : \teiRef[target={https://aoc.media/critique/2021/11/03/reapparitions-et-avatars-de-zapping-zone-installation-fantome-de-chris-marker/}]{https://aoc.media/critique/2021/11/03/reapparitions-et-avatars-de-zapping-zone-installation-fantome-de-chris-marker/} (consulté le 3 septembre 2025).}}. Le monde imaginaire qu’il installe garde ses échardes, n’est pas « imaginaire » au sens d’une confortable évasion, il ne cesse de faire résonner et de relancer la reprise des images photographiques passées. Le paradoxe, que l’on pourrait dire archaïsant, de l’image fixe provoque en fait tout le contraire : un potentiel d’invention constant, comme si Marker refusait l’usure, ne cessant de reprendre des images déjà prises tout en les réinventant. Le monde réel se répare par la photographie et le cinéma, donc par la beauté vitale de l’image technique.}
              \begin{teiCit}[xmlid={cit07-2}]

                \begin{teiQuote}
Le soleil se lève, entre autres. Eh bien c’est comme pour les gens saisis à l’improviste, on ne peut pas se défendre d’un certain sentiment de possession, comme si ces aubes, on les avait… organisées… (un temps.). Et puis je ne sais pas, ce sentiment de rassembler le monde, de le réconcilier, de mettre tous les fuseaux horaires à plat… ça doit faire partie de la nostalgie de l’Eden : qu’il soit partout la même heure. Moi je ne résiste pas à ce genre de films qui vous promène d’une aube à l’autre en disant des trucs comme : il est six heures sur toute la terre\teiNote[xmlid={ftn156},n={75},place={foot}]{\teiP[xmlid={p-479}]{Chris Marker, « Si j’avais 4 dromadaires » \lbrack{}1966\rbrack{}, dans \teiHi[rend={italic}]{Commentaires 2 }\lbrack{}2\teiHi[rend={sup}]{e} éd.\rbrack{}, Paris, Éditions du Seuil, 1967, p. 98.}}… 
\end{teiQuote}
              
\end{teiCit}
              \teiP[xmlid={p29-6}]{L’ubiquité est ici celle de la technique couplée à l’imagination. }
              \teiP[xmlid={p30-6}]{La réticulation simondonienne trouve une traduction artistique dans l’obsession du battement chez Marker\teiNote[xmlid={ftn157},n={76},place={foot}]{\teiP[xmlid={p-480}]{Vincent Jacques, « Images floues, frontières troubles. La pratique photographique de Chris Marker », dans Vincent Jacques (dir.), \teiHi[rend={italic}]{Chris Marker. Photographie}, \teiHi[rend={italic}]{op. cit}., p. 108 et 109. }} : une des toutes premières séquences bouleversantes de \teiHi[rend={italic}]{Sans soleil}, filmée très près, réside dans la conjonction très simple entre des mains posées sur un bastingage de navire japonais et une bande sonore qui reprend le vent et les battements sourds de ce que l’on suppose d’emblée être un cœur, un son entre artifice et réalité qui se transforme encore par le synthétiseur en sorte de son de sous-marin. Le film loue souvent, avec recul toutefois\teiNote[xmlid={ftn158},n={77},place={foot}]{\teiP[xmlid={p-481}]{Marker ne cesse de critiquer l’oubli japonais de la guerre et de la violence qui y fut déployée. }}, la sagesse japonaise (bouddhiste) qui rend hommage aux poupées cassées, qui communient avec toute chose, bénissant même les débris. La référence à Sei Shonagon est connue : Marker se dit touché par le langage de son « copain maniaque » Hayao, car il invente une écriture dont « chacun se servira pour composer sa propre liste de toutes ces choses qui nous font battre le cœur » (\teiHi[rend={italic}]{Sans soleil}), chacun devenant poète. Il n’y a ici aucune naïveté : l’arrêt devant l’horreur est constamment là, l’horreur a un nom comme la beauté absolue (\teiHi[rend={italic}]{Sans soleil}). Si le problème du \teiHi[rend={small-caps}]{xx}\teiHi[rend={sup}]{e} siècle est la « cohabitation des temps » (\teiHi[rend={italic}]{Sans soleil}), Marker parie comme Simondon sur la réversibilité et l’enjambement temporels : la « zone » n’est pas tant geste de zapper au sens de faire défiler distraitement et mécaniquement les flux d’images, c’est plutôt zapper pour enjamber le réel par les images, faire advenir une autre réalité par la force des images techniques, faire advenir ce qui n’est pas encore (\teiHi[rend={italic}]{Sans soleil} toujours\teiNote[xmlid={ftn159},n={78},place={foot}]{\teiP[xmlid={p-482}]{En ce sens, Marker ne me semble pas aller vers l’Atlas Mnemosyne de Warburg : s’ils interrogent certes tous deux thématiquement le temps et la mémoire, travailler les rémanences et la survivance des formes ne se confond pas avec une esthétique du grain informatif et de l’insertion. }}). Et l’on retrouve là quelque chose que l’on avait déjà dans « Si j’avais 4 dromadaires » : « c’est vrai que, quand on regarde autour de soi, c’est l’horreur, c’est la folie, c’est les monstres. Mais il y a déjà… un maquis, une clandestinité du bonheur, une… Sierra Maestra de la Tendresse… quelque chose qui avance, … à travers nous, malgré nous, grâce à nous, quand nous avons la… grâce… et qui annonce, pour on ne sait pas quand, la survivance des plus aimés\teiNote[xmlid={ftn160},n={79},place={foot}]{\teiP[xmlid={p-483}]{Chris Marker, « Si j’avais 4 dromadaires », art. cité, p. 168.}} ». Les points de suspension multipliés sont intentionnels : il est impossible d’asséner cet espoir, il est troué de part en part, il n’y a pas de naïveté politique et historique possible au \teiHi[rend={small-caps}]{xx}\teiHi[rend={sup}]{e} siècle, ce sont « ces petits fragments de guerre enchâssés dans la vie courante » (\teiHi[rend={italic}]{Sans soleil}). L’engagement politique de Marker fait retour, ou plutôt résonne rétroactivement, il n’est pas effacé, mais relancé par l’approfondissement de sa vie de technologue. Comme chez Simondon, la technologie est affaire d’émancipation, de liberté retrouvée par et grâce à la technique. À ce titre, l’image numérique n’est pas le signe d’un renoncement, mais d’une réparation qui change ses modalités et ses exigences. Chris Marker « se déplace et continuera de déplacer sa propre place » : \teiHi[rend={italic}]{Sans soleil} est un manifeste politique qui joue l’écho avec \teiHi[rend={italic}]{Le fond de l’air est rouge}\teiNote[xmlid={ftn161},n={80},place={foot}]{\teiP[xmlid={p-484}]{Jean-Michel Frodon, « Le fond de l’air est rouge, aux tournants du siècle », dans Christine Van Assche, Raymond Bellour et Jean-Michel Frodon (dir.), \teiHi[rend={italic}]{Chris Marker},\teiHi[rend={italic}]{ op. cit.}, p. 286 ; Frodon insiste sur le passage décisif de l’internationalisme à la mondialisation qui accompagne le passage du montage analogique au montage virtuel. }}. }
              \begin{teiCit}[xmlid={cit08-2}]

                \begin{teiQuote}
Le spectateur réussit-il jamais à épuiser les objets qu’il contemple ? Ses errances n’ont pas de fin. Parfois, cependant, après qu’il a sondé un millier de possibilités, il peut avoir l’impression, dans une tension de tous ses sens, de capter un murmure confus. Les images se mettent à rendre des sons et les sons, à leur tour, se font images. Lorsque ce murmure indistinct – le murmure de l’existant (\teiHi[rend={italic}]{the murmure of existence}) – l’atteint, c’est peut-être bien alors qu’il s’approche le plus du but inatteignable\teiNote[xmlid={ftn162},n={81},place={foot}]{\teiP[xmlid={p-485}]{Siegfried Kracauer, \teiHi[rend={italic}]{Théorie du film. La rédemption de la réalité matérielle}, Daniel Blanchard et Claude Orsini (trad.), Philippe Despoix et Nia Perivolaropoulou (présent.), Paris, Flammarion, 2010, p. 246. }}. 
\end{teiQuote}
              
\end{teiCit}
              \teiP[xmlid={p31-6}]{Cette citation nous semble faire sens, car Kracauer imposait justement la beauté vitale de la photographie comme genèse de l’image filmique, comme son noyau réticulaire, si l’on peut dire. Alors même que Kracauer professe une indifférence presque étonnante à l’égard de la technique de l’image, comme si ce qui comptait pour lui était le principe de l’image appareillée sans se préoccuper de ses modalités concrètes, il retrouve pourtant quelque chose de la pensée de Simondon confrontée à celle de Marker : une esthétique de la résonance, du ténu, de l’indéfinie différenciation. En forgeant son concept de « caméra-réalité », il jouait sur le trait d’union pour signifier que la caméra n’est pas caméra \teiHi[rend={italic}]{de}, ni caméra \teiHi[rend={italic}]{face à}, mais articulation indissociable du réel et de son enregistrement, photographie comme déprise et pas seulement comme prise du réel\teiNote[xmlid={ftn163},n={82},place={foot}]{\teiP[xmlid={p-486}]{Voir Carole Maigné, « La radicale caméra-réalité de Siegfried Kracauer », \teiHi[rend={italic}]{Archives de philosophie}, vol. 85, n\teiHi[rend={sup}]{o} 1, 2022, p. 45-65.}}. Le réalisme traversé et construit par la technique, en trait d’union, le réalisme de l’image technique n’implique pas la mainmise de la technique sur le monde, ni un face-à-face, mais une réalité que les appareils modulent de part en part.}
            
\end{teiDiv}
          
\end{teiElement}
        
\end{teiElement}
        \begin{teiElement}[name={group},type={chapter},data-page-title={Beautés jazz(s)},data-page-authors={Pierre Sauvanet@@Université Bordeaux Montaigne, ARTES – UR 24141, France},data-include-href={Ch12\_beautes\_jazz\_Sauvanet.xml},xmlid={chapter-011}]

          \begin{teiElement}[name={front}]

            \begin{teiDiv}[type={titlePage},xmlid={titlepage-013}]

              \teiP[rend={title-main},xmlid={p-487}]{Beautés jazz(s)}
              \teiP[rend={author-aut},xmlid={p-488}]{Pierre \teiHi[rend={small-caps}]{Sauvanet}}
              \teiP[rend={authority\_affiliation},xmlid={p-489}]{Université Bordeaux Montaigne, ARTES – UR 24141, France}
            
\end{teiDiv}
          
\end{teiElement}
          \begin{teiElement}[name={body},xmlid={body-013}]

            \teiP[xmlid={p1-13}]{Première question, en guise de préambule : y a-t-il une beauté non visuelle ? On associe le plus souvent la beauté à la vue. Mais qu’en est-il des autres sens ? Qu’est-ce qu’un beau son, une belle musique ? On sait que pour Plotin, déjà, le beau existe non seulement à travers la vue, mais également l’ouïe – ce que l’on retrouvera, sous une autre forme, dans la notion de « sens intellectualisés » chez Hegel (la vue et l’ouïe en sont, le goût, le toucher et l’odorat n’en sont pas). En tout cas, on ne dit pas une belle odeur, un beau parfum, ni d’ailleurs un beau goût, ou un beau plat, si ce n’est précisément par rapport à la vue. C’est exactement ce que notait Diderot au milieu du \teiHi[rend={small-caps}]{xviii}\teiHi[rend={sup}]{e} siècle, dans ses \teiHi[rend={italic}]{Recherches philosophiques sur l’origine et la nature du beau }: « On dit \teiHi[rend={italic}]{un mets excellent, une odeur délicieuse}, mais non \teiHi[rend={italic}]{un beau mets, une belle odeur}. Lors donc qu’on dit, \teiHi[rend={italic}]{voilà un beau turbot, voilà une belle rose}, on considère d’autres qualités dans la rose et dans le turbot que celles qui sont relatives aux sens du goût et de l’odorat\teiNote[xmlid={ftn1-12},n={1},place={foot}]{\teiP[xmlid={p-490}]{Denis Diderot, \teiHi[rend={italic}]{Recherches philosophiques sur l’origine et la nature du beau} \lbrack{}1752\rbrack{}, dans \teiHi[rend={italic}]{Œuvres esthétiques}, Paris, Bordas, « Classiques Garnier », 1988, p. 418.}} ». De même pour nous : il s’agit de penser une beauté sonore, proprement musicale, qui doive tout au sens de l’ouïe, et rien (ou presque rien) au sens de la vue. On postulera donc qu’une beauté non visuelle, ici musicale, est possible, au sens où (comme le dit également Diderot) elle permet d’éveiller en moi la perception de rapports, et qu’elle se manifeste comme un processus relativement instable plutôt que comme un état stable.}
            \teiP[xmlid={p2-13}]{Depuis Dada, pour qui « la beauté est une sorte de morte », supplantée par « la nouveauté, l’intensité, l’étrangeté » comme le dit Paul Valéry, la notion de beauté aura connu tous les avatars, qu’on est plus habitué à voir renaître dans le domaine visuel dit des « beaux-arts », que ce soit dans l’art moderne ou contemporain. On aura compris qu’on a fait un tout autre choix ici, en s’intéressant de plus près à la notion de beauté dans un type de musique bien précise de la seconde moitié du \teiHi[rend={small-caps}]{xx}\teiHi[rend={sup}]{e} siècle et au-delà, celui de la beauté en jazz, ou plus précisément des beautés jazz(s)\teiNote[xmlid={ftn25-3},n={2},place={foot}]{\teiP[xmlid={p-491}]{Pluriel optionnel et volontairement « barbare », qui évoque aussi le titre que nous avons jugé bon de donner à notre ouvrage à quatre mains avec Colas Duflo, auquel je me permets de renvoyer le lecteur (\teiHi[rend={italic}]{Jazzs }\lbrack{}2003\rbrack{}, Paris, Éditions MF, 2008).}}, à partir essentiellement de deux thèmes qui semblent se répondre, \teiHi[rend={italic}]{Ugly Beauty} de Thelonious Monk (1917-1982) et \teiHi[rend={italic}]{Beauty is a Rare Thing }d’Ornette Coleman (1930-2015). Certes, nous aurions pu là aussi nous intéresser ici aux beautés spécifiquement visuelles du jazz, et notamment aux photographes de renom (d’Herman Leonard à Guy Le Querrec) ou au graphisme des pochettes de disque (du label Blue Note à ECM\teiNote[xmlid={ftn26-2},n={3},place={foot}]{\teiP[xmlid={p-492}]{Pour de nombreux exemples, voir notamment le catalogue de l’exposition \teiHi[rend={italic}]{Le Siècle du jazz}, Daniel Soutif (dir.), Paris, Musée du Quai Branly-Flammarion, 2009.}}). Sans négliger l’apport possible d’une forme de synesthésie entre son et image dans le cas présent (quand on aime un disque en particulier, on a également en tête sa pochette, et les deux semblent souvent indissociables), l’essentiel concernera ici la question de la beauté sonore, la beauté pour l’oreille, donc la forme temporelle bien davantage que la forme spatiale.}
            \teiP[xmlid={p3-13}]{En un mot, notre hypothèse est la suivante : même si elles sont diverses et variées, les beautés jazz(s) ont en commun de ne pas être pures, sereines, canoniques. Au contraire, ce qui fait la beauté du jazz, ou dans le jazz, c’est le trouble, la déchirure, voire la dissonance. Bien entendu, cela paraît encore très général, et il faudra le montrer concrètement à partir de la musique elle-même. Bien entendu encore, on pourra toujours trouver des beautés consonantes et apparemment non dissonantes dans le jazz (Bill Evans, par exemple), mais pour parler comme Nietzsche – et surtout comme Jean-Pierre Moussaron, qui reprenait volontiers Nietzsche pour parler du jazz –, si la forme peut être apollinienne, l’arrière-fond est toujours dionysiaque. Pour le dire d’une formule à l’emporte-pièce : derrière le \teiHi[rend={italic}]{jazz}, il y a toujours le \teiHi[rend={italic}]{blues}.}
            \begin{teiDiv}[type={section1},xmlid={div01-9}]

              \teiHead[xmlid={ugly-beauty-001}]{
                \teiHi[rend={italic}]{Ugly Beauty}
              }
              \teiP[xmlid={p4-13}]{\teiHi[rend={italic}]{Ugly Beauty} est une des dernières compositions de Monk, enregistrée le 14 décembre 1967 pour l’album \teiHi[rend={italic}]{Underground}, et la seule à trois temps de tout son répertoire (on y reviendra). De même, \teiHi[rend={italic}]{Underground} est un des derniers albums de Monk, et le dernier avec son fidèle quartet des années 1960 (Charlie Rouse au saxophone ténor, Larry Gales à la contrebasse, Ben Riley à la batterie). En réalité, à peu près tout le monde s’accorde à dire que c’est son dernier grand album (Monk disparaîtra presque totalement de la scène musicale au début des années 1970, jusqu’à sa mort en 1982 ; il était d’ailleurs un grand spécialiste du silence, au moins depuis la séance en studio du 24 décembre 1954, durant laquelle il opposa à Miles Davis un silence forcené durant le chorus de \teiHi[rend={italic}]{The Man I Love}). Comme le souligne aujourd’hui l’un des meilleurs spécialistes, le pianiste Laurent de Wilde, les enregistrements gravés à cette période « témoignent d’une sève encore intarissable. Tous les tempos y sont balayés avec une égale superbe. Monk ne se fend pas moins, pour \teiHi[rend={italic}]{Underground}, de quatre compositions originales, dont une valse, \teiHi[rend={italic}]{Ugly Beauty}, la première en cinquante ans. \lbrack{}... Mais\rbrack{} combien ironique à ce titre est la récompense obtenue aux Grammy Awards par l’album \teiHi[rend={italic}]{Underground}, pourtant d’une qualité exceptionnelle, non pour son contenu, mais pour sa pochette\teiNote[xmlid={ftn27-2},n={4},place={foot}]{\teiP[xmlid={p-493}]{Laurent de Wilde, \teiHi[rend={italic}]{Monk} \lbrack{}1996\rbrack{}, Paris, Gallimard, « Folio Essais », 2010, p 274-276. Une récente bande dessinée, \teiHi[rend={italic}]{Monk !} (de Youssef Daoudi, Paris, Les Éditions Martin de Halleux, 2018, p. 301 et suiv.) met en scène la fabrication imaginaire de cette pochette, apparemment voulue par le célèbre producteur Teo Macero (qui était aussi celui de Miles Davis à la même époque), et cite également à ce propos le thème \teiHi[rend={italic}]{Ugly Beauty}. Comme l’indique son sous-titre (\teiHi[rend={italic}]{Thelonious, Pannonica… une amitié, une révolution musicale}), la bande dessinée vaut aussi en général pour le graphisme original d’un imaginaire commun entre Monk et sa mécène.}} ». Où l’on retrouve la question laissée en suspens de la beauté sonore par rapport à la beauté visuelle ! Il faut dire que, sur l’ensemble de sa discographie, Monk s’est souvent mis en scène avec une certaine fantaisie pour ses pochettes de disque : qu’on pense à celles de \teiHi[rend={italic}]{Brilliant Corners} (1956, réfléchi sous cinq facettes par un faux miroir), de \teiHi[rend={italic}]{Monk’s Music} (1957, assis dans une petite carriole rouge d’enfant) ou encore de\teiHi[rend={italic}]{ Alone in San Francisco }(1959, en équilibre sur le marchepied du tramway de la ville). }
              \teiP[xmlid={p5-13}]{La pochette de \teiHi[rend={italic}]{Underground}\teiNote[xmlid={ftn28-2},n={5},place={foot}]{\teiP[xmlid={p-494}]{https://en.wikipedia.org/wiki/Underground\_(Thelonious\_Monk\_album).}}, conçue par le designer John Berg (et travaillée avec les photographes Steve Horn et Norman Griner), montre Monk habillé en résistant de la seconde guerre mondiale ; la photo est prise dans son propre appartement, transformé pour l’occasion en refuge imaginaire de la résistance française (voir le graffiti au mur) ; au fond, dans un savoureux capharnaüm indescriptible qu’on ne commentera pas davantage, on distingue, portant une mitraillette, la célèbre Pannonica de Koenigswarter, grande amie de Monk (on y reviendra, là encore).}
              \teiP[xmlid={p6-13}]{\teiHi[rend={italic}]{Ugly Beauty}, donc : faisons le pari que l’oxymore du titre est pertinent et significatif. Derrière des apparences de valse lente et douce, sur un moule classique en jazz de trente-deux mesures de forme AABA, tout le thème semble en réalité construit autour d’une double et durable ambiguïté, à la fois harmonique et rythmique. Sans entrer dans des enjeux musicologiques trop pointus, on peut du moins suggérer ceci : dès le début, la triple répétition de la quinte bémol installe une sorte de doute ou de trouble affirmé (ré-la-la bémol à l’octave do-ré-mi-sol) ; dans le thème, alors que la tonalité d’ensemble paraît commencer en mineur (sur un accord de ré mineur 7), avec un passage par le sol, les mesures 7 et 8 voient s’enchaîner rapidement trois accords (avec passage) de ré bémol majeur 7 à mi bémol mineur 7, tandis que la fin semble s’installer en ré bémol majeur 7, pour mieux recommencer en mineur… Ce sentiment de tension harmonique, d’oscillation, voire de flottement, se retrouve du côté du rythme : dès le début, les arpèges de sept notes suivies d’un silence pourraient faire penser à une mesure binaire et non ternaire, sauf qu’ils sont joués en triolets et répétés trois fois ; dans le thème, on croit s’arrêter souvent, mais la dernière note relance la machine, sur la levée du temps ; puis, lors des chorus, où la batterie (avec la caisse claire jouée aux balais) est un peu plus audible, le batteur ne cesse d’osciller entre trois et quatre temps, sans marquer nécessairement la barre de mesure. Et même, sur l’extrait filmé lors de l’enregistrement en studio, on entend clairement le batteur faire le choix d’une mesure binaire à de nombreux endroits\teiNote[xmlid={ftn29-2},n={6},place={foot}]{\teiP[xmlid={p-495}]{Lien vidéo : \teiRef[target={https://www.youtube.com/watch?v=In16H9J72HY}]{\teiHi[rend={underline}]{https://www.youtube.com/watch?v=In16H9J72HY}}, extrait du remarquable documentaire de Charlotte Zwerin, \teiHi[rend={italic}]{Thelonious Monk. Straight No Chaser} (1988), consulté le 2 juillet 2025.}}. Pourtant, selon certaines sources (Laurent Coq notamment\teiNote[xmlid={ftn30},n={7},place={foot}]{\teiP[xmlid={p-496}]{Voir « Thelonius \lbrack{}\teiHi[rend={italic}]{sic}\rbrack{} Monk, compositeur », table ronde avec Vincent Bessières, Philippe Baudoin, Laurent Coq, Laurent De Wilde et Xavier Prévost, enregistrée à la Cité de la musique, le 3 mars 2012, \teiRef[target={https://pad.philharmoniedeparis.fr/player-conference.aspx\#0985721}]{https://pad.philharmoniedeparis.fr/player-conference.aspx\#0985721} (consulté le 2 juillet 2025).}}), c’est le batteur Ben Riley lui-même qui aurait suggéré à Monk de « tirer » la composition vers le trois temps.}
              \teiP[xmlid={p7-13}]{Attention toutefois à ne pas surinterpréter le titre : il est arrivé souvent à Monk d’intituler ses thèmes de manière volontairement dilettante, comme si cela n’avait strictement aucune importance, en réponse à une demande pressante du producteur ou de l’éditeur (\teiHi[rend={italic}]{Let’s Call This},\teiHi[rend={italic}]{ Think of One},\teiHi[rend={italic}]{ Well You Needn’t} sont de cette veine : \teiHi[rend={italic}]{Appelons-le ceci},\teiHi[rend={italic}]{ Pense à un \lbrack{}titre\rbrack{}},\teiHi[rend={italic}]{ Eh bien tu n’en as pas besoin}, etc.). Dans le cas présent, comme pour d’autres à leur tour (\teiHi[rend={italic}]{Epistrophy}, \teiHi[rend={italic}]{Misterioso}, \teiHi[rend={italic}]{Round Midnight}, etc.), il est pourtant difficile d’imaginer que le titre soit entièrement neutre, non signifiant. À signaler qu’un acteur de la même époque (acteur blanc américain de western bien connu), Jack Palance, répondait, semble-t-il, pour son visage étrange et émacié, au titre de \teiHi[rend={italic}]{ugly beautiful}, ou de \teiHi[rend={italic}]{beautifully ugly}, mais rien n’indique qu’il s’agisse ici d’un clin d’œil de la part de Monk, qui pensait plus vraisemblablement à une sorte d’autoportrait polysémique de sa propre musique, voire – pourquoi pas – à un autoportrait tout court. Autre solution possible de l’énigme, encore plus simple, quoique peu vraisemblable (mais par principe, il faut tout envisager) : dans la mesure où ce thème constitue un hapax dans l’œuvre monkienne (le seul à trois temps), la « laideur » en question pourrait renvoyer à l’étrangeté selon lui du rythme de valse en jazz (dont on rappellera au passage que son usage ne date que de 1957, avec l’album \teiHi[rend={italic}]{Jazz in 3/4 Time} du batteur Max Roach).}
              \teiP[xmlid={p8-13}]{Dans ces conditions, la beauté dont il s’agit ici est-elle vraiment « laide » ? Encore une fois, tout dépend de ce que l’on entend par ce terme. Ce qui est certain, c’est qu’elle n’est pas « jolie », et il est clair que c’est là l’essentiel : la joliesse n’a rien à voir avec l’univers de Monk. Prenons des contre-exemples : harmoniquement, comme des centaines de tubes, on aurait pu tout construire autour d’un accord de ré mineur (ré-fa-la) décliné selon toutes les cadences classiques possibles ; rythmiquement, comme des centaines de tubes, on aurait pu s’en tenir à un bon \teiHi[rend={italic}]{tchac-poum} bien marqué. Ça, ça aurait été joli ! Mais non, il a fallu que notre original s’entête à jouer une quinte bémol trois fois dès le début, creuse toujours l’écart entre ton et demi-ton, et s’entiche d’un mètre en trois-quatre, qui peut tout aussi bien devenir quatre-quatre, six-huit, etc. Dans ces conditions, « laide » voudrait-elle dire « complexe », volontairement voire inutilement compliquée, tarabustée ? Non plus, car ces complications ne sont jamais gratuites : paradoxalement, le thème peut aussi sonner « simple » dans sa lenteur même, en tout cas parfaitement « fluide » sur ses trois temps. Faut-il dire « triste », alors, quand il y a tellement de valses tristes dans le répertoire ? Non plus, car la musique de Monk produit également l’impression d’une certaine gaieté paradoxale, comme lorsque le pianiste lui-même se lève de son instrument et se met à danser… Nous sommes là au cœur de l’esthétique monkienne : l’oxymore de la beauté laide, c’est le non joli qui pourtant sonne bien, c’est le complexe qui pourtant sonne simple, c’est le rythme appuyé qui pourtant sonne fluide, c’est le triste qui pourtant sonne joyeux, bref, c’est la tension qui fait l’harmonie, c’est le silence qui fait le rythme. Si l’on cherchait un curieux équivalent plastique, on pourrait penser à la fois à Kandinsky (pour qui l’art abstrait, dans son exigence spirituelle et sa nécessité intérieure, ne peut pas et ne doit pas être décoratif) et à Matisse (pour qui, au contraire, l’art est de toute façon décoratif, au sens où il embellit le quotidien). En un certain sens, la beauté laide de Monk \teiHi[rend={italic}]{est} et \teiHi[rend={italic}]{n’est pas} décorative : elle l’est, dans la mesure où l’on peut très bien passer ce disque en musique de fond (et le jazz tend aujourd’hui à devenir cette musique en général) ; elle ne l’est pas, dans la mesure où les tensions harmoniques et les brisures rythmiques sont toujours présentes, où ses aspérités et ses complexités sont toujours là, comme des grumeaux dans la pâte musicale.}
              \teiP[xmlid={p9-12}]{Écoutons dans cet esprit deux autres versions de la même \teiHi[rend={italic}]{Ugly Beauty} (sachant que l’on compte aujourd’hui environ cent cinquante versions enregistrées de ce thème !). L’un des plus grands plaisirs du jazz, l’une de ses plus grandes beautés aussi, ce sont les différentes versions d’un même thème, comme autant d’« interprétations » qui portent bien leur nom (en ce qu’elles sont beaucoup plus libres que dans d’autres styles musicaux). Ces deux versions sont radicalement différentes, et même quasiment antithétiques, mais elles confirment en ce sens nos intuitions précédentes : la première (celle du batteur et compositeur Paul Motian, sur l’album hommage \teiHi[rend={italic}]{Monk in Motian}, en 1989) est volontairement hors tempo, comme éthérée, à tel point qu’elle semble flotter dans l’air ; la seconde (celle de Stefano Bollani, pour piano seul, en 2003) condense au contraire les tensions harmoniques dans un cadre rythmique totalement renouvelé, non plus ternaire, mais binaire, et qui réutilise la quinte bémol du début, subtilement altérée, comme une sorte de ligne de basse tout du long du thème, et même des improvisations. En deux mots, une version met en valeur le flottement rythmique, l’autre le frottement harmonique. Dans les deux cas, on retrouve bien quelque chose de notre paradoxale \teiHi[rend={italic}]{Ugly Beauty} en jazz, et quitte à parler de reprises actuelles, il reste à signaler deux parutions récentes portant ce nom : en 2004, un disque du pianiste néerlandais Robert Jan Vermeulen entièrement consacré à Monk ; en 2022, un livre du journaliste américain Phil Freeman défendant la thèse selon laquelle, loin d’être mort, le jazz n’a au contraire jamais été aussi vivant\teiNote[xmlid={ftn31},n={8},place={foot}]{\teiP[xmlid={p-497}]{Phil Freeman, \teiHi[rend={italic}]{Ugly Beauty. Jazz in the 21st Century}, Zero Books, Pennsylvania, John Hunt Publishing, 2022.}}.}
            
\end{teiDiv}
            \begin{teiDiv}[type={section1},xmlid={div02-9}]

              \teiHead[xmlid={beauty-is-a-rare-thing-001}]{
                \teiHi[rend={italic}]{Beauty Is A Rare Thing}
              }
              \teiP[xmlid={p10-12}]{Le premier enregistrement datait du 14 décembre 1967, le suivant date du 2 août 1960. La chronologie n’est pas respectée, car en réalité tout se joue dans l’autre sens : le thème de Monk arrive à la fin de sa carrière quand, largement reconnu au point de faire la couverture du magazine \teiHi[rend={italic}]{Time} (1964), il peut pourtant paraître déjà dépassé par d’autres jeunes jazzmen, dont Ornette Coleman, précisément, qui propose ce titre au début de sa discographie (sur son cinquième album, au nom tout aussi affirmatif que les précédents : \teiHi[rend={italic}]{This Is our Music}\teiNote[xmlid={ftn32},n={9},place={foot}]{\teiP[xmlid={p-498}]{https://en.wikipedia.org/wiki/This\_Is\_Our\_Music\_(Ornette\_Coleman\_album).}}, après \teiHi[rend={italic}]{Something Else !!!!, Tomorrow Is the Question, The Shape of Jazz to Come, Change of the Century}). }
              \teiP[xmlid={p11-12}]{Pour le dire autrement, Ornette à l’époque sonne déjà plus « moderne » que Monk. Quand il débarque en 1959 au Five Spot à New York, avec son quartet sans piano (Don Cherry à la trompette, Charlie Haden à la contrebasse, Billy Higgins à la batterie), il divise aussitôt le monde du jazz (par exemple entre le très classique John Lewis, pianiste du Modern Jazz Quartet, qui va le défendre paradoxalement, et Miles Davis qui va au contraire le détester). }
              \teiP[xmlid={p12-12}]{Il faut dire que le son même du saxophone d’Ornette, en plastique blanc, pose déjà problème à l’époque (sans parler de son jeu à la trompette ou au violon) : que son alto semble rire ou pleurer, il possède un timbre tranchant, et pourtant légèrement voilé, manifestant un lyrisme parfois déchirant, mais toujours en conservant volontairement les petites (ou les grandes) imperfections. On connaît évidemment d’abord Ornette Coleman comme l’inventeur du \teiHi[rend={italic}]{free jazz}, lors de la légendaire séance du 21 décembre 1960 pour le label Atlantic. L’album, sorti en septembre 1961, porte en couverture une peinture de Jackson Pollock, \teiHi[rend={italic}]{White Light} (1954). Quelques mois auparavant, Ornette avait écrit ceci dans ses notes de pochette pour l’album \teiHi[rend={italic}]{Change of the Century }: la musique que nous jouons, « peut-être est-ce quelque chose comme les peintures de Jackson Pollock \lbrack{}\teiHi[rend={italic}]{Maybe it's something like the paintings of Jackson Pollock}\rbrack{} ».}
              \teiP[xmlid={p13-11}]{Le thème en question se trouve en réalité juste avant \teiHi[rend={italic}]{Free Jazz} : l’enregistrement de \teiHi[rend={italic}]{Beauty Is A Rare Thing} date du 2 août 1960 à New York, celui de \teiHi[rend={italic}]{Free Jazz} du 21 décembre, même année, même lieu. On peut donc avancer que \teiHi[rend={italic}]{Beauty }est le dernier enregistrement avant le « free » total, ce qui lui donne peut-être une autre saveur (sachant que Coleman reviendra de toute façon sur une musique moins \teiHi[rend={italic}]{free} après le « \teiHi[rend={italic}]{Free »}, parce qu’il n’avait pas d’autre solution). Thème emblématique, sans doute d’ailleurs moins pour sa musique que pour la symbolique de son titre, \teiHi[rend={italic}]{Beauty Is a Rare Thing }a également fini par désigner, dans l’édition en coffret de 1993 rééditée en 2015, la totalité des enregistrements du saxophoniste sur le label Atlantic entre 1959 et 1961 (une cinquantaine de titres). Au passage, on remarquera que nos deux thèmes occupent respectivement la même place stratégique au sein de leur album de référence, à savoir la deuxième. Pas la première, souvent réservée à un thème au tempo rapide, censé accrocher l’auditeur à la première écoute ; mais la seconde, souvent au tempo lent, qui s’écoute différemment et qui possède une vraie importance dans l’équilibre de l’ensemble. Même si beaucoup de choses les sépareraient par ailleurs (ainsi Monk ne supportant pas la nouveauté du \teiHi[rend={italic}]{free jazz}, lui qui aime tant l’idée même de composition), une autre chose au moins les rassemble, et qu’on entend particulièrement bien dans les deux extraits choisis : l’affirmation de leur musique jusque dans l’imperfection. Monk peut bien jouer un accord avec un ou deux doigts à côté, peu importe (« au piano, il n’y a pas de fausses notes », aurait-il dit) ; Ornette peut bien sembler faire un canard au saxophone, peu importe, au contraire même ! La « beauté » du jazz peut bien être plus ou moins « laide » selon d’autres critères de perfection, c’est l’affirmation de la musique elle-même qui est vitale.}
              \teiP[xmlid={p14-11}]{Une anecdote racontée par Ornette Coleman lui-même en dit long sur notre sujet : « Une fois, j’ai joué pendant dix minutes pour un congrès de l’association des Architectes américains – c’était pour illustrer un débat intitulé “Beauté et laideur”, et moi, on m’a dit que je représentais la laideur\teiNote[xmlid={ftn33},n={10},place={foot}]{\teiP[xmlid={p-499}]{Ornette Coleman, « Tête à tête avec Ornette », \teiHi[rend={italic}]{Jazz Magazine}, n\teiHi[rend={sup}]{o} 125, décembre 1965, p. 30, repris dans Philippe Carles et Jean-Louis Comolli, \teiHi[rend={italic}]{Free jazz Black power}, Paris, Gallimard, « Folio Essais », 2000, p. 383.}}. » L’histoire ne dit pas si les architectes en question avaient conscience de l’humiliation alors infligée, ou bien s’ils appréciaient la musique de Coleman tout en pensant qu’elle était laide, mais tel que l’anecdote est racontée, c’est bien la première hypothèse la plus vraisemblable : littéralement, Ornette Coleman devait \teiHi[rend={italic}]{représenter} la laideur en soi, l’incarner par sa seule présence durant un temps déterminé – « ce qu’il fit sans doute de belle manière avec sa gentillesse coutumière\teiNote[xmlid={ftn34-2},n={11},place={foot}]{\teiP[xmlid={p-500}]{Philippe Méziat, « Ornette Coleman (1930-2015) », \teiHi[rend={italic}]{Citizen Jazz}, 2015, \teiRef[target={https://www.citizenjazz.com/Ornette-Coleman-1930-2015.html}]{https://www.citizenjazz.com/Ornette-Coleman-1930-2015.html} (consulté le 2 juillet 2025).}} », suppose-t-on. Comme si, bien entendu, cette notion même de laideur pouvait être objective, alors qu’elle est en réalité ici, non pas exactement subjective, mais du moins relative – relative à un contexte (historique, esthétique, social, racial, etc.), qui fait que ce ne sont pas tant des sujets qui s’affrontent par leurs goûts différents, mais des groupes constitués ou en train de se constituer comme tels. On sait notamment que, si la musique d’Ornette Coleman est souvent vue par ses détracteurs comme laide, fausse, hybride, violente, dissonante (ce sont les qualificatifs qui reviennent le plus souvent), en revanche le groupe de ses admirateurs n’a cessé, chemin faisant, de croître jusqu’à atteindre des sommets de fanatisme (il y a des fans absolus d’Ornette qui ne jurent que par lui, ce qui est somme toute assez rare en jazz, où les goûts sont souvent éclectiques). Toujours est-il que, pour une lecture racialiste des événements, le fait qu’Ornette Coleman ait pu répondre positivement à une telle demande renvoie à la condition de l’homme noir, qui intègre apparemment sa propre infériorité du point de vue du blanc, mais qui ne cesse de ruser avec elle, avec lui. Ce fait – un jazzman noir acceptant d’incarner une musique laide – apparaît dès lors, au contraire, comme un combat contre l’intégration par l’homme noir de sa propre supposée « laideur », et conduira bientôt à l’affirmation radicale, non seulement d’un \teiHi[rend={italic}]{Black Power}, mais d’une \teiHi[rend={italic}]{Black Beauty} (selon le titre de Miles Davis lui-même, qu’il donne à un de ses disques enregistré en public en 1970, et sorti en 1973).}
              \teiP[xmlid={p15-11}]{Revenons donc au titre, \teiHi[rend={italic}]{Beauty Is A Rare Thing}, « la beauté est une chose rare » : la beauté dont il s’agit ici n’est pas seulement la beauté de la musique, mais pour un musicien comme Ornette Coleman, c’est sans aucun doute une sorte de beauté spirituelle, sacrée, voire mystique\teiNote[xmlid={ftn35-3},n={12},place={foot}]{\teiP[xmlid={p-501}]{Voir Pierre Sauvanet, « Mystique Ornette ? » dans Sylvie Chalaye et Pierre Le Tessier (dir.), \teiHi[rend={italic}]{Sacré Jazz !}, Paris, Passages, « Esthetique(s) Jazz », 2025.}}. Faut-il même aller jusqu’à dire une sorte de beauté « philosophique » ? Pourquoi pas… Quand on sait que, lors d’une tournée en France à l’été 1997, l’artiste avait demandé à son agent de lui trouver un philosophe français avec qui discuter, et qui fut en l’occurrence Jacques Derrida\teiNote[xmlid={ftn36-3},n={13},place={foot}]{\teiP[xmlid={p-502}]{Voir Pierre Sauvanet, « “Qu’est-ce qui arrive ? \teiHi[rend={italic}]{What’s happening?}” Retour sur le malentendu entre Ornette Coleman et Jacques Derrida », \teiHi[rend={italic}]{Epistrophy}, n\teiHi[rend={sup}]{o} 4, « Le jazz, la philosophie et les philosophes », 2019, \teiRef[target={http://www.epistrophy.fr/}]{\teiHi[rend={underline}]{http://www.epistrophy.fr/}} (consulté le 2 juillet 2025).}}, il n’est pas totalement absurde de penser que ce titre puisse être une allusion à la philosophie de Spinoza. Pour être plus précis, la phrase en question n’est autre que : « Tout ce qui est beau est difficile autant que rare\teiNote[xmlid={ftn37-3},n={14},place={foot}]{\teiP[xmlid={p-503}]{Spinoza\teiHi[rend={italic}]{, Éthique} \lbrack{}1677\rbrack{}, Charles Appuhn (trad.), Paris, Garnier-Flammarion, 1965, livre V, prop. XLII, scolie, p. 341.}} ». Ce sont les tout derniers mots de l’\teiHi[rend={italic}]{Éthique }de Spinoza, qui renvoient à une élévation progressive de la pensée tout au long des cinq livres de ce monument philosophique ; en l’occurrence, il s’agit plutôt d’une forme de « salut » de l’être humain par la conscience de sa propre liberté déterminée. Les hommes se croient libres en ce qu’ils sont conscients de leurs actes et ignorants des causes qui les font agir ; le sage du cinquième livre cherche quant à lui à suivre son \teiHi[rend={italic}]{conatus}, à accroître sa conscience, et à remplacer les passions tristes par les passions joyeuses. C’est ce chemin du salut par la liberté de l’âme qui est particulièrement ardu, et rares sont ceux qui y parviennent. On aura donc compris que la beauté en question, qui a fini par s’imposer dans la plupart des traductions (dans la version originale : \teiHi[rend={italic}]{omnia praeclara tam difficilia quam rara funt, }le latin parle en réalité au neutre pluriel de choses \teiHi[rend={italic}]{praeclara}, c’est-à-dire précieuses, excellentes, prééminentes) est plutôt synonyme philosophiquement de souverain bien. Mais tout ceci nous emmènerait trop loin de notre objet\teiNote[xmlid={ftn38-3},n={15},place={foot}]{\teiP[xmlid={p-504}]{Dans un article de blog écrit à la mort de l’artiste, et déjà cité plus haut, Philippe Méziat n’hésite pas à écrire ceci de son côté, auquel on n’est pas loin de souscrire : « Ornette Coleman est profondément spinoziste, et la philosophie de Spinoza pourrait se résoudre, en dernière analyse, en une musique d’inspiration colemanienne. \lbrack{}…\rbrack{} Dans les deux cas de figure, une même tranquillité d’énonciation, mais aussi une même certitude que rien ne doit venir, de l’extérieur, freiner les aspirations du désir (essence de l’homme) et des passions joyeuses. La tristesse ainsi mise au rancart en même temps que la bêtise ne va jamais sans produire quelques réactions du côté de ceux, maîtres ou esclaves, qui en font le support de leur assez pitoyable existence » (art. cité).}}. Retenons l’essentiel, à savoir qu’en l’occurrence, les mots d’Ornette (et de Baruch ?) permettent du moins de questionner le passage de l’esthétique à l’éthique, et réciproquement : le beau n’est pas le joli, ni l’agréable, il engage une exigence, un engagement profond de l’être. Et cela même est excessivement rare, comme on le verra plus tard.}
              \teiP[xmlid={p16-11}]{Écoutons maintenant, en guise de point d’orgue, une autre version en duo du même thème, par Frans Vermeerssen (saxophone alto) et Dion Nijland (contrebasse, jouée à l’archet), enregistrée au conservatoire d’Arnhem (Pays-Bas) en 2019 : elle permet d’entendre de manière épurée les tensions harmoniques entre les deux instruments en présence (le saxophoniste cherchant à imiter les inflexions d’Ornette). Au passage, cette version très récente permet aussi de constater que l’héritage d’Ornette Coleman est toujours bien vivant, et que nous parlons bien, à nouveau, de beautés vitales contemporaines. Comme précédemment avec les cent cinquante versions d’\teiHi[rend={italic}]{Ugly Beauty} (sans compter celles non enregistrées, bien entendu), cette profusion de la reprise dans le jazz reste tout à fait fascinante, et va dans le sens de notre intuition première d’une beauté toujours impure, métissée, diffractée. C’est Colas Duflo qui écrit ceci, à la fin de \teiHi[rend={italic}]{Jazzs }: « Le jazz va contre l’idée d’esthétique pure. \lbrack{}…\rbrack{} Il mélange et déplie en toute impureté tous ces objets pliés que sont les standards, les grilles, les thèmes, les fragments de mélodie, pour en faire des oiseaux de paradis. Toutes les versions d’\teiHi[rend={italic}]{All the Things You Are} (et il y en a !) déplient ce thème d’un certain point de vue. Tous les points de vue sont vrais. Il n’y en a pas de meilleurs. Mais il y en a de plus beaux que d’autres\teiNote[xmlid={ftn39-3},n={16},place={foot}]{\teiP[xmlid={p-505}]{Colas Duflo et Pierre Sauvanet, \teiHi[rend={italic}]{Jazzs},\teiHi[rend={italic}]{ op. cit.}, p. 173 et 174. }}. »}
              \teiP[xmlid={p17-11}]{En ce sens, on associe le plus souvent le jazz à la question de l’improvisation. Or il ne faut pas oublier que l’improvisation ne saurait en elle-même, et prise seule en tant que telle, suffire à définir le jazz. S’il y a de nombreux enregistrements effectivement entièrement improvisés dans le jazz, en revanche il existe de tout aussi nombreuses compositions entièrement écrites, ne laissant que peu ou pas du tout de place à l’improvisation. En fait, l’improvisation n’a de sens que par rapport à la présence ou à l’absence d’une composition. Dans les extraits que nous avons choisis (que ce soit Ornette ou Monk), il y a relativement peu d’improvisation, pour la bonne raison que nous n’avons écouté que l’exposition du thème, et non les chorus. Avec davantage de temps (les deux versions faisant sept minutes chacune), nous aurions pu évidemment découvrir les improvisations qui font suite, sur la même grille harmonique, à l’exposition du thème. Comme on l’a vu plus haut, on sait d’ailleurs qu’Ornette Coleman ira bientôt jusqu’à l’abandon du thème, comme de toute grille harmonique, comme de toute pulsation rythmique, avec l’enregistrement de \teiHi[rend={italic}]{Free Jazz}. À peu de choses près, pourrions-nous donc aller jusqu’à parler de « beauté libre » parmi les beautés jazz(s) ?}
              \teiP[xmlid={p18-11}]{Évidemment, pour un philosophe, une telle appellation est pour le moins connotée, et ce peut être justement l’occasion de revenir, même brièvement, sur un paragraphe de l’histoire de la philosophie qui a déjà fait couler beaucoup d’encre. On sait en effet que, au paragraphe 16 de la \teiHi[rend={italic}]{Critique de la faculté de juger}, Kant distingue beauté libre et beauté adhérente. La seconde est toujours dépendante d’un concept pour pouvoir exister (par exemple, l’architecture dit et dicte au moins en partie ce que le bâtiment doit être, non seulement beau, mais fonctionnel), tandis que la première doit pouvoir s’en passer : par exemple, la beauté d’une fleur est sans pourquoi, mais aussitôt Kant ajoute qu’« on peut encore ranger dans ce genre tout ce que l’on nomme en musique improvisation (sans thème) et même toute la musique sans texte\teiNote[xmlid={ftn40-3},n={17},place={foot}]{\teiP[xmlid={p-506}]{Kant, \teiHi[rend={italic}]{Critique de la faculté de juger }\lbrack{}1790\rbrack{}, Alexis Philonenko (trad.), Paris, Vrin, 1986, § 16, p. 71. Depuis quelques précieux travaux sur le sujet, on connaît mieux aujourd’hui les difficultés d’interprétation de ce redoutable paragraphe. D’une part, Bernard Sève précise en effet que ladite « improvisation » dans le texte français n’est autre que la « fantaisie » (\teiHi[rend={italic}]{Phantasie}) dans le texte allemand : « au sens où l’entend Kant, ce sont les préludes non mesurés de Couperin, les \teiHi[rend={italic}]{Fantasia} des maîtres italiens, peut-être la \teiHi[rend={italic}]{Fantaisie en sol mineur }BWV 542 pour orgue de Bach, ce sont les formes les plus ouvertes, les plus imprévisibles ». Cela dit, « improviser » est bien « un sens moderne de \teiHi[rend={italic}]{phantasieren}, et les fantaisies ou toccatas se donnent comme des improvisations écrites dont l’unité ne peut être que de nature sensible » (\teiHi[rend={italic}]{L’Altérité musicale }\lbrack{}2002\rbrack{}, Paris, Éditions du Seuil, 2013, p. 232 et 233). D’autre part, Sabine Forero-Mendoza et Pierre Montebello précisent à leur tour que « la fantaisie sans thème serait donc, pour Kant, un meilleur exemple de forme pure que la fantaisie à partir d’un thème. Mais la suite de l’énumération qui mentionne “toute la musique sans texte” vient nuancer le propos et élargir considérablement le champ de la référence », du côté de la musique instrumentale en général (\teiHi[rend={italic}]{Kant, son esthétique. Entre mythes et récits}, Dijon, Les Presses du réel, 2013, p. 57 et 58).}} ». En d’autres termes, la musique instrumentale elle-même serait une sorte de « beauté libre », et notamment à travers l’improvisation, qui permet à la musique de ne pas être conditionnée par autre chose qu’elle-même. Attention, bien entendu, à faire la part de l’anachronisme, car il va de soi que Kant ne désigne pas le jazz ici ! De même qu’il ne désigne pas l’art abstrait quand il parle un peu plus haut de motifs décoratifs qui ne représentent rien en eux-mêmes. En même temps, au-delà de cette absurdité chronologique, le musicologue et le mélomane savent bien que la notion même d’improvisation n’est pas limitée dans le temps, et qu’un fil invisible peut tout à fait relier dans l’histoire la basse continue baroque et la ligne de (contre)basse dans le jazz. \teiHi[rend={italic}]{Mutatis mutandis}, donc, et toutes choses égales par ailleurs, nous serions prêts à défendre l’idée qu’il y va bien d’une certaine « beauté libre » dans le jazz de Thelonious Monk ou d’Ornette Coleman – et ce, dans tous les sens du mot.}
            
\end{teiDiv}
            \begin{teiDiv}[type={section1},xmlid={div03-8}]

              \teiHead[xmlid={beaute-et-vitalite-001}]{Beauté et vitalité}
              \teiP[xmlid={p19-11}]{Comme la référence à Spinoza nous l’indiquait déjà, et comme l’idée même de liberté nous y invite maintenant, il est temps de passer de l’esthétique à l’éthique, des enjeux musicaux aux enjeux vitaux. En réalité, quand on pense jazz, quand on vit jazz, c’est à peu de choses près la même chose. Sans avoir cette prétention, mais en assumant ici la première personne du singulier, je sais que l’école du jazz me nourrit sans cesse, et je peux même dire qu’elle nourrit chaque instant de ma vie : non seulement au sens de la musique que j’écoute (pas un jour sans jazz, comme un médicament sans ordonnance médicale pour aller mieux), mais au sens de la musique qui m’inspire (une certaine liberté d’esprit, ou plus concrètement : le principe même de l’improvisation dans le cadre est quelque chose que je pratique régulièrement dans la vie quotidienne).}
              \teiP[xmlid={p20-11}]{On connaît la définition de la beauté selon Stendhal : « une promesse de bonheur\teiNote[xmlid={ftn41-4},n={18},place={foot}]{\teiP[xmlid={p-507}]{Cité notamment par Nietzsche (et en français dans le texte) dans le célèbre paragraphe de \teiHi[rend={italic}]{La Généalogie de la morale }où il oppose les deux conceptions kantiennes et stendhaliennes, emblèmes respectifs d’une esthétique du spectateur et d’une esthétique du créateur : « “Est beau, dit Kant, ce qui provoque un plaisir \teiHi[rend={italic}]{désintéressé}”. – Désintéressé ! Comparez avec cette définition cette autre, d’un véritable “spectateur”, et d’un artiste – Stendhal, qui appelle quelque part la beauté \teiHi[rend={italic}]{une promesse de bonheur}. En tout cas, ici est récusé et rayé le seul aspect du fait esthétique que Kant mette en relief : le \teiHi[rend={italic}]{désintéressement}. Qui a raison, Kant ou Stendhal ? – Assurément, lorsque nos esthéticiens, en faveur de Kant, ne se lassent pas de faire valoir le fait que sous la fascination de la beauté on peut contempler d’une façon “désintéressée” \teiHi[rend={italic}]{même} des statues de femmes nues, on est bien en droit de rire un peu à leurs dépens : – sur ce point délicat, les expériences des artistes sont “plus intéressantes”, et Pygmalion, en tout cas, \teiHi[rend={italic}]{n’était pas} nécessairement un “homme inesthétique” » (Nietzsche, \teiHi[rend={italic}]{La Généalogie de la morale} (1887\rbrack{}\teiHi[rend={italic}]{,} Giorgio Colli et Mazzino Montinari éd., Paris, Gallimard, « Folio Essais », 1966, troisième dissertation, chap. 6, p. 120).}} », citée notamment par Nietzsche dans un développement fameux où il l’oppose au désintéressement de Kant. Si la beauté est bien une promesse de bonheur, on pense ici, pour finir, au livre exceptionnel de la baronne Pannonica de Koenigswarter, \teiHi[rend={italic}]{Les Musiciens de jazz et leurs trois vœux}. Née à Londres en 1913, dans la branche anglaise de la famille Rothschild, ladite baronne fut pendant quelques décennies l’égérie des musiciens de jazz new-yorkais. Ex-épouse d’un millionnaire, elle s’était prise de passion pour la musique noire américaine. Toujours présente dans les clubs, elle devint vite la confidente de quelques-uns des grands noms du jazz, à commencer par Thelonious Monk. Sa maison (\teiHi[rend={italic}]{Cathouse}, pour les « cats », joueurs de jazz, mais aussi les quelque cent vingt chats qui arpentaient la propriété) devint le refuge de tous ceux qui vivaient plus ou moins bien de leur instrument. Durant de nombreuses années, elle recueillit les réponses à une seule et unique question posée à des centaines de musiciens : « Si on t’accordait trois vœux qui devaient se réaliser sur-le-champ, que souhaiterais-tu ? » Elle a d’abord posé la question à Thelonious Monk, qui a répondu :}
              \begin{teiCit}[xmlid={cit01-9}]

                \begin{teiQuote}
1. Que ma musique ait du succès.
\end{teiQuote}
              
\end{teiCit}
              \begin{teiCit}[xmlid={cit02-8}]

                \begin{teiQuote}
2. Que ma famille soit heureuse.
\end{teiQuote}
              
\end{teiCit}
              \begin{teiCit}[xmlid={cit03-8}]

                \begin{teiQuote}
3. Qu’on me donne une amie géniale comme toi !
\end{teiQuote}
              
\end{teiCit}
              \teiP[xmlid={p21-9}]{Et Pannonica d’ajouter : « “Mais Thelonious ! Tout ça, tu l’as déjà !” Il s’est contenté de sourire, et il s’est remis à marcher de long en large\teiNote[xmlid={ftn42-4},n={19},place={foot}]{\teiP[xmlid={p-508}]{Pannonica de Koenigswarter, \teiHi[rend={italic}]{Les Musiciens de jazz et leurs trois vœux}, Paris, Buchet-Chastel, 2006, n\teiHi[rend={sup}]{o} 1, p. 32.}}… » Magnifique philosophie de vie, de la part de quelqu’un qui souhaite au fond ce qu’il a déjà, sans peut-être s’en rendre tout à fait compte lui-même. Parmi les deux cent quatre-vingt-dix-neuf autres réponses, on relèvera également celle d’Ornette Coleman (déjà présent à l’époque), qui ne fait pas dans le détail :}
              \begin{teiCit}[xmlid={cit04-7}]

                \begin{teiQuote}
1. La vie éternelle.
\end{teiQuote}
              
\end{teiCit}
              \begin{teiCit}[xmlid={cit05-6}]

                \begin{teiQuote}
2. L’amour.
\end{teiQuote}
              
\end{teiCit}
              \begin{teiCit}[xmlid={cit06-5}]

                \begin{teiQuote}
3. Le bonheur\teiNote[xmlid={ftn43-4},n={20},place={foot}]{\teiP[xmlid={p-509}]{\teiHi[rend={italic}]{Ibid.}, n\teiHi[rend={sup}]{o} 10, p. 61.}}.
\end{teiQuote}
              
\end{teiCit}
              \teiP[xmlid={p22-9}]{et enfin, celles qui ont trait à la beauté, notamment les deux suivantes :}
              \begin{teiCit}[xmlid={cit07-3}]

                \begin{teiQuote}
1.Justice.
\end{teiQuote}
              
\end{teiCit}
              \begin{teiCit}[xmlid={cit08-3}]

                \begin{teiQuote}
2. Vérité.
\end{teiQuote}
              
\end{teiCit}
              \begin{teiCit}[xmlid={cit09-2}]

                \begin{teiQuote}
3. Beauté \lbrack{}selon Stan Getz\teiNote[xmlid={ftn44-4},n={21},place={foot}]{\teiP[xmlid={p-510}]{\teiHi[rend={italic}]{Ibid.}, n\teiHi[rend={sup}]{o} 136, p. 172.}}\rbrack{}.
\end{teiQuote}
              
\end{teiCit}
              \begin{teiCit}[xmlid={cit10-2}]

                \begin{teiQuote}
Je souhaiterais que la musique devienne encore plus belle qu’elle ne l’est déjà, et pouvoir en écouter sans cesse, en composer sans arrêt \lbrack{}selon Billy Strayhorn, pianiste, compositeur, arrangeur, très proche de Duke Ellington, au point de passer pour son \teiHi[rend={italic}]{alter ego}\teiNote[xmlid={ftn45-4},n={22},place={foot}]{\teiP[xmlid={p-511}]{\teiHi[rend={italic}]{Ibid.}, n\teiHi[rend={sup}]{o} 175, p. 198.}}\rbrack{}.
\end{teiQuote}
              
\end{teiCit}
              \teiP[xmlid={p23-9}]{On aura remarqué que ce dernier se limite volontairement à un seul vœu, comme s’il englobait tout le reste : la beauté de la musique elle-même. Mais de quelle beauté parlons-nous au juste ? En guise de réponse, et de conclusion, il suffirait presque de ne faire qu’un de nos deux titres, en mettant les deux versos de nos pochettes de disque côte à côte : \teiHi[rend={italic}]{Ugly Beauty Is a Rare Thing}. En réalité, la beauté dont nous avons parlé ici, ou plus précisément les beautés jazz(s) au double pluriel, si elles sont relativement « rares » (les amateurs de jazz sont de moins en moins légion), ne sont pas exactement « laides ». On l’a dit en commençant, le type de beauté jazz auquel nous pensons, finalement, est toujours complexe, mélangé, métissé : pas de beauté pure, pas de beauté simple, mais une beauté « tendue » (musicalement, entre deux notes, entre deux silences). Une beauté exigeante qui peut mettre du temps à se révéler, mais qui, une fois présente, ne vous quitte plus. La beauté vitale en jazz, c’est celle qui force Thelonious à se lever de son piano pour danser sur la musique, ventre en avant, pieds en canard, et pourtant léger comme une plume. La beauté vitale en jazz, c’est celle qui fait que je peux pleurer régulièrement, sans raison, en écoutant un thème de Monsieur Thelonious «\teiHi[rend={italic}]{ Sphere }» Monk (son deuxième prénom). Pourtant, je peux bien avouer pour finir que je n’ai pas réellement aimé Monk la première fois que je l’ai écouté – \teiHi[rend={italic}]{idem} avec Ornette Coleman, qui était même bien pire à mes pauvres oreilles d’alors. Et quand on pense que le propre fils de Monk a mis vingt et un ans avant de trouver la musique de son père « belle\teiNote[xmlid={ftn46-4},n={23},place={foot}]{\teiP[xmlid={p-512}]{Voir Laurent de Wilde, \teiHi[rend={italic}]{op. cit.}, p. 286.}} »… Avec le recul, c’est peut-être cela le secret de ces beautés vitales, de celles qui recèlent de la vie comme de celles qui aident à vivre : elles demandent fondamentalement du \teiHi[rend={italic}]{temps}. C’est pourquoi nous terminerons avec Nietzsche, dont cet aphorisme trop peu connu mérite d’être cité en entier :}
              \begin{teiCit}[xmlid={cit11-2}]

                \begin{teiQuote}
Voici ce qui nous arrive dans le domaine musical : il faut avant tout \teiHi[rend={italic}]{apprendre à entendre} une figure, une mélodie, savoir la discerner par l’ouïe, la distinguer, l’isoler et la délimiter en tant qu’une vie en soi : ensuite il faut de l’effort et de la bonne volonté pour la \teiHi[rend={italic}]{supporter}, en dépit de son étrangeté, user de patience pour son regard et pour son expression, de tendresse pour ce qu’elle a de singulier ; – vient enfin le moment où nous y sommes \teiHi[rend={italic}]{habitués}, où nous l’attendons, où nous sentons qu’elle nous manquerait, si elle faisait défaut ; et désormais elle ne cesse pas d’exercer sur nous sa contrainte et sa fascination jusqu’à ce qu’elle ait fait de nous ses amants humbles et ravis, qui ne conçoivent de meilleure chose au monde et ne désirent plus qu’elle-même, et rien qu’elle-même. – Mais ce n’est pas seulement en musique que ceci nous arrive : c’est justement de la sorte que nous avons \teiHi[rend={italic}]{appris à aimer }tous les objets que nous aimons maintenant. Nous finissons toujours par être récompensés pour notre bonne volonté, notre patience, notre équité, notre tendresse envers l’étrangeté, du fait que l’étrangeté peu à peu se dévoile et vient s’offrir à nous en tant que nouvelle et indicible beauté : – c’est là sa \teiHi[rend={italic}]{gratitude} pour notre hospitalité. Qui s’aime soi-même n’y sera parvenu que par cette voie : il n’en est point d’autre. L’amour aussi doit s’apprendre\teiNote[xmlid={ftn47-4},n={24},place={foot}]{\teiP[xmlid={p-513}]{Nietzsche, \teiHi[rend={italic}]{Le gai savoir} \lbrack{}1872\rbrack{}, Giorgio Colli et Mazzino Montinari (éd.), Paris, Gallimard, « Folio Essais », 1982, § 334, p. 223. On trouvera un commentaire de ce texte dans Pierre Sauvanet, \teiHi[rend={italic}]{Les philosophes et l’amour }\lbrack{}1998\rbrack{}, Paris, Ellipses, 2015, p. 149-162.}}.
\end{teiQuote}
              
\end{teiCit}
            
\end{teiDiv}
          
\end{teiElement}
        
\end{teiElement}
        \begin{teiElement}[name={group},type={chapter},data-page-title={La beauté, une expérience socialement définie ?},data-page-subtitle={Sociologie de la perception musicale},data-page-authors={Guilhem Marion@@New York University, Psychology Department, États-Unis},data-include-href={Ch13\_beaute\_experience\_Marion.xml},xmlid={chapter-012}]

          \begin{teiElement}[name={front}]

            \begin{teiDiv}[type={titlePage},xmlid={titlepage-014}]

              \teiP[rend={title-main},xmlid={p-514}]{La beauté, une expérience socialement définie ?}
              \teiP[rend={title-sub},xmlid={p-515}]{Sociologie de la perception musicale}
              \teiP[rend={author-aut},xmlid={p-516}]{Guilhem \teiHi[rend={small-caps}]{Marion}}
              \teiP[rend={authority\_affiliation},xmlid={p-517}]{New York University, Psychology Department, États-Unis}
            
\end{teiDiv}
          
\end{teiElement}
          \begin{teiElement}[name={body},xmlid={body-014}]

            \teiP[xmlid={p1-14}]{Il est évident que les préférences musicales et, plus généralement, les comportements liés à l’écoute de la musique diffèrent d’un individu à l’autre. Les plateformes de streaming peuvent aujourd’hui fournir des bases de données pour étudier ces questions et montrent des différences radicales entre les individus dont une grande partie de la variance peut être expliquée par des facteurs socioculturels et démographiques.}
            \teiP[xmlid={p2-14}]{Tout d’abord, les auditeurs préfèrent la musique à laquelle ils s’identifient culturellement, comme mesurés par l’identification ethnique\teiNote[n={1},place={foot},xmlid={ftn1-13}]{\teiP[xmlid={p-518}]{Jourdan McCrary et Donna Gauthier, « The Effects of Performers’ Ethnic Identities on Preadolescents’ Music Preferences », \teiHi[rend={italic}]{Update: Applications of Research in Music Education}, vol. 14, n\teiHi[rend={sup}]{o} 1, 1995, p. 20-22.}} ou géographique\teiNote[n={2},place={foot},xmlid={ftn58-4}]{\teiP[xmlid={p-519}]{Charlotta Mellander, Richard Florida, Peter J. Rentfrow \teiHi[rend={italic}]{et al}., « The Geography of Music Preferences », \teiHi[rend={italic}]{Journal of Cultural Economics}, vol. 42, 2018, p. 593-618.}}. Deux modèles ont été présentés pour expliquer la relation entre le milieu socioéconomique et les préférences musicales au sein d’un même pays. Premièrement, le modèle de la légitimité culturelle a été présenté et validé en France\teiNote[n={3},place={foot},xmlid={ftn59-4}]{\teiP[xmlid={p-520}]{Pierre Bourdieu, \teiHi[rend={italic}]{La Distinction. Critique sociale du jugement}, Paris, Éditions de Minuit, « Le sens commun », 1979.}} et aux États-Unis\teiNote[n={4},place={foot},xmlid={ftn60-4}]{\teiP[xmlid={p-521}]{Karl F. Schuessler, « Social Background and Musical Taste », \teiHi[rend={italic}]{American Sociological Review}, vol. 13, n\teiHi[rend={sup}]{o} 3, 1948, p. 330-335.}} avant les années 1980. Selon ce modèle, les classes sociales supérieures consomment essentiellement de la musique considérée comme plus sophistiquée ou associée à une valeur culturelle plus élevée (comme la musique classique, la musique contemporaine et l’opéra) et rejettent les autres formes de musique populaire, ce qui relève d’un mécanisme de légitimation de l’appartenance à leur groupe social. Cependant, depuis les années 1990, un autre modèle appelé omnivore/univore a été présenté\teiNote[n={5},place={foot},xmlid={ftn61-4}]{\teiP[xmlid={p-522}]{Richard A. Peterson, « Understanding Audience Segmentation: From Elite and Mass to Omnivore and Univore », \teiHi[rend={italic}]{Poetics}, vol. 21, n\teiHi[rend={sup}]{o} 4, 1992, p. 243-258.}} et validé empiriquement par des études en France\teiNote[n={6},place={foot},xmlid={ftn62-4}]{\teiP[xmlid={p-523}]{Philippe Coulangeon, « Social Stratification of Musical Tastes: Questioning the Cultural Legitimacy Model », \teiHi[rend={italic}]{Revue française de sociologie}, vol. 46, n\teiHi[rend={sup}]{o} 5, 2005, p. 123-154.}}, aux États-Unis\teiNote[n={7},place={foot},xmlid={ftn63-4}]{\teiP[xmlid={p-524}]{Richard A. Peterson et Roger M. Kern, « Changing Highbrow Taste: From Snob to Omnivore », \teiHi[rend={italic}]{American Sociological Review}, vol. 61, n\teiHi[rend={sup}]{o} 5, 1996, p. 900-907.}} et aux Pays-Bas\teiNote[n={8},place={foot},xmlid={ftn64-4}]{\teiP[xmlid={p-525}]{Koen Van Eijck, « Social Differentiation in Musical Taste Patterns », \teiHi[rend={italic}]{Social Forces}, vol. 79, n\teiHi[rend={sup}]{o} 3, 2001, p. 1163-1185.}}. Cette hypothèse alternative affirme que les classes sociales se différencient par la variété de leurs préférences musicales plutôt que par un genre spécifique, avec une tendance des classes supérieures à écouter un éventail plus large de genres musicaux. Des études\teiNote[n={9},place={foot},xmlid={ftn65-5}]{\teiP[xmlid={p-526}]{Christina G. White, \teiHi[rend={italic}]{The Effects of Class, Age, Gender and Race on Musical Preferences: an Examination of the Omnivore/Univore Framework}, thèse, Blacksburg, Virginia Tech, 2001, p. 11.}} utilisant des données entre 1982 et 1997 aux États-Unis montrent une tendance au remplacement du modèle de légitimité culturelle par le modèle omnivore/univore : « l’exclusivité culturelle n’est plus appréciée comme elle a pu l’être dans le passé et elle est plus souvent un signe d’ignorance que de statut\teiHi[rend={sup}]{9}. »}
            \teiP[xmlid={p3-14}]{Cependant, il a aussi été démontré en France en 2003 que les classes supérieures, même si elles se distinguent principalement par leur omnivorité, ont également tendance à écouter plus de musique savante (musique classique et opéra principalement) que les autres\teiNote[n={10},place={foot},xmlid={ftn66-4}]{\teiP[xmlid={p-527}]{Philippe Coulangeon, art. cité.}}, ce qui montre que, même si le modèle le plus explicatif est passé de la légitimité culturelle à l’omnivorité/univorité, les deux modèles cohabitent toujours. L’hypothèse de l’accord social pourrait expliquer pourquoi les individus appartenant à un même groupe social partagent des préférences musicales similaires. Il a en effet été démontré que les auditeurs modifient leurs préférences musicales pour être en accord avec leurs pairs\teiNote[n={11},place={foot},xmlid={ftn67-4}]{\teiP[xmlid={p-528}]{Craig E. Furman et Robert A. Duke, « Effect of Majority Consensus on Preferences for Recorded Orchestral and Popular Music », \teiHi[rend={italic}]{Journal of Research in Music Education}, vol. 36, n\teiHi[rend={sup}]{o} 4, 1988, p. 220-231.}} ou pour éviter les désaccords\teiNote[n={12},place={foot},xmlid={ftn68-3}]{\teiP[xmlid={p-529}]{Harold G. Inglefield, « Conformity Behavior Reflected in the Musical Preference of Adolescents », \teiHi[rend={italic}]{Contributions to Music Education}, n\teiHi[rend={sup}]{o} 1, 1972, p. 56-67.}}, en particulier lorsque les pairs représentent une figure d’autorité en termes de goûts musicaux\teiNote[n={13},place={foot},xmlid={ftn69-3}]{\teiP[xmlid={p-530}]{Jesse Alpert, « The Effect of Disc Jockey, Peer, and Music Teacher Approval of Music on Music Selection and Preference », \teiHi[rend={italic}]{Journal of Research in Music Education}, vol. 30, n\teiHi[rend={sup}]{o} 3, 1982, p. 173-186.}}. Cette idée suit l’idée générale selon laquelle la familiarité socioculturelle augmente les préférences musicales et suscite des réponses émotionnelles plus fortes.}
            \teiP[xmlid={p4-14}]{Les auditeurs ont, en effet, tendance à préférer un contenu musical déjà connu, comme le montre la répétition de courts extraits de musique classique, y compris de compositions tonales et atonales\teiNote[n={14},place={foot},xmlid={ftn70-3}]{\teiP[xmlid={p-531}]{Elizabeth Hellmuth Margulis, « Aesthetic Responses to Repetition in Unfamiliar Music », \teiHi[rend={italic}]{Empirical Studies of the Arts}, vol. 31, n\teiHi[rend={sup}]{o} 1, 2013, p. 45-57.}}, de musique de jazz\teiNote[n={15},place={foot},xmlid={ftn71-3}]{\teiP[xmlid={p-532}]{E. M. Verveer, H. Barry et W. A. Bousfield, « Change in Affectivity with Repetition », \teiHi[rend={italic}]{American Journal of Psychology}, vol. 45, n\teiHi[rend={sup}]{o} 1, 1933, p. 130-134.}}, de musique coréenne\teiNote[n={16},place={foot},xmlid={ftn72-3}]{\teiP[xmlid={p-533}]{Marcia K. Johnson, Jung K. Kim et Gail Risse, « Do Alcoholic Korsakoff’s Syndrome Patients Acquire Affective Reactions? », \teiHi[rend={italic}]{Journal of Experimental Psychology: Learning, Memory, and Cognition}, vol. 11, n\teiHi[rend={sup}]{o} 1, 1985, p. 22-36.}}, pakistanaise\teiNote[n={17},place={foot},xmlid={ftn73-4}]{\teiP[xmlid={p-534}]{Annelies Heingartner et James V Hall, « Affective Consequences in Adults and Children of Repeated Exposure to Auditory Stimuli », \teiHi[rend={italic}]{Journal of Personality and Social Psychology}, vol. 29, n\teiHi[rend={sup}]{o} 6, 1974, p. 719-723.}}, etc. Dans ces expériences, de courts extraits sont répétés plusieurs fois et sont préférés aux versions complètes des mêmes pièces musicales. Au-delà de la reconnaissance exacte d’un fragment musical, la familiarité peut également être définie à l’aide de l’idée de prévisibilité. Les auditeurs, lorsqu’on leur présente un nouveau contenu musical, préfèrent des morceaux contenant des structures musicales proches de celles issues de leur propre culture musicale. Par exemple, les enfants américains préfèrent (en mesurant le temps de regard) les métriques rythmiques occidentales (mesures composées de 3 ou 4 temps) aux métriques balkaniques (en 7, 9 ou 10 temps), tandis que les enfants turcs, qui connaissent à la fois les métriques occidentales et balkaniques, ne manifestent aucune préférence\teiNote[n={18},place={foot},xmlid={ftn74-4}]{\teiP[xmlid={p-535}]{Gaye Soley et Erin Hannon, « Infants Prefer the Musical Meter of their Own Culture: a Cross-Cultural Comparison », \teiHi[rend={italic}]{Developmental Psychology}, vol. 46, n\teiHi[rend={sup}]{o} 1, 2010, p. 286-292.}}. Il a même été démontré que la précision de la mémoire était influencée par la familiarité culturelle : les participants occidentaux se souvenaient mieux des mélodies occidentales que les participants turcs, et inversement pour les mélodies turques. En revanche, les mélodies chinoises (utilisées comme contrôle) entraînaient une faible précision de la mémoire pour les deux groupes\teiNote[n={19},place={foot},xmlid={ftn75-4}]{\teiP[xmlid={p-536}]{Steven M. Demorest, Steven J. Morrison, Münir N. Beken \teiHi[rend={italic}]{et al}., « Lost in Translation: an Enculturation Effect in Music Memory Performance », \teiHi[rend={italic}]{Music Perception}, vol. 25, n\teiHi[rend={sup}]{o} 3, 2008, p. 213-223.}}. Ces résultats ont été répliqués avec des séquences composées spécifiquement pour correspondre aux structures musicales des auditeurs\teiNote[n={20},place={foot},xmlid={ftn76-4}]{\teiP[xmlid={p-537}]{Kat Agres, Samer Abdallah et Marcus Pearce, « Information-Theoretic Properties of Auditory Sequences Dynamically Influence Expectation and Memory », \teiHi[rend={italic}]{Cognitive Science}, vol. 42, n\teiHi[rend={sup}]{o} 1, 2018, p. 43-76.}}.}
            \begin{teiDiv}[type={section1},xmlid={div01-10}]

              \teiHead[xmlid={la-theorie-du-traitement-predictif-001}]{La théorie du traitement prédictif}
              \teiP[xmlid={p5-14}]{Si la perception implique clairement le traitement d’informations sensorielles, les entrées sensorielles ne suffisent pas à rendre compte de l’ensemble des processus perceptifs. En effet, des stimuli sensoriels (par exemple, un son ou une image) peuvent être perçus différemment par différentes personnes\teiNote[n={21},place={foot},xmlid={ftn77-5}]{\teiP[xmlid={p-538}]{David H. Brainard et Anya C. Hurlbert, « Colour Vision: Understanding the Dress », \teiHi[rend={italic}]{Current Biology}, vol. 25, n\teiHi[rend={sup}]{o} 13, 2015, p. R551-R554 ; Daniel Pressnitzer, Jackson Graves, Claire Cambers \teiHi[rend={italic}]{et al}., « Auditory Perception: Laurel and Yanny Together at Last », \teiHi[rend={italic}]{Current Biology}, vol. 28, n\teiHi[rend={sup}]{o} 13, 2018, p. R739-R741.}} ou même par la même personne, mais dans des conditions différentes\teiNote[n={22},place={foot},xmlid={ftn78-5}]{\teiP[xmlid={p-539}]{Joel S. Snyder, Caspar M. Schwiedrzik, A. Davi Vitela \teiHi[rend={italic}]{et al}., « How Previous Experience Shapes Perception in Different Sensory Modalities », \teiHi[rend={italic}]{Frontiers in Human Neuroscience}, vol. 9, 2015, p. 594.}}. Les réseaux sociaux regorgent de tels exemples, dont \teiHi[rend={italic}]{\#theDress} pour la vision et \teiHi[rend={italic}]{\#Laurel \#Yanny} pour l’audition. Une explication de ce phénomène est fournie par la théorie du traitement prédictif\teiNote[n={23},place={foot},xmlid={ftn79-5}]{\teiP[xmlid={p-540}]{Karl J. Friston, Jean Daunizeau, James Kilner \teiHi[rend={italic}]{et al}., « Action and Behavior: A Free-Energy Formulation », \teiHi[rend={italic}]{Biological Cybernetics}, vol. 102, n\teiHi[rend={sup}]{o} 3, 2010, p. 227-260 ; Andy Clark, « Whatever next? Predictive Brains, Situated Agents, and the Future of Cognitive Science », \teiHi[rend={italic}]{Behavioral and Brain Sciences}, vol. 36, n\teiHi[rend={sup}]{o} 3, 2013, p. 181-204.}}. Ce cadre théorique s’articule autour de l’idée que le cerveau développe un modèle du monde qui est utilisé pour prédire les entrées sensorielles et qui est continuellement mis à jour en comparant les \teiHi[rend={italic}]{stimuli} prédits et perçus\teiNote[n={24},place={foot},xmlid={ftn80-4}]{\teiP[xmlid={p-541}]{Horace B. Barlow, « Possible Principles Underlying the Transformation of Sensory Messages », dans Walter A. Rosenblith (dir.), \teiHi[rend={italic}]{Sensory Communication}, Cambridge, MIT Press, 1961, p. 217-234.}}. Ainsi, la perception émerge de l’interaction entre les entrées sensorielles (S) et les attentes ou prédictions internes (P). La comparaison entre les entrées sensorielles et leur prédiction produit une erreur de prédiction (EP = S-P) qui, entre autres fonctions, permet la mise à jour du modèle de prédiction interne lui-même\teiNote[n={25},place={foot},xmlid={ftn81-4}]{\teiP[xmlid={p-542}]{Risto Näätänen, Petri Paavilainen, Teemu Rinne \teiHi[rend={italic}]{et al.}, « The Mismatch Negativity (Mmn) in Basic Research of Central Auditory Processing: A Review », \teiHi[rend={italic}]{Clinical Neurophysiology}, vol. 102, n\teiHi[rend={sup}]{o} 3, 2007, p. 227-260.}}. Par conséquent, la perception est un processus actif par lequel notre cerveau surveille en permanence les statistiques des informations sensorielles entrantes afin (i) d’apprendre et de mettre à jour un modèle interne des régularités du monde qui nous entoure ; et (ii) de prédire, sur la base de ce modèle, les entrées sensorielles afin de moduler leur encodage neuronal et de faciliter leur perception dans des conditions difficiles\teiNote[n={26},place={foot},xmlid={ftn82-4}]{\teiP[xmlid={p-543}]{Matthew K. Leonard, Maxime O. Baud, Matthias J. Sjerps \teiHi[rend={italic}]{et al}., « Perceptual Restoration of Masked Speech in Human Cortex », \teiHi[rend={italic}]{Nature Communications}, vol. 7, n\teiHi[rend={sup}]{o} 1, 2016, p. 1-9.}} ou en biaisant la perception d’images ou de sons ambigus\teiNote[n={27},place={foot},xmlid={ftn83-5}]{\teiP[xmlid={p-544}]{David H. Brainard et Anya C. Hurlbert, « Colour Vision: Understanding the Dress », \teiHi[rend={italic}]{Current Biology}, vol. 25, n\teiHi[rend={sup}]{o} 13, 2015, p. R551-R554 ; Daniel Pressnitzer, Jackson Graves, Claire Cambers \teiHi[rend={italic}]{et al}., art. cité.}}.}
              \teiP[xmlid={p6-14}]{La perception musicale, en raison de ses régularités structurelles, temporelles, spectrale, mélodiques ou harmoniques, offre un paradigme éclairant pour explorer les implémentations des principes du traitement prédictif\teiNote[n={28},place={foot},xmlid={ftn84-5}]{\teiP[xmlid={p-545}]{Stefan Koelsch, Peter Vuust et Karl Friston, « Predictive Processes and the Peculiar Case of Music », \teiHi[rend={italic}]{Trends in Cognitive Sciences}, vol. 23, n\teiHi[rend={sup}]{o} 1, 2019, p. 63-77.}}. La musique peut présenter des régularités dues à la répétition du même motif, ainsi que des régularités qui peuvent être imprévisibles en fonction du seul contexte proximal, mais qui sont néanmoins conformes aux règles d’un style musical ou d’une culture particulière\teiNote[n={29},place={foot},xmlid={ftn85-5}]{\teiP[xmlid={p-546}]{Elizabeth Hellmuth Margulis, « On Repeat », dans \teiHi[rend={italic}]{On Repeat}, Oxford, Oxford University Press, 2014.}}.}
              \teiP[xmlid={p7-14}]{Lorsqu’on écoute de la musique, les prédictions formulées se distinguent à la fois par leur contenu et par le degré de certitude qui leur est associé\teiNote[n={30},place={foot},xmlid={ftn86-5}]{\teiP[xmlid={p-547}]{Stefan Koelsch, Peter Vuust et Karl Friston, « Predictive Processes and the Peculiar Case of Music », \teiHi[rend={italic}]{Trends in Cognitive Sciences}, vol. 23, n\teiHi[rend={sup}]{o} 1, 2019, p. 63-77.}}. Par exemple, dans le cadre d’une progression d’accords tonale (le système harmonique le plus répandu en Occident décrivant l’enchaînement des accords autour d’un centre tonal), il est possible d’anticiper l’accord suivant. Ce processus implique non seulement une prédiction sur le contenu (quel accord exact sera joué ensuite), mais aussi une estimation de la confiance accordée à cette prédiction en fonction du contexte musical. Ainsi, une séquence d’accords inhabituelle peut créer un contexte imprévisible, rendant plus difficile pour l’auditeur de formuler des prédictions précises sur les événements musicaux à venir. La certitude de la prédiction module le gain associé à l’erreur de réponse\teiNote[n={31},place={foot},xmlid={ftn87-5}]{\teiP[xmlid={p-548}]{Karl Friston, « The Free-Energy Principle: A Rough Guide to the Brain? », \teiHi[rend={italic}]{Trends in Cognitive Sciences}, vol. 13, n\teiHi[rend={sup}]{o} 7, 2009, p. 293-301.}}, même dans le contexte de l’écoute de musique naturelle\teiNote[n={32},place={foot},xmlid={ftn88-5}]{\teiP[xmlid={p-549}]{Yi-Fang Hsu, Solene Le Bars, Jarmo A. Hämäläinen \teiHi[rend={italic}]{et al}., « Distinctive Representation of Mispredicted and Unpredicted Prediction Errors in Human Electroencephalography », \teiHi[rend={italic}]{Journal of Neuroscience}, vol. 35, n\teiHi[rend={sup}]{o} 43, 2015, p. 14653-14660.}} et joue donc finalement un rôle dans la manière dont les événements inattendus sont utilisés pour mettre à jour le modèle de prédiction\teiNote[n={33},place={foot},xmlid={ftn89-5}]{\teiP[xmlid={p-550}]{Stefan Koelsch, Peter Vuust et Karl Friston, « Predictive Processes and the Peculiar Case of Music », \teiHi[rend={italic}]{Trends in Cognitive Sciences}, vol. 223, n\teiHi[rend={sup}]{o} 1, 2019, p. 63-77.}}.}
            
\end{teiDiv}
            \begin{teiDiv}[type={section1},xmlid={div02-10}]

              \teiHead[xmlid={donnees-comportementales-et-neuronales-001}]{Données comportementales et neuronales}
              \teiP[xmlid={p8-14}]{Des preuves empiriques du traitement prédictif pendant l’écoute de la musique ont été observées en réponse à des \teiHi[rend={italic}]{stimuli} artificiels contenant des événements plus ou moins attendus (par exemple, un accord de dominante se résolvant soit sur une tonique, soit sur une sixte napolitaine). Lorsqu’on demande aux auditeurs de détecter des changements de timbre (la mélodie est soudainement jouée par un violon plutôt que par la flûte), des temps de réponse plus rapides sont associés à des \teiHi[rend={italic}]{stimuli} musicaux plus attendus\teiNote[n={34},place={foot},xmlid={ftn90-5}]{\teiP[xmlid={p-551}]{Barbara Tillmann, Emmanuel Bigand, Nicolas Escoffier \teiHi[rend={italic}]{et al}., « The influence of musical relatedness on timbre discrimination », \teiHi[rend={italic}]{European Journal of Cognitive Psychology}, vol. 18, n\teiHi[rend={sup}]{o} 3, 2006, p. 343-358.}}. Il a également été démontré que la précision des performances s’améliorait avec la prévisibilité des notes\teiNote[n={35},place={foot},xmlid={ftn91-5}]{\teiP[xmlid={p-552}]{Jamshed J. Bharucha et Keiko Stoeckig, « Priming of Chords: Spreading Activation or Overlapping Frequency Spectra? » \teiHi[rend={italic}]{Perception \& Psychophysics}, vol. 41, 1987, p. 519-524.}}. Il est important de noter que les auditeurs sans formation musicale formelle ont montré des effets similaires à ceux des musiciens expérimentés\teiNote[n={36},place={foot},xmlid={ftn92-5}]{\teiP[xmlid={p-553}]{Emmanuel Bigand et Marion Pineau, « Global Context Effects on Musical Expectancy », \teiHi[rend={italic}]{Perception \& Psychophysics}, vol. 59, n\teiHi[rend={sup}]{o} 7, 1997, p. 1098-1107.}}, ce qui confirme l’hypothèse selon laquelle le traitement prédictif pendant l’écoute de la musique ne nécessite pas de formation musicale formelle.}
              \teiP[xmlid={p9-13}]{De nombreuses études de neuro-imagerie et de neurophysiologie ont été menées pour mettre en évidence les marqueurs neuronaux du codage prédictif pendant l’écoute de la musique. Les premières mesures de ce type ont utilisé des séquences d’accords irréguliers comprenant des accords de sixte napolitaine qui étaient éloignés du contexte harmonique\teiNote[n={37},place={foot},xmlid={ftn93-5}]{\teiP[xmlid={p-554}]{Sakari Leino, Elvira Brattico, Mari Tervaniemi \teiHi[rend={italic}]{et al}., « Representation of Harmony Rules in the Human Brain: Further Evidence from Event-Related Potentials », \teiHi[rend={italic}]{Brain Research}, n\teiHi[rend={sup}]{o} 1142, 2007, p. 169-177.}}. Les accords irréguliers ont déclenché des réponses cérébrales avec une polarité négative, observée au maximum sur les capteurs frontaux droits, et une latence maximale d’environ 150-180 ms, bien que des latences plus longues aient également été observées\teiNote[n={38},place={foot},xmlid={ftn94-5}]{\teiP[xmlid={p-555}]{Stefan Koelsch et Juul Mulder, « Electric Brain Responses to Inappropriate Harmonies During Listening to Expressive Music », \teiHi[rend={italic}]{Clinical Neurophysiology}, n\teiHi[rend={sup}]{o} 113, 2002, p. 862-869.}}. Cette réponse (ERAN) montre que le cerveau réagit en fonction de l’adéquation contextuelle des événements musicaux\teiNote[n={39},place={foot},xmlid={ftn95-5}]{\teiP[xmlid={p-556}]{Stefan Koelsch, « Music-Syntactic Processing and Auditory Memory: Similarities and Differences between Eran and Mmn », \teiHi[rend={italic}]{Psychophysiology}, vol. 46, n\teiHi[rend={sup}]{o} 1, 2009, p. 179-190.}}. Des études plus récentes utilisent des modèles statistiques de la musique pour révéler que les notes les moins prévisibles de mélodies non modifiées génèrent des réponses neuronales plus grandes, compatibles avec l’ERAN\teiNote[n={40},place={foot},xmlid={ftn96-5}]{\teiP[xmlid={p-557}]{Giovanni Di Liberto, Claire Pelofi, Roberta Bianco \teiHi[rend={italic}]{et al}., \teiHi[rend={italic}]{« Cortical Encoding of Melodic Expectations in Human Temporal Cortex »}, \teiHi[rend={italic}]{eLife}, 2020, DOI :\teiRef[target={https://doi.org/10.7554/eLife.51784}]{ 10.7554/eLife.51784}.}}.}
            
\end{teiDiv}
            \begin{teiDiv}[type={section1},xmlid={div03-9}]

              \teiHead[xmlid={enculturation-musicale-apprendre-une-nouvelle-culture-001}]{Enculturation musicale : apprendre une nouvelle culture}
              \teiP[xmlid={p10-13}]{Si nos cerveaux émettent des prédictions en fonction de leur connaissance de notre environnement, comment apprennent-ils à prédire ? S’il est évident que l’apprentissage effréné permet de construire une connaissance d’un domaine, il existe une autre forme d’apprentissage qui se produit sans effort : l’apprentissage implicite. « L’apprentissage implicite est l’apprentissage d’informations complexes de manière non intentionnelle, sans conscience de ce qui a été appris\teiNote[n={41},place={foot},xmlid={ftn97-5}]{\teiP[xmlid={p-558}]{Arthur S. Reber « Implicit Learning and Tacit Knowledge », \teiHi[rend={italic}]{Journal of Experimental Psychology: General}, vol. 118, n\teiHi[rend={sup}]{o} 3, 1989, p. 219.}} ».}
              \teiP[xmlid={p11-13}]{Des exemples tirés de la vie quotidienne sont souvent utilisés, comme apprendre à faire du vélo ou à nager. Cet apprentissage se distingue de l’apprentissage explicite par l’absence de connaissances accessibles consciemment. Les recherches sur l’amnésie montrent souvent que l’apprentissage implicite est intact, mais que l’apprentissage explicite est altéré. Cette forme d’apprentissage est celle mise en avant dans la théorie du traitement prédictif comme forme d’apprentissage des régularités du monde extérieur utilisées pour prédire par la suite. Étant donné que les prédictions émises sont culturellement définies et conformes à la musique issue de la culture des auditeurs, cet apprentissage implicite des structures musicales est appelé \teiHi[rend={italic}]{enculturation musicale}.}
              \teiP[xmlid={p12-13}]{Il a été démontré que les enfants apprennent naturellement la musique de leur propre culture, tout comme ils apprennent leur langue maternelle\teiNote[n={42},place={foot},xmlid={ftn98-5}]{\teiP[xmlid={p-559}]{Catherine E. Snow, « Mothers’ Speech to Children Learning Language », \teiHi[rend={italic}]{Child Development}, vol. 43, n\teiHi[rend={sup}]{o} 2, 1972, p. 549-565.}}, grâce à un mélange d’engagement actif et d’exposition passive et sans nécessiter d’instructions explicites\teiNote[n={43},place={foot},xmlid={ftn99-5}]{\teiP[xmlid={p-560}]{Patricia Shehan Campbell, \teiHi[rend={italic}]{Songs in Their Heads: Music and Its Meaning in Children’s Lives}, Oxford, Oxford University Press, 2010.}}. Des chercheurs ont utilisé des mélodies en métrique occidentale ou balkanique et ont ajouté ou enlevé un temps à certaines mesures afin de disrupter leur structure rythmique. Les enfants de plus d’un an et adultes occidentaux étaient capables de détecter des irrégularités métriques dans des morceaux de musique à métrique occidentale, mais pas dans des morceaux à métrique balkanique. Néanmoins, après seulement une semaine d’écoute de musique balkanique, tous les participants étaient capables de détecter les irrégularités dans les morceaux à métriques occidentales et balkaniques\teiNote[n={44},place={foot},xmlid={ftn100-5}]{\teiP[xmlid={p-561}]{Erin E. Hannon et Sandra E. Trehub, « Tuning in to musical rhythms: Infants learn more readily than adults », \teiHi[rend={italic}]{Proceedings of the National Academy of Sciences}, vol. 102, n\teiHi[rend={sup}]{o} 35, 2005, p. 12639-12643.}}. Il est à noter que les enfants de moins de 12 mois n’étaient capables de détecter les irrégularités dans aucune musique. Il semblerait donc que les mécanismes de prédiction (voire peut-être d’enculturation) ne soient effectifs qu’à l’âge de 12 mois. Il a été récemment démontré qu’un tel apprentissage explicite affectait directement la façon dont les adultes prédisent les notes de musique au sein de mélodies\teiNote[n={45},place={foot},xmlid={ftn101-5}]{\teiP[xmlid={p-562}]{Psyche Loui, David Wessel et Carla Hudson Kam, « Acquiring New Musical Grammars: A Statistical Learning Approach », \teiHi[rend={italic}]{Proceedings of the Annual Meeting of the Cognitive Science Society}, n\teiHi[rend={sup}]{o} 28, 2006, p. 1711-1716. }}.}
            
\end{teiDiv}
            \begin{teiDiv}[type={section1},xmlid={div04-5}]

              \teiHead[xmlid={pourquoi-aimons-nous-la-musique-que-nous-aimons-001}]{Pourquoi aimons-nous la musique (que nous aimons) ?}
              \teiP[xmlid={p13-12}]{En 1971, Daniel Berlyne a proposé un modèle de ces mécanismes en suggérant qu’une forme en U inversé (effet de Wundt) relierait la familiarité et le plaisir musical\teiNote[n={46},place={foot},xmlid={ftn102-5}]{\teiP[xmlid={p-563}]{D. E. Berlyne, \teiHi[rend={italic}]{Aesthetics and Psychobiology}, New York : Appleton-Century-Crofts, 1971.}}. Selon lui, le système de récompense est activé lors de l’exposition à certains motifs musicaux reconnus, mais avec le temps, le système d’aversion s’oppose à cette activation, ce qui conduit à une chute du plaisir ressenti (voir graphique 22). Cet effet a été validé par de nombreuses études\teiNote[n={47},place={foot},xmlid={ftn103-5}]{\teiP[xmlid={p-564}]{Anthony Chmiel et Emery Schubert, « Back to the Inverted-U for Music Preference: A Review of the Literature », \teiHi[rend={italic}]{Psychology of Music}, vol. 6, n\teiHi[rend={sup}]{o} 8, 2017, p. 886-909.}} et l’une d’entre elles, utilisant l’IRM fonctionnelle (procédé permettant de mesurer les zones du cerveau activées), a montré que la musique familière active les régions limbiques et paralimbiques liées à l’émotion ainsi que le circuit de la récompense dans une plus grande mesure que la musique non familière\teiNote[n={48},place={foot},xmlid={ftn104-5}]{\teiP[xmlid={p-565}]{Carlos Silva Pereira, João Teixeira, Patrícia Figueiredo \teiHi[rend={italic}]{et al}., « Music and emotions in the brain: Familiarity matters », \teiHi[rend={italic}]{Plos One}, vol. 6, n\teiHi[rend={sup}]{o} 11, 2011, DOI : \teiRef[target={https://doi.org/10.1371/journal.pone.0027241}]{10.1371/journal.pone.0027241}.}}. Récemment, deux études utilisant un modèle statistique de la musique ont montré un effet de Wundt entre la modélisation computationnelle des attentes musicales et les auto-évaluations du plaisir musical par les participants\teiNote[n={49},place={foot},xmlid={ftn105-5}]{\teiP[xmlid={p-566}]{Benjamin P. Gold, Marcus T. Pearce, Ernest Mas-Herrero \teiHi[rend={italic}]{et al}., « Predictability and Uncertainty in the Pleasure of Music: A Reward for Learning? », \teiHi[rend={italic}]{Journal of Neuroscience}, vol. 39, no 47, 2019, p. 9397-9409 ; Vincent K. M. Cheung, Peter M. C. Harrison, Lars Meyer \teiHi[rend={italic}]{et al}., « Uncertainty and Surprise Jointly Predict Musical Pleasure and Amygdala, Hippocampus, and Auditory Cortex Activity », \teiHi[rend={italic}]{Current Biology}, vol. 29, n\teiHi[rend={sup}]{o} 23, 2019, p. 4084-4092. }}. Une modulation de l’activité dans la région liée à la récompense du striatum ventral induite par l’attente\teiNote[n={50},place={foot},xmlid={ftn106-5}]{\teiP[xmlid={p-567}]{Ofir Shany, Neomi Singer, Benjamin Paul Gold \teiHi[rend={italic}]{et al}., « Surprise-Related Activation in the Nucleus Accumbens Interacts with Music-Induced Pleasantness », \teiHi[rend={italic}]{Social Cognitive and Affective Neuroscience}, vol. 14, n\teiHi[rend={sup}]{o} 4, 2019p. 459-470.}} ou l’incertitude\teiNote[n={51},place={foot},xmlid={ftn107-5}]{\teiP[xmlid={p-568}]{Vincent K. M. Cheung, Peter M. C. Harrison, Lars Meyer \teiHi[rend={italic}]{et al}., art. cité.}} des événements musicaux a également été découverte. Cette littérature semble montrer un lien clair entre l’environnement culturel et les préférences musicales.}
              \begin{teiFigure}[xmlid={figure01-5}]

                \teiHead[xmlid={graphique-22-selon-leffet-wundt-les-evenements-musicaux-generant-un-niveau-intermediaire-de-predictibilite-sont-ceux-qui-engendrent-le-maximum-de-plaisir-001}]{Graphique 22. Selon l’effet Wundt, les événements musicaux générant un niveau intermédiaire de prédictibilité sont ceux qui engendrent le maximum de plaisir.}
                \teiP[rend={caption},xmlid={p-569}]{Source : l’auteur.}
              
\end{teiFigure}
              \teiP[xmlid={p14-12}]{D’après la littérature, il n’est pas évident que les facteurs socioculturels puissent expliquer entièrement la variance des préférences musicales, ce qui signifie qu’il existe probablement des différences individuelles causées par les expériences de vie uniques et la physiologie des individus. La personnalité étant le premier candidat pour de telles différences, Rentfrow et Gosling\teiNote[n={52},place={foot},xmlid={ftn108-5}]{\teiP[xmlid={p-570}]{Peter J. Rentfrow et Samuel D. Gosling, « The Do Re Mi’s of Everyday Life: The Structure and Personality Correlates of Music Preferences », \teiHi[rend={italic}]{Journal of Personality and Social Psychology}, vol. 84, n\teiHi[rend={sup}]{o} 6, 2003, p. 1236-1256.}} ont utilisé le STOMP (\teiHi[rend={italic}]{short test of musical preferences}) et ont observé que les individus ayant une préférence pour les styles réfléchis et complexes (tels que le classique et le jazz) ont tendance à avoir des niveaux plus élevés d’ouverture à l’expérience (mesurés à l’aide du test de personnalité Big Five), à se percevoir comme plus intelligents et moins athlétiques. En revanche, ceux qui préfèrent les styles rythmiques et énergiques (comme la danse électronique et le hip-hop) sont plus susceptibles de présenter des niveaux élevés d’extraversion et d’ouverture à l’expérience, et de se percevoir comme plus athlétiques. Cependant, une méta-analyse (combinant 263 196 participants) sur la relation entre les traits de personnalité et les préférences musicales\teiNote[n={53},place={foot},xmlid={ftn109-5}]{\teiP[xmlid={p-571}]{Thomas et Claudia Mehlhorn, « Can Personality Traits Predict Musical Schaefer Style Preferences? A Meta-Analysis », \teiHi[rend={italic}]{Personality and Individual Differences}, vol. 116, 2017, p. 265-273.}} constate que, bien qu’il y ait eu de petites corrélations pour des styles et des traits particuliers, notamment entre l’ouverture à l’expérience et le soft rock R’n’B, le classique et l’avant-garde, les traits de personnalité ont un pouvoir prédictif limité pour expliquer les différences interindividuelles en matière de préférences musicales. La physiologie pourrait également être un facteur d’aplanissement, mais les quelques études qui ont étudié la relation entre la fréquence cardiaque et la préférence pour certains \teiHi[rend={italic}]{tempi} ont donné des résultats contradictoires. Deux études portant sur les composantes génétiques des préférences musicales (qui se sont révélées très faibles pour les préférences faciales\teiNote[n={54},place={foot},xmlid={ftn110-5}]{\teiP[xmlid={p-572}]{Laura Germine, Richard Russell, P. Matthew Bronstad \teiHi[rend={italic}]{et al.}, « Wilmer. Individual Aesthetic Preferences for Faces Are Shaped Mostly by Environments, Not Genes », \teiHi[rend={italic}]{Current Biology}, vol. 25, n\teiHi[rend={sup}]{o} 20, 2015, p. 2684–2689.}}, mais non absentes) ont généré des résultats contradictoires sur les observations non directes des préférences musicales (par le biais des préférences pour différents domaines comme la mode, la cuisine et la musique\teiNote[n={55},place={foot},xmlid={ftn111-5}]{\teiP[xmlid={p-573}]{Judith Faust, « A Twin Study of Personal Preferences », \teiHi[rend={italic}]{Journal of Biosocial Science}, vol. 6, n\teiHi[rend={sup}]{o} 1, 1974, p. 75-91.}} et des préférences de choix pour les instruments\teiNote[n={56},place={foot},xmlid={ftn112-5}]{\teiP[xmlid={p-574}]{Miriam A. Mosing et Frederik Ullén, « Genetic Influences on Musical Specialization: A Twin Study on Choice of Instrument and Music Genre », \teiHi[rend={italic}]{Annals of the New York Academy of Sciences}, 2018, DOI : 10.1111/nyas.13626.}}). Ces résultats semblent converger vers une explication principalement sociale et acquise des préférences musicales. Néanmoins, le substrat neurocognitif pour l’enculturation musicale au travers du traitement prédictif semble être, quant à lui, hérité et partagé entre de nombreuses espèces animales.}
            
\end{teiDiv}
            \begin{teiDiv}[type={section1},xmlid={div05-5}]

              \teiHead[xmlid={beaute-musicale-001}]{Beauté musicale ?}
              \teiP[xmlid={p15-12}]{Nous avons jusqu’ici résumé la littérature scientifique qui suggère que la préférence musicale est liée à la manière dont notre cerveau prédit la musique et que ces prédictions sont apprises tout au long de notre vie grâce à une exposition passive. Cela explique bien pourquoi chaque auditeur éprouve du plaisir en écoutant des musiques différentes. Cependant, ces connaissances portent sur les émotions, le plaisir et les préférences musicales. Afin de les appliquer à la notion de beauté, nous avons mené une étude auprès de 34 participants à Paris. Nous avons demandé à tous les participants (dont nous avons vérifié qu’ils n’étaient pas familiers avec la musique chinoise) de venir au laboratoire et d’écouter une heure de musique chinoise tout en étant connectés à un système électroencéphalographique. Ce système non invasif mesure l’activité électrique à la surface du cerveau (cortex) et peut être utilisé pour mesurer les prédictions musicales dans le cerveau. Pour chaque chanson, il a été demandé aux participants d’indiquer l’intensité du sentiment de beauté qu’ils ont ressenti. Nous avons ensuite divisé les participants en deux groupes : le premier groupe a été invité à écouter de la musique chinoise (des chansons différentes, mais du même genre) une heure par jour pendant une semaine et l’autre groupe a été invité à écouter des chorals de Bach joués sur le même instrument que la musique chinoise (Guzheng). Les participants ne savaient ni à quel groupe ils étaient assignés ni qu’il y avait différents groupes. Ils sont revenus après une semaine pour suivre la même procédure.}
              \teiP[xmlid={p16-12}]{Le résultat le plus frappant est que nous avons pu observer un changement de direction positif dans l’évaluation de la beauté entre les deux enregistrements pour le groupe qui a écouté de la musique chinoise, mais pas pour celui qui a écouté Bach (légèrement négatif) – voir graphique 23a. Cela signifie donc que le sentiment de beauté est conditionné par la familiarité. Nous n’avons pas pu exposer les participants jusqu’à ce qu’ils soient trop familiers avec la musique chinoise, nous ne pouvons donc pas savoir si les évaluations de la beauté diminuent lorsque les chansons sont trop familières, comme c’est le cas pour les préférences dans l’effet Wundt. L’autre découverte est que cet apprentissage est réversible sur un temps long. Les participants sont revenus une nouvelle fois au laboratoire après deux mois de repos (c’est-à-dire sans avoir écouté de musique différente de celle qu’ils écoutent habituellement) pour renouveler la procédure. Nous avons trouvé l’effet inverse, les notations s’étaient ainsi rapprochées de leur point de départ (graphique 23b). Néanmoins, même après deux mois d’oubli, les résultats sont encore éloignés du point de départ. Par conséquent, la découverte d’une nouvelle musique affecte nos préférences musicales pour longtemps après cette rencontre. Nous avons également demandé aux participants de noter le niveau de plaisir qu’ils ont ressenti ; cela permet ainsi de revenir à la question de la distinction plaisir/beauté, déjà très discutée dans la littérature\teiNote[n={57},place={foot},xmlid={ftn113-5}]{\teiP[xmlid={p-575}]{Aenne Brielmann et Denis Pelli, « Intense Beauty Requires Intense Pleasure », \teiHi[rend={italic}]{Frontiers in Psychology}, vol. 10, 2019, DOI : \teiRef[target={https://doi.org/10.3389/fpsyg.2019.02420}]{10.3389/fpsyg.2019.02420}.}}. Nous avons comparé les mesures de beauté et de plaisir dans notre expérience et avons trouvé une très grande corrélation (r = 87, voir graphique 23c). Ceci signifie que ces deux sensations semblent liées pour les participants. Un nouvel axe de recherche devrait être ouvert afin de comprendre la relation entre beauté et plaisir. Sont-ils différents ? Quels en sont les points de divergences ? Le plaisir est-il nécessaire pour supporter la beauté ? La philosophie a déjà exploré ces questions, néanmoins, il serait intéressant d’en faire une validation empirique à l’aide de données cognitives.}
              \teiP[xmlid={p17-12}]{Notre étude suggère une nouvelle perspective sur la beauté musicale : loin d’être une essence universelle, elle apparaît comme un phénomène culturel façonné par la familiarité. La musique que nous rencontrons tout au long de notre vie sculpte nos attentes et nos surprises, modelant ainsi notre expérience esthétique. Ce cadre propose une rupture avec les conceptions universalistes de la beauté et invite à repenser notre perception du monde sonore. La beauté musicale, tout en reposant sur des mécanismes cognitifs universels, se révèle profondément ancrée dans des normes culturelles. En somme, la beauté n’est pas seulement ressentie ; elle s’apprend, s’oublie, et surtout, se construit au fil de nos expériences.}
              \begin{teiFigure}[xmlid={figure02-4}]

                \teiHead[xmlid={graphiques-23a-b-et-c-changement-dans-levaluation-de-la-beaute-ressentie-apres-avant-une-augmentation-moyenne-superieure-a-0-signifie-que-les-participants-apprecient-davantage-les-morceaux-ces-graphiques-montrent-le-changement-induit-par-la-phase-dexposition-23a-et-la-phase-de-repos-23b-et-illustrent-respectivement-lapprentissage-dune-nouvelle-grammaire-musicale-et-son-declin-apres-deux-mois-dabsence-dexposition-nous-pouvons-voir-que-leffet-de-lexposition-dans-le-groupe-test-a-induit-une-augmentation-des-notes-de-plaisir-qui-ont-ete-partiellement-effacees-apres-la-phase-de-repos-ce-qui-est-une-preuve-evidente-de-lapprentissage-de-la-nouvelle-grammaire-musicale-et-de-son-declin-001}]{Graphiques 23a, b et c. Changement dans l’évaluation de la beauté ressentie (après-avant). Une augmentation (moyenne supérieure à 0) signifie que les participants apprécient davantage les morceaux. Ces graphiques montrent le changement induit par la phase d’exposition (23a) et la phase de repos (23b) et illustrent, respectivement, l’apprentissage d’une nouvelle grammaire musicale et son déclin après deux mois d’absence d’exposition. Nous pouvons voir que l’effet de l’exposition dans le groupe test a induit une augmentation des notes de plaisir qui ont été partiellement effacées après la phase de repos, ce qui est une preuve évidente de l’apprentissage de la nouvelle grammaire musicale et de son déclin.}
                \teiP[rend={caption},xmlid={p-576}]{Source : l’auteur.}
              
\end{teiFigure}
            
\end{teiDiv}
          
\end{teiElement}
        
\end{teiElement}
      
\end{teiElement}
      \begin{teiElement}[name={group},type={section1},xmlid={section1-004}]

        \teiHead[xmlid={l-enjeu-vital-de-la-beaute-001}]{L'enjeu vital de la beauté}
        \begin{teiElement}[name={group},type={chapter},data-page-title={Approche sociologique de la beauté en situation de vulnérabilité},data-page-authors={Gisèle Dambuyant@@Université Sorbonne Paris Nord-CRNS-Inserm-EHESS, IRIS – UMR 8156-U 997, France},data-include-href={Ch14\_approche\_sociologique\_Dambuyant.xml},xmlid={chapter-013}]

          \begin{teiElement}[name={front}]

            \begin{teiDiv}[type={titlePage},xmlid={titlepage-015}]

              \teiP[rend={title-main},xmlid={p-577}]{Approche sociologique de la beauté en situation de vulnérabilité}
              \teiP[rend={author-aut},xmlid={p-578}]{Gisèle \teiHi[rend={small-caps}]{Dambuyant}}
              \teiP[rend={authority\_affiliation},xmlid={p-579}]{Université Sorbonne Paris Nord-CRNS-Inserm-EHESS, IRIS – UMR 8156-U 997, France}
            
\end{teiDiv}
          
\end{teiElement}
          \begin{teiElement}[name={body},xmlid={body-015}]

            \teiP[xmlid={p1-15}]{Si la beauté apparaît relative et subjective, elle est aussi une force, vitalité et pouvoir de transformation pour tout individu. Ressource centrale de l’identité, quelle que soit la situation sociale, on peut se questionner quant à sa forme, voire sa permanence en situation de vulnérabilité. L’approche sociologique différencie cette dernière en zones de fragilité, de précarité et d’exclusion. Reste-t-il encore une place possible pour la beauté en situation de vulnérabilité, et particulièrement, dans des conditions de survie, tant physique que sociale ? }
            \teiP[xmlid={p2-15}]{Nous répondrons à cette problématique en trois temps. Nous étudierons d’abord la beauté contemporaine et ses perspectives, puis la vulnérabilité au \teiHi[rend={small-caps}]{xxi}\teiHi[rend={sup}]{e} siècle en France et nous analyserons enfin la beauté comme ressource en situation de vulnérabilité, car celle-ci, paradoxalement, s’avère encore plus fondamentale en condition de dénuement extrême.}
            \begin{teiDiv}[type={section1},xmlid={div01-11}]

              \teiHead[xmlid={les-evolutions-des-criteres-de-beaute-entre-representations-individuelles-et-valeurs-societales-001}]{Les évolutions des critères de beauté, entre représentations individuelles et valeurs sociétales}
              \teiP[xmlid={p3-15}]{Selon sa définition première, la beauté renvoie à un idéal esthétique, ou à des qualités intellectuelles ou morales, forçant l’admiration. Pour Georges Vigarello, la beauté incarne à la fois l’imaginaire, les valeurs d’une époque et d’une société\teiNote[xmlid={ftn1-14},n={1},place={foot}]{\teiP[xmlid={p-580}]{Georges Vigarello, \teiHi[rend={italic}]{Histoire de la beauté}. \teiHi[rend={italic}]{Le corps et l’art d’embellir de la Renaissance à nos jours}, Paris, Éditions du Seuil, 2004. }}. Par ailleurs, elle distingue les individus dans leurs genres, leurs âges, leurs appartenances sociales et culturelles, par des représentations individuelles et collectives. Dès lors, la beauté doit être considérée pour tous les milieux, pour toutes les cultures et dans toutes ses dimensions : le soin, le traitement fait au corps, l’estime de soi ou l’apparence, configurées par les valeurs sociétales. Comme pour tout concept complexe, la beauté doit être examinée de façon pluridisciplinaire\teiNote[xmlid={ftn24-3},n={2},place={foot}]{\teiP[xmlid={p-581}]{Martine Janner-Raimondi (dir.), \teiHi[rend={italic}]{Oser l’interdisciplinarité. Rencontres croisées autour du vivant}, Paris, Téraèdre, 2022.}}. Si l’on se réfère au substantif latin \teiHi[rend={italic}]{forma}, ou à l’adjectif \teiHi[rend={italic}]{formosus}, on constate que ces notions sont fondamentalement liées à la matérialité et à ses effets. Initialement étudiée par la discipline philosophique, la « beauté » concerne tout autant l’aspect esthétique – elle désigne la « qualité de quelqu’un, de quelque chose qui est beau, conforme à un idéal esthétique » –, que moral : c’est le « caractère de ce qui est digne d’admiration, par ses qualités intellectuelles… Séduisant\teiNote[xmlid={ftn25-4},n={3},place={foot}]{\teiP[xmlid={p-582}]{Définition du dictionnaire Larousse : \teiRef[target={https://www.larousse.fr/dictionnaires/francais/beaut\%C3\%A9/8533}]{https://www.larousse.fr/dictionnaires/francais/beaut\%C3\%A9/8533} (consulté le 4 septembre 2025).}} ». C’est une notion abstraite liée à de nombreux aspects de l’existence humaine qui intéresse toutes les sciences. Historiquement, ce concept est associé à la subjectivité, mais des recherches cognitives récentes en expliquent l’objectivité, et des travaux de neurosciences démontrent que les valeurs esthétiques sont objectivement adaptées et évolutives\teiNote[xmlid={ftn26-3},n={4},place={foot}]{\teiP[xmlid={p-583}]{Norberto M. Grzywacz, professeur de psychologie, de pharmacologie moléculaire et de neurosciences de l’université de Chicago, travaille avec son équipe pluridisciplinaire sur l’expérience de la beauté dans le cerveau. Voir Hassan Aleem and Noberto M. Grzywacz, « The Temporal Instability of Aesthetic Preferences », \teiHi[rend={italic}]{Psychology of Aesthetics Creativity and Arts}, vol. 19, n\teiHi[rend={sup}]{o} 3, p. 505-521.}}.}
              \teiP[xmlid={p4-15}]{Selon les époques et les cultures, les critères de beauté ont évolué, tant pour les hommes (ainsi les cheveux longs des Francs, symbole de puissance, ou le surpoids du bourgeois du \teiHi[rend={small-caps}]{xix}\teiHi[rend={sup}]{e} siècle, synonyme de réussite sociale) que pour les femmes (les larges hanches sont valorisées pour certaines, le corps mince pour d’autres). Si la beauté démontre ses transformations et son hétérogénéité, elle semble bien concerner tout être vivant. Peut-on alors évoquer une beauté vitale ? Pour Frédéric Worms, c’est un paradoxe qu’il faut dépasser pour conceptualiser la beauté comme vitale en tant que telle, car il n’y a pas de vie qui ne soit pas relationnelle et être vivant, c’est s’insérer dans un milieu. En se distanciant par rapport au monde, matériellement, pour en faire un objet de contemplation et de façon relationnelle, par leurs capacités d’attachement, les êtres humains portent en eux la beauté de la vie\teiNote[xmlid={ftn27-3},n={5},place={foot}]{\teiP[xmlid={p-584}]{Voir l’article de Frédéric Worms dans le présent volume. }}. Cette beauté du vivant s’accommode aux normes valorisées dans une société à une période donnée. Elle est donc historiquement et culturellement située. }
              \teiP[xmlid={p5-15}]{En France, en ce début du \teiHi[rend={small-caps}]{xxi}\teiHi[rend={sup}]{e} siècle, la montée de l’individualisme – cette tendance à s’affranchir de toute obligation collective et à ne vivre que pour soi\teiNote[xmlid={ftn28-3},n={6},place={foot}]{\teiP[xmlid={p-585}]{Voir François de Singly et Danilo Martucelli, \teiHi[rend={italic}]{Les sociologies de l’individu}, Paris, Armand Colin, 2009.}} –, autant que le culte de la performance – la recherche du meilleur résultat et de l’accomplissement personnel –, influencent les caractéristiques de la beauté contemporaine\teiNote[xmlid={ftn29-3},n={7},place={foot}]{\teiP[xmlid={p-586}]{Voir Alain Ehrenberg, \teiHi[rend={italic}]{Le culte de la performance}, Paris, Calmann-Lévy, 1991, et Didier Fassin, \teiHi[rend={italic}]{Des maux indicibles. Sociologie des lieux d’écoute}, Paris, La Découverte, 2004.}}. Tout individu doit prendre soin de son corps, anticiper les effets du temps, tout en essayant de le rendre le plus efficient possible. Si les libertés corporelles individuelles existent, elles sont aussi limitées par des normes, qui vont jusqu’aux « tyrannies de l’apparence » et des valeurs sociales, ou encore au « vieillissement intolérable\teiNote[xmlid={ftn30-2},n={8},place={foot}]{\teiP[xmlid={p-587}]{Voir David Le Breton, \teiHi[rend={italic}]{Conduites à risque}, Paris, PUF, 2013.}} ». }
              \teiP[xmlid={p6-15}]{Ainsi, dans la société contemporaine, et dans l’imaginaire collectif, la beauté est associée à la jeunesse, la séduction, l’autonomie et l’intégration. Quelle place la beauté peut-elle donc occuper dans les perspectives de l’évolution de la société plus juste et plus durable définie par l’Organisation des Nations unies (ONU) à l’horizon 2030 ?}
              \teiP[xmlid={p7-15}]{En septembre 2015, 193 États membres de l’ONU ont adopté 17 objectifs de développement durable (ODD) à atteindre par le travail en partenariat, socialement équitables, sûrs d’un point de vue environnemental, économiquement prospères et inclusifs. Parmi ces 17 objectifs, quatre sont environnementaux : la lutte contre le changement climatique (objectif 13), l’accès à l’eau salubre et l’assainissement (objectif 6), la protection de la faune et de la flore aquatiques (objectif 14), la protection de la faune et de la flore terrestres (objectif 15). Quatre objectifs sont économiques : le travail décent et la croissance économique (objectif 8), la promotion de l’innovation et des infrastructures durables (objectif 9), la réduction des inégalités (objectif 10), la production et la consommation responsable (objectif 12). Huit objectifs sont sociaux et concernent particulièrement les missions du travail et de l’intervention sociale\teiNote[xmlid={ftn31-2},n={9},place={foot}]{\teiP[xmlid={p-588}]{Voir Gisèle Dambuyant, « La place du travail social et de l’intervention sociale dans les objectifs de développement durable (ODD) », \teiHi[rend={italic}]{Colloque international du Réseau européen des formations universitaires du travail social }(REFUTS), université de Saragosse, 27 au 29 juin 2022. }}. Ils visent prioritairement un double but : l’éradication de la pauvreté (objectif 1) et la contribution à la justice sociale, en assurant la paix et l’efficacité des institutions (objectif 16), en luttant contre les inégalités de toutes natures. Inégalités alimentaires (objectif 2), éducatives (objectif 4), sexuelles (objectif 5), énergétiques (objectif 7), inégalités de logement (objectif 11) et de santé (objectif 3) auxquelles on pourrait ajouter une inégalité plus récente : la lutte contre la précarité des nouvelles technologies d’information et de communication (NTIC\teiNote[xmlid={ftn32-2},n={10},place={foot}]{\teiP[xmlid={p-589}]{Voir le site de l’ONU : \teiRef[target={https://www.un.org/sustainabledevelopment/fr/objectifs-de-developpement-durable/}]{https://www.un.org/sustainabledevelopment/fr/objectifs-de-developpement-durable/} (consulté le 4 septembre 2025).}}). }
              \teiP[xmlid={p8-15}]{Quelle place, donc, pour la beauté dans ces ODD ? }
              \teiP[xmlid={p9-14}]{La beauté se définit par son caractère digne d’admiration, évoquant alors le sentiment de joie et de plénitude, et renvoyant au sentiment de bonheur et de bien-être. Ainsi, dans la perspective de valeurs sociétales devenant centrales, la beauté prend sa place dans l’objectif 3 : la santé et le bien-être des populations et des travailleurs. La santé est elle-même définie par l’Organisation mondiale de la santé depuis 1946 comme « un état de complet bien-être physique, mental et social » et « ne consiste pas seulement en une absence de maladie ou d’infirmité\teiNote[xmlid={ftn33-2},n={11},place={foot}]{\teiP[xmlid={p-590}]{Voir le site de l’Organisation mondiale de la santé : \teiRef[target={https://www.who.int/fr/}]{https://www.who.int/fr/}, et en particulier sa constitution : \teiRef[target={https://www.who.int/fr/about/governance/constitution}]{https://www.who.int/fr/about/governance/constitution} (consulté le 4 septembre 2025).}} ». La beauté devient ainsi une composante essentielle du bien-être de tout individu, aussi désocialisé soit-il. Comment le bien-être peut-il se décliner en situation de vulnérabilité ? Ceci nous conduit à interroger la notion de beauté et ses réalités contemporaines. }
            
\end{teiDiv}
            \begin{teiDiv}[type={section1},xmlid={div02-11}]

              \teiHead[xmlid={definition-de-la-vulnerabilite-et-place-de-la-beaute-dans-une-situation-de-vulnerabilite-001}]{Définition de la vulnérabilité et place de la beauté dans une situation de vulnérabilité}
              \teiP[xmlid={p10-14}]{Robert Castel associe la vulnérabilité à des environnements de fragilités professionnelles et relationnelles : « Être vulnérable, c’est être exposé à des menaces externes, plus ou moins prévisibles, qui mettent à l’épreuve un certain nombre de ressources détenues par des individus, des groupes et des communautés sur des territoires\teiNote[xmlid={ftn34-3},n={12},place={foot}]{\teiP[xmlid={p-591}]{Robert Castel et Claude Martin (dir.), \teiHi[rend={italic}]{Changements et pensées du changement. Échanges avec Robert Castel}, Paris, La Découverte, 2013, p. 295. De Robert Castel, voir aussi \teiHi[rend={italic}]{Les Métamorphoses de la question sociale. Une chronique du salariat}, Paris, Fayard, 1995.}}. » Ces menaces plus ou moins prévisibles sont associées à des risques qui peuvent amoindrir, voire annihiler, les ressources économiques, sociales et culturelles plus ou moins détenues par chaque personne. }
              \teiP[xmlid={p11-14}]{Ces dernières années, les publics vulnérables sont à la fois plus nombreux et plus hétérogènes. Globalement, les dernières statistiques démontrent une intensification de la pauvreté : 5,3 millions de personnes vivent avec moins de 885 € par mois, le taux de pauvreté a augmenté de 0,5 point en 10 ans (ce qui représente près de 500 000 personnes supplémentaires), un quart des adultes handicapés de moins de 60 ans vivent avec moins de 1 000 € par mois et le taux de pauvreté des 18-29 ans a, quant à lui, progressé de 50 \% en 16 ans\teiNote[xmlid={ftn35-4},n={13},place={foot}]{\teiP[xmlid={p-592}]{\teiHi[rend={italic}]{Rapport sur les inégalités en France}, édition 2021 : \teiRef[target={https://www.inegalites.fr/Le-Rapport-sur-les-inegalites-en-France-edition-2021-vient-de-paraitre}]{https://www.inegalites.fr/Le-Rapport-sur-les-inegalites-en-France-edition-2021-vient-de-paraitre} (consulté le 4 septembre 2025).}}. Ces tendances générales se sont amplifiées avec la crise sanitaire de la Covid-19 et l’on a pu distinguer cinq nouvelles catégories de publics vulnérables. Des personnes en emploi stable, qui ont connu une baisse des ressources ayant des impacts sur leur équilibre budgétaire, des personnes en emploi précaire ou sans emploi, notamment les étudiants en forte difficulté financière ; des personnes en difficulté pour l’accès aux droits, en particulier à cause d’une numérisation croissante des services publics ; des personnes en souffrance, des personnes en forte difficulté psychologique et des aidants eux-mêmes en difficulté\teiNote[xmlid={ftn36-4},n={14},place={foot}]{\teiP[xmlid={p-593}]{Christine Olm, « Étude de faisabilité du baromètre de suivi qualitatif de la pauvreté », Cabinet VIZGET, présentée en séance plénière sur les apports du Conseil national des politiques de Lutte contre la pauvreté et l’exclusion sociale (CNLE), 24 février 2022.}}. Ces évolutions de la vulnérabilité obligent à repenser une nouvelle approche de ce concept en intégrant les diverses formes de difficultés et de souffrances ressenties par les personnes vulnérables. }
              \teiP[xmlid={p12-14}]{Ainsi, et au-delà des contextes sociaux-professionnels, c’est le corps, dans toutes ses dimensions – physique, psychologique et symbolique – qu’il faut prendre en compte. Pour les plus démunis, les sans-abri par exemple, le corps peut devenir la dernière et ultime ressource\teiNote[xmlid={ftn37-4},n={15},place={foot}]{\teiP[xmlid={p-594}]{Gisèle Dambuyant, \teiHi[rend={italic}]{Quand on n’a plus que son corps}, Paris, Armand Colin, 2006.}}. Le corps et son état peuvent alors servir d’indicateurs pour proposer une classification et une nouvelle définition : }
              \begin{teiCit}[xmlid={cit01-10}]

                \begin{teiQuote}
La vulnérabilité est une caractéristique de la condition humaine. Au cours de chaque trajectoire de vie, l’individu est exposé à des risques, corporels et environnementaux, qui le placent, en fonction de ses ressources, en situation de fragilité, de précarité ou d’exclusion. Ainsi majoritairement, la fragilité est une vulnérabilité physique, la précarité, une vulnérabilité sociale et l’exclusion une vulnérabilité globale\teiNote[xmlid={ftn38-4},n={16},place={foot}]{\teiP[xmlid={p-595}]{Gisèle Dambuyant, \teiHi[rend={italic}]{Beauté et vulnérabilité}, dans Clélia Zernik et Justin Jaricot (dir.), \teiHi[rend={italic}]{Abécédaire de la beauté}, Montreuil, Éditions B42, 2022.}}. 
\end{teiQuote}
              
\end{teiCit}
              \teiP[xmlid={p13-13}]{Qu’en est-il de ces zones et de leur composante majeure de vulnérabilité ? }
              \teiP[xmlid={p14-13}]{La fragilité renvoie à « ce qui n’est pas solide, pas résistant ». Appliquée à une personne, elle désigne « sa faible constitution, son peu de résistance physique ou psychologique ». À certains âges extrêmes de la vie (les enfants et les personnes âgées) ou dans certaines situations, comme certaines maladies ou certains handicaps, le corps est vulnérable, car il dépend d’autrui pour vivre. Il existe donc un lien fort entre fragilité et dépendance. La vulnérabilité se centralise sur la fragilité physique. }
              \teiP[xmlid={p15-13}]{La précarité renvoie à l’incertitude « sans garantie de durée, de stabilité, qui peut toujours être remis en cause ». Appliquée à une personne, elle évoque l’insécurité liée aux contextes instables qui la menacent en permanence, renvoyant aux risques tant physiques que sociaux. Dans cet environnement, la personne ressent un sentiment de soi associé à de la honte, à de la culpabilité et à de la perte d’estime de soi. La vulnérabilité s’exprime prioritairement par les contextes sociaux non sécurisés. }
              \teiP[xmlid={p16-13}]{L’exclusion renvoie à « ce qui a été rejeté, chassé du groupe, de l’organisation, de l’institution dont il faisait partie ». La personne exclue se situe « en dehors » de la société, elle est mise à l’écart du fait de sa déviance et se retrouve à la rue ou, du fait de sa délinquance, elle subit l’emprisonnement. Elle s’inscrit majoritairement dans des trajectoires de ruptures et de désocialisation, et elle cumule les manques : sans abri, sans argent, sans famille, sans amis, sans travail… Le corps constitue alors la dernière ressource, la dernière possibilité d’actions fonctionnant en surexploitation et surexposition permanentes, et il se dégrade inexorablement. La personne est située dans un processus de déconstruction identitaire où son intégrité totale est atteinte : physique, psychique et sexuelle. La vulnérabilité est globale. }
              \teiP[xmlid={p17-13}]{Quelle place peut prendre la beauté dans ces situations ? Pour répondre à cette question, nous exemplifierons les analyses par des recherches de terrain, menées notamment auprès de socio-esthéticiens ou socio-esthéticiennes intervenant sur des corps vulnérables\teiNote[xmlid={ftn39-4},n={17},place={foot}]{\teiP[xmlid={p-596}]{Gisèle Dambuyant, \teiHi[rend={italic}]{La socio-esthétique. Prendre soin, soulager et embellir le corps vulnérable}, Toulouse, Érès, 2023.}}. }
            
\end{teiDiv}
            \begin{teiDiv}[type={section1},xmlid={div03-10}]

              \teiHead[xmlid={la-beaute-comme-ressource-en-situation-de-vulnerabilite-001}]{La beauté comme ressource en situation de vulnérabilité}
              \teiP[xmlid={p18-12}]{La beauté impacte l’individu dans toutes ses dimensions, car elle façonne le corps en permettant de construire les représentations de son image et de son estime, par rapport à soi et à l’autre. On se réfère alors à l’identité individuelle qui prend forme, mais se différencie de l’identité collective\teiNote[xmlid={ftn40-4},n={18},place={foot}]{\teiP[xmlid={p-597}]{Voir Nathalie Heinich, \teiHi[rend={italic}]{Ce que n’est pas l’identité}, Paris, Gallimard, 2018. }}. Dès lors, peut-on affirmer que, quelles que soient la situation sociale et la société, la beauté participe à l’identité, dans sa construction et sa permanence ? La question qui se pose alors est celle de la forme que la beauté peut prendre en tant que ressource créée, mobilisée et mobilisable en situation de vulnérabilité ? Peut-on décliner des rapports à la beauté spécifiques pendant l’enfance ou la vieillesse ? Ceux-ci vont-ils se transformer – et de quelles façons ? – pour des personnes en situation de précarité ou d’exclusion ? Reste-t-il une place possible pour la beauté lorsque les besoins vitaux se gèrent de façon incertaine et au quotidien ? }
            
\end{teiDiv}
            \begin{teiDiv}[type={section1},xmlid={div04-6}]

              \teiHead[xmlid={la-beaute-comme-ressource-en-situation-de-fragilite-001}]{La beauté comme ressource en situation de fragilité }
              \teiP[xmlid={p19-12}]{Dans cette zone où la vulnérabilité est prioritairement physique, le rapport à la beauté est favorisé dans des relations sociales de proximité, suscitées par la relation de dépendance à l’autre. Créer et maintenir du lien social, dans une interaction directe, permet de développer la place et le rôle de la beauté, en valorisant le corps et son image, dans la construction identitaire du jeune enfant comme dans la permanence identitaire des adultes. De même, il était essentiel de maintenir des interactions et des accompagnements auprès de travailleurs en situation de handicap, notamment pendant la crise de la COVID-19\teiNote[xmlid={ftn41-5},n={19},place={foot}]{\teiP[xmlid={p-598}]{Voir Gisèle Dambuyant, « Pour une revalorisation des professionnels accompagnant le “travail” des personnes vulnérables. État des lieux et perspectives », \teiHi[rend={italic}]{Empan}, vol. 127, n\teiHi[rend={sup}]{o} 3, 2022, p. 127-133.}}. Faire des soins de socio-esthétiques avec des personnes âgées dépendantes, notamment atteintes de la maladie d’Alzheimer, favorise les souvenirs : « Par le massage, les senteurs, le maquillage, on réactive la sensation de se sentir belle, d’être bien jusqu’au bout et d’exister dans le regard de l’autre\teiNote[xmlid={ftn42-5},n={20},place={foot}]{\teiP[xmlid={p-599}]{Voir Pascal Fellous (réal.), \teiHi[rend={italic}]{Une fontaine en Italie}, 2009, film sur la socio-esthétique.}}. »}
              \teiP[xmlid={p20-12}]{On peut parler ici de la beauté sociale et culturelle, comme support de construction et de continuité identitaires.}
            
\end{teiDiv}
            \begin{teiDiv}[type={section1},xmlid={div05-6}]

              \teiHead[xmlid={la-beaute-comme-ressource-en-situation-de-precarite-001}]{La beauté comme ressource en situation de précarité }
              \teiP[xmlid={p21-10}]{Dans cet environnement de vulnérabilité prioritairement sociale, la personne a un sentiment de soi associé à de la honte, à de la culpabilité et à de la perte d’estime de soi. Dans cette zone, les interactions sociales sont favorisées pour évoquer la place et le rôle de la beauté et ses répercussions sur les échanges avec autrui, notamment lorsque ceux-ci sont réalisés entre des personnes de cultures ethniques différentes. L’acquisition de codes sociaux adaptés au pays facilite les relations sociales et, au-delà, les chances d’intégration. Ainsi, les intervenants sociaux accompagnant des personnes dans leurs parcours migratoires doivent sans cesse prendre en compte des représentations variées, notamment pour les codes vestimentaires, qui valorisent par exemple plus ou moins les couleurs vives, et ajuster leurs prises en charge, notamment auprès des mineurs non accompagnés, pour assurer une protection rapprochée ou distanciée en fonction des problématiques et de l’autonomie des jeunes\teiNote[xmlid={ftn43-5},n={21},place={foot}]{\teiP[xmlid={p-600}]{Voir Gisèle Dambuyant, « \teiRef[target={https://www.cairn.info/revue-empan-2019-4-page-66.htm}]{Accueil et accompagnements socio-éducatifs des mineurs non accompagnés au foyer de l’enfance}. Bouleversements des prises en charge, adaptation des pratiques et complexité des mesures de protection », \teiHi[rend={italic}]{Empan}, vol. 116, n\teiHi[rend={sup}]{o} 4, 2019, p. 67-73.}}.}
              \teiP[xmlid={p22-10}]{On peut parler ici de la beauté sociale et culturelle, comme support de recomposition identitaire. }
            
\end{teiDiv}
            \begin{teiDiv}[type={section1},xmlid={div06-5}]

              \teiHead[xmlid={la-beaute-comme-ressource-en-situation-dexclusion-001}]{La beauté comme ressource en situation d’exclusion}
              \teiP[xmlid={p23-10}]{Être en prison du fait d’actes de délinquance, ou à la rue du fait de la déviance, révèle des réalités d’exclusion. Pourtant, les techniques d’approche et de soin sont à différencier tant les contextes sont opposés : de l’enfermement au leurre de la liberté.}
              \teiP[xmlid={p24-9}]{La personne incarcérée est sous contrôle total de ses actes et des rapports les plus basiques de l’espace et du temps. Malgré tout, les soins de socio-esthétique effectués dans certaines prisons démontrent qu’il existe encore une place pour la beauté, offrant un « droit à aller mieux » pour mieux supporter un quotidien et préparer l’avenir\teiNote[xmlid={ftn44-5},n={22},place={foot}]{\teiP[xmlid={p-601}]{Voir Gisèle Dambuyant, « Les interventions des socio-esthéticiennes auprès des femmes incarcérées », dans \teiHi[rend={italic}]{Guérir, se rétablir, aller mieux… en santé mentale et ailleurs,} colloque international co-organisé par l’AISLF, le CLERSE et le CERIES, université Lille 3, 7-9 janvier 2015 ; aussi présenté au congrès \teiHi[rend={italic}]{Codes}, Institut Pasteur, 3 avril 2015.}}. }
              \teiP[xmlid={p25-8}]{La personne sans-abri est en situation de survie, tant physique que sociale, une situation dans laquelle il s’agit de répondre au quotidien aux besoins les plus primaires : boire, manger, dormir. Cette réalité est faite de violences multifactorielles, environnementales, physiques, symboliques où le corps fonctionne en surexploitation et en surexposition permanentes. Majoritairement « asexués », pour être le plus dissimulés possible, ces corps poursuivent leur anéantissement dans ces conditions d’existence : problèmes d’hygiène et de soins, recours à des pratiques addictives, absence de prise en compte des problèmes de santé. Reste-t-il alors une place possible pour la beauté, car c’est bien dans ces circonstances qu’elle est le plus problématique ? En réalité, elle persiste, mais s’adapte à ces conditions d’existence. Choisir, comme le disent certains, de beaux vêtements, y compris dans un vestiaire d’un espace de solidarité et d’insertion, continuer pour les femmes à se regarder dans les rétroviseurs du camion de maraude, pendant que le travailleur social cherche une place d’hébergement, pour se recoiffer ou se remaquiller, afin, comme elles l’expliquent, de « continuer à être belle même dans la rue\teiNote[xmlid={ftn45-5},n={23},place={foot}]{\teiP[xmlid={p-602}]{Gisèle Dambuyant, « Réalités et projets de vie des sans-abri : lorsque le corps devient l’ultime ressource », \teiHi[rend={italic}]{Bulletin de l’Académie nationale de médecine}, n\teiHi[rend={sup}]{o} 197, 2014, p. 21-32.}} ».}
              \teiP[xmlid={p26-8}]{Tous ces exemples démontrent la continuité de la beauté en situation d’exclusion tout en en soulignant le paradoxe. Si, de façon positive, elle facilite les activités « économiques » comme la prostitution, la beauté représente aussi, de façon négative, un facteur aggravant les risques d’agressions, notamment sexuelles. Dès lors, il existe un contrôle ambivalent de la beauté chez l’acteur qui souhaite tout autant la préserver que la dissimuler, mais on peut attester de la beauté comme ressource centrale dans un univers où il en existe peu.}
            
\end{teiDiv}
            \begin{teiDiv}[type={section1},xmlid={div07-5}]

              \teiHead[xmlid={inversion-de-la-problematique-001}]{Inversion de la problématique }
              \teiP[xmlid={p27-8}]{Ce paradoxe permet d’inverser la problématique précédente : de place « possible », la beauté peut prendre une place « centrale », affirmant son pouvoir, dans toutes les situations de vulnérabilités, et ce, plus particulièrement, en condition d’exclusion. }
              \teiP[xmlid={p28-7}]{On peut alors proposer un parcours global de reconstruction identitaire à partir de la beauté en articulant trois types de pratiques professionnelles sociales : médicales, psychologiques et d’esthétiques. }
              \teiP[xmlid={p29-7}]{Ce parcours de soin par la beauté se décompose en trois étapes. D’abord, en allant au-devant de la personne à la rue, l’approche médico-sociale met le corps à l’abri en reconnaissant l’identité et la beauté humaine, digne de vivre dans un autre environnement. Ensuite, l’approche psychosociale permet une revalorisation de l’estime de soi, favorisant une réappropriation de la beauté sociale et la réalisation d’interactions normalisées. Enfin, l’approche socio-esthétique concerne la restauration de l’image de soi et valorise particulièrement l’identité culturelle et sexuelle. \teiHi[rend={italic}]{In fine,} et en s’inscrivant dans le temps, la beauté peut devenir une ressource centrale pour mettre peu à peu la personne à distance de l’exclusion et l’accompagner vers un univers plus normé, plus sûr et plus valorisant. }
              \teiP[xmlid={p30-7}]{L’approche sociologique permet de démontrer la permanence de la beauté en situation de fragilité, de précarité et d’exclusion. Il existe donc un lien fort entre la vulnérabilité et la beauté. Plus encore, la beauté est un lien entre l’individu vulnérable et le collectif sociétal. Composante centrale de l’identité dans son rapport à soi et aux autres, la place de la beauté perdure, mais diminue au fur et à mesure que la désocialisation s’intensifie. Un horizon d’actions se profile cependant pour la beauté, car elle peut devenir un objet central et légitime pour lutter contre la vulnérabilité. En fonction de la zone de vulnérabilité qu’occupe la personne, la beauté participe à la construction, la continuité, la recomposition voire la reconstruction identitaire et doit être prise en compte à ce titre. La prise en charge des plus vulnérables commence alors par un regard professionnel qui leur restitue leur part d’unicité et d’humanité. La beauté se fait alors vitale, pour tout individu, tout à la fois ligne de départ et ligne d’horizon.}
            
\end{teiDiv}
          
\end{teiElement}
        
\end{teiElement}
        \begin{teiElement}[name={group},type={chapter},data-page-title={La beauté comme baume, après la catastrophe},data-page-authors={Clelia Zernik@@École nationale supérieure des Beaux-arts de Paris, France@@Université PSL, chaire Beauté·s, France},data-include-href={Ch15\_beaute\_comme\_baume\_Zernik.xml},xmlid={chapter-014}]

          \begin{teiElement}[name={front}]

            \begin{teiDiv}[type={titlePage},xmlid={titlepage-016}]

              \teiP[rend={title-main},xmlid={p-603}]{La beauté comme baume, après la catastrophe}
              \teiP[rend={author-aut},xmlid={p-604}]{Clelia \teiHi[rend={small-caps}]{Zernik}}
              \teiP[rend={authority\_affiliation},xmlid={p-605}]{École nationale supérieure des Beaux-arts de Paris, France}
              \teiP[rend={authority\_affiliation},xmlid={p-606}]{Université PSL, chaire Beauté·s, France}
            
\end{teiDiv}
          
\end{teiElement}
          \begin{teiElement}[name={body},xmlid={body-016}]

            \teiP[xmlid={p1-16}]{Dans \teiHi[rend={italic}]{Retour à Tokyo}, Philippe Forest écrit :}
            \begin{teiCit}[xmlid={cit01-11}]

              \begin{teiQuote}
Je ne me sens pas vraiment le cœur d’écrire quoi que ce soit des terrifiantes et magnifiques images avec lesquelles Naoya Hatakeyama témoigne à sa manière d’artiste du tsunami qui a ravagé le village de son enfance et qui a touché sa famille. Il y a une beauté dans l’horreur. C’est ainsi. \lbrack{}…\rbrack{} Le pur ravage que le monde inflige aux hommes fait éclore une sorte de sidérante splendeur parmi un paysage de pleurs\teiNote[n={1},place={foot},xmlid={ftn1-15}]{\teiP[xmlid={p-607}]{Philippe Forest, \teiHi[rend={italic}]{Retour à Tokyo}, Paris, Cécile Defaut, 2014.}}. 
\end{teiQuote}
            
\end{teiCit}
            \teiP[xmlid={p2-16}]{« Terrifiantes et magnifiques », « beauté dans l’horreur ». Que nous disent ces oxymores ?}
            \begin{teiFigure}[xmlid={figure01-6}]

              \teiHead[xmlid={figure-24-couverture-de-louvrage-de-naoya-hatakeyama-kesengawa-lille-editions-light-motiv-2013-courtoisie-de-lauteur-001}]{Figure 24. Couverture de l’ouvrage de Naoya Hatakeyama, \teiHi[rend={italic}]{Kesengawa}, Lille, Éditions Light Motiv, 2013. Courtoisie de l’auteur.}
            
\end{teiFigure}
            \teiP[xmlid={p3-16}]{Dans un ouvrage photographique intitulé \teiHi[rend={italic}]{Kesengawa}\teiNote[n={2},place={foot},xmlid={ftn7-2}]{\teiP[xmlid={p-608}]{Naoya Hatakeyama, \teiHi[rend={italic}]{Kesengawa} \lbrack{}\teiHi[rend={italic}]{La Madeleine}\rbrack{}, Lille, Light Motiv, 2013.}}, évoqué ici par Philippe Forest, Naoya Hatakeyama a regroupé des photographies de sa ville natale, Rikuzentakata, paisible bourgade au bord de la rivière Kesen qui débouche sur la mer non loin, photographies prises avant 2011, puis après 2011 une fois la région dévastée par le tsunami survenu le 11 mars. L’ouvrage est composé de deux parties bien distinctes, séparées d’une page blanche. La première partie regroupe les photographies d’avant 2011, très sereines et solaires, et s’accompagne d’un texte écrit par le photographe pendant son voyage en moto depuis Tokyo le lendemain de la catastrophe pour rejoindre sa famille (sa mère est décédée dans un centre d’évacuation en hauteur qui n’a pas résisté à la force inattendue de la vague) ; la deuxième partie est une succession d’images sans texte, qui révèlent les ruines laissées par le tsunami. Si la page centrale laissée blanche interrompt la lecture du récit et double l’horreur des paysages dévastés d’un mutisme sidéré, l’ouvrage malgré sa division entre un avant heureux et un après horrifié conserve cependant une très grande unité. Cette unité tient à la composition rigoureuse de chaque photographie et pour le dire avec Philippe Forest, à leur très grande beauté. Les mots de Forest soulignent ce scandale de la beauté dans l’horreur, comme si le drame, la catastrophe humaine, la destruction se devaient d’associer la désolation à la laideur. Comment la beauté pourrait-elle coexister avec le drame humain ? Sans doute la beauté ne ressortit pas directement des destructions massives faites par le tsunami ; il existe cependant une beauté des ruines qui convoque une mélancolie spécifique dont l’âme romantique se repaît. Ce n’est pas cet état d’esprit qui préside aux photographies de Naoya Hatakeyama. En effet, ce ne sont pas les ruines en elles-mêmes qui sont belles, mais les photographies. Leur composition est extrêmement soignée et rigoureuse et c’est de cette composition que naît la beauté. La composition s’appuie sur des lignes de coupe très nettes qui segmentent l’image ou la doublent : la ligne d’horizon, la ligne de crête des montagnes, la ligne de signalisation des routes, la ligne du pont, la ligne d’un banc ou d’une canne à pêche. Plus exactement, les images sont souvent composées de trois parties de masses assez identiques : le ciel, le brouillard, la verdure ou l’eau, éléments naturels qui semblent ainsi à l’image s’équilibrer et confèrent ordre et apaisement au paysage. Cet équilibre est accentué par la présence fréquente d’un élément central : un personnage, une barque, un cygne, un pliage ou une bretelle d’autoroute. Quelque chose se noue au cœur, au centre, amarrant l’image et lui donnant une immobilité et comme une éternité. Les images de Naoya Hatakeyama sont donc extrêmement composées, et d’une beauté très épurée. Les surfaces photographiées sont souvent vides. De nombreuses surfaces planes et sans figures, que ce soit la surface de la mer, du ciel ou de la nuit, introduisent comme un aplat uniforme et abstrait dans l’image.}
            \begin{teiFigure}[xmlid={figure02-5}]

              \teiHead[xmlid={figure-25-27-juillet-2004-kesencho-imaizumi-vue-de-la-colline-courtoisie-de-lauteur-001}]{Figure 25. \teiHi[rend={italic}]{27 juillet 2004 : Kesenchô – Imaizumi, vue de la colline}. Courtoisie de l’auteur.}
            
\end{teiFigure}
            \begin{teiFigure}[xmlid={figure03-4}]

              \teiHead[xmlid={figure-26-04-avril-2011-kesencho-imaizumi-vue-de-la-colline-courtoisie-de-lauteur-001}]{Figure 26. \teiHi[rend={italic}]{04 avril 2011 : Kesenchô – Imaizumi, vue de la colline}. Courtoisie de l’auteur.}
            
\end{teiFigure}
            \teiP[xmlid={p4-16}]{Après le 11 mars, il semble dans un premier temps que cette composition très rigoureuse vole précisément en éclats. La première photographie après la page blanche présente en effet un trop-plein de débris, un désordre des couleurs, qui rompent avec la stabilité des compositions précédentes. Toutefois, dès cette première image, il est à souligner la permanence de la structure tripartite de l’image et la large surface du ciel qui fait respirer l’image en dépit de l’étouffement des gravats au sol. Plus encore, des compositions similaires à celles de la première partie, tripartites ou coupées par des lignes de fuite nettes, se retrouvent assez vite : les gravats, l’eau et les montagnes ; les ruines, la forêt et le ciel ; la rivière, la terre et le ciel. De la même manière, on retrouve des photographies organisées selon un élément central, ce qu’il reste du pont, d’une maison. Là non plus, l’image ne bouge pas, comme si elle semblait arrêtée en une simple surface, refermée sur sa rigoureuse structure en aplats. Mer liquide et mer de gravats s’équilibrent et échangent leurs caractéristiques. Avant comme après, les compositions plastiques des photographies sont ponctuées de rimes et de symétries, accentuées par le reflet de l’eau qui dédouble l’image. Les images ont une direction, un pli, mais en même temps une ligne vient toujours la couper comme si l’image était toujours double. Le ciel se reflète sur la mer, sur la rivière. Avant la page blanche, ces deux canoteurs qui rament au même rythme levant la pagaie ensemble et se mirant ensemble dans le reflet de l’eau construisent à la fois une image d’une grande stabilité et sérénité et en même temps poussent le principe du dédoublement à un point tel qu’il semble que la faille soit déjà là, montée à la surface, à l’ultime bord de la catastrophe. Le reflet, c’est un double inversé, modifié. Le feu d’artifice dans le ciel fait écho à l’étoile de mer au sol. L’envers et l’endroit, l’avant et l’après, le haut et le bas, tout se construit autour d’un axe de symétrie, cette cassure qui est aussi celle du récit. En effet, à la structure très nette – l’opposition entre un passé idéalisé et l’horreur du présent – il faut introduire une nuance qui est de taille. Si les premières photos de ruine évoquent un chaos sans ordre et sans espoir, peu à peu au cours du livre des images à nouveau très composées apparaissent, très structurées et évoquant finalement le même apaisement que dans les photographies de la première partie, une même sérénité. Ce sentiment s’appuie sur un très grand sens du cadrage, une très grande beauté des lumières, des miroitements et des reflets, tout à fait similaire à celle de la première partie, et donc qui ne peut pas être sans l’évoquer.}
            \teiP[xmlid={p5-16}]{Aussi, si l’ouvrage se construit bien sur une opposition nette entre deux parties, articulée autour d’une page blanche, lieu d’une ontologie négative de la catastrophe, il y a également derrière cette opposition une continuité, un caractère commun qui est celui de la beauté, sous-tendue par cette présence continue de l’eau et de sa surface en miroir. De fait, la beauté n’est pas tant la caractéristique de la représentation des décombres suite à la catastrophe, mais celle de la continuité du regard du photographe. C’est peut-être également le signe de la continuité du paysage, de son identité propre, de la beauté singulière d’un territoire qui, en dépit de son complet bouleversement, conserve une topologie particulière, un tempérament particulier, lié à la présence de l’eau dans le cas de Rikuzentakata, et que le photographe retrouve au milieu des gravats. On peut se demander cependant pourquoi ce dernier maintient sur son pays natal dévasté cette exigence compositionnelle, cette beauté de surface en dépit de la profonde désolation, et de la catastrophe qui l’a atteint personnellement. Pourquoi maintenir la beauté dans l’horreur ? Comment justifier cette « sidérante splendeur parmi un paysage de pleurs » ?}
            \teiP[xmlid={p6-16}]{Il faut prendre au sérieux ce mouvement instinctif de refus de voir de la beauté dans l’horreur. Et à la gêne que l’on peut éprouver face à ces horreurs magnifiques, Naoya Hatakeyama fait cette réponse elle-même magnifique et terrible : « c’est comme maquiller les morts ».}
            \teiP[xmlid={p7-16}]{On est en droit d’interroger la dimension morale ou immorale de la beauté. Cette beauté des ruines n’a-t-elle pas le parfum immoral de celui qui se repaît de l’horreur, bien à l’abri de ses atteintes ? N’y a-t-il pas quelque chose de malsain à transformer le désastre en spectacle, à esthétiser l’horreur ? Naoya Hatakeyama n’essaie pas de donner un statut particulier à cette beauté. Le spectateur ressent bien qu’il ne s’agit pas d’un effet de beauté, mais d’un besoin originaire que ce soit beau. Naoya Hatakeyama est très conscient du fait qu’il y a beaucoup de sentiment, de nostalgie, dans ses images et que cela passe par le soin apporté au cadrage. Jamais Naoya Hatakeyama n’a eu l’idée de déconstruire la beauté, mais en même temps ce n’est jamais l’effet de beauté qui est recherché. Jamais l’horreur n’est mise au service d’un effet spectaculaire ou séduisant. Bien au contraire, ce sont le cadrage, la beauté de la composition qui confèrent aux photographies leur intériorité, leur épaisseur d’émotion et de sens. S’il faut parler d’instrumentalisation, c’est bien plutôt la beauté qui est instrumentalisée et non l’horreur. La beauté est l’instrument utilisé pour exprimer toute l’horreur et l’effroi du désastre. Plus exactement, cette beauté dit en même temps l’horreur et la réconciliation possible. Pour expliquer son utilisation de la beauté, Hatakeyama rapporte une tradition d’Okinawa. Il est d’usage, à Okinawa, de déterrer les os des morts et de les laver. Par le lavement s’exprime une nostalgie plus pure et sublimée. Hatakeyama raconte qu’en traversant les ruines du tsunami, c’était tellement horrible qu’il demandait à la photographie d’apporter un peu de beauté. C’est, pour lui, un peu comme maquiller les morts. Trouver les gestes, trouver le cadrage qui redonne de la beauté à ce qui l’a définitivement perdue. Les ruines, les os, sont pour Hatakeyama mieux que la disparition totale, c’est un lien vers le passé. Aussi, cette beauté ouvre la surface de l’image pour accompagner le spectateur vers les strates d’histoire personnelle qu’elle recouvre. Et l’eau a cette fonction rituelle de purification et d’apaisement. Si l’eau, la mer et la rivière sont si présentes dans les photographies de Hatakeyama, c’est précisément aussi peut-être pour cette ultime raison : exprimer, au moment même où l’horreur la plus ultime nous touche, ce besoin premier de s’en laver et de se réconcilier. L’eau agit à la manière d’un filtre magique qui d’image en image nous apaise et nous soigne. La photographie n’est plus alors un simple rectangle de papier posé à distance, mais bel et bien un acte, un geste d’apaisement, un rite de purification proposé au spectateur.}
            \teiP[xmlid={p8-16}]{Dans ses derniers travaux, Hatakeyama photographie les différentes étapes de la reconstruction de la région. De manière quasi cyclique, il referme son projet photographique sur des images qui pointent vers le futur de la ville. Là encore, la ruine, le présent sont comme pris en étau entre l’avant et l’après, entre ce qui a été perdu et ce qui se reconstruit. Parfois, Naoya Hatakeyama dit culpabiliser un peu par rapport à d’autres artistes dont les préoccupations sont plus engagées dans l’effort de reconstruction de la société, et lui aussi se demande ce qu’il peut faire en tant qu’artiste. Mais ce qui se trame dans ses images, cette façon de réinscrire le désastre dans un cycle qui englobe un passé et un futur apaisés, est déjà un geste de reconstruction, en ce qu’il a une action évidente de consolation des populations. Et, même si son travail n’est pas de directement aider les gens, il contribue au relèvement de la région en apportant – par cette beauté précisément – l’espoir. Il y a des personnes âgées en particulier qui ont beaucoup apprécié ses photographies, qui ont été très vivement émues et lui ont dit merci. L’eau, après avoir tout englouti, retrouve sa fonction lustrale d’absolution.}
            \teiP[xmlid={p9-15}]{Maurice Pinguet, dans son texte « Transparence et profondeur », compare les films d’Ozu à :}
            \begin{teiCit}[xmlid={cit02-9}]

              \begin{teiQuote}
une eau limpide et profonde, parfaitement tranquille. Si d’un rocher à pic, nous plongeons nos regards dans la mer, il faut que l’eau soit calme et pure pour que sa profondeur devienne visible. Et pour que sa transparence apparaisse, il faut que cette eau soit profonde. Alors notre vision ne s’arrête pas à la surface, elle ne vient pas non plus buter contre le fond trop proche, elle pénètre librement dans l’intimité de l’élément. Sans rencontrer d’obstacle, notre regard va toujours plus loin, nous voyons alors la transparence même, qui n’a ni forme ni couleur, nous saisissons ce qui est sans substance, nous saisissons la profondeur même\teiNote[n={3},place={foot},xmlid={ftn8-2}]{\teiP[xmlid={p-609}]{Maurice Pinguet, « Transparence et profondeur d’Ozu », dans Alain Jouffroy (dir.), \teiHi[rend={italic}]{Écritures japonaises. Cahiers pour un temps}, Paris, Centre Georges Pompidou, 1986, p. 109. }}.
\end{teiQuote}
            
\end{teiCit}
            \teiP[xmlid={p10-15}]{Et cette profondeur, c’est la dimension affective même de ce monde d’objets, l’envers ou la doublure de ce qui est offert à notre regard en surface, ces images d’objets, ces écailles, ces brisures de verre sans épaisseur. L’image ozuienne n’est pas tant un vide qu’une surface, la pellicule externe d’un quotidien vide ou dépeuplé. Si l’image ozuienne est cette surface sans ombre pleinement disponible, qui d’elle-même amoindrit sa propre densité, c’est précisément pour laisser voir en transparence cette doublure affective, qui est la vie même. « Notre attention fascinée par les films d’Ozu plonge dans la transparence de la forme et dans la profondeur de la vie\teiNote[n={4},place={foot},xmlid={ftn9-2}]{\teiP[xmlid={p-610}]{\teiHi[rend={italic}]{Ibid}.}} », écrit Maurice Pinguet. Image plate, géométrique, discrète, précise ou évidente, la surface des films d’Ozu n’atteint cette épure que pour nous laisser plonger dans le mouvement affectif de la vie humaine, ces douleurs et ces tourments, ce que Maurice Pinguet appelle dans un autre texte sur Ozu, « la dialectique de l’amour et du temps ». « Entre les êtres qui se sont aimés, les liens se dénouent s’ils ne se brisent pas, toutes les intensités se nivellent sous l’effet de l’inexorable entropie\teiNote[n={5},place={foot},xmlid={ftn10}]{\teiP[xmlid={p-611}]{Maurice Pinguet, « \teiHi[rend={italic}]{Voyage à Tokyo} d’Ozu », dans \teiHi[rend={italic}]{Le texte Japon. Introuvables et inédits}, réunis et présentés par Michaël Ferrier, Paris, Éditions du Seuil, 2009, p. 148. }}. » Dans la trame du temps du film, dans les dessous de ses lignes géométriques, se meuvent les « cruautés impersonnelles de la vie ». Et si cette profondeur vivante est à ce point lisible, sensible dans les films d’Ozu, au point de nous émouvoir aux larmes, c’est qu’elle s’appuie sur une esthétique de la surface, de l’image transparente et évidente, sur une beauté et une splendeur incontestables.}
            \teiP[xmlid={p11-15}]{Suite à la catastrophe du 11 mars 2011, un grand nombre de conférences, de colloques au Japon ou ailleurs ont posé la question de la rupture radicale que l’événement introduisait dans la pensée ou dans l’image que le pays se faisait de lui-même. Et donc, il me semble que le premier point important à noter est que, si la catastrophe est absolue rupture et non-sens, elle appelle plus que jamais le sens et la figure. Figurer la catastrophe donc, avant d’être possible ou non, est une exigence. La catastrophe est un non-sens qui appelle immédiatement l’expression. Et cette exigence n’est pourtant pas sans soulever cas de conscience et paradoxe. De fait, l’événement catastrophique semble déborder aussi bien les catégories de l’entendement que les limites de l’imagination, ou encore le cadre photographique. Plus encore, c’est la question même du figurable et de la visibilité qu’il pose. Comment rendre compte d’une telle catastrophe dans son ampleur et dans la césure qu’elle fonde ? C’est à ce nouveau défi, inédit et vertigineux, que se trouvent confrontés les artistes qui traitent du 11 mars. Comment répertorier, évoquer, représenter le réel quand il atteint ce stade de démesure et d’innommable ? Si les mots manquent, les images dans leur prolifération indéfinie pourront-elles figurer la catastrophe ? Figurer, c’est-à-dire donner un cadre à ce qui déborde tout cadre, donner une figure à ce qui n’en a pas, assigner une image, à ce qui les multiplie sans fin. Un photographe comme Naoya Hatakeyama lève ou contourne l’aporie en faisant de la coupe plate de l’image une surface vivante et sensible ouverte aux fantômes d’autres temps et hantée de ces différents possibles. Une photographie comme une réserve de sens, une objectivité doublée d’une intériorité, une image doublée d’un geste d’absolution. Comme ce pli de la vague sur la couverture de \teiHi[rend={italic}]{Kesengawa} qui, à la manière d’une doublure de manteau, rappelle la profondeur de cette beauté tout en surface. La facture classique et épurée des images d’Hatakeyama console en appliquant sur le chaos de la destruction le baume de la beauté. D’une manière un peu similaire, la dessinatrice Catherine Meurisse, après l’expérience traumatisante de l’attentat de Charlie Hebdo, se rattache au mot « beauté », comme à une bouée de secours.}
            \begin{teiCit}[xmlid={cit03-9}]

              \begin{teiQuote}
Dans \teiHi[rend={italic}]{La Légèreté}, j’ai formulé une quête de beauté. C’était la première fois que je le faisais, après avoir mûrement réfléchi, après avoir constaté que dans les chroniques qu’écrivait Philippe Lançon dans « Charlie » depuis l’hôpital, très blessé, il parlait de la beauté, et j’étais sciée, je ressentais aussi un besoin de beauté. Alors qu’on est dans l’horreur absolue, dans le chaos, qu’est-ce qui fait que soudain, ce mot arrive… ? Je me suis accrochée à ce mot\teiNote[n={6},place={foot},xmlid={ftn11}]{\teiP[xmlid={p-612}]{Catherine Meurisse, dans \teiHi[rend={italic}]{Cours particulier}, épisode 20, entretien sur France Culture du 21 janvier 2022 : \teiRef[target={https://www.radiofrance.fr/franceculture/podcasts/les-chemins-de-la-philosophie/catherine-meurisse-1787131}]{https://www.radiofrance.fr/franceculture/podcasts/les-chemins-de-la-philosophie/catherine-meurisse-1787131} (consulté le 3 juillet 2025).}}.
\end{teiQuote}
            
\end{teiCit}
            \teiP[xmlid={p12-15}]{Il y aurait donc deux manières de justifier cette horreur magnifique, cette beauté plastique et formelle au sein d’un paysage de larmes.}
            \teiP[xmlid={p13-14}]{Il faudrait tout d’abord dire que la beauté n’est pas coupée de la vie, de ses tourments, de ses mouvements agités. Comme Maurice Pinguet l’a dit au sujet d’Ozu, surface et profondeur sont solidaires et, à la manière d’un dispositif optique, la beauté de surface permet de saisir avec la plus grande acuité, la plus grande clarté, les douleurs de l’existence, l’entropie des sentiments, les déchirements des départs et des disparitions. Ainsi la beauté des images d’Hatakeyama ni ne gomme ni n’émousse la douleur du drame humain et l’horreur de la dévastation. À l’inverse, elle permet, à la manière d’une loupe ou d’un microscope, d’en ressaisir la trame, le fil, les ramifications fines d’une catastrophe intime.}
            \teiP[xmlid={p14-14}]{Mais secondement, et peut-être de manière un peu paradoxale par rapport à la première raison, la beauté est compatible avec l’horreur, non seulement parce qu’elle nous la fait voir, mais également parce qu’elle nous en soulage. La beauté, tel un baume, calme la douleur, restaure l’identité propre d’un territoire à moitié disparu, apaise les tourments liés à la disparition. La beauté agit comme un remède contre la douleur et il faut donner au mot « contre » toute son intensité. La beauté refait ce que la catastrophe a défait. La composition très pure et géométrique des images redonne, malgré le chaos, après le chaos, une continuité à une topographie centrée sur l’eau, redonne à voir l’éclat des lumières, le miroitement des reflets, et fait ainsi revivre le territoire que la catastrophe avait détruit. La beauté relève du soin et de l’attention au tissu du vivant.}
            \teiP[xmlid={p15-14}]{La beauté n’est donc pas un ordre géométrique et transcendant en dehors ou à rebours de la vie, elle est la digue qui donne à la vie toute son intensité et permet à la vie de retrouver un chemin dans le chaos. Tout comme dans la première image après la catastrophe prise par Hatakeyama, la beauté permet de retrouver une ligne, une voie au sein des décombres, et permet à l’image comme aux humains d’à nouveau « tenir debout ».}
          
\end{teiElement}
        
\end{teiElement}
        \begin{teiElement}[name={group},type={chapter},data-page-title={La beauté réellement vitale},data-page-authors={Frédéric Worms@@ENS-PSL, Paris, France},data-include-href={Ch16\_beaute\_reellement\_vitale\_Worms.xml},xmlid={chapter-015}]

          \begin{teiElement}[name={front}]

            \begin{teiDiv}[type={titlePage},xmlid={titlepage-017}]

              \teiP[rend={title-main},xmlid={p-613}]{La beauté réellement vitale}
              \teiP[rend={author-aut},xmlid={p-614}]{Frédéric \teiHi[rend={small-caps}]{Worms}}
              \teiP[rend={authority\_affiliation},xmlid={p-615}]{ENS-PSL, Paris, France}
            
\end{teiDiv}
          
\end{teiElement}
          \begin{teiElement}[name={body},xmlid={body-017}]

            \teiP[xmlid={p1-17}]{La beauté, c’est la vie contemplée : telle est la définition de la beauté à laquelle nous aboutirons ici. }
            \teiP[xmlid={p2-17}]{Mais elle implique de lever aussitôt un malentendu, dont les conséquences sont plus graves qu’on ne pense, et qui fera l’objet de cette étude. }
            \teiP[xmlid={p3-17}]{La vie contemplée : cette expression comporte en effet un risque majeur, celui de nous faire croire que, dans l’expérience de la beauté, la force de la « contemplation » serait telle qu’elle peut nous éloigner de « la vie », qui ne serait plus là que comme image ou, si l’on ose dire, pour la forme. Or, cela n’est pas du tout le cas selon nous, mais c’est bien en revanche une tentation très forte, presque irrésistible, issue de l’expérience même de la « beauté », et qui a plus d’une conséquence majeure.}
            \teiP[xmlid={p4-17}]{Cette tentation va bien au-delà d’un problème portant sur ce qu’on appelle parfois « l’esthétique ». Certes, dans ce domaine déjà, ce qui semble justifier l’attribution de la beauté, à un objet, une personne, une œuvre ou un paysage, semble avoir une force telle qu’elle semble dépasser « la vie », nous élever au-dessus d’elle ! Il n’est pas étonnant qu’il y ait un débat en philosophie pour savoir si le fait de relier le jugement ou le sentiment du beau, à « la vie », n’aurait pas quelque chose de réducteur. C’est un débat important et même essentiel, mais ce n’est pas le plus grave, de ceux que pose la relation de la beauté à la vie. Car il y a plus grave encore, non seulement du côté de l’origine, ou du jugement, mais des effets et des enjeux de la beauté, dans nos vies.}
            \teiP[xmlid={p5-17}]{Relier la beauté à la vie, ce n’est pas seulement se demander si l’on « juge » de la beauté à travers le vivant. Relier la beauté à la vie, c’est se demander si l’expérience de la beauté fait encore et réellement partie de la vie, dans ce qui la définit selon nous, à savoir sa lutte contre le négatif et contre la mort. Il ne s’agit donc pas seulement d’esthétique en un sens étroit. Il s’agit de savoir si la beauté est littéralement et réellement vitale. Et, par là, on entend ceci : la beauté est-elle réellement issue de notre vie, entendue comme lutte contre la mort, et si oui, de quelle façon, puisqu’elle semble d’abord (en nous la donnant à « contempler ») y échapper ? Comment la beauté est-elle vitale, en effet, mais d’abord au sens d’un besoin vital, sans la satisfaction duquel nous ne pourrions pas vivre ? Ce qui peut vouloir dire aussi, comme pour toute fonction vitale, de comporter aussi un risque mortel, celui d’être ambivalente, au moins dans son usage, et de devenir elle-même dangereuse et mortelle. Ne serait-ce pas le cas aussi de la beauté qui, comme tous les vivants humains le savent bien, peut devenir mortelle, et prouverait en quoi elle est « vitale », par ses dangers mêmes ? }
            \teiP[xmlid={p6-17}]{Il ne suffit donc pas de se demander, même si c’est une question essentielle et révélatrice dans toute l’histoire de la philosophie, si les critères du « beau » peuvent être vitaux, en quelque sens esthétique que ce soit, ou s’ils s’y arrachent.}
            \teiP[xmlid={p7-17}]{Il s’agit de se demander d’abord si l’expérience même du beau, comme contemplation de la vie, vient de la vie et, dans ce cas, ce que cela nous apprend sur « la vie » elle-même. }
            \teiP[xmlid={p8-17}]{Et il s’agira ensuite de comprendre comment elle comporte un enjeu vital, face à un contraire, dont nous dirons tout de suite qu’il n’est pas purement esthétique, ou qu’il va plus loin que la seule laideur, même s’il n’est pas réductible à la mort, pour ainsi dire, en général, mais a une force précise qui confirme par contraste le sens même de la beauté dans nos vies. Disons-le tout de suite, nous désignerons ce contraire par le terme précis de \teiHi[rend={italic}]{l’horreur.} Tel serait, selon nous, le véritable contraire de la beauté, confirmant combien elle est plus vitale et nécessaire encore que l’on ne croit. Ces deux exclamations : \teiHi[rend={italic}]{quelle beauté ! quelle horreur !} ne sont-elles pas, dans leur opposition profonde, les pointes les plus extrêmes de nos vies ? }
            \teiP[xmlid={p9-16}]{On voit pourquoi il importe de ne jamais oublier le rapport de la beauté à la vie.}
            \teiP[xmlid={p10-16}]{L’\teiHi[rend={italic}]{origine} vitale de la beauté (et d’abord, en effet, dans des exclamations), les \teiHi[rend={italic}]{critères} et aussi les \teiHi[rend={italic}]{paradoxes} de la beauté vitale (dans la philosophie et l’esthétique), mais surtout les \teiHi[rend={italic}]{enjeux} réellement vitaux de la beauté (jusque dans l’éthique et la politique), telles sont donc les étapes que nous devrons parcourir rapidement ici.}
            \begin{teiDiv}[type={section1},xmlid={div01-12}]

              \teiHead[xmlid={lorigine-vitale-de-la-beaute-la-vie-comme-relation-001}]{L’origine vitale de la beauté : la vie comme relation}
              \teiP[xmlid={p11-16}]{S’il faut revenir sur l’expérience de la beauté dans sa relation pour ainsi dire critique à la vie, c’est pour une raison qui concerne la beauté, bien sûr, sur laquelle on risque de se tromper, mais qui concerne aussi la vie elle-même, pouvant nous conduire à l’abandonner elle aussi à une fausse et dangereuse conception de son sens. Notre hypothèse serait donc celle-ci : c’est qu’il y a quelque chose de si particulier et tout simplement de si fort, dans l’expérience de la beauté que, non seulement elle nous fait oublier son lien avec la vie, mais qu’elle semble impliquer un contresens sur la vie elle-même ! On comprend l’importance de corriger cette erreur.}
              \teiP[xmlid={p12-16}]{Mais quel serait ce contresens qui se glisse selon nous jusque dans les expériences et en fait les plus intenses exclamations relevant de la beauté, des exclamations intenses qui comme telles sont pourtant des indices de sa relation avec la vie ? }
              \teiP[xmlid={p13-15}]{On peut le décrire de manière simple.}
              \teiP[xmlid={p14-15}]{Que se passe-t-il en effet à chaque fois que nous nous exclamons « que c’est beau ! » devant un objet, une œuvre, ou un paysage, à chaque fois que nous disons « tu es beau » ou « tu es belle » devant un visage ou un corps ? }
              \teiP[xmlid={p15-15}]{Une telle exclamation paraît simple en effet à résumer : elle surgit d’une expérience pour nous entraîner vers une autre. Nous étions pris dans une action, un travail, une conversation, que sais-je, « la vie », donc, et nous voici entraînés vers autre chose : la contemplation, qui nous paraît tout de suite (d’un coup, pourrions-nous dire) infinie, inépuisable, d’un objet, d’un être ou du monde ! Nous voici attirés par cette contemplation, irrésistiblement. Et la « beauté » est le nom lui-même le plus intense qu’on peut donner à cette expérience intense qui consiste d’abord en ceci : se tourner vers l’objet, l’observer, l’écouter, le contempler pour lui-même et de manière, encore une fois, qui nous semble infinie. }
              \teiP[xmlid={p16-14}]{Mais, on le voit, cette expérience qui paraît simple est en fait double : elle implique non seulement de se tourner vers la beauté, mais aussi de se détourner de la vie ou même de s’arracher à elle, aussitôt (quoique silencieusement ou implicitement) entendue comme tâche ingrate, bornée et mortelle. Ne dit-on pas que le paysan ne peut pas admirer le paysage, parce qu’il doit y travailler ? Mais couper le paysage du paysan, ce n’est pas seulement couper la contemplation ou la beauté de la vie, c’est réduire celle-ci sans même y penser, dans un arrière-plan aveugle et pour ainsi dire dans le noir ; ce n’est pas seulement s’empêcher de comprendre comment la beauté peut se relier à la vie, mais comment la vie peut s’ouvrir sur la beauté. }
              \teiP[xmlid={p17-14}]{On le voit, la tentation a beau être irrésistible, et peut-être inséparable de l’expérience de la beauté, dans l’exclamation même qu’elle nous arrache, elle n’en comporte pas moins une conséquence et en fait une erreur, particulièrement grave. Car, en réduisant « la vie » à une tâche solitaire, aveugle, sans regard et sans distance, on ne s’empêche pas seulement de comprendre l’expérience de la beauté, on s’empêche de comprendre le vivant lui-même et ce que, dans sa fragilité, il comporte justement et nécessairement de relation, d’ouverture et de création, ce pour quoi il a « besoin » aussi de la beauté et peut souffrir d’en être privé ou de la voir détournée de sa dimension vitale pour devenir, elle aussi, mortelle. C’est donc bien du « vivant » lui-même qu’il faut repartir, pour comprendre comment la beauté peut en sortir, même si le contresens sur la beauté ne cessera de revenir. }
              \teiP[xmlid={p18-13}]{Nous disons : du « vivant » ou des « vivants » car justement, tout est là. Définir « la vie », comme nous le ferons ici, par la lutte contre la mort, c’est justement s’interdire de parler de la vie en général (qui, on le sait, ne meurt pas), mais parler toujours, au contraire, de vivants \teiHi[rend={italic}]{singuliers}, dont la vie est toujours précaire, menacée par une mort bien réelle, qui affrontent donc des dangers toujours plus singuliers eux aussi, dans une « vie » qui est donc par essence non seulement mortelle, mais \teiHi[rend={italic}]{relationnelle}, pour le meilleur et pour le pire. Il y a certes des degrés dans cette relation du vivant au monde et aux autres vivants. Comme l’a montré Bergson dans le premier chapitre de \teiHi[rend={italic}]{Matière et mémoire}\teiNote[xmlid={ftn1-16},n={1},place={foot}]{\teiP[xmlid={p-616}]{Henri Bergson, \teiHi[rend={italic}]{Matière et mémoire} \lbrack{}1896\rbrack{}, Édition Critique, Paris, PUF, « Quadrige », 2009.}}, la perception des êtres et du monde est d’autant plus distante que le contact vital avec eux est moins pressant. La distance mesure en quelque sorte la liberté de l’action, du contact immédiat et peut-être mortel, à un futur ouvert sur une action possible, mais encore liée à la vie. Ce n’est donc certes pas encore une « contemplation » qu’on suppose désintéressée, qui en tout cas s’absorberait dans l’objet pour lui-même, que cette distance de la « perception » pratique ou vitale. Et Bergson lui-même, quant à lui, fera intervenir une rupture radicale entre la perception pratique et la vision esthétique ou artistique, qui doit couper avec elle. Il n’en est pas moins vrai que, dès que l’on pose un vivant, on pose trois réalités et non pas une seule : un être singulier, d’autres êtres singuliers qui peuvent être vitaux ou mortels pour lui, et la \teiHi[rend={italic}]{relation} entre eux qui n’en est pas moins réelle que ces êtres même ! Ce que nous appelons « la vie » n’est pas une force obscure et opaque, mais est déjà une ouverture qui, certes, reste vitale donc polarisée (par le danger ou la joie), mais qui conduit déjà à la perception du monde. Lorsque nous nous exclamons devant la beauté d’un être, nous croyons passer d’un contact obscur à une contemplation pure, mais peut-être nous trompons nous sur l’une \teiHi[rend={italic}]{et sur l’autre} : peut-être la contemplation ne sera-t-elle jamais pure, tandis que la vie, en réalité, était déjà contemplative, \teiHi[rend={italic}]{et aussi expressive et créative}, comme le montre l’exclamation elle-même, exclamation qu’on oublie, \teiHi[rend={italic}]{elle aussi}, dans le sujet qui l’exprime, tant elle semble nous tourner, exclusivement, vers ce qu’il contemple ?}
              \teiP[xmlid={p19-13}]{Mais il faut aller plus loin, et faire une dernière remarque qui nous paraît essentielle. C’est que l’expérience de la beauté ne nous masque pas seulement (par sa force contemplative propre pour ainsi dire) la dimension relationnelle de la vie ou plutôt des vivants en général. Elle nous masque aussi et surtout, la forme que prend cette dimension relationnelle chez certains vivants, et notamment chez les vivants humains ; forme qui, en effet, les ouvre de l’intérieur à la contemplation, à l’expression et même à la création du monde ! La véritable rupture n’est pas entre la contemplation et « la vie », mais entre les degrés des relations mêmes, entre les vivants, ou entre les vivants et le monde ; toute vie étant bien entendu individuelle et relationnelle, mais, par « l’évolution » même des vivants sur notre planète, de manière différenciée et de plus en plus ouverte, avec aussi des risques internes, de plus en plus grands, sur les autres et sur le monde.}
              \teiP[xmlid={p20-13}]{La « beauté » n’est sans doute que le nom le plus intense donné par certaines cultures à une expérience qui est justement celle dont Winnicott a montré, dans toute son œuvre, qu’elle est \teiHi[rend={italic}]{à l’origine même de la culture} et de l’art, à savoir celle de la relation entre les vivants humains singuliers, qui les fait accéder aussi à cet « entre deux » qui est le monde lui-même, « le lieu où nous vivons\teiNote[xmlid={ftn9-3},n={2},place={foot}]{\teiP[xmlid={p-617}]{Donald Winnicott, \teiHi[rend={italic}]{Jeu et réalité} \lbrack{}1971\rbrack{}, Claude Monod (trad.),Paris, Gallimard, 1975.}} ». Tout se passe comme si, entre certains vivants, la relation ouvrait de l’intérieur à la reconnaissance de leur singularité et à l’exploration du monde qui les sépare et les relie, le terme extrême de cette reconnaissance étant justement la notion de beauté, qui masque pourtant sa genèse et sa fonction vitale d’autant plus grandes puisque, sans cette reconnaissance, le sujet lui-même, qui va s’exclamer devant le monde, ne peut pas exister ! Tout se passe alors comme si la première exclamation de beauté n’était pas celle que l’on croit. Ce n’est pas celle que nous appliquons au monde ou à l’art, aux œuvres ou au corps, et même au corps de l’être aimé. C’est d’abord celle dont nous avons fait l’objet : celle qui nous a fait vivre, dans le regard et la voix des autres penchés sur nos berceaux, la suspension elle-même vitale des gestes vitaux apparemment les plus urgents, pour la contemplation du visage singulier. Winnicott montre comment cette création mutuelle des sujets n’est pas seulement, comme dans l’éthologie classique (celle de Konrad Lorentz par exemple) une ruse de la vie pour susciter l’attachement : le « qu’il est mignon » qui permet en fait le soin du nouveau-né et la perpétuation de l’espèce. En fait, la contemplation mutuelle des visages et l’épreuve de ce qui sera appelé leur beauté, et de sa résistance à la destruction, fait accéder aussi de l’intérieur de la relation vitale à la contemplation du monde et à la création des œuvres, pour elles-mêmes ! L’exclamation est découverte et geste à la fois, contemplation et expression, absorption et création, parce qu’entre les vivants s’ouvre le monde et l’objet, si puissamment qu’elle fait oublier son origine, dans cette relation vitale. }
              \teiP[xmlid={p21-11}]{Ainsi, on comprend que la beauté soit le nom le plus élevé de cette contemplation, qui n’arrive jamais à faire totalement oublier son origine et donc aussi sa précarité vitale. Certes, elle s’absorbera à juste titre et infiniment dans la contemplation de l’objet et la création des œuvres, mais son origine et sa précarité relationnelles et vitales ne cesseront de revenir. La preuve en est que la précarité de la beauté est double ou, si l’on veut, on peut la détruire de deux façons : dans l’objet et le monde où nous vivons, mais aussi dans la capacité même du sujet à y accéder et à s’exclamer. Ce n’est pas le travail ou la vie en général qui nous masque la beauté ou nous en prive, mais la destruction bien plus radicale de la subjectivité par la misère, la maladie ou la torture. La nature vitale de la beauté s’atteste à ce que son contraire n’est certes pas la vie, mais qu’elle a le même contraire que la vie : à savoir la laideur ou la mort et, en fait, comme nous le verrons de manière plus précise, ce que nous exprimons sous le terme de l’horreur.}
              \teiP[xmlid={p22-11}]{Ainsi, on voit comment l’exclamation qui porte sur la beauté nous masque non seulement l’origine de la beauté dans la vie, mais l’ouverture de la vie sur la beauté, avec ce qu’elle comporte d’essentiel, et notamment ses enjeux pratiques dans la vie humaine, pour le meilleur et pour le pire. }
              \teiP[xmlid={p23-11}]{On aurait envie, bien sûr, d’aller tout de suite vers ces enjeux, pratiques ou vitaux, entre les humains.}
              \teiP[xmlid={p24-10}]{Mais il est certain que la tentation de couper la beauté de la vie, comme de couper la vie de la beauté (ce qui n’est pas moins grave), qui ne cesse de revenir, a engendré les contresens les plus nombreux de l’histoire de la philosophie et de « l’esthétique », contresens qui à leur tour pèsent non seulement sur nos rapports à la beauté, mais à nos vies d’où elle provient.}
              \teiP[xmlid={p25-9}]{Il faut donc confirmer l’hypothèse que l’on vient d’avancer sur l’origine vitale de la beauté par une esquisse des paradoxes de la beauté vitale, avant de revenir à ces enjeux pratiques dont on comprendra, alors, toute la portée.}
            
\end{teiDiv}
            \begin{teiDiv}[type={section1},xmlid={div02-12}]

              \teiHead[xmlid={le-paradoxe-de-la-beaute-la-contemplation-vitale-001}]{Le paradoxe de la beauté : la contemplation vitale}
              \teiP[xmlid={p26-9}]{Toute l’histoire de la philosophie le montre ou plutôt le confirme : la force de l’expérience de la beauté est telle, qu’on tend à opposer la beauté à la vie, ce à quoi la vie, pourtant, ne cesse pour ainsi dire de résister ! }
              \teiP[xmlid={p27-9}]{C’est même ce qui engendre tous les paradoxes et toutes les querelles de l’esthétique, ou de la philosophie, sur la place du beau dans la nature, dans l’art et dans nos vies. }
              \teiP[xmlid={p28-8}]{Ne serait-il pas simple, pourtant, de tirer de ce que l’on vient de dire un critère de la beauté, qui nous reconduirait tout de suite d’ailleurs vers ses enjeux pratiques ? }
              \teiP[xmlid={p29-8}]{De fait, il nous semble que nous appelons beau ceci : l’objet contemplé pour lui-même et dans sa forme, mais en tant qu’il nous apparaît \teiHi[rend={italic}]{aussi} comme vivant et en relation avec nos vies ; ou, à l’inverse, la vie \teiHi[rend={italic}]{la plus intense}, lorsque cependant nous la rencontrons \teiHi[rend={italic}]{exprimée et contemplée} dans un objet, un corps, ou une œuvre, donc dans une \teiHi[rend={italic}]{autre relation} avec nous ? L’objet le plus fonctionnel, par exemple, lorsqu’il nous apparaît aussi comme vivant et en relation avec notre vie, un mur, une maison ou une table par exemple, quand ils nous semblent comme dessinés et adressés, ne devient-il pas déjà quelque chose comme une \teiHi[rend={italic}]{œuvre} et n’en rendons-nous pas grâce à celui ou celle qui l’a mis sous nos yeux ? Mais, inversement, l’intensité la plus grande de la nature ou du monde, mais aussi d’un être vivant ou humain, lorsque nous la contemplons dans un paysage ou dans une œuvre, un soleil couchant sur la mer ou dans un tableau par exemple, ne nous absorbe-t-elle pas \teiHi[rend={italic}]{comme telle} pour en comprendre la richesse deux fois infinie, l’universalité et en même temps la singularité, avec des degrés qui sont tous les degrés mêmes, les plus disputés et discutés d’ailleurs, du « beau » ? N’y a-t-il pas un besoin de ressentir, de voir et d’exprimer la vie comme la relation qu’elle est, au monde et aux autres vivants, en un sens du mot « sensation » qui n’est pas seulement esthétique ou abstrait, mais vital et profond ? N’est-ce pas là ce qui fait le double pouvoir de la beauté sur nos vies, celui de les absorber, mais aussi de les mobiliser, de les fasciner, mais aussi de les enthousiasmer ? Et ne serait-ce pas là aussi ce qui fait son enjeu non seulement esthétique, mais pratique, jusque dans les usages collectifs ou sociaux, éthiques et politiques que l’on peut en faire ?}
              \teiP[xmlid={p30-8}]{C’est bien ce qu’il nous semble.}
              \teiP[xmlid={p31-7}]{Et pourtant, les philosophies, même les plus grandes, ont eu du mal à le reconnaître, peut-être parce qu’il faut pour cela partir d’une philosophie complète et critique à la fois de la vie et des vivants. }
              \teiP[xmlid={p32-6}]{Certes, les plus profondes ont bien sûr reconnu la marque de cette tension, même si elle devient chez elles un paradoxe : celles qui ont essayé d’arracher la beauté à la vie ont bien vu comment elle y résistait, mais cela devient une tension ; et, en quelque sorte, inversement, chez les philosophies partant de la vie, la beauté semble devenir un problème, mais semble aussi résister, à travers son effet sur les vivants humains ! Il n’est pas question ici, bien sûr, d’explorer tous ces paradoxes de la beauté vitale. Nous n’en prendrons donc que trois exemples, parmi les plus profonds qui soient, qui nous ramèneront aux enjeux vitaux de la beauté, eux aussi plus profonds qu’on n’imagine.}
              \teiP[xmlid={p33-5}]{Il faut bien sûr dire un mot pour commencer de Platon et du \teiHi[rend={italic}]{Banquet} : c’est au cœur de ce dialogue sur l’amour en effet que, par une sorte de nécessité profonde, naît la réflexion sur « le beau » dans sa relation critique avec la vie. Ce paradoxe est porté par « l’amour » lui-même tel que le voit Platon : il est une force vitale, mais creusée par un « manque » qui va entraîner au-delà de la vie. Aimer, c’est une force, mais tournée vers une absence. Que cherchons-nous, que désirons-nous ? C’est cela qui nous donnera le sens de nos vies. Or, la réponse portée par Socrate, qui dit l’avoir entendue d’une prêtresse, Diotime, est la suivante : c’est la beauté. La preuve en est que, dans l’amour d’un « beau corps », la beauté peut et doit se détacher de la vie, de la relation vitale entre les corps, ce détachement conduisant, comme on sait, d’un corps contemplé à la contemplation de la beauté dans plus d’un corps, et finalement au-delà des corps, « en elle-même », dans la lumière de l’être et du bien. Pourtant, on l’oublie, le critère de la beauté reste celui de l’amour : l’enthousiasme de l’enfantement qui échappe à la crainte de la mort (voir 206e-207a puis, \teiHi[rend={italic}]{in fine}, 212a). Ce ne sera certes pas l’enfantement biologique, mais l’enthousiasme des amants, qui se traduit par des discours comme dans le \teiHi[rend={italic}]{Banquet} lui-même, devenant ainsi la preuve en acte, par la beauté de ceux-ci, de sa thèse sur la beauté s’exprimant vitalement, dans ces discours. Extraordinaire démonstration, qui traverse le discours dans tout ce dialogue. Cependant, ce critère finalement vital de la beauté (contre la mort, qui plus est) fait vibrer tous ses autres critères (le rapport à la science, à la contemplation, à la forme et au bien), tout l’arrachement de la beauté à la vie, jusque dans la résistance du corps des amis un peu amants et des convives un peu ivres, malgré tout. Aucune contradiction bien sûr, selon nous : bien plutôt, le signe que la beauté est le nom de la contemplation des vivants et, à travers eux, du monde et de tous les objets du monde, mais aussi, avec ses effets vitaux, depuis l’amour d’un corps jusqu’à celui du savoir, de la sagesse, de la discussion, de la philosophie et de la justice.}
              \teiP[xmlid={p34-5}]{Toute l’histoire de la philosophie va cependant opposer l’un des deux aspects du discours de Diotime, Socrate, et Platon, à l’autre et la beauté à la vie, comme si c’était possible. On peut tracer une grande ligne de fracture entre deux esthétiques : l’une, idéaliste, qui va arracher progressivement le beau au corps, et que l’on retrouve bien sûr chez Hegel, pour qui cela se traduit dans toute l’histoire de l’esprit dans le monde, des pyramides d’Égypte à la littérature universelle, ou chez Alain, qui va, de son côté, de la parade militaire au charme du poème ; l’autre, vitaliste d’une manière ou d’une autre, mais souvent peu critique, réduisant le beau à son effet « sensible », comme le mot esthétique veut certes le dire, ou encore comme la notion de « goût » l’implique, qui renvoie toute l’esthétique à un genre de sensation, le plus vital et le plus subjectif qui soient d’ailleurs, mais qui conduit à ce qu’on appelle parfois une « guerre » du goût : c’est bien le cas, en effet, dans nos sociétés.}
              \teiP[xmlid={p35-5}]{Il n’est pas question ici de parcourir, même à grands traits, cette vaste querelle. Prenons seulement deux des esthétiques vitalistes les plus profondes, justement parce qu’elles partent de la vie dans son opposition à la mort, ou en tout cas dans sa puissance critique, celles de Schopenhauer, prolongé et critiqué par Nietzsche. Pour Schopenhauer, la fonction non pas du beau, mais de l’art, est en effet de nous arracher à la vie, conçue non pas comme plaisir, mais comme souffrance. Dès lors, le « beau » n’est pas le but, c’est l’effet de la contemplation sur nos corps, dont la musique est alors le premier exemple, qui atteste de son arrachement à la vie. Paradoxe profond : l’esthétique est cet effet vital contre la vie même. Et la contemplation sera ici celle du détachement, inspiré des sagesses de l’Inde. Comment aller plus loin à la fois dans le besoin de l’art et dans le refus de son ouverture sur le monde et sur le beau ? Mais comment aussi ne pas ressentir une absence, celle de la contemplation et, sinon de la beauté, du moins de la relation aux êtres qui surgit de la vie elle-même ? Nietzsche ne maintiendra pas tout au long de son œuvre cet espace de la contemplation vitale. Il reviendra lui aussi à une conception de l’art comme force, affirmative cette fois. Et il faudrait déjà, pour le moins, concilier dans un vitalisme critique la fonction affirmatrice et émancipatrice de l’art, son rapport aux deux aspects ambivalents de la vie. Du moins, dans son premier livre, \teiHi[rend={italic}]{La naissance de la tragédie}\teiNote[xmlid={ftn10-2},n={3},place={foot}]{\teiP[xmlid={p-618}]{Friedrich Nietzsche, \teiHi[rend={italic}]{La Naissance de la tragédie} \lbrack{}1870\rbrack{}, dans \teiHi[rend={italic}]{Œuvres}, t. 1, Marc de Launay (dir.), Paris, Gallimard, « Bibliothèque de la Pléiade », 2000.}}, Nietzsche avait-il bien présenté Apollon et Dionysos comme les deux faces de l’art dans la vie, la contemplation, mais aussi l’entraînement, et leur mélange d’ailleurs dans la tragédie. Nietzsche décrit la force de la contemplation dans la beauté apollinienne, qui ne peut cependant rejoindre complètement « la vie » ; et la résistance vitale à « l’horreur » dans les forces artistiques dionysiaques, incapables cependant d’atteindre à la contemplation de la forme. Seule la tragédie peut combiner les deux, à travers la forme donnée à l’horreur et au chant, au chœur et au drame, et c’est la leçon la plus profonde qui soit, dans sa tension même. C’est pour nous, là aussi, un signe de ce que les plus profondes philosophies du beau ou de la vie doivent finalement retrouver leur relation.}
              \teiP[xmlid={p36-5}]{Mais la preuve que cela reste difficile, on la trouvera dans une troisième direction. C’est celle qui a été ouverte par Kant. C’est chez lui que l’on trouvera en effet la relation la plus profonde entre la beauté et la vie, à une réserve fondamentale près, c’est qu’elle n’est pas conçue par lui comme \teiHi[rend={italic}]{réelle} ! Kant voit bien que la vie ouvre sur la beauté qui concentre en elle le sens de la vie : dans ce que l’on contemple où s’unissent toutes nos facultés, ou dans notre expression qui unit aussi toutes nos facultés, au moins dans le « génie ». Mais la vie comme la beauté, dans cette troisième \teiHi[rend={italic}]{Critique}\teiNote[xmlid={ftn11-2},n={4},place={foot}]{\teiP[xmlid={p-619}]{Emmanuel Kant,\teiHi[rend={italic}]{ Critique de la faculté de juger }\lbrack{}1790\rbrack{}, dans \teiHi[rend={italic}]{Œuvres philosophiques}, t. III, Ferdinand Alquié (dir.), Paris, Gallimard, « Bibliothèque de la Pléiade », 1980.}} qui ne nous apprend rien sur les choses, mais seulement sur nous-mêmes, et aussi génialement décrites soient-elles, voici que ce ne sont pas pour Kant des faits réels, mais des interprétations du monde par les humains, qui n’ont pas de vrai contenu, sinon d’indiquer que les humains aspirent à un sens, et pour révéler une « destination » des vivants humains par-delà le monde réel. L’exclamation ou la création de la beauté, l’admiration ou le génie : cela exprime toutes les puissances de la vie du côté de l’objet comme du sujet, mais seulement « pour nous ». Admirable analyse, et frustrante conclusion. Cependant, Kant, en lui enlevant certes à tort tout enjeu théorique, découvre aussi l’enjeu pratique de la beauté. Lorsqu’un humain s’exclame : « c’est beau », il vise l’universel, il considère que cela vaut pour tous. C’est une indication majeure : le beau n’est pas seulement une ouverture sur le monde ni le signe d’une vie qui s’exprime en général ; c’est une force bien précise, \teiHi[rend={italic}]{entre les vivants humains}. Ainsi, Kant, en essayant de dépasser la querelle du beau et de la vie, nous ramène aux enjeux pratiques du beau dans la vie. Et, de fait, la preuve que la beauté est vitale, c’est la manière dont les humains en ont fait et en font un enjeu vital. C’est ce à quoi il faut désormais que l’on vienne.}
            
\end{teiDiv}
            \begin{teiDiv}[type={section1},xmlid={div03-11}]

              \teiHead[xmlid={lenjeu-vital-de-la-beaute-joie-et-horreur-dans-toutes-les-dimensions-de-nos-vies-001}]{L’enjeu vital de la beauté : joie et horreur, dans toutes les dimensions de nos vies}
              \teiP[xmlid={p37-4}]{Un point central résulte de ce qui précède : c’est que la beauté n’est certes pas coupée de la vie, mais est en quelque sorte la norme la plus haute, ou la joie, elle-même vitale, et avec des effets sur nos vies, de sa contemplation. On devrait ajouter d’ailleurs : la contemplation, non pas d’un seul aspect de nos vies, ou de nos vies réduites à leur « aspect » ou à un spectacle, mais bien de toutes leurs dimensions, de sorte que nous jugeons belles aussi des démonstrations scientifiques, des actions humaines, et parfois la vie elle-même (« la vie est belle » « c’est la belle vie », deux exclamations bien différentes, mais avec une intensité convergente), comme si c’était le plus beau compliment (justement) qu’on pouvait leur faire ! Il y a ainsi dans la beauté non pas une coupure avec la vie, mais au contraire sa plus haute intensité, dans toutes ses dimensions. }
              \teiP[xmlid={p38-2}]{Cependant, il est clair aussi que, si la beauté est le nom, dans notre culture, de ce qui, dans toute culture, est la norme la plus haute et joyeuse de la contemplation de la vie, c’est en ce qu’elle a, comme toute norme, un contraire, ou s’oppose (comme le dit, bien sûr, Canguilhem) à quelque chose « de valeur négative » et à du « pathologique ». Et bien sûr, là aussi, nous savons intuitivement que c’est bien le cas : nous savons même que cela peut prendre deux directions.}
              \teiP[xmlid={p39-2}]{La première est celle du négatif, pour ainsi dire, dans la vie elle-même, en tant qu’il s’offre aussi, évidemment, à la contemplation : voir la maladie et la mort, la haine et la guerre, et souvent tout cela ensemble, c’est aller bien au-delà de la laideur que l’on serait d’ailleurs bien en peine de définir. Non, de même que la beauté est la norme de toute la vie contemplée, de même ce qui désigne tout le négatif quand il se donne à voir avec une répulsion extrême, ce n’est pas la seule laideur, en un sens restreint perçu comme une non-conformité à une norme abstraite, c’est bien l’horreur, précisément en ce qu’elle suscite en nous le contraire de la contemplation, une répulsion qui est presque impossible à surmonter et qui est indissociablement vitale, esthétique et éthique. L’horreur : c’est justement ce qu’il est impossible de contempler ; et, bien sûr, cela confirme que la contemplation dont la norme est la beauté a, ou plutôt est, elle-même, une force vitale, c’est-à-dire à la fois précaire et nécessaire, fragile et fondamentale. Il est clair pourtant, et notamment dans le « moment » historique que nous vivons, défini lui-même par l’opposition de la vie et de la mort, il est clair donc qu’une des tensions les plus grandes de la vie et de l’art est de parvenir à affronter l’horreur sans la transmuer en beauté, mais en y voyant un ultime signe de contemplation de la vie et de la joie, à la manière dont Simone Weil parlait de l’\teiHi[rend={italic}]{Iliade} d’Homère comme d’un « poème de la force », c’est-à-dire montrant dans la lumière de la beauté l’horreur nue de la force\teiNote[xmlid={ftn12},n={5},place={foot}]{\teiP[xmlid={p-620}]{Simone Weil, \teiHi[rend={italic}]{L’}Iliade\teiHi[rend={italic}]{ ou le poème de la force} \lbrack{}1941\rbrack{}, dans \teiHi[rend={italic}]{Œuvres}, Paris, Gallimard, 1999.}}. C’est bien sûr l’une des questions suprêmes, non seulement de l’esthétique contemporaine (par exemple dans le récit de vie, de maladie ou de souffrance, y compris bien sûr, sinon même \teiHi[rend={italic}]{d’abord}, éthiques et politiques), mais de la philosophie critique du vivant ou du vitalisme critique, et il faudra y revenir.}
              \teiP[xmlid={p40-2}]{Mais il y a une deuxième direction du négatif qui prouve la réalité vitale de la beauté, à la pointe de l’expérience des vivants humains, donc dans leurs relations éthiques et politiques elles-mêmes : c’est la manipulation elle-même, polarisée ou ambivalente, des normes de la beauté dans la culture. Une œuvre à laquelle nous ne cessons de repenser pour en voir toute l’évidence est un film : \teiHi[rend={italic}]{Freaks} de Todd Browning\teiNote[xmlid={ftn13},n={6},place={foot}]{\teiP[xmlid={p-621}]{Todd Browning (réal.), \teiHi[rend={italic}]{Freaks}, États-Unis, Metro-Goldwyn-Mayer, 1932, 64 min.}}. Il va plus loin en effet que la dénonciation des effets d’une beauté abusivement ou socialement normative, à travers des « canons » imposés de façon explicite ou le plus souvent implicite, la beauté des « publicités » qui transmet aussi toutes les normes et les discriminations d’une société et qui appelle certes un travail critique. Dans ce film, les monstres difformes et les dieux du stade qui coexistent dans une troupe de cirque voient tous les signes se renverser : nous finissons par juger « horribles » à tous les sens du terme la diva qui ressemble à Marlene Dietrich et l’athlétique amant qui la soutient dans sa perversité, et par juger divinement beaux à tous les sens du terme non seulement le couple de nains qu’ils s’acharnent à détruire et à dépouiller, mais tous leurs amis accablés de pathologies si graves et difformes, mais qui restaurent la beauté morale, relationnelle et finalement vitale à travers leurs actions et leurs gestes. L’enjeu vital de la beauté ressort ici pleinement dans sa manipulation sociale ou politique, à laquelle répond le geste critique de l’œuvre d’art. Il ne suffit pas de dénoncer l’usage violent et parfois criminel des normes de beauté, il faut en produire d’autres, qui rayonnent de la beauté des corps jusqu’à la beauté des actes et finalement des relations entre les vivants, et en particulier de façon critique entre les vivants humains.}
              \teiP[xmlid={p41-2}]{Car, nous y revenons bien sûr, de même que l’origine vitale de la beauté se trouvait dans ces relations, de même on y trouvera celle de leur contraire : c’est l’horreur de la \teiHi[rend={italic}]{violation}, comme destruction intime de la subjectivité dans la relation même qui la rend possible, qui surpasse toutes les autres, et qui rend la beauté deux fois impossible, non seulement dans l’objet à contempler, mais bien sûr dans la capacité même du sujet à la contempler. Nous retrouvons ici les deux extrêmes humains de la destruction de la beauté : non seulement dans la laideur et parfois l’horreur du monde, mais aussi dans la déformation et plus radicalement la destruction de notre capacité même à la percevoir et à nous en réjouir. De ce point de vue, le kitsch parfois désolant de la « belle vie » doit aussi être respecté comme une exigence vitale, de même que la capacité parfois nécessaire et parfois si fragile de trouver la vie « belle » même dans l’horreur (et le film lui aussi emblématique, \teiHi[rend={italic}]{La Vie est belle} de Roberto Benigni, dans l’horreur pourtant maintenue des camps de la mort, à travers la préservation par la fiction de la subjectivité de l’enfant, a tout d’une fable sur la beauté vitale\teiNote[xmlid={ftn14-2},n={7},place={foot}]{\teiP[xmlid={p-622}]{Roberto Benigni (réal.), \teiHi[rend={italic}]{La vita è bella}, Italie, Melampo Cinematografica, 1997, 116 min.}}).}
              \teiP[xmlid={p42-2}]{Il n’est pas possible d’aller plus loin ici, tant cela nous entraînerait à la fois vers les puissances de la joie pour elle-même, et dans leur opposition maintenue aux dangers de l’horreur.}
              \teiP[xmlid={p43-2}]{Mais nous pouvons conclure désormais sur la beauté à travers la lecture d’un poème, ou de quelques-uns de ses vers les plus célèbres, que nous extrairons d’un recueil dont le titre ne cessera jamais de sonner comme un défi : \teiHi[rend={italic}]{Les Fleurs du mal}\teiNote[xmlid={ftn15-2},n={8},place={foot}]{\teiP[xmlid={p-623}]{Charles Baudelaire, \teiHi[rend={italic}]{Les Fleurs du mal} \lbrack{}1\teiHi[rend={sup}]{re} éd.\rbrack{}, Alançon, Éditions Auguste Poulet-Malassis, 1857.}}.}
            
\end{teiDiv}
            \begin{teiDiv}[type={section1},xmlid={div04-7}]

              \teiHead[xmlid={la-tout-nest-quordre-et-beaute-luxe-calme-et-volupte-001}]{« Là, tout n’est qu’ordre et beauté / Luxe, calme et volupté »}
              \teiP[xmlid={p44-2}]{Que disent en effet ces deux vers qui rythment comme un refrain \teiHi[rend={italic}]{L’Invitation au voyage} de Baudelaire ?}
              \teiP[xmlid={p45-2}]{Il y a comme trois degrés dans leur vérité.}
              \teiP[xmlid={p46-2}]{Le premier, c’est bien sûr ce qui semble distinguer la beauté de la vie : sa norme formelle, son critère esthétique, « l’ordre » qui résume tout cela ici, et qui s’incarne aussi dans la forme des poèmes et des vers, des rimes et des sons. Oui, la beauté, c’est aussi le contenu et la forme de ce qui est contemplé, dans lequel on se plonge infiniment et pour lui-même, un ordre qui contient tout, dans un spectacle, comme un tableau, et c’est aussi ce que décrit le reste du poème.}
              \teiP[xmlid={p47-2}]{Mais le deuxième degré de cette sagesse est déjà un retour de la vie, dans le deuxième vers : cette contemplation est encore amoureuse et vivante, elle surgit du corps de la femme aimée et d’un désir que la volupté ne résume pas, car il y faut aussi le calme qui apaise les tourments et le luxe qui retrouve son sens parfois nécessaire, en ce qu’il montre que la vie ne réside jamais simplement dans la satisfaction d’un besoin.}
              \teiP[xmlid={p48-2}]{Car tout cela s’oppose à autre chose. }
              \teiP[xmlid={p49-2}]{On invite au voyage, on aspire à aller « là », c’est-à-dire, \teiHi[rend={italic}]{ailleurs}, par opposition, donc, à ce qui dans la vie est aussi horreur, et que tant de poèmes des \teiHi[rend={italic}]{Fleurs du mal} (les bien nommées) affrontent directement, au grand scandale de leurs contemporains : le flétrissement des corps, la charogne, l’épouvante et, bien sûr, la mort. }
              \teiP[xmlid={p50-2}]{L’aspiration à la beauté est vitale, non seulement parce qu’elle se relie au plaisir, mais parce qu’elle affronte le pire et peut tenter de le transmuer (comme dit Baudelaire) en « or », comme si la vie ne pouvait s’opposer à l’horreur de la mort, sinon par la joie.}
            
\end{teiDiv}
          
\end{teiElement}
        
\end{teiElement}
      
\end{teiElement}
      \begin{teiElement}[name={group},type={section1},xmlid={section1-005}]

        \teiHead[xmlid={varia-001}]{Varia}
        \begin{teiElement}[name={group},type={appendix},data-page-title={Un galet au fond d’un lac boueux},data-page-authors={Célia Boutilier@@Université PSL, laboratoire SACRe – EA 7410, France},data-include-href={Ch17\_varia\_un\_galet\_Boutilier.xml},xmlid={appendix-001}]

          \begin{teiElement}[name={front}]

            \begin{teiDiv}[type={titlePage},xmlid={titlepage-018}]

              \teiP[rend={title-main},xmlid={p-624}]{Un galet au fond d’un lac boueux}
              \teiP[rend={author-aut},xmlid={p-625}]{Célia \teiHi[rend={small-caps}]{Boutilier}}
              \teiP[rend={authority\_affiliation},xmlid={p-626}]{Université PSL, laboratoire SACRe – EA 7410, France}
            
\end{teiDiv}
          
\end{teiElement}
          \begin{teiElement}[name={body},xmlid={body-018}]

            \teiP[xmlid={p1-18}]{La beauté peut-être un mot-valise, polysémique, métamorphique, son sens se trouble à mesure que l’on essaie de la saisir, de la définir. Elle est un galet au fond d’un lac boueux. On peut parfois l’apercevoir lorsque l’on cherche bien ou si l’on a de la chance, lorsque le ciel est clair, que l’eau est transparente et que la lumière nécessaire est présente. Essayez de le saisir avec votre main, et au premier contact, la tourbe envahira l’espace qu’il occupe jusqu’à le faire disparaitre dans la pénombre. Sortez le galet du lac, et constatez que l’instant ne se trouve plus en lui.}
            \begin{teiFigure}[xmlid={figure01-7}]

              \teiHead[xmlid={figure-27-platane-dorient-environ-250-ans-circonference-9-30-m-diametre-2-95-m-hauteur-33-m-envergure-30-25-m-classe-arbre-de-la-manche-2005-2022-cerisy-la-salle-celia-boutilier-001}]{Figure 27. \teiHi[rend={italic}]{Platane d’Orient}, environ 250 ans, circonférence 9,30 m, diamètre 2,95 m, hauteur 33 m, envergure 30 × 25 m. Classé « arbre de la Manche 2005 », 2022, Cerisy-la-Salle. © Célia Boutilier.}
            
\end{teiFigure}
            \begin{teiFigure}[xmlid={figure02-6}]

              \teiHead[xmlid={figure-28-les-vaches-de-cerisy-2022-cerisy-la-salle-celia-boutilier-001}]{Figure 28. \teiHi[rend={italic}]{Les vaches de Cerisy}, 2022, Cerisy-la-Salle. © Célia Boutilier.}
            
\end{teiFigure}
          
\end{teiElement}
        
\end{teiElement}
        \begin{teiElement}[name={group},type={appendix},data-page-title={Quelques facéties sur le colloque « Beautés vitales » de Cerisy-la-Salle},data-page-authors={Mathias Dambuyant@@CNRS-EHESS, LAP – UMR 8177, France},data-include-href={Ch18\_varia\_quelques\_faceties\_Dambuyant.xml},xmlid={appendix-002}]

          \begin{teiElement}[name={front}]

            \begin{teiDiv}[type={titlePage},xmlid={titlepage-019}]

              \teiP[rend={title-main},xmlid={p-627}]{Quelques facéties sur le colloque « Beautés vitales » de Cerisy-la-Salle}
              \teiP[rend={author-aut},xmlid={p-628}]{Mathias \teiHi[rend={small-caps}]{Dambuyant}}
              \teiP[rend={authority\_affiliation},xmlid={p-629}]{CNRS-EHESS, LAP – UMR 8177, France}
            
\end{teiDiv}
          
\end{teiElement}
          \begin{teiElement}[name={body},xmlid={body-019}]

            \teiP[xmlid={p1-19}]{
              \teiHi[rend={italic}]{(Sons de cloche)}
            }
            \teiP[xmlid={p2-18}]{Oyez, oyez… salutations à toutes et tous !}
            \teiP[xmlid={p3-18}]{Pendant plus de cinq jours, nous nous sommes rencontrés sur le thème de la beauté vitale…}
            \teiP[xmlid={p4-18}]{Nous avons d’abord été éclairés par les réponses philosophiques : peut-on analyser des beautés au débotté ? Les sociologues et les anthropologues nous ont appris à recentrer notre regard et à trouver beaux les petites bêtes, les pierres, les jardins, les robots, les SDF et même les miss France ! Plus incroyable encore, il y avait de la beauté dans les pans du cerveau, le cerveau des paons, dans la seiche qui encre et dans l’encre qui sèche ; dans le jazz qui fait aimer et l’amour qui fait jaser ; le jour nous prenions des notes, le soir nous les écoutions !}
            \teiP[xmlid={p5-18}]{Merci à toutes et tous pour ce \teiHi[rend={italic}]{beau} colloque !}
          
\end{teiElement}
        
\end{teiElement}
        \begin{teiElement}[name={group},type={appendix},data-page-title={Des traces sur la beauté},data-page-authors={Lena Osseyran@@Metteuse en scène, actrice et architecte},data-include-href={Ch19\_varia\_croquis\_Osseyran.xml},xmlid={appendix-003}]

          \begin{teiElement}[name={front}]

            \begin{teiDiv}[type={titlePage},xmlid={titlepage-020}]

              \teiP[rend={title-main},xmlid={p-630}]{Des traces sur la beauté}
              \teiP[rend={author-aut},xmlid={p-631}]{Lena \teiHi[rend={small-caps}]{Osseyran}}
              \teiP[rend={authority\_affiliation},xmlid={p-632}]{Metteuse en scène, actrice et architecte}
            
\end{teiDiv}
          
\end{teiElement}
          \begin{teiElement}[name={body},xmlid={body-020}]

            \teiP[xmlid={p1-20}]{Ces croquis explorent la beauté intrinsèque du vivant, puisant dans la morphologie de paysages terrestres, marins et aériens, ainsi que dans les vibrations sonores qui les animent. À travers la matérialité du dessin, j’ai cherché à conserver des traces tangibles qui révèlent l’évolution dynamique des formes, où la beauté se manifeste aussi bien sur terre, en mer que dans les airs.}
            \teiP[xmlid={p2-19}]{Ces traits prennent vie à partir des moments partagés lors du colloque, nourris par les interventions sur le vivant, les paysages du château de Cerisy, son jardin, la visite de l’église, les chants et les échanges autour de la notion de beauté.}
            \begin{teiFigure}[xmlid={figure01-8}]

              \teiHead[xmlid={figure-29-graphite-aquarelle-multimedia-001}]{Figure 29. Graphite aquarelle, multimédia.}
            
\end{teiFigure}
          
\end{teiElement}
        
\end{teiElement}
        \begin{teiElement}[name={group},type={appendix},data-page-title={Le Goût de la Beauté},data-page-authors={Keiko Courdy@@Cinéaste et artiste multimédia},data-include-href={Ch20\_varia\_gout\_beaute\_Courdy.xml},xmlid={appendix-004}]

          \begin{teiElement}[name={front}]

            \begin{teiDiv}[type={titlePage},xmlid={titlepage-021}]

              \teiP[rend={title-main},xmlid={p-633}]{Le Goût de la Beauté}
              \teiP[rend={author-aut},xmlid={p-634}]{Keiko \teiHi[rend={small-caps}]{Courdy}}
              \teiP[rend={authority\_affiliation},xmlid={p-635}]{Cinéaste et artiste multimédia}
            
\end{teiDiv}
          
\end{teiElement}
          \begin{teiElement}[name={body},xmlid={body-021}]

            \teiP[xmlid={p1-21}]{\teiHi[rend={italic}]{Le Goût de la Beauté} est un film documentaire écrit et réalisé par Keiko Courdy, Pika Pika films, 2023, 34 minutes.}
            \teiP[xmlid={p2-20}]{La bande annonce peut être visionnée sur le site Film-documentaire : \teiRef[target={https://www.film-documentaire.fr/4DACTION/w\_fiche\_film/70127\_0}]{https://www.film-documentaire.fr/4DACTION/w\_fiche\_film/70127\_0}.}
            \teiP[xmlid={p3-19}]{Le film est disponible sur Vimeo : \teiRef[target={https://vimeo.com/770799324/bb7322d6cf}]{https://vimeo.com/770799324/bb7322d6cf}.}
            \begin{teiFigure}[xmlid={figure01-9}]

              \teiHead[xmlid={figure-30-affiche-du-film-001}]{Figure 30. Affiche du film.}
            
\end{teiFigure}
            \begin{teiFigure}[xmlid={figure02-7}]

              \teiHead[xmlid={figure-31-photogramme-issu-du-film-001}]{Figure 31. Photogramme issu du film.}
            
\end{teiFigure}
          
\end{teiElement}
        
\end{teiElement}
      
\end{teiElement}
      \begin{teiElement}[name={group},type={chapter},data-page-title={Les auteurs},data-include-href={Ch21\_presentation\_auteurs.xml},xmlid={chapter-016}]

        \begin{teiElement}[name={front}]

          \begin{teiDiv}[type={titlePage},xmlid={titlepage-022}]

            \teiP[rend={title-main},xmlid={p-636}]{Les auteurs}
          
\end{teiDiv}
        
\end{teiElement}
        \begin{teiElement}[name={body},xmlid={body-022}]

          \teiP[xmlid={p1-22}]{Catherine \teiHi[rend={small-caps}]{Baroin} est professeure en langue et littérature latines à l’université de Rouen Normandie. Elle est membre de l’ÉRIAC (Équipe de recherche interdisciplinaire sur les aires culturelles) et membre associé d’ANHIMA (Anthropologie et histoire des mondes antiques, UMR 8210). Ses travaux, menés dans une perspective anthropologique, portent sur la mémoire, le rapport des Romains au monde grec, les vêtements et le corps.}
          \teiP[xmlid={p2-21}]{Philippe \teiHi[rend={small-caps}]{Bonnin} est directeur de recherche émérite au CNRS. Architecte et anthropologue, il a créé le réseau de recherche franco-japonais JAPARCHI.}
          \teiP[xmlid={p3-20}]{Célia \teiHi[rend={small-caps}]{Boutilier} est diplômée et doctorante de l’École des Beaux-Arts de Paris-ENS, laboratoire SACRe, université PSL. Sa recherche porte sur la « photographie imaginale » et l’esthétique de la symbiose. \teiRef[target={https://www.xn--cliaboutilier-bhb.com}]{https://www.céliaboutilier.com}.}
          \teiP[xmlid={p4-19}]{Claudine \teiHi[rend={small-caps}]{Cohen}, paléontologue, philosophe et historienne des sciences, est directrice d’études à l’EHESS et directrice d’études cumulante à l’EPHE, université PSL, 3\teiHi[rend={sup}]{e} section SVT, laboratoire Biogéosciences. Ses recherches portent sur l’histoire des sciences de la vie et de la terre et les représentations de la Préhistoire.}
          \teiP[xmlid={p5-19}]{Gisèle \teiHi[rend={small-caps}]{Dambuyant} est maîtresse de conférences HDR en sociologie à l’université Sorbonne-Paris-Nord. Elle s’intéresse à trois champs distincts dans cette discipline : la sociologie du corps et de la précarité, en examinant les conditions d’existence des publics vulnérables ; la sociologie des professions, en analysant le secteur du travail et de l’intervention sociale ; la sociologie des formations, en s’interrogeant sur la place de l’Université française dans ce secteur. }
          \teiP[xmlid={p6-18}]{Mathias \teiHi[rend={small-caps}]{Dambuyant} est docteur en sociologie de l’EHESS et membre associé du laboratoire LAP (Laboratoire d’anthropologie politique).}
          \teiP[xmlid={p7-18}]{Keiko \teiHi[rend={small-caps}]{Courdy} est une cinéaste et artiste multimédia française. Elle écrit, réalise et produit des films et des installations en France et au Japon. Docteure de l’université de Tokyo, après des études de cinéma et de théâtre à l’université Sorbonne-Nouvelle, elle a enseigné les performances multimédias pendant 3 ans à l’université des Arts et du design de Kyoto.}
          \teiP[xmlid={p8-18}]{Carole \teiHi[rend={small-caps}]{Maigné} est professeure ordinaire de philosophie générale et systématique à l’université de Lausanne. Ses recherches portent sur les philosophies allemande et autrichienne des \teiHi[rend={small-caps}]{xix}\teiHi[rend={sup}]{e} et \teiHi[rend={small-caps}]{xx}\teiHi[rend={sup}]{e} siècles. Elle s’intéresse tout particulièrement à la philosophie de la culture (Warburg, Kracauer, Cassirer), à l’esthétique de la photographie et à la philosophie de l’art (formalisme esthétique, école viennoise d’histoire de l’art, Wölfflin, Klein).}
          \teiP[xmlid={p9-17}]{Guilhem \teiHi[rend={small-caps}]{Marion} a une formation musicale du Conservatoire à rayonnement régional de Lyon en composition, écriture et guitare jazz. Il possède aussi une maîtrise de musicologie de l’université Lumière-Lyon-2 ainsi qu’un master d’informatique musicale de l’Institut de recherche et coordination acoustique/musique (IRCAM). Il a soutenu sa thèse au Laboratoire des systèmes perceptifs au sein de l’ENS de Paris. Son travail, entre sciences cognitives, informatique et musicologie, consiste à comprendre les mécanismes cognitifs culturels supportant la perception musicale, la façon dont ils sont construits, encodés dans le cerveau et évoluent au cours de la vie.}
          \teiP[xmlid={p10-17}]{Rémi \teiHi[rend={small-caps}]{Mermet} est docteur en philosophie et en histoire de l’art, ancien postdoctorant auprès de la chaire Beauté·s de l’université PSL, et actuellement chargé de recherches FRS-FNRS à l’UCLouvain Saint-Louis Bruxelles. Son travail porte sur l’histoire et l’actualité de l’esthétique philosophique, dans ses divers rapports aux arts et aux sciences, humaines comme naturelles.}
          \teiP[xmlid={p11-17}]{Didier \teiHi[rend={small-caps}]{Nectoux} est conservateur du musée de Minéralogie de l’École des mines de Paris. Docteur en géologie de l’ingénieur, ingénieur de recherche dans l’industrie, il a exercé en tant qu’enseignant-chercheur durant 20 ans à l’École des mines d’Alès et, ces dernières années, en tant que responsable du musée, ainsi que créateur et responsable de formation d’ingénieurs par apprentissage dans le domaine de l’écoconstruction. Il est l’auteur de plusieurs ouvrages de vulgarisation, d’un MOOC sur les minéraux et dispense de nombreuses conférences à travers la France.}
          \teiP[xmlid={p12-17}]{Léna \teiHi[rend={small-caps}]{Osseyran} est une metteuse en scène, actrice et architecte libanaise, installée en France depuis 2019. Docteure en arts du spectacle, mention théâtre, elle a accompli une thèse en recherche et création au sein du programme doctoral RADIAN (ED Normandie Humanités-Université de Caen Normandie), portant sur \teiHi[rend={italic}]{La représentation des espaces incarnés à travers les corps dévastés et le théâtre contemporain}. Son travail examine les interactions entre l’architecture, le théâtre et les sciences cognitives, en concevant des environnements transdisciplinaires qui interrogent les dynamiques perceptives de l’étrangeté et du mystère inhérents à ces paysages.}
          \teiP[xmlid={p13-16}]{Bertrand \teiHi[rend={small-caps}]{Prévost} est historien de l’art et philosophe, professeur à l’université Bordeaux Montaigne (département d’Arts plastiques) et rattaché à l’unité de recherche Artes. Il a publié plusieurs ouvrages sur l'histoire et la théorie des images, de la Renaissance à l'art contemporain. Ses recherches consistent essentiellement en la construction d’une théorie de l’expression qui trouverait sa matière dans les formes et les pratiques cosmétiques, et son horizon dans une cosmologie esthétique.}
          \teiP[xmlid={p14-16}]{Claudine \teiHi[rend={small-caps}]{Sagaert} enseigne la philosophie en DNMADE (diplôme des métiers d’art et du design). Ses recherches au sein du laboratoire Babel de l’université de Toulon portent sur les représentations du corps dans une approche pluridisciplinaire : philosophie, esthétique et anthropologie. }
          \teiP[xmlid={p15-16}]{Pierre \teiHi[rend={small-caps}]{Sauvanet} est agrégé de philosophie, ancien élève de l’ENS Fontenay-Saint-Cloud, et professeur d’esthétique à l’université Bordeaux Montaigne, où il dirige l’UR ARTES (24141). Ses recherches (qui s’appuient aussi sur une pratique) portent sur une approche philosophique des phénomènes rythmiques. }
          \teiP[xmlid={p16-15}]{Chiara \teiHi[rend={small-caps}]{Vecchiarelli}, docteure en philosophie de l’ENS-PSL, enseigne la philosophie de l’art et l’esthétique à Sorbonne Université, où elle est ATER. Membre du Centre international d’étude de la philosophie française contemporaine (CIEPFC/PhilOfr) de l’ENS, elle a également enseigné à l’ENS Paris-Saclay, à l’université Sorbonne-Nouvelle et à l’université Rennes 2. Ses recherches portent sur la dimension ontogénétique de l’image et les implications à la fois métaphysiques et esthétiques de sa réalité vitale.}
          \teiP[xmlid={p17-15}]{Anne-Lise \teiHi[rend={small-caps}]{Worms }est professeure en langue et littérature grecques anciennes à l’université de Rouen Normandie (ÉRIAC). Ses recherches portent sur les conceptions de la beauté dans l’art, la littérature et la philosophie anciennes, notamment chez Plotin, ainsi que sur les relations entre philosophie et écriture dans l’Antiquité.}
          \teiP[xmlid={p18-14}]{Frédéric \teiHi[rend={small-caps}]{Worms} est professeur de philosophie à l’ENS-PSL. Il travaille sur les relations entre les vivants, ainsi que sur l’histoire de la philosophie française contemporaine, de Bergson à aujourd’hui, et défend un « vitalisme critique ». Il est l’auteur de nombreux ouvrages dont récemment \teiHi[rend={italic}]{Vivre en temps réel }et \teiHi[rend={italic}]{Le vivable et l’invivable} avec Judith Butler. Il a été membre du CCNE pendant huit ans, codirige la collection « Questions de soin » aux PUF, et produit sur France culture l’émission \teiHi[rend={italic}]{Le pourquoi du comment : philosophie}. Il a participé à de nombreuses rencontres à Cerisy-la-Salle et a notamment codirigé \teiRef[target={https://cerisy-colloques.fr/vivant-pub2016/}]{\teiHi[rend={italic}]{Le moment du vivant}} avec Arnaud François en 2012 (PUF, « Philosophie française contemporaine », 2016).}
          \teiP[xmlid={p19-14}]{Clélia \teiHi[rend={small-caps}]{Zernik} est professeure de philosophie de l’art à l’École nationale supérieure des Beaux-arts de Paris. Ses recherches portent sur le cinéma et l’art contemporain japonais, notamment après la catastrophe de mars 2011. Elle est présidente de la chaire Beauté·s de l’université PSL.}
        
\end{teiElement}
      
\end{teiElement}
      \begin{teiElement}[name={group},type={publisher-works},data-page-title={Les colloques de Cerisy},data-include-href={Ch22\_colloques\_de\_Cerisy.xml},xmlid={publisher-works-001}]

        \begin{teiElement}[name={front}]

          \begin{teiDiv}[type={titlePage},xmlid={titlepage-023}]

            \teiP[rend={title-main},xmlid={p-637}]{Les colloques de Cerisy}
          
\end{teiDiv}
        
\end{teiElement}
        \begin{teiElement}[name={body},xmlid={body-023}]

          \begin{teiDiv}[xmlid={div-001}]

            \begin{teiFigure}[xmlid={figure01-10}]

              \teiHead[xmlid={figures-32-001}]{Figures 32.}
            
\end{teiFigure}
            \teiP[xmlid={p1-23}]{Accueillis au \teiHi[rend={bold}]{château de Cerisy-la-Salle} et ses dépendances, monument historique du \teiHi[rend={small-caps}]{xvii}\teiHi[rend={sup}]{e} siècle au cœur du département de la Manche, le \teiHi[rend={bold}]{Centre culturel international de Cerisy} assure la programmation, l’organisation et la publication des \teiHi[rend={bold}]{colloques de Cerisy}. Il est le principal moyen d’action de l’\teiHi[rend={bold}]{Association des amis de Pontigny-Cerisy} (AAPC), \teiHi[rend={bold}]{reconnue d’utilité publique}, dont la mission est de favoriser les \teiHi[rend={bold}]{valeurs intellectuelles et artistiques} en développant les \teiHi[rend={bold}]{échanges culturels et scientifiques internationaux}.}
            \begin{teiDiv}[type={section1},xmlid={div01-13}]

              \teiHead[xmlid={une-aventure-culturelle-et-familiale-001}]{Une aventure culturelle et familiale}
              \teiP[xmlid={p2-22}]{Prolongeant les célèbres \teiHi[rend={bold}]{décades de Pontigny} (1910-1939) initiées par Paul Desjardins en Bourgogne, les \teiHi[rend={bold}]{colloques de Cerisy}, installés en 1952 par Anne Heurgon-Desjardins en Normandie, sont aujourd’hui dirigés par Edith Heurgon, avec le concours de la famille Peyrou-Bas, réunie au sein de la Société civile du château de Cerisy, propriétaire des lieux qu’elle met gracieusement à la disposition de l’Association.}
            
\end{teiDiv}
            \begin{teiDiv}[type={section1},xmlid={div02-13}]

              \teiHead[xmlid={une-experience-de-vie-et-de-pensee-001}]{Une expérience de vie et de pensée}
              \teiP[xmlid={p3-21}]{De Pontigny à Cerisy se poursuit un même projet : offrir la possibilité, dans un cadre prestigieux, \teiHi[rend={bold}]{de vivre et de penser avec ensemble}, dont le caractère unique tient à la \teiHi[rend={bold}]{durée des rencontres}, au « \teiHi[rend={bold}]{génie du lieu} », à \teiHi[rend={bold}]{l’hospitalité} de la famille et de l’équipe du Centre culturel.}
              \teiP[xmlid={p4-20}]{En toute \teiHi[rend={bold}]{indépendance d’esprit} et avec une volonté \teiHi[rend={bold}]{d’ouverture} et \teiHi[rend={bold}]{de brassage} des disciplines, des générations, des nationalités, les \teiHi[rend={bold}]{colloques de Cerisy} accueillent artistes, chercheurs, écrivains, enseignants, étudiants, responsables socioéconomiques et politiques, ainsi que tout public intéressé par les sujets traités. Les \teiHi[rend={bold}]{débats} tiennent un rôle clef pour confronter les points de vue et forger des \teiHi[rend={bold}]{idées neuves}.}
            
\end{teiDiv}
            \begin{teiDiv}[type={section1},xmlid={div03-12}]

              \teiHead[xmlid={une-action-durable-et-renouvelee-001}]{Une action durable et renouvelée}
              \teiP[xmlid={p5-20}]{Depuis 1952, plus de \teiHi[rend={bold}]{950 colloques} ont abordé des domaines très divers (art, littérature, philosophie, psychanalyse, sciences, prospective…). La Normandie y tient une place de choix avec près de 100 rencontres, dont une série prestigieuse sur \teiHi[rend={italic}]{La Normandie médiévale}.}
              \teiP[xmlid={p6-19}]{Près de \teiHi[rend={bold}]{690 ouvrages}, publiés chez des éditeurs variés, sont accessibles grâce, notamment, à la collection \teiHi[rend={bold}]{« Cerisy/Archives »} chez Hermann, qui réédite les colloques épuisés les plus fameux.}
            
\end{teiDiv}
            \begin{teiDiv}[type={section1},xmlid={div04-8}]

              \teiHead[xmlid={un-projet-federateur-et-societal-001}]{Un projet fédérateur et sociétal}
              \teiP[xmlid={p7-19}]{L’\teiHi[rend={bold}]{Association des amis de Pontigny-Cerisy} est ouverte à toute personne intéressée par sa mission et rassemble aujourd’hui environ 1 000 membres. Elle est présidée depuis 2023 par Jean-Louis Bancel, administrée par un conseil de vingt personnes et soutenue par un comité d’honneur rassemblant d’éminentes personnalités intellectuelles.}
              \teiP[xmlid={p8-19}]{La \teiHi[rend={bold}]{Commission de coordination régionale} regroupe, avec l’université de Caen, la DRAC, les collectivités territoriales et les villes partenaires, ainsi que divers acteurs culturels et scientifiques normands. Elle a pour objectif de construire des projets en Normandie et des partenariats locaux.}
              \teiP[xmlid={p9-18}]{Le \teiHi[rend={bold}]{Cercle des partenaires}, créé en 2005, réunit des entreprises, des collectivités territoriales ainsi que des organismes publics et des associations. Il apporte un soutien financier à l’AAPC et prend l’initiative de colloques sur des questions de société et de prospective.}
              \teiP[xmlid={p10-18}]{
                \teiHi[rend={bold}]{Renseignements sur les colloques et publications de Cerisy}
              }
              \teiP[xmlid={p11-18}]{cerisy-colloques.fr – (+33) 2 33 46 91 66}
              \teiP[xmlid={p12-18}]{CCIC, 2, le Château, 50210 Cerisy-la-Salle, France}
              \begin{teiFigure}[xmlid={figure02-8}]

                \teiHead[xmlid={figure-33-001}]{Figure 33.}
              
\end{teiFigure}
            
\end{teiDiv}
            \begin{teiDiv}[type={section1},xmlid={div05-7}]

              \teiHead[xmlid={choix-de-publications-001}]{Choix de publications}
              \teiP[xmlid={p13-17}]{\teiHi[rend={bold}]{Repenser l’Agir moderne (autour d’Armand Hatchuel)}, Eska, 2025}
              \teiP[xmlid={p14-17}]{\teiHi[rend={bold}]{Les Animaux, deux ou trois choses que nous savons d’eux}, Hermann, 2014}
              \teiP[xmlid={p15-17}]{\teiHi[rend={bold}]{L’usage des Ambiances : une épreuve sensible}, Hermann, 2021}
              \teiP[xmlid={p16-16}]{\teiHi[rend={bold}]{De quoi l’Art brut est-il le nom ?}, L’atelier contemporain, 2025}
              \teiP[xmlid={p17-16}]{\teiHi[rend={bold}]{Silvia Baron-Supervielle : le pays de l’écriture}, Kimé, 2025}
              \teiP[xmlid={p18-15}]{\teiHi[rend={bold}]{Brassages planétaires. Jardiner avec Gilles Clément}, Hermann, 2020}
              \teiP[xmlid={p19-15}]{\teiHi[rend={bold}]{Cultures et créations dans les métropoles-monde}, Hermann, 2016}
              \teiP[xmlid={p20-14}]{\teiHi[rend={bold}]{Agencer les multiplicités avec Deleuze}, Hermann, 2019}
              \teiP[xmlid={p21-12}]{\teiHi[rend={bold}]{Du Développement durable aux transitions}, Traversée, Hermann, 2022}
              \teiP[xmlid={p22-12}]{\teiHi[rend={bold}]{L’Écoute des mondes}, Traversée, Herman, 2025}
              \teiP[xmlid={p23-12}]{\teiHi[rend={bold}]{Écrire avec les vivants}, Traversée, Herman, 2022}
              \teiP[xmlid={p24-11}]{\teiHi[rend={bold}]{Écrire pour inventer : à partir des travaux de Jean Ricardou}, Hermann, 2020}
              \teiP[xmlid={p25-10}]{\teiHi[rend={bold}]{L’Enchantement qui revient}, Hermann, 2023}
              \teiP[xmlid={p26-10}]{\teiHi[rend={bold}]{Faire avec le sauvage, renouer avec les vivants}, HDiffusion, 2025}
              \teiP[xmlid={p27-10}]{\teiHi[rend={bold}]{Gestes spéculatifs}, Les Presses du réel, 2015}
              \teiP[xmlid={p28-9}]{\teiHi[rend={bold}]{Humains, animaux, nature : quelle éthique des vertus ?}, Hermann, 2020}
              \teiP[xmlid={p29-9}]{\teiHi[rend={bold}]{Renouveau des Jardins : clés pour un monde durable ?}, Hermann, 2014}
              \teiP[xmlid={p30-9}]{\teiHi[rend={bold}]{Jardins en société}, Traversée, Hermann, 2022}
              \teiP[xmlid={p31-8}]{\teiHi[rend={bold}]{Jardins en politique avec Gilles Clément}, Hermann, 2018}
              \teiP[xmlid={p32-7}]{\teiHi[rend={bold}]{Alexander Kluge, cartographie d’une œuvre plurielle}, Hermann, 2022}
              \teiP[xmlid={p33-6}]{\teiHi[rend={bold}]{Les Lieux qui nous affectent}, Hermann, 2021}
              \teiP[xmlid={p34-6}]{\teiHi[rend={bold}]{André Pieyre de Mandiargues, écrire entre les arts}, PUR, 2025}
              \teiP[xmlid={p35-6}]{\teiHi[rend={bold}]{La Mésologie, un paradigme pour l’anthropocène (A. Berque)}, Hermann, 2018}
              \teiP[xmlid={p36-6}]{\teiHi[rend={bold}]{La Mode comme indiscipline}, B42, 2024}
              \teiP[xmlid={p37-5}]{\teiHi[rend={bold}]{Edgar Morin, les cent premières années}, Hermann, 2024}
              \teiP[xmlid={p38-3}]{\teiHi[rend={bold}]{Futurs de l’Océan, des mers et des littoraux}, Hermann, 2025}
              \teiP[xmlid={p39-3}]{\teiHi[rend={bold}]{Musiques sacrées en Normandie}, Archives de la Manche, 2024}
              \teiP[xmlid={p40-3}]{\teiHi[rend={bold}]{De Pontigny à Cerisy : des lieux pour « penser avec ensemble »}, Hermann, 2011}
              \teiP[xmlid={p41-3}]{\teiHi[rend={bold}]{Valère Novarina, les tourbillons de l’écriture}, Hermann, 2020}
              \teiP[xmlid={p42-3}]{\teiHi[rend={bold}]{Prendre soin : savoirs, pratiques, nouvelles perspectives}, Hermann, 2013}
              \teiP[xmlid={p43-3}]{\teiHi[rend={bold}]{La Région, de l’identité à la citoyenneté}, Hermann, 2016}
              \teiP[xmlid={p44-3}]{\teiHi[rend={bold}]{Sciences de la vie, sciences de l’information}, ISTE, 2017}
              \teiP[xmlid={p45-3}]{\teiHi[rend={bold}]{Victor Segalen, « Attentif à ce qui n’a pas été dit… »}, Hermann, 2018}
              \teiP[xmlid={p46-3}]{\teiHi[rend={bold}]{La Sérendipité. Le hasard heureux}, Hermann, 2011}
              \teiP[xmlid={p47-3}]{\teiHi[rend={bold}]{SIÈCLE : 100 ans de rencontres intellectuelles de Pontigny à Cerisy}, IMEC, 2005}
              \teiP[xmlid={p48-3}]{\teiHi[rend={bold}]{Gilbert Simondon et l’invention du futur}, Klincksieck, 2016}
              \teiP[xmlid={p49-3}]{\teiHi[rend={bold}]{La Traduction dans une société interculturelle}, Hermann, 2022}
              \teiP[xmlid={p50-3}]{\teiHi[rend={bold}]{L’âge de la Transition}, Les petits matins, 2016}
              \teiP[xmlid={p51-2}]{\teiHi[rend={bold}]{Francisco Varela, une pensée actuelle}, Hermann, 2024}
              \teiP[xmlid={p52-2}]{\teiHi[rend={bold}]{Le moment du Vivant}, PUF, 2016}
            
\end{teiDiv}
          
\end{teiDiv}
        
\end{teiElement}
      
\end{teiElement}
    
\end{teiElement}
    \begin{teiElement}[name={back}]

\end{teiElement}
  
\end{teiElement}

\end{teiElement}
\end{lateiDocument}

\cleardoublepage
\pagestyle{plain}
\titleformat{\chapter}[display]{\PURHTitleFont\bfseries\fontsize{16pt}{19pt}\selectfont\centering}{}{10pt}{\MakeUppercase}
\tableofcontents

\end{document}
